% ================================
% \section*{\centering \bf Part III: Spatial Intelligence: From Perception to Generation}
% \addcontentsline{toc}{section}{Part III: Spatial Intelligence: From Perception to Generation}
% ================================





\begin{figure}
    \centering
    % \includegraphics[width=\linewidth, height=12cm]{example-image-golden}
    % \includegraphics[width=0.99\linewidth]{figures/part3.png}
    \includegraphics[width=\linewidth]{figures/part3_opener_v2.pdf}
    
    \caption{Part III in one view. (1) Perceive: a given scene is separated into objects, depth, relational structure, and motion. (2) Reason: the understood scene supports inference and planning, here a path that carries the cup around an obstacle to the shelf. (3) Generate: an intent unfolds into images, 3D scenes, video, motion, and whole worlds, so that generation runs the loop in reverse.}
    \label{fig:p3_main}
\end{figure}

If the previous part asked how space is represented, this one asks what a system does with that representation. Under the definition of Chapter~\ref{chapter:definition}, this use resolves into three capabilities that build on one another: perceiving space, reasoning over it, and generating it, the generative side of prediction. These are less separate subfields than three phases of a single loop running from sensory input to spatial artifact.

Perception comes first because everything downstream depends on it: it converts raw observation into structured spatial evidence, answering what is present, where it lies, and how it moves. Reasoning takes that evidence as given and draws consequences from it, inferring relations, planning routes and actions, and predicting how a configuration will change under intervention. Generation closes the loop by running the process in reverse, mapping an intent (a prompt, a constraint, a partial observation) back into concrete spatial content: an image, a shape, a scene, a motion, an entire simulated world. Read in sequence, the three chapters trace the path from a scene that is given, to one that is understood, to one that is made.

Each of the three rests on the representations built in the previous part and supplies capabilities that the applications of the next will draw on; together they are the core of what makes a system spatially intelligent.






\chapter{Spatial Perception}
\label{chapter:3:sec:6}
\label{sec:perception}
% \label{sec:6}


% \lettrine[lines=3]{\initfamily\textcolor{darkgreen}{S}}{patial perception} 
\FramedInitial{S}{patial} perception is the foundational layer of spatial intelligence: it turns raw sensory input into the structured spatial evidence that everything downstream depends on, answering what is present, where it lies, and how it moves. It is what stands between an agent and a scene, and the reliability of every later inference is bounded by the reliability of what perception delivers.

The chapter follows a single progression, from the elemental to the structured. It begins with recognizing the discrete entities that populate a scene, grounds them in metric geometry, then binds them into relational and scene-level structure, extends that structure across time, and finally broadens it beyond vision by fusing complementary modalities. Each step presupposes the ones before it, so that perception builds from isolated percepts toward a dynamic, multimodal model of space that reasoning and generation can act on.



\section{Spatial Entity Recognition}
\label{sec:entity_recognition}

\FramedInitial{B}{efore} space can be measured or organized, its contents must be named. Spatial entity recognition is that entry point: it answers \emph{what is present}, and every later geometric and relational step builds on the entities it names. Recognition operates at several granularities, and the section is ordered accordingly: detecting discrete object instances, delineating dense regions of pixels or points by segmentation, pinning down sparse keypoints and landmarks, and finally assigning a holistic scene label. What unites them is that they recover identity and category, not yet geometry; where each entity sits and how it is oriented is taken up next.


\subsection{Object Detection}
\label{sec:6-1-2}
% logic: (Xiangtai Li)
% 1, definition
% 2, method before
% 3, method after deep learning
%  % single stage detectors
%  % two-stage detectors
%  % transformer-based detectors. (Detr)
%  % open-vocabulary / open-world detectors.
%  % video object detection methods.
% 4, dataset and benchmark info
% 5, trends and future work.

\subsubsection{Task Definition}
Given a 2D image $I$, \emph{2D object detection} identifies instances of predefined object classes and localizes each one in the image plane. The output is a set of detections
\begin{equation}
\mathcal{D}=\{(c_i, b_i, s_i)\}_{i=1}^{N},
\end{equation}
where $c_i$ is the predicted class, $b_i=(x_i,y_i,w_i,h_i)$ a 2D bounding box parameterized by center and size, and $s_i\in[0,1]$ a confidence score. In short, the task answers \emph{what objects are where}. Figure~\ref{fig:p3_s6_1_2_objdet_overview} summarizes the 2D task and its detector families.

\emph{3D object detection} keeps the same set-of-detections output but places each box in metric 3D space. The input is a LiDAR point cloud $\mathcal{P}=\{\mathbf{p}_j\}_{j=1}^{M}$ with $\mathbf{p}_j\in\mathbb{R}^{3}$ (often with an intensity channel), one or more calibrated images, or both, and each box becomes
\begin{equation}
b_i=(x_i,y_i,z_i,\;l_i,w_i,h_i,\;\theta_i),
\end{equation}
a 3D center, the physical length, width, and height, and a heading angle $\theta_i$ about the vertical axis; indoor benchmarks often use axis-aligned boxes and drop $\theta_i$. Because the box is expressed in the sensor or ego frame rather than the image plane, a 3D detection already carries metric position and heading, which is why driving and robotics stacks consume it directly.
% \end{definition}

\begin{figure}[!t]
    \centering
    \includegraphics[width=\linewidth]{figures/part3/section6/object_detection.pdf}
    \caption{2D object detection: each detected instance is a tuple $(c_i, b_i, s_i)$ of class, bounding box, and confidence score (top), and the progression of detector families with representative methods under each (bottom).}
    \label{fig:p3_s6_1_2_objdet_overview}
\end{figure}

\subsubsection{Methodology}

\paragraph{Before Deep Learning.}
Early detectors~\cite{zou2019object} ran a sliding window over multi-scale image pyramids and classified each candidate region with hand-crafted features and shallow models. Viola--Jones~\cite{viola_jones_2001} combined Haar-like features with an AdaBoost cascade for real-time face detection; HOG features with linear SVMs~\cite{dalal_triggs_2005} became the standard pipeline for pedestrians and generic objects; and deformable part-based models (DPMs)~\cite{felzenszwalb_dpm_2010} represented an object as parts connected by deformable constraints. Manual feature design, dense-window evaluation, and limited robustness to pose, scale, and occlusion capped this line and motivated the move to learned features.

\paragraph{Two-stage Detectors.}
R-CNN~\cite{girshick2014rcnn} first generates category-independent region proposals and then classifies each with a CNN, a large accuracy gain over hand-crafted pipelines bought with high computational cost. Fast R-CNN~\cite{girshick2015fast} shares convolutional computation across proposals through ROI pooling and trains classification and box regression jointly; Faster R-CNN~\cite{ren2015faster} adds a Region Proposal Network so that proposals are nearly free and the whole detector trains end-to-end. Sparse R-CNN~\cite{sun2021sparsercnn} later replaced dense anchors with a small set of learnable proposals, reaching end-to-end detection without a transformer.

\paragraph{One-stage Detectors.}
One-stage detectors drop the proposal step and predict classes and boxes directly from the feature map. YOLO~\cite{redmon2016yolo} casts detection as a single regression over the whole image, trading localization accuracy for speed; YOLOv2 and YOLOv3~\cite{redmon2017yolo9000, redmon2018yolov3} add anchors, multi-scale prediction, and stronger backbones, and later versions~\cite{bochkovskiy2020yolov4, jocher2020yolov5} refine training and deployment. SSD~\cite{liu2016ssd} predicts from multiple feature maps at different scales, and RetinaNet~\cite{lin2017focal} introduces the focal loss to counter the extreme foreground--background imbalance, bringing one-stage accuracy level with two-stage methods. Recent YOLO models move to anchor-free, decoupled heads (YOLOv8~\cite{jocher2023yolov8}) and NMS-free end-to-end training (YOLOv10~\cite{wang2024yolov10}), converging with the set-prediction detectors below.

\paragraph{Transformer-based Detectors.}
DETR~\cite{carion2020end} reformulates detection as direct set prediction with a transformer encoder--decoder, removing anchors and non-maximum suppression at the price of slow convergence and weak small-object performance. Deformable DETR~\cite{Zhu2021DeformableDETR} attends to a sparse set of sampling points, which fixes both; Conditional DETR~\cite{meng2021conditional} conditions decoder attention on spatial queries for faster training; DINO~\cite{zhang2023dino} adds denoising training and better query initialization; and RT-DETR~\cite{zhao2024detrs} brings the family to real-time speed.

\paragraph{Open-Vocabulary Detection.}
Open-vocabulary detection (OVD) recognizes objects outside a fixed label set by aligning region features with text embeddings. The line evolved from embedding distillation to large-scale vision--language pretraining to data-driven vocabulary expansion: ViLD~\cite{gu2021open} and OVR-CNN~\cite{zareian2021ovrcnn} distill or learn a shared region--text embedding space from image--text pairs; GLIP~\cite{li2022grounded} unifies detection with phrase grounding under grounded language--image pretraining; and Detic~\cite{zhou2022detecting} shows that image-level labels over a large vocabulary suffice to train the classifier of an open-vocabulary detector. OWL-ViT~\cite{minderer2022simple}, Grounding DINO~\cite{liu2024grounding}, and YOLO-World~\cite{cheng2024yoloworld} push the same idea toward simpler transfer, tighter grounding, and real-time speed; see the surveys~\cite{zhu2024survey, wu2024towards} for a fuller account.

\paragraph{3D Object Detection.}
3D detectors differ mainly in how they turn sensor data into something a network can process; the detection heads on top largely reuse the 2D families above. On LiDAR, VoxelNet~\cite{zhou2018voxelnet} learns features inside voxels and runs 3D convolutions over the grid, and PointPillars~\cite{lang2019pointpillars} collapses the vertical axis into pillars so that an ordinary 2D CNN on the bird's-eye-view (BEV) plane does the rest at real-time speed. Point-based detectors keep the raw points: PointRCNN~\cite{shi2019pointrcnn} generates proposals from PointNet++ features~\cite{qi2017pointnetplus} of segmented foreground points, PV-RCNN~\cite{shi2020pvrcnn} combines voxel and point features. CenterPoint~\cite{yin2021center} replaces anchors with BEV center heatmaps, the 3D counterpart of anchor-free 2D detection. Indoor scenes are cluttered and object centers often fall in empty space, which favors voting: in VoteNet~\cite{qi2019votenet} surface points vote for object centers before grouping, and 3DETR~\cite{misra20213detr} carries DETR-style set prediction over to point clouds.

Camera-only 3D detection has to infer the depth that LiDAR measures. Monocular detectors regress depth, size, and orientation on top of a 2D detection (Section~\ref{sec:height_distance} treats them as metric estimation). Pseudo-LiDAR~\cite{wang2019pseudolidar} instead converts a predicted depth map into a point cloud and runs a LiDAR detector on it, showing that much of the camera--LiDAR gap came from the representation rather than from depth quality alone. Multi-camera systems lift image features into a shared BEV space, either bottom-up, as Lift-Splat-Shoot~\cite{philion2020lss} does by predicting a depth distribution per pixel and splatting features onto the ground plane, or top-down, as DETR3D~\cite{wang2022detr3d} and BEVFormer~\cite{LiWLXSLQD22} do by projecting 3D queries or BEV cells into the images to sample features. Fusion combines both sensors, from Frustum PointNets~\cite{qi2018frustum}, which crops the point cloud with 2D image detections, to BEVFusion~\cite{liu2023bevfusion}, which merges camera and LiDAR features on a common BEV map. Cube R-CNN~\cite{Omni3D-Brazil0SR0G23} trains a single image-based detector across indoor and outdoor domains on Omni3D. Deployment constraints for driving are discussed in Section~\ref{sec:p4_s9_2}, and video detection, which adds cross-frame association to the same output set, under tracking in Section~\ref{sec:object_tracking}.

\begin{table*}[!t]
\centering
\caption{Summary of object detection datasets across modalities and domains, with image and bounding-box counts per official split; ImageNet VID counts refer to video clips.}
\label{tab:object_detection_dataset}
\fontsize{9}{11}\selectfont
\setlength{\tabcolsep}{1.5mm}
\begin{tabular}{l cc cc cc cc }
\toprule
\multirow{2}{*}{\textbf{Dataset}}   & \multirow{2}{*}{\textbf{Modality, Domain}}   & \multirow{2}{*}{\textbf{Class}}   & \multicolumn{2}{c}{\textbf{train}} & \multicolumn{2}{c}{\textbf{validation}} & \multicolumn{2}{c}{\textbf{test}}  \\
\cmidrule(r){4-5}\cmidrule(r){6-7}\cmidrule(r){8-9}
  & & & \textbf{images} & \textbf{boxes} & \textbf{images} & \textbf{boxes} & \textbf{images} & \textbf{boxes} \\
\midrule
VOC-2007~\cite{pascal-voc-2007} & image, general & 20 & 2,501 & 6,301 & 2,510 & 6,307 & 4,952 & 14,976\\
VOC-2012~\cite{pascal-voc-2012} & image, general & 20 & 5,717 & 13,609 & 5,823 & 13,841 & 10,991 & - \\
ILSVRC-2014~\cite{ImageNet} & image, general & 200 & 456,567 & 478,807 & 20,121 & 55,502 & 40,152 & -\\
MS-COCO-2015~\cite{lin2014microsoft} & image, general & 80 & 82,783 & 604,907 & 40,504 & 291,875 & 81,434 & - \\
MS-COCO-2017~\cite{lin2014microsoft} & image, general & 80 & 118,287 & 860,001 & 5,000 & 36,781 & 40,670 & - \\
Objects365-2019~\cite{shao2019objects365} & image, general & 365 & 600,000 & 9,623,000 & 39,000 & 479,000 & 100,000 & 170,000\\
OpenImages~\cite{kuznetsova2020open} & image, general & 600 & 1,743,042 & 14,610,229 & 41,620 & 303,980 & 125,436 & 937,327 \\
ImageNet VID~\cite{ImageNet} & video, general & 30 & 3,862 & - & 555 & - & 937 & -\\
LVIS~\cite{gupta2019lvis} & image, general & 1,203 & 100,000 & - & - & - & - & - \\
Omni3D~\cite{Omni3D-Brazil0SR0G23} & 3D (image), general  & 98 & 175,573 & - & 19,127 & - & 39,452 & -  \\
\bottomrule
\end{tabular}
\end{table*}

\begin{table*}[!t]
\centering
\caption{Reported COCO performance of representative detector configurations, with evaluation splits, input sizes, and training schedules differing across sources; AP$_S$ / AP$_M$ / AP$_L$ break down by object size (small / medium / large). SSD512* denotes the image-expansion-augmented variant.}
\label{tab:objdet_coco_results}
\fontsize{9}{11}\selectfont
\setlength{\tabcolsep}{1.5mm}
\begin{tabular}{l l c c c c c c c c c}
\toprule
\textbf{Method} & \textbf{Backbone} & \textbf{\#P (M)} & \textbf{FLOPs (G)} & \textbf{AP} & \textbf{AP$_{50}$} & \textbf{AP$_{75}$} & \textbf{AP$_S$} & \textbf{AP$_M$} & \textbf{AP$_L$} & \textbf{Year} \\
\midrule
\multicolumn{11}{l}{\emph{Two-stage detectors}} \\
Faster R-CNN~\cite{ren2015faster}            & R50-FPN     & 41.5 & 207   & 37.4 & 58.1 & 40.4 & 21.2 & 41.0 & 48.1 & 2015 \\
Sparse R-CNN~\cite{sun2021sparsercnn}        & R50-FPN     & 106  & 174   & 45.0 & 64.1 & 49.4 & 26.9 & 47.2 & 59.5 & 2021 \\
\midrule
\multicolumn{11}{l}{\emph{One-stage CNN detectors}} \\
SSD512*~\cite{liu2016ssd}                     & VGG-16      & 27   & 88    & 28.8 & 48.5 & 30.3 & 10.9 & 31.8 & 43.5 & 2016 \\
RetinaNet~\cite{lin2017focal}                & R50-FPN     & 38   & 198   & 36.5 & 55.0 & 39.1 & 19.9 & 40.3 & 48.7 & 2017 \\
YOLOv3~\cite{redmon2018yolov3}               & DarkNet-53  & 62   & 156   & 33.0 & 57.9 & 34.4 & 18.3 & 35.4 & 41.9 & 2018 \\
\midrule
\multicolumn{11}{l}{\emph{Modern YOLO}} \\
YOLOv5-s~\cite{jocher2020yolov5}             & CSPDarkNet  & 7.2  & 16.5  & 37.4 & --   & --   & --   & --   & --   & 2020 \\
YOLOv8-x~\cite{jocher2023yolov8}             & C2f-based    & 68.2 & 257.8 & 53.9 & --   & --   & --   & --   & --   & 2023 \\
YOLOv10-l~\cite{wang2024yolov10}             & C2f/CIB    & 24.4 & 120.3 & 53.2 & --   & --   & --   & --   & --   & 2024 \\
\midrule
\multicolumn{11}{l}{\emph{Transformer set-prediction detectors}} \\
DETR~\cite{carion2020end}                    & R50, 500\,ep & 41   & 86    & 42.0 & 62.4 & 44.2 & 20.5 & 45.8 & 61.1 & 2020 \\
Deformable DETR~\cite{Zhu2021DeformableDETR} & R50, 50\,ep  & 40   & 173   & 43.8 & 62.6 & 47.7 & 26.4 & 47.1 & 58.0 & 2021 \\
DINO~\cite{zhang2023dino}                    & R50, 4-sc., 12\,ep & 47 & 279 & 49.0 & 66.6 & 53.5 & 32.0 & 52.3 & 63.0 & 2023 \\
RT-DETR~\cite{zhao2024detrs}                 & R50          & 42   & 136   & 53.1 & 71.3 & 57.7 & 34.8 & 58.0 & 70.0 & 2024 \\
\bottomrule
\end{tabular}
\end{table*}




\subsubsection{Datasets and Metrics}
Table~\ref{tab:object_detection_dataset} summarizes the benchmarks that shaped the field. \textbf{PASCAL VOC}~\cite{everingham2010pascal} fixed the evaluation protocol and remained the reference through the early deep-learning era despite its modest scale. \textbf{MS COCO}~\cite{lin2014microsoft}, with complex everyday scenes and segmentation masks alongside boxes, then became the standard for both detection and segmentation. \textbf{Open Images}~\cite{kuznetsova2020open}, \textbf{Objects365}~\cite{shao2019objects365}, and \textbf{LVIS}~\cite{gupta2019lvis} scale the vocabulary to hundreds and then over a thousand categories, with LVIS deliberately emphasizing rare classes, which turned it into the canonical long-tail and open-world benchmark.

Evaluation follows the \emph{precision--recall} framework. A detection $(\hat{c}, \hat{b}, \hat{s})$ is correct when it matches a previously unmatched ground-truth instance of class $\hat{c}$, whose box $b$ overlaps $\hat{b}$ at or above the required \emph{Intersection over Union} (IoU) threshold
\begin{equation}
\mathrm{IoU}(\hat{b}, b) = \frac{|\hat{b} \cap b|}{|\hat{b} \cup b|}.
\end{equation}
Sorting the detections of a class by confidence traces a precision--recall curve whose area is the \emph{Average Precision} (AP); the \emph{mean Average Precision} (mAP) averages AP over classes. Benchmarks differ chiefly in the IoU threshold: PASCAL VOC uses a single $\mathrm{IoU} = 0.5$ (AP$_{50}$), whereas COCO averages over ten thresholds from $0.5$ to $0.95$ to obtain its primary AP and reports AP$_{50}$ and AP$_{75}$ as diagnostic slices. COCO further stratifies by object size into AP$_S$ (area $<32^2$ pixels), AP$_M$ ($32^2$--$96^2$), and AP$_L$ ($>96^2$), which exposes the persistent difficulty of small objects.

3D detection is benchmarked mainly on driving data, KITTI~\cite{KITTI-GeigerLU12}, nuScenes~\cite{caesar2020nuscenes}, and the Waymo Open Dataset~\cite{WaymoOpen-SunKDCPTGZCCVHN20}, and indoors on SUN RGB-D~\cite{song2015sunrgbd} and ScanNet~\cite{dai2017scannet}; Omni3D (Table~\ref{tab:object_detection_dataset}) pools image-based 3D annotations from both. The IoU above is computed between 3D boxes or their BEV footprints: KITTI reports 3D and BEV AP at IoU $0.7$ for cars and $0.5$ for pedestrians and cyclists, and indoor benchmarks report mAP at 3D IoU $0.25$ and $0.5$. nuScenes matches detections by center distance on the ground plane ($0.5$, $1$, $2$, and $4$\,m) instead of IoU and folds translation, scale, orientation, velocity, and attribute errors into its nuScenes Detection Score (NDS); Waymo weights AP by heading accuracy (APH).

Table~\ref{tab:objdet_coco_results} collects reported COCO results across the detector families.
On the canonical ResNet-50 backbone, AP has climbed from $37.4$ for Faster R-CNN to $53.1$ for RT-DETR over roughly a decade, with much of the recent gain attributable to transformer-based set prediction (DETR $\rightarrow$ Deformable DETR $\rightarrow$ DINO) and improved denoising-style query initialization.
The AP$_S$ column highlights the well-documented small-object weakness of vanilla DETR ($20.5$); deformable attention and denoising progressively close this gap (Deformable DETR $26.4$, DINO $32.0$, RT-DETR $34.8$).
On the efficiency frontier, the modern YOLO line (YOLOv8-x at $53.9$, YOLOv10-l at $53.2$) matches RT-DETR in accuracy at substantially higher throughput, so the traditional accuracy--speed divide between CNN one-stage detectors and transformer set prediction has largely closed.



% \noindent
\subsubsection{Challenges}
\textbf{The accuracy--efficiency trade-off remains unresolved.}
The transformer-based detectors that top COCO and open-vocabulary leaderboards are still expensive at inference: even with deployment stacks such as TensorRT, their latency and parameter counts rule out most mobile and embedded targets, while the lightweight real-time detectors that fit those budgets trail them in accuracy. The same tension recurs for every recognition task in this section.
\textbf{Open-world detection degrades on the hardest instances.}
Open-vocabulary detectors extend recognition beyond a fixed label set, yet still miss small objects and mislabel rare categories, precisely the long-tail regime that motivated LVIS.
\textbf{Detection inside multimodal large language models is still immature.} Recent systems fold detection into MLLMs as a grounding primitive that ties image regions to language, but even the strongest closed-source models post markedly lower numbers than specialist detectors on standard benchmarks~\cite{zhu2024survey, wu2024towards}.
\textbf{3D detection is bounded by depth.} For camera-only detectors, localization error is dominated by depth error, which grows with distance~\cite{ma2021delving}; LiDAR measures depth directly but its returns thin out with range, so distant and small objects are seen through a handful of points.





\subsection{Segmentation}
\label{sec:segmentation}
% logic: (xiangtai li)
% 1, definition of various segmentation (two or three paragraph)
% 2, methods before deep learning.
% 3, methods using deep learning (FCN-based/U-Net based)
% 4, transformer-based methods.
%    image/video/3D segmentation (transformers)
% 5, datasets.
% 6，trends and challenges.

Segmentation partitions an image, a video, or a 3D scan into meaningful regions. Before deep learning it was a mid-level grouping problem, clustering pixels by similarity somewhere between low-level edge detection and high-level recognition; learned models turned it into dense prediction, with class labels, instance identities, and temporal consistency attached to every region.
Figure~\ref{fig:p3_s6_1_3_segmentation_examples} shows the principal segmentation map types produced on a common scene.

\begin{figure}[!ht]
    \centering
    \includegraphics[width=0.99\linewidth]{figures/part3/section6/segmentation.pdf}
    \caption{Examples of different segmentation maps: semantic, instance, and panoptic on a shared scene.}
    \label{fig:p3_s6_1_3_segmentation_examples}
\end{figure}

\subsubsection{Task Definition}

\textit{Semantic Segmentation}:
Given an input image $I \in \mathbb{R}^{H \times W \times 3}$, \emph{semantic segmentation} assigns a class label to every pixel, producing a dense label map
\begin{equation}
    S \in \{1, \dots, C\}^{H \times W},
\end{equation}
where $C$ is the number of classes; instances of the same class are not distinguished.

\textit{Instance and Panoptic Segmentation}:
\emph{Instance segmentation} detects and segments each object instance separately. The output is a set of labeled masks
\begin{equation}
\mathcal{Y} = \{(m_i, c_i)\}_{i=1}^{N},
\end{equation}
where $m_i \in \{0,1\}^{H \times W}$ is the binary mask of the $i$-th instance and $c_i$ its class; masks may overlap, and only countable \emph{thing} classes are covered. \emph{Panoptic segmentation} keeps the same output form but requires the $N$ masks to form a non-overlapping partition of the image that covers \emph{thing} instances and \emph{stuff} regions alike, each stuff class as a single mask.

\textit{Video Segmentation}:
Given a video clip $V \in \mathbb{R}^{T \times H \times W \times 3}$, \emph{video segmentation} predicts temporally consistent masks across frames, output as a set of mask tubes
\begin{equation}
\mathcal{Y} = \{(m_i, c_i)\}_{i=1}^{N}, \quad m_i \in \{0,1\}^{T \times H \times W},
\end{equation}
each tracking an object or region over time, with semantic, instance, and panoptic variants as above.

\textit{3D Segmentation}:
Given a point cloud $\mathcal{P}=\{\mathbf{p}_j\}_{j=1}^{M}$ with $\mathbf{p}_j\in\mathbb{R}^{3}$ and optional color or intensity, \emph{3D semantic segmentation} labels every point,
\begin{equation}
S\in\{1,\dots,C\}^{M},
\end{equation}
and 3D instance and panoptic segmentation predict point masks $m_i\in\{0,1\}^{M}$ under the same set form and overlap constraints as their image counterparts; voxels or mesh vertices take the place of points for other 3D representations. The definitions carry over unchanged. What changes is the input: the points are unordered, unevenly spaced, and sparse, and that is where 3D methods depart from image methods.


\subsubsection{Methodology}
Figure~\ref{fig:p3_s6_1_3_segmentation_overview} summarizes the method families surveyed below.

\begin{figure}[!t]
    \centering
    \includegraphics[width=\linewidth]{figures/part3/section6/segmentation-method.pdf}
    \caption{Methodological progression of image and video segmentation, from pre-deep-learning methods to foundation and promptable segmenters, with representative methods under each stage.}
    \label{fig:p3_s6_1_3_segmentation_overview}
\end{figure}

% \noindent
\paragraph{Before Deep Learning.}
Classical segmentation grouped pixels either by appearance or by minimizing an energy. The first family includes \emph{thresholding} of pixel intensities~\cite{otsu1979threshold}, \emph{region growing} that merges similar neighbors~\cite{adams1994region}, \emph{$k$-means clustering} in feature space~\cite{macqueen1967kmeans}, and the \emph{watershed} transform, which treats the image as a topographic surface and segments by catchment basins~\cite{vincent1991watersheds}. The second family includes \emph{active contours (snakes)} that evolve a curve toward object boundaries~\cite{kass1988snakes}, \emph{graph cuts} that solve certain segmentation energies to global optimality~\cite{boykov2001graphcuts}, and \emph{Markov} and \emph{conditional random fields} that impose spatial coherence probabilistically~\cite{geman1984mrf,lafferty2001crf}; the last of these carried over directly into early learning-based segmenters.



\paragraph{Fully Convolutional Networks.}
FCNs~\cite{long2015fcn} replace fully connected layers with convolutions, enabling end-to-end dense prediction at arbitrary input size. U-Net~\cite{ronneberger2015unet} adds a symmetric encoder--decoder with skip connections that fuse low- and high-level features, SegNet~\cite{badrinarayanan2017segnet} reuses max-pooling indices for cheap upsampling and boundary recovery, and PSPNet~\cite{zhao2017pspnet} aggregates global and local context through pyramid pooling.

\paragraph{Dilated and Multi-Scale Context Modeling (DeepLab Series).}
The DeepLab series introduces atrous (dilated) convolution to enlarge the receptive field without losing resolution (v1~\cite{chen2015deeplab}), Atrous Spatial Pyramid Pooling for multi-scale context (v2~\cite{chen2017deeplabv2}, refined with batch normalization and stronger backbones in v3~\cite{chen2017deeplabv3}), and an encoder--decoder that recovers boundary detail from low-level features (v3+~\cite{chen2018deeplabv3plus}).

\paragraph{Region-Based Mask Prediction (Mask R-CNN Family).}
Mask R-CNN~\cite{he2017mask} extends Faster R-CNN with a parallel mask branch and ROIAlign, which removes the quantization error of ROI pooling and makes pixel-accurate masks possible. PANet~\cite{liu2018panet}, Hybrid Task Cascade~\cite{chen2019htc}, and PointRend~\cite{kirillov2020pointrend} refine it through better feature flow, cascaded box--mask prediction, and point-wise boundary rendering respectively; this two-stage, region-based family dominated instance segmentation until transformers arrived.



% \noindent
\paragraph{Transformer Backbones and Representation Learning.}
Transformer-based segmentation, following DETR, favors unified query-based formulations that bridge detection and segmentation, predict masks end-to-end without hand-crafted post-processing, and allow joint learning across tasks and modalities. The first step was the backbone: SETR~\cite{zheng2021setr} applies a pure ViT to semantic segmentation with a simple decoder, SegFormer~\cite{xie2021segformer} pairs a lightweight hierarchical transformer with an MLP decoder for a favorable accuracy--efficiency trade-off, and Swin Transformer~\cite{liu2021swin} uses shifted-window attention to produce CNN-like multi-scale features and has become a default segmentation backbone. Self-supervised pretraining such as MAE~\cite{he2022masked} supplies ViT backbones that transfer well to dense prediction.


\paragraph{Query-based Transformer Segmentation.}
Query-based methods cast segmentation as set prediction with object or mask queries. MaX-DeepLab~\cite{wang2021maxdeeplab} is the first purely mask-based panoptic framework without box supervision; MaskFormer~\cite{cheng2021maskformer} unifies semantic and panoptic segmentation by predicting mask embeddings with class labels; K-Net~\cite{zhang2021knet} treats segmentation as kernel prediction with iterative refinement; and Mask2Former~\cite{cheng2022mask2former} adds masked cross-attention and multi-scale decoding, setting a strong baseline across all three tasks. QueryInst~\cite{fang2021queryinst} and Panoptic SegFormer~\cite{li2022panopticsegformer} are variants that respectively generate instance-specific kernels from queries and decouple semantic from instance queries.


\paragraph{Video and Multimodal Transformer-based Segmentation.}
VisTR~\cite{wang2021vistr} predicts spatio-temporal mask tubes end-to-end from object queries; Video K-Net~\cite{li2022videoknet} unifies video semantic, instance, and panoptic segmentation by associating queries across time; and TubeFormer~\cite{wang2022tubeformer} adds tube-level attention for long-term consistency. X-Decoder~\cite{zou2023xdecoder} unifies image segmentation with language-based prediction, while OMGSeg~\cite{OMGSeg} covers image and video segmentation across open-vocabulary and class-agnostic settings within one model.


\paragraph{SAM series.} SAM~\cite{kirillov2023sam} made segmentation promptable: trained on the SA-1B dataset of over one billion masks, it decouples a heavy image encoder from a light prompt-conditioned mask decoder and produces masks from points, boxes, or text with strong zero-shot generalization. SAM~2~\cite{ravi2024sam2} extends this to video, propagating masks across frames with a memory mechanism while keeping interactive control.

\paragraph{3D Point Cloud Segmentation.}
Point-cloud segmenters split by how they impose structure on the points (Section~\ref{sec:rep_3d_point}). Point-based networks consume raw points: PointNet~\cite{qi2017pointnet} applies a shared MLP to every point and pools symmetrically, PointNet++~\cite{qi2017pointnetplus} adds hierarchical neighborhoods, DGCNN~\cite{wang2019dgcnn} builds $k$-nearest-neighbor graphs that are recomputed in feature space at every layer, KPConv~\cite{thomas2019kpconv} defines convolution kernels by points placed in space, and RandLA-Net~\cite{hu2020randlanet} relies on random sampling to scale to scenes of millions of points. Sparse voxel networks such as MinkowskiNet~\cite{choy2019minkowski} and SPVNAS~\cite{tang2020searching} quantize the points and convolve only occupied voxels, which keeps 3D convolution affordable on room- and street-scale scans. Projection methods take the opposite route: RangeNet++~\cite{milioto2019rangenet} unrolls a LiDAR sweep into a range image, runs a 2D segmenter in real time, and projects the labels back onto the points. Transformers entered with Point Transformer~\cite{zhao2021point}, and Point Transformer V3~\cite{wu2024ptv3} replaces exact neighbor search with serialized point orderings so that attention scales to large scenes.

Instance segmentation in 3D retraces the 2D history. PointGroup~\cite{jiang2020pointgroup} predicts per-point offsets toward instance centers and clusters the shifted points, a bottom-up grouping; Mask3D~\cite{schult2023mask3d} applies Mask2Former-style queries to sparse voxel features and predicts instance masks directly. Because 3D labels remain scarce, open-vocabulary 3D segmentation mostly borrows from 2D foundation models: OpenScene~\cite{peng2023openscene} distills CLIP-aligned pixel features into points, and OpenMask3D~\cite{takmaz2023openmask3d} classifies class-agnostic 3D instance masks with CLIP features aggregated over multi-view images.



\begin{table*}[!t]
\centering
\caption{Summary of representative segmentation datasets across 2D and 3D modalities; split sizes count images, scans, or video sequences as indicated. Parenthesized class counts denote common reduced evaluation label sets.}
\label{tab:segmentation_dataset}
\fontsize{9}{11}\selectfont
\setlength{\tabcolsep}{3.7mm}
\begin{tabular}{l c cc ccc }
\toprule
\multirow{2}{*}{\textbf{Dataset}}   & \multirow{2}{*}{\textbf{Purpose}}   & \multirow{2}{*}{\textbf{Classes}} & \multirow{2}{*}{\textbf{Resolution}} & \multicolumn{3}{c}{\textbf{Samples}}  \\
\cmidrule(r){5-7}
% \cmidrule(r){5-6}\cmidrule(r){7-8}
  & & & & \textbf{Train} & \textbf{Valid} & \textbf{Test} \\
\midrule
\rowcolor{lightgrey}\multicolumn{7}{l}{$\bullet$ \textit{2D}}\\
Pascal VOC \cite{everingham2010pascal} & Generic & 21 & Variable & 1464 & 1449 & Private\\
MS COCO \cite{lin2014microsoft} & Generic & 80+ & Variable & 82,783 & 40,504 & 81,434\\
ADE20K \cite{zhou2017ade20k} & Generic & 150 & Variable & 20,210 & 2,000 & 3,000 \\
COCO-stuff \cite{COCO-Stuff-CaesarUF18} & Generic & 172 & Variable & 118,000 & 5,000 & 20,000 \\
IDD \cite{IDD-VarmaSNCJ19} & Urban/Driving & 34 & Variable & 6,993 & 981 & 2,029\\
Dark Zurich \cite{SakaridisDG19} & Urban/Driving & 19 & 1920$\times$1080& 8,377  & 50 & 151\\
ACDC \cite{ACDC-SakaridisDG21} & Driving & 19 & 2048$\times$ 1024& 1,600  & 406 & 2,000\\
PartImageNet \cite{PartImageNet-HeYYKYCLYYY22} & Generic-Part & 158 & Variable & 16,540 & 2,957 & 4,598\\
SiftFlow \cite{SiftFlow-LiuYT09} & Outdoor & 33 & 256$\times$256 & 2,488 & - & 200\\
BSDS300 \cite{BSD-MartinFTM01}& Boundary/Region & -- & Variable & 200 & - & 100 \\
Youtube-Objects \cite{PrestLCSF12} & Generic/Video & 10 & 480$\times$360 & 126 & - & - \\
Cityscapes~\cite{cordts2016cityscapes}  & Urban & 30 (19) & 2048 $\times$1024 & 2975 & 500 & 1525\\
CamVid \cite{CamVid-BrostowFC09} & Urban/Driving & 32 (11) & 960$\times$720 & 367 & 101 & 233 \\
Mapillary Vistas~\cite{neuhold2017mapillary} & Street & 65 & - & 18,000 & - & - \\
\hline
\rowcolor{lightgrey}\multicolumn{7}{l}{$\bullet$ \textit{3D}}\\
NYU-D V2~\cite{silberman2012nyudv2} & Indoor & 40 (13) & 640$\times$480 & 795 & - & 654 \\
SUN-3D~\cite{xiao2013sun3d} & SfM/Video & - & 640$\times$480 & 415 seqs & - & - \\
SUN RGB-D~\cite{song2015sunrgbd} & Indoor & 37 & Variable & 5,285 & - & 5,050 \\
ScanNet~\cite{dai2017scannet} & Indoor & 20 & Variable & 1,201 & 312 & 100 \\
SemanticKITTI~\cite{SemanticKITTI-BehleyGMQBSG19} & LiDAR/Driving & 28 & Point cloud & 19,130 & 4,071 & 20,351 \\
\bottomrule
\end{tabular}
\end{table*}

% \noindent
\subsubsection{Datasets and Metrics}
Table~\ref{tab:segmentation_dataset} summarizes the main benchmarks. \textbf{PASCAL VOC}~\cite{everingham2010pascal} anchored the classical semantic-segmentation era with 20 foreground classes plus background. \textbf{MS COCO}~\cite{lin2014microsoft}, with semantic, instance, box, and panoptic annotations over complex object-centric scenes, is the shared testbed for all three modern task types. \textbf{Cityscapes}~\cite{cordts2016cityscapes} and \textbf{Mapillary Vistas}~\cite{neuhold2017mapillary} specialize to street scenes, the latter stressing long-tail classes and diverse geographies, and \textbf{ADE20K}~\cite{zhou2017ade20k} targets holistic scene parsing over 150 object and stuff classes. The 3D rows of the table cover indoor RGB-D scans (NYU-D V2, SUN RGB-D, ScanNet) and outdoor LiDAR (SemanticKITTI), and S3DIS~\cite{armeni2016s3dis} is a further indoor standard.

Segmentation evaluation reuses IoU (Section~\ref{sec:6-1-2}) with masks in place of boxes. For semantic segmentation the headline metric is the \emph{mean Intersection over Union} (mIoU),
\begin{equation}
\mathrm{mIoU} = \frac{1}{C} \sum_{k=1}^{C} \frac{\mathrm{TP}_k}{\mathrm{TP}_k + \mathrm{FP}_k + \mathrm{FN}_k},
\end{equation}
where $C$ is the number of classes and $\mathrm{TP}_k$, $\mathrm{FP}_k$, $\mathrm{FN}_k$ count the true-positive, false-positive, and false-negative pixels for class $k$.
Because the average is unweighted, mIoU is sensitive to rare classes, which is one motivation behind ADE20K and Mapillary Vistas; pixel accuracy is dominated by frequent classes (\eg, sky, road) and is less informative.

For \emph{instance segmentation}, the COCO mask-AP protocol matches predictions to ground truth by mask IoU and averages AP over IoU thresholds $0.5$--$0.95$.
For \emph{panoptic segmentation}, \emph{Panoptic Quality}~\cite{kirillov2019panoptic} is $\mathrm{PQ} = \mathrm{SQ} \times \mathrm{RQ}$, the average IoU of matched segments times the F1 score over matched, false-positive, and false-negative segments, treating things and stuff uniformly.
For \emph{video object segmentation}, the DAVIS $\mathcal{J}\&\mathcal{F}$ metric averages region similarity ($\mathcal{J}$, mask IoU) and contour accuracy ($\mathcal{F}$, computed on mask boundaries).
3D semantic segmentation reports the same mIoU computed over points or voxels, and 3D instance segmentation on ScanNet reports AP at mask IoU thresholds averaged over $0.5$--$0.95$, with AP$_{50}$ and AP$_{25}$ as looser slices.

\begin{table*}[!t]
\centering
\caption{Reported performance of representative segmentation methods on standard benchmarks: ADE20K mIoU and Cityscapes mIoU (semantic), COCO Mask AP (instance), and COCO PQ (panoptic). Numbers follow each paper's own protocol (split, single-/multi-scale) and are indicative only; ``--'' marks a metric not reported.}
\label{tab:segmentation_results}
\fontsize{9}{11}\selectfont
\setlength{\tabcolsep}{1.4mm}
\begin{tabular}{l l c c c c c}
\toprule
\textbf{Method} & \textbf{Backbone} & \textbf{ADE20K mIoU} & \textbf{Cityscapes mIoU} & \textbf{COCO Mask AP} & \textbf{COCO PQ} & \textbf{Year} \\
\midrule
\rowcolor{lightgrey}\multicolumn{7}{l}{$\bullet$\emph{FCN family and dilated convolution}} \\
FCN-8s~\cite{long2015fcn}                  & VGG-16        & 29.4 & 62.7 & --   & --   & 2015 \\
PSPNet~\cite{zhao2017pspnet}               & R101          & 43.3 & 80.2 & --   & --   & 2017 \\
DeepLab v3+~\cite{chen2018deeplabv3plus}   & Xception-65   & 45.6 & 82.1 & --   & --   & 2018 \\
\midrule
\rowcolor{lightgrey}\multicolumn{7}{l}{$\bullet$ \emph{Region-based instance segmentation}} \\
Mask R-CNN~\cite{he2017mask}               & R50-FPN       & --   & --   & 35.4 & --   & 2017 \\
HTC~\cite{chen2019htc}                     & R50-FPN       & --   & --   & 38.4 & --   & 2019 \\
\midrule
\rowcolor{lightgrey}\multicolumn{7}{l}{$\bullet$ \emph{Transformer backbones}} \\
SETR-MLA~\cite{zheng2021setr}              & ViT-L         & 50.3 & 82.2 & --   & --   & 2021 \\
SegFormer-B5~\cite{xie2021segformer}       & MiT-B5        & 51.0 & 84.0 & --   & --   & 2021 \\
Swin-L + UperNet~\cite{liu2021swin}        & Swin-L        & 53.5 & --   & --   & --   & 2021 \\
\midrule
\rowcolor{lightgrey}\multicolumn{7}{l}{$\bullet$ \emph{Universal query-based segmentation}} \\
MaskFormer~\cite{cheng2021maskformer}      & Swin-L        & 55.6 & --   & --   & 52.7 & 2021 \\
Mask2Former~\cite{cheng2022mask2former}    & Swin-L        & 57.7 & --   & 50.1 & 57.8 & 2022 \\
\bottomrule
\end{tabular}
\end{table*}

Table~\ref{tab:segmentation_results} reports results across these four protocols. The clearest trend is on semantic segmentation, where ADE20K mIoU has roughly doubled in under a decade, from FCN-8s's $29.4$ to Mask2Former's $57.7$; the largest single jumps track two shifts, first to transformer backbones (SETR, SegFormer, Swin-L) and then to universal query-based formulations. That second shift is what unifies the table: MaskFormer is the first entry to report semantic mIoU and panoptic PQ within one framework, and Mask2Former extends the same architecture to all three tasks with competitive results on each. Instance segmentation followed a partly separate path, climbing from Mask R-CNN ($35.4$) and HTC ($38.4$) on R50-FPN to $50.1$ once Mask2Former is paired with a Swin-L backbone, which confirms that transformer backbones lift mask quality at the instance level as well, not only semantic accuracy.


% \noindent
\subsubsection{Challenges}
The cost of self-attention at high resolution, over long video, and on dense 3D data sharpens the accuracy--speed tension noted for detection, and two problems are specific to segmentation.
\textbf{Long-video, multi-object consistency is hard.} Holding spatial precision and temporal coherence together remains difficult in cluttered scenes, particularly for video panoptic segmentation and segmentation-based multi-object tracking, where objects enter, leave, and occlude one another over long horizons.
\textbf{Segmentation is moving toward the open world.} The field is shifting from fixed label sets to class-agnostic, open-vocabulary, and multimodal settings, in which language and other auxiliary signals are folded in to improve spatial reasoning and adaptability, a direction that SAM, X-Decoder, and OMGSeg already anticipate.









\subsection{Keypoint and Landmark Detection}
\label{sec:keypoint_landmark}
Keypoint and landmark detection localizes a set of semantically consistent, precisely positioned points that characterize object geometry in images or videos: eye corners, skeletal joints, or structurally distinctive parts of man-made objects.
Compared with region-based representations, keypoints give a sparser, explicitly geometric encoding that often survives deformation, rotation, and partial occlusion.
Figure~\ref{fig:p3_s6_1_1_keypoint_overview} summarizes the task families and the method families surveyed below.



\subsubsection{Task Definition}
Given an input image $I$, a parametric model $f_{\theta}$ estimates a set of point locations
\begin{equation}
f_{\theta}(I) = \{\mathbf{p}_i\}_{i=1}^{N}, \quad \mathbf{p}_i \in \mathbb{R}^2,
\end{equation}
where each $\mathbf{p}_i$ is the image coordinate of a salient visual or structural element. Landmark detection fixes a predefined semantic identity for each point (a particular eye corner, a particular joint); the two terms are otherwise used interchangeably.

\emph{3D keypoint detection} predicts $\mathbf{p}_i\in\mathbb{R}^{3}$ instead. From images, the depth of each point must be inferred; from depth maps or point clouds $\mathcal{P}$, surface geometry is observed, but occluded or internal landmarks still require inference. For human pose, 3D joints are usually expressed in camera coordinates relative to a root joint (the pelvis), which leaves global depth unconstrained, as a single image cannot fix it.

The task varies along three axes. The input modality sets what geometry is directly available: most detectors work on 2D images (RGB, grayscale, or infrared), where keypoints must be inferred from appearance alone, whereas point clouds expose spatial structure directly; RGB-D or RGB-infrared fusion adds robustness in clutter and poor lighting, and video supplies temporal continuity through motion blur and transient occlusion.
The prediction target ranges from planar image coordinates, which underpin most benchmarks, to 3D keypoints required in robotics and autonomous driving, to 4D trajectories tracked over time.
The application scenario fixes what the points mean. Research has concentrated on humans: body joints for pose estimation, facial landmarks for expression recognition and identity verification, and finger joints for gesture-based interaction and sign language recognition~\cite{ko2019neural}. Keypoints are also defined on everyday items, industrial parts, and vehicles for manipulation, inspection, and driving (\eg, monocular 3D detection from vehicle keypoints~\cite{li2020rtm3d}), on animals for behavior analysis (multi-animal DeepLabCut~\cite{lauer2022multi}), and on anatomical structures in medical imaging, where universal landmark detectors cover several modalities and anatomies~\cite{zhu2021you}.


\begin{figure}[!t]
    \centering
    \includegraphics[width=\linewidth]{figures/part3/section6/keypoint_detection_v2.pdf}
    \caption{Keypoint and landmark detection: representative task families (top) and the methodological progression, mainly of 2D detectors, from hand-crafted feature detectors to generative models, with representative methods under each stage (bottom).}
    \label{fig:p3_s6_1_1_keypoint_overview}
\end{figure}


\subsubsection{Methodology}
Classical pipelines combined local feature detectors with statistical shape models that handled occlusion and deformation through explicit geometric constraints.
Moravec~\cite{moravec1980obstacle}, Harris~\cite{harris1988combined}, and Shi--Tomasi~\cite{shi1994good} corners were followed by feature detectors and descriptors such as SIFT~\cite{lowe2004distinctive}, SURF~\cite{bay2006surf}, FAST~\cite{rosten2006machine}, BRIEF~\cite{calonder2010brief}, and ORB~\cite{rublee2011orb}, which improved scale invariance, speed, and robustness.
For faces and hands, active shape models~\cite{cootes1995asm}, active appearance models~\cite{cootes1998active}, and constrained local models with regularized mean-shift fitting~\cite{saragih2011deformable} provided the parametric frameworks.


Learning-based detection began with supervised descent for face alignment~\cite{xiong2013supervised} and, in 2014, with DeepPose~\cite{toshev2014deeppose}, which regressed human joints directly with a deep network.
Heatmap regression then dominated through 2020: Tompson \etal~\cite{tompson2014joint,tompson2015efficient} coupled convolutional networks with graphical models, convolutional pose machines~\cite{wei2016convolutional} and stacked hourglass networks~\cite{newell2016stacked} refined heatmaps over repeated stages, and HRNet~\cite{wang2020deep} kept a high-resolution branch throughout the network for precise localization.


Multi-person estimation splits into two strategies. Top-down methods detect each person first and then estimate keypoints inside the box, which is accurate but scales with the number of people; bottom-up methods predict all keypoints at once and then group them, which is faster but makes grouping the hard part.
Part Affinity Fields~\cite{cao2017realtime}, consolidated as OpenPose~\cite{CaoHSWS21}, are the canonical bottom-up design; YOLO-Pose~\cite{maji2022yolo}, RTMPose~\cite{jiang2023rtmpose}, and RTMO~\cite{lu2024rtmo} advance real-time multi-person estimation through one-stage and top-down designs.


Vision transformers and foundation models made keypoint detection general-purpose: ViTPose~\cite{xu2022vitpose} set simple transformer baselines for human pose in 2022, and Sapiens~\cite{khirodkar2024sapiens} in 2024 provides a single foundation model for human-centric vision tasks.

3D human pose estimation grew out of these 2D detectors. Lifting methods take detected 2D keypoints as input and regress the 3D joints: a fully connected network on 2D joints alone was a strong baseline~\cite{martinez2017simple}, temporal convolutions over 2D pose sequences (VideoPose3D~\cite{pavllo2019videopose3d}) and spatio-temporal transformers (MotionBERT~\cite{zhu2023motionbert}) use motion to reduce the depth ambiguity, and DiffPose~\cite{gong2023diffpose} casts lifting as diffusion. Direct methods predict 3D from pixels, as in the volumetric heatmaps of Pavlakos \etal~\cite{pavlakos2017coarse}. With a calibrated multi-view rig most of the ambiguity disappears, and learnable triangulation~\cite{iskakov2019learnable} fuses per-view heatmaps into a 3D estimate. Recovering a full body mesh rather than a joint set is taken up with human motion in Section~\ref{sec:p3_s8_4}.


Classical corner detectors were in turn replaced by learned ones: SuperPoint~\cite{detone2018superpoint} learns detection and description jointly under self-supervised homographic adaptation, DISK~\cite{tyszkiewicz2020disk} selects keypoints with a reinforcement-style objective, and SuperGlue~\cite{sarlin2020superglue} matches sparse keypoints; the detector-free LoFTR~\cite{Sun2021LoFTR} instead produces semi-dense correspondences.
These serve as correspondence rather than semantic points and are the workhorses of the structure-from-motion, visual-odometry, and SLAM pipelines in the spatial-localization section that follows.

On point clouds, keypoints mostly serve registration and pose estimation. Intrinsic Shape Signatures~\cite{zhong2009iss} select points whose local neighborhood has a distinctive eigenvalue structure; USIP~\cite{li2019usip} learns, without labels, detectors that stay stable under rigid transformation, and D3Feat~\cite{bai2020d3feat} predicts a detection score and a descriptor for every point, the 3D analogue of SuperPoint. Category-level 3D keypoints are a different kind: Suwajanakorn \etal~\cite{suwajanakorn2018discovery} discover keypoints that are consistent across instances of a category through a relative-pose objective, and kPAM~\cite{manuelli2019kpam} uses such keypoints as the interface for robotic manipulation. Keypoint-based 6D object pose, which predicts the 2D projections of 3D model keypoints and solves PnP, belongs to pose estimation (Section~\ref{sec:object_pose}).

Table~\ref{tab:keypoint_coco_results} reports COCO Keypoints (val2017) results for the CNN, transformer, and real-time families.

\begin{table*}[!t]
\centering
\caption{Reported performance of representative keypoint detection methods on COCO Keypoints (val2017). AP is computed under the OKS-based protocol.}
\label{tab:keypoint_coco_results}
\fontsize{9}{11}\selectfont
\setlength{\tabcolsep}{2.0mm}
\begin{tabular}{l l c c c c c c c c}
\toprule
\textbf{Method} & \textbf{Backbone} & \textbf{Input} & \textbf{\#Params (M)} & \textbf{GFLOPs} & \textbf{AP} & \textbf{AP$_{50}$} & \textbf{AP$_{75}$} & \textbf{AR} & \textbf{Year} \\
\midrule
HRNet-W32~\cite{wang2020deep}  & HRNet-W32   & 256$\times$192 & 28.5  & 7.1   & 74.4 & 90.5 & 81.9 & 79.8 & 2020 \\
HRNet-W48~\cite{wang2020deep}  & HRNet-W48   & 384$\times$288 & 63.6  & 32.9  & 76.3 & 90.8 & 82.9 & 81.2 & 2020 \\
YOLOv5s6-Pose~\cite{maji2022yolo} & YOLOv5-s   & 960$\times$960 & 15.1  & --  & 63.8 & 87.6 & 69.6 & 70.4 & 2022 \\
YOLOv5m6-Pose~\cite{maji2022yolo} & YOLOv5-m & 960$\times$960 & 41.4 & -- & 67.4 & 89.1 & 73.7 & 73.9 & 2022 \\
YOLOv5l6-Pose~\cite{maji2022yolo} & YOLOv5-l & 960$\times$960 & 87.0 & -- & 69.4 & 90.2 & 76.1 & 75.9 & 2022 \\
ViTPose-B~\cite{xu2022vitpose} & ViT-B       & 256$\times$192 & 86    & --  & 75.8 & -- & -- & 81.1 & 2022 \\
ViTPose-L~\cite{xu2022vitpose} & ViT-L       & 256$\times$192 & 307   & --  & 78.3 & -- & -- & 83.5 & 2022 \\
ViTPose-H~\cite{xu2022vitpose} & ViT-H       & 256$\times$192 & 632   & -- & 79.1 & -- & -- & 84.1 & 2022 \\
RTMPose-s~\cite{jiang2023rtmpose} & CSPNeXt-s & 256$\times$192 & 5.5  & 0.68  & 69.6 & -- & -- & -- & 2023 \\
RTMPose-m~\cite{jiang2023rtmpose} & CSPNeXt-m & 256$\times$192 & 13.6  & 1.93  & 73.6 & -- & -- & -- & 2023 \\
RTMPose-l~\cite{jiang2023rtmpose} & CSPNeXt-l & 256$\times$192 & 27.7  & 4.16  & 74.8 & -- & -- & -- & 2023 \\
RTMO-s~\cite{lu2024rtmo} & CSPDarkNet & 640$\times$640 & 9.9 & -- & 67.7 & 87.8 & 73.7 & 71.5 & 2024 \\
RTMO-m~\cite{lu2024rtmo} & CSPDarkNet & 640$\times$640 & 22.6 & -- & 70.9 & 89.0 & 77.8 & 74.7 & 2024 \\
RTMO-l~\cite{lu2024rtmo} & CSPDarkNet & 640$\times$640 & 44.8 & -- & 72.4 & 89.9 & 78.8 & 76.2 & 2024 \\
\bottomrule
\end{tabular}
\end{table*}

\subsubsection{Datasets and Metrics}
Keypoint datasets are mostly embedded in broader pose and action benchmarks rather than built in isolation; Table~\ref{tab:keypoint_dataset_summary} lists the most used.
Human body datasets are the most mature and trace the field's expansion: MPII~\cite{andriluka20142d} and COCO Keypoints~\cite{lin2014microsoft} fixed in-the-wild 2D evaluation, PoseTrack~\cite{andriluka2018posetrack} added temporal association in video, Human3.6M~\cite{IonescuPOS14} supplies controlled multi-view 3D ground truth for 2D-to-3D lifting, DensePose~\cite{guler2018densepose} replaces sparse joints with dense surface correspondences, and CrowdPose~\cite{li2019crowdpose} stresses crowded, occluded scenes.
Face datasets moved from large pose variation (AFLW~\cite{koestinger2011annotated}) and high-density annotation (HELEN~\cite{le2012interactive}) to occlusion (COFW~\cite{burgos2013robust}), the standardized 68-point protocol of 300-W~\cite{sagonas2013300} and its large-pose synthetic extension 300-W-LP~\cite{zhu2016face}, and the 98-point, hard-condition WFLW~\cite{wu2018look}.
Hand datasets contend with fine articulation and self-occlusion through multi-view capture (CMU Panoptic HandDB~\cite{simon2017hand}), synthetic 3D labels (RHD~\cite{zimmermann2017learning}), real 3D pose and shape (FreiHAND~\cite{zimmermann2019freihand}), and interacting hands (InterHand2.6M~\cite{moon2020interhand2}).
Animal-Pose~\cite{cao2019cross} and AP-10K~\cite{yu2021ap} carry the human recipe to animals across species, while ApolloCar3D~\cite{song2019apollocar3d} and KeypointNet~\cite{you2020keypointnet} define keypoints on rigid vehicles and 3D object categories.

\begin{table*}[!t]
\centering
\caption{Representative datasets for keypoint and landmark detection across human body, face, hand, animal, and object domains. ``KP'' denotes the number of annotated keypoints per instance; ``Mod.'' denotes modality (2D image, 3D pose, or multi-view). Sample counts follow the official splits where available; ``--'' marks values not provided by the source paper.}
\label{tab:keypoint_dataset_summary}
\fontsize{9}{11}\selectfont
\setlength{\tabcolsep}{1.0mm}
\begin{tabular}{L{0.22\textwidth} C{0.06\textwidth} L{0.23\textwidth} C{0.11\textwidth} C{0.08\textwidth} L{0.23\textwidth}}
\toprule
\textbf{Dataset} & \textbf{Year} & \textbf{Domain} & \textbf{Mod.} & \textbf{KP} & \textbf{\#Samples (train\,/\,val\,/\,test)} \\
\midrule
\rowcolor{lightgrey} \multicolumn{6}{l}{$\bullet$ \emph{Human body}} \\
MPII~\cite{andriluka20142d}            & 2014 & In-the-wild body & 2D & 16 & 25K images (train+test) \\
COCO Keypoints~\cite{lin2014microsoft} & 2014 & In-the-wild multi-person & 2D & 17 & 57K / 5K / 20K images \\
PoseTrack2018~\cite{andriluka2018posetrack} & 2018 & Video body tracking & 2D & 15 & 593 / 170 / 375 videos \\
Human3.6M~\cite{IonescuPOS14}   & 2014 & Indoor multi-view body & 3D & 32 & 3.6M frames, 11 subjects \\
DensePose~\cite{guler2018densepose}    & 2018 & Dense body surface & 2D & SMPL surface & $\sim$50K annotated persons (COCO subset) \\
CrowdPose~\cite{li2019crowdpose}       & 2019 & Crowded multi-person & 2D & 14 & 10K / 2K / 8K images \\
\midrule
\rowcolor{lightgrey} \multicolumn{6}{l}{$\bullet$ \emph{Face}} \\
AFLW~\cite{koestinger2011annotated}    & 2011 & In-the-wild face & 2D & 21 & 25K faces (22K images) \\
HELEN~\cite{le2012interactive}         & 2012 & High-resolution face & 2D & 194 & 2K / -- / 330 images \\
COFW~\cite{burgos2013robust}           & 2013 & Occluded face & 2D & 29 & 1.3K / -- / 507 images \\
300-W~\cite{sagonas2013300}            & 2013 & Aggregated face & 2D & 68 & 3.1K / -- / 689 images \\
300-W-LP~\cite{zhu2016face}            & 2016 & Large-pose synthetic face & 2D & 68 & 61K images \\
WFLW~\cite{wu2018look}                 & 2018 & In-the-wild face (challenging) & 2D & 98 & 7.5K / -- / 2.5K images \\
\midrule
\rowcolor{lightgrey} \multicolumn{6}{l}{$\bullet$ \emph{Hand}} \\
CMU Panoptic HandDB~\cite{simon2017hand} & 2017 & Multi-view hand & Multi-view & 21 & 14K manual + bootstrapped \\
RHD~\cite{zimmermann2017learning}       & 2017 & Synthetic hand & 2D + 3D & 21 & 41K / -- / 2.7K frames \\
FreiHAND~\cite{zimmermann2019freihand}  & 2019 & Real hand pose + shape & 3D & 21 & 33K / -- / 4K samples \\
InterHand2.6M~\cite{moon2020interhand2} & 2020 & Two-hand interaction & 3D & 42 & 2.6M frames \\
\midrule
\rowcolor{lightgrey} \multicolumn{6}{l}{$\bullet$ \emph{Animal \& Object}} \\
Animal-Pose~\cite{cao2019cross}         & 2019 & 5-class mammal pose & 2D & 20 & $\sim$4K images \\
AP-10K~\cite{yu2021ap}                  & 2021 & 54-species animal pose & 2D & 17 & 10K images \\
ApolloCar3D~\cite{song2019apollocar3d} & 2019 & Vehicle keypoints (driving) & 2D & 66 & 5.3K images \\
KeypointNet~\cite{you2020keypointnet}   & 2020 & Multi-category 3D objects & 3D & 5--23 & 8.3K models \\
\bottomrule
\end{tabular}
\end{table*}

Metrics measure raw or scale-normalized distances between predicted and ground-truth points. The End-Point Error (EPE) is the raw Euclidean distance. The Normalized Mean Error (NME) divides it by a reference length:
\begin{equation}
\mathrm{NME} = \frac{1}{N} \sum_{i=1}^{N} \frac{\lVert \hat{\mathbf{x}}_i - \mathbf{x}_i \rVert_2}{d},
\end{equation}
where $\hat{\mathbf{x}}_i$ and $\mathbf{x}_i$ denote the predicted and ground-truth locations of the $i$-th keypoint, and $d$ is a task-dependent normalization factor (\eg, inter-ocular distance for faces).
The Percentage of Correct Keypoints (PCK) instead thresholds the error,
\begin{equation}
\mathrm{PCK@}\alpha = \frac{1}{N} \sum_{i=1}^{N}
\mathbb{I}\!\left( \lVert \hat{\mathbf{x}}_i - \mathbf{x}_i \rVert_2 \le \alpha \cdot s \right),
\end{equation}
where $\alpha$ controls the tolerance and $s$ is the reference scale, giving an accuracy rate that is easy to compare across methods; the Cumulative Error Distribution (CED) curve and its Area Under the Curve (AUC) summarize behavior across thresholds.

For large-scale benchmarks, Object Keypoint Similarity (OKS) incorporates object scale and keypoint-specific uncertainty,
\begin{equation}
\mathrm{OKS} =
\frac{\sum_i \exp\!\left( -\frac{\lVert \hat{\mathbf{x}}_i - \mathbf{x}_i \rVert_2^2}{2 s^2 \kappa_i^2} \right)\delta(v_i > 0)}
{\sum_i \delta(v_i > 0)}.
\end{equation}
Here $s$ denotes the object scale (the square root of the segmented area), $\kappa_i$ is a per-keypoint normalization constant calibrating the fall-off, and $v_i \in \{0, 1, 2\}$ is the ground-truth visibility flag for keypoint $i$, where $v_i = 0$ indicates unlabeled, $v_i = 1$ indicates labeled but occluded, and $v_i = 2$ indicates labeled and visible.
OKS plays the role of IoU in computing average precision and recall on MS COCO~\cite{lin2014microsoft}. For 3D human pose, the Mean Per-Joint Position Error (MPJPE) averages the Euclidean joint error in millimeters after aligning the root joints; PA-MPJPE applies a Procrustes alignment of scale, rotation, and translation first, and 3DPCK thresholds the error at $150$\,mm, as on MPI-INF-3DHP~\cite{mehta2017monocular}. Point-cloud keypoint detectors are scored by repeatability under rigid transformation and, downstream, by registration recall on benchmarks such as 3DMatch~\cite{zeng20173dmatch}.

\subsubsection{Challenges}
Robustness in unconstrained scenes is still the first problem~\cite{fisal2025reported, guo2025survey}: occlusion, extreme pose, and non-rigid deformation leave partial or ambiguous configurations, and symmetric parts (left versus right limbs, the fingers of a self-occluded hand) are routinely swapped.
Generalization is the second~\cite{kappan2026survey, ma20252d, jiang2022animal}: keypoint definitions differ across datasets even for the same object class, so models trained on one benchmark transfer poorly to another, and the gap widens for animals, objects, and anatomical structures, where annotation is scarce and demographic or species bias enters the training data.
The accuracy--speed tension discussed for detection applies here as well~\cite{khabarlak2021fast, alvarez2021real}: high-precision models are too heavy for embedded targets, and lightweight ones lose accuracy precisely under occlusion and fast motion.










\subsection{Scene Classification}
\label{sec:scene_classification}

% 最开始是简单的scene classification，即简单的分类任务，但是后续变得更加的fine-grained的任务，变成multi-label classification，再更加详细的话，就会变成scene （image/video/3D）captioning task.

Going beyond objects, regions, and keypoints, scene classification (SC) assigns a single semantic label to a whole image describing the environment it portrays~\cite{zeng2021deep}.
Where object recognition names the things in a picture (``car'', ``tree'', ``person''), scene classification names the \emph{where}: ``a busy street'', ``a serene beach'', ``a modern kitchen''.
It asks ``\emph{What kind of place is this?}'', a judgment that rests on the global layout and gist of the image rather than on any single object~\cite{henderson1999high,epstein2005cortical}, and it supplies the holistic prior on which an embodied or virtual agent conditions its later behavior~\cite{liao2016understand,ni2022improved,muhammad2022vision,Zhang2024UrbanVI}.


\subsubsection{Task Definition}

Formally, scene classification learns a mapping $f: I \mapsto y \in \mathcal{Y}$ from an image $I$ to a label $y$ in a predefined category set $\mathcal{Y}$, or, in the multi-label case, to a subset $Y \subseteq \mathcal{Y}$ covering several facets of the scene~\cite{boutell2004learning} (``urban'' for area, ``parking lot'' for place, ``night'' for condition).
Beyond the single-/multi-label split, the task is instantiated by domain, with indoor~\cite{quattoni2009recognizing} and outdoor~\cite{6724252} scenes, remote-sensing imagery~\cite{nogueira2017towards,cheng2017remote}, and driving contexts such as highways and tunnels~\cite{ni2022improved}, and by modality, extending from still images to video~\cite{huang1999integration} and acoustic scenes~\cite{mesaros2018multi}.


\subsubsection{Methodology}

\begin{figure}[!t]
    \centering
    \includegraphics[width=0.99\linewidth]{figures/part3/section6/scene_classification.pdf}
    \caption{Scene classification: four task formulations, from single-label classification to compositional captioning output (top), and the methodological progression from hand-crafted descriptors to multimodal LLMs, with representative methods under each stage (bottom).}
    \label{fig:p3_s6_1_4_scene_classification_overview}
\end{figure}

Figure~\ref{fig:p3_s6_1_4_scene_classification_overview} summarizes the progression from hand-crafted global descriptors through deep features to vision--language models.

Early methods summarized the whole image into one discriminative vector of low-level texture, color, and gradient statistics.
The Spatial Envelope model~\cite{oliva2001modeling} describes a scene by a few perceptual dimensions (naturalness, openness, roughness, expansion, ruggedness), a descriptor commonly called GIST.
Bag-of-Visual-Words (BoVW)~\cite{csurka2004visual} treats the image as a histogram of quantized local features such as SIFT~\cite{lowe2004distinctive}, discarding most spatial information~\cite{7025553}; Spatial Pyramid Matching (SPM)~\cite{lazebnik2006beyond} restores coarse layout by computing BoVW histograms over a pyramid of sub-regions and concatenating them.

Deep networks learn scene features end to end.
CNNs~\cite{krizhevsky2012imagenet,o2015introduction} such as AlexNet~\cite{krizhevsky2012imagenet}, VGG~\cite{simonyan2014very}, ResNet~\cite{he2016deep}, and DenseNet~\cite{huang2017densely}, pre-trained on ImageNet~\cite{ImageNet, krizhevsky2012imagenet}, proved to be strong global feature extractors for scenes, and Places-CNN~\cite{zhou2014learning} re-trained them on the large-scale Places benchmark to establish the standard transfer recipe.
Vision Transformers~\cite{vit} match or exceed CNNs with a class-token embedding that attends over all patches, in the manner of BERT~\cite{devlin2019bert}.


A separate line models layout explicitly, since a stadium is recognized by its global oval while a classroom differs from an office by the mid-scale arrangement of desks and blackboard.
Multi-scale pooling over the CNN feature hierarchy (FPN-style~\cite{lin2017feature}) captures structure at several scales, and scene graphs~\cite{johnson2015image} make objects and their relations explicit, so that propagating over them with graph convolutional networks~\cite{kipf2017semisupervised} yields layout-aware scene features~\cite{Liang2020ADN}.

Vision--language models (VLMs) turn scene classification into open-vocabulary, zero-shot matching.
Contrastive models such as CLIP~\cite{radford2021learning} embed the image and the textual names of all candidate classes in one space and pick the most similar class~\cite{li2023rs}, needing no fine-tuning beyond prompt design~\cite{sahoo2024systematic}; the recipe extends to 3D point clouds~\cite{huang2023clip2point}.
Autoregressive multimodal large language models (MLLMs; LLaVA~\cite{liu2023llava}, Qwen-VL~\cite{bai2023qwen}, Gemma~3~\cite{team2025gemma}) instead generate the label as text in response to an instruction~\cite{Saha_2024_CVPR}.
With rich semantics available for free, the research focus has shifted from one global label toward compositional understanding of what is inside the scene~\cite{chen2023clip2scene}.

\subsubsection{Datasets and Metrics}
Table~\ref{tab:sceneclf_datasets} lists the standard benchmarks. The small early sets, Scene 15~\cite{lazebnik2006beyond} and the indoor MIT 67~\cite{quattoni2009recognizing}, were superseded by SUN 397~\cite{XiaoHEOT10,Xiao2014SUNDE}, which fixed a near-exhaustive category list, and by Places~\cite{zhou2014learning, zhou2017places}, which scaled scene recognition to millions of images across its Places205, Places88, Places365-Standard, and Places365-Challenge subsets. NYU Depth v2~\cite{silberman2012nyudv2} and SUN RGB-D~\cite{song2015sunrgbd} add aligned depth and a scene category per image for indoor RGB-D classification.

\begin{table*}[!t]
\centering
\caption{Representative scene classification datasets. ``Mod.'' denotes modality; ``--'' marks a value not fixed by the source (scene-category subsets vary across protocols).}
\label{tab:sceneclf_datasets}
\fontsize{9}{11}\selectfont
\setlength{\tabcolsep}{2.2mm}
\begin{tabular}{l c l c c l}
\toprule
\textbf{Dataset} & \textbf{Year} & \textbf{Domain} & \textbf{Mod.} & \textbf{\#Classes} & \textbf{\#Images (protocol)} \\
\midrule
Scene 15~\cite{lazebnik2006beyond}                    & 2006 & Indoor + outdoor, grayscale & 2D & 15  & 4,448 \\
MIT Indoor 67~\cite{quattoni2009recognizing}               & 2009 & Indoor                      & 2D & 67  & 15,620 (80 / 20 per class) \\
SUN 397~\cite{XiaoHEOT10,Xiao2014SUNDE}       & 2010 & Indoor + outdoor            & 2D & 397 & 108,754 (50 test per class) \\
Places205~\cite{zhou2014learning}          & 2014 & Indoor + outdoor            & 2D & 205 & 2.5M \\
Places365-Standard~\cite{zhou2017places}   & 2017 & Indoor + outdoor            & 2D & 365 & $\sim$1.8M \\
Places365-Challenge~\cite{zhou2017places}  & 2017 & Indoor + outdoor            & 2D & 365 & $\sim$8M \\
NYU Depth v2~\cite{silberman2012nyudv2}    & 2012 & Indoor                      & RGB-D & -- & 1,449 labeled pairs \\
SUN RGB-D~\cite{song2015sunrgbd}           & 2015 & Indoor                      & RGB-D & -- & 10,335 \\
\bottomrule
\end{tabular}
\end{table*}

Scene classification is evaluated primarily with top-1 accuracy (the fraction of images for which the highest-scoring predicted class matches the ground truth) and top-5 accuracy (the fraction of images where the ground-truth class is among the top five predictions), with the latter especially useful on long-tail benchmarks such as Places-365 and SUN 397.
For multi-label settings, the mean Average Precision (mAP) and class-wise F1 are reported instead, mirroring the protocol used in multi-label image recognition.

\begin{table*}[!t]
\centering
\caption{Reported performance of representative scene classification methods on Places-365 val and SUN 397. ``--'' marks a metric not reported by the corresponding source.}
\label{tab:sceneclf_results}
\fontsize{9}{11}\selectfont
\setlength{\tabcolsep}{5.4mm}
\begin{tabular}{l l c c c}
\toprule
\textbf{Method} & \textbf{Backbone} & \textbf{Places-365 Top-1} & \textbf{SUN 397 Top-1} & \textbf{Year} \\
\midrule
\rowcolor{lightgrey} \multicolumn{5}{l}{$\bullet$ \emph{Supervised}} \\
Places365-AlexNet~\cite{zhou2017places}    & AlexNet   & 53.2 & --   & 2017 \\
Places365-VGG~\cite{zhou2017places}        & VGG-16    & 55.2 & --   & 2017 \\
Places365-ResNet~\cite{zhou2017places}     & R152      & 54.7 & --   & 2017 \\
\midrule
\rowcolor{lightgrey} \multicolumn{5}{l}{$\bullet$ \emph{Vision-language zero-shot}} \\
CLIP~\cite{radford2021learning}            & ViT-B/32  & --   & 63.2 & 2021 \\
\bottomrule
\end{tabular}
\end{table*}

Table~\ref{tab:sceneclf_results} reports results on the two most-cited scene benchmarks. Supervised CNNs on Places-365 cluster in a narrow band of roughly $53$--$55\%$ top-1, with diminishing returns from AlexNet to ResNet, while CLIP reaches $63.2\%$ top-1 on SUN 397 zero-shot~\cite{radford2021learning}, on par with supervised scene classifiers without any target-set training.




\subsubsection{Challenges}
Closed-set accuracy on Places-365, SUN 397, and MIT 67 is now recovered zero-shot by CLIP~\cite{radford2021learning} and by MLLMs (LLaVA~\cite{liu2023llava}, Qwen-VL~\cite{bai2023qwen}, Gemma~3~\cite{team2025gemma}), so ``what kind of place is this'' is no longer the frontier.
The task has climbed a granularity ladder, from one label to attribute sets to grounded, compositional description across image, video, and 3D, and the open problems sit at that richer end, where MLLMs are weaker than their coarse-label accuracy suggests.

\textbf{Fine-grained and long-tail discrimination.} MLLMs still trail dedicated classifiers and CLIP-style models on fine-grained and long-tail categories, where the distinction between, \eg, a ``bistro'' and a ``brasserie'' turns on subtle cues~\cite{NEURIPS2024_5c702404}.

\textbf{Compositional and spatial grounding.} The natural output is no longer a label but a description whose claims can be grounded in specific objects and layouts~\cite{chen2023clip2scene}.
MLLMs describe fluently but ground weakly: objects are referenced vaguely rather than by region or box, and the persistent failure is structural rather than semantic, since models name the constituents of a scene correctly yet err on relative layout, depth ordering, orientation, and cross-view or 3D consistency.
The cause is representational: 2D patch tokens are optimized for semantic alignment with language, not geometry, so fine spatial cues are lost once flattened into the language stream, and spatial-cognition benchmarks show even strong MLLMs degrading sharply as tasks move from naming toward layout and 3D reasoning~\cite{xu2025spatialbench,daxberger2025mmspatial}.

\textbf{Hallucination and faithfulness.} Free-form description reintroduces a failure mode that closed-set classification never had: MLLMs assert objects, attributes, or relations absent from the image, and their reasoning traces rarely point to verifiable spatial evidence~\cite{daxberger2025mmspatial}.






% \newpage

\section{Spatial Localization and Position Estimation}
\label{sec:p3-sec6-2-spatial-localization-and-position-estimation}

\FramedInitial{R}{ecognition} says what is present; it does not say where. Spatial localization and position estimation supply the missing geometry, situating recognized entities in a metric frame and recovering how they are oriented within it. The section climbs a ladder of geometric abstraction: from the per-pixel quantities of depth and surface normals, through metric attributes such as height and distance, to full SE(3) pose, first of the moving camera and then of rigid objects in the scene. Throughout, the concern is the static geometry of a single instant; how these quantities evolve is deferred to spatio-temporal perception.


\subsection{Depth Estimation}
\label{sec:depth_estimation}
Depth estimation predicts, for every pixel, the distance from the camera to the surface it observes, turning a 2D image into the 3D structure that reconstruction, navigation, and manipulation build on. Figure~\ref{fig:p3_s6_2_1_depth_overview} summarizes the input configurations and the method families surveyed below.

\subsubsection{Task Definition}
Given an image $I$ and optional auxiliary observations $\mathcal{V}$ (a second view, a multi-view set, a sparse sensor measurement, or neighboring video frames), a depth estimator predicts a depth map
\begin{equation}
f_\theta(I,\mathcal{V}) = D \in \mathbb{R}_{+}^{H\times W},
\end{equation}
where $D(u)$ is the distance along the optical axis from the camera center to the surface seen at pixel $u$. The task splits along three axes.

\textbf{Sensing.} Active methods emit light (structured light, time-of-flight, LiDAR) and measure its return, providing reliable metric depth even on textureless surfaces but tying accuracy to the sensor and degrading under strong ambient light or at long range. Passive methods infer depth from ordinary images alone; they need no special hardware and have become the dominant research line, at the cost of sensitivity to texture and illumination.

\textbf{Scale.} Metric depth gives distances in physical units, which 3D reconstruction and navigation require. Relative depth recovers only pixel ordering, or depth up to an affine transform; it cannot support metric measurement but is easier to supervise and generalizes better because heterogeneous datasets can be mixed without a consistent scale.

\textbf{Density and Input Configuration.} Classical pipelines produced sparse depth at feature points or semi-dense depth along high-gradient edges; dense per-pixel prediction is now the norm. The input configuration fixes how much geometry is available: active sensing measures depth directly; stereo recovers it from the disparity between two calibrated views; multi-view estimation enforces consistency across many views, which improves robustness to occlusion and matching ambiguity at the price of computation and known camera poses; monocular estimation lacks direct multi-view triangulation constraints and relies on learned statistical, structural, and semantic cues, which is why it is ill-posed in scale yet the most widely applicable setting.

% \paragraph{Extend Tasks}

%  \item Depth Completion (RGB + sparse depth)
%  \item Depth Upsampling
%  \item Depth Refinement
%  \item Depth Fusion

\subsubsection{Methodology}
% 深度估计方法整体脉络可整理为早期依据几何方法直接求解三维的深度数值，而后进入深度学习领域，从早期的有监督方法发展至无监督和自监督方法，其中又由传统的CNN网络转变为更强大的Transformer架构。目前最新的深度估计也在探索具有更好Zero-Shot表现、泛化性能更强的方法。

\begin{figure}[!t]
    \centering
    \includegraphics[width=\linewidth]{figures/part3/section6/depth_estimation_v3.pdf}
    \caption{Depth estimation: four input configurations (top) and five overlapping method families with representative methods on an approximate emergence timeline (bottom). Cost volumes and attention are distinct mechanisms; denoising applies to the generative branch. The depth maps are schematic illustrations, with brighter values indicating nearer surfaces.}
    \label{fig:p3_s6_2_1_depth_overview}
\end{figure}

\paragraph{Geometric Methods.}
Classical depth estimation solves geometry explicitly. In binocular stereo, two rectified cameras with focal length $f$ and baseline $B$ observe a point whose projections differ by the disparity $d = x_1 - x_2$; similar triangles give
\begin{equation}
Z = \frac{fB}{d},
\end{equation}
so depth follows once correspondences are found. Stereo matching is therefore a correspondence problem, usually posed as an energy that balances data fidelity against smoothness; Semi-Global Matching approximates the global optimum by aggregating costs along several 1D paths and remains a standard dense baseline. Structure from Motion (SfM) extends triangulation to many uncalibrated views, recovering camera poses and a sparse point cloud up to global scale, and supplies the poses that Multi-View Stereo (MVS) densifies. MVS typically proceeds by plane sweep: for a reference pixel $u_r$ and a hypothesized depth $d_k$, back-project with the intrinsics $K_r$, transform into a source view with the relative pose $(R, t)$, reproject, and keep the depth whose reprojections are most photometrically consistent across source views,
\begin{equation}
u_s \sim K_s\left(R\, d_k K_r^{-1} u_r + t\right), \qquad
D(u_r) = \arg\min_{d_k} \sum_s \left\| I_s(u_s) - I_r(u_r) \right\|.
\end{equation}
Plane sweep is thus triangulation turned into a search over depth planes, and it is the template that MVSNet~\cite{yao2018mvsnet} and CasMVSNet~\cite{gu2020cascade} later made differentiable as learned cost volumes.

Two single-image cues complete the classical picture. Depth from defocus exploits the thin-lens relation $1/f = 1/z + 1/z'$: a point off the focal plane blurs into a circle of confusion whose diameter grows with its distance from that plane, so measured blur can be inverted for depth. Shape from shading inverts a reflectance model, typically Lambertian, to recover surface orientation from intensity and integrates it into depth; it is treated under surface normals in \S\ref{sec:surface_normal}. Both are ill-posed under unknown illumination, reflectance, or optics and survive today mainly as auxiliary cues.

\paragraph{Active Depth Sensing.}
Active sensors measure depth rather than infer it. Time-of-flight systems time the round trip of emitted light: LiDAR scans the scene with laser beams and returns a sparse but long-range point cloud, whereas ToF cameras flood the scene and recover a dense depth map at every pixel, best suited to short-range indoor use. Structured light projects a coded pattern (the Kinect v1 used a pseudo-random infrared speckle) and triangulates from its deformation; active stereo projects texture onto the scene only to make ordinary stereo matching stable on featureless surfaces.

\paragraph{Early Learning-Based Methods.}
% 早期的机器学习方法
Data-driven depth began in the mid-2000s with monocular estimation. Saxena \etal~\cite{saxena2005learning} regressed depth from hand-crafted multi-scale features with a Markov random field enforcing spatial coherence, and Make3D~\cite{saxena2008make3d} extended the model to the 3D position and orientation of scene patches while releasing a laser-scanned outdoor dataset that served as an early benchmark.

\paragraph{Supervised CNN Regression.}
Eigen \etal~\cite{eigen2014depth} first applied CNNs to monocular depth with a coarse-to-fine two-scale network, and their evaluation protocol became the standard; a follow-up~\cite{eigen2015predicting} showed that the same architecture predicts depth, normals, and semantic labels with minor changes. Deep Convolutional Neural Fields~\cite{liu2015deep} paired CNN unaries with a continuous CRF, FCRN~\cite{laina2016deeper} introduced a fully convolutional encoder--decoder with efficient upsampling of ResNet features, and DORN~\cite{fu2018deep} recast depth as ordinal regression over a spacing-increasing discretization.

\paragraph{Learned Stereo and Multi-View Stereo.}
Learning entered the geometric pipelines first as a replacement for individual components: MC-CNN~\cite{zbontar2015computing} computed stereo matching costs with a Siamese network in place of hand-crafted similarity. Mayer \etal~\cite{mayer2016large} then showed that networks can regress disparity and scene flow end to end, and GC-Net~\cite{kendall2017end} and PSMNet~\cite{chang2018pyramid} embedded the cost volume and its regularization inside the network. The same move reached multi-view geometry with SurfaceNet~\cite{ji2017surfacenet} and MVSNet~\cite{yao2018mvsnet}, which infer depth from several posed images, and R-MVSNet~\cite{yao2019recurrent} and CasMVSNet~\cite{gu2020cascade} improved their memory footprint and resolution.

\paragraph{Self-Supervised and Multi-Task Learning.}
Ground-truth depth is costly, so a second line replaced it with geometric consistency. SfM-Net~\cite{vijayanarasimhan2017sfm} and SfM-Learner~\cite{zhou2017unsupervised} learn depth and camera motion jointly from unlabeled video by minimizing photometric reprojection error between adjacent frames; Monodepth~\cite{godard2017unsupervised} obtains the same signal from stereo pairs, and Monodepth2~\cite{godard2019digging} unifies both sources in one framework. GeoNet~\cite{yin2018geonet} adds optical flow to the joint estimate. Because all of these train a monocular network with stereo or multi-view supervision, the old separation between input configurations began to dissolve. Alongside, multi-task sharing (depth with semantic segmentation~\cite{wang2015towards}), ImageNet-pretrained encoders with structural losses (DenseDepth~\cite{alhashim2018high,deng2009imagenet,huang2017densely}), local planar guidance (BTS~\cite{lee2019big}), and mixed-dataset training for zero-shot transfer (MiDaS~\cite{MiDaS}) each pushed accuracy and generalization.

\paragraph{Transformers and Foundation Models.}
Self-attention~\cite{attention} gives every pixel access to global context, which helps precisely where local convolutions fail: textureless regions, occlusions, and large scenes. DPT~\cite{dpt} established the pattern of a Vision Transformer~\cite{vit} encoder with a light dense decoder, and DepthFormer~\cite{li2023depthformer} combined global attention with local convolutional refinement. Scaling data and model size then produced depth foundation models that are trained once and transferred zero-shot: ZoeDepth~\cite{zoedepth} and Depth Anything~\cite{DepthAnything} learn from large mixed corpora with flexible supervision and generalize across scenes and capture conditions without per-dataset adaptation.

\paragraph{Video and Generative Depth.}
Video adds temporal constraints. ManyDepth~\cite{watson2021temporal} aggregates neighboring frames into a cost volume at test time, whereas MonoViT~\cite{zhao2022monovit} keeps the self-supervised video objective but strengthens single-frame reasoning with a Transformer. A separate line treats depth as conditional generation: VPD~\cite{zhao2023unleashing} showed that text-to-image diffusion backbones carry usable geometric features, Marigold~\cite{ke2024repurposing} reframed monocular depth as conditional denoising and obtained globally coherent maps under hard conditions, and GeoWizard~\cite{fu2024geowizard} extended the recipe to joint depth and normal prediction. The same diffusion-as-perception idea recurs for surface normals and object pose later in this section.

\paragraph{Event-Based Depth.}
% 早期几何驱动 → 学习驱动方法和自监督硬件优化方法 → 综述 (Event-based Stereo Depth Estimation: A Survey 2024)
Event cameras record asynchronous per-pixel brightness changes~\cite{mahowald1992vlsi, lichtsteiner2008128}, offering microsecond latency and high dynamic range where frame cameras fail. Early event-based depth relied on geometric constraints and temporal accumulation~\cite{rebecq2018emvs, ieng2018neuromorphic,zhou2018semi, xie2017event}; learned event stereo~\cite{tulyakov2019learning, zhang2022discrete}, event--image fusion~\cite{mostafavi2021event, ghosh2022multi}, and self-supervised or hardware-aware frameworks~\cite{cho2023learning, de2025derd, cho2024temporal} followed. The line remains a niche complement for extreme dynamics; see the survey~\cite{ghosh2025event}.

\subsubsection{Datasets and Metrics}
Table~\ref{tab:depth_estimation_datasets} groups the main datasets into Type I, which ship explicit depth maps for direct supervision, and Type II, which require depth to be derived from LiDAR, stereo, or SfM. Three waves are visible. Indoor RGB-D collections (NYU Depth v2~\cite{silberman2012nyudv2}, TUM RGB-D~\cite{Sturm2012TUMRGBD}, SUN RGB-D~\cite{song2015sunrgbd}, later ScanNet~\cite{dai2017scannet} and Matterport3D~\cite{chang2017matterport3d}) and the LiDAR-based KITTI~\cite{Geiger2013KITTI,uhrig2017sparsity} fixed the indoor and outdoor protocols that still anchor the field, with Make3D~\cite{saxena2008make3d} as the laser-scanned precursor. Relative-depth supervision from internet photos (DIW~\cite{chen2016single}, ReDWeb~\cite{xian2018monocular}, MegaDepth~\cite{li2018megadepth}) and synthetic data (Scene Flow~\cite{mayer2016large}) then relaxed the need for metric sensors, while ETH3D~\cite{Schoeps2017ETH3D}, iBims-1~\cite{koch2018evaluation}, ApolloScape~\cite{wang2019apolloscape}, and Cityscapes stereo pairs~\cite{cordts2016cityscapes} served precision evaluation and self-supervised training. The current wave is driven by autonomous driving and rendering: multi-sensor fleets (nuScenes~\cite{caesar2020nuscenes}, Waymo Open~\cite{WaymoOpen-SunKDCPTGZCCVHN20}, Argoverse~2~\cite{wilson2023argoverse}, KITTI-360~\cite{liao2022kitti360}, DDAD~\cite{guizilini20203d}) supply millions of LiDAR-calibrated frames, photorealistic synthetic sets (Hypersim~\cite{roberts2021hypersim}, Virtual KITTI~2~\cite{Cabon2020VirtualKITTI2}, UrbanSyn~\cite{UrbanSyn}, TartanAir~\cite{TartanAir}) supply exact ground truth even for transparent surfaces and hard lighting, and DIODE~\cite{vasiljevic2019diode}, BlendedMVS~\cite{yao2020blendedmvs}, ARKitScenes~\cite{baruch2021arkitscenes}, Mannequin~\cite{li2019learning}, and WildRGB-D~\cite{xia2024rgbd} cover high-precision scans, handheld LiDAR, and dynamic or cluttered capture.

\begin{table}[!t]
\centering
\caption{Depth estimation datasets grouped by supervision type: Type~I datasets provide explicit depth maps for direct supervision; Type~II datasets require depth to be derived from raw LiDAR, stereo pairs, or SfM, often through projection or self-supervision.}
\label{tab:depth_estimation_datasets}
\fontsize{9}{11}\selectfont
\setlength{\tabcolsep}{2.5mm}
\begin{tabular}{lccccccr}
\toprule
\textbf{Name} & \textbf{Year} & \textbf{Scene} & \textbf{Type} & \textbf{Metric} & \textbf{Setup} & \textbf{Volume} & \textbf{Data Source} \\
\midrule
\rowcolor{lightgrey} \multicolumn{8}{l}{$\bullet$ \textit{Type I: Datasets with Explicit Depth Maps (Direct Supervision})} \\
Make3D~\cite{saxena2008make3d} & 2009 & Outdoor & Real & Metric & Mono & 534 & Laser Scanner \\
NYU Depth v2~\cite{silberman2012nyudv2} & 2012 & Indoor & Real & Metric & Mono & 1.4K & Kinect (SL) \\
TUM RGB-D~\cite{Sturm2012TUMRGBD} & 2012 & Indoor & Real & Metric & Mono & 39 Seqs & Kinect (SL) \\
SUN RGB-D~\cite{song2015sunrgbd} & 2015 & Indoor & Real & Metric & Mono & 10K & RGB-D Cameras \\
DIW~\cite{chen2016single} & 2016 & Varied & Real & Relative & Mono & 495K & Internet Photos \\
KITTI (Depth)~\cite{uhrig2017sparsity} & 2017 & Outdoor & Real & Metric & Stereo & 93K & LiDAR (Velodyne) \\
ScanNet~\cite{dai2017scannet} & 2017 & Indoor & Real & Metric & Mono & 2.5M & Structure Sensor \\
Matterport3D~\cite{chang2017matterport3d} & 2017 & Indoor & Real & Metric & Mono & 194K & Matterport (RGB-D) \\
ETH3D~\cite{Schoeps2017ETH3D} & 2017 & Both & Real & Metric & Multi-V & 13 Scans & Laser Scanner \\
ApolloScape~\cite{wang2019apolloscape} & 2018 & Outdoor & Real & Metric & Stereo & 140K & LiDAR + Stereo \\
iBims-1~\cite{koch2018evaluation} & 2018 & Indoor & Real & Metric & Mono & 100 & Camera + Laser \\
MegaDepth~\cite{li2018megadepth} & 2018 & Outdoor & Real & Relative & Mono & 130K & Internet + SfM \\
ReDWeb~\cite{xian2018monocular} & 2018 & Varied & Real & Relative & Mono & 3.6K & Internet + Stereo \\
DIODE~\cite{vasiljevic2019diode} & 2019 & Both & Real & Metric & Mono & 26K & Camera + LiDAR \\
BlendedMVS~\cite{yao2020blendedmvs} & 2020 & Both & Syn+Real & Metric & Multi-V & 17K & MVS Reconstruction \\
Hypersim~\cite{roberts2021hypersim} & 2020 & Indoor & Syn & Metric & Mono & 77K & Simulation \\
DDAD~\cite{guizilini20203d} & 2020 & Outdoor & Real & Metric & Multi-V & 76K & HD LiDAR \\
Virtual KITTI 2~\cite{Cabon2020VirtualKITTI2} & 2020 & Outdoor & Syn & Metric & Stereo & 21K & Simulation (Unity) \\
TartanAir~\cite{TartanAir} & 2020 & Varied & Syn & Metric & Stereo & $\sim$1M & Simulation (UE4) \\
ARKitScenes~\cite{baruch2021arkitscenes} & 2021 & Indoor & Real & Metric & Mono & 5K Clips & iPad LiDAR \\
UrbanSyn~\cite{UrbanSyn} & 2023 & Outdoor & Syn & Metric & Mono & 7.5K & Simulation \\
WildRGB-D~\cite{xia2024rgbd} & 2024 & Varied & Real & Metric & Mono & 20K & RGB-D Cameras \\
\midrule
\rowcolor{lightgrey} \multicolumn{8}{l}{$\bullet$ \textit{Type II: Datasets Requiring Depth to be Derived from Raw Geometry (LiDAR, Stereo or SfM)}} \\
KITTI (Raw)~\cite{Geiger2013KITTI} & 2012 & Outdoor & Real & Metric & Stereo & 42K & LiDAR (Velodyne) \\
Cityscapes~\cite{cordts2016cityscapes} & 2016 & Outdoor & Real & Metric & Stereo & 25K & Stereo Camera \\
Scene Flow~\cite{mayer2016large} & 2016 & Varied & Syn & Metric & Stereo & 39K & Simulation (Blender) \\
NuScenes~\cite{caesar2020nuscenes} & 2019 & Outdoor & Real & Metric & Multi-V & 1.4M & LiDAR \\
Waymo Open~\cite{WaymoOpen-SunKDCPTGZCCVHN20} & 2019 & Outdoor & Real & Metric & Multi-V & 200K & LiDAR \\
Mannequin~\cite{li2019learning} & 2019 & Varied & Real & Relative & Mono & 170K & Internet + SfM \\
KITTI-360~\cite{liao2022kitti360} & 2020 & Outdoor & Real & Metric & Stereo & 100K+ & LiDAR + Stereo \\
Argoverse 2~\cite{wilson2023argoverse} & 2022 & Outdoor & Real & Metric & Multi-V & 6M & LiDAR \\
\bottomrule
\end{tabular}

\end{table}

% 经典指标：RMSE、RMSE Log；强调相对偏差时：Abs Rel、Sq Rel；针对尺度模糊、引入不变性：Scale-Invariant Log Error；特定于深度估计，评估相对准确率：δ Threshold Accuracy；相对误差/排序指标主要用于相对深度估计的深度顺序关系
Depth is evaluated with pixel-wise regression metrics. The root mean squared error and its logarithmic variant,
\begin{equation}
\mathrm{RMSE} = \sqrt{ \frac{1}{N} \sum_{i=1}^{N} \left( \hat{d}_i - d_i \right)^2 }, \qquad
\mathrm{RMSE}_{\log} = \sqrt{ \frac{1}{N} \sum_{i=1}^{N} \left( \log \hat{d}_i - \log d_i \right)^2 },
\end{equation}
measure absolute deviation, the log form damping the dominance of far-range errors. Because an error of fixed size matters more at close range, relative errors normalize by the true depth,
\begin{equation}
\mathrm{Abs\ Rel} = \frac{1}{N} \sum_{i=1}^{N} \frac{ \left| \hat{d}_i - d_i \right| }{ d_i }, \qquad
\mathrm{Sq\ Rel} = \frac{1}{N} \sum_{i=1}^{N} \frac{ \left( \hat{d}_i - d_i \right)^2 }{ d_i },
\end{equation}
and Abs Rel is the primary metric on most benchmarks. For monocular prediction, whose scale is ambiguous, the scale-invariant log error discounts a global scale shift,
\begin{equation}
\mathrm{SI\text{-}Log} =
\frac{1}{N} \sum_{i=1}^{N} \left( \log \hat{d}_i - \log d_i \right)^2
-
\frac{1}{N^2} \left( \sum_{i=1}^{N} \left( \log \hat{d}_i - \log d_i \right) \right)^2 .
\end{equation}
Threshold accuracy reports the fraction of pixels within a multiplicative tolerance,
\begin{equation}
\delta =
\frac{1}{N}
\left|
\left\{
i \;\middle|\;
\max\!\left(
\frac{\hat{d}_i}{d_i},
\frac{d_i}{\hat{d}_i}
\right)
< \tau
\right\}
\right|, \qquad \tau \in \{1.25,\, 1.25^2,\, 1.25^3\}.
\end{equation}
Ordinal metrics, which score only the pairwise depth ordering, are used where supervision is relative; they track perceived layout but say nothing about metric precision.

On KITTI, supervised regressors such as BTS~\cite{lee2019big} and DPT~\cite{dpt} reach Abs Rel around $0.06$ with $\delta_{1.25}$ near $0.96$, while the self-supervised Monodepth2~\cite{godard2019digging} reaches Abs Rel near $0.10$ without any ground-truth depth. On NYU Depth v2, ZoeDepth~\cite{zoedepth} and Depth Anything~\cite{DepthAnything} reach $\delta_{1.25}$ of $0.96$ and $0.98$ after fine-tuning, and UniDepth~\cite{UniDepth} matches them zero-shot, which makes transfer a credible alternative to per-benchmark training. Generative models such as Marigold~\cite{ke2024repurposing} match or exceed regression on scale-invariant metrics while producing sharper boundaries, so recent gains come more from representation and data scale than from architectural depth. Boundary sharpness, temporal stability, and latency are still not evaluated in a standard way.

\subsubsection{Challenges}
Monocular depth is ill-posed: distinct 3D scenes project to the same image, so a single view alone does not uniquely determine 3D structure or metric scale and predictions drift across scenes. Scene content adds its own failures. Moving objects break the static-world assumption behind self-supervised and multi-view methods, and in video the result is flicker between frames and scale drift over time. Transparent and reflective materials and textureless regions give weak or misleading appearance cues; thin structures and sharp boundaries lose detail; low light, adverse weather, and underwater imagery degrade models trained on clean data. These scene effects compound into a domain gap: models are tied to the geometry and depth distribution of their training data, the indoor--outdoor divide is especially wide, real depth labels are scarce and noisy, and synthetic data still transfers imperfectly.

Current work attacks these problems from several sides. Depth foundation models~\cite{DepthAnything,UniDepth,zoedepth} supply strong zero-shot priors, and Video Depth Anything~\cite{chen2025video} adds temporal modules to one of them for consistent video depth, while video depth diffusion (DepthCrafter~\cite{hu2025depthcrafter}, ChronoDepth~\cite{shao2025learning}, RollingDepth~\cite{ke2025video}) pursues the same goal generatively. Coupling depth with 3D scene representations (DepthSplat~\cite{xu2024depthsplat}) improves spatial consistency, language cues help anchor metric scale~\cite{zeng2024rsa}, and lightweight architectures such as MonoLENS~\cite{higashiuchi2025monolens} and LMDepth~\cite{long2025lmdepth} target the real-time budgets of driving, robotics, and AR/VR.





\subsection{Surface Normal Estimation}
\label{sec:surface_normal}

Surface normal estimation predicts the local orientation of visible surfaces from images or videos.
Where depth answers ``how far is this point from the camera?'', a normal answers ``which way does the surface face?''.
That makes normals a compact, scale-invariant description of local 3D geometry: a normal field exposes planes, folds, object boundaries, and contact surfaces even when metric depth is ambiguous, which is why it serves reconstruction~\cite{wang2021neus}, relighting and intrinsic decomposition~\cite{munkberg2022extracting,barron2015shape}, physical interaction modeling~\cite{Bear2021Physion}, and egocentric perception~\cite{grauman2022ego4d}.

The learning problem looks simple: from an RGB image, RGB-D frame, multi-view set, or video, output one unit vector per valid pixel.
The inverse problem underneath is badly constrained, since the same intensity can be explained by different combinations of shape, reflectance, and illumination, and camera intrinsics change how image gradients map to 3D orientation.
Progress has therefore come from three directions at once: more and better supervision, stronger geometric inductive biases, and visual priors inherited from large dense-prediction and diffusion models; since 2024 the last of these has turned normal estimation into one output head of broader geometric foundation models.



\subsubsection{Task Definition}
Let $I\in\mathbb{R}^{H\times W\times 3}$ be an image, $\mathcal{V}$ optional auxiliary observations (neighboring frames, stereo or multi-view images, depth, intrinsics, or poses), and $\Omega$ the set of valid pixels.
The goal is to estimate
\begin{equation}
    f_\theta(I,\mathcal{V}) = \hat{\mathbf{N}}, \qquad
    \hat{\mathbf{N}}(u) = \hat{\mathbf{n}}_u \in \mathbb{S}^2,\quad u\in\Omega,
\end{equation}
where each $\hat{\mathbf{n}}_u$ is a unit vector in the camera frame or in a gravity- or world-aligned frame.
Supervision comes from depth sensors, fitted meshes, rendered assets, multi-view reconstruction, or indirectly from photometric and intrinsic-image losses, and camera metadata matters because a pixel's viewing ray depends on focal length and field of view.
Monocular prediction from one RGB image is the canonical benchmark setting; RGB-D variants use measured or predicted depth as input or supervision through the differential relation between depth gradients and normals, multi-view and video variants enforce viewpoint or temporal consistency, object-centric variants target CAD-like objects, hands, faces, and bodies, and panoramic variants predict under equirectangular or fisheye geometry, where pinhole assumptions break.


\begin{figure}[!t]
    \centering
    \includegraphics[width=\linewidth]{figures/part3/section6/surface_normal_estimation.pdf}
    \caption{Surface normal estimation. \textit{Top}: an RGB, video, or multi-view input is mapped by a dense prediction network ((a) ViT-like or (b) diffusion variant) to a per-pixel unit normal field $\hat{\mathbf{n}}\in\mathbb{S}^2$, scored by angular error. \textit{Bottom}: method timelines along three threads: joint depth--normal prediction, intrinsic/shading decomposition, and backbone architecture.}
    \label{fig:p3_s6_2_3_normal_overview}
\end{figure}

\subsubsection{Methodology}
Surface normal estimation has moved from physics-based inverse rendering to supervised dense prediction and, most recently, to foundation-model and diffusion-based geometry estimation.
Figure~\ref{fig:p3_s6_2_3_normal_overview} summarizes the method families; the last two below are grouped as one family in the figure.

\paragraph{Photometric and intrinsic-decomposition methods.}
The classical origin is shape from shading and photometric stereo.
Photometric stereo recovers orientation from several images taken under different lighting; Woodham's formulation~\cite{woodham1980photometric} shows that normals follow once illumination and reflectance are sufficiently constrained.
Single-image shape from shading is harder because a change in intensity may signal a change in orientation, albedo, illumination, or material; SIRFS~\cite{barron2015shape} is the canonical optimization-era method, jointly recovering shape, illumination, and reflectance under natural-image and geometric priors.
These methods identify the physical ambiguities that every later learned model inherits, and physically based rendering datasets such as Hypersim~\cite{roberts2021hypersim}, where shading, materials, depth, and normals are jointly available, keep their logic in play.

\paragraph{Supervised CNN and multi-task dense prediction.}
Eigen and Fergus~\cite{eigen2015predicting} trained one multi-scale convolutional network to predict depth, normals, and semantic labels, establishing that normals can be regressed directly from RGB and that geometric tasks benefit from shared representations.
Later CNN systems improved encoder capacity, decoder resolution, and boundary preservation, usually with normals as one branch of a multi-task perception model, and Taskonomy~\cite{zamir2018taskonomy} and Omnidata~\cite{eftekhar2021omnidata} systematized the paradigm by training dense predictors across many 3D scans and tasks.
A CNN normal estimator runs in one forward pass and drops into real-time reconstruction, AR, or robotics pipelines; its local features, however, blur thin structures and lack the global context needed for consistent orientation across large planes.

\paragraph{Transformer and camera-aware discriminative predictors.}
Vision Transformers let every output pixel attend to the whole image.
DPT~\cite{dpt} showed that Transformer features support strong dense prediction across depth and semantics, and the same global context helps normals: co-planar walls, floors, and object faces may be far apart in the image yet geometrically coupled.
Pure attention discards one geometric fact, that pixels correspond to camera rays; DSINE~\cite{bae2024dsine} injects ray-direction information and shows that camera geometry remains valuable even for large learned backbones.

\paragraph{Diffusion and foundation-model priors.}
A more recent line casts normal estimation as conditional generation or as one head of a general geometric foundation model.
VPD~\cite{zhao2023unleashing} showed that text-to-image diffusion models carry useful features for dense perception, and GeoWizard~\cite{fu2024geowizard} extended the Marigold recipe of \S\ref{sec:depth_estimation} to joint depth and normal prediction under a shared generative prior.
Metric3Dv2~\cite{hu2024metric3dv2} took the discriminative route, training a camera-aware monocular model for metric depth and normals across heterogeneous datasets.
StableNormal~\cite{ye2024stablenormal} addressed a practical weakness of diffusion-based normals: stochastic denoising yields high-quality but high-variance maps, so reducing variance matters as much as improving the mean.
The generative formulation samples uncertain or textureless regions from a learned distribution instead of forcing a single local regression, at the cost of iterative-denoising latency, which motivates distillation and single-step variants.

\paragraph{Temporal, multi-view, and reconstruction-coupled methods.}
Normals matter most when predictions must persist across time or support 3D reconstruction.
Video normal estimation enforces temporal consistency by warping predictions with optical flow or camera poses and penalizing flicker; NormalCrafter~\cite{wang2025normalcrafter} targets temporally consistent sequences for downstream video reconstruction and editing rather than per-frame accuracy alone.
Multi-view settings derive normals from reconstructed meshes or regularize them through cross-view photometric consistency, and neural surface reconstruction such as NeuS~\cite{wang2021neus} consumes normal-like differential geometry implicitly through signed-distance gradients, so normal estimation and reconstruction now supervise each other.

\paragraph{Material- and domain-specific normal estimation.}
A smaller branch revisits what generic models still get wrong.
Transparent, reflective, and laboratory objects violate Lambertian assumptions because their gradients come from refraction, reflection, or the background rather than the surface; TransNormal~\cite{li2026transnormal} targets transparent objects with a dedicated benchmark and model, making explicit that open-world normal prediction cannot be solved by indoor RGB-D supervision alone.



\begin{table}[!t]
\centering
\caption{Representative datasets and benchmarks for surface normal estimation. The ``Normal supervision/use'' column indicates whether normals are directly rendered, derived from depth or meshes, or indirectly related through intrinsic-image supervision. Counts are indicative of commonly used versions and may vary by split or preprocessing protocol.}
\label{tab:normal_dataset}
\fontsize{9}{11}\selectfont
\setlength{\tabcolsep}{4pt}
\begin{tabular}{L{0.15\textwidth}C{0.1\textwidth} C{0.15\textwidth}C{0.15\textwidth}L{0.35\textwidth}}
\toprule
\textbf{Dataset} & \textbf{Year} & \textbf{Domain} & \textbf{Source} & \textbf{Normal supervision / use} \\
\midrule
NYU Depth v2~\cite{silberman2012nyudv2} & 2012 & Indoor & RGB-D & Normals derived from Kinect depth; canonical indoor monocular normal benchmark. \\
SUN RGB-D~\cite{song2015sunrgbd} & 2015 & Indoor & RGB-D & Depth-derived normals across diverse indoor scenes and sensors. \\
ScanNet~\cite{dai2017scannet} & 2017 & Indoor & RGB-D video + mesh & Mesh- or depth-derived normals with camera trajectories; useful for video and reconstruction-linked supervision. \\
Matterport3D~\cite{chang2017matterport3d} & 2017 & Indoor panorama & RGB-D scans + mesh & High-resolution reconstructed geometry; supports room-scale and panoramic normal evaluation. \\
Taskonomy~\cite{zamir2018taskonomy} & 2018 & Indoor & 3D scans & Large multi-task supervision including surface normals, depth, reshading, and semantics. \\
Replica~\cite{straub2019replica} & 2019 & Indoor synthetic / scanned & Rendered geometry & Clean normals from high-quality reconstructed scenes; useful for controlled evaluation. \\
DIODE~\cite{vasiljevic2019diode} & 2019 & Indoor / outdoor & LiDAR depth & Depth-derived normals; stresses cross-domain generalization beyond NYU-style indoor scenes. \\
Hypersim~\cite{roberts2021hypersim} & 2021 & Indoor synthetic & PBR rendering & Ground-truth depth, normals, materials, and lighting for geometry + intrinsic learning. \\
Omnidata~\cite{eftekhar2021omnidata} & 2021 & Indoor multi-source & 3D scans & Large-scale mid-level vision training data with normals and many related tasks. \\
\bottomrule
\end{tabular}
\end{table}


\subsubsection{Datasets and Metrics}
\label{sec:normal_performance}

Table~\ref{tab:normal_dataset} shows that normal supervision is usually inherited from another representation rather than annotated directly.
RGB-D datasets (NYU Depth v2~\cite{silberman2012nyudv2}, SUN RGB-D~\cite{song2015sunrgbd}, ScanNet~\cite{dai2017scannet}) supply measured depth from which local planes or mesh normals are computed, which fixes the indoor evaluation regime but also means the labels are only as good as the depth sensor and are dominated by man-made indoor structure.
Scanned and synthetic collections address coverage and label quality from opposite sides: Matterport3D~\cite{chang2017matterport3d}, Taskonomy~\cite{zamir2018taskonomy}, and Omnidata~\cite{eftekhar2021omnidata} scale supervision across many environments and tasks, whereas Replica~\cite{straub2019replica} and Hypersim~\cite{roberts2021hypersim} give clean rendered normals with physically meaningful lighting at the price of a synthetic-to-real gap.
DIODE~\cite{vasiljevic2019diode} pushes toward indoor--outdoor generalization, yet outdoor normals remain far less standardized than outdoor depth because sparse LiDAR is a poor match for dense orientation supervision; intrinsic-image sets such as IIW~\cite{bell2014intrinsic} and MPI Sintel Intrinsics~\cite{chen2013simple} constrain shading and reflectance for inverse-rendering formulations without being normal benchmarks themselves.
The most visible gap is temporal: although ScanNet and other RGB-D videos contain sequences, there is no widely adopted video normal benchmark, so temporal stability is reported qualitatively or through method-specific warping losses, even though flicker often hurts reconstruction and AR consumers more than a small rise in mean angular error.




Evaluation is angular because normals are unit vectors.
With $\hat{\mathbf{n}}_i,\mathbf{n}_i\in\mathbb{S}^2$ the predicted and ground-truth normals at valid pixel $i$ and $N$ the number of valid pixels, the per-pixel error in degrees is
\begin{equation}
   e_i = \frac{180}{\pi}\arccos\!\left(\mathrm{clip}(\hat{\mathbf{n}}_i^\top \mathbf{n}_i,-1,1)\right),
\end{equation}
summarized by its mean (sensitive to boundary noise and invalid-depth artifacts), its median (robust to a few severe failures), occasionally its RMSE, and by threshold accuracies
\begin{equation}
    \mathrm{Acc}_{\tau} = \frac{1}{N}\sum_{i=1}^{N}\mathbf{1}(e_i < \tau),
    \qquad \tau\in\{11.25^\circ,22.5^\circ,30^\circ\},
\end{equation}
which expose changes in the fraction of practically usable pixels that a small shift in mean error can hide.
Training objectives usually minimize the cosine loss $1-\hat{\mathbf{n}}_i^\top\mathbf{n}_i$, monotonic in angle for unit vectors, but benchmarks report angles for interpretability.
Beyond per-pixel angle, joint depth--normal models are checked for agreement between predicted normals and normals differentiated from depth, reconstruction-oriented work asks whether normal discontinuities align with object boundaries and whether large planes stay coherent, and video methods report temporal consistency by warping predictions between adjacent frames with flow or pose,
\begin{equation}
    \mathrm{TempCons} = \frac{1}{M}\sum_{(i,t)}
    \left\|\hat{\mathbf{n}}_{i}^{\,t} -
    \mathcal{W}_{t+1\rightarrow t}(\hat{\mathbf{n}}^{\,t+1})_i\right\|_2,
\end{equation}
after rotating camera-frame normals into a common frame.
Protocol details move rankings: depth-derived normals are unstable at discontinuities, papers mask invalid pixels differently, and some evaluate in the camera frame while others rotate into a gravity-aligned frame, so the table below records which paper produced each number.

\begin{table}[!t]
    \centering
    \caption{Representative surface-normal estimators on NYU Depth v2 (lower mean / median angular error and higher $\mathrm{Acc}_{11.25^\circ}$ are better). Rows marked $\dagger$ are taken from the re-evaluation in StableNormal~\cite{ye2024stablenormal}, under whose protocol DSINE scores $18.6^\circ$ / $9.9^\circ$ / $56.1\%$; the other rows are as reported by the cited paper (Omnidata as re-evaluated by DSINE~\cite{bae2024dsine}). Numbers across protocols are indicative, not strictly comparable.}
    \label{tab:normal_performance}
    \fontsize{9}{11}\selectfont
    \setlength{\tabcolsep}{2.2mm}
    \begin{tabular}{L{0.26\textwidth} L{0.28\textwidth} c c c c}
    \toprule
    \textbf{Method} & \textbf{Family} & \textbf{Year} & \textbf{Mean}$\downarrow$ & \textbf{Median}$\downarrow$ & $\mathbf{Acc}_{11.25^\circ}\uparrow$ \\
    \midrule
    Eigen \& Fergus~\cite{eigen2015predicting} & Multi-scale CNN, multi-task & 2015 & 20.9 & 13.2 & 44.4 \\
    Omnidata~\cite{eftekhar2021omnidata} & Large-scale mid-level supervision & 2021 & 23.1 & 12.9 & 45.8 \\
    DSINE~\cite{bae2024dsine} & Ray-aware discriminative & 2024 & 16.4 & 8.4 & 59.6 \\
    Metric3Dv2 (ViT-g, zero-shot)~\cite{hu2024metric3dv2} & Geometric foundation model & 2024 & 13.3 & 7.0 & 66.4 \\
    Marigold (normal variant)$^\dagger$~\cite{ke2024repurposing} & Diffusion prior & 2024 & 20.9 & 11.1 & 50.5 \\
    GeoWizard$^\dagger$~\cite{fu2024geowizard} & Joint depth + normal diffusion & 2024 & 20.4 & 11.9 & 47.0 \\
    StableNormal$^\dagger$~\cite{ye2024stablenormal} & Low-variance diffusion & 2024 & 19.7 & 10.5 & 53.0 \\
    \bottomrule
    \end{tabular}
\end{table}

Table~\ref{tab:normal_performance} makes two things visible. Across the reported settings, camera-aware models achieve low angular errors: DSINE roughly halves the median error of the 2015 baseline, and Metric3Dv2 reaches $13.3^\circ$ mean error zero-shot, so projective inductive bias and large camera-diverse pretraining still pay off even with strong backbones. Diffusion-based estimators trail these numbers on NYU Depth v2 under the StableNormal protocol, but they trade accuracy on clean indoor scenes for coherence in textureless or ambiguous regions and for a usable uncertainty, which the single-frame benchmark does not reward; the post-GeoWizard line has accordingly shifted the target from single-frame accuracy toward stable, usable geometry, StableNormal controlling diffusion variance and NormalCrafter~\cite{wang2025normalcrafter} enforcing temporal coherence.

\subsubsection{Challenges}

An RGB image alone does not determine shape, lighting, and reflectance: a dark pixel may be a tilted surface, a dark material, a cast shadow, or a reflection of a dark region, and learned models resolve the ambiguity with priors rather than removing it, which is why estimators look confident on common indoor layouts and fail systematically on rare materials, unusual lighting, and out-of-distribution scenes.
The persistent failure modes follow from this.
\textit{Reflective and transparent surfaces} produce gradients caused by the environment rather than the surface.
\textit{Textureless and weakly lit regions} give little evidence, and networks hallucinate planar sheets.
\textit{Thin structures and depth discontinuities} mix receptive fields and inherit depth-derived label noise, blurring normals around chair legs, plants, railings, and silhouettes.
\textit{Wide-field and panoramic cameras} break models trained under narrow pinhole assumptions, which is what makes ray-aware designs such as DSINE~\cite{bae2024dsine} necessary.
\textit{Dynamic and deformable objects} flicker because most supervision comes from static scans or rigid synthetic scenes.

Dataset bias compounds the problem: canonical benchmarks are indoor and sensor-derived, so outdoor surfaces, hands, faces, deformable and transparent objects, and egocentric manipulation stay underrepresented, and no accepted benchmark jointly measures cross-domain generalization, temporal stability, and downstream reconstruction quality. TransNormal~\cite{li2026transnormal} is a corrective for one slice of the non-Lambertian regime (glass, liquid, mirror, metal, wet roads), which remains far from solved.
The directions now taking shape answer these gaps directly: geometric foundation models that merge depth, normals, segmentation, material, and calibration into one dense-prediction backbone, following Omnidata~\cite{eftekhar2021omnidata} through Metric3Dv2~\cite{hu2024metric3dv2} and GeoWizard~\cite{fu2024geowizard}; diffusion distillation and variance control that keep the coherence and uncertainty of generative models while recovering real-time, deterministic inference~\cite{ye2024stablenormal}; temporally consistent, reconstruction-coupled normals~\cite{wang2025normalcrafter}; and physically grounded models for non-Lambertian materials, where data scale is no substitute for modeling refraction and reflection.



\subsection{Height and Distance Estimation}
\label{sec:height_distance}
Height and distance estimation recovers absolute, real-world lengths: the dimensions of an object (the height of a person), the separation between scene entities (the gap between two vehicles), or the distance traveled by a moving object over a known time interval, from which average speed follows.
Whereas depth estimation (\S\ref{sec:depth_estimation}) recovers a dense per-pixel distance from the sensor to every visible surface, height and distance estimation targets a few metric quantities defined between scene entities rather than relative to the sensor, and is therefore tied to specific object instances and to a downstream measurement requirement.
Radar, ultrasonic, and radio ranging measure such distances directly and are mature; this subsection covers the vision-based problem.

\subsubsection{Task Definition}
The input is one or more images or a video together with a specification of what to measure: the height of an object, the distance between two objects, or the distance an object travels. The output is a length in physical units, or a traveled distance over a known time interval, whose ratio gives average speed (\textbf{speed estimation}).

Every vision-based route to a metric length passes through the camera model, shown in Figure~\ref{fig:p3_s6_2_2_camera_model}. Under the pinhole model a world point $P$ with homogeneous coordinates $\mathbf{P}_w$ projects to the pixel
\begin{equation}
    \mathbf{p} \sim K\,[\,R \mid \mathbf{t}\,]\,\mathbf{P}_w,
\end{equation}
where the extrinsics $(R,\mathbf{t})$ move the point into the camera frame and the intrinsic matrix $K$ (focal length, principal point, pixel aspect) maps normalized image-plane coordinates to pixels. Real lenses deviate from this model. A common forward form of Brown's radial distortion model~\cite{brown1971close} is
\begin{equation}
    \begin{pmatrix} x_{\text{dist}} \\ y_{\text{dist}} \end{pmatrix} =
    \begin{pmatrix} x(1 + k_1 r^2 + k_2 r^4 + k_3 r^6) \\ y(1 + k_1 r^2 + k_2 r^4 + k_3 r^6) \end{pmatrix},
\end{equation}
where $(x,y)$ are undistorted normalized coordinates and $r^2=x^2+y^2$. Tangential terms model sensor--lens misalignment, and camera calibration~\cite{zhang2000flexible} recovers $K$ together with the distortion coefficients. Without calibration, every pixel-to-length conversion inherits an unknown scale.

\begin{figure}[!ht]
    \centering
    \includegraphics[width=0.5\linewidth]{figures/part3/section6/camera-1.png}
    \caption{Pinhole camera model: world coordinates map to camera, image-plane, and pixel coordinates through the extrinsic, projection, and affine matrices; the last two compose into the intrinsic matrix $M_{\mathrm{intrinsic}}$ and all three into the full projection $M$.}
    \label{fig:p3_s6_2_2_camera_model}
\end{figure}

Beyond calibration, geometric pipelines rely on a small set of recurring intermediate quantities, illustrated in Figure~\ref{fig:p3_s6_2_2_subtasks}: vanishing points, where parallel world lines converge and from which camera orientation follows; the horizon line, which fixes the camera's tilt relative to the ground; and a ground-plane homography, which lets lengths on a dominant plane such as a road or floor be measured once a single reference length is known. Stereo pipelines add rectification, which aligns the two image planes so that correspondences lie on the same scanline, and disparity matching. Learning-based pipelines keep object detection (Section~\ref{sec:6-1-2}) as the front end and replace the explicit geometry with regression of depth, 3D center, and dimensions from region features, optionally conditioned on a dense depth map and refined over time when video is available.

\begin{figure}[!t]
    \centering
    \includegraphics[width=\linewidth]{figures/part3/section6/height_estimation_1.pdf}
    \caption{Geometric subtasks supporting height and distance estimation: (a) parallel street edges determine a vanishing point; (b) multiple views of a known planar pattern estimate camera intrinsics and lens distortion; (c) the ground-plane horizon passes through horizontal-direction vanishing points. Street geometry and the separate camera-pitch inset are illustrative.}
    \label{fig:p3_s6_2_2_subtasks}
\end{figure}

\subsubsection{Methodology}
Vision-based height and distance estimation is organized along two axes: the sensor configuration (stereo versus monocular), which decides whether metric depth is available by triangulation, and the inference paradigm (geometry-based versus learning-based), which decides whether scene geometry is solved from projective constraints or inferred from data. Figure~\ref{fig:p3_s6_2_2_method_overview} lays out the resulting families; it draws the uncalibrated case as a third branch for readability, but it is best read as the weakly constrained regime of monocular or stereo inference in which some geometric information is missing and must be supplied by priors.

\paragraph{Stereo Vision.}
A calibrated stereo rig recovers metric depth directly from disparity, $Z = fB/d$ with focal length $f$, baseline $B$, and disparity $d$.
Classical pipelines run calibration, rectification, matching, disparity estimation, and triangulation~\cite{zhang2000flexible,Scharstein2002Stereo,hartley2003multiple}; their accuracy rests on calibration quality and on matching reliability in textureless, reflective, or occluded regions, and object distance is read off by aggregating depth over the detected region. Multi-view stereo adds views for consistency without changing the triangulation principle~\cite{Seitz2006Middlebury,yao2018mvsnet}.
Learned stereo replaces hand-crafted matching costs with learned features and cost volumes~\cite{kendall2017end,chang2018pyramid,lipson2021raftstereo} and, in recent systems, folds distance estimation into stereo 3D object detection, predicting category, 2D extent, depth, and 3D box jointly~\cite{li2019stereorcnn,chen2020dsgn,guo2021ligastereo}. Learning improves robustness, but the metric scale still comes from the calibrated baseline and intrinsics.

\paragraph{Monocular Vision.}
A single image does not fix absolute scale, so monocular methods recover it from explicit geometric constraints, learned priors, or external reference information~\cite{hartley2003multiple,criminisi2000single}.
Classical monocular metrology uses projective geometry: vanishing points of parallel world lines give camera orientation~\cite{caprile1990using,hartley2003multiple}, and a homography between a dominant ground plane and the image enables metric measurement on that plane once one reference length is known~\cite{criminisi2000single,liebowitz1999metric}. The approach is interpretable and accurate in controlled scenes but needs clean geometric cues, reliable calibration, and visible parallel lines or planar supports.
Learning-based monocular methods detect objects in 2D and regress their depth, 3D center, and physical dimensions directly from image features~\cite{chen2016monocular,mousavian2017deep,manei2019monodis,brzozowski2019m3drpn,liu2020smoke,wang2021fcos3d,liu2022monocon,ma2021delving}, implicitly learning the inverse relation between image extent and distance together with perspective, context, and occlusion cues, without an explicit vanishing-point or homography step.

\paragraph{Foundation-Model-Assisted Metric Estimation.}
A newer route uses pretrained metric or metric-aligned depth models as an intermediate representation. DPT, ZoeDepth, Metric3D, and Depth Anything predict dense depth from a single image with strong cross-domain generalization~\cite{dpt,zoedepth,yin2024metric3d,DepthAnything}; pooling the predicted depth inside a detected region or instance mask then gives object distance, and combining masks, depth, and simple geometric constraints gives object height. This removes the need for explicit vanishing-point detection and manual scene calibration and offers a practical path to metric estimation on in-the-wild images.
In the fully uncalibrated case, exact metric recovery from one image remains ill-posed~\cite{hartley2003multiple}; practical systems inject scale through a known reference object, a known camera height, a ground-plane assumption, category-level size priors, or one of the pretrained metric-depth models above~\cite{criminisi2000single,mousavian2017deep,zoedepth,yin2024metric3d}.

\begin{figure}[!t]
    \centering
    \includegraphics[width=\linewidth]{figures/part3/section6/height_estimation_2.pdf}
    \caption{Sources of metric scale in vision-based measurement: (a) calibrated stereo uses focal length, disparity, and a known baseline; (b) monocular metrology transfers a reference height using ground-plane and vertical-direction geometry; (c) a reference-assisted example of weakly calibrated monocular inference uses a known planar rectangle to rectify a shared plane and set scale. Scenes are illustrative; the panels do not imply a universal accuracy ranking.}
    \label{fig:p3_s6_2_2_method_overview}
\end{figure}

\subsubsection{Datasets and Metrics}
Table~\ref{tab:height_distance_estimate_benchmark} summarizes the benchmarks, which fall into three groups by sub-task. Object-size estimation is evaluated on the indoor RGB-D sets CAMERA25 and REAL275~\cite{Wang0HVSG19}, where category-level dimensions are scored jointly with pose on controlled tabletop scenes. Sensor-to-object distance is rarely benchmarked on its own; it is evaluated through the 3D-detection annotations (box centers, sizes, velocities) of the driving datasets KITTI~\cite{Geiger2013KITTI}, nuScenes~\cite{caesar2020nuscenes}, Waymo Open~\cite{WaymoOpen-SunKDCPTGZCCVHN20}, and Argoverse~2~\cite{wilson2023argoverse}, which test metric localization at scale across multi-camera rigs and scenes, with BDD100K~\cite{yu2020bdd100k} as a monocular video source. Speed and displacement estimation is dominated by fixed-camera traffic video: AI City Challenge 2018~\cite{Naphade_2018_CVPR_Workshops}, BrnoCompSpeed~\cite{SochorJSMSHZ17}, UTFPR~\cite{LuvizonNM17}, and QMUL~\cite{moazzam2019image}, plus the UAV-mounted Beihang videos~\cite{KeKLW15}; these are the only benchmarks that score a physically meaningful motion quantity rather than static geometry.

\begin{table*}[!ht]
    \centering
    \caption{Height and distance estimation benchmarks, grouped by the sub-task they score (object size, sensor-to-object distance via 3D boxes, or vehicle speed). BDD100K is included as an auxiliary video source.}
    \label{tab:height_distance_estimate_benchmark}
    \fontsize{9}{11}\selectfont
    \setlength{\tabcolsep}{1.0mm}
    \begin{tabular}{L{0.19\textwidth} C{0.07\textwidth} C{0.15\textwidth} C{0.09\textwidth} C{0.09\textwidth} C{0.12\textwidth} C{0.21\textwidth}}
    \toprule
    \textbf{Dataset} & \textbf{Task} & \textbf{Modality} & \textbf{View} & \textbf{Domain} & \textbf{Scale} & \textbf{Resolution} \\
    \midrule
    CAMERA25 \cite{Wang0HVSG19} & Size  & RGB-D image & Monocular & Mixed-reality indoor & 25K eval. images & - \\
    REAL275 \cite{Wang0HVSG19} & Size  & RGB-D image & Monocular & Indoor & 2,750 test images & - \\
    KITTI 3D~\cite{Geiger2013KITTI} & Distance & RGB + LiDAR     & Monocular     & Vehicle     & 15K frames     & - \\
    nuScenes~\cite{caesar2020nuscenes}     & Size / distance     & RGB + LiDAR + radar     & Multi-view     & Vehicle     & 1,000 scenes & - \\
    Waymo Open Dataset~\cite{WaymoOpen-SunKDCPTGZCCVHN20}     & Size / distance   & RGB + LiDAR     & Multi-view     & Vehicle     & 1,150 segments & - \\
    Argoverse 2 Sensor~\cite{wilson2023argoverse}     & Size / distance  & RGB + LiDAR     & Multi-view     & Vehicle     & 1,000 scenes & - \\
    AI City 2018 Track 1 \cite{Naphade_2018_CVPR_Workshops} & Speed  & Video & Monocular & Vehicle & 27 one-min videos & 1920$\times$1080 @30\,Hz \\
    BrnoCompSpeed \cite{SochorJSMSHZ17} & Speed  & Video & Monocular & Vehicle & 18 videos $\times$ 1 hour & 1920$\times$1080 @50\,Hz\\
    UTFPR \cite{LuvizonNM17} & Speed  & Video & Monocular & Vehicle & 5 hours & 1920$\times$1080 @30\,Hz\\
    QMUL \cite{moazzam2019image} & Speed & Video & Monocular & Vehicle & 1 hour & 360$\times$288 @25\,Hz \\
    Beihang UAV Videos \cite{KeKLW15} & Speed & Image & Monocular & Vehicle & 70 images & 960$\times$540\\
    BDD100K~\cite{yu2020bdd100k}     & Auxiliary     & Video     & Monocular     & Vehicle     & 100K videos & 720p @30\,Hz\\
    \bottomrule
    \end{tabular}
\end{table*}

Each sub-task is scored with the same triplet: an absolute error in physical units, a relative error for cross-range comparison, and a threshold accuracy for downstream binary decisions. For object size, the metric is the mean absolute error in physical units (centimeters for human height, millimeters for category-level pose),
\begin{equation}
    \mathrm{MAE} = \frac{1}{N}\sum_{i=1}^{N} |\hat{s}_i - s_i|,
\end{equation}
with the mean absolute percentage error $\mathrm{MAPE} = \tfrac{1}{N}\sum_i |\hat{s}_i - s_i|/s_i \times 100\%$ alongside for cross-class comparison; when size is scored jointly with pose, the 3D IoU of predicted and ground-truth boxes is thresholded, the NOCS protocol~\cite{Wang0HVSG19} reporting mAP at 3D IoU thresholds of $25\%$ and $50\%$, and pose mAP under joint rotation/translation thresholds such as $(5^\circ,5\,\mathrm{cm})$, $(10^\circ,5\,\mathrm{cm})$, and $(10^\circ,10\,\mathrm{cm})$. Inter-object distance reuses MAE, RMSE, and MAPE on distances $\hat{d}_i$ and the depth-style ratio accuracy $\delta_\tau$ of \S\ref{sec:depth_estimation} with $\tau \in \{1.10,\,1.25\}$; because stereo error grows as $\Delta Z \propto Z^2/(fB)$, KITTI-style protocols~\cite{Geiger2013KITTI} stratify by range (near $<10$\,m, mid $10$--$50$\,m, far $>50$\,m) so that far-field degradation is not averaged away. Speed estimation reports MAE and RMSE in physical speed units, and MAPE as a percentage. AI City Challenge 2018~\cite{Naphade_2018_CVPR_Workshops} reports speed in mph and defines $S_1=\mathrm{DR}(1-\mathrm{NRMSE})$ from vehicle detection rate and speed RMSE min--max normalized across submissions, while BrnoCompSpeed~\cite{SochorJSMSHZ17} also reports median absolute error to reduce sensitivity to outliers. Two protocol details change rankings: per-vehicle averaging (BrnoCompSpeed) weighs trajectories differently from per-frame averaging over localized ground-truth vehicles (AI City), and monocular and uncalibrated methods are sensitive to the assumed focal length, principal point, and reference length, so their reported accuracy is only meaningful together with the calibration assumptions.

Performance is sub-task specific. Category-level size recovery on REAL275~\cite{Wang0HVSG19} remains far from millimeter accuracy in cluttered scenes. On BrnoCompSpeed~\cite{SochorJSMSHZ17}, calibration-aware pipelines reach low single-digit km/h errors on the original sessions but degrade when transferred to unseen camera placements, exposing the brittleness of fixed-camera homography assumptions, and on the AI City 2018 speed track~\cite{Naphade_2018_CVPR_Workshops} the leading entries pair strong detection-and-tracking front ends with multi-frame speed regression, so the S1 score reflects detection recall as much as raw speed accuracy. Across all three sub-tasks the bottleneck is calibration, of the camera and of the scene-plane reference, rather than the regression backbone.

\subsubsection{Challenges}
% 说明这个问题其实在当前的空间理解中，是非常重要的一个任务评估形式，它主要在与空间理解相关的下游任务中，作为一个重要的评估指标，来衡量模型对于空间理解的能力。比如在自动驾驶领域，车辆的速度估计是一个非常重要的任务，它直接关系到车辆的安全性和行驶效率；在室内场景中，物体尺寸的估计也是一个重要的任务，它可以帮助机器人更好地与环境进行交互。因此，准确的高度和距离估计对于提升空间理解模型的实用性和可靠性具有重要意义。
Calibration sensitivity is the dominant open problem: small errors in intrinsics or in the scene-plane reference propagate non-linearly into every recovered length, and methods that treat focal length and reference size as latent variables, or refine calibration online from multi-view consistency, generalize across camera setups where fixed-calibration pipelines fail. The second gap is coverage: benchmarks are confined to indoor tabletop size estimation and vehicle-traffic speed, leaving pedestrian speed in crowds, package dimensions in logistics, and indoor inter-object distance essentially unmeasured. Fusing visual estimates with LiDAR, radar, or inertial sensing would lift robustness in low light and under occlusion, but the vision-only setting treated here still lacks a benchmark that scores metric accuracy jointly with calibration uncertainty.



\subsection{Visual Odometry and Camera Pose Estimation}
\label{sec:visual_odometry}

Visual odometry (VO) and camera pose estimation recover the position and orientation of a camera from visual observations.
They provide the coordinate scaffold that lets a system stitch individual images into a persistent spatial representation: depth, surface normals, scene flow, maps, and 3D reconstructions all become more useful once they can be expressed in a shared camera or world frame.
VO usually means estimating \emph{incremental} motion between consecutive frames, often as a trajectory $\{\mathbf{T}_t\}_{t=1}^{T}$ in $SE(3)$; camera pose estimation is broader and may recover a query image's absolute 6-DoF pose in a known scene, the relative pose between two views, or camera motion and scene geometry jointly from an unordered image collection.
The field was built on geometric optimization and SLAM, and its center of gravity is now shifting toward learned correspondence, differentiable optimization, dense geometric priors, and feed-forward geometric foundation models that predict cameras and geometry together.

\subsubsection{Task Definition}

Let $\mathbf{T}_t\in SE(3)$ denote the $4\times4$ homogeneous camera-to-world transformation for frame $t$, with rotation $\mathbf{R}_t$ and translation $\mathbf{t}_t$.
Given images $\{I_t\}$ (monocular, stereo, RGB-D, multi-view, or video), optional intrinsics $\{\mathbf{K}_t\}$, optional inertial readings, and sometimes a prior map or database, the goal is to estimate relative camera motions $\Delta\mathbf{T}_{i\rightarrow j}=\mathbf{T}_i^{-1}\mathbf{T}_j$ expressed in frame $i$, an absolute pose $\mathbf{T}_q$ for a query image, or a full trajectory together with sparse or dense scene geometry, and occasionally the intrinsics themselves when calibration is unknown.
The variants differ in what is assumed known.
\textbf{Monocular VO} is scale ambiguous without external priors; \textbf{stereo and RGB-D VO} recover metric scale from a known baseline or measured depth; \textbf{visual-inertial odometry} uses IMU measurements to stabilize scale and fast rotation; \textbf{SLAM} adds mapping and loop closure to online odometry; \textbf{SfM} recovers poses and structure from unordered collections; \textbf{map-based relocalization} assumes an existing scene representation and places a query in it; and \textbf{map-free relocalization and feed-forward reconstruction} skip explicit map construction by predicting relative pose or globally aligned point maps directly from images.

\begin{figure}[!t]
    \centering
    \includegraphics[width=\linewidth]{figures/part3/section6/visual_odometry_and_camera_pose_estimation.pdf}
    \caption{Visual odometry and camera pose estimation. \textit{Top}: a video or unordered image set is mapped to relative motion $\Delta\mathbf{T}_{i\rightarrow j}$, an absolute pose $\mathbf{T}_q$ in a known scene, or camera poses $\{\mathbf{T}_t\}$ jointly with scene structure. \textit{Bottom}: method lanes (a) geometry-first, (b) learned and self-supervised VO, (c) hybrid learned--geometric, (d) relocalization, and (e) feed-forward foundation models. Mechanism arrows indicate operations; dated timelines summarize representative methods in publication order, with non-proportional spacing and no implied computational dependency. The dashed warp-feedback path denotes self-supervised training.}
    \label{fig:p3_s6_2_4_vo_overview}
\end{figure}

\subsubsection{Methodology}

Figure~\ref{fig:p3_s6_2_4_vo_overview} summarizes five method families, ordered from explicit geometry to learned geometric priors to feed-forward scene-level geometry.

\paragraph{Geometry-first SfM, VO, and SLAM.}
Classical pipelines detect repeatable features, match them across views, fit epipolar geometry with robust estimators, and refine poses and 3D points by bundle adjustment~\cite{triggs2000bundle,agarwal2010bundle}.
COLMAP~\cite{schonberger2016structure} made this recipe robust enough for internet photo collections, and GLOMAP~\cite{pan2024glomap} revisits the pose-graph and averaging components of global SfM for large-scale reconstruction.
For sequential input, geometric VO and SLAM track camera motion, with SLAM systems additionally maintaining maps and often closing loops: MonoSLAM~\cite{davison2007monoslam}, PTAM~\cite{KleinMurray2007PTAM}, and ORB-SLAM~\cite{mur2015orb} define the feature-based branch, LSD-SLAM~\cite{engel2014lsd} and DSO~\cite{engel2017direct} the direct branch.
Explicit geometric consistency is the family's strength; weak texture, dynamic objects, rolling shutter, illumination shifts, and small-baseline or planar degeneracies are its failure modes.

\paragraph{Learned VO and self-supervised view synthesis.}
End-to-end VO regresses camera motion from image pairs or short windows, DeepVO~\cite{Wang2017DeepVO} being the standard early example with convolutional features and recurrent temporal modeling.
Self-supervised VO instead learns depth and ego-motion together by minimizing photometric reconstruction error between adjacent frames, with SfM-Learner~\cite{zhou2017unsupervised} as the canonical starting point.
This made monocular video supervision scalable, but the geometry remains bound to photometric assumptions (static Lambertian scenes, stable intrinsics, sufficient parallax) and to monocular scale ambiguity unless stereo, IMU, or learned metric priors are added.

\paragraph{Hybrid learned geometry and differentiable optimization.}
The most successful learned systems keep geometric optimization and learn the parts that are hard to hand-design.
Learned features and matchers (SuperPoint~\cite{detone2018superpoint}, SuperGlue~\cite{sarlin2020superglue}, LoFTR~\cite{Sun2021LoFTR}) improve the correspondence front end, and VGG-SfM~\cite{Vggsfm} and ACE0~\cite{acezero} push learned features and scene-coordinate learning into SfM and relocalization.
DROID-SLAM~\cite{teed2021droid} is the representative differentiable-optimization system, coupling dense correspondence, depth, and camera motion through recurrent updates and a dense bundle-adjustment layer, and LEAP-VO~\cite{chen2024leap} stabilizes motion over longer windows through long-term point tracking.

\paragraph{Relocalization and map-based pose estimation.}
Given a query image and a known scene or database, relocalization estimates the query's 6-DoF pose.
PoseNet~\cite{kendall2015posenet} popularized absolute pose regression, which is fast but less accurate than geometry.
Scene-coordinate regression, from Shotton \etal~\cite{Scene_Coordinate} to DSAC~\cite{brachmann2017dsac} and ACE-style relocalizers~\cite{acezero}, predicts a 3D scene coordinate per pixel and solves pose with PnP and RANSAC; hierarchical localization~\cite{sarlin2019coarse} retrieves candidate database images, matches, and verifies geometrically, and remains the high-accuracy reference.
Relative-pose regressors trained on much larger collections, such as Reloc3r~\cite{reloc3r}, revisit the speed and generalization trade-off and show that regression recovers a role when paired with scale and data.

\paragraph{Feed-forward geometric foundation models.}
The main post-2024 shift is toward models that predict scene geometry and camera relations directly, often without SfM initialization.
DUSt3R~\cite{wang2024dust3r} regresses dense 3D point maps from image pairs, providing learned geometry from which correspondences and camera parameters are recovered using geometric post-processing; MASt3R~\cite{mast3r} grounds image matching in this 3D representation, and MASt3R-SfM~\cite{MASt3R_SfM} builds an unconstrained SfM pipeline on it.
MegaSaM~\cite{li2025megasam} and MonST3R~\cite{monst3r} extend the idea to casual dynamic video, where static-scene assumptions fail; CUT3R~\cite{wang2025Cut3R} adds persistent state for continuous 3D perception, Fast3R~\cite{yang2025fast3r} scales feed-forward reconstruction to very large image sets, and VGGT~\cite{wang2025vggt} predicts camera poses, depth, and point maps jointly with a visual-geometry-grounded Transformer.
Geometry does not disappear in this family; it moves from a hand-coded optimization loop into a learned prior that supplies initialization, dense correspondence, or complete camera-and-scene predictions.

\begin{table}[!t]
\centering
\caption{Representative datasets and benchmarks for visual odometry, camera pose estimation, relocalization, and feed-forward scene geometry. Counts and splits vary across protocols.}
\label{tab:vo_pose_datasets}
\fontsize{9}{11}\selectfont
\setlength{\tabcolsep}{4pt}
\begin{tabular}{L{0.15\textwidth}C{0.1\textwidth} C{0.15\textwidth}C{0.15\textwidth}L{0.35\textwidth}}
\toprule
\textbf{Dataset / benchmark} & \textbf{Year} & \textbf{Domain} & \textbf{Input} & \textbf{Primary use} \\
\midrule
KITTI Odometry~\cite{KITTI-GeigerLU12} & 2012 & Outdoor driving & Stereo + GPS/IMU & Classical and learned VO; segment drift evaluation. \\
TUM RGB-D~\cite{Sturm2012TUMRGBD} & 2012 & Indoor & RGB-D + mocap & Indoor RGB-D SLAM and trajectory accuracy. \\
7-Scenes~\cite{Scene_Coordinate} & 2013 & Indoor & RGB-D video & Camera relocalization and scene-coordinate regression. \\
Cambridge Landmarks~\cite{kendall2015posenet} & 2015 & Outdoor landmarks & RGB & Outdoor absolute pose regression and relocalization. \\
EuRoC MAV~\cite{Burri2016EuRoC} & 2016 & Aerial / industrial indoor & Stereo + IMU & Visual-inertial odometry under fast motion. \\
ScanNet~\cite{dai2017scannet} & 2017 & Indoor & RGB-D video + mesh & Learned pose, reconstruction, and dense geometry. \\
MegaDepth~\cite{li2018megadepth} & 2018 & Internet landmarks & SfM supervision & Relative pose, matching, depth, and geometric pretraining. \\
RealEstate10K~\cite{zhou2018stereo} & 2018 & Indoor / real-estate video & Video + recovered poses & Self-supervised video geometry and view synthesis. \\
Virtual KITTI 2~\cite{Cabon2020VirtualKITTI2} & 2020 & Synthetic driving & Rendered stereo / depth / pose & Controlled sim-to-real VO and depth evaluation. \\
DL3DV-10K~\cite{ling2024dl3dv} & 2024 & Large-scale scenes & Multi-view video / images & Generalizable 3D reconstruction and pose-aware training. \\
\bottomrule
\end{tabular}
\end{table}

\subsubsection{Datasets and Metrics}
\label{sec:vo_datasets}

Table~\ref{tab:vo_pose_datasets} shows three kinds of benchmark serving three different questions.
Trajectory benchmarks (KITTI Odometry~\cite{KITTI-GeigerLU12}, TUM RGB-D~\cite{Sturm2012TUMRGBD}, EuRoC MAV~\cite{Burri2016EuRoC}) measure drift under a fixed sensor setup and cover a narrow slice of deployment: urban driving, small rooms, micro-aerial-vehicle flights.
Relocalization benchmarks (7-Scenes~\cite{Scene_Coordinate}, Cambridge Landmarks~\cite{kendall2015posenet}) ask whether a query can be placed in a known scene and favor methods with scene-specific maps, retrieval, or scene-coordinate regression, so they are not comparable to online VO.
Large internet and video corpora (MegaDepth~\cite{li2018megadepth}, RealEstate10K~\cite{zhou2018stereo}, ScanNet~\cite{dai2017scannet}, DL3DV-10K~\cite{ling2024dl3dv}) serve mostly as training data for generalizable geometry models, which need diverse baselines, unknown intrinsics, unordered images, dynamic content, and dense reconstruction targets.

Evaluation is correspondingly layered.
With $\mathbf{T}_i$ and $\hat{\mathbf{T}}_i$ the ground-truth and estimated poses, the absolute trajectory error measures global discrepancy after a rigid or similarity alignment $\mathbf{S}$,
\begin{equation}
    \mathrm{ATE} =
    \sqrt{\frac{1}{N}\sum_{i=1}^{N}
    \left\|\mathrm{trans}\!\left(\mathbf{T}_i^{-1}\mathbf{S}\hat{\mathbf{T}}_i\right)\right\|^2},
\end{equation}
and the relative pose error measures local motion over an interval $\Delta$,
\begin{equation}
    \mathrm{RPE}(\Delta)=
    \frac{1}{N-\Delta}\sum_{i=1}^{N-\Delta}
    \left\|\mathrm{Log}\!\left[
    (\mathbf{T}_i^{-1}\mathbf{T}_{i+\Delta})^{-1}
    (\hat{\mathbf{T}}_i^{-1}\hat{\mathbf{T}}_{i+\Delta})
    \right]\right\|;
\end{equation}
KITTI reports the same idea as translational and rotational drift averaged over fixed-length path segments.
Whether scale is aligned before computing ATE, or scale drift is scored separately, decides whether monocular, stereo, RGB-D, and visual-inertial methods can be compared at all.
Pairwise matching and feed-forward pose methods report rotation error and translation-direction error in degrees, often aggregated as the area under the accuracy curve (AUC) below angular thresholds; relocalization reports the fraction of queries within thresholds such as $(5\,\mathrm{cm},5^\circ)$ indoors or meter-level outdoors, or median translation and rotation error; and feed-forward geometry models add depth error, point-map error, Chamfer distance, registration rate, or view-synthesis quality, since their cameras are judged by whether they support a usable 3D representation.

\begin{table*}[!t]
    \centering
    \caption{Representative methods under two shared protocols. \textit{Top}: monocular SLAM / VO, absolute trajectory error [m] averaged over the 11 EuRoC MAV sequences and the 9 TUM-RGBD freiburg1 sequences, as reported in DROID-SLAM~\cite{teed2021droid} (``fails $k$'' = lost tracking on $k$ sequences, so no average is defined). \textit{Bottom}: multi-view camera pose, AUC@30 on RealEstate10K (unseen) and CO3Dv2, as reported in VGGT~\cite{wang2025vggt}; the learned methods were trained on CO3Dv2 but not RealEstate10K. Numbers are from the cited evaluating paper and are comparable only within a block.}
    \label{tab:vo_pose_performance}
    \fontsize{9}{11}\selectfont
    \setlength{\tabcolsep}{1.5mm}
    \begin{tabular}{L{0.27\textwidth} L{0.23\textwidth} c C{0.14\textwidth} C{0.14\textwidth}}
    \toprule
    \textbf{Method} & \textbf{Family} & \textbf{Year} & \textbf{EuRoC ATE}$\downarrow$ & \textbf{TUM fr1 ATE}$\downarrow$ \\
    \midrule
    ORB-SLAM~\cite{mur2015orb} & Feature-based SLAM & 2015 & fails 2 / 11 & fails 6 / 9 \\
    DSO~\cite{engel2017direct} & Direct sparse odometry & 2017 & 0.601 & -- \\
    ORB-SLAM3~\cite{Campos2021ORBSLAM3} & Feature-based multi-map SLAM & 2021 & fails 1 / 11 & fails 5 / 9 \\
    DROID-SLAM~\cite{teed2021droid} & Differentiable dense BA & 2021 & \textbf{0.022} & \textbf{0.038} \\
    \midrule
    \hiderowcolors \textbf{Method} & \textbf{Family} & \textbf{Year} & \textbf{Re10K AUC@30}$\uparrow$ & \textbf{CO3Dv2 AUC@30}$\uparrow$ \\
    \midrule
    COLMAP + SuperPoint/SuperGlue~\cite{schonberger2016structure,detone2018superpoint,sarlin2020superglue} & Incremental SfM, learned matching & 2016 / 2020 & 45.2 & 25.3 \\
    DUSt3R~\cite{wang2024dust3r} & Pairwise point-map regression & 2024 & 67.7 & 76.7 \\
    MASt3R~\cite{mast3r} & 3D-grounded matching & 2024 & 76.4 & 81.8 \\
    VGGSfM v2~\cite{Vggsfm} & Learned differentiable SfM & 2024 & 78.9 & 83.4 \\
    CUT3R~\cite{wang2025Cut3R} & Continuous 3D perception with state & 2025 & 75.3 & 82.8 \\
    Fast3R~\cite{yang2025fast3r} & Large image-set feed-forward & 2025 & 72.7 & 82.5 \\
    VGGT~\cite{wang2025vggt} & Visual geometry grounded Transformer & 2025 & \textbf{85.3} (93.5 w/ BA) & \textbf{88.2} (91.8 w/ BA) \\
    \bottomrule
    \end{tabular}
\end{table*}

Table~\ref{tab:vo_pose_performance} puts the two regimes side by side. On trajectory benchmarks the lesson is robustness as much as accuracy: the classical systems that are highly accurate when they track lose tracking on a fraction of EuRoC flights and on most TUM freiburg1 sequences, whose rolling shutter, motion blur, and fast rotation defeat feature matching, whereas DROID-SLAM's learned dense correspondence inside a bundle-adjustment loop tracks every sequence at centimeter-level error. On multi-view pose the feed-forward models have overtaken the COLMAP pipeline by a wide margin and keep improving on unseen data (RealEstate10K), with VGGT ahead of DUSt3R-style pairwise regression at a fraction of the runtime, and a short bundle-adjustment refinement still adding several points, which is the hybrid pattern the challenges below return to. The two blocks are not comparable to each other, and that is the point: a method optimized for low drift under a calibrated rig and one optimized for zero-shot sparse-view geometry are judged by different questions.

\subsubsection{Challenges}

Geometric ambiguity sets the floor.
Monocular cameras cannot recover metric scale without priors; pure rotation, planar scenes, and very small baselines make epipolar geometry ill-conditioned; rolling shutter and unknown intrinsics violate the assumptions built into many solvers.
Learned models can hide these failures behind priors but cannot restore missing information, which is why visual-inertial fusion, calibration, and depth priors remain necessary even with large geometry models.
Deployment adds its own: dynamic agents corrupt the static-world assumption behind classical SLAM and bundle adjustment; textureless rooms, repetitive corridors, glass, water, and specular surfaces yield unreliable correspondences; appearance shift across time of day, weather, season, and camera hardware breaks place recognition; and long-horizon operation still needs loop closure, map maintenance, and uncertainty handling, none of which a feed-forward prediction supplies on its own.

The strongest trend is the convergence of SLAM, reconstruction, and foundation geometry. DROID-SLAM~\cite{teed2021droid} showed that optimization can be learned without abandoning geometric structure, and the feed-forward models above predict dense geometry, cameras, and correspondences together; the near-term frontier is hybrid, with feed-forward geometry supplying initialization, dense priors, and generalization and optimization enforcing metric consistency, loop closure, and uncertainty-aware correction.
Two further shifts matter for spatial intelligence. Pose estimation must become \textbf{dynamic-scene aware}, estimating camera and object motion jointly as MegaSaM~\cite{li2025megasam} and MonST3R~\cite{monst3r} begin to do, rather than discarding moving content as outliers; and camera pose should be judged by \textbf{downstream usefulness}, whether it supports stable depth and normal prediction, reconstruction, embodied navigation, AR anchoring, or world-model memory. As pose becomes one output of larger geometric perception systems, the boundary between VO, SLAM, scene reconstruction, and spatial memory will keep blurring.




\subsection{Pose Estimation}
\label{sec:object_pose}
% Author intent (Hao Fei): image to human-pose or 3D pose.
% Scope decision: §6.2.5 covers 6-DoF rigid-object pose;
% sparse landmark detection on humans, hands, and animals is owned by
% §6.1.1, and camera-frame pose is owned by §6.2.4.

6-DoF object pose estimation recovers a rigid transformation
\begin{equation}
\mathbf{T}_o =
\begin{bmatrix}
\mathbf{R} & \mathbf{t}\\
\mathbf{0}^{\top} & 1
\end{bmatrix}
\in SE(3),
\end{equation}
which maps an object's canonical or model frame into the current camera frame, giving a rotation $\mathbf{R}\in SO(3)$ and a translation $\mathbf{t}\in\mathbb{R}^3$ per detected instance.
It is the camera pose problem of \S\ref{sec:visual_odometry} with the reference frame reversed: the camera frame serves as the reference and each rigid object is located within it, using much of the same machinery (2D--3D correspondence, PnP solvers, RANSAC, and render-and-compare refinement) but with object-centric CAD models and pose annotations as supervision.
It is also distinct from the keypoint and landmark detection of \S\ref{sec:keypoint_landmark}: the target is the full 6-DoF pose of a rigid body, not semantic points on an articulated or deformable one.
Object pose is the geometric primitive consumed by robotic manipulation and grasping (revisited in Chapter~\ref{sec:embodied}), AR/VR object overlay, assembly guidance, bin picking, and the metric localization of other agents in driving.

Three axes organize the literature. The \emph{assumption axis} runs from instance-level methods that assume a high-fidelity CAD mesh per object, through category-level methods that assume only a canonical frame per category, to model-free methods that see a previously unseen object through a handful of reference images. The \emph{representation axis} separates direct $SE(3)$ regression from keypoint- or dense-correspondence pipelines that hand pose recovery to PnP, and from render-and-compare refinement of an initial pose. The \emph{foundation-model axis} tracks the shift from per-object training to cross-object generalization built on segmentation, vision, and depth foundation models.

\subsubsection{Task Definition}
The input is an RGB image, an RGB-D image, or a short multi-view burst of a scene containing one or more rigid objects, together with object-specific priors that depend on the setting: a CAD mesh or a small reference image set per object (instance-level and model-free), or a category label with a canonical frame (category-level), optionally with camera intrinsics and an upstream 2D detection or mask.
The output is a pose $\mathbf{T}_o\in SE(3)$ per detected object in the camera frame, augmented in the category-level setting with a metric size vector $\mathbf{s}\in\mathbb{R}_{+}^3$ (9-DoF pose and size).
For symmetric objects the prediction is properly an equivalence class of poses under the object's symmetry group, which the ADD-S and BOP metrics below account for.

\subsubsection{Methodology}

Figure~\ref{fig:p3_s6_2_5_pose_overview} schematizes the task and the three method-family trajectories, which recent methods often straddle.

\begin{figure}[!t]
    \centering
    \includegraphics[width=0.99\linewidth]{figures/part3/section6/pose_estimation.pdf}
    \caption{6-DoF object pose estimation. \textit{Top}: recover $\mathbf{T}_o\in SE(3)$ for each rigid object in an RGB(-D) image, given a CAD mesh, a canonical frame, or a few reference views. \textit{Bottom}: method timelines for (i) instance-level CAD-aware, (ii) category-level (NOCS) canonical-frame, and (iii) novel-object, foundation-era pose estimation.}
    \label{fig:p3_s6_2_5_pose_overview}
\end{figure}

\paragraph{Instance-Level 6-DoF: CAD-Aware Pipelines.}
With a mesh of each target object available, pose estimation becomes alignment of the mesh with image evidence, and synthetic renderings of the CAD models are the main training signal.
PoseCNN~\cite{xiang2017posecnn} set the modern baseline by regressing an object's image-space center and a quaternion rotation with a CNN, and introduced the YCB-Video benchmark.
A second family decouples the network from the geometric solver by predicting projections of sparse 3D keypoints and solving PnP with RANSAC; PVNet~\cite{peng2019pvnet} votes per pixel for each keypoint, which makes it robust to occlusion, while SSD-6D~\cite{kehl2017ssd6d} folds a discretized viewpoint classification into a detector and skips PnP entirely.
DenseFusion~\cite{Wang2019DenseFusion} fuses RGB-D features per pixel and directly regresses poses with iterative refinement. A third family uses dense surface correspondences instead of sparse keypoints: GDR-Net~\cite{wang2021gdrnet} couples dense coordinate regression with a geometry-guided PnP layer, ZebraPose~\cite{su2022zebrapose} encodes the surface with a coarse-to-fine binary code, and SurfEmb~\cite{haugaard2022surfemb} learns continuous correspondence distributions that stay robust under symmetry and occlusion; EPro-PnP~\cite{chen2022epro} replaces the non-differentiable PnP solver with a probabilistic surrogate so the whole pipeline trains end to end.
A fourth family refines by render-and-compare: DeepIM~\cite{li2018deepim} iteratively minimizes the residual between a rendered CAD view and the image, and CosyPose~\cite{Labbe2020CosyPose} extends this to consistent multi-view, multi-object refinement and led the BOP 2020 leaderboard.
The same machinery reaches far from the tabletop, for example spacecraft pose under extreme scale variation~\cite{hu2021widedepthrange}.

\paragraph{Category-Level 6-DoF + Size: Canonical Frame Recovery.}
Without a per-instance mesh but with a category that admits a canonical frame, pose is recovered up to that frame together with three-dimensional object size, a 9-DoF prediction.
NOCS~\cite{Wang0HVSG19} introduced the formulation, mapping every instance of a category into a Normalized Object Coordinate Space and regressing dense correspondences there; its REAL275 benchmark is the standard category-level evaluation.
Later work re-parameterizes the problem: FS-Net~\cite{chen2021fsnet} decouples rotation, translation, and size and speeds up inference through a shape-based pathway, GPV-Pose~\cite{di2022gpvpose} grounds rotation in geometry-guided point-wise voting, SAR-Net~\cite{lin2022sarnet} recovers a deformed shape prior alongside pose to handle intra-class variation, and IST-Net~\cite{liu2023istnet} drops the shape prior and learns the canonical transformation implicitly.
SecondPose~\cite{chen2024secondpose} and AG-Pose~\cite{lin2024agpose} import DINOv2 semantics and instance-adaptive keypoints for further gains on REAL275, and GenPose~\cite{zhang2023genpose} casts category-level pose as score-based diffusion sampling, the precursor of the 9-DoF diffusion model below.
The family's standing assumption, that a meaningful canonical frame exists per category, holds for mugs and cameras but fails for heterogeneous semantic groupings, which motivates the model-free line.

\paragraph{Novel-Object Pose: Model-Based and Model-Free Pipelines.}
The largest recent shift is away from per-object and per-category training toward objects first seen at test time, given only a few reference images or a CAD model.
From a short reference video, OnePose~\cite{sun2022onepose} triangulates a feature point cloud and matches test queries against it, OnePose++~\cite{he2022oneposepp} removes the dependence on repeatable keypoints with a dense-matching front end, and Gen6D~\cite{liu2022gen6d} unifies detection, viewpoint selection, and refinement without per-object fine-tuning.
From a CAD model alone, MegaPose~\cite{labbe2022megapose} trains render-and-compare on millions of synthetic objects and transfers zero-shot, FoundationPose~\cite{wen2024foundationpose} consolidates model-based and model-free operation under large-scale synthetic pretraining and leads the BOP novel-object track, GigaPose~\cite{nguyen2024gigapose} cuts the cost of render-and-compare by sampling templates in a two-DoF rotation space and recovering the remaining parameters from a single patch correspondence, and SAM-6D~\cite{lin2024sam6d} uses the Segment Anything Model~\cite{kirillov2023sam} as a zero-shot detection front end before learned pose matching.
A further sub-line drops pose-specific training altogether and matches frozen foundation features: FoundPose~\cite{ornek2024foundpose} establishes 2D--3D correspondences between DINOv2 patch descriptors of the query and pre-rendered templates, and FreeZe~\cite{caraffa2024freeze} fuses frozen 3D geometric descriptors with 2D foundation features and recovers pose by RANSAC registration; a FreeZe variant took first place for unseen-object localization and detection in the BOP Challenge 2024~\cite{nguyen2024bopchallenge}, overtaking the render-and-compare line.
Diff9D~\cite{liu2024diff9d} reformulates 9-DoF category-level pose as diffusion sampling, the object-pose instance of the diffusion-as-perception pattern seen for depth and normals, which yields calibrated multi-hypothesis predictions (natural for symmetric objects) at the cost of latency.
The foundation-model line has thus largely dissolved the instance-versus-category-versus-model-free distinction; the active question is how cheap the per-test-object supervision can become.

\paragraph{Sensing Modalities.}
Across these axes, the input modality fixes which supervisory signals exist and which symmetries can be resolved. RGB-only pipelines~\cite{xiang2017posecnn,peng2019pvnet,wang2021gdrnet,chen2022epro,sun2022onepose} suit mobile and AR deployment but face scale--depth ambiguity on texture-poor objects; RGB-D fusion (DenseFusion~\cite{Wang2019DenseFusion}, CosyPose~\cite{Labbe2020CosyPose} with depth refinement, FoundationPose~\cite{wen2024foundationpose}) is substantially more accurate on BOP at the cost of a depth sensor; depth-only methods are rare but remain competitive on highly symmetric objects where shading carries little rotational evidence. Monocular depth foundation models (\S\ref{sec:depth_estimation}) are increasingly used as plug-in depth providers at inference, which narrows the RGB / RGB-D gap without an active sensor.

\subsubsection{Datasets and Metrics}
\label{sec:pose_datasets}

Table~\ref{tab:pose_datasets} lists the recurring benchmarks. Instance-level evaluation grew from LINEMOD~\cite{hinterstoisser2012linemod}, the largely texture-poor early benchmark that introduced the ADD metric, through its occlusion-heavy re-annotation LINEMOD-Occluded~\cite{brachmann2014lmo}, which modern methods report as the more informative of the two, to the video setting of YCB-Video~\cite{xiang2017posecnn} and the texture-less, symmetric industrial parts of T-LESS~\cite{Hodan2017TLESS}, where symmetry handling decides the ranking; HomebrewedDB~\cite{kaskman2019hbdb} serves as an out-of-distribution test for methods trained on the BOP core data, and SPEED~\cite{Sharma2019SpacecraftPE} stretches the task to spacecraft under scale and illumination ranges no tabletop set covers. NOCS-REAL275 with its synthetic companion CAMERA~\cite{Wang0HVSG19} is the category-level standard. The BOP Challenge~\cite{hodan2018bop,hodan2020bop,sundermeyer2023bop} consolidates the instance-level datasets under one protocol and has been the field's headline leaderboard since 2017; its 2023 edition distinguishes seen-object and unseen-object settings, with 3D mesh models provided in both. Modern papers report at minimum on LINEMOD-Occluded, YCB-Video, and a BOP-core subset, novel-object methods add the BOP unseen-object track, and category-level methods center on REAL275.

\begin{table*}[!t]
    \centering
    \caption{Representative datasets and benchmarks for 6-DoF object pose estimation. ``--'' marks a value not fixed by the source or varying across protocol versions.}
    \label{tab:pose_datasets}
    \fontsize{9}{11}\selectfont
    \setlength{\tabcolsep}{1.2mm}
    \begin{tabular}{L{0.20\textwidth} c L{0.23\textwidth} C{0.09\textwidth} C{0.13\textwidth} L{0.19\textwidth}}
    \toprule
    \textbf{Dataset} & \textbf{Year} & \textbf{Setting} & \textbf{Objects} & \textbf{Frames} & \textbf{Role} \\
    \midrule
    LINEMOD~\cite{hinterstoisser2012linemod} & 2012 & Instance, texture-poor objects, RGB-D & 15 & $\sim$1.2K / object & Entry benchmark; introduced ADD \\
    LINEMOD-Occluded~\cite{brachmann2014lmo} & 2014 & Instance, heavy mutual occlusion & 8 & re-annotated LM scene & Occlusion stress test \\
    YCB-Video~\cite{xiang2017posecnn} & 2017 & Instance, RGB-D video, moving camera & 21 & $\sim$133K & Sequential pose and tracking \\
    T-LESS~\cite{Hodan2017TLESS} & 2017 & Instance, texture-less, symmetric industrial parts & 30 & -- & Symmetry / texture-poor stress test \\
    HomebrewedDB~\cite{kaskman2019hbdb} & 2019 & Instance, household scenes & -- & -- & Out-of-distribution generalization \\
    NOCS REAL + CAMERA~\cite{Wang0HVSG19} & 2019 & Category-level 9-DoF, RGB-D & 6 categories & $\sim$8K real (+ synthetic) & Category-level standard \\
    SPEED~\cite{Sharma2019SpacecraftPE} & 2019 & Spacecraft, extreme scale and lighting & 1 & -- & Beyond-tabletop deployment \\
    BOP Challenge~\cite{hodan2018bop,hodan2020bop,sundermeyer2023bop} & 2018-- & Unified protocol over LM-O, T-LESS, YCB-V, ITODD, IC-BIN, TUD-L, HB, HOPE & -- & -- & Headline leaderboard; novel-object track since 2023 \\
    \bottomrule
    \end{tabular}
\end{table*}

Each metric answers to a particular failure mode of pose error, and the symmetry-aware variants exist because many objects admit several valid poses.
The Average Distance of Model Points (ADD) is the standard for asymmetric objects: with predicted pose $(\hat{\mathbf{R}}, \hat{\mathbf{t}})$, ground truth $(\mathbf{R}, \mathbf{t})$, and mesh vertices $\mathcal{M}$,
\begin{equation}
    \mathrm{ADD} = \frac{1}{|\mathcal{M}|} \sum_{\mathbf{x} \in \mathcal{M}}
    \left\| (\mathbf{R}\mathbf{x} + \mathbf{t}) - (\hat{\mathbf{R}}\mathbf{x} + \hat{\mathbf{t}}) \right\|,
\end{equation}
and a pose counts as correct when $\mathrm{ADD} < 0.1\,d$ for object diameter $d$, the ADD-0.1d threshold introduced with LINEMOD~\cite{hinterstoisser2012linemod}.
For symmetric objects ADD-S replaces direct vertex correspondence with nearest-neighbor matching,
\begin{equation}
    \mathrm{ADD\text{-}S} = \frac{1}{|\mathcal{M}|} \sum_{\mathbf{x}_1 \in \mathcal{M}}
    \min_{\mathbf{x}_2 \in \mathcal{M}}
    \left\| (\mathbf{R}\mathbf{x}_1 + \mathbf{t}) - (\hat{\mathbf{R}}\mathbf{x}_2 + \hat{\mathbf{t}}) \right\|,
\end{equation}
under the same threshold; the two are reported together so that symmetric and asymmetric objects are diagnosed separately.
For AR overlay the relevant error is in image space, so the 2D projection error averages the pixel distance between mesh vertices projected under the two poses, with 5\,px as the usual threshold.
Category-level evaluation on REAL275 counts a prediction correct when
\begin{equation}
    \angle(\mathbf{R}^{-1} \hat{\mathbf{R}}) < 5^\circ
    \quad \text{and} \quad
    \| \mathbf{t} - \hat{\mathbf{t}} \| < 5\,\mathrm{cm},
\end{equation}
and adds 3D box IoU at $50\%$ and $75\%$ to score pose and size jointly.
The BOP Challenge~\cite{hodan2018bop} averages recall over three symmetry-aware errors, each thresholded at several operating points: the Visible Surface Discrepancy (VSD), which scores only the visible object surface and is therefore robust to occlusion, the Maximum Symmetry-aware Surface Distance (MSSD), the maximum surface-point deviation after selecting the best valid symmetry alignment, and its image-space analogue MSPD,
\begin{equation}
    \mathrm{AR} = \tfrac{1}{3} \bigl( \mathrm{AR}_{\mathrm{VSD}} + \mathrm{AR}_{\mathrm{MSSD}} + \mathrm{AR}_{\mathrm{MSPD}} \bigr),
\end{equation}
averaged over the BOP core datasets.
ADD and ADD-S remain the default for LINEMOD-style instance-level work, BOP-AR for cross-dataset comparison, and projection error for AR; the metrics persist side by side because no single norm captures both strict pose error and the equivalence-class structure that VSD and MSSD make explicit.

\begin{table}[!t]
    \centering
    \caption{Representative 6-DoF object pose results on LINEMOD-Occluded (higher is better). ADD(-S)-0.1d anchors the RGB methods; BOP-AR is reported for the refinement and foundation-era systems. The two columns are different metrics and are not directly comparable.}
    \label{tab:pose_benchmark}
    \fontsize{9}{11}\selectfont
    \setlength{\tabcolsep}{4pt}
    \begin{tabular}{lllllcc}
        \toprule
        \textbf{Method} & \textbf{Family} & \textbf{Modality} & \textbf{Train assumption} & \textbf{Year} & \textbf{ADD(-S)-0.1d} $\uparrow$ & \textbf{BOP-AR} $\uparrow$ \\
        \midrule
        PoseCNN~\cite{xiang2017posecnn}             & Instance / direct         & RGB   & per-object   & 2017 & 24.9 & --   \\
        PVNet~\cite{peng2019pvnet}                  & Instance / keypoint+PnP   & RGB   & per-object   & 2019 & 40.8 & --   \\
        GDR-Net~\cite{wang2021gdrnet}               & Instance / coord. reg.    & RGB   & per-object   & 2021 & 62.2 & --   \\
        % 需核查: 63.3 is CosyPose's RGB single-view LM-O result (BOP 2020); the RGB-D refined variant scores higher.
        CosyPose~\cite{Labbe2020CosyPose}           & Instance / refinement     & RGB   & per-object   & 2020 & --   & 63.3 \\
        SurfEmb~\cite{haugaard2022surfemb}          & Instance / dense corr.    & RGB-D & per-object   & 2022 & --   & 75.8 \\
        % 需核查: MegaPose LM-O BOP-AR value and modality to be confirmed against the BOP leaderboard.
        MegaPose~\cite{labbe2022megapose}           & Novel-object / render\&comp. & RGB-D & cross-object & 2022 & --   & 58.3 \\
        FoundationPose~\cite{wen2024foundationpose} & Novel-object / found.       & RGB-D & cross-object & 2024 & --   & 78.8 \\
        \bottomrule
    \end{tabular}
\end{table}

Table~\ref{tab:pose_benchmark} reports LINEMOD-Occluded, the most widely used instance-level benchmark under the BOP protocol.
Among RGB methods scored by ADD(-S)-0.1d, explicit geometry steadily closes the gap over direct regression: PoseCNN's 24.9 rises to 40.8 with per-pixel keypoint voting (PVNet) and to 62.2 with geometry-guided dense coordinate regression (GDR-Net).
On the symmetry-aware BOP-AR, the render-and-compare CosyPose (63.3) and the dense-correspondence SurfEmb (75.8) are now matched or exceeded by novel-object systems without per-object training, FoundationPose reaching 78.8 against MegaPose's earlier 58.3; this is the 2023--2024 shift that made novel-object pose a deployment target rather than a research curiosity.
The gap between the two families is widest on the symmetric objects of T-LESS, where losses that model symmetry (BOP-AR-aware objectives, ADD-S optimization) generalize better than those assuming a single asymmetric pose.

\subsubsection{Challenges}

\textbf{The onboarding bottleneck has shifted from pose to detection.}
Moving to novel objects relocated the limiting stage rather than removing it.
On the BOP Challenge 2024~\cite{nguyen2024bopchallenge}, 6-DoF \emph{localization} of unseen objects has nearly caught up with the seen-object regime (the strongest training-free entry, a FreeZe variant~\cite{caraffa2024freeze}, trails the best seen-object method by a few points of average recall), whereas \emph{detection} of unseen objects still lags its seen-object counterpart by roughly a third and is now the largest error source in the novel-object stack.
The open problem is zero-shot detection and segmentation of a never-seen instance in clutter, together with the cost of \emph{onboarding} it: BOP 2024 added a track in which an object is registered from a short reference video, Any6D~\cite{lee2025any6d} registers from a single RGB-D anchor image, and how few views, or whether a single image or a text prompt, suffice is the defining frontier question. The same challenge pivoted toward high-resolution AR/VR headset recordings, a signal that the benchmark-realism gap, including articulated and deformable objects, open-vocabulary targets named rather than modeled (POPE~\cite{fan2023pope}), and egocentric capture, is now a first-order concern.

\textbf{Absolute metric scale remains ambiguous for model-free targets.}
The true size of an unseen object is under-determined from RGB alone, since scale and depth trade off along the viewing ray.
Category-level 9-DoF methods absorb this by regressing size in a canonical frame~\cite{Wang0HVSG19}, and model-free systems can use a depth sensor or another metric reference, while CAD-based methods can use the model's known extent; without such information, scale is the residual unknown separating a plausible pose from a usable one, which Any6D~\cite{lee2025any6d} makes explicit by coupling pose recovery with scale alignment, and monocular depth priors inherit the affine ambiguity of \S\ref{sec:depth_estimation}.

\textbf{Accuracy trades sharply against latency.}
The methods that top the novel-object leaderboard are the least suited to closed-loop use: on BOP 2024~\cite{nguyen2024bopchallenge} the most accurate unseen-object method runs in tens of seconds per image, close to an order of magnitude slower than both the fastest seen-object pipelines and the practical sub-second entries, which buy speed with accuracy.
Heavy segmentation, render-and-compare, and diffusion priors (Diff9D~\cite{liu2024diff9d}) are the cause, which is why one-correspondence template matching (GigaPose~\cite{nguyen2024gigapose}), distillation, and selective invocation of expensive components have become load-bearing for robotic and AR loops at frame rate.

\textbf{Symmetry and geometric ambiguity are irreducible.}
Symmetric and texture-less parts of the kind that fill T-LESS~\cite{Hodan2017TLESS} admit a continuous or discrete family of photometrically indistinguishable poses, and any loss assuming a unique target diverges on them; transparent and specular surfaces violate the brightness-constancy and depth-sensing assumptions behind correspondence and RGB-D fusion; rotation parameterization adds degeneracies that continuous 6D or axis-angle encodings reduce but do not remove; and heavy occlusion, the regime of LINEMOD-Occluded~\cite{brachmann2014lmo}, starves the estimator of the evidence needed to resolve any of them.
Generative multi-hypothesis pose (Diff9D~\cite{liu2024diff9d}) predicts a distribution instead of a point, but turning that distribution into a \emph{calibrated} confidence a planner can trust in safety-critical manipulation and driving is still largely unaddressed.

Across these problems, pose pipelines increasingly consume foundation-model priors as plug-in modules (promptable segmentation for zero-shot detection, monocular depth for inference-time geometry) rather than training pose-specific backbones, and the boundary between pose estimation, 6-DoF tracking, and neural reconstruction is dissolving, as FoundationPose~\cite{wen2024foundationpose} folds estimation and tracking into one loop and the reconstruction discussion of Chapter~\ref{sec:generation} pursues further. Promptable, real-time, uncertainty-aware pose is the target the next generation of methods converges on.



% \newpage

\section{Spatial Structure Construction}
\label{sec:p3-sec6-3-spatial-structure-construction}

\FramedInitial{W}{ith} entities recognized and their positions fixed, what remains is how they hang together. Spatial structure construction turns independently localized entities into an organized configuration: which objects relate to which, how those relations compose into a scene, and how the scene as a whole is laid out. The section builds outward in scope, from pairwise spatial relations, through the scene graphs that assemble them, to room-level layout; it then lets structure take on time and scale, as the event structure of action detection and the whole-environment structure of mapping and SLAM. Structure in real perception is seldom purely static, and what the section yields is the explicit spatial configuration that downstream reasoning and planning act on.






\subsection{Spatial Relation Extraction}
\label{sec:spatial_relation_extraction}

% Task definition
%  - Motivation & Importance
%  - Formal Task Definition
%  - Subtasks, if have:

% Part 2: Method taxonomy.
% modality
%

% Task Definition
%  - Motivation & Importance
% ,liu2024edge，liu2025boost
Spatial Relation Extraction (SRE) identifies the relative spatial configuration of objects, such as ``table-left-of-chair'' or ``car-behind-bus''.
It turns detections or segmentations into spatial relation knowledge that feeds reasoning, scene graph generation~\cite{xu2025llava,li2024pixels,wu2025universal}, autonomous driving~\cite{malawade2022roadscene2vec,zhou2023pit,liu2024edge,liu2025boost}, and embodied AI~\cite{SongBTTSB25,zhang2025vlabench}.
It is narrower than Scene Graph Generation (SGG), which runs the full pipeline from detection to relation prediction and includes non-spatial predicates (``person-holding-book''); here SRE covers only geometric, topological, and projective spatial predicates.


%  - Formal Task Definition
\subsubsection{Task Definition}
Given a visual scene (a 2D image, a video, a point cloud, or an RGB-D stream) and two target object instances specified by tags, boxes, or masks, SRE outputs their spatial relationship either as a triplet $(s, p, o)$ with $p$ drawn from a spatial predicate vocabulary, or as a natural-language sentence describing the relation and its context.
Three subtasks differ in output granularity. \textbf{Visual Spatial Classification (VSC)}~\cite{yang2019spatialsense, zhao2024acm} predicts a predicate from a closed label set (``on'', ``in'', ``next to'') or verifies whether a given predicate holds for the object pair. \textbf{Visual Spatial Description (VSD)}~\cite{zhao2022visual, zhao2023generating} generates a free-form sentence that can carry attributes and context beyond the bare predicate. \textbf{Quantitative spatial understanding}~\cite{cheng2024spatialrgpt,chen2024spatialvlm} replaces qualitative predicates with metric quantities (distances, absolute sizes, 3D orientations), which navigation and manipulation require. The field has moved through these three in order, from closed-set classification to generation to metric quantification.


% Part 2: Method taxonomy.
% (Focus on giving a comprehensive but structured overview; no need to cover every paper in depth.) 规则+学习
\subsubsection{Methodology}

Figure~\ref{fig:p3_s6_3_1_sre_overview} summarizes the task and the three method families surveyed below.

\begin{figure}[!t]
    \centering
    \includegraphics[width=\linewidth]{figures/part3/section6/spatial_relation_extraction.pdf}
    \caption{Spatial relation extraction. \textit{Top}: given a visual scene (image, video, point cloud, or RGB-D) and two target objects $(O_1, O_2)$, predict (a) a closed-set relation label (VSC), (b) a free-form description (VSD), or (c) a quantitative metric estimate. \textit{Bottom}: (i) rule-based, (ii) learning-based decoupled, and (iii) MLLM-based mechanisms with method timelines (chronological spacing is not proportional). Metric outputs require depth or scale information.}
    \label{fig:p3_s6_3_1_sre_overview}
\end{figure}

\paragraph{Rule-based approaches.}
Rule-based pipelines~\cite{malawade2022roadscene2vec,zhao2022visual,zhou2024hktsg} first run off-the-shelf perception (a 2D detector with a bird's-eye-view transform~\cite{liu2008bird}, or a 3D scene extractor~\cite{nie2020total3dunderstanding}) to obtain each object's 3D position, size, orientation, and appearance in a camera-centered frame whose vertical axis is aligned with gravity, and then apply deterministic geometric rules to the centroids to produce $(O_1, p, O_2)$ triplets~\cite{zhao2022visual}.
Let $(x_{O_i}, y_{O_i}, z_{O_i})$ denote the centroid of object $O_i$ in this shared camera-centered frame, with $x$ pointing right and $y$ pointing upward. Then
\begin{equation}
\begin{aligned}
\Delta x &= |x_{O_1} - x_{O_2}|, \quad
\Delta y = |y_{O_1} - y_{O_2}|, \quad
\Delta z = |z_{O_1} - z_{O_2}|, \\
&\Delta y > \tau \,\land\, y_{O_1} > y_{O_2} \;\Rightarrow\; \text{``above''}, \\
&\Delta x > \Delta z \,\land\, \Delta x > \tau \,\land\, x_{O_1} < x_{O_2} \;\Rightarrow\; \text{``left''}, \\
&\Delta x > \Delta z \,\land\, \Delta x > \tau \,\land\, x_{O_1} > x_{O_2} \;\Rightarrow\; \text{``right''},
\end{aligned}
\end{equation}
where $\tau$ is a small threshold (\eg, $0.2$\,m).
Such rules are cheap but brittle: 2D-to-3D back-projection misplaces occluded or oddly posed objects (a rear vehicle whose centroid lands ambiguously on the BEV plane), fixed rules cannot express functional or contextual relations (``person riding bicycle'', a book more likely on a shelf than under it), and every new domain needs its rules rewritten.


\paragraph{Learning-based approaches.}
Zhao \etal~\cite{zhao2022visual} decoupled spatial relation generation from generic captioning and trained it with a joint visual--spatial objective; 3DVSD~\cite{zhao2023generating} moved the same task into 3D scene modeling and added structure-aware sampling that holds up under occlusion.

\paragraph{MLLM-based approaches.}
Multimodal large language models carried the line from image-level predicates toward finer and more metric grounding.
SpatialVLM~\cite{chen2024spatialvlm} used monocular-depth-derived spatial QA supervision to train a vision-language model so the model answers metric queries such as ``How far is A from B?'', outperforming GPT-4V on spatial VQA.
SpatialBot~\cite{cai2024spatialbot} exposes pixel-level depth queries through a depth API and trains progressively, translating spatial reasoning into higher robotic grasping success.
SpatialRGPT~\cite{cheng2024spatialrgpt} sharpens granularity from the whole image to regions, accepting boxes or masks and auto-generating 8.7M region-level spatial samples through scene-subgraph sampling (S$^3$).
SpatialLLM~\cite{ma2025spatialllm} pairs a hybrid CLIP--DINOv2 encoder with staged data injection to hold qualitative and quantitative capabilities in one model.


% Part 3: Benchmarks & metrics.
\subsubsection{Datasets and Metrics}
Table~\ref{tab:sre_dataset} lists the benchmarks. They differ mainly in how relations were collected: manual labeling of robot-captured lab scenes in SpatialAwareRobot~\cite{wang2025spatial}; adversarial crowdsourcing in SpatialSense~\cite{yang2019spatialsense}, where annotators propose relations that fool a baseline and only those count, which makes the dataset deliberately hard; the merge of SpatialSense and Visual Genome~\cite{krishna2017visual} into VSD~\cite{zhao2022visual}, whose VSD-v1 version contains over 143K descriptions; the 34-preposition, English--French annotations of SpatialVOC2K~\cite{belz2018spatialvoc2k}; and the automatically generated, depth-aware SpaceSGG~\cite{xu2025llava}, released as 10K scene descriptions, 20K question--answer pairs, and 20K multi-turn conversations over 20K scenes.

\begin{table}[!t]
\centering
\caption{Benchmark datasets for spatial relation extraction. ``--'' marks an open vocabulary without a fixed category count.}
\label{tab:sre_dataset}
\fontsize{9}{11}\selectfont
\setlength{\tabcolsep}{3mm}
\begin{tabular}{lccccc}
\toprule
\textbf{Dataset} & \textbf{Scene Type} & \textbf{Modality} & \textbf{Annotation} & \textbf{Samples} & \textbf{Relation Categories} \\
\midrule
SpatialAwareRobot~\cite{wang2025spatial} & Robotics & Image & Manual & 1,000 & 7 \\
SpatialSense~\cite{yang2019spatialsense} & Generic & Image & Manual & 11,569 & 9 \\
VSD~\cite{zhao2022visual} & Generic & Image & Manual & 29,272  & -- \\
SpatialVOC2K~\cite{belz2018spatialvoc2k} & Generic & Image & Manual & 2,026  & 34 \\
SpaceSGG~\cite{xu2025llava} & Generic & Image & Automated & 20,000  & --  \\
\bottomrule
\end{tabular}
\end{table}

Each subtask has its own evaluation regime, and the metrics are not comparable across them.
VSC is scored as classification: SpatialSense~\cite{yang2019spatialsense} measures binary accuracy for given object--predicate pairs, with the driving-scene systems Roadscene2Vec~\cite{malawade2022roadscene2vec} and HKT-SG~\cite{zhou2024hktsg} reporting per-relation F1 or accuracy on their own data; results on SpatialSense's adversarial split lag well behind those on conventional crowdsourced data, which is the contextual-relation difficulty taken up under Challenges.
VSD is scored as generation: VSDNet~\cite{zhao2022visual} and 3DVSD~\cite{zhao2023generating} report BLEU-4, METEOR, and CIDEr on the VSD datasets, and LLaVA-SpaceSGG~\cite{xu2025llava} reports QA accuracy on SpaceSGG.
Quantitative understanding is the newest regime: SpatialVLM~\cite{chen2024spatialvlm} and SpatialRGPT~\cite{cheng2024spatialrgpt} evaluate metric estimates using human-annotated spatial questions and answers (SpatialVLM) or questions derived from ground-truth 3D cuboid annotations (SpatialRGPT-Bench), SpatialLLM~\cite{ma2025spatialllm} reports combined qualitative and quantitative VQA accuracy, and SpatialBot~\cite{cai2024spatialbot} adds robotic grasp success as an embodied proxy.
MLLM-based models dominate the qualitative leaderboards but still fall short on metric-scale reasoning, which motivates the unified representations discussed next.


\subsubsection{Challenges}
Spatial relations are not purely geometric. ``A cup on a table'' is plausible and ``a cup on a wall'' is not, yet the two differ only in object identity, so rule-based and purely geometric methods cannot capture functional or commonsense constraints; a model needs to know what objects are for and where they tend to be, not just their boxes or point coordinates.
Metric grounding is the second gap. MLLMs handle contextual and functional predicates well but judge distances, sizes, and angles poorly; injecting geometry and depth during pre- or post-training, or exposing dedicated spatial tools, is the current route, and the longer-term goal is a representation that holds discrete predicates (``near'') and continuous quantities (``2 meters away'') together, since navigation and manipulation need both.
Generalization is the third. SpatialSense~\cite{yang2019spatialsense} and VSD~\cite{zhao2022visual} cover a small slice of real scenes, and models trained on them degrade under new object appearances, scene complexity, lighting, and occlusion; SpatialVLM~\cite{chen2024spatialvlm} and its successors generalize better, but unseen and dynamic environments remain largely untested.





\subsection{Scene Graph Generation}
\label{sec:scene_graph_generation}
% - `Writing task`: 任务定义（TSG、VSG、PSG、HOI / HRI (egocentric)） → 表达与评测 → 方法类别 → 数据与协议 → 难点与趋势
% - Content:
%- Table:  SGG datasets
% - Figure:  三种SG示意图
% https://github.com/ChocoWu/Awesome-Scene-Graph-Generation


Scene graph generation (SGG)~\cite{xu2017scene,chang2021comprehensive,li2024scene} parses visual content into a graph of triplets, each joining a subject and an object through a predicate, for example $(\text{person}, \text{rides}, \text{horse})$.
The graph records not only which objects are present but how they interact, which is what retrieval~\cite{johnson2015image,miao2024scenegraphloc}, captioning~\cite{zhong2020comprehensive,lu2025benchmarking}, autonomous driving~\cite{zhou2022toward,liu2023learning,malawade2022roadscene2vec}, and robotics~\cite{nguyen2025generating,werby2024hierarchical,garg2024robohop} consume downstream.


%  - Formal Task Definition
\subsubsection{Task Definition}
Given a visual scene, SGG produces a graph
\begin{equation}
\mathcal{G} = \{(s_i, p_i, o_i)\}_{i=1}^{M}, \qquad s_i, o_i \in \mathcal{O},\; p_i \in \mathcal{P},
\end{equation}
where each entity in $\mathcal{O}$ is an object instance with a localization (a bounding box, or a mask in the panoptic setting) and a class label, and $p_i$ is a predicate from a relation vocabulary $\mathcal{P}$ describing how $s_i$ and $o_i$ interact.

% 任务分类：
% 1. Textual Scene Graph Generation
% 2. 2D Image Scene Graph Generation
% 2.1  Panoptic Scene Graph Generation
% 3. Spatio-Temporal (Video) Scene Graph Generation
% 4. 3D/ Scene Graph Generation
% 5. Universal Scene Graph Generation
% 6. HOI
% 7. Ego-SGG

The task is instantiated by input modality. \textbf{Image SGG (ISG)}~\cite{krishna2017visual,zellers2018neural} on static 2D images is the original and most studied setting and anchors most benchmarks. \textbf{Video SGG (VSG)}~\cite{teng2021target, cong2021spatial} adds time, so relations become dynamic and temporally extended, and temporal context disambiguates interactions a single frame leaves open. \textbf{Textual SGG (TSG)}~\cite{li2023factual,choi2022scene} parses a sentence (``A person is riding a horse in a field'') into a graph, treating SGG as structured language parsing. \textbf{3D SGG (3DSG)}~\cite{wald2020learning,wu2021scenegraphfusion} operates on point clouds or meshes and models relations in metric 3D space, free of the perspective ambiguity of images. \textbf{Universal SGG (USG)}~\cite{wu2025universal} aims at one graph from several simultaneous modalities (image, video, 3D, text).
Two further variants change the level of detail rather than the modality. \textbf{Human--Object Interaction detection (HOI)}~\cite{li2019transferable, wang2020learning} restricts the subject to humans and has become a field of its own because of pose variability and the subtlety of intent. \textbf{Panoptic SGG (PSG)}~\cite{yang2022panoptic,yang2023panoptic} grounds nodes on panoptic segmentation masks instead of boxes, so relations can involve ``stuff'' (sky, road) and are localized without the imprecision of boxes.

\subsubsection{Methodology}
% 1.按模态分（Image、video、textual、3D、Universal Scene Graph Generation）
Figure~\ref{fig:p3_s6_3_2_sgg_overview} illustrates the five modality settings and the two specializations with a representative input and output each.

\begin{figure}[!t]
    \centering
    \includegraphics[width=\linewidth]{figures/part3/section6/scene_graph.pdf}
    \caption{Scene graph generation by input modality (top) and sub-task setting (bottom). \textit{Top}: image, video, text, 3D, and universal (multimodal) SGG, each with a sample input and output graph. \textit{Bottom}: (i) human--object interaction (HOI), restricted to human subjects at box level, and (ii) panoptic SGG (PSG), grounding ``things'' and ``stuff'' on pixel masks.}
    \label{fig:p3_s6_3_2_sgg_overview}
\end{figure}

% 2.two-stage —> one-stage
\paragraph{Two-stage methods.}
Two-stage methods~\cite{xu2017scene,zellers2018neural,li2021bipartite,lin2020gps,li2023compositional,zhou2024hierarchical} were the standard for years: a pre-trained detector such as Faster R-CNN~\cite{ren2015faster} produces boxes and per-object features, and a separate relation model predicts a predicate for each pair. Specializing each stage made the graphs accurate, but detector errors propagate uncorrected into relation prediction, the two models are trained and tuned separately, and object interactions are only modeled after detection has already committed.

\paragraph{One-stage methods.}
One-stage methods~\cite{im2024egtr,wang2024oed,newell2017pixels,liu2021fully} fold detection and relation prediction into one end-to-end network. Pixel-based~\cite{newell2017pixels} and fully convolutional~\cite{liu2021fully} formulations came first; after DETR~\cite{carion2020end}, most one-stage SGG adopted set prediction with learned queries, and the methods differ in what the queries represent. RelTR~\cite{cong2023reltr} uses paired subject--object queries so the Transformer predicts relations between pairs directly; SGTR~\cite{li2022sgtr} decouples subject, object, and predicate queries and composes them, which lets each component be predicted on its own; and PSGTR~\cite{yang2022panoptic} uses triplet queries that emit scene-graph triplets without a detector at all.

% 3.引入先验知识，LLM
\paragraph{Prior knowledge.}
Because predicates are long-tailed and many relations are not visually decidable, SGG has long imported priors. Language priors~\cite{liang2017deep,lu2016visual,cui2018context} supply semantic embeddings that help rare predicates; statistical priors~\cite{yu2017visual,zhou2023learning,zhou2024hierarchical} model object-pair co-occurrence; knowledge graphs such as ConceptNet~\cite{speer2017conceptnet} add commonsense~\cite{zhang2024hiker,chen2023more}; and LLMs bring world knowledge~\cite{chen2024scene,kim2024llm4sgg}.

\paragraph{Beyond images.}
The non-image settings reuse these designs and add modality-specific machinery. Video SGG tracks object tubes and aggregates relations over time: TRACE~\cite{teng2021target} decomposes relation prediction by temporal segment and STTran~\cite{cong2021spatial} applies a spatio-temporal Transformer over frame-level graphs. Textual SGG is a parsing problem, for which the FACTUAL parser~\cite{li2023factual} fine-tunes a sequence-to-sequence model to emit faithful triples. 3D SGG predicts relations from point-cloud or mesh segments with graph networks~\cite{wald2020learning}, and SceneGraphFusion~\cite{wu2021scenegraphfusion} does so incrementally from an RGB-D stream so the graph grows with the reconstruction. Panoptic SGG is served by the triplet-query PSGTR~\cite{yang2022panoptic} and its video extension~\cite{yang2023panoptic}, and USG~\cite{wu2025universal} shares one graph decoder across image, video, 3D, and text inputs.

% TODO add HOI，Traffic SG dataset.

\subsubsection{Datasets and Metrics}


\begin{table}[!t]
  \centering
\caption{Benchmark datasets for scene graph generation.}
\label{tab:SGG_data}
\fontsize{8}{10}\selectfont
\setlength{\tabcolsep}{2.5mm}
  \begin{tabular}{lcccccc}
    \toprule
    \textbf{Dataset} & \textbf{Modality} &\textbf{ Obj. Class} & \textbf{BBox} & \textbf{Rela. Class} & \textbf{Triplets} & \textbf{Instances} \\
    \midrule
    Visual Phrase~\cite{Sadeghi2011VisualPhrases} & Image & 8 & 3,271 & 9 & 1,796 & 2,769 \\
    Scene Graph~\cite{johnson2015image} & Image & 266 & 69,009 & 68 & 109,535 & 5,000 \\
    VRD~\cite{lu2016visual} & Image & 100 & - & 70 & 37,993 & 5,000 \\
    Open Images v7~\cite{kuznetsova2020open} & Image & 600 & 3,290,070 & 31 & 374,768 & 9,178,275 \\
    Visual Genome~\cite{krishna2017visual} & Image & 5,996 & 3,843,636 & 1,014 & 2,347,187 & 108,077 \\
    GQA~\cite{hudson2019gqa} & Image & 1,740 & - & 330 & - & 113,018 \\
    VrR-VG~\cite{liang2019vrr} & Image & 1,600 & 282,460 & 117 & 203,375 & 58,983 \\
    UnRel~\cite{peyre2017weakly} & Image & - & - & 18 & 76 & 1,071 \\
    SpatialSense~\cite{yang2019spatialsense} & Image & 3,679 & - & 9 & 13,229 & 11,569 \\
    SpatialVOC2K~\cite{belz2018spatialvoc2k} & Image & 20 & 5,775 & 34 & 9,804 & 2,026 \\
    PSG~\cite{yang2022panoptic} & Image (panoptic) & 133 & - & 56 & - & 49K \\
    AUG~\cite{li2024aug} & Image (overhead) & 76 & - & 61 & - & - \\
    STAR~\cite{li2025star} & Image (satellite) & 48 & 219,120 & 58 & 400,795 & 31,096 \\
    ReCon1M~\cite{yan2025recon1m} & Image (satellite) & 60 & 859,751 & 64 & 1,149,342 & 21,392 \\
    SkySenseGPT~\cite{luo2024skysensegpt} & Image (satellite, instruction) & - & - & - & - & - \\
    Traffic Scene Graph~\cite{zhou2022toward} & Image (traffic) & 2,266 & - & 4,272 & - & 451 \\
    ImageNet-VidVRD~\cite{shang2017video} & Video & 35 & - & 132 & 3,219 & 1,000 \\
    VidOR~\cite{shang2019annotating} & Video & 80 & - & 50 & - & 10,000 \\
    Action Genome~\cite{ji2020action} & Video & 35 & 0.4M & 25 & 1.7M & 10,000 \\
    AeroEye~\cite{nguyen2024cyclo} & Video (drone) & 56 & - & 384 & - & 2.2M \\
    PVSG~\cite{yang2023panoptic} & Video (panoptic) & 126 & - & 57 & 4,587 & 400 \\
    ASPIRe~\cite{nguyen2024hig} & Video (interlacements) & - & - & 4.5K & - & 1.5K \\
    Ego-EASG~\cite{rodin2024action} & Video (egocentric) & 407 & - & 235 & - & - \\
    IESG~\cite{liu2023learning} & Video (traffic) & 9 & - & 16 & - & - \\
    3DSSG~\cite{wald2020learning} & 3D & 528 & - & 39 & - & 48K \\
    PSG4D~\cite{yang20234d} & 4D & 46 & - & 15 & - & - \\
    4D-OR~\cite{ozsoy20224d} & 4D (operating room) & 12 & - & 14 & - & - \\
    MM-OR~\cite{ozsoy2025mm} & 4D (operating room) & - & - & - & - & - \\
    EgoExOR~\cite{ozsoy2025egoexor} & 4D (operating room) & 36 & - & 22 & 568,235 & - \\
    TSG Bench~\cite{yang2025llm} & Text & - & - & - & 11,820 & 4,289 \\
    FACTUAL~\cite{li2023factual} & Text, Image & 4,042 & - & 1,607 & 40,149 & 40,369 \\
    DiscoSG-DS~\cite{lin2025discosg} & Text, Image & 4,018 & - & 2,033 & 68,478 & 8,830 \\
    \bottomrule
  \end{tabular}
\end{table}

Table~\ref{tab:SGG_data}\footnote{The table is extended from the GitHub repository: \url{https://github.com/ChocoWu/Awesome-Scene-Graph-Generation}} lists the benchmarks by modality.
The image benchmarks run from the small precursors Visual Phrase~\cite{Sadeghi2011VisualPhrases}, Scene Graph~\cite{johnson2015image}, and VRD~\cite{lu2016visual}, which first exposed the long-tailed predicate distribution, to Visual Genome~\cite{krishna2017visual}, Open Images~\cite{kuznetsova2020open}, and GQA~\cite{hudson2019gqa}, with VrR-VG~\cite{liang2019vrr} and UnRel~\cite{peyre2017weakly} carved out to test visually relevant and zero-shot relations respectively.
Visual Genome is the primary benchmark but its raw release (33,877 object and 40,480 relationship synsets) is noisy and heavily long-tailed, so the community works with cleaned subsets: the broad cleaning in Table~\ref{tab:SGG_data} keeps roughly 6K object and 1K predicate classes, and the de-facto SGG benchmark VG150 restricts these to the 150 most frequent object and 50 most frequent predicate classes.
Spatial-only relations are covered by SpatialSense and SpatialVOC2K (see the preceding subsection), panoptic grounding by PSG~\cite{yang2022panoptic}, and non-standard viewpoints by the overhead AUG~\cite{li2024aug}, the satellite STAR~\cite{li2025star}, ReCon1M~\cite{yan2025recon1m}, and instruction-based SkySenseGPT~\cite{luo2024skysensegpt}, and the traffic-scene graphs of~\cite{zhou2022toward}.
Video benchmarks began with ImageNet-VidVRD~\cite{shang2017video} and the trajectory-annotated VidOR~\cite{shang2019annotating}, then specialized to human activity (Action Genome~\cite{ji2020action}), drone footage (AeroEye~\cite{nguyen2024cyclo}), panoptic video (PVSG~\cite{yang2023panoptic}), fine topological interlacements (ASPIRe~\cite{nguyen2024hig}), egocentric capture (Ego-EASG~\cite{rodin2024action}), and imminent traffic events (IESG~\cite{liu2023learning}).
3DSSG~\cite{wald2020learning} is the standard 3D benchmark; PSG4D~\cite{yang20234d} adds time, and the operating-room sets 4D-OR~\cite{ozsoy20224d}, MM-OR~\cite{ozsoy2025mm}, and EgoExOR~\cite{ozsoy2025egoexor} bring multimodal (RGB-D, audio, robot logs) and egocentric surgical scenes.
On the text side, TSG Bench~\cite{yang2025llm} evaluates LLMs on graph generation from narratives, while FACTUAL~\cite{li2023factual} and DiscoSG-DS~\cite{lin2025discosg} pair images with text for factually consistent, cross-modal graphs.

Evaluation separates what is given from what is predicted. In \textbf{Predicate Classification (PredCls)} boxes and object labels are given and only predicates are predicted; in \textbf{Scene Graph Classification (SGCls)} boxes are given and both object labels and predicates are predicted; in \textbf{Scene Graph Generation (SGGen, also SGDet)} everything is predicted, and a triplet counts only if both boxes overlap their ground truth at IoU $\geq 0.5$, with \textbf{Phrase Detection} relaxing this to the union box and SGGen+ also crediting correctly detected singleton nodes. Panoptic SGG replaces box IoU with mask IoU.
The standard metric is recall of ground-truth triplets among the top-$K$ predictions,
\begin{equation}
    R@K = \frac{|Top_{K} \cap GT|}{|GT|},
\end{equation}
under a graph constraint (one predicate per pair) or without it. Because $R@K$ is dominated by frequent predicates, \textbf{mean Recall@K} ($mR@K$) averages recall per predicate class and has become the long-tail metric of record; the gap between $R@K$ and $mR@K$ drove much of the debiasing literature~\cite{lin2020gps,li2021bipartite,zhou2024hierarchical}. Weighted triplet recall, micro- versus macro-averaging over images, the stricter $mAP_{rel}$ / $mAP_{phr}$ of the Open Images challenge, and retrieval-based R-Precision are reported in specific settings.

Each modality has its own benchmark and protocol, so leaderboards are not comparable across types: video SGG uses trajectory-based evaluation on ImageNet-VidVRD and VidOR, while Action Genome commonly uses frame-level PredCls, SGCls, and SGDet; textual SGG uses set F1 over triples on FACTUAL and LLM-judged action coverage on TSG Bench; 3D SGG uses object-, predicate-, and triplet-level recall on 3DSSG; panoptic SGG uses mask-IoU $R@K$ on PSG and PVSG. Table~\ref{tab:sgg_performance} therefore reports only the image setting, where VG150 SGDet is shared.

\begin{table}[!t]
\centering
\caption{Representative image SGG methods on Visual Genome (VG150), scene graph detection (SGDet) with graph constraint, as reported in the comparison of EGTR~\cite{im2024egtr}. Two-stage methods use a Faster R-CNN detector with a ResNeXt-101-FPN backbone; one-stage methods are DETR-based with the backbone shown. ``--'' = not reported.}
\label{tab:sgg_performance}
\fontsize{9}{11}\selectfont
\setlength{\tabcolsep}{3.0mm}
\begin{tabular}{llccccc}
\toprule
\textbf{Method} & \textbf{Family / backbone} & \textbf{Year} & \textbf{R@50} & \textbf{R@100} & \textbf{mR@50} & \textbf{mR@100} \\
\midrule
Neural Motifs~\cite{zellers2018neural} & Two-stage, X-101-FPN & 2018 & 32.1 & 36.9 & 5.5 & 6.8 \\
BGNN~\cite{li2021bipartite}            & Two-stage, X-101-FPN & 2021 & 31.0 & 35.8 & 10.7 & 12.6 \\
SGTR~\cite{li2022sgtr}                 & One-stage, DETR-R101 & 2022 & 24.6 & 28.4 & 12.0 & 15.2 \\
RelTR~\cite{cong2023reltr}             & One-stage, DETR-R50  & 2023 & 27.5 & -- & 10.8 & -- \\
EGTR~\cite{im2024egtr}                 & One-stage, Deformable DETR-R50 & 2024 & 30.2 & 34.3 & 7.9 & 10.1 \\
\bottomrule
\end{tabular}
\end{table}

Table~\ref{tab:sgg_performance} shows the trade-off that still shapes image SGG. The two-stage Motifs baseline keeps the highest overall recall, but its mean recall is the lowest, which is the frequent-predicate bias in numbers; BGNN's debiasing doubles mean recall at a small cost in $R@K$. The one-stage DETR-style models start several recall points lower, since they must detect and relate in one pass with a weaker backbone, yet SGTR and RelTR reach the best mean recall of the group, and EGTR closes most of the overall-recall gap to two-stage methods while running end to end.


\subsubsection{Challenges}
% 1.单模态SGG -> 统一SGG
% 2.与MLLM结合 构建结构化mllm
% 3.单粒度/层次化场景图 -> 无限粒度场景图
% 4.下游具身应用

\paragraph{From single-modality to universal SGG.}
Methods are built per modality, yet real scenes are multimodal, with images carrying visual detail and text carrying abstract concepts and events. A universal scene graph~\cite{wu2025universal} would express all of them in one cross-modal representation; the hard part is identifying the same objects, events, and relations across modalities and suppressing cross-modal noise, and granularity has to follow the application rather than a fixed ontology, since robotic interaction or driving may need relations at a level of detail that image-level graphs do not provide.

\paragraph{Integrating SGG with MLLMs.}
MLLMs are strong at language and weaker at structured visual understanding. Scene graphs offer a way in: pretraining objectives that mask and predict graph nodes (objects, attributes, relations) across image, video, and text inputs would teach a language model to read structure rather than only the next token, and the scene-graph-guided MLLMs of the preceding subsection~\cite{xu2025llava} are an early instance.

\paragraph{From images to embodied use.}
Most SGG analyzes static images or videos, but the graph is most useful when it is kept live: robots update it as they act~\cite{nguyen2025generating,werby2024hierarchical,garg2024robohop}, and driving systems read relations among lanes, signs, and pedestrians from it~\cite{zhou2022toward,liu2023learning,malawade2022roadscene2vec}. As these uses grow, the scene graph shifts from a product of perception to the state on which decision and interaction operate.





\subsection{Layout Estimation and Room Parsing}
\label{sec:layout_estimation}

Layout estimation recovers the structural geometry of an indoor space, the corners, walls, floor, and ceiling that bound it, as a simplified 3D ``wireframe'' called the room layout.
Room parsing extends this by also detecting, localizing, and labeling the objects inside (furniture and fixtures), giving a complete description of the scene.
For outdoor scenes, the analogous problem is recovering a road layout in bird's-eye view from a street image~\cite{YangLLYMHP21}.
Figure~\ref{fig:p3_s6_3_3_layout_overview} shows the inputs and output representations for the structural part of the task.

\begin{figure}[!t]
    \centering
    \includegraphics[width=0.99\linewidth]{figures/part3/section6/layout_parsing.pdf}
    \caption{Inputs and output representations of structural layout estimation. Top: (i) indoor perspective images, (ii) 360$^\circ$ panoramas, and (iii) outdoor street views. Bottom: (a) polygonal room layouts and (b) schematic indoor/outdoor surface-label maps, illustrating representations that may be predicted separately or jointly. Object-level parsing and outdoor bird's-eye-view outputs are not depicted.}
    \label{fig:p3_s6_3_3_layout_overview}
\end{figure}

\subsubsection{Task Definition}
Given an image $\mathbf{I}\in\mathbb{R}^{H\times W\times 3}$ (a perspective view, a $360^\circ$ panorama, or a multi-view or video capture), layout estimation recovers the dominant surfaces that bound the navigable volume of a room as a \emph{room layout}
$\mathcal{L}=\{P_k\}_{k=1}^{K}$, where each $P_k$ is a planar polygon in 3D; equivalently the output may be plane parameters $(\mathbf{n}_k,d_k)$ or a per-pixel structural-surface label map. Under a vertical-wall model, the layout compresses to a floor-plan polygon with floor and ceiling heights $y_f, y_c$. Outdoor road-layout methods instead predict road regions and vehicle occupancy in bird's-eye view.
Room parsing adds per-object instance masks $\{m_i\}$ or boxes with class labels, drawn from a fixed ontology or, in the open-vocabulary setting, from free-form language.

\subsubsection{Methodology}

\paragraph{Pre-deep-learning approaches.}
Classical room layout estimation rests on the \textbf{Manhattan assumption}~\cite{CoughlanY99}: the scene is built from surfaces aligned with three mutually orthogonal directions, as walls, floors, and ceilings in most built environments are.
Hedau \etal~\cite{HedauHF09} introduced the task, modeling a cluttered indoor image jointly as a 3D box layout (the room as if empty) and per-pixel surface labels (visible walls, floor, ceiling, and objects).
Three mutually orthogonal vanishing points are first estimated with a robustified variant of Rother's method~\cite{Rother02}, in which a line segment $l$ votes for a candidate vanishing point $p$ with
\begin{equation}
v(l,p) = |l| \cdot \exp\!\left(-\frac{\alpha(l,p)^2}{2\sigma^2}\right),
\end{equation}
where $\alpha(l,p)$ is the angle between $l$ and the line from the midpoint of $l$ to $p$, $|l|$ is the segment length, and $\sigma$ is a fixed bandwidth; candidates are selected greedily, removing the lines that supported each selected point.
The vanishing points fix the orientation of the box; sampling pairs of rays through two of them generates box-layout hypotheses, whose remaining edges follow from the third vanishing point, and a learned structured regressor scores the hypotheses and selects the layout.

\paragraph{Deep learning approaches.}
Mallya and Lazebnik~\cite{MallyaL15} first used a fully convolutional network to predict informative edge maps, then kept the classical hypothesis-generation-and-ranking step, scoring hypotheses with features from the learned edges together with line membership and geometric context.
Later deep estimators are best organized by input modality and geometric representation, which vary independently: panorama input does not imply a cuboid output, and non-cuboid geometry need not break the Manhattan assumption. Figure~\ref{fig:p3_s6_3_3_methodology_overview} groups the representative methods of 2015--2021 along these lines.

\begin{figure}[!t]
    \centering
    \includegraphics[width=\linewidth]{figures/part3/section6/layout_parsing_method.pdf}
    \caption{Representative layout-estimation methods, 2015--2021. Each lane pairs a core-mechanism illustration with a publication timeline: (i) perspective corner and edge prediction, (ii) panorama representations, (iii) beyond-cuboid floor plans, and (iv-a, iv-b) indoor plane geometry and outdoor cross-view prediction. Timeline spacing is not proportional; arrows indicate chronology, not method dependencies.}
    \label{fig:p3_s6_3_3_methodology_overview}
\end{figure}

\textbf{(i) Perspective-image corner and edge predictors.}
RoomNet~\cite{lee2017roomnet} regresses an ordered set of room corners with a recurrent encoder--decoder conditioned on one of 11 predefined cuboid topologies, and LayoutNet~\cite{ZouCSH18} predicts corner and boundary probability maps and fits a cuboid to them. The topology constraint simplifies reconstruction, but a single perspective view often leaves part of the room outside the frame.

\textbf{(ii) Panorama-based methods.}
A $360^\circ$ equirectangular panorama removes the field-of-view truncation of a perspective image, though not occlusion. LayoutNet~\cite{ZouCSH18} predicts 2D corner and boundary maps from a gravity-aligned panorama; HorizonNet~\cite{sun2019horizonnet} compresses the layout into three 1D per-column signals (floor-wall boundary, ceiling-wall boundary, and wall-wall corner indicator), adds Pano Stretch augmentation, and reconstructs general Manhattan layouts as well as cuboids; HoHoNet~\cite{sun2021hohonet} learns latent horizontal features decodable into layout, depth, and semantics; and LED2-Net~\cite{wang2021led2net} supervises the layout through differentiable depth rendering, comparing depth rendered from the predicted layout with the target geometry.

\textbf{(iii) Beyond-cuboid and Atlanta-world methods.}
Floor-plan prediction separates the number of walls from the constraints on their orientation. DuLa-Net~\cite{yang2019dulanet} combines equirectangular and ceiling--floor projections to estimate a floor plan under a Manhattan prior, which admits non-rectangular rooms with orthogonal walls. AtlantaNet~\cite{pintore2020atlantanet} projects a gravity-aligned panorama onto horizontal planes above and below the camera and predicts a room footprint and height under the Atlanta-world model, which keeps walls vertical but allows arbitrary horizontal orientation, including curved footprints. Both still assume vertical walls, so slanted walls and ceilings remain out of reach.

\textbf{(iv) Plane-based indoor geometry and outdoor cross-view prediction.}
GeoLayout~\cite{0005ZZ20} predicts pixel-level plane parameters, infers the depth map of each dominant plane, and reads layout boundaries off the intersections of those depth maps, a geometric alternative to regressing polygon vertices; the accompanying Matterport3D-Layout dataset covers cuboid and non-cuboid rooms. For outdoor scenes, Yang \etal~\cite{YangLLYMHP21} use a cross-view transformation with cycle-consistency constraints to predict road layout and vehicle occupancy in bird's-eye view from a monocular front-facing image.

\subsubsection{Datasets and Metrics}
Table~\ref{tab:layout_estimation_room_parsing_dataset} summarizes the benchmarks. Perspective-image evaluation grew from Hedau~\cite{HedauHF09} (308 web and LabelMe images with box layouts and surface labels) through its augmented successor Hedau+ and the more cluttered SUNbox drawn from SUN2012~\cite{XiaoHEOT10}, both built by Mallya and Lazebnik~\cite{MallyaL15} to train FCN edge predictors, to LSUN Room Layout~\cite{yu2015lsun}, the standard perspective benchmark with corner positions and cuboid topologies on over 5K images. Corridor and video data serve navigation: Yang~\cite{YangMS16} labels ground versus wall in corridor images, FreeYourCamera~\cite{FurlanMSLS13} records six-degree-of-freedom smartphone walks that frequently violate box and Manhattan assumptions, and MichiganIndoorCorridor~\cite{TsaiXLK11} provides non-Manhattan corridors with glass and reflections. Panoramic and 3D evaluation rests on PanoContext~\cite{zhang2014panocontext} (cuboid layouts and furniture for about 500 SUN360 panoramas, the standard testbed for LayoutNet, HorizonNet, and LED2-Net), Stanford 2D-3D-S~\cite{armeni2017joint} (70,496 perspective RGB images and 1,413 RGB panoramas, with aligned depth, normals, semantics, and 3D meshes over 270 rooms), and Matterport3D-Layout~\cite{0005ZZ20}, derived from Matterport3D~\cite{chang2017matterport3d} with RANSAC-fitted plane ground truth for cuboid and non-cuboid rooms.

\begin{table*}[!t]
    \centering
    \caption{Layout estimation and room parsing benchmarks. Scale is in images unless noted.}
    \label{tab:layout_estimation_room_parsing_dataset}
    \fontsize{9}{11}\selectfont
    \setlength{\tabcolsep}{2.0mm}
    \begin{tabular}{L{0.20\textwidth} c C{0.08\textwidth} C{0.11\textwidth} C{0.12\textwidth} L{0.27\textwidth}}
    \toprule
    \textbf{Dataset} & \textbf{Year} & \textbf{Scene} & \textbf{Modality} & \textbf{Scale} & \textbf{Annotation} \\
    \midrule
    Hedau \cite{HedauHF09} & 2009 & Room & Image & 308 & Box layout + surface labels \\
    Hedau+ \cite{MallyaL15} & 2015 & Room & Image & 337 & Box layout (Hedau + 75 bedroom images) \\
    SUNbox \cite{MallyaL15} & 2015 & Room & Image & 543 & Box layout, cluttered SUN2012 scenes \\
    LSUN Room Layout \cite{yu2015lsun} & 2015 & Room & Image & 5,394 & Corners + 11 cuboid topologies \\
    Yang \cite{YangMS16} & 2016 & Corridor & RGB-D image & 967 & Ground / non-ground masks \\
    FreeYourCamera \cite{FurlanMSLS13} & 2013 & Mixed & Video (6-DoF) & 10 seq.\ (84 frames) & Layout under non-box, non-Manhattan motion \\
    MichiganIndoorCorridor \cite{TsaiXLK11} & 2011 & Corridor & Video & 11 seq.\ ($\sim$3,300 frames) & Ground + three walls on first frame \\
    PanoContext \cite{zhang2014panocontext} & 2014 & Room & Panorama & $\sim$500 & Cuboid layout + furniture \\
    Matterport3D-Layout \cite{0005ZZ20} & 2020 & Room & RGB-D image & 7,360 & Plane parameters, cuboid + non-cuboid \\
    Stanford 2D-3D-S \cite{armeni2017joint} & 2017 & Building & RGB-D & 70,496 views; 1,413 panoramas & Depth, normals, semantics, mesh \\
    \bottomrule
    \end{tabular}
\end{table*}

Metrics follow the output representation. Perspective benchmarks (Hedau, LSUN) report \emph{pixel error}, the fraction of pixels whose predicted surface label differs from the ground truth, and \emph{corner error}, the mean distance between predicted and true corners normalized by image size; the edge-map methods additionally report the contour-detection F-scores at the optimal dataset scale (ODS, one threshold for the whole set) and optimal image scale (OIS, the best threshold per image). Panoramic and 3D benchmarks (PanoContext and panoramic layout subsets of Stanford 2D-3D-S and Matterport3D) report 3D and 2D \emph{room IoU} between the predicted and ground-truth layout volumes or floor plans, with frequency-weighted IoU and pixel accuracy as segmentation-style complements. A method trained for one input type rarely transfers to the other without re-training, so perspective and panoramic results are not compared directly, and the outdoor road-layout protocol~\cite{YangLLYMHP21} is separate again.

\begin{table}[!t]
    \centering
    \caption{Panorama-based layout estimation, 3D IoU (\%) on PanoContext, Stanford 2D-3D-S, and the 2,295-panorama Matterport3D layout subset (overall), as reported in the comparison of LED2-Net~\cite{wang2021led2net}. Higher is better.}
    \label{tab:layout_performance}
    \fontsize{9}{11}\selectfont
    \setlength{\tabcolsep}{1.5mm}
    \begin{tabular}{L{0.16\textwidth} L{0.28\textwidth} c C{0.10\textwidth} C{0.13\textwidth} C{0.14\textwidth}}
    \toprule
    \textbf{Method} & \textbf{Representation} & \textbf{Year} & \textbf{PanoContext} & \textbf{Stanford 2D-3D-S} & \textbf{Matterport3D} \\
    \midrule
    LayoutNet~\cite{ZouCSH18}              & 2D corner + boundary maps      & 2018 & 74.48 & 76.33 & 75.82 \\
    DuLa-Net~\cite{yang2019dulanet}         & Dual projection, Manhattan plan & 2019 & 77.42 & 79.36 & 75.05 \\
    HorizonNet~\cite{sun2019horizonnet}     & 1D column representation       & 2019 & 82.17 & 79.79 & 78.73 \\
    AtlantaNet~\cite{pintore2020atlantanet} & Top-down footprint, Atlanta world & 2020 & 78.76 & 82.43 & 80.02 \\
    LED2-Net~\cite{wang2021led2net}         & Differentiable depth rendering & 2021 & \textbf{82.75} & \textbf{83.77} & \textbf{81.52} \\
    \bottomrule
    \end{tabular}
\end{table}

Table~\ref{tab:layout_performance} compares the panorama methods under one protocol. The 1D column representation of HorizonNet gained about eight IoU points over LayoutNet's 2D maps on PanoContext, and the geometric supervision of LED2-Net adds a further point or two on every benchmark; AtlantaNet exceeds HorizonNet on the Matterport3D subset, whereas DuLa-Net remains below it, so these results do not show a uniform advantage for beyond-cuboid methods.

\subsubsection{Challenges}
The geometric priors that make layout estimation tractable also bound it. Cuboid models cannot express alcoves or L-shaped plans; Manhattan models keep orthogonal walls even when the footprint is free-form; the Atlanta-world model of AtlantaNet~\cite{pintore2020atlantanet} frees wall orientation but not wall tilt; and plane-based methods such as GeoLayout~\cite{0005ZZ20} approximate curved geometry only with many planes. Slanted enclosures, split levels, and the outdoor bird's-eye-view target of~\cite{YangLLYMHP21}, which is a different output space altogether, all sit outside the standard formulations.
Clutter and occlusion are the second obstacle: furniture hides the corners, boundaries, and vanishing lines that both classical and learned methods read, so estimators must hallucinate hidden structure from layout priors, and confidently wrong boundaries without uncertainty are the typical failure.
Generalization follows from the data: benchmarks are indoor, mostly rooms and corridors, and models trained on one set degrade on another's appearance, architecture, and lighting, while sparse, unposed multi-view capture breaks the pose assumptions of classical pipelines; the feed-forward geometry models of \S\ref{sec:visual_odometry} are the natural starting point for layout from such input.





\subsection{Action Detection}
\label{sec:action_detection}
% \wsq{人体动作行为识别与时序定位}

Action detection (AD) finds human actions in untrimmed video: for each action instance it returns the time interval $[t_s, t_e]$ in which it occurs and its category $c$. Where action recognition classifies a clip that has already been trimmed to one action, detection must also decide when actions start and end, in streams that contain many of them.
It belongs to spatial structure construction because it organizes the entities and relations of a scene along the time axis into events, the temporal counterpart of the scene graphs above and the basis of video SGG and human--object interaction.
Figure~\ref{fig:p3_s6_3_4_ad_overview} illustrates the task.


\begin{figure}[!ht]
    \centering
    \includegraphics[width=\linewidth]{figures/part3/section6/action_detection.pdf}
    \caption{Action detection: each action instance in an untrimmed video is localized to $[t_s,t_e]$ and labeled with class $c$ (top). \textit{Bottom}: three task illustrations: (i) action classification of a trimmed clip, (ii) temporal action localization (illustrated with one interval), and (iii) temporal action detection of multiple instances.}
    \label{fig:p3_s6_3_4_ad_overview}
\end{figure}

\subsubsection{Task Definition}
\label{sec:6_3_4_definition}
Given an untrimmed video (or a continuous multi-sensor stream), AD outputs a set of tuples $(t_s, t_e, c)$ with start time, end time, and category for every action instance.
The following terms emphasize different aspects of the task. \textbf{Action classification (AC)} takes a trimmed clip and predicts one category, with no localization. \textbf{Temporal action localization (TAL)} predicts start and end times in untrimmed video, and may include multiple instances; the term is often used interchangeably with TAD. \textbf{Temporal action detection (TAD)} emphasizes joint localization and classification of multiple instances per video, possibly overlapping or concurrent, across temporal scales, each to be localized and classified.


\begin{figure}[!ht]
    \centering
    \includegraphics[width=0.99\linewidth]{figures/part3/section6/action_detection_method_1.pdf}
    \caption{Representative methods and mechanisms for (i) action classification, (ii) temporal localization with boundary estimation and alternative supervision signals, and (iii) temporal detection with anchor-based, anchor-free, boundary-based, and query-based formulations. Method timelines show publication order, not computational dependencies; spacing is not proportional. Same-year methods are grouped. The localization and detection panels emphasize different aspects of closely related tasks.}
    \label{fig:p3_s6_3_4_methodology_overview}
\end{figure}

\subsubsection{Methodology}

Figure~\ref{fig:p3_s6_3_4_methodology_overview} organizes the methods by the three settings; each builds on the previous one's representations.

\paragraph{Action Classification.}
Two-stream CNNs~\cite{simonyan2014two} model appearance from RGB frames in one stream and motion from dense optical flow in another, fusing the two at the output. 3D CNNs~\cite{diba2017temporal} convolve over space and time jointly, C3D~\cite{tran2015learning} and I3D~\cite{carreira2017quo} being the standard examples, at the cost of a large parameter count that R(2+1)D~\cite{tran2018closer} and P3D~\cite{qiu2017learning} reduce by factoring 3D convolution into 2D spatial and 1D temporal steps. TSN~\cite{wang2016temporal} samples a few snippets sparsely across a long video and supervises at video level, which scales long-range temporal modeling cheaply, and Transformers~\cite{girdhar2019video,yang2022recurring} model global spatio-temporal dependencies among sampled tokens beyond the local reach of convolution.

\paragraph{Temporal Action Localization.}
Early TAL ran a classifier over dense, fixed-size sliding windows, which was both slow and blind to the wide range of action durations. Structured segment networks~\cite{zhao2017temporal} instead model a proposal's starting, ongoing, and ending stages and score its completeness, which refines boundary regression. Boundary-sensitive methods~\cite{lin2018bsn} go further and predict start and end probabilities at every snippet, then pair high-confidence boundary points into proposals; this flexible boundary generation is the component on which later detection systems rest.

Because precise boundary annotation is expensive, weakly supervised TAL (WSTAL)~\cite{shou2018autoloc,zhang2021cola,paul2018w} learns from video-level labels alone. The task becomes multiple-instance learning: a video is positive if at least one segment contains the action, class activation sequences are derived from video features, and foreground segments are thresholded into proposals. Trained only with a classification loss, such models fire on the most discriminative moments rather than the whole action, so proposals are incomplete and high confidence does not imply accurate boundaries. Point-level supervision (P-TAL)~\cite{xia2024realigning,lee2021learning,fu2022compact} adds one labeled frame per instance, which tends to fall near the temporally most salient, roughly central moment, and uses it to realign confidence and infer complete boundaries from that single anchor.


\paragraph{Temporal Action Detection.}
TAD methods differ in how they generate and score temporal proposals.
Anchor-based detection~\cite{gao2017turn} transplants Faster R-CNN and SSD to the time axis: a dense set of multi-scale temporal anchors is classified and regressed, first in two stages (S-CNN~\cite{shou2016temporal}) and then in one (SSAD~\cite{lin2017single}); fixed anchors produce many redundant proposals and limit boundary precision, which anchor-free variants such as AFSD~\cite{lin2021learning} address by regressing boundaries from a temporal center representation.
Boundary-based detection~\cite{xiong2017pursuit,lin2018bsn} is bottom-up: start and end probabilities are predicted over the whole timeline and boundary matching scores every candidate pairing, which gives flexible, precise boundaries; TSA-Net~\cite{gong2019scale} adds multi-dilation temporal convolution for multi-scale context and BGCN~\cite{chen2021boundary} graph convolution over co-occurring proposals.
Query-based detection~\cite{wu2021towards,shi2022react,kim2023self,zhu2024dual} follows DETR: a small set of learned queries attends to global video features in a Transformer encoder--decoder and decodes the action set in parallel, removing anchors, NMS, and hand-designed matching at the cost of heavier computation.
Across the three, anchors give dense parallel scoring but coarse boundaries, boundary methods give precise boundaries but complex matching, and query methods give the simplest end-to-end pipeline; all three keep moving toward longer-range context, continuous temporal representations, and less frame-level annotation.

\subsubsection{Datasets and Metrics}
Table~\ref{tab:action_detection_dataset} lists the benchmarks by role. UCF101~\cite{soomro2012ucf101} and HMDB51~\cite{kuehne2011hmdb} are trimmed-clip classification sets; AVA~\cite{gu2018ava} adds spatio-temporal boxes on movie footage. THUMOS'14~\cite{jiang2014thumos}, whose 20 classes carry dense temporal boundaries, and ActivityNet~\cite{caba2015activitynet}, with 200 classes over 20K untrimmed videos, are the standard TAL and TAD benchmarks, extended in scale and granularity by HACS Segment~\cite{zhao2019hacs} and FineAction~\cite{liu2022fineaction}. MultiTHUMOS~\cite{yeung2018every}, Charades~\cite{sigurdsson2016hollywood}, and TVSeries~\cite{de2016online} supply dense, multi-label, overlapping actions in sports, homes, and television; Breakfast~\cite{kuehne2014language}, 50Salads~\cite{stein2013combining}, MPII Cooking 2~\cite{rohrbach2016recognizing}, and COIN~\cite{tang2019coin} cover procedural cooking and instructional activity for segmentation and step detection; and MEVA~\cite{corona2021meva} (thousands of hours of surveillance) and MoVi~\cite{ghorbani2021movi} (motion-captured actors) sit at the large-scale and controlled extremes.

\begin{table}[!t]
\centering
\caption{Benchmark datasets for action detection.}
\label{tab:action_detection_dataset}
\fontsize{9}{11}\selectfont
\setlength{\tabcolsep}{5.5pt}
\begin{tabular}{lccccc}
\toprule
\multirow{2}{*}{\textbf{Dataset}} & \textbf{Action} & \multirow{2}{*}{\textbf{Videos}} & \textbf{Duration} & \textbf{Multi-label} & \multirow{2}{*}{\textbf{Domain}} \\
& \textbf{Categories} & & \textbf{(hours)} & \textbf{(Per-Frame)} & \\
\midrule
UCF101~\cite{soomro2012ucf101}           & 101  & 13,320 & 27 & --  & Sports, Daily Activities\\
HMDB51~\cite{kuehne2011hmdb}           & 51   & 6,849  & -- & --  & Movies \\
AVA~\cite{gu2018ava}              & 80   & 430    & 108 & Yes & Movies \\
THUMOS'14~\cite{jiang2014thumos}        & 101 (20 w/ boundaries)  & 18,394 & --  & No & Sports \\
MultiTHUMOS~\cite{yeung2018every}  & 65                         & 413             & --                       & Yes         & Sports                 \\
Breakfast~\cite{kuehne2014language}    & 48                         & 1,712           & 77                    & No          & Cooking                \\
50Salads~\cite{stein2013combining}     & 17                         & 50              & 5.3                       & No          & Cooking                \\
MPII cooking 2~\cite{rohrbach2016recognizing} & 59                        & 273             & 27          & No          & Cooking                \\
TVSeries~\cite{de2016online}     & 30                         & 27              & --                   & Yes         & TV Series              \\
HACS Segment~\cite{zhao2019hacs} & 200                        & 50,000             & --                       & No          & Daily Activities       \\
ActivityNet-1.3~\cite{caba2015activitynet}  & 200  & 19,994 & 849   & No & Daily Activities \\
Charades~\cite{sigurdsson2016hollywood}         & 157  & 9,848  & --    & Yes         & Daily Activities \\
MoVi~\cite{ghorbani2021movi}             & 20   & --     & 9    & -- & Daily Activities \\
MEVA~\cite{corona2021meva}             & 37   & --     & 9,300  & -- & Surveillance \\
COIN~\cite{tang2019coin}         & 180 tasks                  & 11,827          & 476                       & No          & Daily Activities       \\
FineAction~\cite{liu2022fineaction}      & 106  & 16,732 & --    & Yes         & Daily Activities \\
\bottomrule
\end{tabular}
\end{table}

Each setting has its own metric. Action classification reports Top-1 accuracy, $\mathrm{Acc}_{\mathrm{Top}\text{-}1} = \frac{1}{N} \sum_{i=1}^{N} \mathbb{I}(\hat{y}_i = y_i)$, with Top-$k$ on large label sets; temporal boundaries are given, so localization is not assessed. Localization quality is measured by the temporal intersection over union of predicted and ground-truth intervals,
\begin{equation}
    \mathrm{tIoU}(I_p, I_g) =
    \frac{|I_p \cap I_g|}{|I_p \cup I_g|},
\end{equation}
a prediction counting as correct above a threshold $\tau$; TAL reports average precision at one or several thresholds. TAD scores detections $(I_p, \hat{y}, s)$ with confidence $s$ by mean average precision over classes and thresholds,
\begin{equation}
    \mathrm{mAP} = \frac{1}{|\mathcal{C}|}
    \sum_{c \in \mathcal{C}}
    \frac{1}{|\{\tau_m\}|}
    \sum_{\tau_m} \mathrm{AP}_c(\tau_m),
\end{equation}
which penalizes misclassification, loose boundaries, and duplicates together; THUMOS'14 uses $\tau \in \{0.3,\dots,0.7\}$ and ActivityNet $\{0.5,\dots,0.95\}$, and proposal methods also report average recall at a fixed proposal budget. Weakly supervised methods are often evaluated on ActivityNet-1.2 rather than 1.3, so their numbers are not interchangeable with fully supervised ones.

\begin{table*}[!t]
    \centering
    \caption{Representative methods under the two shared protocols. \textit{Top}: action classification, Top-1 accuracy (\%) on UCF101 (split 1, two-stream RGB+flow where applicable, with ImageNet or Kinetics pretraining as in each paper). \textit{Bottom}: temporal action detection on THUMOS'14, mAP (\%) at tIoU $=0.5$ and averaged over tIoU $\in\{0.3,\dots,0.7\}$, as reported by the cited papers and the comparisons in AFSD~\cite{lin2021learning}, CoLA~\cite{zhang2021cola}, and DualDETR~\cite{zhu2024dual}; ``weak'' = video-level labels only. ``--'' = not reported.}
    \label{tab:ad_performance}
    \fontsize{9}{11}\selectfont
    \setlength{\tabcolsep}{2.0mm}
    \begin{tabular}{L{0.22\textwidth} L{0.30\textwidth} c C{0.14\textwidth} C{0.16\textwidth}}
    \toprule
    \textbf{Method} & \textbf{Family} & \textbf{Year} & \textbf{UCF101 Top-1}$\uparrow$ & \\
    \midrule
    Two-Stream CNN~\cite{simonyan2014two} & RGB + optical-flow streams & 2014 & 88.0 & \\
    TSN~\cite{wang2016temporal}           & Sparse segment sampling    & 2016 & 94.2 & \\
    R(2+1)D~\cite{tran2018closer}         & Factorized 3D convolution  & 2018 & 97.3 & \\
    I3D~\cite{carreira2017quo}            & Inflated 3D convolution    & 2017 & \textbf{98.0} & \\
    \midrule
    \hiderowcolors \textbf{Method} & \textbf{Family} & \textbf{Year} & \textbf{mAP@0.5}$\uparrow$ & \textbf{Avg.\ mAP@[0.3:0.7]}$\uparrow$ \\
    \midrule
    S-CNN~\cite{shou2016temporal}     & Anchor-based, two-stage          & 2016 & 19.0 & -- \\
    SSN~\cite{zhao2017temporal}       & Segment-structure proposals      & 2017 & 29.8 & -- \\
    BSN~\cite{lin2018bsn}             & Boundary-sensitive proposals     & 2018 & 36.9 & 36.8 \\
    AutoLoc~\cite{shou2018autoloc}    & Weak, outer--inner contrast      & 2018 & 21.2 & -- \\
    CoLA~\cite{zhang2021cola}         & Weak, snippet contrastive        & 2021 & 32.2 & -- \\
    AFSD~\cite{lin2021learning}       & Anchor-free, salient boundary    & 2021 & 55.5 & 52.0 \\
    DualDETR~\cite{zhu2024dual}       & Query-based (dual DETR)          & 2024 & \textbf{70.4} & \textbf{66.8} \\
    \bottomrule
    \end{tabular}
\end{table*}

Table~\ref{tab:ad_performance} shows how far the two settings have diverged. Trimmed-clip classification on UCF101 was effectively solved once Kinetics-pretrained 3D networks arrived, with I3D and R(2+1)D within a point of each other near 98\%. Detection on THUMOS'14 tells a different story: mAP@0.5 moved from 19.0 (S-CNN) to 36.9 (BSN) over the anchor-to-boundary transition, jumped to 55.5 with anchor-free salient-boundary features (AFSD), and reached 70.4 with query-based decoding (DualDETR), while the best weakly supervised method (CoLA, 32.2) matches fully supervised accuracy of only a few years earlier, which quantifies what video-level labels still cost.


\subsubsection{Challenges}
Fully supervised TAD depends on densely annotated boundaries, which scale poorly to long, loosely structured real-world video; weakly and semi-supervised methods relax this with video-level labels or sparse temporal cues~\cite{zhao2025pose,ghoddoosian2023weakly}, at the accuracy cost the table above makes visible.
Actions are also still detected in isolation. Coupling TAD with structured temporal representations such as scene graphs or event schemas~\cite{wu2025universal,lu2025benchmarking} would let a system infer boundaries, temporal relations, and higher-level activity structure jointly, which matters most in instructional video, driving, and human--robot interaction.
Finally, benchmarks and detectors remain visual and closed-set, whereas real action understanding draws on audio, language, gaze, and context and must handle actions outside a fixed vocabulary~\cite{zhou2024learning,bao2025exploiting,liu2026synergistic}; multimodal temporal detection, where audio or language sharpens boundaries and disambiguates semantics, and open-vocabulary formulations that loosen the classification-centric setup are the directions large multimodal video corpora now make feasible.





\subsection{Mapping and SLAM}
\label{sec:mapping_slam}
Mapping builds a spatial representation of an unknown environment from streaming sensor observations, together with an estimate of the agent's trajectory through it.
Doing so from on-board sensors requires the agent to localize itself inside the partial map it is still building: the chicken-and-egg problem known as \emph{Simultaneous Localization and Mapping (SLAM)}. SLAM resolves it by estimating trajectory and map jointly from the same sensor stream, and observations of previously visited places supply the external references that bound the drift of purely proprioceptive odometry (IMU, wheel encoders). The result is a globally consistent spatial model that supports localization, path planning, obstacle avoidance, and exploration.
Figure~\ref{fig:p3_s6_3_5_slam_modules} depicts the pipeline from sensors through front-end odometry and back-end optimization to the final map.



\begin{figure}[!ht]
    \centering
    \includegraphics[width=\linewidth]{figures/part3/section6/slam.pdf}
    \caption{Example SLAM outputs: real-time visual SLAM (left), a keyframe trajectory over a sparse feature map with tracked image features in the inset, and LiDAR SLAM (right), a dense point-cloud map.}
    \label{fig:p3_s6_3_5_slam_overview}
\end{figure}

\subsubsection{Task Definition}
The input is a stream of exteroceptive observations (camera images, RGB-D frames, LiDAR, or radar), optionally fused with proprioceptive readings (IMU, wheel odometry, velocity and steering commands) and priors such as a coarse map, an initial pose, or GPS.
The output is a map plus the estimated trajectory. Maps lie on an abstraction--detail spectrum: a \emph{topological map} is a graph of distinctive places and their traversable connections, with no metric layout; a \emph{sparse metric map} stores registered keypoints or landmarks and the camera-pose graph that observes them, enough for accurate localization but not for geometry; and a \emph{dense metric map} reconstructs a complete surface or volume (point cloud, mesh, TSDF, or occupancy grid) for obstacle avoidance and 3D reconstruction. Two dense representations have become prominent in the past three years: \emph{neural implicit maps} that store geometry and appearance in network weights (an MLP in iMAP~\cite{Sucar2021iMAP}, a voxel-grid encoder in NICE-SLAM~\cite{Zhu2022NICESLAM}) and \emph{3D Gaussian Splatting (3DGS) maps} that represent the scene as explicit, differentiable 3D Gaussians~\cite{Keetha2024SplaTAM,Yan2024GSSLAM,Matsuki2024MonoGS}.


\subsubsection{Methodology}

SLAM systems are classified by the exteroceptive sensor and by the density of the map they build.
\textbf{LiDAR SLAM} measures range directly with a stable error model; KISS-ICP~\cite{Vizzo2023KISSICP} and KISS-SLAM~\cite{Guadagnino2025KISSSLAM} show that a minimal pipeline delivers robust real-time 3D LiDAR odometry and mapping at city scale. \textbf{Visual SLAM} uses monocular, stereo, or RGB-D cameras, exploiting texture for place recognition, with PTAM~\cite{KleinMurray2007PTAM} and the ORB-SLAM family~\cite{mur2015orb,Campos2021ORBSLAM3} as the foundational systems. \textbf{Multi-sensor SLAM} tightly fuses cameras, IMUs, and/or LiDAR (VINS-Mono, OKVIS, FAST-LIO2), and \textbf{event-camera SLAM} uses asynchronous per-pixel brightness changes for high-speed and high-dynamic-range operation~\cite{Huang2026EventSLAMSurvey}.
By map density, sparse keypoint-and-pose-graph systems (ORB-SLAM) sit next to direct methods that minimize photometric error over semi-dense (LSD-SLAM~\cite{engel2014lsd}) or sparse (DSO~\cite{engel2017direct}) image samples, dense geometric systems that fuse point clouds, meshes, or TSDFs (KinectFusion, ElasticFusion), neural implicit systems (iMAP~\cite{Sucar2021iMAP}, NICE-SLAM~\cite{Zhu2022NICESLAM}, Co-SLAM, NICER-SLAM), and 3DGS systems that render and optimize Gaussians jointly with poses (SplaTAM~\cite{Keetha2024SplaTAM}, MonoGS~\cite{Matsuki2024MonoGS}, SplatSLAM).
Read historically, the field moved from feature-based geometry (PTAM, ORB-SLAM) to direct methods (LSD-SLAM, DSO), to tightly coupled multi-sensor systems (OKVIS, VINS-Mono, ORB-SLAM3), to learning-augmented pipelines (DROID-SLAM~\cite{teed2021droid}, SuperPoint and SuperGlue matching~\cite{detone2018superpoint,sarlin2020superglue}), to neural implicit and then 3DGS maps, and most recently to 3D geometry foundation models (DUSt3R~\cite{wang2024dust3r}, MASt3R~\cite{mast3r}, MASt3R-SLAM~\cite{Murai2025MASt3RSLAM}).

\begin{figure}[!ht]
    \centering
    \includegraphics[width=\linewidth]{figures/part3/section6/slam_method.pdf}
    \caption{SLAM module data flow. Sensor observations feed odometry and place recognition; odometry constraints and verified revisit constraints enter back-end optimization. Mapping registers raw observations using optimized poses and landmarks, producing a map and trajectory. The blue edge in the pose graph is a loop-closure constraint between previously visited poses, not an execution-loop arrow. Map representations are alternatives.}
    \label{fig:p3_s6_3_5_slam_modules}
\end{figure}

Whatever the sensor or map, a SLAM system runs the four modules of Figure~\ref{fig:p3_s6_3_5_slam_modules} in a loop.
The \textbf{front end} is visual (or LiDAR) odometry, the incremental motion estimate between consecutive readings treated as a stand-alone task in \S\ref{sec:visual_odometry}; because each pose is estimated relative to the last, its errors accumulate as drift.
\textbf{Back-end optimization} corrects that drift globally. Modern back ends pose the problem as factor-graph optimization: robot poses and landmarks are nodes, and measurements between them (visual reprojection, IMU pre-integration, GPS, loop-closure relative poses) are factors; solving for all nodes at once yields a consistent trajectory and map, factor graphs absorb heterogeneous sensors naturally, and incremental solvers such as iSAM2 update the estimate in real time, with GTSAM, g\textsuperscript{2}o, and OKVIS as representative libraries.
\textbf{Loop closing} supplies the back end with its strongest constraints by recognizing that the agent has returned to a previously visited place, comparing current observations against the history and, on a match, adding a constraint between two temporally distant poses that pulls the accumulated drift of the intervening trajectory back into alignment.
\textbf{Mapping} finally places the sensor data in the global frame using the optimized trajectory, in whatever representation the application needs; for neural implicit and 3DGS maps it also minimizes a rendering loss (photometric or perceptual such as LPIPS) between observed images and views rendered from the current map, optimizing geometry, appearance, and poses jointly~\cite{Keetha2024SplaTAM}.

\begin{table*}[!th]
\centering
\caption{Mapping and SLAM benchmarks. Scale counts sequences unless noted.}
\label{tab:mapping_slam_dataset}
\fontsize{9}{11}\selectfont
\setlength{\tabcolsep}{1.2mm}
\begin{tabular}{L{0.17\textwidth} C{0.10\textwidth} C{0.16\textwidth} C{0.11\textwidth} C{0.13\textwidth} C{0.07\textwidth} C{0.13\textwidth}}
\toprule
\textbf{Dataset} & \textbf{Task} & \textbf{Modality} & \textbf{Agent} & \textbf{Scene} & \textbf{Scale} & \textbf{Ground Truth} \\
\midrule
COLD \cite{PronobisC09} & VSLAM & mono + omni + odom & wheeled & lab & 76 & laser \\
TUM RGB-D \cite{Sturm2012TUMRGBD} & VSLAM & RGB-D & handheld, wheeled & office, work shed & 39 & MoCap \\
KITTI \cite{KITTI-GeigerLU12} & VSLAM, VO & stereo + laser & car & suburban, highway & 22 & GPS/OXTS \\
ICL-NUIM \cite{HandaWMD14} & VSLAM, VO & RGB-D & synthetic & home, office & 8  & simulator \\
EuRoC \cite{Burri2016EuRoC} & MSLAM & stereo + IMU & drone & work shed & 11  & Vicon MoCap/laser \\
SceneNet \cite{McCormacHLD17} & VSLAM & RGB-D & synthetic & home & 15K  & simulator \\
TUM VI \cite{Schubert2018TUMVI} & MSLAM, VO & stereo + IMU & handheld & campus buildings & 28  & MoCap \\
Bonn RGB-D \cite{PalazzoloBLGS19} & VSLAM & RGB-D & handheld & lab & 24  & MoCap \\
InteriorNet \cite{LiSMCTYHTL18} & MSLAM & RGB-D + IMU & synthetic & home & 15K  & simulator \\
OpenLORIS-Scene \cite{ShiLZTTLZSQSGW020} & MSLAM & RGB-D stereo + IMU + odom & wheeled & indoor & 22  & MoCap/laser \\
Replica \cite{straub2019replica}        & dense       & RGB-D     & synthetic      & indoor       & 18    & simulator \\
ScanNet++ \cite{yeshwanth2023scannet++}  & dense, sem  & RGB-D     & handheld       & indoor       & 460   & laser \\
TartanAir \cite{TartanAir}      & MSLAM, VO   & stereo + IMU & synthetic  & diverse      & 1M+ frames   & simulator \\
ETH3D SLAM \cite{ETH3DSLAM}             & VSLAM, VO   & stereo + global shutter & handheld & indoor/outdoor &  -- & LiDAR \\
DSEC \cite{gehrig2021dsec}              & event SLAM  & stereo events & car       & driving      & 53    & LiDAR + GPS \\
CU-Multi \cite{albin2025cu}            & collab.\ SLAM & RGB-D + LiDAR + IMU & wheeled & outdoor & 8 runs & RTK-GNSS \\
\bottomrule
\end{tabular}
\end{table*}

\subsubsection{Datasets and Metrics}
Table~\ref{tab:mapping_slam_dataset} groups the benchmarks by what they stress. The classical trajectory sets are TUM RGB-D~\cite{Sturm2012TUMRGBD} (handheld indoor RGB-D with motion-capture ground truth), KITTI~\cite{KITTI-GeigerLU12} (stereo and LiDAR driving), and the synthetic ICL-NUIM~\cite{HandaWMD14} with exact ground truth, while COLD~\cite{PronobisC09} targets topological localization under changing conditions. Visual-inertial systems are tested on EuRoC~\cite{Burri2016EuRoC} (micro-aerial vehicle), TUM VI~\cite{Schubert2018TUMVI}, and OpenLORIS-Scene~\cite{ShiLZTTLZSQSGW020}, the last built for lifelong SLAM in changing indoor scenes, and Bonn RGB-D~\cite{PalazzoloBLGS19} isolates dynamic objects. Large synthetic corpora (SceneNet~\cite{McCormacHLD17}, InteriorNet~\cite{LiSMCTYHTL18}, Replica~\cite{straub2019replica}, TartanAir~\cite{TartanAir}) supply photorealistic training and controlled evaluation, with Replica the standard testbed for neural implicit and 3DGS SLAM; ScanNet++~\cite{yeshwanth2023scannet++} and ETH3D SLAM~\cite{ETH3DSLAM} add laser-scanned geometry for dense reconstruction quality; and DSEC~\cite{gehrig2021dsec} and CU-Multi~\cite{albin2025cu} extend evaluation to event cameras and multi-robot data association. AirSwarm~\cite{li2025airswarm} is a multi-UAV research platform rather than a benchmark dataset.

Evaluation covers trajectory, map, rendering, and system cost. Trajectory accuracy uses the absolute trajectory error and relative pose error defined in \S\ref{sec:visual_odometry}, after aligning the estimated trajectory to the ground truth by a rigid or, for monocular systems, similarity transform; RMSE is the usual summary, with mean, median, and standard deviation also reported.
Dense maps are scored by \emph{accuracy} (mean distance from reconstructed points to the nearest ground-truth point) and \emph{completeness} (the reverse), summarized by an F-score at a distance threshold, with Chamfer distance common for neural and 3DGS maps. Neural implicit and 3DGS maps add novel-view-synthesis metrics (PSNR, SSIM, LPIPS) on held-out views of the same trajectory. Memory footprint, tracking and mapping frame rates, and end-to-end latency decide embedded deployment~\cite{Guadagnino2025KISSSLAM}.
Mature evaluation tools such as EVO\footnote{\url{https://github.com/MichaelGrupp/evo}} and the RPG Trajectory Evaluation toolbox\footnote{\url{https://github.com/uzh-rpg/rpg_trajectory_evaluation}} implement these protocols.








\subsubsection{Challenges}
Deployment often violates common SLAM assumptions of static scenes, stable illumination, abundant texture, and ample computation. Dynamic objects induce correspondences that drift with their own motion; illumination change, reflective and transparent surfaces, and repetitive texture create perceptual aliasing that confuses data association and place recognition alike; and resource-constrained platforms must run ever heavier front and back ends. Current work answers each directly. Moving features are detected and excluded through semantic segmentation, motion-consistency checks, or learned masks; learned matchers (SuperPoint~\cite{detone2018superpoint}, SuperGlue~\cite{sarlin2020superglue}, LoFTR~\cite{Sun2021LoFTR}) replace hand-crafted descriptors and tolerate larger viewpoint and illumination changes; neural implicit and 3DGS maps are optimized jointly with poses and yield photorealistic reconstructions usable for inspection, telepresence, and policy learning; and the two-view geometry priors of DUSt3R~\cite{wang2024dust3r} and MASt3R~\cite{mast3r} remove the fragile hand-engineered front end, already enabling real-time dense monocular SLAM at the robustness of MASt3R-SLAM~\cite{Murai2025MASt3RSLAM}.

Two broader shifts follow. Lifelong SLAM aims at months of continuous operation by combining relocalization, map update, and place re-identification against changing environments, and is now an evaluation track of its own; collaborative multi-agent SLAM fuses observations across robots for faster coverage and graceful degradation, increasingly together with event cameras and heterogeneous sensors for high-speed, low-light, and high-dynamic-range conditions. What remains open is concrete: 3DGS maps still do not scale to city-level memory and rendering budgets, learned components carry a sim-to-real gap, and incorrect loop closures under strong perceptual aliasing are still the failure that breaks every back end.





% \newpage



\section{Spatio-Temporal Perception}
\label{sec:p3-sec6-4-spatio-temporal-perception}

\FramedInitial{M}{ost} of what precedes has described a single instant. Spatio-temporal perception adds the missing axis, asking not only where entities are but how they move, when events happen, and which identities persist. 
Its three concerns rise from the pixel through the object to the event: scene flow recovers dense, per-point motion between frames; object tracking maintains the identity of an entity as it moves through space and time; and temporal localization fixes when an action or event occurs along a sequence.
Together they give the largely static perception of the preceding sections a temporal thread, the continuity on which motion understanding and long-horizon perception depend.






\subsection{Scene Flow Estimation}
\label{sec:scene_flow}
Scene flow is the dense three-dimensional motion of every visible surface point between two instants: optical flow lifted out of the image plane and joined to the change in depth. The image-based pipelines surveyed here take stereo video, monocular video, or a multi-view sequence and produce per-pixel 3D motion; some refinement modules additionally consume estimated or measured depth. Dynamic scene reconstruction~\cite{park2021nerfies}, motion-aware rendering~\cite{li2021neural}, physical interaction modeling~\cite{Bear2021Physion}, and egocentric world understanding~\cite{grauman2022ego4d} all consume it.

\subsubsection{Task Definition}
Given consecutive frames $I_t$ and $I_{t+1}$ (optionally with a stereo partner or several views, and with intrinsics and relative poses when known), scene flow estimation returns a field $\mathbf{v}(\mathbf{x}) \in \mathbb{R}^3$ that displaces the surface point imaged at pixel $\mathbf{x}$ in frame $t$ to its position at $t+1$, in camera or world coordinates. In the stereo setting this is equivalently a triple $(d_t, d_{t+1}, \mathbf{f})$ of disparity in both frames plus optical flow, which is how the KITTI benchmark represents and scores it. In image-only pipelines, structure and motion are inferred from appearance, parallax, and temporal correspondence. Geometry-aware refinement modules may instead take depth estimated upstream or measured by RGB-D sensors; point-cloud-only methods are outside the scope here.

\subsubsection{Methodology}

\begin{figure}[!ht]
    \centering
    \includegraphics[width=0.99\linewidth]{figures/part3/section6/scene_flow_estimation.pdf}
    \caption{Image- and video-based scene flow. Top row: observation settings and supervision signals. Bottom row: (1) variational / piecewise-rigid energy minimization, (2) stereo and temporal CNNs, (3) recurrent SE(3) refinement with image and depth pairs, (4) monocular estimation using photometric and geometric constraints, and (5) related motion-field approaches to dynamic view synthesis. The columns illustrate distinct mechanisms rather than successive stages of one pipeline.}
    \label{fig:p3_s6_4_3_method_overview}
\end{figure}

Methods differ along two axes: what they observe (stereo, the classical and still dominant setting; monocular video, the self-supervised frontier; multi-view rigs) and what supervises them (synthetic ground truth, photometric and geometric consistency, or view synthesis). Figure~\ref{fig:p3_s6_4_3_method_overview} groups them into five families along these axes.

\paragraph{Variational and energy-based methods.}
\citet{vedula1999three} introduced the term and the problem: recover 3D structure and 3D motion jointly from a calibrated camera array. \citet{huguet2007variational} formulated it using stereo image pairs at two consecutive times, solving disparity and flow jointly in one variational energy; \citet{wedel2011stereoscopic} decoupled the two into a stereo stage and a flow stage, which made real-time estimation feasible; and \citet{basha2013multi} returned to multi-view input under a view-centered energy. The strongest classical formulation is the piecewise-rigid scene model of \citet{vogel20153d} (PRSF/PRSM), which segments the scene into rigidly moving planar patches and was the leading entry on KITTI Scene Flow at the benchmark's release.

\paragraph{End-to-end stereo and temporal CNNs.}
\citet{mayer2016large} recast disparity and flow as supervised regression (DispNet, FlowNet) and released FlyingThings3D, Monkaa, and Driving, the synthetic corpora that made this training regime possible. SENSE~\cite{jiang2019sense} shares one encoder across disparity, optical flow, occlusion, and semantic heads; DRISF~\cite{ma2019deep} keeps the piecewise-rigid prior of Vogel \etal\ but implements it as an instance-level decoder over learned features. DRISF surpassed both listed variational baselines on KITTI Scene Flow, while SENSE improved on OSF but not PRSM (Table~\ref{tab:scene_flow_performance}), shifting the bottleneck from optimization to supervised data.

\paragraph{Recurrent geometric refinement.}
RAFT~\cite{teed2020raft} introduced all-pairs correlation volumes with GRU-based iterative updates for 2D optical flow; RAFT-3D~\cite{teed2021raft} uses image and depth pairs to iteratively refine a per-pixel SE(3) field with rigid-motion embeddings. Depth can be obtained from stereo or RGB-D input; it is an input to the refinement rather than necessarily inferred jointly from appearance. Its scene-flow results are included in Table~\ref{tab:scene_flow_performance}.

\paragraph{Monocular and self-supervised methods.}
Mono-SF~\cite{brickwedde2019mono} was the first monocular entry on KITTI Scene Flow, resolving scale with a learned single-view depth prior. Self-Mono-SF~\cite{hur2020self} jointly predicts disparity and scene flow with a single decoder, using unlabeled stereo video during training and monocular image pairs at test time; stereo geometry provides the scale cue, while photometric and 3D consistency losses provide self-supervision. A separate test-time optimization approach for LiDAR point clouds (Neural Scene Flow Prior~\cite{li2021neuralsfp}) lies outside the image-based scope of this subsection.

\paragraph{Implicit-field and view-synthesis-coupled methods.}
A parallel line never targets scene flow as output but obtains a dense 4D motion field as a byproduct of dynamic view synthesis. Neural Scene Flow Fields~\cite{li2021neural} trains a dynamic NeRF together with an explicit scene-flow MLP; D-NeRF~\cite{pumarola2021d} represents motion as a deformation field into a canonical frame; Nerfies~\cite{park2021nerfies} extends this to casually captured monocular video; and 4D Gaussian Splatting~\cite{wu20244d} attaches time-parameterized motion to explicit Gaussian primitives. These methods are evaluated by rendering quality rather than by motion error, a split returned to under Challenges below.

\subsubsection{Datasets and Metrics}

Table~\ref{tab:scene_flow_datasets} lists the image-based benchmarks. Synthetic sets provide dense, exact ground truth at scale (FlyingThings3D as the standard pretraining corpus; Sintel and Spring for film-grade motion, with Spring at HD resolution and super-resolved ground truth; TartanAir for environment diversity), while KITTI Scene Flow 2015~\cite{menze2015object}, built on the KITTI suite~\cite{KITTI-GeigerLU12}, is the only real-world benchmark and carries sparse LiDAR-derived labels on 200 training and 200 test scenes. LiDAR-based scene-flow splits of nuScenes~\cite{caesar2020nuscenes}, Waymo Open~\cite{WaymoOpen-SunKDCPTGZCCVHN20}, and Argoverse~2~\cite{wilson2023argoverse}, and the event-camera flow of DSEC~\cite{gehrig2021dsec}, serve the point-cloud and event communities and are not image-based benchmarks in the sense used here.

\begin{table}[t]
\centering
\fontsize{9}{11}\selectfont
\setlength{\tabcolsep}{1.0mm}
\caption{Image-based scene-flow benchmarks. Counts are indicative of the commonly used releases.}
\label{tab:scene_flow_datasets}
\begin{tabular}{L{0.26\textwidth} c C{0.10\textwidth} C{0.11\textwidth} C{0.14\textwidth} L{0.24\textwidth}}
\toprule
\textbf{Dataset} & \textbf{Year} & \textbf{Type} & \textbf{Input} & \textbf{Scale} & \textbf{Ground truth} \\
\midrule
KITTI Scene Flow 2015~\cite{menze2015object} & 2015 & Real driving & Stereo video & 200 train / 200 test scenes & Sparse disparity ($t$, $t{+}1$) + flow from LiDAR and CAD fits \\
FlyingThings3D / Monkaa / Driving~\cite{mayer2016large} & 2016 & Synthetic & Stereo video & $\sim$25K / 8.6K / 4.4K frames & Dense disparity, flow, disparity change \\
MPI Sintel~\cite{sintel} & 2012 & Synthetic film & Stereo video & 1,041 train frames, 23 seq. & Dense flow, depth, disparity \\
Spring~\cite{Spring} & 2023 & Synthetic film & Stereo video & $\sim$6K frames, 47 seq., 1920$\times$1080 & Dense stereo, flow, scene flow (super-resolved) \\
TartanAir~\cite{TartanAir} & 2020 & Synthetic & Stereo video & $>$1M frames, 30 environments & Dense depth, flow, pose; indoor and outdoor, weather \\
\bottomrule
\end{tabular}
\end{table}

KITTI scores scene flow in image space. A pixel's disparity in frame $t$ (D1), its disparity in frame $t{+}1$ (D2), and its optical flow (Fl) are each counted as outliers when the error exceeds both 3\,px and 5\% of the ground-truth value, and a scene-flow outlier (SF) is any pixel that fails at least one of the three; the benchmark reports outlier percentages over all pixels and separately over background and foreground. Point-level evaluation on synthetic data instead uses the 3D end-point error
\begin{equation}
    \mathrm{EPE}_{3D} = \frac{1}{N} \sum_{i=1}^{N} \| \hat{\mathbf{v}}_i - \mathbf{v}_i \|_2,
\end{equation}
with $\hat{\mathbf{v}}_i$ and $\mathbf{v}_i$ the predicted and true motion of point $i$, together with accuracy ratios at fixed thresholds (the fraction of points with error below 0.05\,m or 5\%, and below 0.1\,m or 10\%) and an outlier ratio (error above 0.3\,m or 10\%), conventions inherited from point-cloud scene flow.

Table~\ref{tab:scene_flow_performance} collects KITTI Scene Flow 2015 test results for the explicit-output families. RAFT-3D reduces SF-all error from 8.97\% (PRSM) to 5.77\%; the table does not separate foreground and background results. The monocular rows are a different regime: Mono-SF, with a learned depth prior, sits at 23\% SF-all, and Self-Mono-SF without any ground truth at 49.5\% (33.9\% after fine-tuning), the price of giving up stereo. The implicit-field family has no row because it reports rendering quality (PSNR, LPIPS) on its own scene sets rather than motion error.

\begin{table}[!t]
    \centering
    \fontsize{9}{11}\selectfont
    \setlength{\tabcolsep}{2.2mm}
    \caption{KITTI Scene Flow 2015 test set, outlier percentages (lower is better) over all pixels. Stereo rows from the RAFT-3D comparison~\cite{teed2021raft}; monocular rows from the Mono-SF~\cite{brickwedde2019mono} and Self-Mono-SF~\cite{hur2020self} papers (``ft'' = fine-tuned on KITTI). D1 / D2: disparity in frames $t$ / $t{+}1$; Fl: optical flow; SF: scene flow.}
    \label{tab:scene_flow_performance}
    \begin{tabular}{L{0.30\textwidth} L{0.26\textwidth} C{0.07\textwidth} C{0.07\textwidth} C{0.07\textwidth} C{0.07\textwidth}}
    \toprule
    \textbf{Method} & \textbf{Family} & \textbf{D1-all} & \textbf{D2-all} & \textbf{Fl-all} & \textbf{SF-all} \\
    \midrule
    \rowcolor{lightgrey} \multicolumn{6}{l}{$\bullet$ \textit{Stereo}}\\
    OSF~\cite{menze2015object}          & Energy-based, rigid objects   & 5.79 & 7.77 & 7.83 & 10.23 \\
    PRSM~\cite{vogel20153d}             & Energy-based, piecewise rigid & 4.27 & 6.79 & 6.68 & 8.97 \\
    SENSE~\cite{jiang2019sense}         & Shared-encoder CNN            & 2.22 & 5.89 & 7.64 & 9.55 \\
    DRISF~\cite{ma2019deep}             & CNN + rigid instances         & 2.55 & 4.04 & 4.73 & 6.31 \\
    RAFT-3D~\cite{teed2021raft}         & Recurrent SE(3) refinement    & 1.81 & 3.67 & 4.29 & 5.77 \\
    \midrule
    \rowcolor{lightgrey} \multicolumn{6}{l}{$\bullet$ \textit{Monocular}}\\
    Mono-SF~\cite{brickwedde2019mono}   & Learned depth prior           & 16.32 & 19.59 & 12.77 & 23.08 \\
    Self-Mono-SF~\cite{hur2020self}     & Self-supervised               & 34.02 & 36.34 & 23.54 & 49.54 \\
    Self-Mono-SF-ft~\cite{hur2020self}  & Semi-supervised, ft           & 22.16 & 25.24 & 15.91 & 33.88 \\
    \bottomrule
    \end{tabular}
\end{table}

\subsubsection{Challenges}

The supervision landscape is lopsided. Every dense-ground-truth benchmark in Table~\ref{tab:scene_flow_datasets} is synthetic and stereo, so monocular methods are scored only against KITTI's sparse labels, and models pretrained on FlyingThings3D or Sintel still lose accuracy on real driving footage after fine-tuning, most visibly on foreground motion where synthetic random-object trajectories differ from vehicle kinematics. All current benchmarks also emphasize rigid or random-object motion; dense ground truth for non-rigid humans, hands, and animals does not exist at scale, which leaves the family of implicit-field methods that handle such motion best without a quantitative yardstick.

That yardstick problem is the second challenge. Explicit-output methods~\cite{vogel20153d,jiang2019sense,teed2021raft} are scored by KITTI's outlier rates; implicit-field methods~\cite{li2021neural,pumarola2021d,wu20244d} by PSNR and LPIPS of rendered views. No protocol evaluates both on the same input, so the two paradigms are compared only qualitatively, and whether a motion field is good depends on whether scene flow is the end (the KITTI leaderboard) or a means (reconstruction, novel-view synthesis, physical reasoning). The cost profiles diverge as well: RAFT-3D needs several GRU iterations per frame, variational methods minutes per frame, and implicit fields minutes to hours per scene.

Foundation priors are reshaping the monocular line. Depth Anything~\cite{DepthAnything} and DINOv2~\cite{dinov2} supply the per-pixel depth and semantic features that Self-Mono-SF had to learn from scratch, and SAM~2~\cite{ravi2024sam2} segmentation together with the object-pose estimators of Section~\ref{sec:object_pose} makes an object-centric formulation (a per-object SE(3) trajectory plus residual non-rigid motion) practical, a bridge to the tracking problem of Section~\ref{sec:object_tracking}. As implicit-field scene flow merges into 4D reconstruction, with 4D Gaussian Splatting~\cite{wu20244d} the early example, evaluation is likely to follow the downstream task rather than the KITTI leaderboard.
% TODO: anchor §8.x once §8 reconstruction subsections stabilize








\subsection{Object Tracking}
\label{sec:object_tracking}

Tracking localizes entities in every frame and associates the detections across frames so that each entity keeps one identity through the sequence. Where detection and segmentation describe a single instant, tracking yields trajectories, the representation that autonomous driving~\cite{luo2021exploring,chiu2021probabilistic}, surveillance~\cite{ojha2015imageobjectsurveillance}, AR/VR~\cite{chen2021hightracking}, and robotics~\cite{Deitke2020RoboTHOR} build on.

\begin{figure}[!ht]
    \centering
     \includegraphics[width=0.95\linewidth]{figures/part3/section6/object_tracking.pdf}
    \caption{Object-tracking subtasks: (a) single object tracking (SOT) from an initial box, (b) multiple object tracking (MOT) with identity-consistent boxes, (c) multiple object tracking and segmentation (MOTS) with per-frame masks, (d) 3D multi-object tracking (3D-MOT) of oriented boxes from LiDAR or camera input, (e) pose and landmark tracking, and (f) panoptic tracking with semantic classes for stuff and persistent instance IDs for things. Repeated IDs denote the same entity across frames; arrows indicate increasing time. Scene images and overlaid states are illustrative.}
    \label{fig:tracking_task}
\end{figure}

\subsubsection{Task Definition}
Given a frame sequence $\{I_t\}_{t=1}^{T}$ (a video, or RGB-D, LiDAR, or multi-view streams), a tracker returns a set of trajectories $\{\mathcal{T}_k\}$, each $\mathcal{T}_k = \{(t, b_t^k)\}$ a time-indexed series of states $b_t^k$ with one identity $k$ held fixed over the sequence. The state is a bounding box in the basic setting and a segmentation mask, a 3D box, or a set of keypoints in the variants.

The variants, illustrated in Figure~\ref{fig:tracking_task}, differ in how the targets are specified and what state is tracked. Single object tracking (SOT)~\cite{fan2019lasot,muller2018trackingnet} follows one target given its first-frame box; multiple object tracking (MOT)~\cite{milan2016mot16,dendorfer2020mot20,yu2020bdd100k} tracks every instance of given categories, typically pedestrians and vehicles, under identity constraints; MOTS~\cite{voigtlaender2019mots} and video object segmentation replace boxes with masks; 3D-MOT~\cite{AB3DMOT,yin2021center,caesar2020nuscenes,WaymoOpen-SunKDCPTGZCCVHN20} outputs world-coordinate 3D boxes from LiDAR, RGB-D, or multi-view input; pose tracking~\cite{andriluka2018posetrack,iqbal2017posetrack,samanta2016data} and landmark tracking~\cite{zafeiriou20173d} follow skeletal joints or facial keypoints; and panoptic tracking combines pixel-level semantic classes for stuff with temporally consistent instance identities for things. Category-specific trackers for pedestrians~\cite{stadler2021improving}, vehicles~\cite{van2001real}, or animals~\cite{zhang2023animaltrack} specialize any of these.


\subsubsection{Methodology}
Figure~\ref{fig:task_tracking_method} organizes the methods into four families around image-plane box and mask tracking (SOT, MOT, MOTS, VOS). Pose, landmark, and panoptic tracking reuse the same propagate-or-associate machinery with a different state representation and are covered at the benchmark level in Table~\ref{tab:video_object_tracking_dataset}.

\paragraph{Propagation-based methods.}
Because a target moves little between consecutive frames, its previous state is a strong prior for the current one, and the oldest tracking paradigm exploits this by predicting the next state from the current one and correcting it against the new frame. Correlation filters realize this in the frequency domain: MOSSE~\cite{bolme2010visual} learns a minimum-output-sum-of-squared-error filter from a single frame and updates it online at hundreds of frames per second; KCF~\cite{henriques2014high} and MKCF~\cite{tang2015multi} add kernelization and multiple-kernel learning, respectively; ECO~\cite{DanelljanBKF17} adds factorized convolution operators, a compact generative sample model, and conservative updates. These trackers are fast enough for embedded use but drift under large appearance change and clutter. MDNet~\cite{NamH16} answered with online discriminative learning, a shared CNN with per-sequence domain branches that separate target from background, at a high per-frame cost. Siamese trackers moved the representation learning offline: SiamFC~\cite{BertinettoVHVT16} correlates a template embedding with the search-region embedding in one forward pass, and SiamRPN~\cite{LiYWZH18}, ATOM~\cite{DanelljanBKF19}, and DiMP~\cite{BhatDGT19} add region-proposal regression, IoU-prediction refinement, and a discriminatively learned prediction head. Transformer trackers then replaced correlation with attention: TransT~\cite{ChenYZ0YL21} fuses template and search features with self- and cross-attention, STARK~\cite{0002PF0L21} predicts the box with an encoder--decoder that also carries temporal state, and OSTrack~\cite{YeCMSC22} joins feature extraction and relation modeling in a single stream over concatenated template and search tokens, the current baseline design.

Mask propagation applies the same idea at the pixel level for video object segmentation~\cite{pont20172017,xu2018youtube}, with memory of past frames replacing a single template. XMem~\cite{ChengS22XMem} keeps multi-scale spatiotemporal memories that remain stable over videos of more than 10{,}000 frames; AOT~\cite{YangWY21aot} and DeAOT~\cite{YangY22adeaot} encode object identities as tokens so that many objects propagate in one pass; and SAM~2~\cite{ravi2024sam2} adds a streaming memory to the Segment Anything Model~\cite{kirillov2023sam}, giving promptable tracking that transfers across VOS benchmarks.


\begin{figure}[!t]
    \centering
    \includegraphics[width=0.99\linewidth]{figures/part3/section6/object_tracking_method.pdf}
    \caption{Four method families for object tracking: (i) propagation-based, (ii) detect-then-associate, (iii) joint detection--association, and (iv) universal models with a shared backbone and task heads for SOT, MOT, VOS, and MOTS.}
    \label{fig:task_tracking_method}
\end{figure}

\paragraph{Detect-then-associate methods.}
The MOT standard pipeline detects objects in every frame and then solves an assignment problem between detections and existing tracks. SORT~\cite{bewley2016sort} predicts track positions with a Kalman filter~\cite{welch1995kalman} and matches by IoU with the Hungarian algorithm~\cite{Kuhn10}; DeepSORT~\cite{WojkeBP17} adds an appearance embedding so that tracks survive occlusion. Track-RCNN~\cite{voigtlaender2019mots}, JDE~\cite{WangZLLW20jde}, and FairMOT~\cite{ZhangWWZL21FairMOT} train the re-identification embedding jointly with the detector; association remains a separate step, which keeps them in this family. All of these discard low-confidence detections, and since occluded targets are exactly those with low scores, trajectories fragment where they are hardest to keep. ByteTrack~\cite{ZhangSJYWYLLW22} associates nearly every detection and uses motion similarity to the existing tracks to reject the false positives among the low-score boxes, a change in association alone that raised MOT17 HOTA by several points (Table~\ref{tab:tracking_performance}). In 3D, the same pipeline reappears: AB3DMOT~\cite{AB3DMOT} runs a 3D Kalman filter with Hungarian matching on LiDAR detections, and CenterPoint~\cite{yin2021center} tracks 3D object centers through predicted velocities, with nuScenes~\cite{caesar2020nuscenes} using AMOTA/AMOTP and Waymo~\cite{WaymoOpen-SunKDCPTGZCCVHN20} using MOTA/MOTP.

\paragraph{Joint detection--association methods.}
End-to-end trackers fold association into the detector. CenterTrack~\cite{ZhouKK20} represents objects as center points and regresses their offset to the previous frame, so that association is a local lookup; TrackFormer~\cite{MeinhardtKLF22} introduces tracking-by-attention, in which track queries carried from the previous frame and new object queries are decoded jointly by a Transformer; MOTR~\cite{zeng2022motr} extends DETR~\cite{carion2020end} with track queries that persist across the whole sequence and are supervised with a sequence-level loss, which is what gives it its advantage on DanceTrack~\cite{SunCJYBKL22}, where appearance is uninformative and motion is non-linear. The price is training cost and data: these models need long annotated sequences and remain behind detect-then-associate trackers on MOT17 (Table~\ref{tab:tracking_performance}).

\paragraph{Universal models.}
A single model for SOT, MOT, VOS, and MOTS is the most recent direction. UniTrack~\cite{wang2021different} showed that one self-supervised appearance representation serves propagation and association tasks alike; Unicorn~\cite{yan2022unicorn} trains one network end to end for all four tasks; UTT~\cite{MaSZ0XYY22} unifies SOT and MOT in a Transformer; UniVS~\cite{li2024univs} casts every video-segmentation task as prompt-driven segmentation; and OneTracker~\cite{HongYZLZGJCLCZ24} adapts a foundation tracker to RGB, depth, thermal, and event inputs with light adapters. In practice, SAM~\cite{kirillov2023sam} combined with a memory tracker (XMem, DeAOT) already provides promptable multi-object segmentation and tracking from points, boxes, scribbles, or text, and methods such as UASTrack~\cite{UASTrack} and Un-Track~\cite{WuZRV0PGT24UNTrack} extend unified tracking to RGB paired with depth, thermal, or event inputs. Language enters as a specification signal as well: LaSOT~\cite{fan2019lasot} pairs every sequence with a sentence, TNL2K~\cite{wang2021tnl2k} is built for text-initialized tracking, and LaMOT~\cite{li2025lamot} brings language-guided specification to the multi-object setting.

\begin{table}[!ht]
\centering
\caption{Datasets for object tracking with bounding boxes (Bbox), segmentation masks (Mask), or human poses. ``Frames / Boxes'' gives annotated frames, or annotated boxes where reported (MOT16/17, TrackingNet, YouTube-BB). MOT17 reuses its 14 videos with three public detection sets. For pose datasets, the annotation column gives the setting and the number of body landmarks.}
\label{tab:video_object_tracking_dataset}
\fontsize{9}{11}\selectfont
\setlength{\tabcolsep}{2.2mm}
\begin{tabular}{lcccccc}
\toprule
 \textbf{Name} & \textbf{Videos} & \textbf{Frames / Boxes} & \textbf{Categories} & \textbf{Year} & \textbf{Attributes} & \textbf{Annotations}\\
 \midrule
\rowcolor{lightgrey} \multicolumn{7}{l}{$\bullet$ \textit{Single Object Tracking}} \\
VIVID~\cite{collins2005open} &   9 &  16,274 &  - & 2005 & - & Bbox\\
OTB100~\cite{wu2013online} &   98 &  58,610 &  -  & 2015 & 11 & Bbox\\
ALOV300~\cite{smeulders2013visual} &   314 &  89,364 &   64 & 2013 & - & Bbox\\
TC128~\cite{liang2015encoding} &   128 &  55,652 &  - &  2015  & 11 & Bbox\\
VOT16~\cite{roffo2016visual} &   60 &  21,455 &  - & 2016 & - & Bbox\\
DAVIS-16~\cite{perazzi2016benchmark} & 50 &  3,455 & 4 & 2016 & 15& Mask \\
UAV20L~\cite{benchmark2016uav} &   20 &  58,670 &   -  & 2016 & 12 & Bbox\\
UAV123~\cite{benchmark2016uav} &   123 &  113,476 &   9  & 2016 & 12 & Bbox\\
NUS-PRO~\cite{nus_pro} &   365 &  135,305 &  8 & 2016 & - & Bbox\\
VOT17~\cite{alatan2017visual} &   60 &  21,356 & 24 & 2017 & - & Bbox\\
YouTube-BB~\cite{real2017youtube} &   380,000 &  5.6M &  23  & 2017  & - & Bbox\\
NFS~\cite{kiani2017need} &   100 &  383,000 &   17  & 2017 & 9 & Bbox\\
TrackingNet~\cite{muller2018trackingnet} &   30,132 &  14.2M &   27  & 2018 & 15  & Bbox\\
LaSOT~\cite{fan2019lasot}  &  1,400  & 3.52M &  70 & 2019 & 14 & Bbox\\
TNL2K~\cite{wang2021tnl2k} &  2,000  &  1.24M &  - & 2021 & 17 & Bbox\\
GOT-10k~\cite{huang2021got}   & 10,000 & 1.5M & 563  & 2021 & 6 & Bbox\\
\midrule
\rowcolor{lightgrey} \multicolumn{7}{l}{$\bullet$ \textit{Multiple Object Tracking}} \\
MOT16~\cite{milan2016mot16} &   14 &  292,733 & 1 & 2016  & - & Bbox\\
MOT17~\cite{milan2016mot16} &   14 &  564,228 & 1 &   2017 & -  & Bbox\\
BDD100K~\cite{yu2020bdd100k} &   1,600 &  318K &  8  & 2020  & - & Bbox\\
\midrule
\rowcolor{lightgrey} \multicolumn{7}{l}{$\bullet$ \textit{Video Object Segmentation}} \\
DAVIS-17~\cite{pont20172017} & 150 & 10,459 & - & 2017 &-  & Mask\\
YouTube-VOS~\cite{xu2018youtube} & 4,453 & - &  94 & 2018  & -  & Mask\\
\midrule
\rowcolor{lightgrey} \multicolumn{7}{l}{$\bullet$ \textit{RGB-D Object Tracking}} \\
PTB~\cite{song2013tracking} &  100 &  21,542 & - & 2013 & 5 &  Bbox \\
STC~\cite{xiao2017robust} &  36 & 9,009 & -  & 2017 & 12 &  Bbox \\
CDTB~\cite{lukezic2019cdtb} &  80 &  101,956 & - & 2019 &  13 &  Bbox  \\
DepthTrack~\cite{yan2021depthtrack} & 200 &  294,591 & - & 2021  &  15 &  Bbox \\
ARKitTrack~\cite{zhao2023arkittrack} & 300 & 229.7K & 144 & 2023 & 16 & Mask \\
RGBD1K~\cite{zhu2023rgbd1k} &  1,050&  717,900 & - & 2023 &  15 &  Bbox  \\
\midrule
\rowcolor{lightgrey} \multicolumn{7}{l}{$\bullet$ \textit{Human Pose Tracking}} \\
VideoPose2.0~\cite{sapp2011parsing} & 44 & 1,286 &- &  2011 & - & Single-person (-) \\
Poses in the Wild~\cite{cherian2014mixing} & 30 & 830 & - & 2014 & - & Single-person (7) \\
ICDPose~\cite{samanta2016data} & 60 & 2,700 & - & 2016 &- &  Single-person (-)\\
CMU Panoptic~\cite{joo2017panoptic} & 65 & 297K & - & 2017 & - & Multi-person (-) \\
Multi-Person PoseTrack~\cite{iqbal2017posetrack} & 60 & - & - & 2017 &- &  Multi-person (14)\\
PoseTrack18~\cite{andriluka2018posetrack} & 1,138 & 66,374 & - & 2018 &- &  Multi-person (15)\\
\bottomrule
\end{tabular}
\end{table}

\subsubsection{Datasets and Metrics}

Table~\ref{tab:video_object_tracking_dataset} groups the benchmarks by what is tracked. The SOT line runs from short, category-agnostic clips (OTB100, TC128, the VOT challenges) to large, long-term, category-diverse training sets (TrackingNet, LaSOT, GOT-10k, whose test classes are disjoint from training) and to condition-specific tests for aerial footage (UAV123), high frame rate (NFS), and language initialization (TNL2K). MOT benchmarks are pedestrian-only in MOT16/17 and MOT20 and multi-class driving footage in BDD100K; the 3D-MOT counterparts are the nuScenes and Waymo tracking tasks. DAVIS and YouTube-VOS supply mask-level ground truth, the latter with categories unseen at training time; the RGB-D sets add depth and have grown from tens of sequences (PTB, STC) to over a thousand (RGBD1K); and pose tracking moved from single-person upper-body clips (VideoPose2.0, Poses in the Wild, ICDPose) to multi-person full-body sequences with person identities (PoseTrack18, CMU Panoptic).

Each line has its own metrics. SOT uses the success plot: with $\mathrm{SR}(\tau)$ the fraction of frames whose predicted-box IoU exceeds $\tau$, the headline number is its area under the curve,
\begin{equation}
    \mathrm{AUC} = \int_0^1 \mathrm{SR}(\tau)\, d\tau ,
\end{equation}
reported with normalized precision (the fraction of frames whose center error, divided by the target size, is below a threshold) on LaSOT and TrackingNet, and with average overlap (mean IoU over frames) and $\mathrm{SR}_{0.5}$, $\mathrm{SR}_{0.75}$ on GOT-10k; the VOT challenges use expected average overlap (EAO), which folds accuracy and failure rate into one expectation~\cite{roffo2016visual}. MOT scores detection and identity jointly. MOTA~\cite{bernardin2008evaluating} charges false negatives, false positives, and identity switches against the number of ground-truth objects,
\begin{equation}
    \mathrm{MOTA} = 1 - \frac{\mathrm{FN}+\mathrm{FP}+\mathrm{IDS}}{\mathrm{GT}} ,
\end{equation}
and so is dominated by detection; IDF1 is the F1 score of identity-level true positives, $\mathrm{IDF1} = 2\,\mathrm{IDTP} / (2\,\mathrm{IDTP} + \mathrm{IDFP} + \mathrm{IDFN})$, and emphasizes identity preservation; HOTA~\cite{luiten2021hota}, $\sqrt{\mathrm{DetA}\cdot\mathrm{AssA}}$ averaged over localization thresholds, balances the two and has become the primary MOT metric. VOS reports the Jaccard index $\mathcal{J} = |M \cap G| / |M \cup G|$ of predicted and true masks and the contour F-measure $\mathcal{F}$ from boundary precision and recall~\cite{perazzi2016benchmark}, averaged as $\mathcal{J}\&\mathcal{F}$.

Table~\ref{tab:tracking_performance} reports the methods of the Methodology above on the benchmarks each line treats as standard. Within SOT the step from a discriminative head over correlation features (DiMP) to attention-based fusion (TransT, STARK) is worth about ten points of LaSOT AUC, and the one-stream design of OSTrack adds six more over TransT. On MOT17 the detect-then-associate trackers lead, and ByteTrack's gain over FairMOT comes from association alone; the end-to-end trackers (CenterTrack, TrackFormer, MOTR) trail on MOT17 but MOTR leads ByteTrack by 6.5 HOTA on DanceTrack~\cite{zeng2022motr}, where appearance cues are weak. VOS accuracy on DAVIS-17 has moved from the mid-80s (XMem, DeAOT) to above 90 $\mathcal{J}\&\mathcal{F}$ with SAM~2, which is also the first method trained on a corpus far larger than the benchmarks themselves.

\begin{table}[!t]
    \centering
    \fontsize{9}{11}\selectfont
    \setlength{\tabcolsep}{2.2mm}
    \caption{Representative trackers on the standard benchmarks of each line. SOT: AUC on LaSOT and TrackingNet, AO on GOT-10k (from the OSTrack comparison~\cite{YeCMSC22}). MOT: MOT17 test, private detections (from the MOTR~\cite{zeng2022motr} and ByteTrack~\cite{ZhangSJYWYLLW22} comparisons). VOS: $\mathcal{J}\&\mathcal{F}$ on DAVIS-17 val and $\mathcal{G}$ on YouTube-VOS 2019 val (from the SAM~2 comparison~\cite{ravi2024sam2}).}
    \label{tab:tracking_performance}
    \begin{tabular}{L{0.30\textwidth} L{0.26\textwidth} C{0.10\textwidth} C{0.11\textwidth} C{0.10\textwidth}}
    \toprule
    \rowcolor{lightgrey} \multicolumn{5}{l}{$\bullet$ \textit{Single object tracking}}\\
    \textbf{Method} & \textbf{Design} & \textbf{LaSOT AUC} & \textbf{TrackingNet AUC} & \textbf{GOT-10k AO} \\
    \midrule
    DiMP~\cite{BhatDGT19}         & Discriminative prediction head   & 56.9 & 74.0 & 61.1 \\
    TransT~\cite{ChenYZ0YL21}      & Transformer feature fusion       & 64.9 & 81.4 & 67.1 \\
    STARK~\cite{0002PF0L21}        & Encoder--decoder, temporal state & 67.1 & 82.0 & 68.8 \\
    OSTrack-384~\cite{YeCMSC22}    & One-stream Transformer           & 71.1 & 83.9 & 73.7 \\
    \midrule
    \rowcolor{lightgrey} \multicolumn{5}{l}{$\bullet$ \textit{Multiple object tracking (MOT17 test)}}\\
    \hiderowcolors \textbf{Method} & \textbf{Family} & \textbf{HOTA} & \textbf{MOTA} & \textbf{IDF1} \\
    \midrule
    CenterTrack~\cite{ZhouKK20}            & Joint, center offsets        & 52.2 & 67.8 & 64.7 \\
    TrackFormer~\cite{MeinhardtKLF22}      & Joint, tracking-by-attention & --   & 74.1 & 68.0 \\
    MOTR~\cite{zeng2022motr}               & Joint, persistent queries    & 57.8 & 73.4 & 68.6 \\
    FairMOT~\cite{ZhangWWZL21FairMOT}      & Detect + joint Re-ID         & 59.3 & 73.7 & 72.3 \\
    ByteTrack~\cite{ZhangSJYWYLLW22}       & Detect + low-score association & 63.1 & 80.3 & 77.3 \\
    \midrule
    \rowcolor{lightgrey} \multicolumn{5}{l}{$\bullet$ \textit{Video object segmentation}}\\
    \hiderowcolors \textbf{Method} & \textbf{Design} & \textbf{DAVIS-17 $\mathcal{J}\&\mathcal{F}$} & \textbf{YT-VOS 2019 $\mathcal{G}$} & \\
    \midrule
    XMem~\cite{ChengS22XMem}       & Multi-store memory               & 86.0 & 85.6 & \\
    DeAOT~\cite{YangY22adeaot}     & Identity tokens, decoupled propagation & 86.2 & 86.1 & \\
    SAM~2~\cite{ravi2024sam2}      & Promptable, streaming memory     & 90.7 & 89.3 & \\
    \bottomrule
    \end{tabular}
\end{table}

\subsubsection{Challenges}
Identity under occlusion remains the central failure. Motion models and appearance embeddings bridge short occlusions, but long occlusions and dense crowds still produce identity switches and fragmented tracks~\cite{ZhangSJYWYLLW22,SunCJYBKL22}, and memory-based propagation drifts or forgets over long videos even when each step is accurate~\cite{ChengS22XMem}. HOTA and IDF1 measure the outcome but not robustness to occlusion length or the calibration of track confidence, so the metrics reward the easy part of the problem.

Trackers are also tuned to a few curated benchmarks and transfer poorly to UAV footage, thermal and event imagery, adverse weather, and multi-camera surveillance; foundation-model trackers such as OneTracker~\cite{HongYZLZGJCLCZ24} and the RGB-X methods UASTrack~\cite{UASTrack} and Un-Track~\cite{WuZRV0PGT24UNTrack} address this directly but inherit the biases of their pretrained backbones. Long streams add memory and compute limits on top of drift.

The accuracy of Transformer trackers comes at a cost that keeps them off embedded hardware, while the lightweight correlation-filter line gives up robustness; and the unified models that would amortize that cost across tasks meet negative transfer between objectives with different demands, mask precision against identity association in particular~\cite{yan2022unicorn}. Language and segmentation priors widen what a tracker can be asked to do, at the price of prompt sensitivity that current benchmarks do not measure.







\subsection{Temporal Localization}
\label{sec:temporal_localization}

Temporal localization answers \emph{when} something happens in a long, untrimmed video. The condition can be a category from a closed set, which gives temporal action localization (TAL) as treated in Section~\ref{sec:action_detection}, or a free-form sentence, which gives temporal sentence grounding (TSG): return the segment that best matches ``the person opens the microwave''. This subsection sets the two side by side, with the weight on the text-conditioned problem, which turns localization from exhaustive detection into retrieval and is what video search and interactive navigation~\cite{hendricks2017mcn,gao2017tall,krishna2017densecap} rely on.

\begin{figure}[!ht]
    \centering
    \includegraphics[width=\linewidth]{figures/part3/section6/temporal_localization.pdf}
    \caption{Temporal localization methods. Top row (TAL, category-conditioned): (i) window-based, (ii) two-stage proposal-based, (iii) one-stage anchor-free, and (iv) Transformer set prediction. Bottom row (TSG, text-conditioned): (a) cross-modal matching, (b) 2D temporal-map reasoning with a proposal-graph alternative, and (c) Transformer set prediction. Hatched map cells are invalid start--end pairs; graph nodes represent candidate moments. Coarse-to-fine corpus retrieval is shown separately as a scope extension, not a mandatory stage of set prediction. Scene frames and scores are illustrative.}
    \label{fig:p3_s6_4_2_method_overview}
\end{figure}

\subsubsection{Task Definition}
\textbf{Temporal Action Localization (TAL)} takes an untrimmed video $V=\{x_t\}_{t=1}^{T}$, where $x_t$ may be a visual, audio, or multimodal feature, and detects \emph{all} instances of $K$ predefined action classes, returning $\mathcal{A}=\{a_i\}_{i=1}^{N}$ with $a_i=(t_i^{\mathrm{start}},\, t_i^{\mathrm{end}},\, c_i,\, s_i)$: start and end times, a class label $c_i\in\{1,\ldots,K\}$, and a confidence $s_i$ (the $(t_s, t_e, c)$ tuples of Section~\ref{sec:action_detection}).

\textbf{Temporal Sentence Grounding (TSG)}, also called natural language video localization or video moment retrieval, pairs $V$ with a sentence $q$ and returns a ranked list of moments $\mathcal{M}(q,V)=\{m_j\}_{j=1}^{M}$, $m_j=(\tau_j^{\mathrm{start}},\,\tau_j^{\mathrm{end}},\, r_j)$, with boundaries and a relevance score $r_j$; usually one segment is the answer. TAL is thus an exhaustive detection task over a closed label set, TSG a targeted retrieval task guided by open-ended language.

Both are studied under several settings, which Figure~\ref{fig:p3_s6_4_2_task_settings} lays side by side. Offline methods see the whole video; online methods infer causally over the streamed prefix and are judged on low-latency detection of action starts~\cite{shou2018odas,gao2019startnet}. Supervision ranges from frame-accurate boundaries to video-level labels only (weakly supervised TAL via multiple-instance learning, Section~\ref{sec:action_detection}; weakly supervised TSG by aligning the two modalities during training~\cite{ma2020vlanet}), with self-supervised video pretraining~\cite{tong2022videomae,wang2023videomaev2} and open-vocabulary or zero-shot variants~\cite{liberatori2024t3al} reducing the labeled data needed. TSG additionally varies in scope, from a single video~\cite{gao2017tall,hendricks2017mcn} to corpus-level retrieval that first finds the video and then the moment (VCMR)~\cite{lei2020tvr}, and in output, with QVHighlights~\cite{lei2021qvhighlights} asking for a saliency ranking of highlight clips alongside the grounded moment.

\begin{figure}[!ht]
    \centering
    \includegraphics[width=\linewidth]{figures/part3/section6/temporal_localization_task_settings.pdf}
    \caption{Six settings of temporal localization. Top row (TAL): (a) offline detection of $K$ classes; (b) online detection of action starts~\cite{shou2018odas,gao2019startnet}; (c) weakly supervised training from video-level labels. Bottom row (TSG): (d) single-video sentence grounding~\cite{gao2017tall,hendricks2017mcn}; (e) corpus-level moment retrieval~\cite{lei2020tvr}; (f) joint grounding and highlight detection~\cite{lei2021qvhighlights}. The online panel separates event onset, detection time, and the current observation boundary; future frames are unavailable. Dashed intervals in the weak-supervision panel denote inferred hypotheses, not boundary annotations. Scene examples, confidence values, and saliency ranks are illustrative.}
    \label{fig:p3_s6_4_2_task_settings}
\end{figure}

\subsubsection{Methodology}
Figure~\ref{fig:p3_s6_4_2_method_overview} traces both tasks along the same path from hand-designed pipelines to end-to-end Transformers.

The TAL row is the detection lineage of Section~\ref{sec:action_detection}: sliding windows over hand-crafted features such as improved dense trajectories~\cite{wang2013idt}, then two-stage proposal-plus-classification systems (R-C3D~\cite{xu2017rcc3d}, SSN~\cite{zhao2017temporal}, TAL-Net~\cite{chao2018talnet}) with dedicated boundary-sensitive proposal generators (BSN~\cite{lin2018bsn}, BMN~\cite{lin2019bmn}, G-TAD~\cite{xu2020gtad}), then anchor-free dense detectors and query-based set prediction. Three developments go beyond that treatment. Dense single-stage detectors were scaled by ActionFormer~\cite{zhang2022actionformer}, a multi-scale Transformer encoder over the feature pyramid, and TriDet~\cite{shi2023tridet}, which replaces self-attention with a convolutional scalable-granularity-perception module and a trident boundary head; both still regress boundaries at every position and rely on non-maximum suppression. DETR-style set prediction removes anchors and suppression: RTD-Net~\cite{tan2021rtdnet} applies a relaxed Transformer decoder to proposal generation and TadTR~\cite{liu2022tadtr} decodes the full action set end to end~\cite{liu2022e2etad,kim2024tetad}. And masked video pretraining~\cite{tong2022videomae,wang2023videomaev2} now supplies the backbones that these heads are fine-tuned on.

TSG began with cross-modal matching. CTRL~\cite{gao2017tall} embeds the sentence and candidate clips in a shared space and ranks clips, with boundary regression on top; MCN~\cite{hendricks2017mcn} does the same for DiDeMo with global video context. Later models reason over the candidates jointly: MAN~\cite{zhang2019man} refines moments iteratively through a graph over proposals, 2D-TAN~\cite{zhang2020twodtan} scores every (start, end) pair on a two-dimensional temporal-adjacency map so that neighboring candidates inform each other, and VSLNet~\cite{zhang2020vslnet} drops proposals altogether and predicts the start and end positions as span extraction, in analogy to reading comprehension. Moment-DETR~\cite{lei2021qvhighlights} brought set prediction to grounding, decoding moments and saliency scores directly from an encoder--decoder, while TVR~\cite{lei2020tvr} extended grounding to corpus scale with a coarse-to-fine search over subtitles and frames.

\subsubsection{Datasets and Metrics}

Table~\ref{tab:temporal_localization_datasets} lists the benchmarks; the TAL rows are those of Section~\ref{sec:action_detection}, retained for comparison. The TSG benchmarks move from a closed domain (TACoS~\cite{tacos2013}, cooking) to open-domain sentences over indoor activity (Charades-STA~\cite{gao2017tall}), Flickr videos with multiple references per query (DiDeMo~\cite{hendricks2017mcn}), and dense paragraph-level captions (ActivityNet Captions~\cite{krishna2017densecap}); YouCook2~\cite{zhou2017youcook2} adds procedural structure, TVR~\cite{lei2020tvr} multimodal (subtitle) context and corpus-level retrieval, and QVHighlights~\cite{lei2021qvhighlights} saliency labels for joint highlight detection.

\begin{table}[t]
\centering
\caption{Representative datasets for temporal localization; ``Task'' distinguishes category-conditioned TAL from text-conditioned TSG. Counts are indicative of commonly used versions and splits.}
\label{tab:temporal_localization_datasets}
\small
\setlength{\tabcolsep}{6pt}
\begin{tabularx}{\linewidth}{@{}l c c c X@{}}
\toprule
\textbf{Task} & \textbf{Dataset} & \textbf{Year} & \textbf{Scale / Classes} & \textbf{Temporal annotations} \\
\midrule
TAL & THUMOS14~\cite{jiang2014thumos} & 2014 & 200 train + 213 test videos / 20 cls & Frame-accurate, untrimmed \\
TAL & ActivityNet v1.3~\cite{caba2015activitynet} & 2016 & $\sim$20k videos / 200 cls & Untrimmed, diverse \\
TAL & HACS Segments~\cite{zhao2019hacs} & 2019 & 50k videos / 200 cls; 139k segments & Clip+segment, dense \\
TAL & EPIC-KITCHENS-100~\cite{damen2020epic100} & 2020 & 700 egocentric videos / verb-noun labels & Egocentric, narration-aligned \\
TAL & MultiTHUMOS~\cite{yeung2018every} & 2015 & $\sim$400 videos / 65 cls; $\sim$38.7k ann. & Dense, multi-label \\
TAL & FineAction~\cite{liu2022fineaction} & 2022 & 17k videos / 106 cls; 103k instances & Fine-grained, precise \\
\midrule
TSG & TACoS~\cite{tacos2013} & 2013 & 127 cooking videos; $\sim$18.8k query-segment pairs & Cooking, aligned \\
TSG & Charades-STA~\cite{gao2017tall} & 2017 & $\sim$6.6k videos; 16,128 pairs & Charades-based pairs \\
TSG & DiDeMo~\cite{hendricks2017mcn} & 2017 & $\sim$10k Flickr videos; $\sim$40k queries & Multi-reference, ranking \\
TSG & ActivityNet Captions~\cite{krishna2017densecap} & 2017 & 20k videos; $\sim$100k sentences & Dense captions \\
TSG & YouCook2~\cite{zhou2017youcook2} & 2018 & 2k untrimmed cooking videos / 89 recipes & Cooking procedures \\
TSG & TVR~\cite{lei2020tvr} & 2020 & 21.8k clips (60--90 s) from 6 TV shows; 109k queries & Subtitles, multimodal \\
TSG & QVHighlights~\cite{lei2021qvhighlights} & 2021 & $\sim$10k YouTube videos; query+saliency & Query-saliency, highlights \\
\bottomrule
\end{tabularx}
\end{table}

Temporal IoU between a predicted and a ground-truth segment is the primitive for both tasks, and TAL is scored by mAP at fixed tIoU thresholds (per-threshold on THUMOS14, averaged over $[0.5{:}0.95]$ on ActivityNet) and by AR@AN for proposal quality~\cite{anet2019proposals}, as defined in Section~\ref{sec:action_detection}. TSG uses retrieval metrics instead. For queries $\mathcal{Q}$, ranked predictions $\hat{m}_j$, and reference segments $G(q)$,
\begin{equation}
\mathrm{R}@n,\mathrm{IoU}{=}\theta \,=\, \frac{1}{|\mathcal{Q}|}\sum_{q\in\mathcal{Q}} \mathbf{1}\!\left(\max_{1\le j\le n}\;\max_{g\in G(q)} \mathrm{tIoU}(\hat{m}_j,g)\ge \theta\right)
\end{equation}
is the fraction of queries for which some top-$n$ prediction overlaps a reference by at least $\theta$ (commonly $n\in\{1,5\}$, $\theta\in\{0.3,0.5,0.7\}$), and $\mathrm{mIoU} = \frac{1}{|\mathcal{Q}|}\sum_{q} \max_{g\in G(q)} \mathrm{tIoU}(\hat{m}_1,g)$ scores the top prediction alone; QVHighlights adds mAP over tIoU thresholds~\cite{lei2021qvhighlights}. Recall at a threshold binarizes overlap and ignores the quality of the ranking, which motivated ranking-sensitive alternatives such as AxIoU~\cite{togashi2022axiou}, and annotator disagreement about where a described moment begins and ends puts a ceiling on how precisely any of these can be measured.

Table~\ref{tab:temporal_localization_performance} gives R@1 results for the TSG methods above; TAL numbers are in the table of Section~\ref{sec:action_detection}. Relational reasoning over candidates (2D-TAN) and span prediction (VSLNet) improved the listed Charades-STA R@1 from CTRL's 23.6/8.9 to 39.7/23.3 and 54.2/35.2, respectively, and Moment-DETR adds a further step on Charades-STA without proposals or suppression; its QVHighlights numbers are the reference point for joint grounding and highlight detection.

\begin{table}[!t]
    \centering
    \fontsize{9}{11}\selectfont
    \setlength{\tabcolsep}{1.2mm}
    \caption{Temporal sentence grounding, R@1 at tIoU 0.5 / 0.7 (\%). Charades-STA and ActivityNet Captions values as reported in the 2D-TAN~\cite{zhang2020twodtan}, VSLNet~\cite{zhang2020vslnet}, and Moment-DETR~\cite{lei2021qvhighlights} papers (visual features differ across methods as in the originals; MAN uses I3D features); QVHighlights test values from Moment-DETR without pretraining, with mAP averaged over tIoU thresholds.}
    \label{tab:temporal_localization_performance}
    \begin{tabular}{L{0.26\textwidth} L{0.26\textwidth} C{0.13\textwidth} C{0.15\textwidth} C{0.11\textwidth}}
    \toprule
    \textbf{Method} & \textbf{Design} & \textbf{Charades-STA} & \textbf{ActivityNet Captions} & \textbf{QVHighlights} \\
     & & R@1 0.5 / 0.7 & R@1 0.5 / 0.7 & R@1 0.5 / 0.7 (mAP) \\
    \midrule
    CTRL~\cite{gao2017tall}             & Clip--sentence matching      & 23.6 / 8.9  & 29.0 / 10.3 & -- \\
    MAN~\cite{zhang2019man}             & Iterative graph adjustment   & 46.5 / 22.7 & --          & -- \\
    2D-TAN~\cite{zhang2020twodtan}      & 2D temporal-adjacency map    & 39.7 / 23.3 & 44.5 / 26.5 & -- \\
    VSLNet~\cite{zhang2020vslnet}       & Span prediction              & 54.2 / 35.2 & 43.2 / 26.2 & -- \\
    Moment-DETR~\cite{lei2021qvhighlights} & Set prediction              & 53.6 / 31.4 & --          & 52.9 / 33.0 (30.7) \\
    \bottomrule
    \end{tabular}
\end{table}

\subsubsection{Challenges}

Boundaries are soft. Actions span sub-second gestures to multi-minute activities, which no single temporal resolution serves, and annotators disagree about where an instance starts and ends; a sentence adds its own ambiguity about which part of an activity it describes. The tIoU-thresholded metrics treat these fuzzy boundaries as exact, so part of the measured error is irreducible~\cite{togashi2022axiou}, and probabilistic boundary modeling with calibrated confidences is the natural response.

The two tasks are converging. A closed label set limits TAL, while TSG is bounded by the vocabulary of its query sets; text prompts in place of class labels~\cite{liberatori2024t3al} and set prediction for moments~\cite{lei2021qvhighlights} point to a single architecture that grounds both categories and sentences. Doing this over minutes of video, across camera viewpoints (egocentric footage in particular), with audio, narration, or noisy subtitles as further evidence~\cite{lei2020tvr,lei2021qvhighlights}, and causally for streaming use~\cite{shou2018odas} is where the current systems fall short: memory-efficient long-context encoders and corpus-scale coarse-to-fine retrieval are the two ingredients most actively developed.








% \newpage

\section{Multimodal Spatial Perception}
\label{sec:p3-sec6-5-multimodal-spatial-perception}
% \wsq{将视觉与语言、音频、触觉等模态结合，增强空间感知范围与粒度。}

\FramedInitial{T}{he} preceding sections perceive space largely through vision. Yet vision alone is often incomplete: it is blind past occluders, silent about what lies outside the frame, and mute on which entity a person actually means. Multimodal spatial perception closes these gaps by grounding spatial understanding in signals beyond the image. The section takes two complementary couplings: audio-visual spatial localization, which uses sound to place events that vision may miss, and vision-language grounding, which ties linguistic reference to the specific entity or region it denotes. Neither adds a new spatial primitive; each supplies a channel of evidence that makes perception more complete where a single modality falls short.


\subsection{Audio-Visual Spatial Localization}
\label{sec:av_localization}
% \wsq{声源定位与空间目标匹配（如声音中找人）}

Humans localize a sound by combining auditory cues (interaural level and time differences, spectral shaping) with visual context (lip motion, object dynamics)~\cite{blauert1997spatial}. Audio-visual spatial localization (AV-SL) asks a machine to do the same: find the pixel region, object, or 3D point that a sound comes from, under reverberation, noise, occlusion, and several simultaneous sources~\cite{arandjelovic2017look,owens2018audio}. The two modalities are complementary. Binaural or array audio supplies directional cues, while distance requires additional geometric, near-field, or learned evidence; vision says which entities are plausible sounders and separates acoustically confusable ones (two similar instruments, only one playing). Spatial audio for AR/VR~\cite{morgado2018self}, teleconferencing~\cite{tao2021talknet}, and audio-guided navigation~\cite{chen2020soundspaces} are the main consumers.

\subsubsection{Task Definition}
The input is an audio stream $A$ (mono, binaural, or multichannel, optionally ambisonic~\cite{zotter2019ambisonics}) time-aligned with a video $V$ (RGB, depth, or multi-view) of the same interval; the audio may contain zero, one, or several sources, and the video may show silent objects that look like sounders. The output localizes the sounding entity in a form that depends on the variant: a 2D heatmap or mask, a face or body box with a speaking state, an event class with its temporal extent and direction of arrival (DOA, azimuth and elevation) and optionally distance~\cite{politis2023dataset}, a 2D location with depth, or a 3D point or trajectory in camera or world coordinates. For visible-region localization, off-screen sources or audio with no visible source should yield a null prediction; 3D acoustic localization may still estimate an off-screen source.

Five variants are in common use. Visual sound source localization (VSSL) highlights the sounding region of a frame or clip as a class-agnostic heatmap, in single- and multi-source settings and increasingly with negative or off-screen audio as a robustness test. Active speaker detection (ASD) decides who is speaking and localizes the talker as a face track, often under streaming constraints~\cite{tao2021talknet}. Audio-visual sound event localization and detection (AV-SELD) couples what and when with where: event class, temporal extent, DOA, and in recent benchmarks distance~\cite{politis2023dataset}. The 2.5D variant adds depth or range to a 2D localization from monocular depth~\cite{DepthAnything}, stereo, or echoes~\cite{gao2020visualechoes}; and 3D DOA estimation returns a spatial bearing, whereas source-position or trajectory estimation additionally uses range or triangulation constraints from calibrated cameras, arrays, or egocentric rigs~\cite{grauman2022ego4d}.

\subsubsection{Methodology}
Figure~\ref{fig:p3_s6_5_1_method_overview} maps the five variants onto four method families.

\begin{figure}[!t]
    \centering
    \includegraphics[width=\linewidth]{figures/part3/section6/audio_video_spatial_localization.pdf}
    \caption{Audio-visual spatial localization. \textit{Top}: task outputs for (i) VSSL, (ii) ASD, (iii) AV-SELD, (iv) 2.5D localization, and (v) 3D localization and tracking. \textit{Bottom}: method families (A) cross-modal correspondence, (B) audio--lip synchrony, (C) joint spatial event prediction, and (D) geometric lifting or fusion; neutral connectors mark principal task--family associations, not computational flow. Directional arrows inside panels indicate prediction or information flow. The 2.5D panel links a selected image location to depth; 3D position additionally requires range or triangulation constraints. Scene imagery and overlays are illustrative.}
    \label{fig:p3_s6_5_1_method_overview}
\end{figure}

\paragraph{Cross-Modal Correspondence (VSSL).}
VSSL is learned almost entirely without location labels, from the co-occurrence of sound and image. An audio encoder maps the clip to an embedding $\mathbf{a}$, a visual encoder maps the frame to a feature map $\mathbf{v}(\mathbf{x})$, and the localization map is the per-position similarity
\begin{equation}
    S(\mathbf{x}) = \frac{\mathbf{a}^{\top}\mathbf{v}(\mathbf{x})}{\|\mathbf{a}\|\,\|\mathbf{v}(\mathbf{x})\|},
\end{equation}
trained so that matching audio-image pairs score higher than mismatched ones~\cite{senocak2018learning,owens2018audio}; Sound of Pixels~\cite{zhao2018sound} instead mixes two clips and asks the network to separate them conditioned on pixel features, which forces the same alignment. Later work targets the failure modes of co-occurrence. LVS~\cite{chen2021localizing} mines hard negatives inside the same image (regions that look like sounders but do not match the audio) and introduces the VGG-SS localization benchmark; deep multimodal clustering~\cite{hu2019deep} and coarse-to-fine mixture objectives~\cite{qian2020multiple} handle several sources at once; and SLAVC~\cite{mo2022closer} shows that earlier methods depend on early stopping and adds momentum encoders and extreme visual dropout so that the alignment does not collapse onto background.

\paragraph{Audio--Lip Alignment (ASD).}
ASD runs a face detector and tracker and then decides, per track and per frame, whether the visible mouth produced the audio. The decision rests on audio-visual synchrony: a visual encoder over the lip region and an audio encoder over the spectrogram are compared across time~\cite{afouras2020self}, and TalkNet~\cite{tao2021talknet} adds long-term temporal modeling of both streams with cross-attention, which is what separates a brief pause from a non-speaker. Active Speakers in Context~\cite{alcazar2020active} scores each candidate jointly with the other faces in the frame, and ASDNet~\cite{kopuklu2021how} makes the three stages explicit (per-frame audio-visual encoding, inter-speaker relation modeling, temporal modeling), with each stage trained for its role.

\paragraph{Joint Event and DOA Prediction (AV-SELD).}
SELDnet~\cite{adavanne2018sound} set the template: a convolutional recurrent network over multichannel spectrogram features with two heads, one classifying active events per frame and one regressing their DOA. ACCDOA~\cite{shimada2021accdoa} folds the two into one output, a vector per class whose direction is the DOA and whose length is the activity, which removes the need to balance two losses; the DCASE challenge tracks~\cite{politis2020overview} standardized the protocol, and STARSS23~\cite{politis2023dataset} added video and source distance.

\paragraph{Geometric Lifting (2.5D and 3D).}
Lifting a 2D localization to metric space uses geometry rather than learned correspondence alone. VisualEchoes~\cite{gao2020visualechoes} shows that echoes of an emitted chirp carry depth: the time and spectrum of the return, fused with the image, improve monocular depth estimation and give the 2.5D variant its range. Audio-visual spatial alignment~\cite{morgado2020learning} learns representations by predicting whether 360$^\circ$ video and ambisonic audio are spatially consistent, a pretext that encodes direction, and self-supervised spatial audio generation~\cite{morgado2018self} learns the inverse mapping from 360$^\circ$ video to ambisonics. Full 3D systems combine classical array processing (time differences of arrival from GCC-PHAT~\cite{knapp1976generalized}, steered-response-power beamforming~\cite{dibiase2001robust}) with calibrated cameras and triangulation, and in simulation (SoundSpaces~\cite{chen2020soundspaces,chen2022soundspaces}) with the agent's own motion.

\subsubsection{Datasets and Metrics}

Table~\ref{tab:p3_s6_5_1_avsl_tasks} pairs each variant with its benchmarks and metrics; there is no shared benchmark across variants because their outputs differ.

\begin{table}[!ht]
\centering
\caption{AV-SL task variants with representative datasets, outputs, and metrics. Newer VSSL test suites add negative or off-screen audio to measure robustness.}
\label{tab:p3_s6_5_1_avsl_tasks}
\small
\begin{tabular}{@{}L{0.12\textwidth} L{0.34\textwidth} L{0.20\textwidth} L{0.24\textwidth}@{}}
\toprule
\textbf{Task} & \textbf{Representative Datasets} & \textbf{Primary Outputs} & \textbf{Typical Metrics} \\
\midrule
VSSL (2D) & Flickr-SoundNet~\cite{senocak2018learning}, MUSIC~\cite{zhao2018sound}, VGG-SS~\cite{chen2021localizing} (from VGG-Sound~\cite{chen2020vggsound}), AVSBench~\cite{zhou2022avs} & heatmap / mask & cIoU, AUC; AP with negatives \\
ASD & AVA-ActiveSpeaker~\cite{roth2020ava} & talker boxes + speaking state & mAP \\
AV-SELD & STARSS23~\cite{politis2023dataset}, DCASE SELD tracks~\cite{politis2020overview} & event + time + DOA (+ distance) & F-score, error rate, DOA error, distance error \\
2.5D & Replica / Matterport3D echo splits~\cite{gao2020visualechoes} & 2D + depth & AbsRel, RMSE \\
3D DOA / trajectory & STARSS23 positions~\cite{politis2023dataset}, SoundSpaces~\cite{chen2020soundspaces} & DOA / 3D point / trajectory & angular error, positional error \\
\bottomrule
\end{tabular}
\end{table}

VSSL is scored by thresholding the localization map and computing consensus intersection over union (cIoU) against the annotations; reported metrics include the fraction of test images with cIoU above 0.5 (localization accuracy) and the area under the IoU-threshold curve (AUC)~\cite{senocak2018learning}; the extended VGG-SS and Flickr protocols add clips whose audio has no visible source and report detection-style average precision requiring both visible-source presence and correct localization~\cite{chen2021localizing,mo2022closer}. ASD uses detection-style mAP over face tracks with speaking labels~\cite{roth2020ava}. SELD combines detection (F-score, error rate) with localization (DOA error, localization recall)~\cite{politis2020overview}, while STARSS23 additionally provides source-distance annotations~\cite{politis2023dataset}; 2.5D inherits the depth metrics of Section~\ref{sec:depth_estimation}, and 3D settings report angular and positional error.

Table~\ref{tab:p3_s6_5_1_avsl_performance} collects the two variants with a stable leaderboard. On VSSL the negative-audio protocol changes the ranking: LVS leads the earlier methods on the standard Flickr test but drops to 9.8 AP when silent clips are mixed in, whereas SLAVC holds both. ASD on AVA-ActiveSpeaker has moved from 79.2 mAP for the dataset baseline to above 93 with explicit temporal and inter-speaker modeling, and the remaining differences among top systems are about latency and streaming rather than accuracy. AV-SELD has no such table: each DCASE edition changes the data and the metric, and distance-aware evaluation is not directly comparable to the DOA-only tracks that SELDnet and ACCDOA were evaluated on.

\begin{table}[!t]
    \centering
    \fontsize{9}{11}\selectfont
    \setlength{\tabcolsep}{1.2mm}
    \caption{Visual sound source localization and active speaker detection. VSSL: cIoU (\%) on the Flickr-SoundNet and VGG-SS test sets, and AP (\%) on their extended versions with negative audio, all models trained on VGG-Sound 144k, as reported in the SLAVC comparison~\cite{mo2022closer}. ASD: mAP (\%) on AVA-ActiveSpeaker val, from the TalkNet~\cite{tao2021talknet} and ASDNet~\cite{kopuklu2021how} papers.}
    \label{tab:p3_s6_5_1_avsl_performance}
    \begin{tabular}{L{0.26\textwidth} L{0.26\textwidth} C{0.09\textwidth} C{0.09\textwidth} C{0.09\textwidth} C{0.09\textwidth}}
    \toprule
    \rowcolor{lightgrey} \multicolumn{6}{l}{$\bullet$ \textit{Visual sound source localization}} \\
    \textbf{Method} & \textbf{Design} & \textbf{Flickr cIoU} & \textbf{VGG-SS cIoU} & \textbf{Ext.\ Flickr AP} & \textbf{Ext.\ VGG-SS AP} \\
    \midrule
    Attention10k~\cite{senocak2018learning} & Audio-guided attention   & 42.3 & 18.5 & --   & -- \\
    DMC~\cite{hu2019deep}                   & Multimodal clustering    & 55.6 & 23.9 & --   & -- \\
    LVS~\cite{chen2021localizing}           & Hard negatives           & 71.9 & 34.4 & 9.8  & 5.2 \\
    SLAVC~\cite{mo2022closer}               & Momentum encoder, dropout & 83.2 & 37.2 & 51.6 & 33.0 \\
    \midrule
    \rowcolor{lightgrey} \multicolumn{6}{l}{$\bullet$ \textit{Active speaker detection (AVA-ActiveSpeaker val, mAP)}}\\
    \hiderowcolors \textbf{Method} & \textbf{Design} & \textbf{mAP} & & & \\
    \midrule
    AVA baseline~\cite{roth2020ava}         & Audio-visual fusion      & 79.2 & & & \\
    ASC~\cite{alcazar2020active}            & Inter-speaker context    & 87.1 & & & \\
    TalkNet~\cite{tao2021talknet}           & Long-term temporal, cross-attention & 92.3 & & & \\
    ASDNet~\cite{kopuklu2021how}            & Three-stage architecture & 93.5 & & & \\
    \bottomrule
    \end{tabular}
\end{table}

\subsubsection{Challenges}
Co-occurrence is not causation. A model that has learned which objects usually make which sounds will localize a guitar whether or not it is playing, and the negative-audio protocols~\cite{chen2021localizing,mo2022closer} exist because the standard tests rewarded exactly this; the open problem is grounding on the acoustic evidence itself, which matters most for off-screen and silent cases where the right answer is no localization.

Going from 2D saliency to metric 3D requires semantics and geometry together. Depth foundation models~\cite{DepthAnything} (Section~\ref{sec:depth_estimation}) and self-supervised visual features~\cite{dinov2} now supply priors that AV-SL pipelines previously learned from scratch, and STARSS23~\cite{politis2023dataset} provides joint event, time, DOA, and distance annotations, supporting distance-aware extensions; multi-object objectives for overlapping events are the current frontier there.

Multiple simultaneous sources, identity through occlusion, and on-screen/off-screen transitions over long videos~\cite{qian2020multiple,tian2018audio} are the AV-SL form of the identity problem of Section~\ref{sec:object_tracking} and the pose-continuity problem of Section~\ref{sec:object_pose}. Streaming use adds causal encoders, low-latency fusion, calibration drift, and graceful fallback when one modality is missing~\cite{tao2021talknet,kopuklu2021how}.
% TODO: anchor §8.x once §8 (Spatial Reconstruction and Generation) subsections stabilize --- AV-SL is a natural consumer of 4D scene reconstruction (spatial audio rendering, dynamic-scene physical reasoning) per the §6.4.1 (Scene Flow) §Challenges ¶3 convergence narrative.



\subsection{Vision-Language Grounding}
\label{sec:vl_grounding}
% \wsq{``把手放在右边的杯子上'' -> 找到右边的杯子; 与 RefCOCO / Refer-Youtube-VOS 等任务相关}

Vision-language grounding ties a linguistic reference to the element of a scene it denotes. Given \textit{``Put your hand on the small mug on the right side of the table behind the blue book''}, a grounded system must resolve the noun phrase, its attributes, and its spatial relations to one particular mug in the image; this referential link is what instruction following, multimodal dialogue, and human-robot interaction build on.

\subsubsection{Task Definition}
The input is a scene $S$ (an image, a video, or a 3D environment) and a query $Q$, from a bare noun phrase (``the red cup'') to a sentence with attributes, relations, actions, or an implicit goal (``the person who is waving'', ``something to open the box with''). The output $\mathcal{R}(Q,S)$ is the referred region in the form the variant requires: a box, a mask, a set of regions (possibly empty), a per-frame mask track, or a 3D box. Two formulations recur through the methods below: ranking a set of candidate regions $r$ by a cross-modal score, or regressing the region directly,
\begin{equation}
    \hat{r} = \arg\max_{r \in \mathcal{C}(S)} s\big(\phi(Q), \psi(S, r)\big)
    \qquad \text{or} \qquad
    \hat{r} = f_\theta(S, Q),
\label{eq:p3_s6_5_2_grounding}
\end{equation}
with $\phi$ and $\psi$ the text and region encoders and $\mathcal{C}(S)$ a proposal set.

\begin{figure}[!ht]
    \centering
    \includegraphics[width=0.99\linewidth]{figures/part3/section6/visual_language_graounding_task.pdf}
    \caption{Vision-language grounding task variants: (i) REC / RES~\cite{yu2016modeling,mao2016generation}; (ii) generalized REC / RES with zero, one, or multiple targets~\cite{liu2023gres}; (iii) reasoning grounding~\cite{lai2024lisa}; (iv) video referring and reasoning grounding~\cite{luo2026drV}; (v) 3D and embodied grounding~\cite{chen20263d}.}
    \label{fig:p3_s6_5_2_method_overview_task}
\end{figure}

Figure~\ref{fig:p3_s6_5_2_method_overview_task} lays out the five variants, which differ in scene modality, output granularity, and query complexity.

\textbf{Referring expression comprehension and segmentation (REC / RES).} One image, one expression, one target, returned as a box (REC) or a mask (RES). RefCOCO and RefCOCO+~\cite{yu2016modeling} (short phrases, the latter without location words) and RefCOCOg~\cite{mao2016generation} (longer descriptions), all using MS~COCO~\cite{lin2014microsoft} images, with expressions collected through ReferItGame~\cite{kazemzadeh2014referitgame} or Mechanical Turk, have been the standard benchmarks for a decade; Flickr30k Entities~\cite{plummer2015flickr30k,young2014image} grounds every entity mention of a caption at once.

\textbf{Generalized REC / RES (GREC / GRES).} An expression may match zero, one, or several objects (``the red carrots'' versus ``the unicorn''), which single-target methods cannot represent. gRefCOCO~\cite{liu2023gres}, Ref-ZOM~\cite{hu2023beyond}, PhraseCut~\cite{wu2020phrasecut}, and OmniLabel~\cite{schulter2023omnilabel} evaluate this setting, and MaskGroups-HQ~\cite{cao2025refer} generalizes the prompt itself to arbitrary groups of masks specified by text and reference masks. The output is a set, so the metric must be as well (see Datasets and Metrics below).

\textbf{Reasoning grounding.} Some queries require inference beyond noun matching: ``the loser'' needs a judgment about the people shown, ``turn on the TV'' needs the remote. CLEVR-Ref+~\cite{liu2019clevr} diagnoses this on rendered scenes with ground-truth reasoning programs; ReasonSeg~\cite{lai2024lisa} poses implicit, multi-step queries on real images with mask output, extended by MMR~\cite{jang2025mmr} (194K implicit queries, multi-target and part-level), ReasonSeg-Inst~\cite{yang2023lisa++} (instance level), and MC-Bench~\cite{xu2024mcbench} (instances across image pairs). Ref-Adv~\cite{dong2026ref} keeps the classic REC format but makes expressions minimally sufficient, negation-heavy, and paired with same-category distractors.

\textbf{Video referring and reasoning grounding.} The referent must be found in every frame through motion, appearance change, and occlusion, and the query may name it by what it does (``the man waving his hand''). Refer-Youtube-VOS~\cite{seo2020urvos}, Refer-DAVIS17~\cite{khoreva2018video}, A2D- and JHMDB-Sentences~\cite{gavrilyuk2018actor}, MeViS~\cite{ding2023mevis}, ReVOS~\cite{yan2024visa}, and ReasonVOS~\cite{bai2024one} add language to existing segmentation and actor-action datasets; Ref-VPS~\cite{bagchi2025refereverything} extends the referent from objects to processes such as smoke spreading or waves breaking, which have no object boundary at all. Tracking machinery is usually folded in to carry identity across frames~\cite{fang2025decoupled}.

\textbf{3D and embodied grounding.} The scene is a point cloud or a set of views, and the query refers to an object in 3D (``the chair next to the window''). ScanRefer~\cite{chen2020scanrefer} requires detecting the referent as a 3D box and Multi3DRefer~\cite{zhang2023multi3drefer} allows zero to many targets; ReferIt3D~\cite{achlioptas2020referit3d} provides the object instances and scores selection among them; Instruct3D~\cite{he2024segpoint} and Multi3DRes~\cite{wu20243d} ask for point-level masks from instruction-style queries. Embodied grounding makes the output an action: an agent told ``go pick up the red cup on the right'' must identify the cup and reach it~\cite{zhen20243d}.

\subsubsection{Methodology}

\begin{figure}[!t]
    \centering
    \includegraphics[width=0.99\linewidth]{figures/part3/section6/visual_language_graounding_method.pdf}
    \caption{Vision-language grounding method families: (i) proposal ranking, (ii) compositional reasoning, (iii) end-to-end grounding, (iv) grounded pre-training, (v) foundation models, and (vi) tool-augmented agents. Panels (i)-(iv) correspond to families (A)-(D) in the text, (v) to (E)-(F), and (vi) to (G); family (H), generation for grounding, is not depicted.}
    \label{fig:p3_s6_5_2_method_overview_method}
\end{figure}

The methods below fall into eight families, (A) to (H); Figure~\ref{fig:p3_s6_5_2_method_overview_method} illustrates (A) through (G). They are trajectories rather than eras: the oldest paradigm has resurfaced at foundation-model scale.

\paragraph{(A) Proposal Ranking.}
The first grounding models instantiate the ranking form of Eq.~\eqref{eq:p3_s6_5_2_grounding}: detect region proposals, embed them and the query in a joint space, and pick the best-scoring region. Deep Fragment Embeddings~\cite{karpathy2014deep} aligned sentence fragments with detected regions, and scene-graph CRFs~\cite{johnson2015image} grounded structured queries to regions for retrieval. Proposal quality bounded these methods and context-dependent expressions defeated them. The same two-stage shape has returned with far stronger proposals: RAS~\cite{cao2025refer} has a segmentation foundation model propose candidate masks and a mask-centric multimodal LLM \emph{select} the group that satisfies the prompt, and this non-autoregressive selection beats mask generation in both accuracy and latency.

\paragraph{(B) Compositional Reasoning.}
Modular networks parse the expression and route its parts to dedicated modules. CMN~\cite{hu2017modeling} decomposes an expression into a subject--relationship--object triplet handled by localization and relationship modules; MAttNet~\cite{yu2018mattnet} soft-parses phrases into appearance, location, and relationship attention modules; dynamic graph attention~\cite{yang2019dynamic} reasons over the linguistic structure on a scene graph. These models disambiguate compositional queries (``the man behind the woman in red'') but inherit every parsing error.

\paragraph{(C) End-to-End Grounding.}
The regression form of Eq.~\eqref{eq:p3_s6_5_2_grounding} removes proposals: a single network fuses the sentence embedding with the image feature map and predicts the box. One-Stage Grounding~\cite{yang2019fast} did this with a detector head, TransVG~\cite{deng2021transvg} with a pure Transformer in which a learnable token attends over visual and textual tokens and regresses coordinates. On the segmentation side, CRIS~\cite{wang2022cris} transfers CLIP's image-text alignment to the pixel level with text-to-pixel contrastive learning and made end-to-end RES the standard formulation.

\paragraph{(D) Grounded Pre-Training.}
Phrase-region alignment learned at scale turned grounding into an open-vocabulary capability. MDETR~\cite{kamath2021mdetr} trains a DETR-style set predictor on text-aligned detection corpora so that every box is tied to a span of the caption; GLIP~\cite{li2022grounded} reformulates detection as phrase grounding and scales to 27M region-caption pairs, giving zero-shot transfer; RegionCLIP~\cite{zhong2022regionclip} adds region-level contrastive alignment, and Grounding DINO~\cite{liu2024grounding} deep language fusion at every stage of a DINO detector, which made it the default open-vocabulary detector of later tool chains (family (G)). The residual weakness is the \emph{entity shortcut}: models latch onto object names and ignore the description. DesCo~\cite{li2023desco} counters it in the data, pre-training on LLM-generated rich descriptions with context-sensitive negative queries, which raises description-conditioned detection on LVIS and OmniLabel.

\paragraph{(E) Grounding MLLMs.}
The first multimodal LLMs built for grounding gave the language backbone a region interface. Kosmos-2~\cite{peng2023kosmos} emits location tokens that hyperlink phrases to boxes inside the generated text; Ferret~\cite{you2023ferret} and Ferret-v2~\cite{zhang2024ferret} encode boxes, points, and free-form regions as inputs at any resolution. Mask output comes from attaching a decoder: LISA~\cite{lai2024lisa} decodes a dedicated segmentation token's embedding into a mask, which is what makes reasoning segmentation possible, and LISA++~\cite{yang2023lisa++} extends it to instances within a dialogue; GLaMM~\cite{rasheed2024glamm} unifies phrase grounding, referring segmentation, and free-form conversation under a grounded-conversation objective; Sa2VA~\cite{yuan2025sa2va} couples a LLaVA-style backbone to SAM~2~\cite{ravi2024sam2} and carries dense grounding from images to video referring segmentation. Moondream~3's segmentation extension~\cite{reid2026moondream} dispenses with the decoder and writes the mask autoregressively as a vector path, refined by a lightweight head and trained with reinforcement learning on mask overlap.

\paragraph{(F) Generalist MLLMs.}
A parallel line keeps no grounding component at all: the model writes coordinates as plain text, and grounding is one capability among many. Florence-2~\cite{xiao2024florence} is the compact case, a sequence-to-sequence model trained on the 5.4B-annotation FLD-5B corpus that casts captioning, detection, grounding, and polygon-vertex segmentation into one text interface, with mask accuracy below the decoder-based models of family (E). The frontier generalists, Gemini~\cite{reid2024gemini}, Qwen2.5-VL and Qwen3-VL~\cite{bai2025qwen2.5vl}, InternVL~2.5 and InternVL~3~\cite{chen2024expanding,zhu2025internvl3}, and Seed1.5-VL~\cite{guo2025seed1.5vl}, now match or exceed the specialists on the RefCOCO family from prompting alone (Table~\ref{tab:p3_s6_5_2_grounding_performance}); Llama~3~\cite{grattafiori2024llama} folds phrase-to-box and phrase-to-mask annotations into its multimodal pre-training but reports no REC numbers.

\paragraph{(G) Tool-Augmented Agents.}
A third line delegates localization to specialist models under an outer control loop, which has moved from fixed pipelines to generated programs to learned tool use. Grounded SAM~\cite{ren2024grounded} chains Grounding DINO boxes into SAM~\cite{kirillov2023sam} masks without training; ViperGPT~\cite{suris2023vipergpt} has a code LLM write a Python program that calls perception APIs, so that logic and arithmetic are symbolic and perception is delegated, with interpretable intermediate steps. SegAgent~\cite{zhu2025segagent} treats mask annotation as a sequential decision process, imitating the click sequences of human annotators on an interactive segmenter; DeepEyes~\cite{zheng2025deepeyes} internalizes the tool, using reinforcement learning to make the model emit coordinates that trigger zoom-in crops inside its own chain of thought. SAM~3 Agent~\cite{carion2025sam} wraps SAM~3, whose concept prompts return every instance matching a noun phrase, with an MLLM that proposes phrases, inspects the returned masks, and selects those satisfying the full expression, which gives zero-shot referring and reasoning segmentation without grounding-specific training.

\paragraph{(H) Generation for Grounding.}
The newest family treats a large text-to-image or text-to-video generator as a grounding model, on the premise that generative pre-training already induces a fine-grained map from language to spatial and temporal structure. VD-IT~\cite{zhu2024exploring} uses a frozen text-to-video diffusion model as the feature extractor for referring video segmentation and finds its features more temporally consistent and better aligned with language than discriminative video backbones; ReferEverything~\cite{bagchi2025refereverything} keeps the generator intact and fine-tunes it to predict mask latents decoded by the frozen VAE, which transfers zero-shot to rare objects, ``stuff'', and the amorphous processes of Ref-VPS. Vision Banana~\cite{gabeur2026image} instruction-tunes a state-of-the-art image generator to render the task output as an image, segmentation as painting the referred region a prescribed color, and matches specialist systems zero-shot while keeping its generative ability. Grounding quality in this family rises with the underlying generator.

\subsubsection{Datasets and Metrics}

Table~\ref{tab:p3_s6_5_2_grounding_tasks} lists the benchmarks by variant. Annotation cost rises steeply from a 2D box to a per-frame mask to a 3D box, so image-level REC/RES corpora dominate in size and maturity, video referring sets are an order of magnitude smaller, and 3D corpora smaller still. Within the image benchmarks a generational split has opened. RefCOCO/+/g now serve mainly as training corpora: specialists and generalist MLLMs report Acc@0.5 above 90\% on them, an audit found label errors in 14\%, 24\%, and 5\% of RefCOCO, RefCOCO+, and RefCOCOg~\cite{chen2025revisiting}, and much of the score survives shuffling the word order of the expression~\cite{dong2026ref}. The evaluation-oriented replacements are built to expose what saturation hides: Ref-L4~\cite{chen2025revisiting} (cleaned labels, long expressions, a 365-category vocabulary), Ref-Adv~\cite{dong2026ref} (shortcut-suppressed, where the same models fall to roughly 60\%), HC-RefLoCo~\cite{wei2024hcrefloco}, GigaGrounding~\cite{ma2024gigaground}, and GroundingME~\cite{li2026groundingme}, which scores discrimination, spatial reasoning, limited visibility, and the rejection of ungroundable queries, and on which the best model resolves 45\% of samples while most score zero on rejection.

\begin{table}[!ht]
\centering
\caption{Common vision-language grounding tasks with representative datasets and typical evaluation metrics; several tasks share datasets (\eg, RefCOCO underlies both REC and RES).}
\label{tab:p3_s6_5_2_grounding_tasks}
\fontsize{9}{11}\selectfont
\setlength{\tabcolsep}{4pt}
\begin{tabular}{L{0.18\textwidth} L{0.34\textwidth} L{0.18\textwidth} L{0.22\textwidth}}
\toprule
\textbf{Task} & \textbf{Representative Datasets} & \textbf{Primary Outputs} & \textbf{Typical Metrics} \\
\midrule
REC / RES & RefCOCO/+/g, Flickr30k Entities, Ref-L4, Ref-Adv, GroundingME & box / mask & Acc@0.5; gIoU / cIoU (RES) \\
GREC / GRES & gRefCOCO, OmniLabel, PhraseCut, Ref-ZOM, MaskGroups-HQ & multi-box / multi-mask & F1, N-acc, gIoU / cIoU, AP$_{\text{descr}}$ \\
Reasoning grounding & CLEVR-Ref+, ReasonSeg, ReasonSeg-Inst, MMR, MC-Bench & mask (or box) & gIoU / cIoU, Acc@0.5 \\
Video referring & MeViS, Refer-YT-VOS, Refer-DAVIS17, ReVOS, A2D-Sent., Ref-VPS & per-frame mask & $\mathcal{J}\&\mathcal{F}$; mAP (A2D-Sent.) \\
3D / Embodied & ScanRefer, ReferIt3D, Multi3DRefer, Instruct3D & 3D box / 3D mask & Acc@0.25/0.5, F1@0.5, mIoU \\
\bottomrule
\end{tabular}
\end{table}

Every metric here is built on the IoU of predicted and reference regions (Section~\ref{sec:object_tracking} gives the mask form). REC reports Accuracy@0.5, the fraction of expressions whose box reaches IoU $\ge 0.5$, with stricter thresholds and the threshold-averaged mAcc~\cite{chen2025revisiting} increasingly added; phrase grounding, where several proposals may be returned, reports Recall@$k$. For set-valued outputs (GREC / GRES) the prediction is compared with the ground-truth set through per-expression F1 over boxes or masks~\cite{hu2023beyond} and a no-target accuracy for expressions with no referent~\cite{liu2023gres}, and OmniLabel's AP$_{\text{descr}}$ computes average precision over free-form descriptions including negative ones~\cite{schulter2023omnilabel}. Segmentation reports two aggregates that differ in weighting: gIoU (also called mIoU) averages per-sample IoU and so weights every expression equally, whereas cIoU accumulates intersection and union over the whole set before dividing and so weights large objects more; reasoning and referring segmentation papers report both~\cite{lai2024lisa}, often with Precision@$X$ for the distribution of mask quality and a boundary IoU for edge accuracy~\cite{reid2026moondream}. Video referring segmentation uses the $\mathcal{J}\&\mathcal{F}$ composite of Section~\ref{sec:object_tracking}.

Table~\ref{tab:p3_s6_5_2_grounding_performance} reports the two variants on which the families can be compared. On REC the step from proposal ranking and modular reasoning (MAttNet) to end-to-end regression (TransVG) and then to grounded pre-training (MDETR, Grounding DINO) is about 14 points of RefCOCO Acc@0.5, and the generalist MLLMs (Qwen2.5-VL, InternVL~2.5) sit at or above the best specialist without benchmark-specific training. On RES the grounding MLLMs with a mask decoder have passed the end-to-end specialists (CRIS to Sa2VA is eleven cIoU points on RefCOCO), and Sa2VA's video columns show the same model carrying over to referring video segmentation; on ReasonSeg, LISA-7B reaches 52.9 gIoU after fine-tuning where open-vocabulary segmenters stay below 30~\cite{lai2024lisa}. The tool-augmented and generative families are evaluated zero-shot on subsets of these benchmarks and are not directly comparable here. What the table cannot show is cross-variant breadth: methods strong on image REC routinely degrade on video, reasoning, and 3D grounding, and no single architecture leads all five variants.

\begin{table}[!ht]
    \centering
    \fontsize{9}{11}\selectfont
    \setlength{\tabcolsep}{1.5mm}
    \caption{Representative grounding methods. REC: Acc@0.5 (\%) on RefCOCO val / testA / testB and RefCOCOg val, from the Grounding DINO~\cite{liu2024grounding}, Ferret~\cite{you2023ferret}, and Qwen2.5-VL~\cite{bai2025qwen2.5vl} comparisons. RES: cIoU (\%) on the RefCOCO / RefCOCO+ / RefCOCOg val splits and $\mathcal{J}\&\mathcal{F}$ on MeViS and Ref-DAVIS17, from the Sa2VA comparison~\cite{yuan2025sa2va} (CRIS from its own paper, ResNet-101).}
    \label{tab:p3_s6_5_2_grounding_performance}
    \begin{tabular}{L{0.26\textwidth} L{0.22\textwidth} C{0.10\textwidth} C{0.10\textwidth} C{0.10\textwidth} C{0.10\textwidth}}
    \toprule
    \rowcolor{lightgrey} \multicolumn{6}{l}{$\bullet$ \textit{Referring expression comprehension (Acc@0.5)}} \\
    \multirow{2}{*}{\textbf{Method}} &  \multirow{2}{*}{\textbf{Family}} & \multicolumn{3}{c}{\textbf{RefCOCO}} & \textbf{RefCOCOg} \\
    & & val & testA & testB &  val \\
    \midrule
    MAttNet~\cite{yu2018mattnet}             & (B) modular attention      & 76.7 & 81.1 & 70.0 & 66.6 \\
    TransVG~\cite{deng2021transvg}           & (C) end-to-end Transformer & 81.0 & 82.7 & 78.4 & 68.7 \\
    MDETR~\cite{kamath2021mdetr}             & (D) grounded pre-training  & 86.8 & 89.6 & 81.4 & 81.6 \\
    Grounding DINO-L~\cite{liu2024grounding} & (D) grounded pre-training  & 90.6 & 93.2 & 88.2 & 86.1 \\
    Ferret-13B~\cite{you2023ferret}          & (E) grounding MLLM         & 89.5 & 92.4 & 84.4 & 85.8 \\
    Qwen2.5-VL-72B~\cite{bai2025qwen2.5vl}   & (F) generalist MLLM        & 92.7 & 94.6 & 89.7 & 89.9 \\
    InternVL2.5-78B~\cite{chen2024expanding} & (F) generalist MLLM        & 93.7 & 95.6 & 92.5 & 92.7 \\
    \midrule
    \rowcolor{lightgrey} \multicolumn{6}{l}{$\bullet$ \textit{Referring segmentation (cIoU) and video referring segmentation ($\mathcal{J}\&\mathcal{F}$)}}\\
    \hiderowcolors \textbf{Method} & \textbf{Family} & \textbf{RefCOCO} & \textbf{RefCOCO+} & \textbf{RefCOCOg} & \textbf{MeViS / Ref-DAVIS17} \\
    \midrule
    CRIS~\cite{wang2022cris}                 & (C) end-to-end, CLIP       & 70.5 & 62.3 & 59.9 & -- \\
    LISA-7B~\cite{lai2024lisa}               & (E) mask-decoder MLLM      & 74.9 & 65.1 & 67.9 & -- \\
    GLaMM-7B~\cite{rasheed2024glamm}         & (E) grounded conversation  & 79.5 & 72.6 & 74.2 & -- \\
    VISA-13B~\cite{yan2024visa}              & (E) video reasoning MLLM   & --   & --   & --   & 44.5 / 70.4 \\
    Sa2VA-8B~\cite{yuan2025sa2va}            & (E) MLLM + SAM~2           & 81.9 & 76.5 & 78.9 & 46.9 / 75.2 \\
    \bottomrule
    \end{tabular}
\end{table}

\subsubsection{Challenges}

\textbf{Ambiguity and interaction.} ``The guy on the right'' depends on the viewpoint, and ``the red cup'' presumes one salient red cup. People resolve this through shared context or a follow-up question; grounding models make an arbitrary choice or fall back on a dataset bias such as the largest instance. Recognizing that an expression matches several entities, abstaining under uncertainty, and asking for clarification are outside almost every current benchmark, which is one-shot (one query, one image), while real use is a dialogue with memory of what was grounded before.

\textbf{Reasoning and commonsense.} ``The tool needed to open the box'' cannot be resolved from the image alone, and adversarial evaluation shows that much of current REC accuracy rests on shortcut matching rather than reasoning~\cite{dong2026ref}. LLM-based models carry the knowledge but may answer ``a key'' when none is visible, and vision-only models lack it entirely; reasoning-driven segmentation~\cite{lai2024lisa} and agentic active perception~\cite{zheng2025deepeyes} are the two bridges. Video adds temporal and causal structure (``the man who picks up the phone and then leaves'', ``the car that caused the accident'') that current video grounding, which attends mostly to appearance and simple motion, does not capture.

\textbf{Open world and 3D.} Open-vocabulary detectors and CLIP-based models have blind spots on rare objects, fine-grained distinctions (viola versus violin), proper nouns, and figurative references (``the blue crown'' for a cap), and per-category analyses of large multimodal models show a strong bias toward common training categories~\cite{chen2025revisiting}. Description-conditioned pre-training~\cite{li2023desco} and generative pre-training~\cite{bagchi2025refereverything,gabeur2026image} are the routes to long-tail coverage that do not require labeling it. Functional references (``the driver's seat'') also imply spatial context that only 3D understanding resolves, and the 3D-language work surveyed above~\cite{hong20233d,zhen20243d} will grow with augmented reality and robotics.

\textbf{Efficiency and evaluation.} The strongest systems are multi-model tool chains or LLMs with vision adapters and do not fit AR glasses or robots; compact models~\cite{xiao2024florence,reid2026moondream} and non-autoregressive selection~\cite{cao2025refer} show that accuracy does not require scale, but the accuracy-efficiency frontier is unmapped. Evaluation has the saturation problem described under Datasets and Metrics above, and the new benchmarks, each with its own protocol, risk a cycle of leaderboard chasing; a single protocol spanning image, video, 3D, and dialogue, scoring rejection and clarification as well as localization, does not yet exist.



\section{Summary and Outlook}
\label{sec:p3_s6_6}

This chapter followed spatial perception from the elemental to the structured: recognizing the entities in a scene, fixing them in metric geometry, binding them into relations, layouts, and maps, carrying that structure through time, and supplementing vision with sound and language. Three observations recur across these tasks.

\textbf{Priors fill gaps that observation leaves open.} Monocular depth, surface normals, and camera motion are underdetermined by a single image: scale, reflectance, and the geometry behind an occluder are not fixed by the pixels (§\ref{sec:depth_estimation}, §\ref{sec:visual_odometry}). Learned models increasingly resolve this ambiguity with strong priors, which makes them confident on familiar scenes and fragile on transparent or reflective materials, moving objects, adverse conditions, and domains their training data under-represent. Distinguishing what a percept observes from what a prior supplies is a requirement the later chapters inherit, and Chapter~\ref{chapter:p5_discussions} returns to it as a question of observability.

\textbf{Foundation models move the bottleneck rather than remove it.} Open-vocabulary detectors and segmenters, zero-shot scene classifiers, and 3D geometry foundation models have widened what perception covers without task-specific training, yet the hardest cases persist at the long tail: small and rare objects, unseen instances in clutter, and metric judgments of distance, size, and angle, on which multimodal models remain weak even when their semantic predictions are strong (§\ref{sec:spatial_relation_extraction}). The accuracy--efficiency tension noted for detection recurs in every task of the chapter.

\textbf{Evaluation rewards the easier part of each problem.} Many benchmarks are synthetic, saturated, or confined to a few curated scenes; tracking metrics do not measure robustness to occlusion length (§\ref{sec:object_tracking}), temporal-localization metrics treat fuzzy boundaries as exact, and grounding benchmarks pose a single query against a single image rather than an interaction (§\ref{sec:vl_grounding}).

What the chapter delivers is structured spatial evidence: entities with positions, relations, trajectories, and maps. Chapter~\ref{sec:reasoning} takes that evidence as input and asks what can be inferred and planned from it, and Chapter~\ref{sec:generation} runs the mapping in the opposite direction.
\newpage
\chapter{Spatial Reasoning}
\label{sec:reasoning}

\FramedInitial{S}{patial} reasoning is the step that turns perceived structure into decisions: given detections, depths, poses, and a task, it works out the relations, the causes, and the actions that follow, the why and how of spatial interaction rather than merely the what and where. It is where perception becomes goal-directed, supporting inference, planning, and problem-solving in physical and virtual environments.

The chapter is organized along the spectrum by which spatial reasoning has actually developed, from explicit symbolic models to learned and foundation-model-driven ones. At one end, reasoning runs over hand-built representations of geometry, physics, and logic, whose structure is inspectable and whose conclusions can be certified. In the middle, it emerges from data, as neural models acquire spatial competence directly rather than being told the rules. More recently, large pre-trained models supply reasoning as a general capability to be elicited and steered. Hybrid approaches recombine these, pairing the semantic reach of foundation models with the precision and verifiability of symbolic and learned components; a closing survey of benchmarks then records how the field measures spatial reasoning as it shifts from passive observation toward active, situated interaction.

\section{Rule-based Symbolic Spatial Reasoning}
\label{sec:p3_s7_1}

\FramedInitial{R}{ule-based} symbolic spatial reasoning covers the methods that operate over \emph{explicit, model-based representations} of space (geometric primitives, physical laws, graph connectivity, formal constraints, action rules, or causal mechanisms) rather than over learned latent representations alone. What sets this family apart is that reasoning happens inside a model that can be inspected: its state and transitions are legible, so claims about feasibility, optimality, or safety can be backed by certificates rather than empirical averages, and when a plan fails it fails legibly, as an infeasible inverse kinematics, a violated constraint, a missing precondition, or an inconsistent causal attribution, not as an opaque correlation.

The section surveys six kinds of explicit structure, ordered from the most concrete to the most abstract: geometric and algebraic models of position, orientation, motion, and feasibility; physical models of contact, friction, and stability; topological and graph abstractions of space that turn planning into search; formal constraint systems whose solvers return certified plans or proofs of infeasibility; rule-based and programmatic logic that sequences and supervises actions; and causal models that explain, predict counterfactually, and prescribe interventions.
% Table~\ref{tab:p3_s7_1_paradigm_comparison} at the end of the section compares the six.

These paradigms are not isolated alternatives. Modern robotic systems routinely compose them: geometric solvers serve as feasibility oracles inside constraint-based search; behavior trees wrap symbolic plans for fault tolerance; causal action languages compile to SAT; topological maps anchor both classical and learning-based navigation policies. The unifying thread is that each paradigm carries explicit structure (a state, an operator, a constraint, a rule, or a mechanism) that can be inspected, composed, and reasoned about.

\subsection{Geometric and Algebraic Reasoning}
\label{subsec:p3_s7_1_1}
% - Coordinate transformation, workspace computation, FoV estimation
% - Euclidean, projective geometry, Jacobian-based reasoning
% - Applications: reachability analysis, sensor coverage, grasp range

Geometric and algebraic reasoning uses explicit mathematical models of position, orientation, shape, and motion to turn local measurements (image points, depths, joint angles, odometry) into quantities that can be analyzed and verified: where things are in a common frame, what a sensor can see, what a mechanism can reach, and whether a grasp will hold. The same models extend to time-varying settings, where relationships evolve under kinematic and differential constraints and the question becomes not only where entities can be but how they can move. Precision, interpretability, and domain independence have made this the common substrate of robotics~\cite{Murray1994RoboticManipulation,Lynch2017ModernRobotics,siciliano2010robotics}, AR/VR~\cite{hartley2003multiple,azuma1997survey}, and autonomous systems~\cite{Thrun2005Probabilistic,sola2018micro,LaValle2006Planning}, and the source of the hard constraints that learned components are checked against. Figure~\ref{fig:p3_s7_1_1_geo_algebra_abstract} gives the abstract formulation.

\begin{figure}[!ht]
    \centering
    \includegraphics[width=0.99\linewidth]{figures/part3/section7/711_geo.pdf}
    \caption{Abstract task formulation of geometric and algebraic reasoning: (i) local spatial measurements and system models are mapped by (ii) a core of reusable symbolic formalisms to (iii) analyzable, verifiable spatial outputs.}
    \label{fig:p3_s7_1_1_geo_algebra_abstract}
\end{figure}

\subsubsection{Task Definition}
Given measurements $\{\mathbf{z}_k\}$ expressed in local frames (pixels, ranges, joint readings, pose increments) and a system model $\Theta$ (sensor calibration, kinematic parameters, environment geometry), geometric reasoning computes a world-frame quantity and a feasibility verdict,
\begin{equation}
    \mathbf{x} = g(\{\mathbf{z}_k\};\, \Theta), \qquad
    \psi(\mathbf{x};\, \Theta) \in \{\mathrm{true}, \mathrm{false}\},
\end{equation}
where $\mathbf{x}$ may be a pose, a reconstructed surface, a workspace or visibility region, or a trajectory, and $\psi$ is a chosen predicate, such as reachability, visibility, or force closure, checked against the model. In the time-varying case the inputs are indexed over time and $\mathbf{x}$ becomes a trajectory with reachable sets and feasibility certificates along it.

\subsubsection{Methodology}

\paragraph{Transformation and Pose Representation.}
Rigid-body relationships between sensors, objects, and mechanisms are elements of the special Euclidean group
\begin{equation}
\mathbf{T} =
\begin{bmatrix}
\mathbf{R} & \mathbf{t} \\
\mathbf{0}^\top & 1
\end{bmatrix} \in SE(3),
\end{equation}
where $\mathbf{R} \in SO(3)$ is a rotation and $\mathbf{t} \in \mathbb{R}^3$ a translation; composing them moves measurements across frames. Early robotics parameterized serial chains through the Denavit--Hartenberg convention~\cite{Denavit1955Kinematic,Paul1981Robot}, which gives forward kinematics from a minimal set of link parameters. Lie-group formulations instead write motion as exponentials of screw axes, the product-of-exponentials representation~\cite{Murray1994RoboticManipulation,Lynch2017ModernRobotics}, which is coordinate-free and behaves well where the DH parameterization has discontinuities. Quaternions~\cite{Shoemake1985QuaternionCurves,Horn1987Orientation} and dual quaternions~\cite{Daniilidis1999DualQuat,horn2021online} give compact, numerically stable encodings of rotation and screw motion for calibration and sensor fusion.

\paragraph{Projective and Multi-View Geometry.}
Reasoning from images must recover 3D structure and camera motion from 2D measurements under perspective distortion and scale ambiguity. The constraint between corresponding points in two views is the epipolar condition
\begin{equation}
\mathbf{x}'^\top \mathbf{F} \mathbf{x} = 0,
\end{equation}
with $\mathbf{F}$ the fundamental matrix. The eight-point algorithm~\cite{LonguetHiggins1981Scene} solves for $\mathbf{F}$ linearly, Hartley's normalization~\cite{Hartley1997EightPoint} made it numerically stable, and Nistér's five-point solver~\cite{Nister2004FivePoint} reaches the minimal case for calibrated cameras. Zhang's homography-based calibration~\cite{zhang2000flexible} removed the need for special rigs; the field is consolidated in~\cite{hartley2003multiple}, and structure-from-motion systems such as COLMAP~\cite{schonberger2016structure} scale it to thousands of images, with RANSAC~\cite{Fischler1981RANSAC} enforcing the algebraic constraints against outliers throughout.

\paragraph{Jacobian and Differential Formulations.}
For motion rather than static configuration, the map from configuration to task space,
\begin{equation}
\mathbf{x} = f(\mathbf{q}),
\end{equation}
is differentiated to relate velocities,
\begin{equation}
\dot{\mathbf{x}} = \mathbf{J}(\mathbf{q}) \dot{\mathbf{q}},
\end{equation}
and the Jacobian $\mathbf{J}(\mathbf{q})$ carries the information about feasibility, singularities, and controllability along a motion. Whitney's resolved motion rate control~\cite{Whitney1969ResolvedMotion} used it to command Cartesian velocities without closed-form inverse kinematics; Nakamura's damped least squares~\cite{Nakamura1986IKSingularity} trades exactness for stability near singularities; redundancy resolution~\cite{Liegeois1977Redundancy,Siciliano1991Redundancy} uses null-space motion for secondary objectives; and Khatib's operational-space formulation~\cite{Khatib1987OperationalSpace} separates task and null-space dynamics so that objectives are written directly in task space.

\paragraph{Visibility and Coverage Geometry.}
What a sensor at pose $p$ can observe is the set
\begin{equation}
\mathcal{V}(p) = \{ x \in \mathcal{E} \mid \text{line}(p,x) \cap \mathcal{O}
= \emptyset \},
\end{equation}
with $\mathcal{E}$ the environment and $\mathcal{O}$ the obstacles. The Art Gallery theorem~\cite{Chvatal1975ArtGallery,ORourke1987ArtGallery} bounds the number of viewpoints needed for full coverage; polygon visibility~\cite{Lee1983Visibility} and shortest-path queries~\cite{Guibas1987ShortestPath} made the computations algorithmic; and Tarabanis \etal~\cite{Tarabanis1995SensorPlanning} cast sensor planning as geometric optimization over fields of view modeled as cones, frusta, or polygons.

\paragraph{Workspace and Reachability Analysis.}
The reachable workspace of a mechanism is
\begin{equation}
\mathcal{R} = \{ \mathbf{x} \mid \exists\, \mathbf{q} \in \mathcal{C}_{\mathrm{free}}
\text{ such that } \mathbf{x} = f(\mathbf{q}) \},
\end{equation}
with $\mathcal{C}_{\mathrm{free}}$ the feasible configuration space. Roth~\cite{Roth1975Performance} and Gosselin~\cite{Gosselin1990Singularity} added performance and dexterity measures so that manipulability is characterized along with reachability; Merlet~\cite{Merlet1998Workspace} classified workspace types; Zlatanov \etal~\cite{Zlatanov1994Constraints} analyzed constraint-induced singularities; and sampling-based reachability maps~\cite{lu2026richmap} trade closed-form guarantees for empirical coverage in high-dimensional systems.

\paragraph{Contact and Grasp Geometry.}
Whether a set of contacts can hold an object is decided in wrench space. Ferrari and Canny~\cite{Ferrari1992Grasp} model the forces each contact can apply as a convex set of wrenches $\mathcal{W}$; the grasp is force-closed if
\begin{equation}
\mathbf{0} \in \text{int}\,\text{conv}(\mathcal{W}),
\label{eq:p3_s7_1_1_force_closure}
\end{equation}
so that any external wrench can be resisted. GraspIt!~\cite{Miller2004GraspIt} made these metrics computable for arbitrary hands and objects; task-oriented measures~\cite{Lin2015TaskGrasp} weight the wrench space by what the task requires; and Dex-Net~\cite{Mahler2017DexNet} samples analytic grasps at scale to train data-driven grasp predictors that keep the geometric criterion as their label.

\subsubsection{Benchmarks}
Geometric reasoning is evaluated against geometrically grounded ground truth: camera poses and intrinsics, depth maps and meshes, object models, grasp annotations. The pose, odometry, SLAM, and multi-view reconstruction benchmarks (TUM RGB-D~\cite{Sturm2012TUMRGBD}, EuRoC~\cite{Burri2016EuRoC}, TUM-VI~\cite{Schubert2018TUMVI}, KITTI~\cite{KITTI-GeigerLU12}, ETH3D~\cite{Schoeps2017ETH3D}, Tanks and Temples~\cite{knapitsch2017tanks}, DTU~\cite{Aanaes2016DTU}, Middlebury~\cite{Seitz2006Middlebury,Scharstein2002Stereo}, HPatches~\cite{Balntas2017HPatches}) and the long-term localization suites under seasonal and day-night change (Aachen, RobotCar Seasons, CMU Seasons~\cite{Sattler2018Benchmarking}, InLoc~\cite{Taira2018InLoc}, 4Seasons~\cite{wenzel2021seasons}) are those of Section~\ref{sec:visual_odometry} and Section~\ref{sec:mapping_slam}, where their protocols and headline results are tabulated. Specific to the reasoning problems above are BAL~\cite{agarwal2010bundle}, which standardizes large bundle-adjustment problems for solver evaluation; the Active Vision Dataset~\cite{Johns2016AVD}, ScanNet~\cite{dai2017scannet}, and HM3D~\cite{Ramakrishnan2021HM3D} for next-best-view, coverage, and exploration protocols; Cornell~\cite{Jiang2011Cornell}, Jacquard~\cite{Depierre2018Jacquard}, and GraspNet-1Billion~\cite{Fang2020GraspNet} for grasp quality (rectangle labels, simulated grasp trials, and dense 6-DoF annotations with seen / unseen / novel splits respectively), with the YCB object set~\cite{Calli2015YCB} as the shared physical test set; and MoveIt~\cite{Cohen2012MoveIt}, whose reachability tooling is the de facto harness for inverse kinematics and motion-planning evaluation.

\subsubsection{Challenges and Trends}
The formulations above assume exact geometry, and real sensing is noisy, incomplete, and subject to calibration drift~\cite{Sturm2012TUMRGBD,Burri2016EuRoC}; bundle adjustment~\cite{agarwal2010bundle} and reachability analysis scale poorly with problem size; many methods fail near singularities or degenerate feature configurations~\cite{Nakamura1986IKSingularity,Hartley1997EightPoint}; the frameworks are strongest for static rigid bodies and approximate deformable and human-centric geometry~\cite{SicilianoKhatib2016}; and evaluation is concentrated on a few canonical datasets and manipulators, which limits cross-domain comparison.

The response has been to keep the geometry and wrap it in probabilistic or learned machinery: factor-graph SLAM and robust bundle adjustment attach uncertainty to the symbolic model (Section~\ref{sec:mapping_slam}); learned correspondence priors are refined by RANSAC~\cite{Fischler1981RANSAC}; and the long-term benchmarks~\cite{wenzel2021seasons,Sattler2018Benchmarking} push the same pipelines toward continual adaptation under changing conditions.




\subsection{Physical Reasoning}
\label{subsec:p3_s7_1_2}
% - Concerned with physical laws governing spatial environments (mechanics, dynamics, collisions, friction, stability, etc.).
% - Newtonian mechanics, quasi-static models, collision simulation
% - Applications: grasp stability, push/pull prediction, object stacking

Physically-aware spatial reasoning predicts how a spatial configuration evolves under physical law: rigid-body dynamics, contact and friction, and the conditions for static stability. Will a grasp hold? How will an object move when pushed? Will a stack of blocks stand or fall? Answering these from a model of mechanics, rather than from observed correlations, is what lets a robot anticipate the outcome of an interaction before committing to it; Figure~\ref{fig:physical-reasoning} shows the three canonical cases.

\subsubsection{Task Definition}
Let the environment at time $t$ be the state
\begin{equation}
    \mathbf{s}_t = (\mathbf{q}_t, \mathbf{v}_t, \theta),
\end{equation}
with generalized coordinates $\mathbf{q}_t \in \mathbb{R}^n$, velocities $\mathbf{v}_t=\dot{\mathbf{q}}_t$, and physical parameters $\theta$ (masses and inertias in the mass matrix $M(\mathbf{q})$, friction coefficients $\{\mu_k\}$, restitution, geometry). An action $\mathbf{u}_t$ (a push, a grasp) enters as a generalized force $\boldsymbol{\tau}_{\mathrm{ext},t+1}$, and the next state follows from a \emph{contact-implicit time-stepping} update:
\begin{equation}
\begin{aligned}
M(\mathbf{q}_{t})\,\frac{\mathbf{v}_{t+1}-\mathbf{v}_{t}}{\Delta t}
+ \mathbf{h}(\mathbf{q}_{t},\mathbf{v}_{t+1})
&= \boldsymbol{\tau}_{\mathrm{ext},t+1}
+ \sum_{k=1}^{K}\Big(J_{n,k}^\top \lambda_{n,k} + J_{t,k}^\top \boldsymbol{\lambda}_{t,k}\Big), \\
\mathbf{q}_{t+1}&=\mathbf{q}_{t}+\Delta t\,\mathbf{v}_{t+1},
\end{aligned}
\end{equation}
where $\mathbf{h}$ collects Coriolis, centrifugal, and gravitational terms, and for each potential contact $k$ with gap function $\phi_k(\mathbf{q})$ and Jacobians $J_{n,k},J_{t,k}$ the normal and tangential contact forces satisfy non-penetration, complementarity, and Coulomb friction:
\begin{equation}
\begin{aligned}
\phi_k(\mathbf{q}_{t+1}) \ge 0,\quad \lambda_{n,k}&\ge 0,\quad \lambda_{n,k}\,\phi_k(\mathbf{q}_{t+1})=0,\\
\|\boldsymbol{\lambda}_{t,k}\| &\le \mu_k\,\lambda_{n,k}.
\end{aligned}
\label{eq:p3_s7_1_2_contact}
\end{equation}
Physical reasoning is the computation, or approximation, of this transition $\mathbf{s}_{t+1}=\mathcal{T}(\mathbf{s}_t,\mathbf{u}_t)$ and of the queries derived from it (stability, collision, reachability). When motion is slow enough that inertial terms vanish, the \emph{quasi-static} form
\begin{equation}
\sum_{k=1}^{K}\Big(J_{n,k}^\top \lambda_{n,k} + J_{t,k}^\top \boldsymbol{\lambda}_{t,k}\Big) + \boldsymbol{\tau}_{\mathrm{ext}} - \mathbf{h}(\mathbf{q},\mathbf{0}) \approx \mathbf{0},
\qquad \|\boldsymbol{\lambda}_{t,k}\| \le \mu_k\,\lambda_{n,k},
\end{equation}
reduces reasoning to an equilibrium feasibility test, which is the setting of the classical analytic models of grasping, pushing, and stacking.

\begin{figure}[t]
  \centering
  \includegraphics[width=0.95\linewidth]{figures/part3/section6/physical_reasoning.pdf}
  \caption{Physics-aware spatial reasoning in robotics: (a) grasp stability from contact forces and friction cones, (b) planar pushing under quasi-static assumptions, and (c) stacking stability under equilibrium and support constraints.}
  \label{fig:physical-reasoning}
\end{figure}

\subsubsection{Methodology}
Model-based physical reasoning is carried out either numerically, in a physics engine that integrates the dynamics above, or analytically, with closed-form criteria for the quasi-static cases. Both rest on the same Newtonian foundation and differ in whether fidelity or interpretability is the priority.

\paragraph{Physics Engines.}
MuJoCo~\cite{todorov2012mujoco}, Bullet~\cite{coumans2021}, ODE, and DART~\cite{DART2023} integrate rigid-body dynamics with contact by time stepping, and they resolve contact in one of two ways. Penalty methods treat a contact as a compliant spring-damper that resists interpenetration,
\begin{equation}
f_{n,k} = \max\!\left(0, k_p\,\delta_k + k_d\,\dot{\delta}_k\right),
\end{equation}
with penetration depth $\delta_k$, stiffness $k_p$, and damping $k_d$, and regularize friction similarly; they are cheap and simple, but stiff parameters destabilize the integration and soft ones introduce artificial compliance, which is acceptable in graphics and approximate planning but not in control. Constraint-based methods enforce Eq.~\eqref{eq:p3_s7_1_2_contact} directly as complementarity constraints~\cite{trinkle1997dynamic}, solved per step as a linear or nonlinear complementarity problem~\cite{cottle2009linear}, a convex quadratic program~\cite{lin1987iterative}, or a second-order cone program for exact friction cones~\cite{alizadeh2003second} (projected Gauss--Seidel and sequential-impulse schemes are the usual implementations). They give physically consistent, nearly penetration-free contact forces at higher cost and with harder scaling in the number of contacts, so engines in practice hybridize: hard constraints on normal forces, regularized tangential friction, warm starting, contact aggregation, and stabilization terms.

\paragraph{Analytic Quasi-Static Models.}
For the three settings of Figure~\ref{fig:physical-reasoning}, closed-form criteria replace simulation. At a point contact the admissible forces form the friction cone
\begin{equation}
\mathcal{F}_k = \{(f_{n,k},\mathbf{f}_{t,k}) \mid f_{n,k} \ge 0, \ \|\mathbf{f}_{t,k}\|\le \mu_k f_{n,k}\},
\end{equation}
which a polyhedral approximation turns into linear constraints. Mapping each contact's cone through its moment arm into wrench space and taking the convex hull gives the grasp wrench space, and the grasp is force-closed when the origin lies strictly inside it, Eq.~\eqref{eq:p3_s7_1_1_force_closure}~\cite{bicchi2000robotic}; form closure is the stronger, purely geometric condition in which contact geometry alone prevents motion, and the Ferrari--Canny $\epsilon$-metric (the distance from the origin to the hull boundary) grades grasp quality. For planar pushing the limit surface relates the applied wrench to the resulting slip, and its ellipsoidal approximation yields a motion cone of object velocities reachable from a given push, the basis of stable pushing and closed-loop push planning. For stacking, stability under gravity requires that the gravitational wrench be balanced by the supporting contacts~\cite{groth2018shapestacks}: without friction, the center of mass must project into the support polygon; with friction, the contacts' friction cones must generate a wrench cone containing the negative of the gravitational wrench.

\subsubsection{Benchmarks}
PHYRE~\cite{bakhtin2019phyre} poses 2D physics puzzles in which an object must be placed so that simulated physics reaches a goal, scored by the log-weighted area under the success-rate curve as the attempt budget increases (AUCCESS) within and across puzzle templates; ShapeStacks~\cite{groth2018shapestacks} labels 3D stacks as stable or unstable; PlasticineLab~\cite{huang2021plasticinelab} is a differentiable-physics suite for deformable manipulation; and the video suites IntPhys~\cite{riochet2018intphys}, CLEVRER~\cite{Yi2020CLEVRER}, and Physion~\cite{Bear2021Physion} test prediction and counterfactual judgment of physical outcomes, with ComPhy~\cite{chen2022comphy} asking for hidden properties such as mass to be inferred from motion. Grasp quality under quasi-static analysis is evaluated on the grasping benchmarks of §\ref{subsec:p3_s7_1_1} (Cornell, YCB, GraspNet-1Billion). Metrics are predictive: position or trajectory error against the simulated outcome, success rate (stable stack, successful push or grasp), binary stability accuracy, and physical-plausibility terms such as momentum conservation and constraint violation.

\subsubsection{Challenges and Trends}
The parameters that matter most, mass, friction, elasticity, internal stability, are rarely observable. They must be inferred from interaction, from how far an object moves when pushed or how it responds to a collision; ComPhy~\cite{chen2022comphy} tests exactly this, and most methods still either assume the parameters known, as classical simulation does, or approximate them in ways that fail outside the lab. Latent-dynamics learning~\cite{battaglia2016interaction} and system identification~\cite{an1985estimation,muralidhar2020phynet} are the two established routes.

The simulation-to-reality gap is the second limit. Small errors in restitution, micro-contact behavior, or unmodeled environmental factors compound once a policy leaves the engine, and fidelity trades against speed: high-fidelity simulation is too slow for large-scale training or real-time control, lightweight approximation too inaccurate. Hybrid physics-plus-learning models, domain adaptation between simulated and real dynamics, and uncertainty-aware transfer each address one side of this trade-off; none closes it.

Perception is the third. Geometry, material, and object state come from noisy, incomplete sensing, and a small error in shape, pose, or contact location propagates into a large error in predicted force or stability. Systems that couple vision to a physics engine, Galileo~\cite{wu2015galileo} and PhyNet~\cite{muralidhar2020phynet} among them, show that physical priors improve perception and that physical consistency can constrain it in return, but end-to-end pipelines in which the two regularize each other remain the exception.



\subsection{Topological and Graph-based Reasoning}
\label{subsec:p3_s7_1_3}
% - Roadmap construction, A*, Dijkstra, Voronoi diagrams
% - Applications: navigation planning, environment mapping, multi-agent pathfinding

Topological and graph-based reasoning replaces continuous geometry with discrete connectivity, so that path planning, exploration, and coordination become combinatorial search over nodes and edges: places, corridors, and routes rather than metric trajectories. The abstraction is what makes large-scale, dynamic, or partially known environments tractable, where maintaining globally consistent geometry is hard but connectivity and reachability are stable. Figure~\ref{fig:p3_s7_1_3_topo_graph_abstract} gives the abstract formulation.

\subsubsection{Task Definition}
Given a configuration space $\mathcal{C}$ with free subset $\mathcal{C}_{\mathrm{free}} \subseteq \mathcal{C}$ (known a priori or inferred online from sensor streams), robot constraints, and optional priors such as floor plans or road networks, most methods build a planning graph $G=(V,E,\mathbf{w})$ in which vertices $V$ are places, sampled configurations, or abstract states, edges $E$ are feasible transitions, and $\mathbf{w}: E \to \mathbb{R}_{\geq 0}$ assigns metric or semantic costs (distance, time, energy, risk); in some methods the graph is implicit and expanded locally during search. The output is a collision-free path $\pi=\langle v_0=s,\dots,v_k=g\rangle$ of minimal accumulated cost, a topological map of place connectivity (often linked to local metric submaps), or, with several agents, conflict-free trajectories $\{\pi_i(t)\}_{i=1}^N$ that minimize makespan or total cost under vertex and edge conflict constraints.


\subsubsection{Methodology}

\paragraph{Visibility and Line-of-Sight Graphs.}
Connecting configurations that can see each other reduces shortest-path planning to graph search with interpretable waypoints. The visibility graph of Lozano-P{\'e}rez and Wesley~\cite{LozanoPerez1979Visibility} links obstacle, start, and goal vertices and recovers exact Euclidean shortest paths in polygonal environments by Dijkstra or A*; the configuration-space construction via Minkowski sums~\cite{LozanoPerez1983Cspace} extends it to robots of finite size by growing the obstacles instead of shrinking the robot; continuous-Dijkstra algorithms settle the complexity at $O(n\log n)$~\cite{Hershberger1999EuclideanShortest,Mitchell2000Handbook}. Grid planners bring line-of-sight into the search itself: Field D* interpolates across cell faces for smooth, replannable paths~\cite{Ferguson2006FieldD}, Theta* updates parents by visibility to produce near-Euclidean paths on the fly~\cite{Nash2007Theta}, and Jump Point Search prunes symmetric grid paths by jumping to points where successor choices change~\cite{Harabor2011JPS}.


\begin{figure}[!ht]
    \centering
    \includegraphics[width=0.99\linewidth]{figures/part3/section7/713_topo.pdf}
    \caption{Abstract task formulation of topological and graph-based reasoning: (i) free space, sensing, start and goal states, and constraints are turned by (ii) construction and search over a graph $G=(V,E,w)$ into (iii) decision-ready graph artifacts such as collision-free paths, topological maps, coordinated multi-agent trajectories, and coverage plans. These are alternative task outputs. In the multi-agent example, rows denote vertices and columns denote successive time steps; horizontal segments represent waiting.}
    \label{fig:p3_s7_1_3_topo_graph_abstract}
\end{figure}


\paragraph{Clearance-Oriented Skeletons.}
Where safety matters more than length, free space is retracted to a skeleton that keeps connectivity and maximizes distance to obstacles. Voronoi diagrams and medial axes give maximal-clearance paths, with Fortune's sweepline~\cite{Fortune1987Voronoi} for efficient planar construction; the generalized and hierarchical Voronoi graphs~\cite{Choset2000HGVG,Choset2000GVGExploration} let a robot follow ridges with its sensors and build the clearance roadmap incrementally while detecting junctions and loops; and distance-transform updates~\cite{Felzenszwalb2012EDT,Lau2013DynamicEDT} maintain the skeleton in dynamic environments for real-time replanning.

\paragraph{Decomposition Graphs and Coverage.}
Discretizing space into cells trades geometric fidelity for simplicity and completeness. Occupancy grids~\cite{Moravec1985SonarMaps,Elfes1989Occupancy} are probabilistic cell maps updated from range data on which planning is grid search; exact decompositions serve coverage, with boustrophedon decomposition~\cite{Choset1998Boustro} partitioning free space into monotone cells and ordering them through an adjacency graph for complete, low-redundancy sweeps; and the coverage literature is systematized in~\cite{Galceran2013CoverageSurvey}.

\paragraph{Sampling-Based Roadmaps and Trees.}
In high-dimensional configuration spaces, the graph is discovered by sampling. Probabilistic roadmaps~\cite{Kavraki1996PRM} sample and connect configurations in a preprocessing phase amortized over many queries and are probabilistically complete; RRT and RRT-Connect~\cite{LaValle1998RRT,Kuffner2000RRTConnect} grow trees toward unexplored regions for single-query planning under kinematic and kinodynamic constraints; PRM* and RRT*~\cite{Karaman2011RRTStar} add connection-radius and rewiring rules that make them asymptotically optimal; and FMT*~\cite{Janson2015FMT} performs lazy dynamic programming over a batch of samples, while BIT*~\cite{Gammell2015BIT} combines heuristic search with progressively denser sample batches for anytime improvement.

\paragraph{Optimal, Any-Angle, and Incremental Graph Search.}
On a given graph, Dijkstra's algorithm and A*~\cite{Dijkstra1959,Hart1968Astar} are the optimal single-source and heuristic searches; LPA*~\cite{Koenig2004LPA} and D* Lite~\cite{Koenig2002DstarLite} repair a shortest path incrementally as edge costs change; and the any-angle planners above (Field D*, Theta*) improve path quality, while JPS accelerates optimal search on its underlying grid.

\paragraph{Multi-Agent and Hierarchical Graph Reasoning.}
With several agents on one graph, coordination is search over a time-expanded network. Conflict-Based Search~\cite{Sharon2015CBS} splits the problem into a high-level search over conflict constraints and low-level single-agent planning and is complete and optimal; ICTS~\cite{Sharon2013ICTS} searches joint cost allocations backed by individual searches; ECBS and EECBS~\cite{Barer2014ECBS,Li2021EECBS} are bounded-suboptimal with focal heuristics and better conflict handling; and SAT or ILP compilations~\cite{Surynek2016SATMAPF,Yu2016OptimalMAPF} reach global optimality on moderate fleets at the cost of solver scalability. Hierarchical and multi-resolution planning~\cite{Botea2004HPA} and contraction hierarchies~\cite{Geisberger2008CH} precompute abstraction layers that answer shortest-path queries on continental road networks in well under a millisecond, with standardized maps and instances~\cite{Sturtevant2012Benchmarks,MovingAI2025} for comparison.

\subsubsection{Benchmarks}
The MovingAI suite~\cite{Sturtevant2012Benchmarks,MovingAI2025} supplies standardized 2D maps from games, cities, and synthetic generators for single-agent search (path cost, runtime, memory, node expansions) and the MAPF instances on which CBS, ICTS, ECBS, and compilation-based solvers are compared by makespan, sum of costs, and success under time and memory limits, with the ICAPS MAPF tracker~\cite{Stern2019MAPFsurvey} as the community leaderboard. OMPL's Planner Arena~\cite{OMPL2025} compares PRM, RRT, FMT, and BIT variants on success rate, planning time, path quality, and anytime behavior, and the DIMACS shortest-path challenge~\cite{DIMACS2025} evaluates query latency, preprocessing cost, and optimality on continental road graphs. Topological mapping and pose-graph SLAM are stressed by place recognition under appearance change: RADISH~\cite{radish2003} provides classic multi-sensor logs, and Oxford RobotCar~\cite{Maddern2017RobotCar}, NCLT~\cite{Carlevaris2016NCLT}, and Nordland~\cite{olid2018singleviewplace} span seasons and weather, with KITTI and EuRoC (Section~\ref{sec:visual_odometry}) for urban and indoor pose graphs.

\subsubsection{Challenges and Trends}
Graph density sets the trade-off: sparse graphs lose connectivity or optimality, dense ones lose real-time response, and the problem sharpens in high-dimensional and lifelong settings~\cite{Karaman2011RRTStar,Koenig2002DstarLite}. Perceptual aliasing, drifting pose graphs, and dynamic obstacles complicate planning under partial information~\cite{Thrun2005Probabilistic,Stachniss2006Mapping}, and on the topological SLAM benchmarks the binding difficulty is not search but appearance robustness, with cross-season and cross-weather queries far harder than same-condition retrieval. Multi-agent coordination strains conflict resolution and time-expanded representations as fleets grow, while fairness and deadlock avoidance must be kept~\cite{Sharon2015CBS,Yu2016OptimalMAPF,Stern2019MAPFsurvey}.

The directions are structured rather than end-to-end: hierarchical and multi-resolution planning that abstracts and refines~\cite{Botea2004HPA,Geisberger2008CH}, risk-aware and formally verified planning under uncertainty~\cite{Luders2010ChanceConstraints,zhang2023safety}, bounded-suboptimal and distributed coordination for responsiveness~\cite{Barer2014ECBS,Li2021EECBS}, and neuro-symbolic methods that enrich the graph with learned perception, heuristics, and semantics~\cite{kimura-etal-2021-neuro}.





\subsection{Constraint-based Spatial Solving}
\label{subsec:p3_s7_1_4}
% - Continuous geometric/motion feasibility: identifying configurations or trajectories that satisfy geometric, dynamic, collision, and temporal constraints under a given goal.
% - CSPs, SAT/SMT, ILP/MILP, QP/NLP, IK solvers, feasible region modeling
% - Applications: manipulation planning, multi-robot coordination, assembly validation

Many robotics problems are not about finding \emph{a} plan but one that \emph{must} satisfy strict requirements: no collisions, no joint-limit violations, stable contacts, correct timing, sometimes guaranteed dynamics or safety margins. Constraint-based spatial solving states the problem as a collection of constraints, what must hold and what must not, and hands it to a solver that either returns a satisfying solution or proves that none exists. Compared with sampling- and heuristic-driven planning, this buys certificates of feasibility, an explanation of why a candidate fails, and, under an explicit model, optimality or bounded-suboptimality guarantees, which is what manipulation in clutter, multi-robot coordination in shared space, and assembly demand. Figure~\ref{fig:p3_s7_1_4_constraint_solving_abstract} gives the abstract formulation.

\begin{figure}[!ht]
    \centering
    \includegraphics[width=0.99\linewidth]{figures/part3/section7/714_constraint.pdf}
    \caption{Abstract task formulation of constraint-based spatial solving: (i) decision variables and constraints, (ii) a solver core that satisfies or optimizes them, and (iii) feasible trajectories, schedules, and assembly sequences, with feasibility or infeasibility certificates and optimality-gap bounds where supported by the solver and problem formulation. The output types are alternatives; failure of a local solver alone is not an infeasibility proof.}
    \label{fig:p3_s7_1_4_constraint_solving_abstract}
\end{figure}

\subsubsection{Task Definition}
Let $x=(\mathbf{q}_{1:T},\mathbf{u}_{1:T},\mathbf{z})$ collect continuous states and controls together with discrete mode and logic variables $\mathbf{z}$. Constraint-based solving seeks
\begin{equation}
\text{find } x\in\mathcal{X}\quad\text{s.t.}\quad g_i(x)\le 0,\ \ h_j(x)=0,\ \ \phi(\mathbf{z},x)=\top,
\end{equation}
where the inequalities $g_i$ encode geometric and dynamic feasibility (collision, joint, and dynamics limits), the equalities $h_j$ encode kinematics and contact conditions, and $\phi$ encodes task logic such as preconditions and effects. The optimization variant minimizes an objective $f(x)$ as well. The output is a feasible $x$ (a collision-free and dynamically admissible trajectory, a task-robot schedule, an assembly sequence), with an objective value or optimality gap when optimizing, or a certificate that no such $x$ exists.

\subsubsection{Methodology}
The solver families differ in the variables and constraints they handle and the guarantees they return; practical systems are rarely pure, combining discrete reasoning over actions and modes with continuous optimization over motions and contacts, and calling fast geometric oracles (collision checking, inverse kinematics) inside the solver loop.

\paragraph{Constraint Programming (CSP/CP).}
When the decisions are combinatorial (which object, which placement, which grasp) but each must pass a geometric test, constraint programming encodes the choices as variables under constraints and interleaves symbolic branching with propagation, pruning, and geometric checks. Lozano-P{\'e}rez and Kaelbling~\cite{LozanoPerez2014ConstraintTAMP} cast task-and-motion planning as constraint satisfaction over action skeletons, pruning symbolic branches as soon as a geometric feasibility test fails; Lagriffoul \etal~\cite{Lagriffoul2012ConstraintPropagation} propagate and filter continuous domains before any expensive collision or IK call, so that many failures become cheap inconsistencies; and Dantam \etal~\cite{Dantam2018IncrementalTAMP} interleave discrete search with incremental feasibility queries that reuse information across related constraints.

\paragraph{SAT/SMT Compilations.}
Compiling planning and feasibility into Boolean formulas, or into formulas over richer theories such as real arithmetic, puts clause learning, incremental solving, and conflict-driven pruning to work. dReal~\cite{Gao2013dReal} decides nonlinear real arithmetic up to a numerical tolerance ($\delta$-completeness), which gives certified answers where exact decidability is impossible; SMT-based planning for hybrid systems~\cite{Shoukry2016SMTPlanning} reasons first over an abstraction and adds constraints as counterexamples appear; SAT encodings of multi-agent pathfinding~\cite{Surynek2016SATMAPF} obtain optimal schedules on time-expanded graphs (the MAPF problem itself is treated in §\ref{subsec:p3_s7_1_3}); and SMT over linear real arithmetic extends the encodings to continuous time and to robustness against execution delays~\cite{Kolarik2023MAPFSMT,Tan2024RobustMAPF}.

\paragraph{Mixed-Integer Formulations (MICP/MILP).}
When discrete mode choices (which region, which contact mode, which precondition) interact with continuous geometry and dynamics, binary variables encode the nonconvex choice and convex constraints capture feasibility within a mode. IRIS~\cite{Deits2014IRIS} inflates large convex obstacle-free regions that serve as certified safe sets, turning collision avoidance into membership constraints with explicit margins; Graphs of Convex Sets~\cite{Marcucci2024GCSShortest} compile motion planning through clutter into a shortest-path problem over convex nodes and edges that convex optimization, with mixed-integer selection where needed, solves with certificates; and later work constructs convex sets directly in configuration space and in space-time~\cite{Marcucci2023SciRobotics,Dai2024CIRIS}. In each case a globally nonconvex feasibility problem is replaced by a sequence of convex ones without giving up the guarantee.

\paragraph{Trajectory Optimization.}
Even with the task fixed, an executable motion must satisfy kinematic limits, collision avoidance, smoothness, and often dynamics, which is a continuous, usually nonconvex, constrained optimization whose local optimum comes with optimality conditions and sensitivities. CHOMP~\cite{Zucker2013CHOMP} deforms an initial path by functional gradient descent on a signed-distance field; STOMP~\cite{Kalakrishnan2011STOMP} replaces the gradient with stochastic perturbations to handle non-smooth costs and local minima; TrajOpt~\cite{Schulman2014TrajOpt} solves a sequence of convex programs with hinge-loss collision penalties; and GPMP2~\cite{Mukadam2018GPMP2} writes planning as inference on a sparse factor graph with Gaussian-process priors, which gives scalability and incremental replanning. Increasingly these optimizers run inside the safe sets of IRIS and GCS, which supply initialization, constrain the search, or verify the result.

\paragraph{Inverse Kinematics as a Constraint Subproblem.}
``Is there a joint configuration that reaches this pose under limits and collisions?'' is the most frequent feasibility query in any hybrid planner, so IK solvers are built as fast oracles. IKFast~\cite{Diankov2010IKFast} derives closed-form solutions and generates optimized code for repeated real-time queries; TRAC-IK~\cite{Beeson2015TRACIK} combines complementary numerical strategies to raise the solve rate and avoid stalling the planner above it; and global IK by mixed-integer programming~\cite{Dai2019GlobalIK} returns certifiable solutions for full chains when a feasibility proof is needed.

\paragraph{Hybrid Logic--Geometric Planning.}
In long-horizon problems the discrete decision (what to do next) and the continuous one (how to do it) cannot be separated, and task-and-motion planners close the loop between them. Logic-Geometric Programming~\cite{toussaint2015lgp,Toussaint2017LGPcoop} couples symbolic action sequences to bound-constrained nonlinear programs through a multi-bound tree search, so that continuous feasibility and cost decide which symbolic branches are worth expanding; PDDLStream~\cite{garrett2020pddlstream} integrates black-box geometric generators (grasps, placements, motions) into symbolic planning by optimistic refinement, requesting continuous witnesses only when needed while remaining complete under suitable stream assumptions; R-LGP~\cite{Ly2023RLGP} biases the search with reachability guidance; and COAST~\cite{Vu2024COAST} uses constraints to communicate between task planner and stream generators, cutting wasted geometric proposals by an order of magnitude in planning time.

\subsubsection{Benchmarks}
Evaluation asks whether a method returns solutions that satisfy strict feasibility conditions and, increasingly, whether it can show that it did (margins, certificates, optimality gaps) under a comparable oracle budget of collision checks and IK calls. OMPL's benchmark suite and Planner Arena~\cite{Sucan2012OMPL} score planners on standardized problems with feasibility diagnostics and solution cost; the platform-independent TAMP benchmarks of Lagriffoul \etal~\cite{Lagriffoul2018TAMPBench} specify symbolic operators together with IK and collision feasibility so that a jointly feasible task-motion plan is required; and the assembly benchmarks RAMP~\cite{Collins2023RAMP} (real-world assemblies requiring geometric fit, stability under gravity, and grasp or fixture feasibility), ASAP~\cite{Tian2024ASAP} (large-scale 3D assembly sequences validated by physics-based feasibility checks), and FurnitureBench~\cite{Heo2023FurnitureBench} (long-horizon real-world furniture assembly scored by success rate and completion time) stress feasibility across geometry, stability, and sequencing at once. Multi-robot coordination under space-time collision constraints uses the MAPF benchmarks of §\ref{subsec:p3_s7_1_3}.

\subsubsection{Challenges and Trends}
The central difficulty is integrating a discrete task structure with continuous, highly nonconvex feasibility at scale and with certificates. Task-level choices (modes, grasps, placements, contact sets, schedules) interact combinatorially with geometric and dynamic constraints and fragment the feasible set; mixed-integer formulations grow expensive with the number of binary decisions~\cite{Deits2014IRIS,Marcucci2023SciRobotics}, trajectory optimizers return locally feasible motions that may violate a symbolic assumption such as the intended contact mode~\cite{Zucker2013CHOMP,Schulman2014TrajOpt}, and symbolic planners propose sequences that are geometrically unrealizable. Information flows poorly in both directions: continuous solvers rarely return explanations (active sets, separating hyperplanes, unsatisfiable cores) that the symbolic layer could use to prune, and symbolic abstractions are often too coarse to keep the continuous subproblems tractable~\cite{garrett2020pddlstream,toussaint2015lgp}. On the hybrid and assembly benchmarks this coupling, not raw search speed, is what separates planners: those that interleave symbolic branching with continuous oracles waste fewer feasibility checks and degrade more gracefully across domains.

Assurance under uncertainty is the second difficulty. Real systems need machine-checkable feasibility and fast infeasibility proofs, but nonlinear dynamics, contact, and hybrid interactions make strong certificates hard; convex surrogates and bounded-horizon relaxations buy speed at the price of completeness or tightness; execution requires real-time replanning with expensive oracles in the loop; and robustness to noise and model error calls for explicit risk and temporal-logic constraints~\cite{Blackmore2011ChanceConstraints,Raman2014STL,Ames2019CBF}. Evaluations still under-report certificate quality, oracle counts, and failure explanations.

The trend is toward structure with certificates: graphs of convex or geodesically convex safe sets that reduce global nonconvex feasibility to sequences of convex subproblems with explicit proofs~\cite{Marcucci2024GCSShortest,Dai2024CIRIS}, coupled to trajectory-optimization and MPC back-ends that return sensitivities and constraint activations for the search above them; chance constraints, temporal logic, and control barrier functions for persistent safety~\cite{Blackmore2011ChanceConstraints,Ames2019CBF}; GPU acceleration for real-time replanning; and learning that guides branching, cutting, or warm starts without discarding the guarantee~\cite{Ly2023RLGP,Vu2024COAST}.





\subsection{Rule-based and Programmatic Reasoning}
\label{subsec:p3_s7_1_5}
% - Discrete task logic and sequential decision-making: using rules/operators to specify when and in what order actions should be performed.
% - FSM, PDDL, STRIPS, Behavior Trees, If-then rules
% - Applications: multi-step task planning, logic-guided assembly, agent behavior switching

Rule-based and programmatic reasoning is the explicit, human-auditable layer that decides \emph{when} an action is allowed, \emph{in what order} actions happen, and \emph{what to do next} when conditions change, using symbolic operators, transition rules, or program-like control structures. In a robot it is the supervisory layer that sequences reusable skills (navigate, grasp, insert) and invokes geometric solvers or learned policies only when their preconditions hold, which is what makes a decision explainable (why the detour when the corridor was blocked), verifiable, and able to carry commonsense and safety constraints through a long task. Figure~\ref{fig:p3_s7_1_5_rule_programmatic_abstract} gives the abstract formulation.

\begin{figure}[t]
    \centering
    \includegraphics[width=0.99\linewidth]{figures/part3/section7/715_program.pdf}
    \caption{Abstract task formulation of rule-based and programmatic reasoning: (i) symbolic task structure, (ii) four representations for executing decision logic, and (iii) executable plans, reactive policies, and monitored execution traces with replanning after a disturbance.}
    \label{fig:p3_s7_1_5_rule_programmatic_abstract}
\end{figure}

\subsubsection{Task Definition}
A symbolic planning problem is a tuple $\Pi=\langle \mathcal{S},\mathcal{A},s_0,G\rangle$ with symbolic states $\mathcal{S}$, an operator library $\mathcal{A}$, an initial state $s_0$, and a goal $G$; each operator $a\in\mathcal{A}$ carries a precondition and effects,
\begin{equation}
a=\big(\mathrm{pre}(a),\ \mathrm{eff}^{+}(a),\ \mathrm{eff}^{-}(a)\big),\qquad s'=\big(s\setminus\mathrm{eff}^{-}(a)\big)\cup\mathrm{eff}^{+}(a),
\end{equation}
and is applicable in $s$ iff $\mathrm{pre}(a)\subseteq s$. The planning output is a sequence $\pi=\langle a_1,\dots,a_k\rangle$ whose successive application to $s_0$ reaches a state satisfying $G$; the reactive variant instead returns a policy $\mu$ that maps the current observation and symbolic state to a behavior. Depending on the formalism, the inputs also carry resources, temporal constraints, and exogenous events.

\subsubsection{Methodology}
The formalisms split by whether they \emph{plan} (search for a correct sequence before acting) or \emph{react} (choose a behavior continually at run time); deployed systems usually do both, with a planner proposing the sequence and an execution structure (FSM, behavior tree, rule engine) monitoring conditions and triggering recovery.

\paragraph{Classical Symbolic Planning (STRIPS/PDDL and Heuristic Search).}
STRIPS~\cite{Fikes1971STRIPS} fixed the representation: planning as search over symbolic states with action schemas, so that a plan can be checked against the model and the continuous motion layer sits behind a clean abstraction boundary. PDDL~\cite{McDermott1998PDDL} standardized domain and problem specification so that planners could be compared across navigation, logistics, and assembly, with extensions (ADL, numeric fluents, temporal constructs) that capture robotic requirements without abandoning explicit semantics. Graphplan~\cite{Blum1997PlanningGraph} builds a planning graph whose mutual-exclusion constraints prune incompatible action combinations early, and FF~\cite{Hoffmann2001FF} made forward search fast with relaxed-plan heuristics that ignore delete effects, which capture ordering constraints such as unstacking before stacking in object-rich domains. Task-and-motion planning (§\ref{subsec:p3_s7_1_4}) extends these planners with geometric feasibility checks~\cite{Kaelbling2011HierarchicalNow,Garrett2021TAMPReview,Zhao2024TAMPsurvey}; the common pattern is to keep the symbolic reasoning explicit, delegate feasibility to motion and IK solvers, and wrap the plan in an execution structure with monitors and recovery.

\paragraph{Finite-State Machines and Reactive Mode Switching.}
An FSM switches among a few well-defined behaviors (navigate, align, grasp) with explicit transition guards, which makes it transparent and fast but prone to state explosion as contingencies multiply. Statecharts~\cite{Harel1987Statecharts} add hierarchy and concurrency, so that complex controllers are nested reusable submachines with precise semantics; Brooks' subsumption architecture~\cite{Brooks1986Subsumption} layers behaviors so that higher ones override lower ones, a blueprint for composing navigation and avoidance without global planning, systematized as behavior-based robotics~\cite{Arkin1998BehaviorBased}; teleo-reactive programs~\cite{Nilsson1994TeleoReactive} evaluate conditions continuously and fire the highest-priority satisfied rule; and ROS tooling such as FlexBE~\cite{conner2023flexbe} provides hierarchical, event-driven state machines at deployment scale.

\paragraph{Behavior Trees.}
A behavior tree re-evaluates its conditions on every tick and chooses among structured alternatives (sequence, fallback, parallel), so that long-horizon execution survives disturbances without committing to one fixed sequence; subtrees encapsulate reusable routines, guards and fallbacks trigger immediate recovery, and the composition semantics fit perception and motion skills. BTs originated in game AI~\cite{Isla2005Halo2}; Marzinotto \etal~\cite{Marzinotto2014BT} formalized their semantics and their relation to other switching structures, and Colledanchise and {\"O}gren~\cite{Colledanchise2017BT,Colledanchise2018BTbook} showed which compositions preserve safety and robustness and consolidated the design principles. Compiling or wrapping a symbolic plan into a BT keeps the nominal task structure while exposing runtime guards and recovery branches~\cite{mart2021behaviourtree}; ARIAC systems pair PDDL planning with BT execution monitors to handle failures, occlusions, and defective parts~\cite{liu2024autonomous}; and event-driven BTs~\cite{Agis2020EDBT} react to non-periodic stimuli rather than waiting for the next tick.

\paragraph{Logic-based Execution and Production Rule Systems.}
A production system keeps a working memory of facts and fires forward-chaining rules when their conditions hold, which suits event-condition-action policies, mode switching, and safety enforcement (no-go zones, mutual exclusions, priorities) and depends on fast incremental matching. OPS5~\cite{Forgy1981OPS5} set the model; the Rete algorithm~\cite{Forgy1982Rete} caches partial matches so that large rule sets update in milliseconds as facts change; CLIPS~\cite{Riley1991CLIPS} made the toolchain portable; and Soar~\cite{Laird1987Soar} integrated large rule bases with planning and execution in one architecture. In robotics, KnowRob~\cite{tenorth2013knowrob} and CRAM~\cite{beetz2010cram,beetz2023cram} join symbolic knowledge to rule-driven execution for reasoning about actions, effects, and everyday manipulation, and LTLMoP~\cite{finucane2010ltlmop} compiles temporal-logic specifications into correct-by-construction reactive controllers whose safety and sequencing constraints are explicit at the program level.

\subsubsection{Benchmarks}
The International Planning Competition supplies the canonical PDDL domains and problems across classical, temporal and numeric, probabilistic, and HTN tracks, compared by plan quality, makespan, and runtime~\cite{McDermott1998PDDL}, with \texttt{planning.domains}~\cite{PlanningDomains} mirroring and extending the sets for reproducible evaluation. Long-horizon embodied suites test execution under partial observability and variation: BEHAVIOR-100 in iGibson and BEHAVIOR-1K in OmniGibson~\cite{Srivastava2022BEHAVIOR,li2024behavior1k} specify household activities at the predicate level, ALFRED~\cite{shridhar2020alfred} and TEACh~\cite{Padmakumar2022TEACh} score instruction-conditioned action sequences, and RoboTHOR~\cite{Deitke2020RoboTHOR} and the Habitat challenge~\cite{szot2021habitat2} pair high-level decision logic with navigation policies under SPL. ARIAC~\cite{liu2024autonomous} injects faulty parts, occlusions, and failures into kitting and assembly, and FurnitureBench~\cite{Heo2023FurnitureBench} and RoboCup@Home~\cite{Wisspeintner2009RoboCupHome} stress recovery and explainability on real tasks; qualitative solvers over RCC-8~\cite{Randell1992RCC} and Allen relations~\cite{janhunen2019allen} are tested on consistency problems, and composite pipelines reuse the TAMP and OMPL harnesses of §\ref{subsec:p3_s7_1_4}.

\subsubsection{Challenges and Trends}
The results on these benchmarks divide along the closed-world boundary. On IPC-style domains with fully specified preconditions and effects, classical planners remain strong; on long-horizon embodied tasks the limit is perception noise and grounding ambiguity rather than search, and a logically applicable action may be geometrically unrealizable because of a collision, an IK failure, or a missing grasp. Hand-authored operators, rules, and recovery logic also grow combinatorially with object count and contingency. Integrated task-and-motion planning (§\ref{subsec:p3_s7_1_4}) answers the coupling problem by interleaving symbolic search with feasibility checks or compiling symbolic commitments into optimization back-ends~\cite{LozanoPerez2014ConstraintTAMP,garrett2020pddlstream,toussaint2015lgp}, and hierarchical abstraction and compositional skill models contain the search while keeping it auditable~\cite{Garrett2021TAMPReview,Zhao2024TAMPsurvey}.

Robustness to noise and nonstationarity is why the industrial and instruction-conditioned benchmarks reward compositions rather than single formalisms: a behavior tree around a symbolic plan makes disturbances and recoveries explicit through runtime guards, fallbacks, and safety-preserving composition~\cite{Marzinotto2014BT,Colledanchise2017BT,Colledanchise2018BTbook}, and Rete-style engines~\cite{Forgy1982Rete} supply fast event-condition-action policies over streaming facts, the pattern ARIAC entries converge on~\cite{liu2024autonomous}.

Learned perception and language models increasingly propose the high-level procedure or heuristic while the symbolic structure stays as the safety shell that enforces preconditions, effects, and constraints: SayCan~\cite{ahn2022icanisay} grounds language in robot skills by value-based scoring, and plan-to-BT compilation~\cite{mart2021behaviourtree} exposes parallelism, guards, and recovery at run time. Domain modeling, heuristic guidance, and recovery synthesis are being automated piece by piece while the programmatic structure is retained for auditing, debugging, and verification.



\subsection{Causal Reasoning}
\label{subsec:p3_s7_1_6}
% - Cause--effect relationships between spatial events (``why'' and ``what-if'').
% - Structural Causal Models (SCM), Counterfactual SCM, PC Algorithm, GES, LiNGAM
% - Task dependency planning, event explanation, intervention & prevention, multi-robot coordination

Causal reasoning moves from describing \emph{what} happened to explaining \emph{why} it happened and predicting \emph{what would have happened} under a different action. For a robot the relevant outcomes are chains of spatial events (a collision, a delay, a deadlock, an unstable grasp, a near-miss in shared space), and making the causal assumptions explicit, rather than leaving them implicit in ``actions cause state changes'', is what lets the system attribute a failure to its causes rather than its symptoms, evaluate counterfactual alternatives such as a different approach direction or timing, and, with several agents, separate primary causes from downstream effects when choosing an intervention. Figure~\ref{fig:p3_s7_1_6_causal_reasoning_abstract} gives the abstract formulation.

\begin{figure}[!th]
    \centering
    \includegraphics[width=0.99\linewidth]{figures/part3/section7/716_causal.pdf}
    \caption{Abstract task formulation of causal reasoning for spatial intelligence: (i) event streams, candidate variables, prior knowledge, and interventional data feed (ii) methods that discover and verify causal structure, yielding (iii) causal models, explanations, counterfactual answers, and intervention or prevention strategies. The method families serve different roles rather than forming a serial pipeline. Verification returns a satisfying witness or an unsatisfiability proof/core where supported; counterfactual scenes keep the objects fixed while changing the blue agent's action.}
    \label{fig:p3_s7_1_6_causal_reasoning_abstract}
\end{figure}

\subsubsection{Task Definition}
The inputs are spatiotemporal event streams (agent actions, object states, contact and visibility relations, environment variables) mapped to causal variables by measurement functions (\texttt{in\_contact}, \texttt{blocked}, \texttt{priority\_violated}), optional prior knowledge from physics or commonsense, and, when available, interventional data in which a controllable variable was set deliberately. A structural causal model over these variables,
\begin{equation}
    \mathcal{M} = \langle \mathcal{G}, \{f_i\}, P(U) \rangle, \qquad X_i := f_i(\mathrm{Pa}(X_i), U_i),
\end{equation}
consists of a directed acyclic graph $\mathcal{G}$, a mechanism $f_i$ that generates each variable from its parents $\mathrm{Pa}(X_i)$ and exogenous noise $U_i$, and a noise distribution $P(U)$~\cite{Pearl2009Causality}. The outputs are the model itself when it must be learned, the explanation of an observed outcome (which variables caused it), the answer to an interventional query $P(Y \mid do(X=x))$ or a counterfactual one ($Y_{x}$ for a unit whose actual trace was observed), and an intervention or prevention strategy that changes controllable causes to reach a desired spatial outcome.

\subsubsection{Methodology}
The methods differ in where the causal structure comes from (assumed, learned from data, or written as domain knowledge) and in what they are asked (intervention effects, counterfactuals, responsibility, verification); in spatial domains they also differ in how they use locality, temporal precedence, and network connectivity as inductive bias.

\paragraph{Structural Causal Models and Counterfactuals.}
An SCM makes mechanisms explicit (``clearance determines collision risk'', ``traffic density sets delay'') and separates intervention from observation: the $do(\cdot)$ operator replaces a structural equation instead of conditioning on a value, which is what is needed when, say, congestion and delay are both driven by time of day. Spatial pipelines typically start from a graph fixed by temporal precedence and locality priors (edges among nearby regions, interacting agents, adjacent cells), estimate the mechanisms from data, and validate them against implied conditional independences and refutation tests~\cite{Peters2017Elements}. Counterfactuals follow the abduction-action-prediction recipe: infer the exogenous noise consistent with the observed trace, apply the hypothetical intervention, predict the outcome~\cite{Pearl2009Causality}. DoWhy and DoWhy-GCM~\cite{DoWhy2020,Bloebaum2024DoWhyGCM} implement identification, estimation, and refutation for this workflow on real spatiotemporal data.

\paragraph{Causal Discovery: Constraint- and Score-Based.}
When the graph is unknown but observational data are plentiful, constraint-based methods read structure off conditional independences: the PC algorithm builds the skeleton by independence tests and orients what it can, returning a Markov equivalence class, and FCI~\cite{Spirtes2000CPS} handles latent confounding and selection bias by returning a partial ancestral graph. Score-based methods instead search graph space under an objective such as BIC: GES~\cite{Chickering2002GES} greedily adds and then removes edges over equivalence classes, and its forward-backward structure often recovers mesoscale organization such as bottlenecks that drive downstream congestion. Spatial data strain both. Large conditioning sets are expensive and statistically weak, so practice restricts conditioning to a radius or $k$-hop neighborhood, forbids edges that violate adjacency or reachability priors, and uses order-independent variants (PC-stable~\cite{Colombo2014PCStable}) whose high-dimensional consistency is understood~\cite{Kalisch2007PC}; ambiguous orientations are then resolved by targeted interventions. causal-learn~\cite{zheng2024causal} and the TETRAD / Py-Tetrad implementations of FGES~\cite{Ramsey2018Tetrad,Ramsey2023PyTetrad} provide these pipelines at facility and city scale.

\paragraph{Independent-Noise Discovery (LiNGAM).}
In linear-Gaussian settings many edge directions are unidentifiable from independences alone. LiNGAM~\cite{Shimizu2006LiNGAM} assumes a linear structural model with independent non-Gaussian noise, under which the causal ordering is identifiable, and DirectLiNGAM~\cite{Shimizu2011DirectLiNGAM} estimates it by regression rather than ICA. Spatiotemporal signals need care because autocorrelation violates i.i.d. sampling assumptions and aggregation may change the noise distribution; temporal blocking, prewhitening, and locality constraints are the usual remedies, and the recovered ordering then scaffolds an SCM for intervention analysis.

\paragraph{Causal Action Languages and Plan-Space Causal Links.}
Much causal structure in robotics is designed knowledge rather than discovered: actions have direct and indirect effects, and the system must reason about change under incomplete information. The action language C\textsuperscript{+}~\cite{Giunchiglia2004CPlus} writes causal theories of action with inertia and nonmonotonic semantics, compactly describing how fluents such as \texttt{occupied(c)} or \texttt{visible(a,o)} evolve; CCALC~\cite{Lifschitz2000CCalc} compiles them to SAT so that prediction, diagnosis, and plan synthesis use the same solvers; and multi-valued fluents and external geometric predicates (collision checks, IK) are being integrated without abandoning the formal semantics~\cite{giunchiglia2023causal}. In plan space, partial-order causal-link planning~\cite{McAllester1991SNLP} represents a plan as partially ordered actions plus links $A \xrightarrow{p} B$ (action $A$ establishes precondition $p$ for $B$), which makes threats explicit and repairable by promotion, demotion, or binding constraints; UCPOP~\cite{Penberthy1992UCPOP} extended the representation to ADL and conditional effects and VHPOP~\cite{Younes2004VHPOP} added heuristics and reachability analysis. In spatial domains a causal link is often a physical dependency, ``unblock door'' enabling ``move cart'', which bridges symbolic causality and spatial feasibility.

\paragraph{Execution and Verification.}
A causal plan must be realized in time and checked. Simple temporal networks~\cite{Dechter1991STN} enforce difference constraints between events with polynomial consistency and dispatchable forms~\cite{Tsamardinos1998Dispatch}, and STNUs with dynamic-controllability analysis~\cite{Morris2001STNU} handle contingent durations, so that ordering and clearance constraints stay satisfiable under bounded disturbance. The SAT/SMT machinery of §\ref{subsec:p3_s7_1_4} enters here as an explanation device: when a schedule is unsatisfiable, the unsat core is a set of conflicting commitments, a concrete statement of which decisions caused the failure~\cite{Surynek2016SATMAPF,Yu2016OptimalMAPF}; and for hybrid dynamics $\delta$-complete SMT (dReal)~\cite{Gao2013dReal} and dReach~\cite{Kong2015dReach} verify safety and synthesize interventions with numerically sound proofs or counterexamples.

\subsubsection{Benchmarks}
No single benchmark covers the causal pipeline, and the suites stress different links of it. CLEVRER~\cite{Yi2020CLEVRER} pairs synthetic collision videos with descriptive, explanatory, predictive, and counterfactual questions, and CLEVRER-Humans~\cite{CLEVRERHumans2024} adds human-annotated causal event graphs against which model attributions can be compared; the video suites of §\ref{subsec:p3_s7_1_2} (PHYRE~\cite{bakhtin2019phyre}, Physion and Physion++~\cite{Bear2021Physion,tung2023physion}) score intervention selection and latent-property inference, and CRaFT~\cite{Saha2022CRaFT}, ContPhy~\cite{Zheng2024ContPhy}, NExT-QA~\cite{Xiao2021NExTQA}, AGQA~\cite{GrundeMcLaughlin2021AGQA,Gupta2023AGQA2}, and SUTD-TrafficQA~\cite{Xu2021TrafficQA} ask causal questions in natural language about synthetic and real-world videos, including traffic scenes. Embodied platforms with predicate-level state (BEHAVIOR-1K, ALFRED, iGibson 2.0~\cite{Li2021iGibson2}, ManiSkill2~\cite{Gu2023ManiSkill2}, RLBench~\cite{james2020rlbench}; §\ref{subsec:p3_s7_1_5}) test whether inferred dependencies improve intervention selection and recovery; the MAPF suites and Flatland~\cite{Mohanty2020Flatland} (§\ref{subsec:p3_s7_1_3}) test reasoning over space-time constraints in which one agent's decision causes another's delay; and ARCH-COMP~\cite{ARCHCOMP2024} with the dReal / dReach problem sets~\cite{Gao2013dReal,Kong2015dReach} evaluate causal safety reasoning on hybrid dynamics by proof or counterexample and solve time. Discovery itself is evaluated on recovery of the Markov equivalence class under faithfulness, with scalability bounded by conditioning-set size.

\subsubsection{Challenges and Trends}
Identification is brittle. Latent confounding, selection bias, and temporal aggregation can make a query non-identifiable from logs, and estimates become unstable when locality and precedence assumptions fail~\cite{Pearl2009Causality,Peters2017Elements}; constraint- and score-based discovery is expensive on dense telemetry with strong short-range dependence and sensitive to test strength, ordering, and regularization~\cite{Spirtes2000CPS,Colombo2014PCStable,Kalisch2007PC,Chickering2002GES}; and LiNGAM's non-Gaussianity is exactly what smoothed trajectories and aggregated traffic signals lack~\cite{Shimizu2006LiNGAM,Shimizu2011DirectLiNGAM}. The response is structure-aware discovery that constrains search by topology and locality and then resolves the remaining ambiguity by targeted robot interventions, connecting learning to actuation~\cite{Spirtes2000CPS,Chickering2002GES}.

Causal reasoning is only as good as its variables, and in embodied settings contact, occlusion, blockage, and object state must be grounded from noisy perception. The errors of that mapping propagate into SCM estimates and counterfactual conclusions~\cite{Peters2017Elements,DoWhy2020}. Knowledge-enabled frameworks such as KnowRob and CRAM (§\ref{subsec:p3_s7_1_5}) show the shape of the interface, perception populating an explicit, interpretable model whose parameters are then estimated statistically, and the trend is toward such perception-to-symbol interfaces that tolerate sensor noise~\cite{tenorth2013knowrob,beetz2010cram,beetz2023cram}.

Between a causal recommendation and a safe execution there is a verification-synthesis gap: the recommendation must become a temporally controllable schedule and a certified controller, and execution uncertainty invalidates a symbolic causal plan unless temporal networks and hybrid verification are integrated with it~\cite{Dechter1991STN,Morris2001STNU}. SCM toolchains are therefore being embedded in decision stacks so that counterfactual policies compile into action-language or plan-space structures and then into dispatchable networks~\cite{Bloebaum2024DoWhyGCM,Giunchiglia2004CPlus,Morris2001STNU}, with SAT/SMT certificates and unsat-core explanations supplying the accountability~\cite{Surynek2016SATMAPF,Gao2013dReal,Kong2015dReach}; evaluation that scores explanation quality, interventional effectiveness, and certification cost on the same task is what the field still lacks.



% \subsection{Summary}
% \label{subsec:p3_s7_1_summary}

% The six paradigms in §\ref{sec:p3_s7_1} build spatial intelligence on explicit, inspectable structure (geometric, mechanical, topological, formal, programmatic, or causal) rather than on learned correlations alone, and that choice buys four things. Plans, schedules, and explanations are artifacts that can be debugged and verified rather than distributions to be sampled. Each paradigm increasingly produces machine-checkable certificates (force-closure conditions, optimality witnesses, $\delta$-complete SMT proofs, unsat cores) on which safety-critical deployment can rely. Explicit interfaces (operators, constraints, rules, causal variables) let the paradigms compose with one another and with learned components without losing their guarantees. And where the model captures the relevant structure, classical methods generalize from a compact specification rather than from a large corpus.

% Table~\ref{tab:p3_s7_1_paradigm_comparison} sets the six side by side; they are complementary rather than redundant, each targeting a structure that the others presuppose (geometry beneath physics), wrap (behavior trees around symbolic plans), or extend (causal mechanisms over constraint solutions).

% \begin{table}[t]
%     \centering
%     \footnotesize
%     \caption{Comparison of the six foundational paradigms of rule-based symbolic spatial reasoning covered in §\ref{sec:p3_s7_1}.}
%     \label{tab:p3_s7_1_paradigm_comparison}
%     \begin{tabularx}{\linewidth}{@{}l X X X X@{}}
%         \toprule
%         \textbf{Paradigm} & \textbf{Core representation} & \textbf{Output / artifact} & \textbf{Notable strength} & \textbf{Main limitation} \\
%         \midrule
%         §\ref{subsec:p3_s7_1_1} Geometric \& algebraic         & $SE(3)$ poses, projective geometry, Jacobians, contact wrenches                  & Frame-aligned coordinates, workspace / reachability sets, feasibility certificates      & Closed-form precision; coordinate-free formulations                          & Brittle near singularities; assumes exact geometry \\
%         §\ref{subsec:p3_s7_1_2} Physical                       & Rigid-body dynamics, contact / friction laws, quasi-static equilibrium           & Force-closure verdict, motion cone, support-polygon test                                & Grounded in mechanics; analytically interpretable                            & Sim-to-real gap; latent parameters (mass, friction) often unobserved \\
%         §\ref{subsec:p3_s7_1_3} Topological / graph-based      & Roadmap / decomposition / sampling graph $G=(V,E,\mathbf{w})$                    & Collision-free path $\pi$, topological map, multi-agent trajectories                    & Combinatorial tractability; mature solvers (A*, CBS, CH)                     & Density vs.\ optimality trade-off; appearance robustness across seasons \\
%         §\ref{subsec:p3_s7_1_4} Constraint-based               & Decision variables + collision / IK / dynamics / logic constraints               & Certified-feasible plan, optimality gap, infeasibility witness                          & Strong correctness guarantees; certificate-producing                         & Mixed-integer scaling; symbol-grounding from perception \\
%         §\ref{subsec:p3_s7_1_5} Rule-based / programmatic      & STRIPS / PDDL operators, FSMs, Behavior Trees, production rules                  & Action plan, reactive policy, monitored execution trace                                 & Auditable; reactive recovery; modular composition                            & Modeling burden; brittle under perception noise without grounding \\
%         §\ref{subsec:p3_s7_1_6} Causal                         & Structural causal models, temporal networks (STN / STNU)                         & Causal model, $do(\cdot)$-effect estimate, counterfactual answer, intervention strategy & Counterfactual / responsibility queries; identifiability theorems            & Grounding causal variables; latent confounding; LiNGAM needs non-Gaussian noise \\
%         \bottomrule
%     \end{tabularx}
% \end{table}

% The same design philosophy exposes the limitations. Explicit models must be authored and maintained, they presume that the relevant structure is observable, and they degrade gracelessly when reality violates their idealizations (perception noise, latent parameters, unmodeled contact, drifting causal mechanisms). These are the conditions under which data-driven methods excel: learned pipelines absorb unobservable variables and perception noise into end-to-end objectives and generalize across regimes that would be intractable to specify by hand.

% The rest of Chapter~\ref{sec:reasoning} takes up spatial reasoning from that complementary angle: §\ref{sec:p3_s7_2} on reasoning learned from data, §\ref{sec:p3_s7_3} on reasoning elicited from foundation models, §\ref{sec:p3_s7_4} on hybrids of the two with the classical paradigms, and §\ref{sec:p3_s7_5} on how the field measures all of them. Throughout, the classical paradigms of §\ref{sec:p3_s7_1} reappear as feasibility oracles, safety shells, structural priors, and explainability layers, and the most effective systems are those that compose the two regimes: explicit structure where it is authoritative, learned structure where it is unavoidable.


% \newpage

\section{Learning-based Spatial Reasoning}
\label{sec:p3_s7_2}

\FramedInitial{L}{\textbf{earning-based spatial reasoning}} refers to the class of methods that acquire and perform spatial reasoning from data, rather than from the hand-authored geometric, mechanical, or symbolic models surveyed in §\ref{sec:p3_s7_1}.
The six subsections that follow are ordered by the \emph{locus of the structural prior}: it begins inside the learned representation (latent feature interaction, explicit programs, hard-coded symmetries, optimization layers in the forward pass), moves to the agent's action-conditioned behavior (reinforcement learning with spatial inductive biases), and ends with joint decision-making across coordinated agents.



\subsection{Implicit Neural Spatial Reasoning}
\label{subsec:p3_s7_2_1}



We define \emph{implicit neural spatial reasoning} as the family of neural approaches in which spatial relations are encoded in learned feature representations, attention patterns, latent object interactions, or architectural inductive biases, rather than specified through hand-written rules, symbolic programs, or fully explicit relational graphs.
\emph{Implicit} does not mean that these models never use boxes, masks, object proposals, or task-level labels; it means that the reasoning procedure itself is never given as a symbolic program, so the model must learn spatial dependencies from end-to-end or weakly structured supervision.
Several families below (latent scene graphs, capsule networks) expose partly inspectable structure, yet remain implicit in that their reasoning representations are learned rather than fixed by hand-written symbolic rules.
Figure~\ref{fig:p3_s7_2_1_implicit_formulation} shows the representative tasks.


\begin{figure}[!th]
    \centering
    \includegraphics[width=0.99\linewidth]{figures/part3/section7/implicit_neural_spatial_reasoning_task.pdf}
    \caption{Representative tasks in which spatial relations must be inferred without explicit symbolic supervision: (1) visual question answering over compositional scenes, (2) image captioning with spatial attention, and (3) visual relational reasoning.}
    \label{fig:p3_s7_2_1_implicit_formulation}
\end{figure}




\subsubsection{Task Definition}
Given one or more images $I$, optionally paired with a language query or task goal $q$ (and possibly boxes, masks, or proposals as features), an implicit reasoner computes a latent spatial representation and a task output,
\begin{equation}
  z=f_\theta(I,q),\qquad y=h_\psi(z),
\end{equation}
where $z$ encodes the relations relevant to the task and $y$ is an answer label, a caption, or a relation label.
In a purely implicit setting, supervision reaches the model through $y$; semi-explicit variants may additionally use object or relation labels without prescribing a symbolic reasoning program.
The methods below differ in the mechanism by which $z$ is formed (Figure~\ref{fig:p3_s7_2_1_implicit_method}).




\begin{figure}[!th]
    \centering
    \includegraphics[width=0.99\linewidth]{figures/part3/section7/implicit_neural_spatial_reasoning_method.pdf}
    \caption{Four mechanisms for learning spatial representations: (a) relational aggregation (Relation Network, GNN message passing), (b) spatial attention with multi-scale fusion (FPN), (c) conditional modulation (FiLM), and (d) geometry-aware inductive biases (capsule routing, equivariant features). Blue marks selected features, messages, and operations. The mechanisms can be combined in one model, and attention and FPN support spatial representation rather than forming complete reasoning procedures.}
    \label{fig:p3_s7_2_1_implicit_method}
\end{figure}

\subsubsection{Methodology}

\paragraph{Relational Aggregation over Latent Objects.}
The most direct approach treats a set of feature vectors as latent ``objects'' and learns to aggregate their interactions.
\textbf{Relation Networks (RN)}~\cite{santoro2017simple} score every pair: given objects $O=\{o_1,\dots,o_n\}$ (\eg, feature vectors from spatial regions of a CNN map) and an optional query $q$,
\begin{equation}
  \text{RN}(O, q) = f_{\phi}\!\left(\sum_{i,j} g_{\theta}(o_i, o_j, q)\right),
\end{equation}
where $g_{\theta}$ produces a query-conditioned relation embedding for each pair and $f_{\phi}$ maps the sum to the task output.
Considering all pairs forces the network to learn pairwise spatial and semantic relations, at $O(n^2)$ cost in the number of objects.
\textbf{Graph neural networks and latent scene graphs} generalize this from a complete graph to a learned, sparser topology, with proposals or feature regions as nodes and learned edges as spatial dependencies; iterative message passing for scene-graph generation~\cite{xu2017scene} alternates between refining node features and updating relation embeddings, learning the topology end-to-end.
Such models are the clearest \emph{semi-explicit} case in this subsection: the graph is inspectable and often trained with relation annotations (\eg, from Visual Genome), yet the intermediate relational representation is learned rather than specified by a program, and its quality is bounded by the upstream proposal stage.

\paragraph{Spatial Attention and Multi-Scale Feature Fusion.}
A second family controls \emph{where} and \emph{how strongly} the model attends across the spatial layout, and fuses evidence across scales, rather than scoring all pairs.
CNN spatial attention weights locations in a feature grid, from the hard and soft attention of \citet{xu2015show} for captioning and grounded VQA~\cite{zhu2016visual7w} to SCA-CNN's joint spatial and channel attention~\cite{chen2017sca}; vision transformers make relating image parts intrinsic through self-attention over patches~\cite{vit}, with locality-aware variants (Swin~\cite{liu2021swin}, PVT~\cite{wang2021pyramid}, Twins~\cite{chu2021twins}) trading global receptive field for cost; and Feature Pyramid Networks~\cite{lin2017feature} add a top-down pathway so that a relation coupling a fine object with a coarse region (a car ``on'' a street) can be read from one map.
These designs are representation backbones that preserve and aggregate spatial information for the reasoning built on top of them, not reasoning modules in themselves.

\paragraph{Conditional Feature Modulation.}
A third family conditions one representation on another rather than aggregating or attending.
In \textbf{FiLM (Feature-wise Linear Modulation)}~\cite{perez2018film}, a conditioning vector (\eg, an encoded query or global context) predicts affine parameters $(\gamma, \beta)$ that are applied feature-wise to a target map, so the same backbone can be steered per task.
For a query like ``What is to the left of the green cube?'', a FiLM layer can learn to amplify features to the left of regions identified as the green cube, making reasoning adaptive and context-dependent.

\paragraph{Geometry-Aware Inductive Biases.}
A fourth family bakes geometric regularities into the architecture, trading data-driven flexibility for sample efficiency where the relevant geometry is known.
\textbf{Capsule Networks}~\cite{sabour2017dynamic} replace scalar activations with vector-valued capsules whose orientation encodes instantiation parameters (pose, deformation) and whose magnitude encodes presence; \emph{dynamic routing} decides how lower-level capsules compose into higher-level ones by agreement in predicted pose, a stronger part-whole bias than translation-only convolution, though capsules have been hard to scale beyond small digit and shape datasets.
Equivariant and steerable convolutions~\cite{cohen2016group}, which commute with a chosen transformation group, are the subject of §\ref{subsec:p3_s7_2_3}; both share the cost that the symmetry must be chosen in advance.


\subsubsection{Benchmarks}

CLEVR~\cite{johnson2017clevr} is the canonical benchmark, pairing rendered scenes with multi-step relational questions whose ground-truth program structure is known; Sort-of-CLEVR~\cite{santoro2017simple} isolates a controlled relational vocabulary.
The decisive stress tests are the \emph{systematic-generalization} splits: CLEVR-CoGenT~\cite{johnson2017clevr} holds out attribute-composition combinations, and GQA~\cite{hudson2019gqa} scales compositional questions to real images.
Open-domain VQA v2~\cite{goyal2017making} and Visual7W~\cite{zhu2016visual7w} probe whether implicit reasoning transfers beyond synthetic templates to natural images and free-form language.
On CLEVR, RN reaches 95.5\% and FiLM 97.7\%; Table~\ref{tab:p3_s7_2_2_clevr_gqa} in §\ref{subsec:p3_s7_2_2} places these against the program-structured reasoners, which close the remaining gap only on in-distribution questions.


\subsubsection{Challenges and Trends}
The \emph{interpretability challenge} is intrinsic: because relations live in distributed activations rather than named symbols, it is hard to audit which relations a model actually exploits, and attention-map visualizations are only an indirect proxy.
\emph{Compositional generalization} is the most striking failure: even at 97--98\% CLEVR accuracy, the same models drop sharply on CLEVR-CoGenT~\cite{johnson2017clevr} and GQA~\cite{hudson2019gqa}, revealing dataset-specific shortcuts rather than transferable spatial concepts; the related \emph{shortcut-learning} concern in open-domain VQA, where language priors and object co-occurrence dominate the signal, is the imbalance VQA~v2's balanced design targets~\cite{goyal2017making}.
\emph{Scalability} divides the families: RN-style pairwise modeling is $O(n^2)$ in the number of objects, attention backbones trade receptive field against compute~\cite{liu2021swin,wang2021pyramid}, and GNNs inherit the brittleness of the proposal stage~\cite{xu2017scene}.
The more fundamental limit is the \emph{2D-to-3D} gap: these methods reason over 2D feature grids and image-plane relations, whereas much spatial reasoning depends on depth ordering, occlusion, viewpoint, and metric layout, and strong 2D backbones and vision--language models degrade sharply once relations hinge on 3D geometry~\cite{daxberger2025mmspatial,xu2025spatialbench}.

Two responses run through the rest of this section.
The explicit-structured methods of §\ref{subsec:p3_s7_2_2} reintroduce programs and graphs over learned features, recovering inspectability while keeping the representational power of deep networks; the geometry-aware methods of §\ref{subsec:p3_s7_2_3} answer the 2D-to-3D gap directly by building the symmetry group into the architecture, where it is known.
Whether web-scale pretraining alone closes the systematic-generalization gap, or explicit structure remains necessary, is the open question that §\ref{sec:p3_s7_3} takes up.








\subsection{Explicit Structured Neural Reasoning}
\label{subsec:p3_s7_2_2}


% IEP/EE 本质上是 program-structured / neuro-symbolic 方法，更适合放在下一节 explicit structured reasoning

Explicit structured neural reasoning decomposes a spatial problem into an inspectable intermediate structure, a program, a scene graph, a sequence of reasoning states, or executable code that calls vision tools, and obtains the answer by executing or traversing that structure over visual representations.
Unlike the rule-based reasoning of §\ref{sec:p3_s7_1}, the structure is learned or inferred from data rather than written by hand; unlike the implicit reasoners of §\ref{subsec:p3_s7_2_1}, it exposes spatial predicates directly, so that an inspectable edge or a named operation can represent containment, adjacency, support, or relative position, making them auditable and composable instead of entangled in a feature space.
Figure~\ref{fig:p3_s7_2_2_explicit_formulation} previews the representative tasks.

\begin{figure}[!th]
    \centering
    \includegraphics[width=0.99\linewidth]{figures/part3/section7/explicit_structured_neural_reasoning_task.pdf}
    \caption{Representative tasks for explicit structured neural reasoning, in which the reasoning trace is inspectable: (1) program-style VQA on compositional scenes, (2) state-machine traversal of a probabilistic scene graph, and (3) execution of an LLM-generated visual program.}
    \label{fig:p3_s7_2_2_explicit_formulation}
\end{figure}




\subsubsection{Task Definition}
Given a visual observation $\mathbf{I}$ (a single image or a set of views) and a query or goal $q$, the system constructs an intermediate structure and executes it,
\begin{equation}
  \mathcal{S}=\mathrm{Parse}_\phi(\mathbf{I},q),\qquad y=\mathcal{E}(\mathcal{S};\mathbf{I}),
\end{equation}
where $\mathcal{S}$ may be a program layout $P=(m_1,\dots,m_k)$, a recurrent sequence of reasoning states, a scene graph $G=(V,E)$ with a finite-state traversal over its nodes, or a code trace that calls external tools, and the executor $\mathcal{E}$ realizes $\mathcal{S}$ through neural modules, symbolic operators, graph transitions, or tool calls.
The output $y$ may be an answer, a localized region, a predicted relation, or an action.
The defining property is therefore not a specific program syntax but the presence of an inspectable $\mathcal{S}$ that mediates between perception and prediction, which is what makes the reasoning path explicit and controllable.




\subsubsection{Methodology}

The four families in Figure~\ref{fig:p3_s7_2_2_explicit_formulation_method} differ in how much structure they expose, enough to make reasoning inspectable and compositional, but not so rigid that a perception error, a parser error, or a missing primitive causes deterministic failure.
Each exposes a different facet of space: recurrent architectures a temporal trace of attentional focus, scene-graph methods object-centric relations, program-based methods named operations such as filtering, relating, counting, and comparing, and LLM-generated programs tool-call traces that combine detection, segmentation, localization, and external knowledge.


\paragraph{Structured Recurrent Reasoning Architectures.}
This family imposes a step-wise reasoning structure in the architecture itself, so the model is fully differentiable and trained end-to-end from input--output pairs while still exposing a step-by-step trace.
The Memory, Attention, and Composition (MAC) network~\cite{hudson2018compositional} chains $p$ recurrent cells, each of which attends to the question to choose the operation for the current step (control unit), extracts the relevant visual information by spatial attention guided by that control signal (read unit), and integrates it into a reasoning state passed to the next cell (write unit); the attention maps of successive read units provide a transparent trace of where the model is looking, without a rigid symbolic program.
MAC is thus a boundary case: explicit at the level of reasoning dynamics, implicit at the level of spatial representation, since the operations run over dense feature maps rather than named objects or relations.
Because the spatial logic stays in the learned weights, it is prone to statistical shortcuts (if ``red spheres'' co-occur with ``cubes'' in training, the model may learn that correlation rather than ``left of''), and soft attention maps are often diffuse, so ``to the left'' bleeds into irrelevant regions in cluttered scenes.
\begin{figure}[!th]
    \centering
    \includegraphics[width=0.99\linewidth]{figures/part3/section7/explicit_structured_neural_reasoning_method.pdf}
    \caption{Four method families for explicit structured neural reasoning, ordered by the intermediate structure $\mathcal{S}$ each exposes: (A) the MAC recurrent reasoning cell~\cite{hudson2018compositional}; (B) the Neural State Machine~\cite{hudson2019learning}; (C) program-based reasoning, a parse-to-execute pipeline whose executor ranges from neural modules (NMN~\cite{andreas2016neural}, IEP~\cite{johnson2017inferring}) through the symbolic executor of NS-VQA~\cite{yi2018neural} to the quasi-symbolic executor of NS-CL~\cite{mao2019neuro}; and (D) LLM-as-program-generator pipelines that call vision tools (\eg,~ViperGPT~\cite{suris2023vipergpt}, VisProg~\cite{gupta2023visual}, Code-as-Policies~\cite{liang2023code}).}
    \label{fig:p3_s7_2_2_explicit_formulation_method}
\end{figure}

\paragraph{Structured Scene Representations.}
A related line strengthens the structural prior on the reasoning substrate.
Where MAC reasons over a soft grid of features, the Neural State Machine (NSM)~\cite{hudson2019learning} reasons over a probabilistic semantic scene graph and performs differentiable state transitions over it: a concept-detection module (Mask R-CNN) identifies entities, each node $v_i$ is initialized with semantic properties over a learned concept vocabulary, each edge $e_{ij}$ carries a distribution over spatial and semantic relation types (left of, behind, on top of, holding), and the query is answered by propagating a distribution over nodes through query-conditioned transitions.
More broadly, scene graphs act as a generic structured substrate: a downstream reasoner can propagate probabilistic states over the graph (as in NSM), attend over it as a soft relational prior, or search it for graph-based question answering, in each case obtaining an inspectable cut-point between perception and reasoning.
The price is that reasoning quality becomes tightly coupled to object detection, attribute grounding, and relation prediction, a dependency we return to in the challenges.


\paragraph{Modular and Program-Based Reasoning.}
A third family translates the query itself into an executable program, and its members differ mainly in the nature of the \emph{executor}.
The Neural Module Network (NMN)~\cite{andreas2016neural} is the canonical \emph{neural-executor} instance: an external semantic parser converts a question such as ``What color is the llama?'' into a layout, \eg,~\texttt{classify[color](attend[llama])}, whose slots are instantiated by distinct learnable modules (attention, re-attention, classification, combination, measurement) trained jointly to maximize the likelihood of the gold answer.
Because NMN relies on an off-the-shelf parser, IEP~\cite{johnson2017inferring} replaces it with a sequence-to-sequence program generator that emits the program as a sequence of functions, trained jointly with the execution modules via REINFORCE~\cite{williams1992simple} or by maximizing answer likelihood over a distribution of programs.
Moving further toward symbols, Neuro-Symbolic VQA (NS-VQA)~\cite{yi2018neural} adopts a \emph{symbolic executor}: a scene parser turns the image into a symbolic table, after which the program reduces to deterministic functional operators, yielding fully interpretable execution at the cost of a near-perfect-perception assumption.
The Neuro-Symbolic Concept Learner (NS-CL)~\cite{mao2019neuro} occupies the middle ground with a \emph{quasi-symbolic} executor: a neural perception module produces object-centric representations and a semantic parser emits symbolic programs, but the executor relaxes discrete operations such as filtering and relating into differentiable ones, so concept learning and semantic parsing can be trained from question--answer pairs, while the object detector is pretrained using bounding-box annotations.

\paragraph{LLM-Generated Visual Programs.}
The most active recent extension lets a large language model emit the program directly, replacing the seq2seq parser of IEP or the hand-crafted parser of NMN with a \emph{tool-based executor}.
ViperGPT~\cite{suris2023vipergpt} and VisProg~\cite{gupta2023visual} pair a frozen LLM with a Python-style API of vision tools (detectors, captioners, attribute classifiers) that the generated code invokes, making the approach compositional and zero-shot without per-task training.
Code-as-Policies~\cite{liang2023code} pushes the same program-synthesis paradigm into embodied control, treating the LLM-emitted program as a closed-loop policy over perception and actuation primitives; it is included here as an embodied extension of the paradigm rather than as a VQA method.
The interpretability gain is real, since the emitted code is human-readable, but these systems re-inherit the brittleness of the underlying LLM (hallucinated API calls, ungrounded tool routing) as a failure mode the earlier hand-coded DSLs were designed to avoid.

\paragraph{Unifying View: Parser, Representation, Executor, and Supervision.}
The four families are points in a shared design space with four axes.
The \emph{parser} that maps a query to structure may be hand-built (NMN), learned end-to-end (IEP, NS-CL), or a frozen LLM (ViperGPT, VisProg); the \emph{representation} over which reasoning runs may be a dense feature map (MAC), an object-centric scene graph (NSM, and the scene-graph-generation literature~\cite{xu2017scene} supervised at scale by Visual Genome~\cite{krishna2017visual}), or a symbolic scene table (NS-VQA); the \emph{executor} may be neural, symbolic, quasi-symbolic, or a code interpreter; and the \emph{supervision} may be answer-only, program-supervised, or, in tool-based systems, augmented with execution traces and tool feedback.
The inspectable structure between perception and answer is the common thread: it provides a natural cut-point at which symbolic and learned components can be swapped without retraining the rest of the stack, and the same four axes reappear in the foundation-model reasoners of §\ref{sec:p3_s7_3}.

\paragraph{Beyond VQA: 3D, Temporal, and Embodied Extensions.}
The same design pattern extends beyond static images: referring-expression grounding has been lifted from 2D images to 3D scenes~\cite{achlioptas2020referit3d}, where object-centric and spatial-relation graphs supply the substrate for reasoning about support, containment, proximity, and navigability, and in video and embodied settings the intermediate structure becomes a temporal scene graph or a task plan that links perception to action (§\ref{sec:p3_s7_3}, Part~IV).
These extensions also expose the central limitation of current explicit methods: most available program primitives and graph predicates are still designed for static 2D relations, whereas viewpoint change, object permanence, physical interaction, and multi-step navigation need 3D, temporal, and causal operators.


\subsubsection{Benchmarks and Performance}

The CLEVR family~\cite{johnson2017clevr} of §\ref{subsec:p3_s7_2_1} is the canonical testbed here as well, with the added property that each question is generated alongside a ground-truth functional program, which gives both program-supervised and answer-only training signals; MAC, IEP, NS-CL, and NS-VQA all converge to the high-90s on it, making it attractive but saturated.
GQA~\cite{hudson2019gqa} is the dominant open-domain benchmark: it pairs natural-image scenes with compositional questions whose ground-truth scene graphs (from Visual Genome~\cite{krishna2017visual}) make scene-graph-based reasoners directly applicable, and its out-of-distribution splits expose the systematic-generalization gap hidden by saturated CLEVR scores.
VQA v2~\cite{goyal2017making} is a general VQA benchmark, while the cited visual-program systems ViperGPT and VisProg report task-specific comparisons on GQA.
Referring-expression grounding benchmarks (RefCOCO and RefCOCO+~\cite{yu2016modeling}, RefCOCOg~\cite{mao2016generation}, and the earlier ReferItGame~\cite{kazemzadeh2014referitgame}) are equally relevant, since they test whether a system can bind spatial language (``the object left of / behind / next to \dots'') to a localized region even when the output is a box or mask rather than an answer.
Embodied program-as-policy systems such as Code-as-Policies are evaluated by task success in robotic or simulated manipulation environments and should not be compared against VQA accuracy.
Table~\ref{tab:p3_s7_2_2_clevr_gqa} collects the headline numbers.

\begin{table}[!ht]
    \centering
    \caption{Implicit and explicit structured reasoners on CLEVR and GQA (overall accuracy, \%, higher is better). CLEVR numbers use the evaluation splits reported in the original papers; NS-CL reports validation accuracy after training on 5K images; GQA numbers are test-dev or test-standard as reported. ViperGPT and VisProg run zero-shot with a frozen LLM, and VisProg is evaluated on a curated test-dev subset, so its number is not directly comparable. RN and FiLM are the implicit reference points of §\ref{subsec:p3_s7_2_1}.}
    \label{tab:p3_s7_2_2_clevr_gqa}
    \fontsize{9}{11}\selectfont
    \setlength{\tabcolsep}{1.5mm}
    \begin{tabular}{L{0.22\textwidth} L{0.26\textwidth} L{0.20\textwidth} C{0.10\textwidth} C{0.10\textwidth}}
        \toprule
        \textbf{Method} & \textbf{Structure exposed} & \textbf{Supervision} & \textbf{CLEVR} & \textbf{GQA} \\
        \midrule
        RN~\cite{santoro2017simple}            & none (pairwise latent)              & answers                       & 95.5 & --   \\
        FiLM~\cite{perez2018film}              & none (feature modulation)           & answers                       & 97.7 & --   \\
        MAC~\cite{hudson2018compositional}     & recurrent control / read / write    & answers                       & 98.9 & 54.1 \\
        IEP~\cite{johnson2017inferring}        & generated program, neural executor  & programs (700k) $+$ answers   & 96.9 & --   \\
        NS-CL~\cite{mao2019neuro}              & program, quasi-symbolic executor    & boxes (detector) $+$ answers (curriculum) & 98.9 & --   \\
        NS-VQA~\cite{yi2018neural}             & program, symbolic executor          & scene tables $+$ 270 programs & 99.8 & --   \\
        NSM~\cite{hudson2019learning}          & probabilistic scene graph, automaton & answers                      & --   & 63.2 \\
        ViperGPT~\cite{suris2023vipergpt}      & LLM-generated code, tool executor   & none (zero-shot)              & --   & 48.1 \\
        VisProg~\cite{gupta2023visual}         & LLM-generated code, tool executor   & none (zero-shot)              & --   & 50.5 \\
        \bottomrule
    \end{tabular}
\end{table}


\subsubsection{Challenges and Trends}
Exposing programs, graphs, state transitions, or tool calls buys interpretability but creates new failure modes, and the numbers above locate them.
The first bottleneck is \textbf{parsing}: if a question is mapped to the wrong program or traversal path, even a correct executor produces the wrong answer.
The second is \textbf{symbol grounding and underspecification}: spatial predicates such as \texttt{Left}, \texttt{Behind}, \texttt{Near}, or \texttt{Between} depend on viewpoint, scale, object extent, and task context, so many spatial queries do not determine a unique program, a single deterministic DSL operator is often too rigid, and new predicates demand manual re-engineering of the executor.
The third is \textbf{error propagation}: mistakes by the detector, scene parser, relation classifier, or program generator are inherited by the executor and can cause deterministic failure, which is why NSM reaches 63\% on GQA while the same family saturates CLEVR; on natural images the bottleneck has moved from reasoning execution to perception, scene-graph construction, language grounding, and long-tail relation coverage.
A fourth, more structural tension is the \textbf{differentiability--interpretability trade-off}: the more symbolic the executor, the more auditable but the harder to train end-to-end; the more differentiable, the more robust but the less transparent.
Because the value of explicit structure lies in the intermediate trace, \textbf{evaluation} should move beyond final-answer accuracy to whether the parsed program is faithful to the query, whether graph edges correspond to real relations, whether tool calls are necessary and grounded, and whether the execution trace actually supports the answer.

Most explicit systems also reason over static 2D layouts or object graphs rather than 3D geometry, temporal dynamics, or physical causality: they can answer whether $A$ is left of $B$, but rarely simulate how a scene changes under motion, contact, containment, support, or occlusion, which is exactly the mental-simulation capacity human-level spatial reasoning requires.
LLM-generated programs widen the scope by replacing fixed parsers with flexible code generation, trading a small accuracy gap on saturated benchmarks for zero-shot generality across VQA, grounding, and knowledge-grounded VQA, but they add hallucinated tool calls, invalid APIs, ungrounded intermediate variables, and the absence of a verifiable execution loop.
The emerging direction is therefore not a return to purely symbolic reasoning but \emph{verifiable hybrid} systems, in which learned perception produces object-centric representations, structured programs or graphs provide an inspectable scaffold, and neural or tool-based executors absorb the variability of real visual inputs, the direction of §\ref{sec:p3_s7_3} and §\ref{sec:p3_s7_4}.









\subsection{Geometry-aware Neural Spatial Reasoning}
\label{subsec:p3_s7_2_3}
Where the implicit and explicit methods above learn spatial relationships from data, geometry-aware approaches build the relevant symmetries into the network itself: an object rotated, translated, or re-indexed keeps its identity, and the architecture is constructed so that this holds by design rather than being learned from many transformed examples.
The gain is sample efficiency and robustness under the chosen transformations; the cost is that the symmetry group has to be chosen in advance.
Figure~\ref{fig:p3_s7_2_3_geometry_formulation} pairs three symmetry-driven task settings with the method families that supply the matching inductive bias.

\begin{figure}[t]
    \centering
    \includegraphics[width=\linewidth]{figures/part3/section7/geometry-aware_neural_spatial_reasoning.pdf}
    \caption{Symmetry-driven tasks (top) and geometric neural mechanisms (bottom). Top: (i) rotation-invariant classification, (ii) SE(3)-equivariant grasping, shown with a planar rotation, and (iii) a molecular scalar invariant to rigid motion and atom reindexing. Bottom: (a) STN~\cite{jaderberg2015spatial}, (b) G-CNN~\cite{cohen2016group}, (c) steerable CNN~\cite{cohen2017steerable}, (d) PointNet~\cite{qi2017pointnet}, and (e) TFN / SE(3)-Transformer~\cite{thomas2018tensorfield,fuchs2020se}. Blue marks transformations, geometric features, and preserved outputs; all examples are schematic.}
    \label{fig:p3_s7_2_3_geometry_formulation}
\end{figure}

\subsubsection{Task Definition}
The input $x$ lives in a space with a known geometric structure (a 2D image, a 3D point cloud, a molecular graph), optionally with a query $q$, and the network $M_{\theta,G}$ is designed to respect a chosen group $G$ of transformations: translations for standard CNNs, rotations SO(3), rigid motions SE(3), or permutations of set elements.
Its output $\mathbf{z}$ is either invariant or equivariant to the action of $G$.

\begin{quotebox}
\noindent\textbf{Invariance vs.\ equivariance.} An \emph{invariant} representation stays fixed under the group action, $M_{\theta,G}(g\cdot x)=M_{\theta,G}(x)$, and deliberately discards pose; an \emph{equivariant} representation co-transforms with the input, $M_{\theta,G}(g\cdot x)=\rho(g)\cdot M_{\theta,G}(x)$, where $\rho(g)$ is the representation of the group action on the feature space, and thereby preserves it. In either case geometric consistency is a native constraint that the architecture guarantees by construction.
\end{quotebox}

\subsubsection{Methodology}

Invariance is the right target when the output should not depend on orientation or position (classification), equivariance when pose and relative arrangement are the point (manipulation, keypoints, instance segmentation, molecular modeling); the layers below realize each.

\paragraph{Invariance by Pooling and Augmentation.}
The baseline strategies work without changing the architecture.
Augmenting the training set with rotated, flipped, or scaled copies pushes a standard CNN toward assigning the same label to all variants, but buys only an approximate invariance and requires examples of every transformation to be tolerated; global average pooling~\cite{lin2013network} at the end of a feature extractor discards the ``where'' of each feature and keeps only its presence, which is the main reason classification CNNs are robust to translation.
Neither guarantees invariance to anything beyond translation.

\paragraph{Active Alignment with Spatial Transformer Networks.}
Rather than ignoring transformations, a Spatial Transformer Network (STN)~\cite{jaderberg2015spatial} learns to undo them: a small localization network regresses the parameters of an affine transformation that it predicts will bring the object into a canonical pose, a grid generator turns these parameters into a sampling grid, and a differentiable sampler produces the rectified feature map, so the localization network is trained by backpropagation from the task loss alone.
The module can be inserted at any point in a network and acts as learned attentional alignment, normalizing a region of interest for the layers downstream; the invariance it provides is learned, not guaranteed.

\paragraph{Permutation Invariance over Point Sets.}
A separate symmetry arises when the input is an unordered set rather than a grid: the model must be invariant to the ordering of its elements.
DeepSets~\cite{zaheer2017deep} shows that any permutation-invariant function over a set can be written as $f(X)=\phi\!\left(\sum_{x\in X}\psi(x)\right)$ for suitable $\phi$ and $\psi$, encode each element independently, sum, then decode, which turns set-level invariance from an ad-hoc choice into a constructive recipe and makes explicit which aggregators preserve invariance and which silently break it; the same framework yields permutation-\emph{equivariant} layers for per-element outputs (per-point labels, per-atom properties).
PointNet~\cite{qi2017pointnet} is the canonical instance on raw point clouds, a shared multilayer perceptron per point followed by symmetric max pooling, and PointNet++~\cite{qi2017pointnetplus} adds locality by repeatedly grouping points within a metric ball and applying the same shared-MLP-plus-pool pattern at multiple scales.
Permutation invariance comes for free in this family; rigid-motion invariance does not, which is the gap the equivariant methods below address.

\paragraph{From Translation to Group Equivariance.}
A standard convolution is translation-equivariant (shifting the image shifts the feature map) but not rotation-equivariant: rotating the input changes the feature map unpredictably.
Group equivariant convolutions (G-CNNs)~\cite{cohen2016group} generalize the operation to a larger discrete group such as the four planar rotations: a single filter is rotated into a set of transformed copies, all are correlated with the input, and the feature map gains an orientation dimension, so a feature that reads ``edge at $45^\circ$'' becomes ``edge at $135^\circ$'' at the corresponding location when the input rotates by $90^\circ$.
Equivariance is therefore not learned but imposed, layer by layer.

\paragraph{Steerable Filters and Equivariance in Every Layer.}
Steerable CNNs~\cite{cohen2017steerable} replace the group-indexed filter copies by filters expressed in a steerable basis, so that the response at a transformed orientation follows by linear combination of basis responses rather than by explicit rotation; the general E(2) framework~\cite{weiler2019general} extends this to continuous rotations and reflections with circular-harmonic bases under one configurable interface, and library efforts such as e3nn~\cite{geiger2022e3nn} do the same for E(3) / SE(3).
Full equivariance also requires that pooling, nonlinearities, and normalization respect the group: a standard ReLU applied to the components of a geometric vector breaks it (activations are therefore applied to invariant quantities such as the vector norm), and channel-wise batch normalization breaks the coupling between channels, which motivates equivariant normalization~\cite{meng2024geometric}.
The guarantee is only as strong as its weakest layer; one non-equivariant operation anywhere in the stack silently undoes the symmetry the rest preserve.

\paragraph{Equivariant Attention and Continuous SE(3) Symmetries.}
A more recent line lifts the treatment from convolutions to attention, with queries, keys, and values that transform predictably under a chosen Lie group.
The SE(3)-Transformer~\cite{fuchs2020se} is the canonical instance, building each attention layer from tensor-field representations so that rotating the input scene rotates the whole stack of outputs consistently; Tensor Field Networks~\cite{thomas2018tensorfield} use the same machinery without an attention head and have been especially influential in molecular modeling, where SE(3) equivariance is a problem constraint rather than an optional bias.
Both parameterize the kernel as a combination of spherical-harmonic basis functions, which is computationally heavy, so a parallel line pursues lighter equivariance: E($n$)-equivariant graph networks (EGNN)~\cite{satorras2021egnn} achieve E($n$)-equivariant message passing from relative distances and coordinate updates alone, and Vector Neurons~\cite{deng2021vector} lift scalar point-cloud features to 3D vectors, obtaining SO(3) equivariance as a drop-in primitive, each trading representational richness for scalability.
With the libraries above, trying equivariance on a new task now costs about as much engineering as a convolutional baseline; the choice of group still has to be made by the designer.


\subsubsection{Benchmarks}

Benchmarks in this area are constructed to expose behavior under symmetry transformations.
MNIST-rot is the standard 2D testbed, separating methods that learn rotation robustness through augmentation from those that build it in; ModelNet40~\cite{wu20153d} evaluates point-cloud classification under three protocols, aligned training and test (z/z), training with azimuthal rotations and testing under arbitrary SO(3) rotations (z/SO(3)), and arbitrary rotations in both (SO(3)/SO(3)), which is where augmentation, STN-style alignment, and PointNet's permutation-invariant pooling are separated from genuine rotation handling; ScanObjectNN~\cite{uy2019scanobjectnn} adds real scan noise and occlusion.
For equivariance as a hard constraint, QM9~\cite{ramakrishnan2014quantum} (small-molecule property prediction) and interatomic-potential tasks are the reference, where the SE(3)-equivariant models NequIP~\cite{batzner2022e3} and MACE~\cite{batatia2022mace} set the current points for data efficiency and accuracy.
Table~\ref{tab:p3_s7_2_3_symmetry} reports the two ladders on which the method families can be compared directly.

\begin{table}[!ht]
    \centering
    \caption{Symmetry benchmarks. MNIST-rot: test error (\%, lower is better); plain CNN and P4CNN figures from \citet{cohen2016group}, E(2)-steerable figure from \citet{weiler2019general}. ModelNet40: classification accuracy (\%, higher is better) under the train/test rotation protocols z/z, z/SO(3), and SO(3)/SO(3), all as re-evaluated by \citet{deng2021vector}.}
    \label{tab:p3_s7_2_3_symmetry}
    \fontsize{9}{11}\selectfont
    \setlength{\tabcolsep}{3.0mm}
    \begin{tabular}{L{0.30\textwidth} C{0.14\textwidth} C{0.11\textwidth} C{0.11\textwidth} C{0.13\textwidth}}
        \toprule
        \textbf{Method} & \textbf{MNIST-rot err.} & \textbf{z/z} & \textbf{z/SO(3)} & \textbf{SO(3)/SO(3)} \\
        \midrule
        Plain CNN (Z2CNN)~\cite{cohen2016group}                & 5.03 & --   & --   & --   \\
        P4CNN (G-CNN)~\cite{cohen2016group}                    & 2.28 & --   & --   & --   \\
        E(2)-steerable CNN~\cite{weiler2019general}            & 0.68 & --   & --   & --   \\
        PointNet~\cite{qi2017pointnet}                         & --   & 85.9 & 19.6 & 74.7 \\
        PointNet++~\cite{qi2017pointnetplus}                   & --   & 91.8 & 28.4 & 85.0 \\
        VN-PointNet~\cite{deng2021vector}                      & --   & 77.5 & 77.5 & 77.2 \\
        VN-DGCNN~\cite{deng2021vector}                         & --   & 89.5 & 89.5 & 90.2 \\
        \bottomrule
    \end{tabular}
\end{table}


\subsubsection{Challenges and Trends}
The most fundamental difficulty is choosing the symmetry group: invariance to a group that is too large (SO(3) when only SO(2) is task-relevant) destroys discriminative information, invariance to one too small leaves the model brittle to the symmetries the task actually exhibits, and the trade-off has no purely data-driven solution~\cite{cohen2016group,cohen2017steerable}.
Cost scales with the group: classical equivariant architectures index filters by group elements, so parameters and memory grow with the size of the chosen group, limiting how freely reflections, multi-scale rotations, or higher-order tensor representations can be included~\cite{weiler2019general}.
The assumption of \emph{exact} symmetry is also brittle: point clouds with self-occlusion, partial-view RGB, and noisy LiDAR scans do not realize an SO(3)-orbit, so an architecture that is mathematically equivariant can still fail when the input violates the assumed structure.
Table~\ref{tab:p3_s7_2_3_symmetry} shows where the two targets diverge: on 2D classification both hand-coded invariance and learned equivariance saturate, so data efficiency rather than accuracy becomes the discriminating axis, while on ModelNet40 the invariance-only PointNet family collapses under unseen rotations and the equivariant family holds its accuracy, because the task requires geometric information to be \emph{preserved}, not discarded.
Finally, hard equivariance sits uneasily with the attention-based representations that dominate the rest of the stack (§\ref{subsec:p3_s7_2_1}, §\ref{subsec:p3_s7_2_2}, vision--language pretraining): self-attention is permutation-equivariant but not rotation-equivariant, and bolting an equivariant layer onto an attention pipeline often degrades the inductive bias that motivated the choice.

Equivariant attention is the main compromise on offer, and frame averaging~\cite{puny2022frameaveraging} moves in the opposite direction, turning any off-the-shelf backbone into an invariant or equivariant one by symmetrizing its outputs over a small input-dependent set of group frames, which decouples the choice of symmetry from the choice of architecture and sidesteps the parameter cost of group-indexed filters.
These threads are pulling geometry-aware reasoning back toward the rest of the chapter: vision--language and vision--language--action models attach equivariant heads where 3D pose and geometric consistency matter, as in grasp policies and pose-conditioned generation, and primitives first devised for object classification are finding second lives in world models, embodied control, and 3D generation.




\subsection{Differentiable Optimization Layers for Spatial Reasoning}
\label{subsec:p3_s7_2_4}

Where the three preceding subsections put structure into the learned representation, differentiable optimization layers put it into the computation: selected feed-forward operations are replaced by the solution of a constrained optimization problem, so that hard logical, geometric, or physical constraints are enforced inside the forward pass while the whole system stays end-to-end differentiable~\cite{gould2022exploiting}.
In spatial reasoning this is how a network enforces distances, collision avoidance, or alignment, or solves a sub-problem such as pose estimation, within a learnable pipeline.
Figure~\ref{fig:p3_s7_2_4_differentiable_optimization_formulation} previews the representative tasks (top row) and the method families (bottom row).

\begin{figure}[!ht]
    \centering
    \includegraphics[width=\linewidth]{figures/part3/section7/differentiable_optimization_layers_for_spatial_reasoning.pdf}
    \caption{Differentiable optimization layers for spatial reasoning. Top, tasks: (1) QP relaxation of binary assignment for Sudoku~\cite{amos2017optnet,wang2019satnet}, (2) differentiable convex MPC~\cite{amos2018differentiable,agrawal2019differentiable}, and (3) bundle adjustment as differentiable nonlinear least squares for end-to-end SLAM~\cite{teed2021droid}. Bottom, methods: (1) OptNet~\cite{amos2017optnet}, (2) CVXPY Layers~\cite{agrawal2019differentiable,diamond2016cvxpy}, (3) Deep Declarative Networks and differentiable geometry solvers~\cite{gould2021deep,campbell2020solving}, and (4) SATNet-style differentiable SAT layers~\cite{wang2019satnet,shindo2024learning}.}
    \label{fig:p3_s7_2_4_differentiable_optimization_formulation}
\end{figure}


\subsubsection{Task Definition}
Given an input parameter vector $\theta$ produced by the preceding layers, a differentiable optimization layer computes
\begin{equation}
    z^\star = \operatorname*{arg\,min}_{z \in \mathcal{Z}} f(z; \theta) \;\; \text{s.t.}\;\; g(z;\theta) \leq 0, \; h(z;\theta) = 0,
\end{equation}
where $f$ is an objective function, $g,h$ denote inequality and equality constraints, and $\mathcal{Z}$ is a feasible domain.
The forward pass computes the optimal solution $z^\star(\theta)$, while the backward pass relies on implicit differentiation of the optimality conditions to obtain $\tfrac{\partial z^\star}{\partial \theta}$; the layer is declarative, in that the model specifies what to solve rather than how to compute it.
Under this view, QP solvers~\cite{amos2017optnet}, convex optimization layers~\cite{agrawal2019differentiable}, differentiable geometry solvers~\cite{teed2021droid}, and differentiable logic program solvers~\cite{wang2019satnet,shindo2024learning} are all instantiations of the same layer, tailored to different domains of spatial or symbolic reasoning.


\subsubsection{Methodology}

\paragraph{QP and Convex Optimization Layers.}
\citet{amos2017optnet} introduced OptNet, which treats the solution of a quadratic program as a layer: the forward pass solves
\begin{equation*}
    z^\star = \operatorname*{arg\,min}_z \tfrac{1}{2} z^{\top} Q z + q^{\top} z
\end{equation*}
subject to affine constraints whose matrices can depend on previous network outputs, and the backward pass differentiates the KKT optimality conditions exactly, with a custom GPU interior-point solver that handles many small QPs in parallel and reuses its factorization for the backward pass.
OptNet learned to solve Sudoku from data where conventional architectures could not, which showed that implicit layers can handle discrete or constrained reasoning problems.
\citet{agrawal2019differentiable} generalized the idea to any disciplined convex program: CVXPY Layers~\cite{diamond2016cvxpy} canonicalize the problem into conic form, solve it, and differentiate through the cone program's optimality conditions, so a constrained least-squares, a constrained logistic regression, or a robot-motion-limit constraint can be inserted as a layer without deriving gradients by hand.

\paragraph{Implicit Differentiation.}
All of these layers share one mechanism.
Rather than unrolling an iterative solver, one differentiates the closed-form optimality (KKT) conditions, which for a convex program comprise primal feasibility, dual feasibility, complementarity, and stationarity~\cite{DeepImplicitLayers}; differentiating this system and solving one linear system yields $\frac{\partial z^\star}{\partial \theta}$ exactly, without back-propagating through solver iterations.
Where OptNet hand-codes this backward pass, Deep Declarative Networks~\cite{gould2021deep} apply automatic differentiation to the KKT system itself.
The same principle extends to root-finding, fixed-point, and equilibrium layers via the implicit function theorem; deep equilibrium models~\cite{bai2019deep} treat an entire weight-tied network as a single fixed-point layer differentiated this way, and differentiable physics engines and planners propagate gradients through the equilibrium or optimum conditions of a physical system.

\paragraph{Differentiable Geometry Solvers.}
A subclass is tailored to geometric problems that are themselves argmin problems, typically minimizing a reprojection error: camera pose from 2D--3D correspondences, bundle adjustment, inverse kinematics.
\citet{campbell2020solving} introduced a differentiable blind PnP solver that estimates 3D pose as a layer; \citet{teed2021droid} built a differentiable bundle-adjustment module into a SLAM system, solving for camera poses and 3D points by iterative optimization while letting gradients flow; EPro-PnP~\cite{chen2022epro} recasts PnP as a probabilistic layer with a learnable distribution over poses; and Theseus~\cite{pineda2022theseus} provides a reusable library for differentiable nonlinear least squares (Gauss--Newton, Levenberg--Marquardt) that underlies bundle adjustment, pose-graph optimization, and inverse kinematics.
These layers keep the domain-specific optimizer but make it differentiable, so the whole vision or robotics system trains end-to-end while the geometry stays correct and constrained, for instance so that estimated object poses agree across views; \citet{gould2021deep} refer to such constructs as declarative nodes and demonstrate them on visual Sudoku and robust geometry estimation.

\paragraph{Differentiable Logic Layers.}
Beyond continuous optimization, the same construction extends to symbolic reasoning.
SATNet~\cite{wang2019satnet} integrates a differentiable satisfiability solver as a layer: the forward pass solves a relaxed SAT problem (a quadratic unconstrained binary optimization) to produce assignments consistent with logical constraints, and the backward pass differentiates implicitly through the solver's fixed-point conditions; \citet{shindo2024learning} extend differentiable logic programs to abstract visual reasoning benchmarks such as Raven's Progressive Matrices.
These layers belong to the same class of constraint-based implicit modules as the QP and convex layers, specialized to combinatorial and relational dependencies.


\subsubsection{Benchmarks}
Differentiable optimization layers have been validated across symbolic, geometric, and control settings.
Sudoku~\cite{amos2017optnet,wang2019satnet} and its visual variant~\cite{gould2022exploiting} test whether a layer models constraint satisfaction and propagates gradients through discrete reasoning; Raven's Progressive Matrices (RPM / I-RAVEN)~\cite{shindo2024learning} test whether differentiable logic captures relational structure in visual analogy puzzles; cart-pole, inverted-pendulum, and reaching tasks in MuJoCo and OpenAI Gym~\cite{agrawal2019differentiable} test whether differentiable MPC~\cite{amos2018differentiable} or convex layers improve control while enforcing physical feasibility; and maze navigation and obstacle avoidance test QP-based collision-avoidance layers by whether the agent reaches its target without violating safety constraints.
Evaluation combines task metrics with criteria specific to the embedded solver: geometric accuracy (translation / rotation error in pose estimation, trajectory or map error for a bundle-adjustment layer), task success rate (puzzles solved, feasible plans produced) together with the count or magnitude of violated constraints and, for physics-informed layers, residual energy or momentum drift, and solver cost and gradient stability (forward / backward runtime, iterations to converge, and sensitivity of training to approximate versus exact implicit gradients).
Table~\ref{tab:p3_s7_2_4_perf} lists the representative layers by solver type.

\begin{table}[t]
    \centering
    \small
    \caption{Representative differentiable-optimization layers, organized by solver / layer type (\cf\ Figure~\ref{fig:p3_s7_2_4_differentiable_optimization_formulation}).}
    \label{tab:p3_s7_2_4_perf}
     \fontsize{9}{11}\selectfont
    \setlength{\tabcolsep}{1.0mm}
    \begin{tabular}{L{0.26\textwidth} L{0.10\textwidth} L{0.28\textwidth} L{0.28\textwidth}}
        \toprule
        \textbf{Method} & \textbf{Year} & \textbf{Solver / layer type} & \textbf{Task / benchmark} \\
        \midrule
        OptNet~\cite{amos2017optnet}                 & 2017 & QP layer, KKT implicit diff.            & Sudoku from data            \\
        Differentiable MPC~\cite{amos2018differentiable} & 2018 & convex MPC layer                    & planning / control          \\
        CVXPY Layers~\cite{agrawal2019differentiable} & 2019 & disciplined convex / cone layer        & control, geometry           \\
        SATNet~\cite{wang2019satnet}                 & 2019 & differentiable MAXSAT relaxation       & visual Sudoku               \\
        Deep Declarative Nets~\cite{gould2021deep}   & 2021 & implicit-function-theorem nodes        & robust geometry, Sudoku     \\
        DROID-SLAM~\cite{teed2021droid}              & 2021 & differentiable bundle adjustment       & monocular / stereo SLAM     \\
        Theseus~\cite{pineda2022theseus}             & 2022 & differentiable nonlinear least squares & pose-graph / BA / IK        \\
        EPro-PnP~\cite{chen2022epro}                 & 2022 & probabilistic PnP layer                & monocular 6-DoF pose        \\
        \bottomrule
    \end{tabular}
\end{table}




\subsubsection{Challenges and Trends}
Solving an optimization problem inside a layer is more expensive than a convolution or a matrix product, and the forward pass becomes the bottleneck as problem size grows, which is exactly the regime of SLAM, bundle adjustment, and motion planning.
OptNet's batched GPU solver shows how problem structure buys tractability for many small problems, but global route planning or large-scale 3D reconstruction still overwhelms existing solvers, forcing approximate solves, warm starts, or hierarchical decomposition to keep each optimization manageable.

Gradient stability is the second difficulty.
If the solution is non-unique or depends discontinuously on the input, as at boundary optima, at active-set changes in a linear program, or in the presence of discrete variables, the implicit gradient is unstable or undefined; strict convexity guarantees a well-defined Jacobian, and real problems frequently violate it.
Entropy smoothing, small regularizers, and perturbation methods restore usable gradients in many cases, and in others the backward pass has to be designed by hand.

The third is physical consistency.
An optimization layer enforces only the constraints it is given: if the formulation omits a friction constraint in a contact model, the network exploits the gap and converges to implausible solutions, and because the layer sits inside a trainable system, the surrounding network can learn to feed it inputs that satisfy the constraints superficially while violating their intended meaning.
The benefits of interpretability and correctness therefore depend on the optimization model reflecting the real dynamics, which in practice means careful constraint modeling, additional penalties, or iterative refinement of the solver design; current work refines solvers, numerical stability, and integration with deep-learning frameworks toward layers that are both theoretically grounded and deployable in spatial systems.



\subsection{Reinforcement Learning with Spatial Inductive Biases}
\label{subsec:p3_s7_2_5}

Where the previous four subsections constrain the learned representation, reinforcement learning shifts the locus of structure to the \emph{action-conditioned} setting, in which an agent must reason about spatial relationships as a consequence of its own behavior over a long horizon.
Many spatial tasks, such as navigation, exploration, manipulation, and embodied interaction, cannot be reduced to static perception or single-step prediction; they require reasoning about how actions transform the agent's position, orientation, and relations to surrounding objects over time, under geometric constraints, partial observability, and delayed reward.
Generic RL policies built on unstructured function approximators struggle to generalize across unseen layouts, so a growing body of work introduces \emph{spatial inductive biases}, explicit planning iterations, map-like latent representations, differentiable search operators, or relational state abstractions, that constrain the hypothesis space of policies or value functions to reflect the compositional and geometric structure of spatial decision-making.


\subsubsection{Task Definition}
At each step $t$ the agent receives an observation $o_t$ from the environment $\mathcal{E}$ (RGB / RGB-D frames, egocentric maps, point clouds, language goals) together with its action history $a_{<t}$, chooses an action $a_t$ from a continuous control space or a discrete set $\mathcal{A}=\{1,\ldots,K\}$, and receives a scalar reward $r_t$.
The policy or value function is parameterized with a spatial structural prior $S$ (a differentiable planning module, a learned map or belief representation, a relational or graph encoder, an equivariant policy head),
\begin{equation}
  \pi_\theta(a_t \mid o_t, a_{<t};\, S),\qquad
  \max_\theta\ \mathbb{E}\!\left[\sum_{t=0}^{\infty} \gamma^t r_t\right],
\end{equation}
so that $S$ shapes \emph{which} policies and value functions the agent considers in the first place, not only how it generalizes after the fact; the optimization itself remains a standard RL objective (Bellman, policy gradient, actor--critic, or preference-based).

\begin{figure}[!ht]
    \centering
    \includegraphics[width=0.7\linewidth]{figures/part3/section7/rl_overview.pdf}
    \caption{The general framework in RL: the agent acts on the environment, receives an observation and a reward, and updates its policy from this feedback.}
    \label{fig:rl_overview}
\end{figure}


\subsubsection{Methodology}

Figure~\ref{fig:rl_overview} shows the loop every method below shares: the agent observes, acts, receives a reward and a new state, and updates its policy to maximize long-term return.
Learning a good policy requires balancing \textbf{exploration} of untried actions against \textbf{exploitation} of the actions currently believed best, the exploration--exploitation trade-off that underlies all of RL.
The subsections that follow walk through the optimization backbones in the order the field developed them, value-based, policy-based, actor--critic, PPO, and the recent preference-based and group-relative objectives, and note at each step where spatial structure enters.


\begin{figure}[!ht]
    \centering
    \includegraphics[width=0.70\linewidth]{figures/part3/section7/value-based-rl.pdf}
    \caption{Workflow of value-based RL.}
    \label{fig:value_based_rl}
\end{figure}

\paragraph{Value-Based RL.}
Value-based methods learn a value function that estimates the expected cumulative return of acting under a policy $\pi$.
The optimal state-action value function is
\begin{equation}
Q^*(s,a) = \max_\pi \mathbb{E}_\pi \left[ \sum_{t=0}^{\infty} \gamma^t r_t \,\middle|\, s_0 = s, a_0 = a \right],
\end{equation}
where $\gamma \in (0,1)$ is the discount factor, and it satisfies the \textit{Bellman equation}: if the optimal values $Q^*(s',a')$ of the next state are known for every action $a'$, the best action now is the one that maximizes $r+\gamma Q^*(s',a')$,
\begin{equation}
    Q^*(s, a) = \mathbb{E}_{s' \sim \mathcal{E}} [r+\gamma \max_{a'} Q^{*}(s', a') | s, a].
\end{equation}
Iterating this identity as an update, $Q_{i+1}(s, a) = \mathbb{E}_{s' \sim \mathcal{E} }[r+\gamma\; \max_{a'} \; Q_{i}(s', a') | s, a]$, converges to $Q^*$, but a table of values indexed by every state is useless in high-dimensional observation spaces, so a function approximator $Q(s,a;\theta)\approx Q^*(s,a)$ is learned instead, a neural network in the deep-RL setting (a Q-network).
Training minimizes a sequence of regression losses whose target is the bootstrapped Bellman estimate,
\begin{equation}
    L_i(\theta_i) = \mathbb{E}_{s, a \sim \rho(\cdot)} [(y_i - Q(s, a; \theta_i))^{2}],
\end{equation}
with $y_i = \mathbb{E}_{s' \sim \mathcal{E}} [r + \gamma \; \max_{a'} Q(s', a';\theta_{i-1}) | s, a]$ computed from the previous parameters $\theta_{i-1}$, held fixed, and $\rho(s,a)$ the distribution of visited state--action pairs; the greedy policy $\pi(s;\theta) = \arg\max_{a} Q(s, a;\theta)$ then follows from the learned values.

Deep Q-networks (DQN)~\cite{mnih2015human} made this scale to high-dimensional observations with experience replay and a slowly updated target network; their convergence is analyzed theoretically by \citet{fan2020theoreticalQlearning}, and Figure~\ref{fig:value_based_rl} shows the resulting workflow.
Beyond control, value-based RL drives \emph{active visual perception}: \citet{caicedo2015active} cast object localization as a sequential decision process in which a Q-network iteratively transforms an attention window until it tightly bounds the target, turning detection into a learned spatial search policy, a formulation since applied to active detection, visual search, and region selection for fine-grained recognition, where the agent learns \emph{where to look} rather than scanning exhaustively.
In spatial domains the limitation of vanilla DQN is that any planning behavior must be learned implicitly from data, which generalizes poorly to novel layouts, so several works embed planning structure in the value estimate itself.
Value Iteration Networks (VIN)~\cite{tamar2016valueiterationnetworks} unroll Bellman updates as convolutions over a spatial grid, a differentiable approximation of value iteration; Gated Path Planning Networks (GPPN)~\cite{lee2018gatedpath} replace the recurrence with gated units to improve optimization stability and long-range dependencies; and under partial observability QMDP-Net~\cite{karkus2017qmdpnet} integrates belief updates and approximate POMDP planning (QMDP)~\cite{littman1995learningQMDP} into one differentiable network.
In each case the learned value function is constrained to follow spatially structured dynamic programming rather than arbitrary function approximation.


\begin{figure}[!ht]
    \centering
    \includegraphics[width=0.7\linewidth]{figures/part3/section7/policy-based-rl.pdf}
    \caption{Workflow of policy-based RL.}
    \label{fig:policy_based_rl}
\end{figure}

\paragraph{Policy-Based RL.}
Policy-based methods~\cite{williams1992simple,zhang2021sampleefficientreinforcement} parameterize a policy $\pi_\theta(a|s)$ directly and maximize expected return,
\begin{equation}
J(\theta) = \mathbb{E}_{\pi_\theta} \left[ \sum_{t=0}^{\infty} \gamma^t r_t \right],
\end{equation}
by gradient ascent on
\begin{equation}
\nabla_\theta J(\theta) = \mathbb{E}_{\pi_\theta} \left[ \nabla_\theta \log \pi_\theta(a_t|s_t) \, R_t \right],
\end{equation}
where $R_t = \sum_{k=t}^{\infty} \gamma^{\,k-t} r_k$ is the return from time-step $t$ (Figure~\ref{fig:policy_based_rl}).
Optimizing action distributions directly suits spatial decision problems in which exploration and long-horizon credit assignment are central~\cite{chaplot2018active}, and policy gradients train embodied agents to actively gather visual evidence: Embodied Question Answering~\cite{das2018embodied} optimizes a navigation policy so that the agent moves to vantage points from which it can answer questions about a 3D scene.
The spatial competence, however, comes from the architecture rather than the objective; policy-gradient agents generalize across layouts when the policy carries a learned map~\cite{gupta2017cognitivemapping}, a grounded spatial representation~\cite{janner2018representationlearning}, memory, or a relational encoder~\cite{zambaldi2019relationaldeep}.


\begin{figure}[!ht]
    \centering
    \includegraphics[width=0.8\linewidth]{figures/part3/section7/AC-based-rl.pdf}
    \caption{Workflow of Actor-Critic-based RL.}
    \label{fig:actor_critic_based_rl}
\end{figure}

\paragraph{Actor-Critic RL.}
Actor--critic methods~\cite{degris2012modelfreerl} combine the two by jointly learning a policy $\pi_{\theta}(a|s)$ (the actor, which controls behavior) and a value function $V_\phi(s)$ (the critic, which scores the actions taken), as in Figure~\ref{fig:actor_critic_based_rl}.
The critic provides a low-variance estimate of the advantage,
\begin{equation}
A_{\phi}(s_t,a_t) \approx r_t + \gamma V_{\phi}(s_{t+1}) - V_{\phi}(s_t),
\end{equation}
or, when a Q-function is learned as well, $A_{\phi}(s_t,a_t) = Q_{\phi}(s_t,a_t) - V_{\phi}(s_t)$, and the advantage replaces the raw return $R_t$ in the policy-gradient update.
The asynchronous advantage actor-critic (A3C)~\cite{mnih2016asynchronous} made this stable and scalable and became a standard backbone for embodied navigation and spatial control: target-driven visual navigation~\cite{zhu2017target} conditions an A3C agent on an image of the goal and reaches semantically specified targets in unseen indoor scenes, and Neural SLAM~\cite{zhang2017neuralslam} gives an A3C agent a differentiable map memory.
The separation between actor and critic matches the structure of spatial tasks, with the actor responsible for geometry-aware decisions and the critic supplying global, long-horizon value signals, which is why actor--critic methods dominate navigation benchmarks that require memory, planning, and partial observability~\cite{mirowski2016learningnavigate}.



\begin{figure}[!ht]
    \centering
    \includegraphics[width=0.8\linewidth]{figures/part3/section7/ppo-based-rl.pdf}
    \caption{Workflow of PPO-based RL.}
    \label{fig:ppo_based_rl}
\end{figure}

\paragraph{PPO-Style Policy Optimization.}
Proximal Policy Optimization (PPO)~\cite{schulman2017proximalpo} is an actor--critic algorithm that stabilizes training by constraining each policy update through a clipped surrogate objective,
\begin{equation}
L^{\text{PPO}}(\theta) =
\mathbb{E} \left[
\min \left(
r_t(\theta) A_t,
\text{clip}(r_t(\theta), 1-\epsilon, 1+\epsilon) A_t
\right)
\right],
\end{equation}
where $r_t(\theta)$ is the probability ratio between the new and old policies (Figure~\ref{fig:ppo_based_rl}).
Its robustness and sample efficiency made PPO the de facto backbone for training agents in large-scale photorealistic 3D simulators such as Gibson~\cite{Xia2018Gibson} and Matterport3D~\cite{chang2017matterport3d}, and in manipulation suites such as panda-gym\footnote{\url{https://github.com/qgallouedec/panda-gym}}.
Scaled to 2.5 billion frames with DD-PPO~\cite{wijmans2019ddppo}, it essentially solves PointGoal navigation on Habitat (99.7\% success and 0.97 SPL on Gibson validation), and the same backbone underpins vision-and-language navigation~\cite{anderson2018vision}, object-goal navigation~\cite{chaplot2020object}, audio-visual navigation~\cite{chen2020soundspaces}, and planning over uncertain topological maps~\cite{beeching2020learninguncertaintopological}.

\begin{figure}[!ht]
    \centering
    \includegraphics[width=0.8\linewidth]{figures/part3/section7/dpo-based-rl.pdf}
    \caption{Workflow of DPO-based RL.}
    \label{fig:dpo_based_rl}
\end{figure}

\paragraph{Preference-Based Policy Optimization (DPO).}
Preference-based learning replaces the reward signal with comparisons.
Direct Preference Optimization (DPO)~\cite{rafailov2023directpo} learns a policy from preference pairs without an explicit reward model or online rollouts (Figure~\ref{fig:dpo_based_rl}): for two completions $(y_1,y_2)$ of an input $x$, a labeler marks the preferred one, $y_w \succ y_l \mid x$, and by expressing the implied reward of a Bradley--Terry preference model~\cite{bradley1952rankbt} through the optimal policy itself, DPO collapses the reward-learning-then-RL pipeline into a single maximum-likelihood objective on $\pi_\theta$:
\begin{equation}
    \mathcal{L}_{\text{DPO}}(\pi_{\theta};\pi_{\text{ref}}) = - \mathbb{E}_{(x,y_w, y_l)\sim \mathcal{D}} \big[ \log \sigma \big(\beta\log\frac{\pi_{\theta}(y_w|x)}{\pi_{\text{ref}}(y_w|x)} - \beta\log\frac{\pi_{\theta}(y_l|x)}{\pi_{\text{ref}}(y_l|x)}\big)\big],
\end{equation}
where $\pi_{\text{ref}}$ is a fixed reference policy and $\beta$ controls the deviation from it.
Developed for language-model alignment, DPO applies to spatial reasoning whenever supervision takes the form of trajectory preferences, plan rankings, or human judgments of spatial behavior~\cite{liang2025structuredpreference,huang2024making,qiu2025spatialPreference}, which arise naturally in navigation, manipulation, and interactive tasks where dense rewards are hard to specify.



\begin{figure}[!ht]
    \centering
    \includegraphics[width=0.8\linewidth]{figures/part3/section7/grpo-based-rl.pdf}
    \caption{Workflow of GRPO-based RL.}
    \label{fig:grpo_based_rl}
\end{figure}
\paragraph{Group-Based Policy Optimization (GRPO).}
Group Relative Policy Optimization (GRPO)~\cite{guo2025deepseekr1} removes the learned critic and replaces it with a group-relative baseline computed from multiple sampled trajectories (Figure~\ref{fig:grpo_based_rl}).
For each question $q$, it samples a group of outputs $\{o_1, o_2, \cdots, o_G\}$ from the old policy $\pi_{\theta_{\text{old}}}$ and optimizes $\pi_\theta$ by maximizing
\begin{equation}
\begin{split}
     \max_{\pi_\theta} &\;
    \mathbb{E}_{q \sim D_{\text{rl}}, \{o_i\}_{i=1}^{G} \sim \pi_{\theta_{\text{old}}}(O|q)}  \\
    & \Bigg[ \frac{1}{G}
        \sum_{i=1}^G
        \text{min}(\frac{\pi_\theta(o_i)}{\pi_{\theta_{\text{old}}}(o_i)} A_i, \text{clip}(\frac{\pi_\theta(o_i)}{\pi_{\theta_{\text{old}}}(o_i)}, 1-\epsilon, 1+\epsilon) A_i )
        - \beta \, \mathbb{D}_{\mathrm{KL}}\big(\pi_\theta \| \pi_{\mathrm{ref}}\big)
    \Bigg],
\end{split}
\end{equation}
where $\epsilon$ and $\beta$ are hyperparameters and the advantage $A_i$ is computed from the rewards $\mathbf{r}=\{r_1,\dots,r_G\}$ that a reward model assigns to the group, normalized within the group only:
\begin{equation}
    A_i = \frac{r_i - \text{mean}(\mathbf{r})}{\text{std}(\mathbf{r})}.
\end{equation}
Eliminating the value network reduces memory and stabilizes training.
Although explored first for language models, the formulation transfers to spatial decision-making wherever several candidate plans or trajectories can be scored against one another, and it has become the standard recipe for reinforcement fine-tuning of spatial vision--language models on visual reasoning, navigation, multi-step assembly, and video spatial reasoning~\cite{tan2025reasonrft,gao2026octonav,tang2025legoPuzzles,ouyang2025spacer}.


\paragraph{Architectural Spatial Inductive Biases.}
The backbones above are domain-agnostic; what makes a policy \emph{spatial} is the prior $S$ in its architecture, and three families recur across the examples already met.
\textbf{Map-based representations} give the agent an explicit allocentric memory: cognitive maps~\cite{gupta2017cognitivemapping} and Neural Map~\cite{parisotto2018neural} write observations into a spatially indexed latent grid, while Neural SLAM~\cite{zhang2017neuralslam} and Active Neural SLAM~\cite{chaplot2020learning} couple differentiable mapping with a hierarchical exploration policy.
\textbf{Differentiable planning modules} embed dynamic programming in the network, as in the VIN / GPPN / QMDP-Net line discussed under value-based RL, or learn a value-equivalent model and plan by lookahead search, as in MuZero~\cite{schrittwieser2020mastering}.
\textbf{Relational abstractions} replace flat feature maps with object- or graph-structured state, letting attention or message passing capture relations that transfer across layouts~\cite{zambaldi2019relationaldeep}.
In each case the RL objective is unchanged; the spatial competence comes from the prior, which is why these architectures generalize across geometries where unstructured policies overfit.


\subsubsection{Benchmarks}

Two complementary benchmark families shape the empirical landscape.
Embodied navigation and exploration benchmarks stress long-horizon planning under partial observability in photorealistic 3D scenes: Habitat~\cite{savva2019habitat} is the dominant platform, paired with real-world scans from Matterport3D~\cite{chang2017matterport3d} and Gibson~\cite{Xia2018Gibson} for large, scene-diverse PointGoal and ObjectGoal navigation that exercises geometry-aware exploration and cross-scene generalization; SoundSpaces~\cite{chen2020soundspaces} adds audio-visual cues as a multimodal test of spatial state estimation; and SOON~\cite{zhu2021soon} extends vision-and-language navigation to scenario-oriented object goals in unseen environments.
Symbolic and relational benchmarks test whether structured priors yield combinatorial and compositional generalization: PDDLGym\footnote{\url{https://github.com/tomsilver/pddlgym}} builds Gym environments from PDDL domains with relational observations and actions, and lightweight gridworld suites such as MiniGrid\footnote{\url{https://github.com/Farama-Foundation/Minigrid}} support controlled evaluation of planning, memory, and compositional generalization; the differentiable planners (VIN, GPPN, QMDP-Net) are evaluated on procedurally generated gridworld mazes and POMDPs of increasing size, where the question is whether the architectural prior generalizes to unseen layouts rather than memorizing training mazes.
Across both families the choice of RL algorithm matters less than the choice of prior: PPO has converged as the standard optimizer for embodied navigation, and differences between agents mostly reflect the map, planning module, or relational encoder on the policy, while on gridworld benchmarks the methods that embed planning structure consistently dominate unstructured DQN, with the largest gains where the task is most planning-shaped.


\subsubsection{Challenges and Trends}

Sample inefficiency comes first.
Spatial tasks involve long horizons, sparse rewards, and combinatorial state spaces, so na\"ive exploration is prohibitively expensive; architectural priors such as learned maps, planning modules, and relational representations improve efficiency~\cite{tamar2016valueiterationnetworks,zhang2021sampleefficientreinforcement}, but current RL systems still need orders of magnitude more interaction than humans or classical planners to acquire robust spatial competence.

Generalization and transfer come second.
Many agents perform well in their training environments and degrade sharply in unseen layouts, novel geometries, or altered dynamics, which suggests that much of the learned behavior rests on superficial correlations rather than abstract spatial structure; structured state representations, world models, and modular policy design~\cite{chaplot2018active,gupta2017cognitivemapping} help, but how to enforce spatial invariances and compositional generalization in RL in a principled way remains open.

Reward specification is the third.
Spatial tasks often lack dense, well-defined rewards, especially when success depends on long-term plan quality, safety, or efficiency rather than immediate outcomes.
Preference-based methods such as DPO~\cite{rafailov2023directpo} and group-relative objectives such as GRPO~\cite{guo2025deepseekr1} learn from comparative feedback instead, at the price of reward-model bias, preference inconsistency, and the difficulty of aligning learned preferences with the true task objective in complex spatial domains.

The directions that follow from these three are hybrid methods that integrate learning with explicit planning or search~\cite{tamar2016valueiterationnetworks,karkus2017qmdpnet}, structure-aware representations (object-centric, relational, graph-based) that scale spatial reasoning beyond grid-based and navigation-centric benchmarks~\cite{zambaldi2019relationaldeep}, and evaluation protocols that weight out-of-distribution generalization, robustness to environmental variation, and reasoning quality over performance on fixed benchmarks.




\subsection{Multi-agent Spatial Reasoning}
\label{subsec:p3_s7_2_6}


Where the previous five subsections study a single learning agent, multi-agent spatial reasoning (MASR) moves to the \emph{joint} regime: multiple embodied or abstract agents perceive, communicate, and act within a shared environment to accomplish cooperative or competitive goals, and the structural prior now governs decentralized decisions over a shared spatial state.
MASR links low-level sensorimotor streams (proprioception, occupancy grids, trajectories) to relational decision processes (joint planning, role assignment, coalition formation); its core problem is \emph{credit assignment and coordination} under spatial constraints, that is, how agents with partial, decentralized observations reason over nearby entities, infer latent intent, and synthesize actions that produce coherent team-level behavior.
Deep learning has moved the field from rule-driven coordination to end-to-end policies that integrate perception, communication, and control under partial observability and non-stationarity~\cite{lowe2017maddpg,foerster2018coma,rashid2018qmix}, with autonomous driving and interactive motion forecasting over road geometry (WOMD~\cite{ettinger2021womd}, nuScenes~\cite{caesar2020nuscenes}), multi-robot exploration, formation control, and patrol under attrition~\cite{goeckner2024magec}, and real-time strategy and team-sports simulation~\cite{samvelyan2019smac,kurach2020grf} as its representative domains.
We stress at the outset that the canonical algorithms surveyed below (VDN, QMIX, COMA, MADDPG) are themselves spatially agnostic: they solve a generic cooperative Dec-POMDP and see geometry only through the environment and the reward.
Explicit spatial structure enters chiefly through the benchmarks and through the graph-based family, where the interaction graph $\mathcal{G}_t$ makes neighborhood relations a first-class part of the policy.

\begin{figure}[t]
    \centering
    \includegraphics[width=0.95\linewidth]{figures/part3/section7/multi-agent_spatial_reasoning.pdf}
    \caption{Multi-agent spatial reasoning tasks (top) and mechanisms (bottom). Top: (a) cooperative micromanagement under partial observation (SMAC~\cite{samvelyan2019smac}), (b) interactive motion forecasting (WOMD, nuScenes~\cite{ettinger2021womd,caesar2020nuscenes}), and (c) multi-robot patrol after agent loss~\cite{goeckner2024magec}. Bottom: (d) centralized critic (MADDPG, COMA~\cite{lowe2017maddpg,foerster2018coma}), (e) value decomposition (VDN, QMIX~\cite{sunehag2018vdn,rashid2018qmix}), (f) graph coordination, and (g) attention across agents and time. Graph and sequence architectures serve prediction or control and do not by themselves imply decentralized execution. Red marks the lost agent; dashed trajectories are candidate futures and dashed graph links are intermittent interactions. Examples are schematic.}
    \label{fig:p3_s7_2_6_masr_formulation}
\end{figure}


\subsubsection{Task Definition}
We cast cooperative multi-agent spatial reasoning as a graph-structured Dec-POMDP
\begin{equation}
\mathcal{M}=(\mathcal{V},\{\mathcal{A}_i\},\mathcal{S},P,R,\{\mathcal{O}_i\},O,\rho_0,\gamma,T),
\end{equation}
with a (possibly time-varying) interaction graph $\mathcal{G}_t=(\mathcal{V},\mathcal{E}_t)$ derived from state (\eg, proximity or visibility) or estimated locally.
At time $t$, agent $i$ receives a partial egocentric observation $o_t^i\!\sim\!O(\cdot\mid s_t,i)$ (RGB-D, LiDAR, occupancy, kinematics) and optional neighbor messages $m_t^i$; we treat $\tilde{o}_t^i=(o_t^i,m_t^i,\mathcal{G}_t^i)$ as the input available to agent $i$, with any bandwidth cost included in $R$.
Agents act with decentralized, history-dependent policies
\begin{equation}
a_t^i \sim \pi_\theta^i(\cdot \mid h_t^i),\quad
h_t^i=\mathrm{upd}(h_{t-1}^i,\tilde{o}_t^i),
\end{equation}
the joint policy factorizes as $\pi_\theta(\mathbf a_t\mid \mathbf h_t)=\prod_i \pi_\theta^i(a_t^i\mid h_t^i)$, and the objective is
\begin{equation}
\max_{\theta}\ \mathbb{E}_{s_1\sim\rho_0,\ \pi_\theta,\ P}
\Big[\sum_{t=1}^{T}\gamma^{t-1}R(s_t,\mathbf a_t)\Big],\qquad \mathbf a_t=(a_t^1,\dots,a_t^n).
\end{equation}
In predictive variants, auxiliary heads output multi-agent forecasts or spatial fields with confidences for planning or losses; these are \emph{not} environment actions.
Settings differ along two further axes: online methods act in streaming settings under tight latency, whereas offline methods learn from logged interaction data (\eg, driving datasets) with batch updates and conservative behavior constraints; and centralized training with decentralized execution (CTDE) exploits global state to train joint value functions or critics while enforcing decentralized policies at test time~\cite{foerster2018coma,rashid2018qmix}, whereas fully decentralized learning never accesses global state, trading sample efficiency for deployability.


\subsubsection{Methodology}
Figure~\ref{fig:p3_s7_2_6_masr_formulation} summarizes the tasks above and the method families discussed below.

\paragraph{Classical Paradigms.}
Before learning, coordination relied on potential fields, formation keeping via graph Laplacians, and optimal multi-agent pathfinding with conflict-based search~\cite{Sharon2015CBS}, treated in §\ref{subsec:p3_s7_1_3}; these methods offer strong guarantees but struggle with perception noise, combinatorial scaling, and model mismatch in open worlds.

\paragraph{CTDE and Value Decomposition.}
Two lines made CTDE work.
Centralized critics resolve credit assignment and non-stationarity while keeping execution decentralized: COMA uses a counterfactual baseline computed by a centralized critic for decentralized actors~\cite{foerster2018coma}, and MADDPG extends deterministic policy gradients with agent-specific centralized critics~\cite{lowe2017maddpg}.
Value factorization instead decomposes the joint value: VDN sums per-agent utilities to approximate $Q_{\text{tot}}$ for cooperative tasks~\cite{sunehag2018vdn}, and QMIX generalizes the sum to a monotonic mixing network that preserves individual-global argmax consistency with more representational capacity~\cite{rashid2018qmix}.
These approaches made MARL scalable on SMAC~\cite{samvelyan2019smac}, where the value-decomposition line dominated the early cooperative leaderboard while the centralized-critic methods anchored other regimes (MADDPG on the particle environments, COMA on StarCraft micromanagement); the monotonic constraint and homogeneous-role assumptions limit expressivity, as the challenges below detail.

\paragraph{Graph Neural Networks for Coordination.}
Graph-based MARL encodes agents as nodes and interactions as edges, and message passing or attention over the neighborhood learns role-specialized coordination~\cite{jiang2018gcrl}; this family improves sample efficiency and robustness under topology changes and communication limits, and MAGEC~\cite{goeckner2024magec} shows resilience to agent attrition and communication disturbance in multi-robot patrol by combining GNN encoders with PPO-style multi-agent optimization.
Graph MARL targets a different axis from the leaderboard methods, robustness to topology disturbances rather than headline return, which exposes a gap between coordination-quality benchmarks (SMAC, MPE) and deployment-quality ones (multi-robot patrol, driving forecasting on WOMD and nuScenes) that standardized suites still under-represent.

\paragraph{Sequence Models and Transformers.}
Casting MARL as sequence modeling applies attention over multi-entity, multi-horizon contexts: the Multi-Agent Transformer~\cite{wen2022mat} maps the sequence of agents' observations to a sequence of actions with an encoder--decoder, using the multi-agent advantage decomposition to turn joint policy search into sequential decisions with linear complexity in the number of agents.
Transformer controllers, often hybridized with graph layers, address long-range temporal and relational dependencies, unify perception and control, and support offline RL through trajectory modeling; learned memory alleviates partial observability, while quadratic attention and credit assignment over long horizons remain open.


\subsubsection{Benchmarks}
Benchmarks span two regimes: simulation suites that stress cooperation, partial observability, and communication learning under controlled physics and reward shaping, and real-world multi-agent trajectory datasets that test interactive prediction and motion forecasting at deployment scale.
Table~\ref{tab:MASR_datasets} summarizes the most widely used resources.

\begin{table}[t]
\centering
\caption{Common MASR benchmarks spanning simulated coordination to real-world multi-agent trajectory datasets.}
\label{tab:MASR_datasets}
\fontsize{9}{11}\selectfont
\setlength{\tabcolsep}{3mm}
\begin{tabular}{@{}lccc@{}}
\toprule
\textbf{Dataset / Suite} & \textbf{Year} & \textbf{Scale} & \textbf{Action/Task Space} \\
\midrule
SMAC \cite{samvelyan2019smac} & 2019 &
  22+ micro scenarios & Discrete, decentralized \\
Multi-Agent Particle Env.\ (MPE) \cite{lowe2017maddpg} & 2017 &
  10+ 2D templates & Continuous/discrete \\
GRF \cite{kurach2020grf} & 2020 &
  Dozens of tasks / full-game & Continuous/discrete  \\
nuScenes \cite{caesar2020nuscenes} & 2020 &
  1,000 scenes $\times$ \SI{20}{s} @ \SI{2}{Hz} &
  Trajectory prediction \\
WOMD \cite{ettinger2021womd} & 2021 &
  $10^{5}$ scenes $\times$ \SI{20}{s} @ \SI{10}{Hz} &
  Trajectory prediction \\
Neural MMO \cite{mmo20aamas, MMO23nips} & 2020 &
  Persistent worlds; $O(10^2)$ agents & Long-horizon survival \\
\bottomrule
\end{tabular}
\end{table}

Evaluation depends on whether the setting is control or prediction.
Control is scored by mean episodic return $\mathbb{E}_{\pi}\!\left[\sum_t r_t\right]$, success or win rate, time-to-goal, collision and rule-violation rates, fairness or load-balancing indices for cooperative logistics, and communication overhead; prediction and planning are scored by spatial accuracy, from displacement errors over the forecast horizon (minADE, minFDE, miss rate on WOMD and nuScenes) to intersection-over-union between predicted and ground-truth occupied regions or trajectory tubes.


\subsubsection{Challenges and Trends}

Multi-agent spatial reasoning faces structural difficulties that distinguish it from the single-agent settings of §\ref{subsec:p3_s7_2_1}--§\ref{subsec:p3_s7_2_5}, organized here along four axes within the classical multi-agent paradigm.

\paragraph{Credit Assignment under Partial Observability and Long Horizons.}
The decentralized-execution constraint of the Dec-POMDP forces each agent to act on its own history $h^i_t$ alone, so relating a local action $a^i_t$ to a delayed team-level reward $r_T$ across a horizon $T$ produces an irreducibly noisy policy gradient.
Counterfactual baselines such as COMA \cite{foerster2018coma} marginalize over alternative actions of agent $i$ while holding the rest of $\mathbf{a}_t$ fixed and reduce gradient variance substantially, but they still rely on a centralized critic that, at scale, becomes a moving target for itself.
Value decomposition (VDN, QMIX) \cite{sunehag2018vdn,rashid2018qmix} side-steps the centralized critic by factorizing $Q_\mathrm{tot}$ across agents, yet the monotonicity constraint $\partial Q_\mathrm{tot}/\partial Q_i \ge 0$ that guarantees decentralized \emph{argmax} consistency cannot represent strictly non-monotonic cooperation games, a representational ceiling made explicit by the later non-monotonic factorizations QTRAN, weighted QMIX, and QPLEX, and visible on cooperative matrix games and SMAC's hardest scenarios~\cite{son2019qtran,rashid2020wqmix,wang2021qplex,samvelyan2019smac}.
Under sparse reward shaping, policies trained against these signals tend to chase short-term spatial cues (``move toward the closest enemy'', ``follow the nearest teammate'') rather than team-level objectives.

\paragraph{Scaling Laws of Joint Reasoning and Communication.}
Sequence-model multi-agent Transformers and dense graph message passing both incur $\mathcal{O}(n^2)$ complexity in the number of agents, plus additional quadratic terms across temporal contexts, which collides with the real-time latency budgets of multi-robot \cite{goeckner2024magec} and autonomous-driving deployments \cite{ettinger2021womd,caesar2020nuscenes}.
Sparse, learned interaction graphs are the natural mitigation, but they introduce their own credit-assignment difficulty: the edges $\mathcal{E}_t$ are themselves part of the learned policy, and gradients flowing through hard edge selection are biased or high-variance.
Communication bandwidth is a related gap: in the formulation above any message $m^i_t$ should carry a cost that appears in $R$, but the standard simulators (SMAC \cite{samvelyan2019smac}, MPE \cite{lowe2017maddpg}, GRF \cite{kurach2020grf}) idealize communication as free and noiseless, leaving little benchmark surface against which to push bandwidth-aware coordination.

\paragraph{Non-stationarity, Heterogeneous Roles, and Partner Generalization.}
Even in fully cooperative settings, intra-team policies co-evolve, so from the perspective of any single agent the environment is non-stationary \cite{lowe2017maddpg}, the failure mode the centralized-critic family of MADDPG / COMA \cite{lowe2017maddpg,foerster2018coma} was introduced to address.
The problem sharpens when teammates are heterogeneous, may drop out (the patrol-with-attrition setting of MAGEC \cite{goeckner2024magec}), or are drawn from an unseen partner distribution at evaluation time, the zero-shot coordination problem.
Persistent open-world simulators such as Neural MMO \cite{mmo20aamas,MMO23nips} push along this axis with procedurally generated maps and long-horizon survival, but standardized protocols for partner generalization and role specialization remain immature compared with the single-agent supervised benchmarks of §\ref{subsec:p3_s7_2_1} and §\ref{subsec:p3_s7_2_3}.

\paragraph{Ambiguous Spatial Semantics and Sim-to-Real Transfer.}
The egocentric observations $o^i_t$ that drive decentralized policies are partial and noisy: occluded teammates, uncertain affordances, and depth-discontinuity-induced map errors distort both prediction and control.
Driving benchmarks such as WOMD \cite{ettinger2021womd} and nuScenes \cite{caesar2020nuscenes} expose the prediction side of this gap, but the closed-loop control side (whether the same forecasts survive when an agent acts on its own predictions) is comparatively under-evaluated.
The sim-to-real gap on multi-robot systems is wider still: perception noise and actuation latency interact non-linearly with the coordination signal in ways that simulators rarely model, and resilience under realistic communication disturbance and agent attrition has only recently received first-class benchmark coverage \cite{goeckner2024magec}.

Three routes follow from these axes.
Topology-aware sparse interaction (learned edges, hierarchical groupings, or attention masks parameterized by $\mathcal{G}_t$) preserves long-range coordination while scaling sub-quadratically, with the GNN+PPO patrol architecture of MAGEC \cite{goeckner2024magec} as an early instantiation and graph encoders combined with sequence-model attention over short temporal windows as the natural next step.
The monotonicity ceiling of VDN / QMIX \cite{sunehag2018vdn,rashid2018qmix} can be relaxed by attention-weighted mixers, weighted projection schemes, or dueling-style decompositions that keep individual-global \emph{argmax} consistency while admitting non-monotonic cooperation, so that COMA-style counterfactual credit assignment \cite{foerster2018coma} no longer depends on a joint-action critic that becomes the bottleneck.
And explicit objectives for agent attrition and communication disturbance \cite{goeckner2024magec}, paired with failure-centric curricula and hooks for formal verification, move resilience from an emergent property of training to a first-class objective, decisive in safety-critical multi-robot and driving deployments \cite{ettinger2021womd,caesar2020nuscenes} where worst-case performance dominates the deployment decision; closed-loop evaluation suites that jointly report team return, safety, calibration, and latency are the missing measurement substrate, and the benchmarks for partner generalization, bandwidth-aware coordination, and sim-to-real transfer that would play the role SMAC and GRF played for CTDE are still under construction.





% \subsection{Summary}
% \label{subsec:p3_s7_2_summary}

% Read from top to bottom, the six paradigms move the structural prior from the representation to behavior: implicit reasoning leaves it latent, explicit and geometry-aware reasoning lifts it into a program or a symmetry group, optimization layers place it in the forward pass, and the two reinforcement-learning subsections let it shape the policy of one agent and then of a team.
% Table~\ref{tab:p3_s7_2_paradigm_comparison} summarizes the six along the axes used for the symbolic paradigms in §\ref{sec:p3_s7_1}.

% \begin{table}[t]
%     \centering
%     \fontsize{8}{10}\selectfont
%     \setlength{\tabcolsep}{2pt}
%     \caption{Comparison of the six paradigms of learning-based spatial reasoning covered in §\ref{sec:p3_s7_2}, on the same axes as the comparison table of §\ref{sec:p3_s7_1}; the \emph{locus} column gives where each paradigm injects its structural prior.}
%     \label{tab:p3_s7_2_paradigm_comparison}
%     \begin{tabularx}{\textwidth}{@{}p{2.9cm} p{3.3cm} p{3.3cm} X X@{}}
%         \toprule
%         \textbf{Paradigm} & \textbf{Structural prior (\emph{locus})} & \textbf{Canonical inputs / outputs} & \textbf{Strengths} & \textbf{Limitations} \\
%         \midrule
%         §\ref{subsec:p3_s7_2_1} Implicit neural reasoning      & Latent prior: attention, relation networks, FiLM modulate the learned feature space                & Pixel / feature inputs $\to$ answer, caption, or relation label                                  & End-to-end fluency on perception-heavy benchmarks; minimal modeling burden                & Limited inspectability; compositional generalization gaps; shortcut learning under sparse supervision \\
%         §\ref{subsec:p3_s7_2_2} Explicit structured reasoning  & Program- or graph-level prior: parsed layouts, scene graphs, neural-module / LLM-generated programs & Image $+$ compositional question $\to$ executable program $\to$ stepwise inspectable answer       & Compositional generalization; auditability; modular reuse of perception tools              & Parser / program-generator quality is a hard ceiling; brittle on long-tail vocabulary; supervision-hungry \\
%         §\ref{subsec:p3_s7_2_3} Geometry-aware reasoning       & Symmetry-group prior: invariance / equivariance to a chosen group of input transformations         & Image / point cloud / molecule $\to$ classification, pose, or property prediction                & Sample efficiency; out-of-distribution rotation / translation robustness                    & Group choice is a silent error mode; classical equivariant filters scale with $\lvert G \rvert$; brittle when inputs violate the assumed symmetry \\
%         §\ref{subsec:p3_s7_2_4} Differentiable optim.\ layers  & Optimization prior: a constrained convex / SAT problem is solved inside the forward pass            & Image / pose / observation $\to$ optimization-layer call $\to$ optimal solution $+$ KKT gradients & Certified feasibility / structure on the prediction; end-to-end training preserves geometry & Solver wall-clock dominates the forward pass; biased / high-variance gradients near constraint boundaries \\
%         §\ref{subsec:p3_s7_2_5} RL with spatial priors          & Action-conditioned prior: differentiable planning, learned maps, relational policies               & Observation $+$ action history $\to$ action under MDP / POMDP; reward signal                      & Long-horizon credit assignment; generalization across spatial layouts; embodied tasks       & Sample inefficiency vs.\ supervised paradigms; reward-shaping fragility; sim-to-real gap \\
%         §\ref{subsec:p3_s7_2_6} Multi-agent spatial reasoning   & Joint action-conditioned prior: Dec-POMDP, CTDE, graph-MARL, sequence-model coordination           & Per-agent egocentric obs.\ $+$ messages $+$ interaction graph $\to$ joint decentralized actions   & Team-level coordination; topology-aware robustness; resilience under attrition              & Credit assignment under partial observability $+$ non-stationarity; $\mathcal{O}(n^2)$ scaling; partner generalization immature \\
%         \bottomrule
%     \end{tabularx}
% \end{table}

% Every paradigm in the table trades the certificates of §\ref{sec:p3_s7_1} (feasibility proofs, force-closure verdicts, $do(\cdot)$-effect estimates) for end-to-end trainability, and each keeps a slice of the classical structure as its prior; which slice it keeps determines both its strengths and its blind spots.
% As the locus of that prior migrates from representation to behavior, the evaluation regime migrates with it, from offline benchmark accuracy at the top of the table to closed-loop team return at the bottom.
% What no single row supplies is \emph{cross-locus composition}: a geometry-aware backbone (§\ref{subsec:p3_s7_2_3}), an inspectable program-level reasoner (§\ref{subsec:p3_s7_2_2}), a differentiable feasibility layer (§\ref{subsec:p3_s7_2_4}), and an action-conditioned policy (§\ref{subsec:p3_s7_2_5}--§\ref{subsec:p3_s7_2_6}) stacked into one pipeline without losing the guarantee each contributes.
% That composition problem is the bridge to §\ref{sec:p3_s7_3}, where the prior is repositioned inside a large pre-trained model, and to §\ref{sec:p3_s7_4}, where it is distributed across a hybrid of classical and learned components.

% \newpage


\section{Foundation Model-driven Spatial Reasoning}
\label{sec:p3_s7_3}

\FramedInitial{F}{oundation} model-driven spatial reasoning turns to large pre-trained models (LLMs, vision-language models, vision-language-action models, and large vision models) as the engine of spatial inference. The motivation is empirical: such models display a striking facility with spatial concepts and relations even without task-specific fine-tuning, and a large body of work now recruits them as the core of systems for spatial question answering, planning, and explanation.

What separates these approaches is the locus and granularity at which the foundation model intervenes. At one end the model is the reasoning core itself, answering spatial queries directly (LLM/VLM-centric reasoning) or additionally emitting low-level actions (VLA-based reasoning). Its reasoning can be drawn out step by step, as intermediate deductions in language (chain-of-thought) or as an executable program synthesized from the task (code and DSL generation). It can be coupled to external machinery such as planners, renderers, simulators, and perception modules (tool-augmented reasoning), or grounded in external spatial memory (retrieval- and memory-augmented reasoning). And it can be handed control of a closed interaction loop with the environment (agentic reasoning) or paired with a learned neural simulator of dynamics that serves as the spatial substrate (world-model-based reasoning).


\subsection{LLM/VLM-centric Spatial Reasoning}
\label{subsec:p3_s7_3_1}
    % - Focus: the llm/vlm model for spatial reasoning
% [202511] Deconstructing Spatial Intelligence in Vision-Language Models
% [ICML2025] Why Is Spatial Reasoning Hard for VLMs? An Attention Mechanism Perspective on Focus Areas

Large language and vision-language models have become a default substrate for spatial reasoning, not because they compute geometry explicitly, but because large-scale pretraining leaves them with a usable, if imperfect, prior over how objects are arranged, oriented, and related. This subsection examines that prior directly: what spatial competence text-only LLMs acquire from language alone, how visual grounding sharpens it in VLMs, and where both still fall short of the metric, viewpoint-consistent reasoning that physical tasks demand. The methods are organized by the intervention used to elicit or strengthen this competence, from prompting to architectural change.

\begin{figure}[t]
    \centering
    \includegraphics[width=\linewidth]{figures/part3/section7/LLM_VLM-centric_spatial_reasoning.pdf}
    \caption{LLM/VLM-centric spatial reasoning, illustrated by the query ``Which object is left of the lamp?'': (i) text-only or visual inputs with optional spatial evidence, (ii) a pretrained LLM/VLM with four complementary, not sequential, intervention strategies (prompt, training, data, model), and (iii) the answer, other task outputs, and evaluation categories. Shading marks the queried region, not a measured attention map.}
    \label{fig:p3_s7_3_1_llm_vlm_formulation}
\end{figure}

\subsubsection{Task Definition}
Spatial reasoning, which embodied AI~\cite{driess2023palm, zhu2024spa}, autonomous driving~\cite{bogdoll2025muvo,zhou2024hktsg}, and spatial content generation~\cite{yu2025wonderworld, liu2025agentic} all depend on, is cast here as inference over learned representations rather than explicit geometric computation: a pre-trained model $\mathcal{M}_\theta$ answers a spatial query $q$ from context $c$,
\begin{equation}
  y=\mathcal{M}_\theta(q, c),\qquad c\in\{\varnothing,\ I,\ I+\text{aux}\},
\end{equation}
where the context may be absent (a text-only LLM), an image or image set $I$ (a VLM), or an image together with auxiliary evidence such as a scene graph, region features, or depth, and the output $y$ is a relation, a metric estimate, an orientation, or a navigation plan.
For text-only LLMs the spatial prior is built from statistical patterns in massive corpora~\cite{yamada2023evaluating,li2024advancing,wang2024picture}: such models infer that ``sky'' is above ``ground'', but the reasoning is \emph{ungrounded}, with no mapping to physical reality or visual signal, and therefore prone to physically impossible inferences.
VLMs ground the same prior by aligning the LLM with a visual encoder, and spatial competence then extends beyond recognizing objects and their intrinsic attributes to reasoning about relations among objects and to \emph{spatial extrapolation}, inferring unseen states or running mental simulations from what is perceived~\cite{liu2025deconstructing,chen2024spatialvlm}.




\subsubsection{Methodology}
% Training-free Prompting
% Model-Centric Architecture
% Explicit Spatial Injection
% Data-Centric Approaches

Interventions against the spatial deficiencies of foundation models fall into four levels, in increasing order of invasiveness: prompt engineering leaves the model untouched, training strategies update its weights, data-centric augmentation changes what it learns from, and model-centric enhancement alters the architecture.

\paragraph{Prompt Engineering.}
Prompting elicits latent spatial ability at inference time without weight updates.
Textual prompting applies chain-of-thought reasoning~\cite{liao2024reasoning} so that the model decomposes a spatial task into intermediate steps (identifying reference objects before deducing relative positions), and inserting scene graphs or coordinates extracted by expert models into the prompt raises accuracy further~\cite{mitra2024compositional}.
Visual prompting guides attention through markers superimposed on the image itself (masks, coordinate axes, keypoints), and in zero-shot settings such explicit visual cues can outperform fine-tuned models~\cite{yang2023set,liu20253daxisprompt}.

\paragraph{Training Strategies.}
Contrastive pretraining objectives induce a ``bag-of-concepts'' behavior that under-supervises spatial structure, so stronger supervisory signals are introduced to enforce spatial alignment.
Auxiliary losses such as image reconstruction or coordinate prediction compel the model to attend to pixel-level spatial detail rather than global semantics alone~\cite{wang2024reconstructive}.
Reinforcement learning with spatial reward functions (PPO or GRPO~\cite{wu2025reinforcing,shen2025fine}; view-consistency rewards~\cite{wang2025svqa}) trains models to perform mental visual simulation or to produce more precise spatial descriptions, moving spatial training from imitation toward preference-based optimization.

\paragraph{Data-Centric Augmentation.}
Because high-quality spatially annotated data is scarce, augmentation drives much of the progress.
Existing 2D/3D datasets are re-annotated with explicit spatial relations, metric distances, or orientations~\cite{wang2024all}, and simulation engines synthesize data with perfect geometric ground truth, which is especially effective for extrapolation tasks such as mental rotation and path planning where real-world annotation is impractical~\cite{tang2024sparkle,jiang2022pseudo}.

\paragraph{Model-Centric Enhancement.}
Architectural changes target the fine-grained visual-spatial features that mainstream encoders miss.
Because CLIP-style encoders~\cite{radford2021learning} prioritize global semantics over local geometry, higher-resolution encoders or fused features from multiple visual experts (SAM~\cite{kirillov2023sam, ravi2024sam2, carion2025sam}, DINOv2~\cite{dinov2}) sharpen the perception of shape and position; region-of-interest features or scene graphs supplied as additional input tokens decouple objects from their relationships and tighten fine-grained alignment~\cite{guo2024regiongpt}; and for 3D reasoning, models ingest point clouds, voxel grids, or depth maps directly~\cite{wangspatial} or lift 2D features into implicit 3D representations from multi-view images~\cite{zhu2024llava}.
A recent alternative internalizes 3D awareness through auxiliary supervision instead of extra inputs: SpatialSV~\cite{tang2026spatialsv} forces hidden visual features to predict explicit 3D targets (depth, ray maps, point clouds) during training, injecting geometric priors at no inference cost.



\subsubsection{Benchmarks}
% [EgoOrientBench]: Is 'right' right? enhancing object orientation understanding in multimodal large language models through egocentric instruction tuning

% [GeoMeter]: Understanding depth and height perception in large visual-language models
% [CV-Bench]: Cambrian-1: A fully open, vision-centric exploration of multimodal llms
% [What's Up]: What's" up" with vision-language models? investigating their struggle with spatial reasoning
% [SEED-Bench]: Seed-bench: Benchmarking multimodal llms with generative comprehension
% [SRBench]: Mind the gap: Benchmarking spatial reasoning in vision-language models
% [MINDCUBE]: Spatial mental modeling from limited views
% [RealWorldQA]: Grok-1.5v vision preview
% [OmniSpatial]: Omnispatial: Towards comprehensive spatial reasoning benchmark for vision language models

% \cite{jung2025right,azad2025understanding,tong2024cambrian,KamathHC23a,li2023seed,stogiannidis2025mind,yin2025spatial,jia2025omnispatial}

Evaluation has moved from general VQA to benchmarks that probe spatial capability along four levels~\cite{liu2025deconstructing}.
Intrinsic-attribute benchmarks test perceptual precision on single objects, with EgoOrientBench~\cite{jung2025right} evaluating object orientation from an egocentric viewpoint and GeoMeter~\cite{azad2025understanding} metric distance and dimension, and they show that even state-of-the-art models are weak at reference-free perception.
Relational benchmarks, the most populated class, test the parsing of spatial relations through vision-centric questions such as relative position (CV-Bench~\cite{tong2024cambrian}, What's Up~\cite{KamathHC23a}), and this is where general-purpose VLMs peak.
Deep-spatial-logic benchmarks go beyond visible relations: SRBench~\cite{stogiannidis2025mind} evaluates mental rotation and navigation, and MindCube~\cite{yin2025spatial} the construction of an internal 3D mental model from multiple views.
Comprehensive suites such as RealWorldQA~\cite{xai2024grok15v} and OmniSpatial~\cite{jia2025omnispatial} cover all levels at once; a wider coverage map of these static probes appears in Table~\ref{tab:spatial_reasoning_benchmark} of §\ref{subsec:p3_s7_3_7}.



\subsubsection{Challenges and Trends}
Three obstacles limit what foundation models can be made to do spatially.
Mainstream vision encoders such as CLIP optimize for global semantic alignment and capture coarse semantics at the expense of fine-grained geometric detail~\cite{wang2023makes,tong2024eyes}, and the contrastive objective compounds this by inducing the ``bag-of-concepts'' behavior that under-supervises structural relations among objects~\cite{zhou2025logic,wang2025spatialclip}.
Training data with precise metric annotation (quantitative distance, 3D coordinates) is scarce, and existing datasets rarely support the extrapolation tasks that require reasoning beyond direct observation~\cite{liu2023visual,agrawal2023stupd}.
Spatial reference is ambiguous: ``to the left'' may mean the viewer's left (egocentric), the object's left (allocentric), or another agent's left, and VLMs seldom switch frame of reference from context, which misaligns interpretation in embodied use~\cite{liu2023visual,lee2025perspective}.

Two directions respond to these.
Reinforcement learning with process rewards such as GRPO moves training from static VQA toward interaction, letting the model acquire spatial logic by trial and error in simulation or the real world~\cite{wu2025reinforcing}.
World-model-informed pretraining attacks the 2D-to-3D gap directly, learning physical regularities (gravity, collision) alongside image-text correspondence from synthetic corpora with exact geometric ground truth~\cite{yang2026mindjourney,tang2024sparkle,chen2024spatialvlm}; the use of such world models at inference time is the subject of §\ref{subsec:p3_s7_3_8}.




\subsection{VLA-based Spatial Reasoning}
\label{subsec:p3_s7_3_2}
    % - Vision-language-action model
% [Arxiv202509] Pure Vision-Language-Action (VLA) Models A Comprehensive Survey
% [Arxiv202505] Vision-Language-Action Models Concepts, Progress, Applications and Challenges

Vision-language-action (VLA) models push spatial reasoning past description and into execution: the same latent competence that lets a model answer where an object sits must now drive an end effector toward it without violating depth, contact, or stability constraints. Spatial reasoning here is therefore inseparable from the agent's body and the physics of its workspace, and a plan that reads correctly on paper still fails if it ignores reachability or collision. The methods below are organized by their action-generation paradigm (autoregressive, diffusion-based, reinforcement-learning-driven, and hybrid).

\begin{figure}[t]
    \centering
    \includegraphics[width=\linewidth]{figures/part3/section7/VLA_spatial_reasoning.pdf}
    \caption{VLA-based spatial reasoning as a perception--action loop: (i) embodied inputs, (ii) a VLA policy that generates actions autoregressively, by diffusion (denoising $A_K$ to $A_0$), or with a hybrid, plus optional RL refinement, and (iii) an action representation executed by a controller, with environment feedback closing the loop. $\Delta p,\Delta R, g, q^{*}$ are end-effector translation and rotation increments, gripper command, and joint targets. The robot scenes illustrate the initial state and intended outcome, not experimental observations.}
    \label{fig:p3_s7_3_2_vla_formulation}
\end{figure}

\subsubsection{Task Definition}

A VLA model extends the reasoning core of §\ref{subsec:p3_s7_3_1} with an action head: given a visual observation $o_t$, a language instruction $\ell$, and the robot's proprioceptive state $s_t$, it emits an executable action chunk
\begin{equation}
  a_{t:t+H}\sim\pi_\theta(\cdot\mid o_t,\ \ell,\ s_t),
\end{equation}
where each action encodes end-effector increments, gripper commands, joint targets, or pose waypoints that a controller converts into motion~\cite{zhang2025pure,sapkota2025vision}.
Spatial reasoning under this paradigm is \emph{embodied grounding}: language tokens must be anchored not only to visual features but to the three-dimensional geometry, affordances, and constraints of the workspace, so that the generated actions respect depth, stability, and contact rather than merely naming where objects are.



% 模型在具身环境中，跨越视觉感知和自然语言指令，理解、定位并操纵物理实体的能力 。

\subsubsection{Methodology}
VLA methods for spatial reasoning fall into four classes by their action-generation and reasoning paradigm.

% 自回归
\paragraph{Autoregression-Based Methods.}
Autoregressive methods cast spatial manipulation as next-token prediction over a unified token stream of discrete action tokens.
Gato~\cite{reed2022generalist} showed that heterogeneous modalities can be jointly tokenized and trained in one autoregressive model; PointVLA~\cite{li2025pointvla} injects point-cloud representations into a pretrained VLA for lightweight 3D reasoning without redesigning the architecture; and TraceVLA~\cite{zheng2024tracevla} adds visual trace prompting, explicit spatiotemporal cues that improve long-horizon spatial awareness and task consistency during generation.

% 扩散
\paragraph{Diffusion-Based Methods.}
Diffusion methods model action generation as conditional denoising, which suits trajectories that must respect physical constraints and uncertainty.
SE(3)-DiffusionFields~\cite{urain2022se} runs the diffusion process in $\mathrm{SE}(3)$ pose space and jointly optimizes grasping and motion planning; 3D Diffuser Actor~\cite{ke20243d} conditions the denoiser on structured 3D scene representations so that manipulation trajectories stay geometrically consistent; and ForceVLA~\cite{yu2025forcevla} brings force and tactile feedback in through a mixture-of-experts design for contact-rich manipulation.

% 强化学习
\paragraph{Reinforcement Learning-Driven Methods.}
Reinforcement learning refines VLA policies through interaction with the environment, particularly on long-horizon tasks (the broader treatment of RL with spatial inductive biases is §\ref{subsec:p3_s7_2_5}).
VIP~\cite{ma2022vip} learns a goal-conditioned, action-independent value function that implicitly encodes spatial distance to assess task progress; NaVILA~\cite{cheng2024navila} couples VLA representations with RL policies so that quadruped robots navigate and avoid dynamic obstacles on difficult terrain; and SafeVLA~\cite{zhang2025safevla} adds safety critics that constrain high-risk behavior during manipulation.

% 混合(自回归+扩散)

\paragraph{Hybrid Architectures.}
Hybrids combine generative mechanisms and add explicit spatial representations to strengthen 3D reasoning.
3D-VLA~\cite{zhen20243d} joins autoregressive language modeling with diffusion-based action prediction inside a generative 3D world model; VoxPoser~\cite{huang2023voxposer} has an LLM compose 3D voxel value maps, separating goal interpretation from geometric planning; ReKep~\cite{huang2024rekep} models spatiotemporal dependencies as relational keypoint constraints for precise manipulation; and RoboPoint~\cite{yuan2024robopoint} predicts affordance maps that encode where actions are feasible, as a prior for downstream planning and control.


\subsubsection{Benchmarks}
Evaluation rests on two complementary resources: generalization-oriented benchmarks in simulation and large-scale real-world embodied datasets.
VIMA-Bench~\cite{jiang2023vima} measures spatial generalization through a four-level protocol that runs from random object placement to fully novel task compositions; ALFRED~\cite{shridhar2020alfred} poses long-horizon household tasks that combine high-level goals with low-level instructions; and CALVIN~\cite{mees2022calvin} and LOTA-Bench~\cite{choi2024lota} test language-conditioned manipulation across environments, scoring generalization to unseen objects and configurations.
Among datasets, Open X-Embodiment~\cite{o2024open} is the most comprehensive, integrating 22 robot datasets from 21 institutions and 527 skills to enable cross-platform generalization; BridgeData~\cite{ebert2021bridge} spans 71 tasks in 10 environments and shows that joint training on limited out-of-domain data raises success rates substantially; and RH20T~\cite{fang2024rh20t} adds vision, force, and audio across 147 tasks, suited to one-shot learning and experience-based policy transfer.
Across these suites performance falls sharply along every generalization axis (L1 to L4 in VIMA-Bench, seen to unseen splits in ALFRED, single to multi-step chains in CALVIN), while data scale and embodiment diversity are the most reliable lever for cross-platform transfer and multimodal sensing (force, audio) becomes decisive where vision alone underdetermines the contact configuration.



\subsubsection{Challenges and Trends}
The generalization drop above reflects a \textbf{semantic--physical misalignment}: VLA systems inherit their competence from LLMs and VLMs that excel at semantic abstraction and cross-modal alignment but have had little exposure to physical interaction, so instructions are often interpreted correctly and executed incorrectly.
The same gap appears as \textbf{pseudo-interaction}: actions are generated from statistical co-occurrence in the training data rather than from an understanding of environment dynamics or causality, so when the environment is unfamiliar or changes quickly the model cannot revise its strategy from sensory feedback.

Two further limits are practical.
\textbf{Long-horizon consistency} is unsolved for autoregressive and diffusion policies alike: over extended action sequences the spatial representation drifts or hallucinates, and existing embodied datasets, including Open X-Embodiment, are dominated by tabletop manipulation and grasping, so environmental diversity and large unstructured scenes are under-represented.
\textbf{Inference latency} compounds across autoregressive decoding of large Transformers, and for high-frequency tasks such as dynamic grasping or mobile navigation the compute and memory demands remain prohibitive on embedded platforms.

The directions that follow are a unified token stream over vision, language, and action that lets VLA systems grow into proto world models (the world-model view is developed in §\ref{subsec:p3_s7_3_8}); active perception and causal reasoning, in which the agent probes its environment and updates its policy from feedback, as the answer to pseudo-interaction; and virtual-real data synthesis (high-fidelity simulation, synthetic generation, multi-robot data sharing) as the route to embodied data at the scale that text corpora provided for language models.


\subsection{Chain-of-Thought for Spatial Reasoning}
\label{subsec:p3_s7_3_3}
\label{ssec:cot-reasoning}

Spatial reasoning poses unique challenges for large language models (LLMs), as it requires the manipulation of geometric relations, perspective-taking, and dynamic transformations: skills that are not purely symbolic but grounded in embodied cognition.
Chain-of-Thought (CoT) prompting~\cite{wei2022chain}, originally developed as a mechanism to elicit intermediate reasoning steps in arithmetic and commonsense tasks, has therefore been imported into spatial reasoning not as a generic recipe for longer answers, but as a way to externalize the intermediate spatial commitments that a model must keep consistent: which object is left of which, which frame of reference is active, which action changes the scene, and which hypothesized 3D configuration explains the observation.
The evidence is mixed.
In some settings, explicit rationales help decompose a spatial query into auditable subgoals; in others, they induce a fluent but ungrounded verbal trace that makes spatial performance worse~\cite{kancheti2026cotdegrades}.

\begin{figure}[t]
    \centering
    \includegraphics[width=\linewidth]{figures/part3/section7/chain-of-thought_spatial_reasoning.pdf}
    \caption{Spatial reasoning tasks (top) and intermediate-trace structures (bottom). Top: (a) textual deduction in a north-up frame, (b) visual counting, where the boxed red cube is the only one left of the green sphere, and (c) embodied planning revised by new observations. Bottom: CoT, ToT, and GoT differ in trace topology; PoT and CoA ground the trace in program execution and environment interaction. $\tau=(t_1,\ldots,t_L)$ is the intermediate trace and $\hat y=g(c,q,\tau)$ the final answer. Scenes and trace diagrams are illustrative, not recorded model computation.}
    \label{fig:p3_s7_3_3_cot_formulation}
\end{figure}


\subsubsection{Task Definition}
The task is to answer a spatial query $q$ about a scene or environment, given context $c$ that may be purely textual, visual, 3D, temporal, or embodied.
A CoT-style reasoner augments the answer with an intermediate trace
\begin{equation}
    \tau = (t_1,\ldots,t_L), \qquad \hat{y}=g(c,q,\tau),
\end{equation}
where each $t_\ell$ records a partial spatial commitment: a relation, coordinate estimate, subgoal, viewpoint change, object binding, or predicted action effect.
The trace is useful only insofar as it improves the grounded answer $\hat{y}$ or makes the failure mode inspectable; many apparent ``reasoning steps'' are not grounded in metric or perceptual evidence and merely rationalize a final guess, so the central question is not whether a model can produce more text, but whether the trace helps preserve spatial invariants across transformations, viewpoints, time steps, and actions.

\subsubsection{Methodology}
CoT, Tree-of-Thought (ToT)~\cite{yao2023tot}, and Graph-of-Thought (GoT)~\cite{besta2024got} differ mainly in the topology of the intermediate trace: CoT is a single path through partial spatial commitments; ToT branches over alternative hypotheses or subplans and uses a value function to prune; GoT allows shared sub-derivations, such as reusing a computed room layout or object graph across several queried relations.
All three search over partial reasoning states with a proposal mechanism, an evaluator $V$, and a halting rule $H$,
\begin{equation}
    \tau^\star=\arg\max_{\tau} V(\tau)\quad\text{s.t.}\quad t_{\ell+1}\in\mathrm{Propose}(t_{\le \ell}),\ \ H(\tau)=1,
\end{equation}
and differ in whether the search structure is a path, a tree of alternatives, or a graph of reusable sub-results, and in whether the trace stays verbal or becomes executable through programs (Program-of-Thought, PoT~\cite{chen2022pot}) or environment-changing actions (Chain-of-Action, CoA~\cite{pan2024coa}).
The use cases below organize how these traces are deployed.

\paragraph{Text-Only Spatial Deduction.}
In linguistic spatial puzzles, route descriptions, and relative-position QA, the model must maintain a latent arrangement without visual feedback.
CoT can expose intermediate bindings (\eg, ``A is north of B; C is west of A'') and reduce one-shot relational confusion~\cite{hu2023chain}, while ToT-style branching can preserve competing layouts until enough constraints have been applied.
The persistent failure mode is coordinate drift: even with explicit steps, LLMs lose track of the origin, orientation, or accumulated displacement in textual mazes and grid-navigation tasks~\cite{zhang2025mitigating}, so text-only CoT is most reliable when the trace is tied to a compact symbolic artifact (a table, coordinate sketch, or adjacency list) rather than to unconstrained prose.

\paragraph{Multimodal Scene Reasoning.}
In image- or video-grounded QA, a CoT trace must bind language to perceptual evidence: which object instance is referred to, what relation is visible, which is occluded, and which viewpoint assumption is made.
Free-form rationales can verbalize a relation that the image does not support, so recent work moves toward grounded intermediate artifacts (structured layouts, object bindings, diagrams, or executable module calls): SCoT-style planners use layout-coordinate traces to guide diffusion models on the generation side~\cite{chen2026spatialcotgeneration}, and program-mediated systems call perception modules rather than describe them~\cite{hsu2025makes}, the paradigm of §\ref{subsec:p3_s7_3_4}.
Here ToT and GoT serve less as generic search algorithms than as hypothesis-management devices that let a model compare object correspondences, camera frames, or depth-order explanations before committing to an answer.

\paragraph{Embodied Planning and VLA Reasoning.}
In navigation, manipulation, and vision-language-action tasks, spatial reasoning unfolds over time, and the trace must connect perceptual state, subgoal choice, action feasibility, and feedback from the environment.
CoT-style subgoals support long-horizon planning in navigation~\cite{lin2025navcot,liu2025navr1}, embodied instruction following~\cite{lu2023thinkbot,mu2023embodiedgpt}, and VLA models that expose intermediate visual or action commitments: CoT-VLA pairs reasoning with predicted visual states~\cite{zhao2025cotvla}, ACoT-VLA formulates the trace in action space through coarse action intents and reference trajectories~\cite{zhong2026acotvla}, and scene graphs stabilize the trace by giving the model a persistent substrate for object identities and relations~\cite{zhang2025embodiedvsr}.

\subsubsection{Benchmarks}
Evaluation mirrors the three use cases.
Spatial puzzles, relative-position QA, and textual maze navigation are evaluated with prompt-based CoT on the suites of~\cite{hu2023chain,zhang2025mitigating}; visual question answering and scene-grounded compositional queries are evaluated under CoT with program execution (§\ref{subsec:p3_s7_3_4}), where the rationale's correctness is mediated by program output rather than by the rationale string; and long-horizon planning with sensory feedback is evaluated through CoT-augmented navigation (NavCoT~\cite{lin2025navcot}, NavR1~\cite{liu2025navr1}), CoT-augmented instruction following (ThinkBot~\cite{lu2023thinkbot}, EmbodiedGPT~\cite{mu2023embodiedgpt}), visual- and action-trace supervision for VLA models~\cite{zhao2025cotvla,zhong2026acotvla}, and scene-graph-augmented embodied reasoning~\cite{zhang2025embodiedvsr}.


\subsubsection{Challenges and Trends}
The transfer of CoT to spatial intelligence is uneven.
\textbf{Coordinate hallucination} persists under simple textual CoT: even with ToT-style beam expansion, prompt-based planners mistrack absolute positions across multi-step chains in maze and grid-navigation problems~\cite{zhang2025mitigating}.
\textbf{Ungrounded supervision} can make things worse: textual CoT applied without a structured or rationale-grounded training signal degrades spatial VLM performance, because the model is forced to verbalize a reasoning trace it does not possess~\cite{ji2025enhancing,kancheti2026cotdegrades}, especially when the prompt rewards plausible explanation length over geometric correctness.
\textbf{Search cost} scales super-linearly in solution depth for ToT and GoT (and in the branching factor for ToT), which limits real-time or on-device deployment.
\textbf{Trace--action mismatch} arises in embodied settings when a VLA model produces a coherent verbal rationale while its action head follows a different latent policy; action-aware CoT aligns rationales with action tokens and visual predictions~\cite{zhong2026acotvla}, but reliable attribution between the visible trace and the executed behavior remains unresolved.

\emph{Grounded rationale supervision} pairs CoT with explicit spatial signals (scene graphs~\cite{zhang2025embodiedvsr}, depth maps, 3D coordinate frames, action traces, or layout-coordinate traces for spatial generation~\cite{chen2026spatialcotgeneration}) so the rationale is verifiable rather than free-form.
A complementary line makes the trace \emph{execution-grounded}: PoT-style traces compile and run against a geometric engine or vision--language tool, with execution outputs serving as a truth signal, the natural bridge from CoT to code / DSL generation in §\ref{subsec:p3_s7_3_4}; for embodied agents the parallel bridge runs from CoT to CoA and the agentic orchestration loops of §\ref{subsec:p3_s7_3_7}.


\subsection{Code/DSL Generation for Spatial Reasoning}
\label{subsec:p3_s7_3_4}
Code/DSL generation treats spatial reasoning as program synthesis: an LLM or MLLM translates a natural-language task into an executable program, and a more deterministic system (a robot control API, a physics simulator, a library of vision functions) executes it.
Relative to the explicit structured reasoners of §\ref{subsec:p3_s7_2_2}, the specialized parser and fixed module set are replaced by the world knowledge and reasoning ability of the LLM, so the operation set is no longer confined to a predefined vocabulary; relative to the chain-of-thought traces of §\ref{subsec:p3_s7_3_3}, the trace is a program whose execution against an external environment supplies a truth signal that a free-form rationale cannot.
Figure~\ref{fig:p3_s7_3_4_code_formulation} shows the two operating modes.

\begin{figure}[t]
    \centering
    \includegraphics[width=\linewidth]{figures/part3/section7/code_DSL_spatial_reasoning.pdf}
    \caption{Code/DSL generation for spatial reasoning in two modes: understanding (top), where a counting query becomes a program over perception calls whose call trace is inspectable, and planning (bottom), where an instruction becomes robot API calls whose execution feedback drives program revision. Function names are illustrative pseudocode, not a specific system's API, and the scenes do not show experimental execution.}
    \label{fig:p3_s7_3_4_code_formulation}
\end{figure}


\subsubsection{Task Definition}
Given a high-level, often ambiguous, natural-language query $q$ and access to an API $\mathcal{A}$ (a robot's control interface, a programming-language interpreter, or a set of vision-language functions), the LLM acts as a high-level planner that emits a program and offloads execution to the external system $\mathcal{E}$,
\begin{equation}
  P=\mathrm{LLM}(q,\mathcal{A}),\qquad y=\mathrm{Exec}(P;\mathcal{E}),
\end{equation}
where $P$ is written in a general-purpose language such as Python or in a domain-specific language and composed of calls that manipulate objects, query spatial relations, or control an agent, and the output $y$ is a change in the environment's state (a robot moving an object), an answer to the query, or both.

\subsubsection{Methodology}

\paragraph{Code/DSL for Understanding.}
For static scenes the LLM generates code that dissects a visual query by composing calls to perception modules, so the program is itself an explicit, verifiable reasoning chain.
VisProg~\cite{gupta2023visual} prompts an LLM to produce a program in a high-level DSL that sequentially invokes modules for image understanding, image manipulation, and knowledge retrieval, implemented either heuristically in Python or by neural models such as OWL-ViT~\cite{minderer2022simple} and Stable Diffusion~\cite{rombach2022ldm}.
ViperGPT~\cite{suris2023vipergpt} instead exposes two classes, \texttt{ImagePatch} and \texttt{VideoSegment}, whose methods (\texttt{find}, \texttt{exists}, \texttt{compute\_depth}, \texttt{verify\_property}) each call a pretrained model, and generates a Python script that chains these calls to answer a complex query about an image or video, composing existing vision tools zero-shot.
Both depend on a static, human-defined DSL; VADAR~\cite{marsili2025visual} removes that limit by synthesizing the API implementations dynamically, in an API-generation stage followed by program synthesis, and LEFT~\cite{hsu2023s} makes the approach learnable by having the LLM generate differentiable logic programs that chain trainable perception networks, so concept grounding can be learned where no pretrained module exists.
The same recipe lifts to 3D: LLM-Grounder~\cite{yang2024llmgrounder} decomposes a grounding query into Python code that calls 3D grounding tools such as OpenScene~\cite{peng2023openscene} or LERF~\cite{kerr2023lerf} to locate objects in a scene.

\paragraph{Code/DSL for Planning.}
Beyond static understanding, the generated code can serve as a policy that acts in an environment, with the LLM as high-level planner and the interpreter handling deterministic low-level execution.
Code as Policies~\cite{liang2023code} has the LLM write Python that calls a robot's control API, so that ``put the apple in the bowl'' becomes \textit{robot.pick('apple')} followed by \textit{robot.place('bowl')}, with implicit steps and preconditions supplied by the model's common sense.
In simulation, Voyager~\cite{wang2023voyager} explores Minecraft by writing, executing, and refining its own code: it prompts the LLM for code to achieve a goal, receives environment feedback (success, failure, new inventory items), debugs accordingly, and adds successful routines to a growing library of reusable skills.


\subsubsection{Benchmarks}
Code/DSL generation for spatial reasoning is typically evaluated under three benchmark families that mirror the two-mode split of the Methodology above. For 2D scene understanding, methods are evaluated on compositional visual question answering (GQA~\cite{hudson2019gqa}), referring expression comprehension (RefCOCO / RefCOCO+ / RefCOCOg~\cite{yu2016modeling,mao2016generation,kazemzadeh2014referitgame}), and image-editing benchmarks. For 3D scene understanding, evaluation has converged on 3D referring expression and 3D visual grounding benchmarks (\eg,\ ScanRefer~\cite{chen2020scanrefer}, Multi3DRefer~\cite{zhang2023multi3drefer}, ScanQA~\cite{azuma2022scanqa}), which exercise both the perception-API composition and the long-tail of object-attribute reasoning. For embodied planning, the canonical evaluation suites are ALFWorld~\cite{shridhar2020alfworld} and VirtualHome~\cite{puig2018virtualhome} (long-horizon household tasks) and Minecraft (open-world goal acquisition, as exemplified by Voyager~\cite{wang2023voyager}); the headline metric in this family is task-success rate over multi-step goals rather than single-step accuracy. Across all three families, the recurring methodological observation is that code-generation methods inherit the accuracy of their underlying perception or control modules, so benchmark performance tracks the module suite more tightly than it tracks the LLM backbone.

\subsubsection{Challenges and Trends}
The interpretable trace comes with three recurring limits.
A \textbf{fixed, human-authored DSL} (ViperGPT's \texttt{find} / \texttt{compute\_depth} / \texttt{verify\_property}) restricts the model to a predefined operation set, so spatial concepts outside the API are systematically out of reach; VADAR's dynamic API synthesis~\cite{marsili2025visual} is a response but does not solve the coverage problem.
\textbf{Brittle execution with silent failure} follows from the dependence on perceptual modules (OWL-ViT~\cite{minderer2022simple}, OpenScene~\cite{peng2023openscene}, LERF~\cite{kerr2023lerf}): an upstream error propagates through the program and yields a plausible but wrong answer with no diagnostic signal.
\textbf{Sample-inefficient debug loops} limit self-refinement agents such as Voyager~\cite{wang2023voyager} and Code as Policies~\cite{liang2023code}, whose LLM calls per task scale poorly in high-dimensional or stochastic environments.

Two directions follow.
PoT-style execution-grounded rationales are being folded back into chain-of-thought training, with the generated program serving as both rationale and verification signal, and \emph{learnable DSLs} such as LEFT~\cite{hsu2023s} replace fixed perception with differentiable logic programs whose components are trainable, pointing toward a DSL that adapts under task supervision rather than being hand-engineered.


\subsection{Tool-augmented Reasoning}
\label{subsec:p3_s7_3_5}
\label{ssec:tool-augmented}
Tool-augmented reasoning in spatial intelligence is the paradigm in which an LLM or VLM acts not as an isolated solver but as an orchestrator over a catalogue of external operators.
These operators may be planners, renderers, CAD systems, physics simulators, retrieval engines, perception modules, or metric estimators.
The section boundary is deliberately operational: tools are treated as black boxes with input--output contracts, and the central research problem is the orchestration policy that decides when to call them, how to parameterize the call, and how much to trust the returned result.
When the executable program itself becomes the reasoning trace, the detailed treatment belongs to code / DSL generation (§\ref{subsec:p3_s7_3_4}); when persistent external memory becomes the learned object, it belongs to retrieval- and memory-augmented reasoning (§\ref{subsec:p3_s7_3_6}); when the loop becomes closed-loop control, it becomes agentic spatial reasoning (§\ref{subsec:p3_s7_3_7}); and when the tool is a learned predictive simulator, the specialized discussion is world-model-based reasoning (§\ref{subsec:p3_s7_3_8}).

\begin{figure}[t]
    \centering
    \includegraphics[width=\linewidth]{figures/part3/section7/tool_augmented_spatial_reasoning.pdf}
    \caption{Tool-augmented spatial reasoning as iterative orchestration over black-box tools: (i) scene, instruction, and tool catalogue $\mathcal{T}$, (ii) an LLM/VLM that selects tool $i_t$ with arguments $q_t$, receives result $r_t$, and appends to history $h_t$ until STOP, and (iii) the answer or plan with its tool-call trace. Execution, representation, validation, and memory are tool roles, not sequential stages. Solid and dashed boxes mark the selected cube and target compartment; the scene is a generated illustration with annotated, not detected, overlays.}
    \label{fig:p3_s7_3_5_tool_formulation}
\end{figure}

\subsubsection{Task Definition}
A tool-augmented reasoning system is formally defined as a tuple $(\mathcal{M}_\theta, \mathcal{T}, \pi_\phi)$, where $\mathcal{M}_\theta$ denotes the reasoning core, $\mathcal{T} = \{T_1, \dots, T_k\}$ is a catalogue of external tools, and each tool $T_i : \mathcal{I}_i \rightarrow \mathcal{O}_i$ is exposed only through a structured input--output interface.
The orchestration policy $\pi_\phi$ governs tool selection, argument construction, response parsing, retry / fallback behavior, and feedback integration.
Given an input $x \in \mathcal{X}$, the system maintains a time-ordered reasoning history $h_t$ (initially $h_0 = \varnothing$) and proceeds by the iterative update
\begin{equation}
\begin{aligned}
s_t &= \mathcal{M}_\theta(x, h_{t-1}), &
(i_t, q_t) &= \pi_\phi(s_t), \\
r_t &= T_{i_t}(q_t), &
h_t &= h_{t-1} \mathbin{\Vert} (s_t, i_t, q_t, r_t),
\end{aligned}
\end{equation}
where $s_t$ is the intermediate reasoning state, $(i_t, q_t)$ is the selected tool-and-query pair, $r_t$ is the returned result, and $\Vert$ denotes appending to the interaction history.
The loop halts when the policy or model emits a STOP decision, after which an output decoder $g$ produces the final answer $y = g(h_T)$.
Both the reasoning core and the tool internals are black-box; what is exposed, optimized, and evaluated is the interface history and the policy that controls it.

\subsubsection{Methodology}
Two design choices interact: the \emph{tool catalogue} determines which black-box capabilities are available, and the \emph{orchestration policy} determines how they are sequenced, checked, and folded back into the model's spatial answer.
The catalogue roles below are complementary, not sequential.

\paragraph{Execution-Augmented Reasoning.}
Execution tools bridge the semantic--geometric gap between high-level language and low-level feasible action~\cite{zhang2025scaling}: they expose callable primitives such as motion planners, robot skills, affordance estimators, navigation controllers, or simulator actions, and the orchestrator translates a spatial intention into a sequence of calls, monitors whether the result satisfies the goal, and recovers when a tool fails.
Representative systems include affordance-grounded planning, where language-model outputs are modulated by value functions for skill feasibility~\cite{ahn2022icanisay}; API- or code-mediated manipulation, where perception and control primitives are composed through generated calls~\cite{liang2023code,huang2023voxposer}; and action-token models that fold robot actions into the model vocabulary~\cite{brohan2022rt,zitkovich2023rt2,driess2023palm}.
SpatialClaw~\cite{cho2026spatialclaw} studies the action interface itself: a stateful Python kernel pre-loaded with perception and geometry primitives lets the agent execute one analysis step, observe intermediate outputs, and adapt the next call rather than committing to a full static program.
Closed-loop feedback and hierarchical decomposition~\cite{huang2022innermonologue,singh2023progprompt,yang2024vlmtamp,chen2024autotamp,bhat2024brainbodyllm} push this family toward agentic orchestration.

\paragraph{Representation-Augmented Reasoning.}
Representation tools externalize spatial working memory through sketches, 3D renders, CAD systems, scene graphs, game engines, or visualization buffers, motivated by the observation that transformers struggle to manipulate fine-grained geometry internally~\cite{hu2024visualsketchpad}.
Externalization reduces hallucination by moving spatial computation out of parametric memory~\cite{hu2024visualsketchpad} and enables iterative self-correction through visual feedback loops~\cite{wu2025reinforcing}; the claim is that spatial reasoning improves when models can \textit{see what they think}, as in Vision-Augmented Prompting~\cite{xiao2024enhancing}, which interleaves image synthesis with chain-of-thought so that textual and visual reasoning can be reconciled.
Recent spatial-tool agents make the evidence acquisition active: PERIA~\cite{li2026peria} pairs perception tools that expose symbolic or spatial evidence with interaction tools that manipulate the visual context, learning multi-tool behavior from synthesized trajectories and reinforcement learning, and S-Agent~\cite{dai2026sagent} casts the VLM as a semantic planner over a hierarchy of spatial tools that lift 2D object evidence into 3D geometric support while accumulating scene and agent memory across multi-view images or video.

\paragraph{Validation-Augmented Reasoning.}
Validation tools check whether intermediate outputs are logically, causally, geometrically, or physically consistent: the orchestrator translates a candidate answer or plan into a verifiable representation and queries a planner, solver, judge, or simulator.
% symbolic planners
Generated spatial plans can be encoded as constraint-satisfaction problems or PDDL specifications and validated by symbolic planners or SMT solvers~\cite{liu2023llmp,hao2023reasoning,yang2023coupling}, which checks preconditions, geometric relations, and causal effects;
% physics simulators as verifiers
physics simulators such as MuJoCo~\cite{todorov2012mujoco} and Isaac Sim\footnote{\url{https://developer.nvidia.com/isaac/sim}} act as physical oracles for the feasibility of proposed movements, collisions, and trajectories~\cite{liu2022mindseye,huang2023voxposer,chen2024spatialvlm};
% Neuro-symbolic methods
neuro-symbolic systems integrate Answer Set Programming solvers that translate natural language into formally verifiable logical rules for relational and topological coherence~\cite{yang2020neurasp,premsri2024neurosymbolic};
% plan-critic
and plan-critic mechanisms~\cite{lei2025stma,wang2025multiagent} or an LLM-as-a-Judge that evaluates and refines the plans of a primary planner close a self-correcting loop~\cite{hariharan2025plan}.


\paragraph{Memory-Augmented Reasoning.}
Memory tools expose external storage and retrieval through a black-box query interface: the orchestrator asks for relevant prior observations, spatial subgraphs, object histories, or task episodes and decides how to incorporate what is returned, which addresses the \textit{``crisis of context''} in long-horizon robotic tasks whose sensory streams exceed the finite context window of an LLM.
At this level the memory may be a hierarchical cognitive architecture such as RoboMemory or an open-vocabulary 3D scene graph such as ConceptGraphs (both treated in §\ref{subsec:p3_s7_3_6}), a classical 3D or dynamic scene graph~\cite{armeni20193dsg,rosinol20203ddsg}, or a system-level memory manager: MemOS~\cite{li2025memos} introduces paging and scheduling that swap spatial contexts in and out of the reasoning core, treating the LLM as a CPU and the tool as RAM, and spatio-temporal retrieval engines such as ReMEmbR~\cite{anwar2025remembr} and MemoTime~\cite{tan2025memotime} index video history in vector databases so that an agent can answer ``where did I leave the keys five minutes ago?''.

\paragraph{Predictive Tools and World-Model Feedback.}
Some catalogues include tools that simulate counterfactual or future spatial states rather than observe or verify the present one: generative video models, learned dynamics models, and predictive latent models are called as expensive but powerful instruments for ``what would happen if'' or ``what would this look like from another viewpoint''~\cite{liu2022mindseye,yang2025cambrians}.
Here they are one class of callable tool; their feedback forms, imagination cost, and fidelity are the subject of §\ref{subsec:p3_s7_3_8}.


\subsubsection{Benchmarks}
Evaluation asks whether the orchestrator selects useful tools and converts their outputs into better spatial decisions.
For execution-augmented tool use the canonical suites are SayCan~\cite{ahn2022icanisay} (LLM-grounded household manipulation), ALFWorld~\cite{shridhar2020alfworld} and VirtualHome~\cite{puig2018virtualhome} (long-horizon text-and-embodied environments), and Voyager-on-Minecraft~\cite{wang2023voyager} (open-world skill accretion through code-as-policy).
For spatial-reasoning-specific tool use, SpaceTools~\cite{chen2025spacetools} targets the metric gap between VLM perception and physical manipulation, and Cambrian-S~\cite{yang2025cambrians} introduces the spatial-supersensing framework that exercises tool augmentation from perception (Stage~I) to predictive world modeling (Stage~IV).
Validation- and memory-augmented reasoning are evaluated on the memory and agentic benchmarks of §\ref{subsec:p3_s7_3_6} and §\ref{subsec:p3_s7_3_7}, and predictive tools on those of §\ref{subsec:p3_s7_3_8}, where the central question is whether imagined evidence improves the final answer enough to justify its cost.


\subsubsection{Challenges and Trends}
The recurring limitations originate in the LLM--tool interface rather than in any individual tool.
The \textbf{exploration bottleneck} is combinatorial: each tool adds an action dimension and tool composition multiplies them, so naive reinforcement learning over tool use is inefficient; Double Interactive Reinforcement Learning (DIRL)~\cite{chen2025spacetools} warm-starts from expert demonstrations before open exploration, but the scaling of the action space itself is unresolved.
The \textbf{metric gap} moves rather than disappears: VLMs excel at semantic perception but fail at the precise metric estimation manipulation requires, and offloading geometry to depth estimators and pose solvers lets the LLM reason over measurements instead of hallucinations~\cite{chen2025spacetools,zhang2025scaling}, yet query formulation, output parsing, and error attribution at the interface become the new precision bottleneck.
\textbf{Interface calibration} compounds this: tool outputs arrive with heterogeneous units, confidence scores, coordinate frames, and failure modes, and an orchestrator that cannot judge whether a returned depth map, solver result, retrieved memory, or imagined rollout is trustworthy can perform worse than a direct model response.

\emph{Learned orchestration} is one direction: policies are moving from hand-written call rules toward imitation- or reinforcement-learned controllers that decide which tool is worth its latency and uncertainty budget~\cite{chen2025spacetools}.
A parallel shift is the \emph{System-1 / System-2 architectural split}, in which the neural backbone handles fast, intuitive perception (System~1) while the toolset handles slow, deliberative, verifiable reasoning (System~2)~\cite{yang2025neurosymbo}, achieving both flexibility and rigor.


\subsection{Retrieval- \& Memory-augmented Spatial Reasoning}
\label{subsec:p3_s7_3_6}
\label{ssec:memory_spatial_reasoning}

Retrieval- and memory-augmented spatial reasoning is best understood through a \emph{Planner-Memory-Retriever} decomposition. The Planner $\mathcal{P}$ decides what spatial information is needed for the current instruction; the Memory Bank $\mathcal{M}$ stores an external, updateable state of the world; and the Retriever $\mathcal{R}$ selects the subset of that memory that should condition the next reasoning or action step. Unlike one-shot retrieval-augmented generation, spatial memory must also support \emph{write} and \emph{update}: an agent may discover that a door is blocked, that an object has moved, or that a previous topological belief was wrong, and the memory state must change accordingly.

The memory bank is therefore not merely a text corpus. In spatial domains, $\mathcal{M}$ may contain vector embeddings of views and instructions, metric maps and voxel belief states, dynamic 3D scene graphs, frontier snapshots, past trajectories, user preferences, and brain-inspired episodic archives. The central problem is to make these heterogeneous memories queryable: semantic retrieval must locate relevant objects and events, while geometric retrieval must preserve relations such as containment, support, reachability, and path connectivity.

\begin{figure}[t]
    \centering
    \includegraphics[width=\linewidth]{figures/part3/section7/retrieval_memory_augmented_spatial_reasoning.pdf}
    \caption{Retrieval- and memory-augmented spatial reasoning: (i) a query and a memory bank of vector records, spatial scene graphs, and episodic archives, (ii) retrieval and update via spatial RAG, dynamic 3D scene graphs, or layered memory, and (iii) alternative output forms rather than joint outputs of one query. In the example, the keys were on the kitchen counter five minutes ago and are in the hallway now. Scenes and retrieval diagrams are illustrative, not measured similarity scores.}
    \label{fig:p3_s7_3_6_memory_formulation}
\end{figure}

\subsubsection{Task Definition}
Where §\ref{subsec:p3_s7_3_5} treats memory as one black-box tool among others, here the \emph{external memory itself} is the central learning target. Formally, a retrieval- \& memory-augmented system is the triple $(\mathcal{P}, \mathcal{M}, \mathcal{R})$, where $\mathcal{P}$ is a reasoning core (an LLM, VLM, or embodied planner), $\mathcal{M}$ is an external memory store, and $\mathcal{R}: \mathcal{Q} \times \mathcal{M} \to \mathcal{C}$ maps a query $q_t$ and the current memory state to a relevant context fragment $c_t$. Given an instruction $x$, an observation stream $o_t$, and a reasoning history $h_t$ (initially $h_0 = \varnothing$), the system proceeds by the iterative update
\begin{equation}
\begin{aligned}
q_t &= \mathrm{Query}_{\mathcal{P}}(x, o_t, h_{t-1}), &
c_t &= \mathcal{R}(q_t, \mathcal{M}_{t-1}), &
s_t &= \mathcal{P}(x, o_t, c_t, h_{t-1}), \\
\mathcal{M}_t &= \mathrm{Write}(\mathcal{M}_{t-1}, o_t, s_t, c_t), &
h_t &= h_{t-1} \cup \{(o_t, q_t, c_t, s_t)\},
\end{aligned}
\end{equation}
where $\mathrm{Query}_{\mathcal{P}}$ extracts the planner's information need, $\mathcal{R}$ returns the retrieved context, and $\mathrm{Write}$ inserts, revises, or invalidates memory entries using the newest observation and decision state. The decomposition brings three concrete advantages: \emph{scalability}, by letting agents operate over large environments while fetching only task-relevant subgraphs or vectors; \emph{personalization}, through targeted retrieval of user profiles, historical preferences, and episodic interactions; and \emph{explainability}, since decisions can be traced to the specific memories, rules, or map fragments that were retrieved.

\subsubsection{Methodology}
Three families of memory architecture instantiate the $(\mathcal{P}, \mathcal{M}, \mathcal{R})$ triple at different levels of abstraction and scale.

\paragraph{Retrieval-Augmented Generation for Navigation and Instruction.}
Vision-and-Language Navigation (VLN) demands a tight coupling of visual perception, language understanding, and action generation, and retrieval-augmented generation (RAG) bridges high-level planning and local observation by externalizing spatial context; the design question is what the memory bank should store and how the retriever should index it.

\emph{Textual abstraction vs.\ visual fidelity.} NavRAG~\cite{wang2025navrag} argues for hierarchical abstraction: a ``Hierarchical Scene Description Tree'' converts the 3D environment into structured text (Region $\rightarrow$ Object $\rightarrow$ Attribute), so the retriever can access information at the granularity the query requires, at the cost of visual nuance. VisRAG~\cite{yu2024visrag} takes the opposite stance with a multimodal pipeline that bypasses text parsing entirely: a VLM retrieves images directly from the query, preserving texture, lighting, and style that text-only retrieval discards.

\emph{Privileged maps vs.\ egocentric mental simulation.} A second line replaces privileged global maps with human-like episodic memory. GoViG~\cite{wu2025govig} generates instructions solely from egocentric observations of start and goal states, alternating between predicting the next visual view and generating the corresponding instruction, so that the navigational narrative is built on a hallucinated path; retrieval shifts from static facts to potential futures.

\emph{Hybridizing spatial and semantic logic.} Retrieval in physical spaces requires logic that pure semantic search cannot provide (\eg,\ ``near'', ``inside'', ``along''). Spatial-RAG~\cite{yu2025spatial} combines sparse spatial filtering (SQL-based constraints) with dense semantic matching, formulating retrieval as multi-objective optimization, which is what grounds CoT reasoning~\cite{liu2023retrieval} so that multi-step plans respect geometric feasibility.

\paragraph{Dynamic 3D Scene Graphs and Subgraph Retrieval.}
As environments grow, flat memory representations such as point clouds become intractable, and hierarchical abstraction through 3D scene graphs (3DSGs) has taken over: Hydra~\cite{hughes2022hydra} builds metric-semantic scene graphs in real time, and open-vocabulary successors embed vision--language features directly into graph nodes~\cite{gu2024conceptgraphs}, turning scene graphs from rigid structural maps into queryable vector stores. Because global graphs overwhelm planner context windows, processing has moved from holistic graph analysis to sparse subgraph retrieval, through semantic pruning~\cite{rana2023sayplan} or RAG-based mechanisms~\cite{xie2024embodiedrag} that isolate task-relevant topological neighborhoods. The convergence of graph and vector retrieval matters for spatial memory in particular: graph edges preserve physical relations, while vector neighborhoods preserve open-vocabulary semantic similarity.

\paragraph{Unified and Brain-Inspired Memory Architectures.}
RAG provides access to information; the organization of that information is equally critical, and recent architectures borrow the brain's division of labor. RoboMemory~\cite{lei2025robomemory} integrates working, short-term, and long-term memory in one cognitive framework so that the agent iteratively refines its model of the environment, and 3D-Mem~\cite{yang20253dmem} takes a visual-centric route, using ``frontier snapshots'' to manage exploration and hallucinate potential future states. 3DLLM-Mem~\cite{hu20253dllmmem} grounds the integration of retrieved memory into the inference stream with a memory fusion layer that cross-attends from the current instruction into the stored 3D memory bank:
\begin{equation}
f^Q_{\text{fuse}} = \mathrm{Softmax}\!\left( \frac{f^Q_t (f^K)^\top}{\sqrt{C}} \right) f^V,
\end{equation}
where $f^Q_t$ are the query features of the current instruction and $f^K$, $f^V$ the keys and values derived from the memory bank. Addressing specific failure modes, VLM$^2$~\cite{liu2025vision} introduces a dual-memory module for the semantic--geometric misalignment of video-based spatial reasoning, and Vision-to-Geometry~\cite{cai2025vision} incorporates geometric cues such as field-of-view coverage to prevent redundant re-exploration.



\subsubsection{Benchmarks}
Progress in memory-augmented spatial reasoning requires separating \emph{memory-specific} tests from broad spatial probes.

\paragraph{Memory-specific evaluation.}
3DMem-Bench~\cite{hu20253dllmmem} serves as a standard for assessing long-term persistence: comprising over 26{,}000 trajectories, it challenges agents with Spatial--Temporal QA (\eg,\ ``What color was the bed in the room you visited five minutes ago?''), requiring the retrieval of past observations that are no longer in the active context. Moving beyond static environments, EmbodiedBench~\cite{yang2025embodiedbench} tests long-horizon planning where objects are moved or manipulated during the task, benchmarking the \emph{write/update} capability of memory systems, their ability to overwrite obsolete data without catastrophic forgetting. OpenEQA~\cite{majumdar2024openeqa} offers a complementary episodic-memory setting by asking questions about an environment the agent has previously explored, making it useful for measuring whether retrieved experience rather than immediate perception supports the answer.

\paragraph{Broad spatial probes as diagnostic complements.}
Benchmarks such as SpaCE-10~\cite{gong2025space10}, SITE~\cite{wang2025site}, and the static probes collected in Table~\ref{tab:spatial_reasoning_benchmark} (OmniSpatial~\cite{jia2025omnispatial}, Spatial-DISE~\cite{huang2025spatialdise}, Open3DVQA~\cite{zhang2025open3dvqa}) evaluate spatial orientation, perspective taking, metric estimation, open-space 3D relations, and intrinsic/extrinsic static or dynamic reasoning. They are valuable for diagnosing what information a memory system should retain (viewpoint-dependent relations, object-centric attributes, multi-step multi-view dependencies), but they are not memory benchmarks unless the protocol hides necessary information in a prior episode and requires retrieval from $\mathcal{M}$: a model can fail a broad spatial probe because it lacks geometric perception, because it retrieved the wrong memory, or because the memory was never written in the first place.


\subsubsection{Challenges and Trends}
The recurring limitations originate in the Planner-Memory-Retriever interface rather than in any individual component.
\textbf{Query formation under unbounded context} is the central bottleneck, not memory size alone: the act of deciding what to retrieve requires compressing the current task, viewpoint, uncertainty, and history into a finite representation, and memory paging and scheduling alleviate but do not remove the need for a learned retrieval policy.
\textbf{Write/update consistency} is where append-only memories become dangerous: agents need mechanisms for reconciling conflicting observations, invalidating stale object locations, and representing uncertainty when the world may have changed since the last visit; EmbodiedBench-style protocols make the issue visible, but robust update semantics remain open.
\textbf{Integration mechanics} are unresolved between in-context learning, which concatenates retrieved text to the prompt (simple but token-inefficient), and cross-attention fusion, which injects retrieved 3D features into attention layers (3DLLM-Mem~\cite{hu20253dllmmem}; expressive but architecture-coupled), with no clear winner on the fidelity--modularity trade-off.
\textbf{Modality mismatch} has no general theory: RAG-over-text (NavRAG~\cite{wang2025navrag}) discards visual nuance, RAG-over-images (VisRAG~\cite{yu2024visrag}) loses structural priors, and when to retrieve text, images, or 3D representations remains task-dependent.

Memory is shifting \emph{from static maps to dynamic semantic belief states}: systems store probabilities, semantic labels, and interaction histories rather than coordinates, enabling agents to reason under uncertainty about where an object \emph{likely} resides given commonsense and past observation, as RoboMemory~\cite{lei2025robomemory} typifies.
Graph- and vector-based retrieval are converging into hybrid architectures that use vector search to identify graph entry nodes and topological traversal for neighborhood reasoning.
Most abstractly, \emph{meta-cognitive memory} such as Meta-Memory~\cite{mao2025meta} generates task-specific ``cognitive maps'' and retrieves \emph{strategies} rather than raw data, pointing toward learnable retrieval policies and bridging into the agentic orchestration of §\ref{subsec:p3_s7_3_7}.


\subsection{Agentic Spatial Reasoning}
\label{subsec:p3_s7_3_7}

Agentic spatial reasoning asks how a model organizes \emph{interaction} with space. The core object is not a single chain-of-thought prompt, tool call, memory store, VLA policy, or world model, but an orchestration policy that repeatedly decides what to perceive, what to infer, what to do, and how to revise its state after feedback. In this closed-loop view, a wall is not only a geometric plane but also a constraint on motion; a cup is not only an object category but also a potential affordance; and a failed action is not only an error but also evidence that should reshape the next plan.

``Agentic'' is therefore a composition principle: the rationales, programs, tools, memory, VLA policies, and world models of the preceding subsections become resources inside a perceive--reason--act--reflect loop, and the research question is how the loop is scheduled and corrected: when should the agent gather more evidence, retrieve memory, call a geometric tool, simulate a future, ask another agent, or execute an action?

\begin{figure}[t]
    \centering
    \includegraphics[width=\linewidth]{figures/part3/section7/agentic_spatial_reasoning.pdf}
    \caption{Agentic spatial reasoning as closed-loop orchestration. Top: (i) instruction, observation, robot state, tools, and memory, (ii) the perceive--reason--act--reflect loop under planning core $\mathcal{P}$ and policy $\pi$, and (iii) complementary, not simultaneous, output forms; the door-state sequence is a prediction and the failed counter search is feedback from a different step. Bottom: (a)--(d) mechanisms that support the loop. The scene is a generated illustration, not an experimental observation.}
    \label{fig:p3_s7_3_7_agentic_formulation}
\end{figure}

\subsubsection{Task Definition}
Whereas §\ref{subsec:p3_s7_3_5} treats tools as black-box operators and §\ref{subsec:p3_s7_3_6} treats external memory as the central learning target, agentic spatial reasoning composes both into a closed-loop cycle in which the orchestration policy itself becomes the learning target. Formally, an agentic spatial reasoning system is the tuple $(\mathcal{P}, \mathcal{S}, \mathcal{A}, \mathcal{F}, \pi)$, where $\mathcal{P}$ is a planning core (an LLM or VLM), $\mathcal{S}$ is the spatial state space (perception, internal map, history), $\mathcal{A}$ is the action space (sensorimotor commands or sub-goals), $\mathcal{F}$ is a reflection / critic function that scores proposed actions or rationales, and $\pi$ is the orchestration policy. Given an instruction $x$, the system maintains a history $h_t$ (initially $h_0 = \varnothing$) and proceeds by the iterative perceive--reason--act--reflect cycle
\begin{equation}
\begin{aligned}
o_t &= \mathrm{Perceive}(\text{env}), &
s_t &= \mathcal{P}(o_t, h_{t-1}), \\
a_t &\sim \pi(\cdot \mid s_t), &
r_t &= \mathcal{F}(s_t, a_t, o_{t+1}), \\
h_t &= h_{t-1} \cup \{(o_t, s_t, a_t, r_t)\},
\end{aligned}
\end{equation}
where $\mathrm{Perceive}$ extracts the current spatial observation, $\pi$ samples a proposed action, $\mathcal{F}$ supplies a reflection signal (verbal critique, planner-critic verdict, or value estimate), and the history accumulates all four. This formulation makes explicit that ``agentic'' is not a new capability but rather a structural promotion: the orchestration policy $\pi$ becomes the central learning target rather than the planner $\mathcal{P}$ or any individual tool.

\subsubsection{Methodology}
Four orchestration problems organize the methods: grounding observations into executable state, closing the feedback loop, coordinating across agents, and using action and world-model modules without reducing agency to any one of them.

\paragraph{Grounding Observations into Executable State.}
The first responsibility of an agentic loop is to convert an instruction and observation into a state that can support action, preserving referential consistency between the agent's internal model and the physical environment.
Probabilistic affordance alignment, as in SayCan~\cite{ahn2022icanisay}, PaLM-E~\cite{driess2023palm}, and RT-2~\cite{zitkovich2023rt2}, weights the likelihood of an instruction (``pick up the apple'') by the geometric feasibility of the action in the current scene; spatial reasoning stays implicit, and the orchestration layer must decide whether that is sufficient or whether to request an explicit representation.
Explicit intermediate representations supply the precision implicit alignment lacks: SpatialCoT~\cite{liu2025spatialcot} compels the VLM to output metric coordinates inside its chain of thought, a ``think-coordinate-act'' cycle that creates a verifiable spatial bottleneck; Embodied-R1~\cite{yuan2025embodiedr1} uses reinforcement learning to refine a pointing representation so that the model outputs 3D reference points (spatial anchors) that decouple high-level planning from motor control; and procedural querying treats the environment as a queryable database, through dynamically synthesized APIs (VADAR, §\ref{subsec:p3_s7_3_4}) or a stateful code kernel that executes one cell at a time and revises the next from the observed outputs (SpatialClaw, §\ref{subsec:p3_s7_3_5}).
Beyond the immediate percept, agents ground reasoning in persistent structured abstractions: object-centric 3D scene graphs such as ConceptGraphs (§\ref{subsec:p3_s7_3_6}) let the agent reason over topological relations (``behind'', ``supported by'') rather than raw geometry, and MapGPT~\cite{chen2024mapgpt} converts topological maps into textual prompts that give the LLM a global view of connectivity for long-horizon planning.

\paragraph{Closing the Feedback Loop.}
Agentic spatial reasoning is a temporal \emph{process} rather than a single inference, and the methodological shift is from producing one rationale to maintaining a loop that observes the consequences of its own choices.
The chain-of-thought paradigm (§\ref{subsec:p3_s7_3_3}) improved reasoning depth but remained open-loop; ReAct~\cite{yao2022react} closed the loop by interleaving reasoning traces with executable actions so that spatial outcomes are observed before the next step is planned. Observation alone is often insufficient for correction, so Reflexion~\cite{shinn2023reflexion} introduced ``verbal reinforcement learning'', in which the agent writes explicit critiques of its failures into short-term context, and DEPS~\cite{wang2023describe} forces the agent to \emph{explain} failures in geometric terms before choosing new sub-goals, turning retrying into debugging.
For coherence across long horizons, agents need memory beyond the context window: Voyager (§\ref{subsec:p3_s7_3_4}) stores successful spatial behaviors as executable code in a vector database and retrieves them when a similar context recurs, consolidating transient experience into procedural knowledge.
Reflection can also precede execution. Reasoning via Planning (RAP)~\cite{hao2023reasoning} and Tree of Thoughts~\cite{yao2023tot} explore branching spatial possibilities in a latent space; STMA~\cite{lei2025stma} ties such simulation to dynamic knowledge graphs with planner--critic loops that validate whether a proposed path is topologically valid; and vision-language world models such as VLWM~\cite{chen2025vlwm} roll out semantically described future states so that the orchestration policy can compare hypothetical futures before committing.

\paragraph{Coordinating Multiple Spatial Agents.}
In shared environments spatial reasoning becomes distributed, and movement has to be negotiated across perspectives and intents.
Multi-agent reinforcement learning induces spatial cooperation without explicit communication, as in the SMAC formations learned by decentralized actor--critic and value-decomposition methods~\cite{samvelyan2019smac} (§\ref{subsec:p3_s7_2_6}), but there each agent treats the others as dynamic obstacles and has no symbolic model of its partners.
A \emph{spatial theory of mind} adds the ability to model another agent's geometric viewpoint, transforming egocentric perception into an allocentric frame (``my left'' is ``your right''): RoCo~\cite{mandi2024roco} uses foundation models for embodied dialog in which agents disambiguate spatial references by simulating the partner's visual field, and ProAgent~\cite{zhang2024proagent} shows agents communicating to update a shared map and synchronize movement, converging on a common ``language of space''.

\paragraph{Orchestrating Action and World-Model Modules.}
Generalist VLA policies and world models are components of agentic systems, but they do not by themselves define agency; they become agentic when an orchestration policy uses them inside the loop, with a VLA module executing a sub-goal, a world model predicting the consequence of a candidate action, and a reflection module deciding whether to retry, replan, or gather more evidence.
What makes an action representation useful to a loop is that it exposes intermediate commitments the loop can inspect, revise, or hand to a controller: beyond the unified action interfaces of RT-2~\cite{zitkovich2023rt2}, Gato~\cite{reed2022generalist}, and Octo~\cite{team2024octo}, world-language-action models~\cite{yang2026wla} jointly predict textual subtasks, subgoal images, and robot actions, and ACoT-VLA~\cite{zhong2026acotvla} reasons through coarse action intents before fine-grained control.
Cross-embodiment transfer (§\ref{subsec:p3_s7_3_2}; RoboCat's self-improvement cycles adapt to new bodies~\cite{bousmalis2023robocat}) becomes an orchestration problem too, since the same high-level plan may require different perception checks, affordance tests, and recovery routines on different bodies.
World models (§\ref{subsec:p3_s7_3_8}), from DreamerV3~\cite{hafner2023mastering} and Genie~\cite{genie3} to WorldPilot-style prediction inside planning loops~\cite{lin2026worldpilot}, serve the loop as cognitive simulators: the agentic use is not video generation for its own sake but closed-loop decision support, in which the agent imagines, critiques, and then acts.





\subsubsection{Benchmarks}
Evaluating agentic spatial reasoning requires benchmarks that go beyond static visual question answering: a useful protocol tests whether the agent can gather missing observations, choose tools or actions, use memory across long horizons, and revise its plan after feedback, and static spatial probes should be read as component diagnostics rather than direct evidence of closed-loop agency.
Two protocols exercise the loop directly: iVISPAR~\cite{mayer2025ivispar} evaluates VLMs as agents in an interactive sliding-tile setting with visual 2D, visual 3D, and text inputs, where the answer depends on multi-step state transitions rather than one-shot recognition, and TIR-Bench~\cite{li2025tirbench} requires agents to use tools (zoom, crop, rotate) to manipulate images while solving spatial puzzles.
Memory-centered settings (OpenEQA~\cite{majumdar2024openeqa}, §\ref{subsec:p3_s7_3_6}) and open-3D probes such as Open3DVQA~\cite{zhang2025open3dvqa} test what the loop must retain, and the broad visual-spatial probes (SpatialEval~\cite{wang2024picture}, VSI-Bench~\cite{yang2025vsibench}, OmniSpatial~\cite{jia2025omnispatial}, Spatial-DISE~\cite{huang2025spatialdise}) identify weak spatial primitives along the axes of static vs.\ dynamic content, single vs.\ multi-view, ego vs.\ exo perspective, and seen vs.\ unseen queried information; Table~\ref{tab:spatial_reasoning_benchmark} maps the coverage.

\begin{table*}[!ht]
  \caption{Representative spatial-reasoning benchmarks discussed in §\ref{subsec:p3_s7_3_7}. Perspective marks whether ego and exo viewpoints are explicitly represented; Interactive marks protocols that exercise action or tool selection rather than only answer generation; Query Target is whether the queried information is visible (Seen), hidden (Unseen), or an interactive state (State). \textit{Undeclared} marks attributes the source paper does not state.}
  \label{tab:spatial_reasoning_benchmark}
  \centering
  \resizebox{\textwidth}{!}{
  \begin{tabular}{@{}lccccccccc}
    \toprule
      \multirow{2}{*}{\textbf{Benchmark}}
      & \multirow{2}{*}{\textbf{QA / Tasks}}
      & \multirow{2}{*}{\textbf{Modality}}
      & \multirow{2}{*}{\textbf{Scene}}
      & \multicolumn{2}{c}{\textbf{View Type}}
      & \multicolumn{2}{c}{\textbf{Perspective}}
      & \multirow{2}{*}{\textbf{Interactive}}
      & \multirow{2}{*}{\textbf{Query Target}} \\
      \cmidrule(lr){5-6}
        \cmidrule(lr){7-8}
        & & &
        & \textbf{Single} & \textbf{Multi}
        & \textbf{Ego} & \textbf{Exo}
        & & \\
    \midrule
    3DSRBench~\cite{ma20253dsrbench} & 2{,}772 & Image & Static & \checkmark  & \crossmark & \crossmark  &  \checkmark & \crossmark & Seen\\
    InternSpatial-Bench~\cite{deng2025internspatial} &  6{,}008 & Image & Static & \checkmark  & \crossmark & \crossmark  &  \checkmark & \crossmark & Seen \\
    OmniSpatial~\cite{jia2025omnispatial} & $\sim$8{,}400 & Image & Static & \checkmark & \crossmark & \checkmark & \checkmark  & \crossmark & Seen\\
    Spatial-DISE~\cite{huang2025spatialdise} & 559 & Image & Mixed & \checkmark & \checkmark & \multicolumn{2}{c}{\cellcolor{lightgray}\textit{Taxonomy}} & \crossmark & Seen\\
    Open3DVQA~\cite{zhang2025open3dvqa} & 9{,}000 & 3D / Sim & Static & \checkmark & \crossmark & \checkmark & \checkmark & \crossmark & Seen\\
    iVISPAR~\cite{mayer2025ivispar} & Puzzles & 2D/3D/Text & Dynamic & \checkmark & \crossmark & \multicolumn{2}{c}{\cellcolor{lightgray}\textit{N/A}} & \checkmark & State\\
    Spatial457~\cite{wang2025spatial457} & 23{,}752 & Image & Static & \checkmark & \crossmark & \multicolumn{2}{c}{\cellcolor{lightgray}\textit{Undeclared}}  & \crossmark & Seen\\
    All-Angles Bench~\cite{yeh2025seeing} & $\sim$2{,}100 & Image & Static  & \crossmark & \checkmark & \multicolumn{2}{c}{\cellcolor{lightgray}\textit{Undeclared}} & \crossmark  & Seen \\
    EmbSpatial-Bench~\cite{du2024embspatial} & 3{,}640 & Image & Static & \checkmark  & \crossmark & \checkmark & \crossmark & \crossmark &  Seen \\
    MMSI-Bench~\cite{yang2025mmsi} & 1{,}000 & Image & Static & \crossmark & \checkmark  & \multicolumn{2}{c}{\cellcolor{lightgray}\textit{Undeclared}} & \crossmark & Seen\\
    MindCube~\cite{yin2025spatial} & 21{,}154 & Image & Static & \crossmark & \checkmark & \multicolumn{2}{c}{\cellcolor{lightgray}\textit{Undeclared}}  & \crossmark & Unseen \\
    VSI-Bench~\cite{yang2025vsibench} & 5{,}130 & Video & Static & \checkmark & \crossmark & \checkmark & \crossmark & \crossmark & Seen\\
    VLM4D~\cite{zhou2025vlm4d} & $\sim$1{,}800 & Video & Dynamic & \checkmark & \crossmark &  \checkmark & \checkmark & \crossmark & Seen \\
    STI-Bench~\cite{li2025sti} & 2{,}064 & Video & Dynamic & \checkmark & \crossmark & \checkmark & \checkmark & \crossmark &  Seen\\
    DSI-Bench~\cite{zhang2025dsi} & $\sim$1{,}700 & Video & Dynamic & \checkmark & \crossmark & \multicolumn{2}{c}{\cellcolor{lightgray}\textit{Undeclared}} & \crossmark &  Seen\\
    \bottomrule
  \end{tabular}
  }
\end{table*}

Table~\ref{tab:spatial_reasoning_benchmark} should be read as a coverage map, not as a leaderboard for agentic systems. Most entries evaluate static answer generation or localized visual-spatial primitives; only interactive protocols such as iVISPAR and tool-use settings directly exercise action selection and feedback. Claims about agentic spatial reasoning should therefore report loop-level metrics (observation gathering, action choice, recovery after failure, memory update, and final task success) rather than single-turn accuracy alone.

\subsubsection{Challenges and Trends}
Boundaries are collapsing between language and action (VLAs), between robot bodies (cross-embodiment), and between perception and imagination (world models), and two difficulties gate progress toward generalist spatial agents.
The \textbf{sim-to-real fidelity gap}: high-fidelity simulators (predictive world models, generative video) accelerate training but introduce systematic biases that physical execution exposes, and the cost--precision frontier for transfer remains open.
\textbf{Long-horizon spatiotemporal coherence}: across multi-hour episodes the agent's internal map drifts, and episodic memory and reflection reduce but do not eliminate the drift, because the bottleneck is the finite context window of the planning core.

\emph{Explicit orchestration over observation--state--action--reflection tuples} is one direction: VLA and WLA systems are beginning to expose intermediate subgoals, action intents, and predicted states that a loop can inspect rather than hiding all control inside a single forward pass. A second is \emph{mental simulation via generative world models}: DreamerV3-style, Genie-style, VLWM-style, and WorldPilot-style models let agents verify spatial logic in imagination before acting, shifting the paradigm from ``trial-and-error'' toward ``trial-in-imagination'' while keeping the agentic focus on when to simulate, when to act, and when to revise.


\subsection{World-Model-based Spatial Reasoning}
\label{subsec:p3_s7_3_8}


% Test-Time Scaling with World Models for Spatial Reasoning
% When and How Much to Imagine: Adaptive Test-Time Scaling with World Models for Visual Spatial Reasoning
% From 2D to 3D Cognition: A Brief Survey of General World Models
% Vagen: Reinforcing world model reasoning for multi-turn vlm agents
% Scaling and Beyond: Advancing Spatial Reasoning in MLLMs Requires New Recipes
% SpatialDreamer: Incentivizing Spatial Reasoning via Active Mental Imagery
% Think with 3D: Geometric Imagination Grounded Spatial Reasoning from Limited Views



% 在这里想要重新梳理一下应该描述的内容：
% 在这个章节中更想强调的是如何使用world model as a tool来实现spatial reasoning task，而不是去讨论world model本身的训练和发展。也就是说，world model是一个工具，spatial reasoning task是一个目标。我们要定义的是如何用这个工具去实现这个目标。
% 并且需要在文中也要提到的是，world model的具体构建和训练是在后续的章节8.6 World Modeling中

% world 作为一个工具其核心的能力就是模拟了世界，因此能够在spatial reasoning过程中提供一些反馈，这个反馈可以是多种形式的，比如生成未来的图片，或者是生成未来的状态，或者是生成未来的动作序列。然后这些反馈可以被用来提升spatial reasoning。
% 另外，world model 也可以得到 latent representation，比如利用JEPA得到的 latent representation，这些latent representation也可以被用来提升spatial reasoning。


A world model is a neural simulator of how world states transition, conditioned on an agent's actions and/or external multi-modal context.
Trained on large-scale episode trajectories, it acquires a generalizable, latent model of environment dynamics.
This subsection treats that model as a \emph{tool}: the object of study is the spatial-reasoning task, and the world model is the instrument used to solve it.
Its defining capability is that it \emph{simulates} the world and can therefore return a predictive feedback signal, a rollout of imagined futures or a dynamics-aware latent, that a reasoner consumes to answer a spatial query or to plan.
We accordingly organize the discussion by \emph{how the world model is queried and what feedback it returns}, rather than by how it is built.
The complementary question of how world models are themselves constructed, trained, and scaled is the subject of §\ref{sec:worldmodeling} (World Modeling); here we take a pretrained world model as given.
Figure~\ref{fig:p3_s7_3_8_world_model_formulation} schematizes this tool-use formulation; its two usage modes, understanding-oriented and planning-oriented, cut across the feedback forms surveyed below.

\begin{figure}[!t]
    \centering
    \includegraphics[width=0.99\linewidth]{figures/part3/section7/world_model_based_spatial_reasoning.pdf}
    \caption{World-model-based spatial reasoning read as tool use: (i) a spatial query $q$ with context $c$ and optional actions $a_{0:T-1}$; (ii) a pretrained world model (constructed and trained as in §\ref{sec:worldmodeling}) whose predictive feedback $f$ takes four possibly overlapping forms, namely (a) future observations, (b) future states, (c) action-conditioned trajectories, and (d) predictive latents; and (iii) a readout $y=g(c,f)$ that gives an inferred relation (understanding) or a selected plan (planning). The robot scenes are illustrative, not experimental observations.}
    \label{fig:p3_s7_3_8_world_model_formulation}
\end{figure}

\subsubsection{Task Definition}
% 这里应该是定义world model如何帮助spatial reasoning task，而不是什么是world model。也就是说，world model是一个工具，spatial reasoning task是一个目标。我们要定义的是如何用这个工具去实现这个目标。

The task addressed in this subsection is spatial reasoning itself:
given a query $q$ about a scene together with the available context $c$ (a textual description, one or more observations, or a scene specification), a reasoner must produce a spatial answer $y$, either an inferred spatial property $o$ (\eg, a relation, a structure, or a physical attribute) or a plan / policy $\pi$ that acts on the scene.
What distinguishes the methods surveyed here is \emph{how} they arrive at $y$: rather than mapping $c$ to $y$ in a single feed-forward pass, the reasoner consults a \emph{world model} as an auxiliary instrument that can imagine how the scene would look or evolve, and reasons over that imagined evidence.
Concretely, the world model is a pretrained simulator exposing a distribution
$\mathbb{P}(s_0, s_1, \dots, s_T \mid a_0, a_1, \dots, a_{T-1}, c)$
over world states $s_t$, or over their partial observations $\operatorname{proj}(s_t)$ when full states are unavailable, optionally conditioned on an action sequence $\left\{a_t\right\}_{t=0}^{T-1}$ (null when the model is action-free).
A world-model-based reasoner treats $\mathbb{P}$ as a queryable tool: given the query, it (i) conditions $\mathbb{P}$ on the context $c$ and, when planning, on a candidate action sequence; (ii) draws predictive feedback $f$ from $\mathbb{P}$, a rollout of future observations, future states, or action-conditioned trajectories, or a dynamics-aware latent; and (iii) applies a readout $g$ that maps the feedback back to the spatial answer, $y = g(c, f)$.
The world model supplies the imagined evidence $f$, but producing the spatial answer $y$ remains the objective, and the fidelity of $f$ bounds how much the world model can help.





\subsubsection{Methodology}

We treat the world model as a pretrained instrument and organize methods by \emph{the form of predictive feedback it returns} to the reasoner.
Two families emerge: \emph{generative feedback}, in which the model is rolled out to synthesize explicit futures that the reasoner inspects, and \emph{latent-representation feedback}, in which the model's dynamics-aware embedding is consumed directly without decoding; the understanding-versus-planning distinction cuts across both.

\paragraph{Generative Feedback via Future Observations.}
The most direct use of a world model is to synthesize future or counterfactual \emph{pixels} and read spatial evidence off them, an analysis-by-synthesis strategy in which the reasoning outcome $o$ (a spatial relationship, structure, or physical property) is whatever best reproduces the observed or imagined imagery under $\mathbb{P}$.
Predating deep generative models, Galileo~\cite{wu2015galileo} instantiates this idea with a physics engine in place of a neural simulator: it tracks object velocities across raw videos and recovers 3D shape, mass, and friction as the parameters that best reproduce the observed motion, so that the projected state $\operatorname{proj}(s_t)$ is the object velocity $v_t$ and the videos themselves define the transition distribution $\mathbb{P}(v_0, \dots, v_T \mid c)$.
Modern instances replace the engine with an Internet-scale image or video diffusion prior: DragAPart~\cite{li2024dragapart}, Puppet-Master~\cite{li2024puppetmaster}, ATOP~\cite{vora2025articulate}, and DreamArt~\cite{lu2025dreamart} fine-tune such priors into articulation simulators $\mathbb{P}(I_0, \dots, I_T \mid c)$ whose frames $I_t = \operatorname{proj}(s_t)$ render an articulated object under a hypothesized configuration, and recover the 3D articulation (\eg, its motion axes) by aligning the rendered hypothesis with the synthesized video.
A second, inference-time use synthesizes \emph{unobserved viewpoints} to resolve spatial questions that a single view underdetermines: MindJourney~\cite{yang2026mindjourney} lets a vision-language model ``imagine'' the scene from a new camera pose with a view-conditioned world model before answering, turning the world model into a test-time reasoning aid rather than a training-time simulator.
Because such imagination is expensive and not always helpful, AVIC~\cite{yu2026avic} learns \emph{when} and \emph{how much} to imagine, gating world-model invocation on the sufficiency of the current visual evidence and matching or exceeding fixed-imagination strategies at a fraction of the world-model calls.
Pushing imagination from a fixed sweep to an \emph{active} policy, SpatialDreamer~\cite{cao2025spatialdreamer} closes the loop: a reinforcement-learned agent chooses which viewpoints to imagine through the world model and then grounds its answer in the imagined evidence, under a geometric policy-optimization objective that rewards geometrically consistent rollouts.

\paragraph{Generative Feedback via Future States.}
Rather than decode pixels, a second line rolls the world model out in a \emph{structured or higher-dimensional state space}, trading photorealism for geometric consistency.
TesserAct~\cite{zhen2025tesseract} learns a 4D (3D-plus-time) world model whose rollouts stay spatially coherent enough to support robotic manipulation, and persistent-3D world models~\cite{zhou2025learning} maintain an explicit 3D memory so that revisited regions remain consistent across long horizons.
State-space feedback is attractive when the downstream query is itself geometric, since it sidesteps the appearance-modeling burden that pixel rollouts carry.

\paragraph{Generative Feedback via Action-Conditioned Trajectories.}
When the goal is to \emph{act}, the world model is queried as a simulator: candidate action sequences are rolled out and scored by the favorability of the resulting states, yielding a policy $\pi:(s_{\leq t}, a_{<t}) \mapsto a_t$ that steers the episode toward a preferred terminal state $s_T \in S_{\text{good}}$ (the model-based-RL view; the broader treatment of RL with spatial inductive biases is §\ref{subsec:p3_s7_2_5}).
The paradigm originates with Ha and Schmidhuber~\cite{ha2018world}, who introduced the term ``world model'' in its current sense, training a generative recurrent model of future observations and using it as an internal simulator for policy search, and with Dreamer~\cite{hafner2019dreamer}, which learned the simulator in a compact latent space and optimized the policy purely on imagined rollouts.
Scaling this recipe, UniSim~\cite{yang2023learning} builds a video-diffusion simulator from which real-world manipulation policies can be extracted, and DreamGen~\cite{jang2025dreamgen} uses video world models as trajectory generators, pseudo-labeling the synthesized rollouts with an inverse-dynamics model to train policies that generalize to new behaviors and environments.
Most recently, VAGEN~\cite{wang2025vagen} turns this feedback back onto the reasoner itself, using reinforcement learning to \emph{supervise} a vision-language agent's internal world-model reasoning: it decomposes that reasoning into state estimation (``what is the current state?'') and transition modeling (``what comes next?'') and rewards accurate latent prediction.

\paragraph{Latent-Representation Feedback.}
A final family dispenses with generation entirely and consumes the world model's \emph{predictive latent} as the spatial signal.
Joint-embedding predictive architectures learn to predict the representation of a future or masked observation rather than its pixels: I-JEPA~\cite{assran2023jepa} establishes the recipe for images, and V-JEPA~2~\cite{assran2025vjepa} extends it to video, yielding dynamics-aware embeddings that support understanding, prediction, and planning without ever rendering a frame, while VL-JEPA~\cite{chen2025vl} ports the idea to vision-language inputs.
A related line keeps the latent explicitly \emph{geometric}: 3DThinker~\cite{chen20253dthinker} performs 3D mental imagery during reasoning by generating a 3D latent and aligning it with that of a 3D foundation model, letting a vision-language model reason about spatial relations from only a few views without any 3D input or explicit rendering.
Because the feedback is an embedding rather than an image, latent world models avoid the appearance-synthesis cost and the pixel-space memory drift that limit generative rollouts, at the price of a signal that is harder to inspect directly; they are an increasingly prominent, query-efficient alternative for spatial reasoning.

\subsubsection{Benchmarks}
Evaluation follows the usage mode.
Planning-oriented use is scored on the embodied suites of §\ref{subsec:p3_s7_3_2} and on long-horizon household environments such as ALFWorld~\cite{shridhar2020alfworld} and VirtualHome~\cite{puig2018virtualhome}, where the question is whether policies extracted from imagined rollouts transfer to the real environment.
Understanding-oriented use is scored on the static spatial probes that single views underdetermine, VSI-Bench~\cite{yang2025vsibench} and MindCube~\cite{yin2025spatial} among them, under test-time-scaling protocols that vary how many imagined views the reasoner is allowed and whether a verifier selects among them~\cite{jha2025probing,yu2026avic}.
The world model itself is assessed on video-generation quality (FVD, FID) together with task-specific physics-consistency probes, and surveys of general world models track the field's shift from 2D toward 3D-consistent simulation~\cite{xie2025worldsurvey}, while recent analyses argue that advancing MLLM spatial reasoning demands purpose-built recipes and data rather than scale alone~\cite{zhang2025scaling}.
No single benchmark currently exercises a world model along perception fidelity, action conditioning, and long-horizon spatial memory at once, a gap the challenges below make explicit.


\subsubsection{Challenges and Trends}
Reading the world model as a tool clarifies where the difficulties lie: some are inherited from the instrument itself, and are properly the concern of §\ref{sec:worldmodeling}, while others are specific to how a reasoner queries it. The most active questions of 2025--2026 are no longer whether a world model \emph{can} imagine, but whether its imagined evidence can be \emph{trusted} and whether consulting it is worth the cost.

Imagination is not a free lunch. Systematic probes report that world-model imagination improves spatial reasoning in only a minority of cases, and that adding imagined views yields non-monotonic, sometimes negative accuracy gains at steep token and latency cost~\cite{yu2026avic,jha2025probing}; the sheer amount of imagination is an unreliable proxy for accuracy~\cite{zhang2025scaling}, which recasts ``when and how much to imagine'' as a first-class test-time-scaling problem, with adaptive gating~\cite{yu2026avic} an early answer and a predictive account of \emph{which} queries need the world model still missing.
Faithfulness is the companion problem: a world model can render fluent futures that are physically or relationally wrong, so likelihood under $\mathbb{P}$ is necessary but not sufficient, and when the imagined evidence lacks fidelity verification buys little, since heuristic verifiers give limited calibration gains and an information bottleneck sets in, as a reasoner cannot recover a signal the imagination never contained~\cite{jha2025probing}. Grounding imagined evidence to the actual scene and calibrating trust in it remain open; physics-consistency objectives and geometry-grounded latents~\cite{chen20253dthinker} are partial responses.

The reasoner can be no more reliable than the simulator it queries. Even state-of-the-art models such as Genie 3~\cite{genie3}, despite markedly longer-horizon consistency, still drift in spatial memory, and consistency mechanisms~\cite{zhou2025learning,li2025vmem} narrow but do not close this gap, with hallucination and long-horizon error compounding over rollout length; raising this ceiling is the province of world-model construction (§\ref{sec:worldmodeling}), and the reasoning methods surveyed here can only mitigate its symptoms.

The feedback forms also trade off along different axes, with no principled account of which a given query should use: pixel rollouts are inspectable but expensive and prone to appearance drift; state or 4D rollouts are geometrically consistent but require structured simulators; latent feedback is cheap but opaque; action rollouts commit early to a policy; and most systems commit to one form a priori. Extending world-model-based reasoning beyond pixels (\eg, to forces~\cite{gillman2025force}, tactile, or audio modalities) is gated by the absence of internet-scale pre-training data in those modalities.

Two shifts define the current frontier. \emph{From external tool to trained-in reasoning}: rather than invoking a frozen simulator at inference time, recent work \emph{trains} the reasoner's use of the world model, whether by reinforcement learning that supervises an internal state-estimation and transition-modeling loop~\cite{wang2025vagen} or that rewards geometrically consistent active imagination~\cite{cao2025spatialdreamer}, or by distilling the world model's imagined dynamics directly into a vision-language model so that no simulator need be queried at test time~\cite{zhang2026world2vlm}. \emph{From pixels to predictive latents}: the imagined branch is increasingly consumed as a latent process whose hidden states guide reasoning, rather than as a faithfully rendered video~\cite{assran2023jepa,assran2025vjepa,chen20253dthinker}. Both shifts move the field away from ever-larger, faithfully-rendered pixel simulators and toward cheaper, more trustworthy, query-efficient feedback that a reasoner can afford to consult, and can be trained to consult well.









% \subsection{Summary}
% \label{subsec:p3_s7_3_summary}

% The eight subsections vary \emph{where} and \emph{how} the foundation model intervenes in the reasoning loop: as the reasoning core itself (§\ref{subsec:p3_s7_3_1}--§\ref{subsec:p3_s7_3_2}), as a generator of intermediate rationales or programs (§\ref{subsec:p3_s7_3_3}--§\ref{subsec:p3_s7_3_4}), as an orchestrator over tools, memory, and a closed interaction loop (§\ref{subsec:p3_s7_3_5}--§\ref{subsec:p3_s7_3_7}), or beside a learned simulator that supplies the spatial substrate (§\ref{subsec:p3_s7_3_8}).
% Table~\ref{tab:p3_s7_3_paradigm_comparison} summarizes the eight along the axes used for the symbolic paradigms in §\ref{sec:p3_s7_1} and the learning-based paradigms in §\ref{sec:p3_s7_2}.

% \begin{table}[t]
%     \centering
%     \fontsize{8}{10}\selectfont
%     \setlength{\tabcolsep}{2pt}
%     \caption{Comparison of the eight paradigms of foundation-model-driven spatial reasoning covered in §\ref{sec:p3_s7_3}, on the same axes as the comparison tables of §\ref{sec:p3_s7_1} and §\ref{sec:p3_s7_2}.}
%     \label{tab:p3_s7_3_paradigm_comparison}
%     \begin{tabularx}{\textwidth}{@{}p{2.9cm} p{3.3cm} p{3.3cm} X X@{}}
%         \toprule
%         \textbf{Paradigm} & \textbf{Locus of foundation-model intervention} & \textbf{Canonical inputs / outputs} & \textbf{Strengths} & \textbf{Limitations} \\
%         \midrule
%         §\ref{subsec:p3_s7_3_1} LLM/VLM-centric reasoning   & Pre-trained model \emph{is} the reasoning core; intervention via prompt / train / data / model-centric enhancement & Text-only or image$+$query $\to$ qualitative relation, distance, orientation, navigation plan & Zero-shot deployment; rapid prompt iteration; benefits from any backbone improvement                  & Ungrounded under text-only; metric-precision gap under VLM; reference-ambiguity in embodied use \\
%         §\ref{subsec:p3_s7_3_2} VLA-based reasoning         & Pre-trained model also emits low-level actions (action tokens / deltas / waypoints)                                & Image $+$ proprioception $+$ instruction $\to$ executable action sequence                            & End-to-end semantic-to-motor pipeline; cross-embodiment positive transfer via large data            & Semantic--physical misalignment; pseudo-interaction; inference-latency under autoregressive decoding \\
%         §\ref{subsec:p3_s7_3_3} Chain-of-thought (CoT \& variants) & Pre-trained model, with intermediate textual rationale as the prior; five topologies (CoT/ToT/GoT/PoT/CoA) unified as guided search & Query $\to$ rationale trace $\to$ grounded answer (qualitative relation, coordinate, plan, action) & Multi-step decomposition; verifier-friendly intermediate states; search-quality tunable               & Spatial-coordinate hallucination under simple textual CoT; rationale not grounded in truth signal; ToT/GoT cost scales super-linearly in depth \\
%         §\ref{subsec:p3_s7_3_4} Code / DSL generation         & Pre-trained model generates an executable program; environment supplies the truth signal                          & Query$+$API $\to$ program $P$ $\to$ executed output (answer with trace, or environment change)   & Interpretable rationale trace; execution-grounded; composes existing perception or control tools     & Reliance on a fixed human-authored DSL; brittle code execution under upstream-module error; sample-inefficient debug loops \\
%         §\ref{subsec:p3_s7_3_5} Tool-augmented reasoning      & Pre-trained model orchestrates a tool catalogue $\mathcal{T}$ via policy $\pi_\phi$, along execution / representation / validation / memory axes & Query$+$observation$+$tool set $\to$ history-conditioned answer $+$ tool-call trace             & Modular extension of perceptual / geometric / verification capacity; orchestration policy as the learning target & Combinatorial exploration in tool-action space; metric gap in the LLM--tool interface; conceptual overlap with §\ref{subsec:p3_s7_3_4} / §\ref{subsec:p3_s7_3_6} \\
%         §\ref{subsec:p3_s7_3_6} Retrieval- \& memory-augmented & External memory bank itself is the central learning target; Planner--Memory--Retriever triple $(\mathcal{P}, \mathcal{M}, \mathcal{R})$ & Query$+$scene observation$+$memory $\to$ action plan, retrieved subgraph, retrieved past observation & Kilometer-scale environment coverage; personalization; explainable trace of retrieved memories       & Stateless inference under unbounded context; ICL vs.\ cross-attention fusion unresolved; modality mismatch in retrieval (text vs.\ image vs.\ 3D) \\
%         §\ref{subsec:p3_s7_3_7} Agentic reasoning             & Orchestration policy $\pi$ over the perceive--reason--act--reflect cycle becomes the central learning target; tools (§\ref{subsec:p3_s7_3_5}) and memory (§\ref{subsec:p3_s7_3_6}) composed into closed-loop control & Instruction$+$observation $+$ tools $+$ memory $\to$ sensorimotor action sequence, accumulated skill library, updated world model & Closed-loop error correction; long-horizon procedural-knowledge accumulation; explicit Spatial Theory-of-Mind in multi-agent settings & Sim-to-real fidelity gap; long-horizon spatiotemporal coherence; convergence with §\ref{subsec:p3_s7_3_3}--§\ref{subsec:p3_s7_3_6} blurs the agentic boundary \\
%         §\ref{subsec:p3_s7_3_8} World-model--based reasoning  & Learned generative simulator $\mathbb{P}(s_0,\dots,s_T \mid a_0,\dots,a_{T-1}, c)$ replaces / augments the foundation model as the spatial substrate & Context $c$ ($+$ actions if conditioned) $\to$ inferred spatial properties (understanding) or policies $\pi$ (planning) & Predictive-physics reasoning; inference-time mental simulation; 4D / object-centric structural priors & World-model upper bound (Genie-style long-term memory loss); modality coverage gated by pre-training data; likelihood-vs.-physics gap (likelihood necessary but not sufficient) \\
%         \bottomrule
%     \end{tabularx}
% \end{table}

% Down the table the locus of intervention moves from \textbf{model-as-core} (§\ref{subsec:p3_s7_3_1}--§\ref{subsec:p3_s7_3_2}) through \textbf{model-as-rationale-generator} (§\ref{subsec:p3_s7_3_3}--§\ref{subsec:p3_s7_3_4}) and \textbf{model-as-orchestrator} over external components (§\ref{subsec:p3_s7_3_5}--§\ref{subsec:p3_s7_3_7}) to a \textbf{learned simulator as the reasoning substrate} (§\ref{subsec:p3_s7_3_8}), and the principal learning target migrates with it, from model weights to prompt structure to the orchestration policy $\pi$ to the world model $\mathbb{P}$ itself.
% The paradigms are cumulative rather than parallel: from §\ref{subsec:p3_s7_3_3} onward each adds one dimension (programs, tools, memory, reflection, world models) on top of the previous ones, so the ``agentic'' boundary of §\ref{subsec:p3_s7_3_7} lies where the orchestration policy itself becomes a learning target rather than where any new capability appears, which is an open research question rather than a settled definition.
% The open problem that cuts across all eight is the \emph{precision--grounding asymmetry}: the foundation models that supply semantic competence (§\ref{subsec:p3_s7_3_1}--§\ref{subsec:p3_s7_3_4}) lack the metric precision physical manipulation requires, while the external components that supply precision (§\ref{subsec:p3_s7_3_5}--§\ref{subsec:p3_s7_3_8}) inherit the LLM--tool interface as a new bottleneck.
% That asymmetry is the bridge to §\ref{sec:p3_s7_4}, where foundation-model-driven and symbolic or learning-based components are composed to close the gap from both sides at once.

% \newpage
\section{Hybrid Spatial Reasoning}
\label{sec:p3_s7_4}

\FramedInitial{C}{onventional} approaches to spatial reasoning exhibit fundamental limitations.
Systems that rely solely on linguistic or symbolic reasoning typically lack perceptual grounding in real-world visual and physical structures (\eg, geometry, physics, spatial continuity), making them inadequate for handling cross-view transformations, temporal evolution, or physically coherent scenarios.
Conversely, purely perceptual or vision-based systems, such as those combining visual recognition with rule-based geometric reasoning, often fall short in performing high-level abstract reasoning, multi-step planning, long-term memory integration, or interactive decision-making.
These limitations have motivated a growing body of work advocating for a hybrid paradigm of spatial reasoning: integrating heterogeneous modules with distinct capabilities (perception, world modeling, memory, and action) so as to jointly leverage low-level spatial and physical cues together with high-level cognitive reasoning capabilities.
A perceptual or physics-informed module supplies geometry, relations, affordances, and dynamics; a symbolic or neural module plans, infers, and abstracts over them; memory and history extend the pipeline to multi-step, multi-view, and temporally extended tasks; and because no single modality, task, or dataset carries the whole system, the composition transfers across vision-only, language-only, embodied, and multimodal settings.

% 在空间推理 (spatial reasoning) 的研究中，传统单一范式存在显著局限性。以仅依赖语言/符号推理 (symbolic reasoning) 的系统为例，它们往往缺乏对真实世界视觉／物理结构 (geometry / physics / continuity) 的感知与空间记忆，难以处理跨视角 (cross-view)、跨时 (temporal) 和物理连贯 (physical consistency) 的复杂情境；另一方面，仅依赖感知或视觉模块 (例如视觉识别 + 规则几何推理) 的系统，则可能缺乏高阶抽象推理 (abstract reasoning)、规划 (planning)、长程记忆 (long-term memory) 以及交互式决策 (interactive decision) 的能力。
% 因此，很多工作主张一种 “混合 (hybrid)”空间推理 (spatial reasoning) 的研究范式 — 即通过整合异构 (heterogeneous) 模块 (perception, world-model, reasoning, memory, action) —— 以同时兼顾空间结构感知、几何 / 物理一致性、跨视角／跨模态理解、长期规划与推理。
% 这种混合范式的优点包括但不限于：
% 感知 + 抽象推理：底层视觉 / 物理世界模型 (world-model) 提供空间结构与物体关系 (spatial relationships / affordances / dynamics)，而高层符号 /语言 /逻辑推理 (symbolic or neural reasoning) 提供规划、决策、归纳／演绎等能力；
% 长期记忆与交互能力：通过将 memory / history / context 纳入 world-model + reasoning pipeline，使系统能够处理长期任务 (multi-step, multi-view, temporally-extended)；
% 泛化与灵活性 (generalization & flexibility)：混合模块使系统不依赖单一任务或单一数据集 (\eg pure 3D supervision, or pure NL reasoning)，从而更具拓展性 (scalability) 与适应性 (adaptability)。


\subsection{Task Definition}
\label{subsec:p3-sec7-4-task-definition}
A hybrid spatial reasoning system is a pair $(\{M_i\}_{i=1}^{k}, \pi)$: a set of $k$ heterogeneous modules $M_i \in \{M_\text{percept}, M_\text{world}, M_\text{symb}, M_\text{mem}, M_\text{plan}, M_\text{act}\}$, each contributing a distinct inductive bias (perception, world model, symbolic constraints, memory, planning, or action), and an orchestration function $\pi : \mathcal{S} \to \{1, \dots, k\}$ that selects which module acts on the current system state $s_t \in \mathcal{S}$ (perception, internal map, history). Starting from the input $s_0 = x$, the system evolves as
\begin{equation}
  s_{t+1} = M_{\pi(s_t)}(s_t), \qquad y = \mathrm{Read}(s_T),
  \label{eq:p3_s7_4_hybrid}
\end{equation}
and the three composition styles found in practice are special cases of Eq.~\eqref{eq:p3_s7_4_hybrid}: a \emph{pipeline} fixes $\pi$ to a constant module order, an \emph{iterative} system lets $\pi$ revisit modules until a verifier accepts, and an \emph{end-to-end} system makes every $M_i$ differentiable so that the whole trajectory is trained jointly. Hybrid spatial reasoning is the paradigm in which $\pi$ and the module set $\{M_i\}$ are designed together to minimize the precision-grounding asymmetry: the foundation models of §\ref{sec:p3_s7_3} supply semantic competence, and the symbolic and learning-based components of §\ref{sec:p3_s7_1}--§\ref{sec:p3_s7_2} supply metric precision and verifiable structure. Existing systems therefore differ along three axes: \emph{which} modules they compose, \emph{how} they compose them (pipeline, iterative, or end-to-end), and how the orchestration function is realized (hand-coded rules or a learned policy).

\subsection{Methodology}
\label{subsec:p3-sec7-4-methodology}

Five hybridization patterns recur in the literature, distinguished by which module the foundation model is paired with and by what supplies the truth signal.

\begin{figure}
    \centering
    \includegraphics[width=\linewidth]{figures/part3/section7/foundation_rl_hybrid.pdf}
    \caption{The foundation-model + reinforcement-learning hybrid pattern: an LLM/VLM hypothesis generator coupled with an RL controller and with a physics engine or symbolic solver for verification. The RL controller's policy (written $\pi$ in the figure) is $\pi_{\mathrm{RL}}$ in the text; the bottom bar is the orchestration function $\pi$ of Eq.~\eqref{eq:p3_s7_4_hybrid}.}
    \label{fig:p3_s7_4_foundation_rl_hybrid}
\end{figure}

\paragraph{Foundation Models with Reinforcement Learning.}
A prominent hybrid couples a foundation model's semantic competence with reinforcement learning's long-horizon decision-making, often closed by a verifier that supplies executable structure (Figure~\ref{fig:p3_s7_4_foundation_rl_hybrid}).
The foundation model proposes hypotheses, sub-goals, or candidate action sequences; an RL policy $\pi_{\mathrm{RL}}$ commits to and executes them over extended horizons; and the resulting reward, optionally augmented by a physics engine or symbolic checker, updates both the policy and the orchestration.
Embodied-R~\cite{zhao2025embodied} instantiates this division of labor: a large vision--language model handles perception, while a small language model trained by RL, with a reward that scores the logical consistency between its reasoning trace and its answer, performs the slow, deliberative spatial reasoning, evaluated on UrbanVideo-Bench and VSI-Bench.
The group-relative fine-tuning recipes of §\ref{subsec:p3_s7_2_5} (Reason-RFT, OctoNav, SpaceR~\cite{tan2025reasonrft,gao2026octonav,ouyang2025spacer}) are the same hybrid with the verifier folded into a verifiable reward and are not repeated here.
The foundation model contributes broad priors and RL supplies credit assignment and exploration; the cost is the sample inefficiency and reward-specification difficulty inherited from RL (§\ref{subsec:p3_s7_2_5}).

\paragraph{Foundation Models with Explicit Physical Simulation.}
A second pattern pairs the foundation model with a physics engine, so that the truth signal is a simulated trajectory rather than a learned score; the engine side is the physical reasoning of §\ref{subsec:p3_s7_1_2}.
The LLM proposes hypotheses, structural descriptions, or programs, the simulator executes them under conservation laws, kinematic constraints, contact dynamics, and reachability limits, and the discrepancy between simulation and observation is fed back for another round.
LLMPhy~\cite{cherian2024llmphy} is the clearest instance: the LLM writes Python programs that embed estimates of latent physical parameters (friction, stiffness, mass) and of the scene layout, a physics engine executes them, and the reconstruction error against the observed video is returned as feedback, so that building the digital twin becomes a black-box optimization driven by the LLM.
DiffGen~\cite{jin2024diffgen} makes the loop differentiable, generating robot demonstrations by descending gradients through a differentiable simulator, a differentiable renderer, and a vision--language model until the rendered outcome matches the language instruction in embedding space.
PhysVLM~\cite{zhou2025physvlm} moves the physics into the input instead: a robot's reachable workspace is abstracted into a robot-agnostic space-physical reachability map that the VLM consumes alongside the image, trained on Phys100K and evaluated on the EQA-phys benchmark.
At the far end, ChronoLLM~\cite{wang2025chronollm} fine-tunes LLMs to write PyChrono multibody-simulation scripts, from a single pendulum to a full vehicle on deformable terrain; here the model authors the simulation rather than reasoning inside it, which places it at the boundary with the code-generation paradigm of §\ref{subsec:p3_s7_3_4}.
What the simulator contributes is an inductive bias that stays valid for unseen objects and configurations, so the hybrid extrapolates where a purely statistical model would not.

\paragraph{Foundation Models with Symbolic Reasoning.}
A third pattern inserts an explicit symbolic layer (logical constraints, relational predicates, or formal planning specifications) between natural-language reasoning and executable spatial inference, replacing the hand-crafted parsers of §\ref{sec:p3_s7_1} with foundation-model-generated symbolic structure.
Two sub-patterns differ in when the symbolic machinery is used.
In the first it is called at inference time: NS3D~\cite{hsu2023ns3d} has an LLM translate a 3D referring expression into a first-order logic program that a neural executor grounds against the scene's objects and relations on ReferIt3D, and LEFT~\cite{hsu2023s} (§\ref{subsec:p3_s7_3_4}) makes the same programs differentiable so that concept grounding is learned rather than hand-specified.
In the second, symbolic knowledge enters only at training time as a constraint on a neural model: \citet{premsri2024neurosymbolic} train language models for reasoning over spatial language with spatial logic rules as auxiliary losses, and RelateSeg~\cite{lin2025neuro} extracts relations such as $\langle\text{cat}, \text{to-right-of}, \text{person}\rangle$ from open-vocabulary segmentation outputs, writes them as first-order logic, and relaxes them with fuzzy logic into a single auxiliary loss that improves mIoU on four benchmarks without adding parameters.
The Scene Language~\cite{zhang2025scenelanguage} (§\ref{subsec:p3_s8_2_5}) is the representation such hybrids increasingly target: programs for structure, words for semantics, embeddings for appearance.
The symbolic layer buys consistency and compositional guarantees; its price is a closed predicate vocabulary that the neural side constantly exceeds, a point taken up under Challenges and Trends below.

\paragraph{Foundation Models with Spatial Memory and World Models.}
A fourth pattern gives the foundation model a persistent spatial state that it does not have to carry in its context window: an explicit memory (episodic trajectories, topological maps, 3D layouts) or a learned world model that predicts how the scene changes under action.
The hybrid question is the interface between the two.
MSNav~\cite{liu2025msnav} splits zero-shot vision-and-language navigation into a memory module, a spatial module, and a decision module, with a small fine-tuned LLM extracting the landmarks and relations the spatial module needs, and reaches state-of-the-art zero-shot results on R2R and REVERIE.
MindExplore~\cite{li2025towardsvla} stacks a reasoning layer (task-specific chains of thought that emit meta-actions) over an acting layer (a mixture of policy experts with a multimodal diffusion policy) and uses a memory mechanism as the feedback channel between the two, running at 30~FPS in dynamic sand environments.
The 3D spatial memories of §\ref{subsec:p3_s7_3_6} (3DSPMR~\cite{cai2025vision}, Embodied-RAG~\cite{xie2024embodiedrag}, Meta-Memory~\cite{mao2025meta}) and the world-model coupling of §\ref{subsec:p3_s7_3_8} are the memory- and simulator-side components these hybrids draw on; MindJourney's test-time gain on SAT~\cite{yang2026mindjourney} (about eight points of top-1 accuracy from imagined viewpoints alone, with the VLM untouched) is the cleanest quantitative evidence that a world model used as a tool closes part of the gap.
In each case the foundation model keeps high-level inference, instruction following, and semantic grounding, and offloads persistence, geometry, and dynamics to the dedicated module.

\begin{figure}[!ht]
    \centering
    \includegraphics[width=0.8\linewidth]{figures/part3/section7/agentic_hybrid_orchestration.pdf}
    \caption{Illustration of the agentic hybrid orchestration pattern: a foundation-model controller dynamically selects perception, world-modeling, symbolic-reasoning, memory, and action tools within a closed-loop perceive--reason--act--reflect cycle.}
    \label{fig:p3_s7_4_agentic_hybrid}
\end{figure}

\paragraph{Agentic Orchestration of Hybrid Components.}
The fifth pattern makes the orchestration function itself the object of design: a foundation model acts as a deliberative controller that decides, at each step, which of the perception, world-model, symbolic, memory, and action modules to invoke, within a closed perceive--reason--act--reflect loop (Figure~\ref{fig:p3_s7_4_agentic_hybrid}).
It differs from the agentic reasoning of §\ref{subsec:p3_s7_3_7} in what is being scheduled: there, tools and memory; here, additionally symbolic engines and physics simulators, with the scheduler $\pi$ of Eq.~\eqref{eq:p3_s7_4_hybrid} itself a candidate for learning.
VADAR~\cite{marsili2025visual} (§\ref{subsec:p3_s7_3_4}) is the program-synthesis version of such a controller, generating the API it then calls.
``See it. Say it. Sorted''~\cite{zhang2025seeitsayit} is a training-free version for sketch-to-diagram generation: a critic VLM proposes a few qualitative relational edits, several LLMs synthesize candidate SVG programs, and a judge VLM selects among them, which keeps global constraints such as alignment and connectivity intact where numerical estimation would drift.
\citet{manh2025mindmeetspace} arrive at the same module inventory from neuroscience, arguing that such a controller should sit over multimodal sensing, multisensory integration, egocentric--allocentric conversion, a cognitive map, spatial memory, and spatial reasoning.
The gain is compositionality and long-horizon consistency that no static architecture provides; the cost is the latency and attribution problems discussed below.

% 之后
% 具体地，今年来，研究者越来越倾向于将大型基础模型 (foundation models) 与强化学习 (reinforcement learning, RL) 结合，以构建具备 spatial reasoning 能力的混合系统。
% 例如，最新工作 Embodied-R: Collaborative Framework for Activating Embodied Spatial Reasoning in Foundation Models via Reinforcement Learning (简称 Embodied-R) 就是一个典型代表。该工作将大型视觉—语言模型 (VLM) 用于感知 (visual observation) —— 提取视觉特征与空间线索；同时用较小的语言模型 (LM) 作为推理 /决策模块 (reasoner / planner)，并通过强化学习训练 (设计 reward 考量 think-answer 逻辑一致性) 实现 “slow-thinking” (慢思考) 的空间推理。
% 这些工作表明：将基础模型 (VLM / LLM) 与 world-model + RL 结合，是当前实现混合空间推理的最有前景路径之一 — 尤其适合应对复杂、动态、物理 /几何约束强的现实场景 (例如 3D 布局、具身导航、物体交互、场景理解等)。
% 此外，从更理论／宏观层面来看，有分析指出当前基础模型 (foundation models) 在决策 (decision-making) 与空间 /世界智能 (world-model) 方面存在结构性空缺。因此，通过 RL、world-model、模块化 /混合架构 (hybrid architecture) 来弥补这一空缺，有助于未来构建真正具有空间智能 (spatial intelligence) 的通用系统 (general-purpose agents / embodied AI / spatially-aware AI)。

%   通过将混合架构纳入智能体 (agent) 或规划 + 行动 (planning + action) 回路 (loop) 中，我们可以构造更为通用、灵活、可扩展 (scalable) 的空间智能系统。这种 “Agentic Hybrid Orchestration” 的设计核心包括以下几个关键维度：
% LLM / VLM 作为 Orchestrator (协调者／指挥者)：高层语言 /符号 /规划模块 (LLM) 负责决定 “何时感知 (perceive)”、 “如何构建或更新 world-model (scene graph / layout / spatial representation)”、 “何时进行推理 /规划 (plan)” 以及 “何时发起行动 /布局 /生成 / manipulation”。
% 工具选择 (tool selection) 与顺序 (sequencing)：系统可以预定义或动态维护一组 “工具 (tools)” —— 比如视觉感知模块、world-model builder、符号 /语言推理模块、布局 /渲染 模块 —— LLM 负责根据当前任务与上下文选择合适工具并安排调用顺序。
% 反馈驱动 (feedback-driven) 的迭代 (iterative) 推理 (reasoning)：通过 RL (或其他优化方式) 将最终目标 (task success, spatial consistency, physical plausibility, layout realism / utility) 的反馈 (reward / loss) 传播回 orchestrator 和各工具调用决策。
% 这种 Agentic Hybrid Orchestration 的好处是，它将混合架构系统转化为一个动态、交互式、多轮 (multi-step) 的代理 (agent) 回路 (loop)，从而具备长期规划、空间结构一致性、物理可行性、泛化 与适应性 — 正是混合空间推理 (Hybrid Spatial Reasoning) 的目标。

\subsection{Benchmarks}
\label{subsec:p3-sec7-4-benchmarks}
Hybrid systems are evaluated end to end on the benchmark of their target task rather than module by module, and the benchmark is often one the paper introduces.
Among the systems above, Embodied-R reports on UrbanVideo-Bench and VSI-Bench, MindJourney on SAT, NS3D on ReferIt3D, MSNav on R2R and REVERIE, and PhysVLM on RoboVQA and OpenEQA alongside its own EQA-phys; LLMPhy ships TraySim, CLEVRER-LLMPhy, and Real-TraySim, MindExplore ships SandGo-1k and SandThink-21k, and DrawingBench~\cite{kim2025drawingbench} is built specifically for agentic hybrids, scoring an LLM's sequence of low-level mouse actions against eight rule-based criteria for the spatial configuration it was asked to draw.
Physics-side components are still checked against PHYRE and ShapeStacks (§\ref{subsec:p3_s7_1_2}) and symbolic components against the constraint- and rule-based suites of §\ref{subsec:p3_s7_1_4}--§\ref{subsec:p3_s7_1_5}, but no current benchmark attributes an end-to-end score to the modules that produced it; that gap is taken up under Challenges and Trends below and in the cross-paradigm survey of §\ref{sec:p3_s7_5} (Table~\ref{tab:benchmark_comparison}, whose 15-row agentic subset is Table~\ref{tab:spatial_reasoning_benchmark} of §\ref{subsec:p3_s7_3_7}).

\subsection{Challenges and Trends}
\label{subsec:p3-sec7-4-challenges-and-trends}
Hybrid spatial reasoning has produced architectures that combine complementary inductive biases in principled ways, but several recurring limitations originate in the module-composition framework itself rather than in any individual module.

\textbf{Frame, Scale, and Geometric Consistency.}
Unlike generic modular systems, spatial hybrids must reconcile heterogeneous \emph{spatial} representations across modules: perception emits camera or egocentric frames, map and memory modules operate in world or allocentric coordinates, symbolic reasoners abstract to topological relations, and physics engines expect metric units. Misaligned reference frames, scales, or conventions silently corrupt downstream reasoning, and because each module reports confidence in its own terms, small upstream pose, depth, or scale errors propagate and compound across the pipeline without a coherent notion of uncertainty. Maintaining frame, scale, and geometric consistency is thus a first-order obstacle with no analogue in non-spatial hybrid systems.

\textbf{Interfaces, Latency, and Attribution.}
The strength of a hybrid system is bounded by its worst module interface. Symbolic-to-neural translation is particularly brittle when the symbolic side uses closed-vocabulary predicates while the neural side proposes open-vocabulary ones, and the interface design becomes a hidden first-class hyperparameter that current literature seldom isolates. Each module invocation also adds wall-clock latency that compounds across a perceive--reason--act--reflect loop, so high-frequency manipulation and navigation may not be reachable through orchestration alone. And when a hybrid fails, attributing the failure to a specific module is non-trivial; this complicates iteration and is often cited as a practical reason why end-to-end models can outperform hybrid stacks on aggregate benchmarks despite offering weaker per-module guarantees.

The orchestration function $\pi$ is migrating from hand-coded rules to \emph{learnable orchestration policies} trained with reinforcement or imitation learning, which blurs the boundary with §\ref{subsec:p3_s7_3_7}.
\emph{Differentiable hybrid stacks}, built from differentiable physics simulators, differentiable optimization layers (§\ref{subsec:p3_s7_2_4}), and differentiable retrieval, make end-to-end training of a hybrid viable and remove the historical interface penalty, offering a path through the precision--grounding asymmetry that motivates this section.
Evaluation is changing too, toward the \emph{holistic and diagnostic}: in response to the attribution problem, emerging practice pairs aggregate task metrics with per-module diagnostic probes and counterfactual module-swap ablations that isolate where a hybrid system succeeds or fails, a direction consolidated in the cross-paradigm benchmark survey of §\ref{sec:p3_s7_5}.

% \newpage
z\section{Spatial Reasoning Benchmarks}
\label{sec:p3_s7_5}

\FramedInitial{E}{arly} spatial-reasoning benchmarks positioned the model as a detached observer. CLEVR~\cite{johnson2017clevr} and VSR~\cite{liu2023visual} ask for relative geometric relations from a static third-person viewpoint, and their 2D inputs leave depth ambiguous, a weakness that 3DSRBench~\cite{ma20253dsrbench}, MMSI-Bench~\cite{yang2025mmsi}, DSI-Bench~\cite{zhang2025dsi}, and STI-Bench~\cite{li2025sti} now probe directly.
SQA3D~\cite{ma2022sqa3d}, OpenEQA~\cite{majumdar2024openeqa}, and EmbSpatial-Bench~\cite{du2024embspatial} instead place the model inside the scene, so that it must self-localize before it can answer, and ScanQA~\cite{azuma2022scanqa}, ScanRefer~\cite{chen2020scanrefer}, and Multi3DRefer~\cite{zhang2023multi3drefer} ground queries in point clouds and RGB-D data, where the answer is a coordinate rather than a relation.
The reasoning demanded has deepened in parallel: IntPhys~\cite{riochet2018intphys} and Physion~\cite{Bear2021Physion} test whether a model expects gravity, collision, and friction to hold over time, and OpenEQA and RealWorldQA~\cite{xai2024grok15v} ask functional questions (``Can I navigate through this gap?'', ``Where would I find a cloth to clean a spill?'') in which geometry matters only through what it affords.

Data provenance splits the field. Simulators, the sources behind SpatialVLM~\cite{chen2024spatialvlm}, AI2-THOR~\cite{kolve2017ai2thor}, and Habitat~\cite{szot2021habitat2}, yield unlimited, perfectly annotated data; real-world collections such as ScanNet~\cite{dai2017scannet} and OpenEQA supply the clutter, occlusion, and sensor noise that simulation does not reproduce.
The resulting practice is sim-to-real: pre-train on large synthetic corpora for general spatial logic, then fine-tune on the scarcer real data.
Annotated 3D data remain scarce and fine-grained metric estimation remains hard for large models; both limits recur throughout the table below.

% \wsq{总结现有的有关spatial reasoning 相关的benchmarks}


\subsection{Comparative Analysis}
\label{subsec:p3-sec7-5-comparative-analysis}

Table~\ref{tab:benchmark_comparison} compiles 61 spatial-reasoning benchmarks across four categories (physical-focused, rooted in §\ref{subsec:p3_s7_1_2}; vision-based, overlapping with §\ref{subsec:p3_s7_3_1}; 3D-focused, rooted in §\ref{subsec:p3_s7_3_6} and §\ref{subsec:p3_s7_3_8}; and robot-environment, rooted in §\ref{subsec:p3_s7_3_7} and §\ref{sec:p3_s7_4}).
A narrower 15-row subset focused specifically on agentic spatial reasoning is provided in Table~\ref{tab:spatial_reasoning_benchmark} of §\ref{subsec:p3_s7_3_7}; this subset filters Table~\ref{tab:benchmark_comparison} to image-and-video benchmarks that exercise agentic deployments.

\begin{table*}[!t]
  \caption{Spatial-reasoning benchmarks surveyed in this chapter, grouped by what they test. \texttt{Number} is the evaluation-set size where the source paper documents one (train / test when both are given, scene counts for simulators, $\sim$ for approximate figures). \texttt{Perspective} marks explicitly designed ego (first-person) and exo (third-person) settings; \textit{Undeclared} marks attributes the source paper does not state. \texttt{Data Source} and \texttt{Task Type} follow the source papers.}
  \label{tab:benchmark_comparison}
  \centering
  \fontsize{9}{10}\selectfont
  \setlength{\tabcolsep}{1.4mm}
  \resizebox{\linewidth}{!}{
  \begin{tabular}{@{}lcccccccc}
    \toprule
      \multirow{2}{*}{\textbf{Benchmark}}
      & \multirow{2}{*}{\textbf{Number}}
      & \multirow{2}{*}{\textbf{Modality}}
      & \multicolumn{2}{c}{\textbf{View Type}}
      & \multicolumn{2}{c}{\textbf{Perspective}}
      & \multirow{2}{*}{\textbf{Data Source}}
      & \multirow{2}{*}{\textbf{Task Type}} \\
      \cmidrule(lr){4-5}
        \cmidrule(lr){6-7}
        & &
        & \textbf{Single} & \textbf{Multi}
        & \textbf{Ego} & \textbf{Exo}
        & \\
    \midrule
    \rowcolor{lightgrey}\multicolumn{9}{c}{$\bullet$ \textit{Physical Focused Spatial Reasoning}} \\
    {IntPhys}~\cite{riochet2018intphys} & 15k / 4.3k & 3D   & \checkmark & \crossmark    & \crossmark & \checkmark   & Synthetic    & VoE plausibility\\
    {IntPhys 2}~\cite{bordes2025intphys2}& 1,416    & 3D   & \checkmark & \crossmark    & \crossmark & \checkmark       & Synthetic    & VoE plausibility\\
    {Physion}~\cite{Bear2021Physion}  & $\sim$1.2k  & 3D    & \checkmark & \crossmark  & \crossmark & \checkmark   & Synthetic & Physical prediction / choice \\
    {Physion$++$}~\cite{tung2023physion} & $\sim$10k  & 3D  & \checkmark & \crossmark   & \crossmark & \checkmark  & Synthetic  & Physical prediction / choice\\
    {Physical Concepts Dataset}~\cite{piloto2020learn}  & 300k    & 3D   & \checkmark & \crossmark   & \crossmark & \checkmark  & Synthetic   & classification\\
    {Physics-IQ}~\cite{motamed2025generative} & 396  & Video     & \checkmark & \crossmark   & \crossmark & \checkmark  & Real   & Generative physical prediction \\
    {SPhyR}~\cite{siedler2025sphyr} & 1,296 & Image & \checkmark & \crossmark & \multicolumn{2}{c}{\cellcolor{lightgray}\textit{Undeclared}} & Synthetic & Physical prediction / infill \\
    \hline
    \rowcolor{lightgrey}\multicolumn{9}{c}{$\bullet$ \textit{Vision-Based Spatial Reasoning}} \\
    CLEVR \citep{johnson2017clevr} & 853k & Image & \checkmark & \crossmark & \crossmark & \checkmark & Simulation & VQA \\
    VSR \citep{liu2023visual}& 10,972 & Image & \checkmark & \crossmark & \crossmark & \checkmark & Real & VQA \\
    What's Up \citep{KamathHC23a}  & $\sim$8,000 & Image & \checkmark & \crossmark & \crossmark & \checkmark & Real & VQA\\
    RocketScience~\citep{hoehing2025RocketScience} & 482 & Image & \checkmark & \crossmark  & \crossmark & \checkmark & Real & VQA \\
    SpaRE \citep{abs-2504-20648} & 3.4M & Image &  \checkmark & \crossmark & \crossmark & \checkmark & Real/Sim & VQA \\
    RealWorldQA~\cite{xai2024grok15v} & 765 & Image & \checkmark & \crossmark & \checkmark & \crossmark & Real & VQA \\
    3DSRBench~\cite{ma20253dsrbench} & 2,772 & Image & \checkmark  & \crossmark & \crossmark  &  \checkmark & Real/Sim & VQA\\
    InternSpatial-Bench~\cite{deng2025internspatial} &  6,008 & Image  & \checkmark  & \crossmark & \crossmark  &  \checkmark  & Real & VQA \\
    OmniSpatial~\cite{jia2025omnispatial} & $\sim$8,400 & Image & \checkmark & \crossmark & \checkmark & \checkmark   &  Real & VQA\\
    Spatial457~\cite{wang2025spatial457} & 23,752 & Image & \checkmark & \crossmark & \multicolumn{2}{c}{\cellcolor{lightgray}\textit{Undeclared}}  & Real & VQA\\
    All-Angles Bench~\cite{yeh2025seeing} & $\sim$2,100 & Image  & \crossmark & \checkmark & \multicolumn{2}{c}{\cellcolor{lightgray}\textit{Undeclared}} & Real  & VQA \\
    EmbSpatial-Bench~\cite{du2024embspatial} & 3,640 & Image & \checkmark  & \crossmark & \checkmark & \crossmark &  Real & VQA \\
    MMSI-Bench~\cite{yang2025mmsi} & 1,000 & Image & \crossmark & \checkmark  & \multicolumn{2}{c}{\cellcolor{lightgray}\textit{Undeclared}} &  Real & VQA\\
    MindCube~\cite{yin2025spatial} & 21,154 & Image & \crossmark & \checkmark & \multicolumn{2}{c}{\cellcolor{lightgray}\textit{Undeclared}}  & Simulation  & VQA \\
    BLINK~\cite{fu2024blink} & 3,807 & Image & \checkmark & \crossmark & \crossmark & \checkmark & Real & VQA \\
    NuScenes-SpatialQA~\cite{tian2025nuscenesspatialqa} & 3.3M+ & Video & \checkmark  & \checkmark  & \checkmark  & \checkmark  & Real & VQA \\
    VSI-Bench~\cite{yang2025vsibench} & 5,130 & Video & \checkmark & \crossmark & \checkmark & \crossmark & Real  & VQA\\
    VLM4D~\cite{zhou2025vlm4d} & $\sim$1,800 & Video & \checkmark & \crossmark &  \checkmark & \checkmark &  Real & VQA \\
    STI-Bench~\cite{li2025sti} & 2,064 & Video  & \checkmark & \crossmark & \checkmark & \checkmark & Real & VQA\\
    DSI-Bench~\cite{zhang2025dsi} & $\sim$1,700 & Video  & \checkmark & \crossmark & \multicolumn{2}{c}{\cellcolor{lightgray}\textit{Undeclared}} & Real & VQA\\
    SITE~\cite{wang2025site} & 8,068  & Video & \checkmark & \crossmark & \multicolumn{2}{c}{\cellcolor{lightgray}\textit{Undeclared}} & Real/Sim & VQA \\
    RoboSpatial \citep{SongBTTSB25} & 3M & Image/3D & \checkmark & \crossmark & \checkmark & \crossmark & Real/Sim & VQA \\
    OSD \citep{cheng2024spatialrgpt} & \textit{Undeclared} & Image & \checkmark & \crossmark & \crossmark & \checkmark & Real & VQA \\
    SpatialVLM \citep{chen2024spatialvlm} & 2B & Image & \checkmark & \crossmark & \crossmark & \checkmark & Real/Sim & VQA \\
    ProVision \citep{abs-2412-07012} & 65k & Image & \checkmark & \crossmark & \crossmark & \checkmark & Simulation & VQA \\
    WhatIsWhere \citep{LiangZZZ16} & 6,000 & Image & \checkmark & \crossmark & \crossmark & \checkmark & Real & VQA \\
    SURDS~\cite{guo2024surds} & 9,250 & Image & \checkmark & \crossmark & \checkmark & \crossmark & Real & VQA \\
    EarthSpatialBench~\cite{xu2026earthspatial} & 325k+ & Image & \checkmark & \crossmark & \crossmark & \checkmark & Real & VQA \\
    PureSpace~\cite{li2026purespace} & 1,200 & Image & \checkmark & \crossmark & \multicolumn{2}{c}{\cellcolor{lightgray}\textit{Undeclared}} & Simulation & VQA \\
    SAT~\cite{ray2024sat} & 175k / 150 & Image & \crossmark & \checkmark & \checkmark & \checkmark & Simulation & VQA \\
    SpatiaLQA~\cite{xie2026spatialqa} & 9,605 & Image & \checkmark & \crossmark & \crossmark & \checkmark & Real & VQA \\
    Q-Spatial++~\cite{wang2025svqa} & 101 & Image & \checkmark & \crossmark & \crossmark & \checkmark & Real & VQA \\
    \hline
    \rowcolor{lightgrey}\multicolumn{9}{c}{$\bullet$ \textit{3D Focused Spatial Reasoning}} \\
    ScanQA~\cite{azuma2022scanqa} & 41,363 & 3D & \crossmark & \checkmark & \crossmark & \checkmark & Real (ScanNet) & 3D VQA \\
    SQA3D~\cite{ma2022sqa3d} & 33,460 & 3D & \checkmark & \crossmark & \checkmark & \crossmark & Real (ScanNet) & 3D VQA \\
    MMScan QA~\cite{lyu2024mmscan} &  295K & 3D+Image & \checkmark & \checkmark & \crossmark & \checkmark & Real (EmbodiedScan) & 3D VQA \\
    OpenEQA~\cite{majumdar2024openeqa} & 1,632 & Video/3D & \checkmark & \crossmark & \checkmark & \crossmark & Real (HM3D, ScanNet) & Embodied QA (EM / A splits) \\
    Space3D-Bench~\cite{szymanska2024space3dbench} & 1,000 & 3D & \checkmark & \crossmark & \checkmark & \crossmark  & Real(Replica) & 3D QA \\
    Scan2Cap~\cite{chen2021scan2cap} & 39,775 & 3D & \crossmark & \checkmark & \crossmark & \checkmark & Real (ScanNet) & 3D Captioning \\
    MMScan Captioning~\cite{lyu2024mmscan} & 172k & 3D+Image & \checkmark & \checkmark & \crossmark & \checkmark & Real (EmbodiedScan) & 3D Captioning \\
    ScanRefer~\cite{chen2020scanrefer} & 51,583 & 3D & \crossmark & \checkmark & \crossmark & \checkmark & Real (ScanNet) & 3D Grounding \\
    Multi3DRefer~\cite{zhang2023multi3drefer} & 42,735 & 3D & \crossmark & \checkmark & \crossmark & \checkmark & Real (ScanNet) & 3D Grounding \\
    SR3D~\cite{achlioptas2020referit3d} & 83,572 & 3D & \crossmark & \checkmark & \crossmark & \checkmark & Real (ScanNet) & 3D Grounding \\
    NR3D~\cite{achlioptas2020referit3d} & 41,503 & 3D & \crossmark & \checkmark & \crossmark & \checkmark & Real (ScanNet) & 3D Grounding \\
    SpaceEra++~\cite{guan2026spaceera} & 925k & Video/3D & \crossmark & \checkmark & \checkmark & \crossmark & Synthetic & 3D VQA \\
    Ego3D-Bench~\cite{gholami2025ego3d} & 8,600+ & Image & \crossmark & \checkmark & \checkmark & \crossmark & Real & VQA \\
    SpatialScore~\cite{wu2025spatialscore} & 28,093 (Hard: 1,400) & Image/3D & \crossmark & \checkmark & \multicolumn{2}{c}{\cellcolor{lightgray}\textit{Undeclared}} & Real/Sim & VQA \\
    iVISPAR~\cite{mayer2025ivispar} & 300 & 3D Env & \checkmark & \crossmark & \crossmark & \checkmark & Simulation & Interactive reasoning \\
    \hline
    \rowcolor{lightgrey}\multicolumn{9}{c}{$\bullet$ \textit{Robot}} \\
    AI2-THOR~\cite{kolve2017ai2thor} & 120 (Scenes) & 3D Env & \crossmark & \checkmark & \checkmark & \crossmark & Simulation & Nav/Manip \\
    TDW~\cite{gan2020threedworld} & \textit{Infinite} & 3D Env & \crossmark & \checkmark & \checkmark & \checkmark & Simulation & Physical Sim \\
    Habitat 2.0~\cite{szot2021habitat2} & 100 (Scenes) & 3D Env & \crossmark & \checkmark & \checkmark & \crossmark & Simulation & Navigation \\
    iGibson 2.0~\cite{Li2021iGibson2} & 15 (Scenes) & 3D Env & \crossmark & \checkmark & \checkmark & \crossmark & Simulation & Interactive Nav \\
    Habitat 3.0~\cite{puig2023habitat3} & 211 (Scenes) & 3D Env & \crossmark & \checkmark & \checkmark & \crossmark & Simulation & Social Nav \\
    R2R~\cite{anderson2018vision} & 21,567 & Image/Path & \crossmark & \checkmark & \checkmark & \crossmark & Real & Navigation \\
    RxR~\cite{ku2020roomacrossroom} & 126,069 & Image/Path & \crossmark & \checkmark & \checkmark & \crossmark & Real & Navigation \\
    \bottomrule
  \end{tabular}
  }
\end{table*}

\emph{A modality migration} runs from the top of the table to the bottom: the physical-focused benchmarks rely on synthetic 3D renderings, the vision-based block is image- and video-heavy with growing 3D extensions, the 3D-focused block is built on point clouds and RGB-D scans, and the robot block has moved largely to 3D environments.
\emph{The Perspective axis is imbalanced}: a large fraction of vision-based benchmarks carry \textit{Undeclared} ego/exo annotations, the physical-focused benchmarks lean exo (a third-person observer judging physical plausibility), and the robot block is entirely ego; the field is migrating from spectator-style observation to situated cognition faster than benchmark annotations record it.
\emph{The Data-Source spectrum is bifurcated} between predominantly synthetic or simulated (top block and robot block) and predominantly real (middle blocks), with the Real/Sim benchmarks (SpaRE, 3DSRBench, RoboSpatial, SpatialVLM, SITE) as the bridge that the sim-to-real training practice of the opening relies on.

\subsection{Future Directions}
\label{subsec:p3-sec7-5-future-directions}

\paragraph{Unified Spatial-Cognition Benchmark.}
No single benchmark in Table~\ref{tab:benchmark_comparison} exercises a model along all of perception fidelity, symbolic precision, memory persistence, and long-horizon action conditioning simultaneously. 
A unified benchmark covering all four §\ref{sec:reasoning} paradigms (§\ref{sec:p3_s7_1} rule-based, §\ref{sec:p3_s7_2} learning-based, §\ref{sec:p3_s7_3} foundation-model-driven, §\ref{sec:p3_s7_4} hybrid) would close the evaluation-attribution challenge flagged under Challenges and Trends in §\ref{sec:p3_s7_4}.

\paragraph{Cross-Modal Evaluation.}
The current axes (Modality, View, Perspective, Data Source, Task Type) under-specify the cross-modal demand of contemporary spatial reasoning. 
Future benchmarks should explicitly probe modality-to-modality transfer (\eg, 2D-trained models evaluated on 3D queries; text-conditioned policies evaluated on video-grounded tasks), so that the modality-mismatch challenge flagged in §\ref{subsec:p3_s7_3_6} (text-vs.-image-vs.-3D retrieval) becomes a first-class evaluation axis rather than a per-paper case study.

\paragraph{Embodied Long-Horizon Coverage.}
The Robot block in Table~\ref{tab:benchmark_comparison} covers navigation thoroughly but under-represents long-horizon manipulation and multi-room interactive tasks. 
The §\ref{subsec:p3_s7_3_7} 15-row subset (Table~\ref{tab:spatial_reasoning_benchmark}) further restricts attention to image-and-video benchmarks; the embodied long-horizon suites ALFWorld and VirtualHome (§\ref{subsec:p3_s7_3_4}) and cross-embodiment evaluation on Open X-Embodiment (§\ref{subsec:p3_s7_3_2}) are the nearest existing coverage, and standardizing on them would close the long-horizon-spatiotemporal-coherence challenge flagged under Challenges and Trends in §\ref{subsec:p3_s7_3_7}.

\paragraph{From Descriptive to Interactive Spatial Reasoning.}
The Task Type column of Table~\ref{tab:benchmark_comparison} is dominated by visual question answering: the overwhelming majority of benchmarks ask a model to \emph{describe} static spatial relations (``where is the object''), and only the small Robot block evaluates reasoning that is conditioned on the agent's own actions.
Yet the trajectory traced in this chapter, from passive observation to situated interaction, calls for benchmarks that test \emph{interactive} spatial reasoning: predicting how an action changes the spatial state, judging affordances and reachability (``can I fit, reach, or grasp this''), and reasoning counterfactually about manipulation outcomes.
Such evaluation couples the action-conditioned reasoning of §\ref{subsec:p3_s7_2_5} with the embodied hybrid pipelines of §\ref{sec:p3_s7_4}, extending the descriptive benchmarks surveyed here toward the manipulation- and affordance-grounded settings that embodied deployment ultimately demands.



\section{Summary and Outlook}
\label{sec:p3_s7_6}

This chapter followed spatial reasoning along the spectrum on which it developed: inference over explicit geometric, physical, topological, logical, and causal models (§\ref{sec:p3_s7_1}); reasoning acquired from data (§\ref{sec:p3_s7_2}); reasoning elicited from large pre-trained models (§\ref{sec:p3_s7_3}); hybrids that combine these components (§\ref{sec:p3_s7_4}); and the benchmarks that measure them (§\ref{sec:p3_s7_5}).

\textbf{Precision and grounding sit at opposite ends of the spectrum.} Rule-based methods are inspectable and can certify feasibility or safety, but they assume exact, fully specified inputs and a closed vocabulary of predicates, and they break down where perception is noisy or a task falls outside the authored rules. Learning-based methods acquire spatial competence from data, yet often reach high benchmark accuracy through dataset-specific shortcuts that do not survive novel compositions (§\ref{subsec:p3_s7_2_1}), and they need far more interaction than classical planners. Foundation models bring broad semantic competence, but their visual encoders retain only coarse geometry, textual reasoning mistracks coordinates over multi-step chains (§\ref{subsec:p3_s7_3_3}), and vision-language-action models often interpret an instruction correctly and execute it incorrectly (§\ref{subsec:p3_s7_3_2}). Hybrid systems aim to pair the semantic reach of the last family with the metric precision and verifiability of the first two.

\textbf{The failures concentrate at interfaces.} Across tool-augmented, memory-augmented, and hybrid reasoning, the recurring limitations arise less within any single component than where components meet: in forming a query that a memory or tool can answer, in translating between closed and open vocabularies, and in keeping coordinate frames and scales consistent across modules (§\ref{sec:p3_s7_4}).

\textbf{Evaluation lags the shift toward situated reasoning.} The field is moving from questions posed to a detached observer toward agents that must localize themselves, gather evidence, and act, yet no current benchmark exercises perception fidelity, symbolic precision, memory persistence, and long-horizon action together, and few attribute an end-to-end score to the components that produced it (§\ref{sec:p3_s7_5}).

Reasoning in this chapter takes the evidence of Chapter~\ref{chapter:3:sec:6} as given. Chapter~\ref{sec:generation} turns to the inverse problem of producing spatial content, and its world models (§\ref{sec:p3_s8_6}) supply the simulators that world-model-based reasoning (§\ref{subsec:p3_s7_3_8}) consults.

\newpage


\chapter{Spatial Reconstruction and Generation}
\label{sec:generation}


\FramedInitial{S}{patial} reconstruction and generation is the productive counterpart to spatial perception and reasoning: rather than analyzing a scene that is given, it synthesizes spatial content, recovering or creating the geometry, appearance, and dynamics of objects, scenes, and worlds from conditions such as text, images, partial observations, or explicit constraints.
Where perception (Chapter~\ref{chapter:3:sec:6}) maps observations into structured spatial representations and reasoning (Chapter~\ref{sec:reasoning}) draws inferences over them, generation closes the loop by mapping intent back into concrete spatial artifacts, a capability central to embodied agents, content creation, simulation, and world modeling.

The chapter is organized along an axis of increasing output structure and dimensionality, from flat images to volumetric objects and scenes, to video, embodied motion, and full four-dimensional space-time content, and finally to world models and to programmatic generation that emits interpretable, executable structure. What holds this range together is a single recurring backbone: an encoder maps conditions into a latent, a generator samples structure, and a decoder or renderer produces the final artifact, specialized to each output domain. Read this way, the subfields are less seven separate problems than one conditional-generation abstraction instantiated against progressively richer codomains.


\section{2D Spatial Generation}
\label{sec:p3_s8_1}

\FramedInitial{T}{he} 2D spatial generation encompasses a broad spectrum of tasks that require synthesizing or editing images under explicit spatial constraints.
These constraints may be \emph{user-specified} (bounding boxes, masks, sketches, keypoints, depth maps) or \emph{derived} from scene understanding (layout graphs, semantic maps, relational structures), and the overarching goal is to produce images that are not only visually realistic but also \emph{faithful to the specified spatial configuration}: objects appearing in the correct locations, at the correct scales and orientations, while remaining globally coherent in lighting, perspective, and style.
2D generation is where the most mature \emph{semantic-to-pixel prior} in contemporary AI is taught to honor \emph{spatial structure}: it is the proving ground on which large pre-trained visual priors are forced to localize objects, follow geometric scaffolds, respect occlusion ordering, and obey region- or instruction-level edits, all without falling out of the implicit grammar of natural images.

The section follows the 2D pipeline in three steps that share one latent-generator framework.
§\ref{subsec:p3_s8_1_1} reviews foundational text-to-image generation, the semantic-to-pixel prior on which everything else here is built, and the three paradigms (adversarial, autoregressive, and diffusion with its flow-matching generalization) that parameterize it.
§\ref{subsec:p3_s8_1_2} adds a structural condition (layouts, masks, edges, depth, pose, scene graphs, region phrases) and asks how it is injected into a frozen or co-trained generator.
§\ref{subsec:p3_s8_1_3} replaces the empty canvas with an existing image and a mask, so that the generator must change what lies inside the mask while preserving what lies outside it.
% A short Summary (§\ref{subsec:p3_s8_1_summary}) compares the three along shared axes.

\subsection{Text-to-image Generation}
\label{subsec:t2i}
\label{subsec:p3_s8_1_1}

Text-to-image (T2I) generation synthesizes an image from a natural-language description.
For the rest of this chapter its importance lies less in the images themselves than in the prior it provides: a high-capacity mapping from compositional language to the visual space of objects, surfaces, layouts, and viewpoints, learned once at web scale and then reused.
The spatial controllers of §\ref{subsec:p3_s8_1_2} and the editors of §\ref{subsec:p3_s8_1_3} impose constraints on this prior rather than replacing it, and the text-to-3D pipelines of §\ref{sec:p3_s8_2} distill it through score distillation or multi-view consistency objectives; whatever relational competence the prior has, or lacks, therefore propagates downstream.
The field has passed through three architectural paradigms, from adversarial training to autoregressive token modeling to iterative denoising, with diffusion and its flow-matching generalization now setting the standard for quality and prompt fidelity~\cite{rombach2022ldm,peebles2023dit}.


\subsubsection{Task Definition}
We formulate text-to-image generation in a paradigm-agnostic form that exposes the structural choices shared by modern systems.
Let $\mathcal{X}$ denote the space of tokenized text prompts and $\mathcal{Y}=[0,1]^{H\times W\times 3}$ the pixel-space image domain.
Almost all modern T2I systems operate not directly on $\mathcal{Y}$ but on a compressed \emph{latent space} $\mathcal{Z}\subset\mathbb{R}^d$ provided by an autoencoder pair
\begin{equation}
E:\mathcal{Y}\rightarrow\mathcal{Z},\qquad D:\mathcal{Z}\rightarrow\mathcal{Y},\qquad D\!\circ\!E\approx \mathrm{id}_{\mathcal{Y}},
\label{eq:t2i_autoencoder}
\end{equation}
where $E$ encodes a training image into a low-dimensional code and $D$ decodes a code back into a pixel image.
A text encoder $\phi:\mathcal{X}\rightarrow\mathcal{E}$ maps the prompt into a conditioning embedding $\phi(\mathbf{x})\in\mathcal{E}$.
The core learning problem is then to train a conditional latent sampler
\begin{equation}
G_\theta:\mathcal{E}\rightarrow\mathcal{Z},\qquad \hat{\mathbf{z}}\sim p_\theta\!\big(\mathbf{z}\,\big|\,\phi(\mathbf{x})\big),
\label{eq:t2i_generator}
\end{equation}
so that, after decoding, the resulting image $\hat{\mathbf{y}}=D(\hat{\mathbf{z}})$ is both photorealistic and semantically aligned with the prompt $\mathbf{x}$.
The three paradigms covered below differ chiefly in how the conditional distribution $p_\theta(\mathbf{z}\mid\phi(\mathbf{x}))$ is parameterized and optimized; the autoencoder pair $(E,D)$ and the text encoder $\phi$ are largely shared design choices across paradigms.
Figure~\ref{fig:p3_s8_1_1_framework} schematizes this paradigm-agnostic pipeline and the substrate that the three paradigms share.

\begin{figure}[!ht]
\centering
\includegraphics[width=\linewidth]{figures/part3/section8/t2i_framework_and_paradigms.pdf}
\caption{The latent-generator view of text-to-image generation. Top: the shared pipeline of Eqs.~\eqref{eq:t2i_autoencoder} and~\eqref{eq:t2i_generator}; the encoder $E$ and the dashed paths are used only in training. Bottom: the sampler $G_\theta$ as (a) a GAN, (b) an autoregressive token model, and (c) diffusion or flow matching. Not every architecture uses a separate autoencoder; the corgi is a schematic illustration.}
\label{fig:p3_s8_1_1_framework}
\end{figure}


\subsubsection{Methodology}

Three generative paradigms, taken in the order in which each dominated T2I, parameterize the sampler $G_\theta$ of Eq.~\eqref{eq:t2i_generator} differently; the bottom row of Figure~\ref{fig:p3_s8_1_1_framework} draws them against the common pipeline.
A second set of advances refines the sampler without redefining it: backbone evolution, sampling acceleration, and the migration of leading systems to proprietary stacks.
Structural conditioning interfaces, the fourth such refinement, are the subject of §\ref{subsec:p3_s8_1_2} and are not repeated here.

\paragraph{Generative Adversarial Networks.}\label{para:p3_s8_1_1_gan}
GAN-based T2I models instantiate $G_\theta$ as a feed-forward generator $G_\theta(\mathbf{n},\phi(\mathbf{x}))\!\to\!\hat{\mathbf{z}}$ driven by an auxiliary noise vector $\mathbf{n}\sim\mathcal{N}(0,I)$; the conditional distribution $p_\theta(\mathbf{z}\mid\phi(\mathbf{x}))$ is defined only implicitly, as the pushforward of $\mathbf{n}$ through $G_\theta$.
Training follows the two-player game of \citet{goodfellow2014gan}, pitting the generator against a discriminator $D_\psi$,
\begin{equation}
\min_\theta \max_\psi\ \mathbb{E}_{\mathbf{y}\sim p_\text{data}}\!\big[\log D_\psi(\mathbf{y},\phi(\mathbf{x}))\big] + \mathbb{E}_{\mathbf{n}}\!\big[\log\!\big(1-D_\psi(D(G_\theta(\mathbf{n},\phi(\mathbf{x}))),\phi(\mathbf{x}))\big)\big],
\label{eq:t2i_gan_objective}
\end{equation}
in which $D_\psi$ scores both photorealism and text-image consistency.
The result is sharp high-frequency detail and single-pass inference, bought at the price of unstable training and mode collapse.
Within this framework, StackGAN~\cite{stackgan2017} generated a low-resolution image first and refined it with a second GAN conditioned on both the text and the coarse image; AttnGAN~\cite{attngan2018} let the generator attend to individual words while rendering different regions, with an auxiliary image-text matching loss (DAMSM) that reduced attribute errors such as wrongly colored objects; XMC-GAN~\cite{xmcgan2021} replaced that loss with cross-modal contrastive learning and improved multi-object prompts.
By 2021 diffusion models had overtaken these systems on COCO, chiefly because they train stably and cover more of the image distribution, and research attention moved on; the GAN emphasis on high-frequency detail survives in hybrid and adversarially fine-tuned samplers.

\paragraph{Autoregressive Models.}\label{para:p3_s8_1_1_ar}
Autoregressive (AR) models discretize the latent space: the autoencoder $E$ is realized as a vector-quantized variational autoencoder (VQ-VAE) whose output is a sequence of discrete tokens $\mathbf{z}=(z_1,\dots,z_N)\in\mathcal{V}^N$ over a learned codebook $\mathcal{V}$, and $G_\theta$ is factorized in causal raster order,
\begin{equation}
p_\theta\!\big(\mathbf{z}\,\big|\,\phi(\mathbf{x})\big) \;=\; \prod_{n=1}^{N} p_\theta\!\big(z_n\,\big|\,z_{<n},\,\phi(\mathbf{x})\big),
\label{eq:t2i_ar_factorisation}
\end{equation}
trained by maximum likelihood on the next-token cross-entropy.
T2I thus becomes sequence-to-sequence translation, and the transformer that performs it inherits the compositional competence of language models: DALL-E~\cite{dalle2021} modeled text and image tokens as a single sequence and showed zero-shot generalization across prompt types.
The cost is decoding latency, one forward pass per token, which grows with resolution.
Parti~\cite{parti2022} scaled the transformer to 20B parameters, generated at $256\times256$, and reached $1024\times1024$ with a convolutional super-resolution stage; Muse~\cite{chang2023muse} abandoned strict sequential order for masked token prediction, filling in many tokens per step over a handful of rounds, and matched diffusion quality on several benchmarks at a fraction of the sampling time.
AR generation remains attractive where images are produced jointly with text (interleaved documents, storyboards) and a single sequence model handles both.

\paragraph{Diffusion Models.}\label{para:p3_s8_1_1_diffusion}
Diffusion models realize $G_\theta$ as iterative denoising rather than a single feed-forward pass.
A fixed \emph{forward} process corrupts the clean latent $\mathbf{z}_0=E(\mathbf{y})$ along a schedule $(\alpha_t,\sigma_t)_{t=0}^{T}$, $\mathbf{z}_t=\alpha_t \mathbf{z}_0+\sigma_t\epsilon$ with $\epsilon\sim\mathcal{N}(0,I)$ and $\mathbf{z}_T\!\approx\!\mathcal{N}(0,I)$; a noise-prediction network $\epsilon_\theta$ learns to invert it under the denoising score-matching objective
\begin{equation}
\min_\theta\ \mathbb{E}_{t,\mathbf{z}_0,\epsilon}\!\left[\left\|\epsilon-\epsilon_\theta\!\big(\mathbf{z}_t,\,t,\,\phi(\mathbf{x})\big)\right\|_2^2\right],
\label{eq:t2i_ldm_objective}
\end{equation}
which is the diffusion-specific realization of the training objective implicit in Eq.~\eqref{eq:t2i_generator}.
At inference the model starts from $\mathbf{z}_T\sim\mathcal{N}(0,I)$, applies $\epsilon_\theta$ for $t=T,\dots,1$, and decodes $\hat{\mathbf{y}}=D(\hat{\mathbf{z}}_0)$; each step corrects a small error, which is why the final samples are both detailed and faithful.
Running this process in the latent space of Eq.~\eqref{eq:t2i_autoencoder} rather than in pixels, as Latent Diffusion Models~\cite{rombach2022ldm} did (\eg, $4\times64\times64$ instead of $3\times512\times512$ for a 512px image), cut memory and compute enough to train on commodity hardware and to release Stable Diffusion openly.
Conditioning the denoiser on embeddings from a pretrained language-image model, as in GLIDE~\cite{glide2022} and DALL$\cdot$E~2~\cite{dalle2_2022}, improved prompt alignment, and classifier-free guidance~\cite{ho2022classifierfree} supplied a single knob that trades diversity against fidelity by extrapolating the conditional prediction away from the unconditional one.
The remaining weakness is sampling cost, dozens of sequential network evaluations per image; higher-order ODE samplers and the distillation methods below address it.
Diffusion or flow-matching cores now underlie most leading systems, with the native multimodal autoregressive stacks discussed under proprietary systems the main exception.

\paragraph{Flow Matching and Rectified Flow.}\label{para:p3_s8_1_1_fm}
Rather than a fourth paradigm, flow matching is diffusion reformulated in continuous time.
It parameterizes $G_\theta$ as a learned \emph{velocity field} $v_\theta(\mathbf{z}_t,t,\phi(\mathbf{x}))$ instead of a noise predictor: training picks the straight interpolation $\mathbf{z}_t=(1-t)\mathbf{z}_0+t\mathbf{z}_1$ between a clean latent $\mathbf{z}_0=E(\mathbf{y})$ and a Gaussian sample $\mathbf{z}_1\!\sim\!\mathcal{N}(0,I)$ and regresses the velocity,
\begin{equation}
\min_\theta\ \mathbb{E}_{t,(\mathbf{z}_0,\mathbf{z}_1)}\!\left[\left\|v_\theta(\mathbf{z}_t,t,\phi(\mathbf{x}))-(\mathbf{z}_1-\mathbf{z}_0)\right\|_2^2\right],
\label{eq:t2i_fm_objective}
\end{equation}
and inference integrates the probability-flow ODE $\mathrm{d}\mathbf{z}_t/\mathrm{d}t=v_\theta(\mathbf{z}_t,t,\phi(\mathbf{x}))$ from $t=1$ back to $t=0$.
The formulation was introduced concurrently by~\citet{lipman2023flow} (Flow Matching) and~\citet{liu2023rectifiedflow} (Rectified Flow); \citet{albergo2023building} gave the unifying stochastic-interpolant view in which diffusion is the special case of a noisier interpolation schedule.
Three consequences matter in practice. Straight-line interpolations yield nearly straight ODE trajectories, so matched quality is reached in noticeably fewer steps (genuine one-step generation still needs a reflow or distillation stage); the objective is a single regression with no noise-schedule design and no per-schedule guidance recalibration; and the loss scales smoothly with model size and resolution.
Stable Diffusion~3~\cite{sd3rft2024} scaled the recipe to T2I on a DiT-style transformer backbone and reported the best GenEval scores among open models together with favorable human preference on Parti-prompts; the Flux family followed the same recipe.
Flow matching composes with every refinement below as readily as a noise-predictor core does and has become the default continuous-time formulation in 2024-onwards systems.

\paragraph{Backbone Evolution: From U-Net to Diffusion Transformers.}
Early diffusion models used the U-Net for its multiscale feature fusion and locality bias.
The Diffusion Transformer (DiT) of~\citet{peebles2023dit} replaced it with a stack of transformer layers over latent patch tokens, giving every denoising step a global receptive field, and showed that FID improves monotonically with depth, width, and token count along predictable scaling laws, reaching FID 2.27 on class-conditional ImageNet at $256^2$.
Attention cost grows quadratically with resolution, so hybrid CNN-transformer designs and efficient attention remain active; the DiT lineage nonetheless underlies SD3 and most 2024-onwards systems, and it points toward a single transformer that generates several modalities, differing only in tokenization and training data.

\paragraph{Sampling Acceleration and Distillation.}
Fifty or more denoising steps rule out interactive use and low-power deployment.
\citet{salimans2022progressive} distilled a $T$-step teacher into a $T/2$-step student and iterated the halving to reach $1$--$4$ steps.
Consistency Models~\cite{song2023consistency} train a network to map a noisy latent at any timestep directly to the clean image, either from scratch with a consistency loss or by distilling a pretrained diffusion model, and Latent Consistency Models~\cite{lcm2023} applied the idea to Stable Diffusion's latent space, matching $50$-step quality in $2$--$4$ steps.
All of these are post-training procedures that leave the architecture untouched, so any strong backbone can be made fast after the fact.

\paragraph{Proprietary and Commercial Systems.}
Leading systems have migrated to proprietary, industrially scaled stacks of three kinds.
Native multimodal models such as OpenAI's GPT~Image~1.5~\cite{openai2025images} and Google's Gemini~3~Pro~Image (Nano~Banana~Pro) run a unified transformer over interleaved text and image tokens and emphasize semantic reasoning and multi-turn editing; latent-diffusion iterations such as Imagen~4 scale the LDM recipe with stronger text encoders and curated data; preference-tuned systems such as Midjourney~v7~\cite{midjourney_v7_alpha} and FLUX.2~[max]~\cite{bfl_flux2_max} fine-tune strong priors on human-preference data (FLUX.2 on the rectified-flow recipe of §\ref{para:p3_s8_1_1_fm}; Midjourney's architecture is undisclosed).
Each reduces to one of the three paradigms plus some of the refinements just listed; because architectural and training detail is withheld, public benchmarking and reproduction still rest on the open SD / DiT stack.

\subsubsection{Benchmarks}
Three independent things are usually conflated under ``benchmark'': the training corpora on which models are fit, the evaluation suites of prompts that probe particular capabilities, and the metrics that score the resulting images.
A model trained on LAION-5B can be evaluated on GenEval prompts and scored by FID, CLIPScore, or an LLM judge, and disagreement between those scores is itself informative.

\paragraph{Training Corpora.}
Corpus scale has driven much of T2I progress.
Early work fit fine-grained sets with ten captions per image, CUB-200-2011 (11,788 bird images)~\cite{wah2011cub} and Oxford-102 Flowers (8,189 images)~\cite{nilsback2008flowers}; MS-COCO Captions~\cite{lin2014microsoft} (123,287 images, five captions each) became the standard general-scene corpus and, for years, the evaluation set as well, a conflation that the suites below were designed to end.
Current models train on web-scale image-text pairs: LAION-5B~\cite{schuhmann2022laion5b} releases 5.85 billion CLIP-filtered pairs, and proprietary corpora of similar scale, less documented, underlie DALL-E~2 and Parti.
Scale brought duplication, demographic bias, unsafe content, and licensing disputes with it, and filtered subsets such as LAION-Aesthetics and curation pipelines are now part of every training recipe.

\paragraph{Evaluation Benchmark Suites.}
An evaluation suite is a curated prompt set plus a verification protocol: prompts are adversarial or capability-stratified, reference images are usually absent, and scoring is automatic (an object detector, a VQA model, or an LLM/MLLM judge) rather than reconstruction-based.
Table~\ref{tab:p3_s8_1_1_eval_benchmarks} lists the main suites since 2022.
The earliest, DrawBench~\cite{saharia2022photorealistic} and PartiPrompts~\cite{parti2022}, are human-pairwise prompt sets released with Imagen and Parti; HRS-Bench~\cite{bakr2023hrsbench} and HEIM~\cite{lee2023heim} are holistic, scoring 13 and 12 axes respectively; TIFA~\cite{hu2023tifa}, GenEval~\cite{ghosh2023geneval}, and T2I-CompBench~\cite{huang2023t2icompbench} introduced automatic verification by VQA, object detection, and BLIP-VQA plus UniDet, and the latter two are now reported alongside FID and CLIPScore in most papers.
For spatial intelligence the salient suite is VISOR~\cite{gokhale2022visor}, which isolates two-object relations (left/right, above/below) over $25{,}280$ prompts and finds that otherwise strong models frequently render the wrong configuration.
DPG-Bench~\cite{hu2024ella} targets dense, long compositional prompts, GenAI-Bench~\cite{li2024genaibench} pairs compositional prompts with VQAScore judging, and WISE~\cite{niu2025wise} asks whether generated images agree with real-world facts (flag colors, landmarks, anatomy), a knowledge axis orthogonal to compositional binding.

\begin{table}[!ht]
\centering
\fontsize{9}{11}\selectfont
\setlength{\tabcolsep}{1.0mm}
\caption{Representative evaluation benchmark suites for text-to-image generation, ordered approximately by year of release. ``Scoring'' indicates the dominant verification mechanism: detector- or VQA-based automatic checking, learned preference models, MLLM/LLM-as-judge rubric scoring, or aggregated human preference. Several suites combine multiple scoring channels.}
\label{tab:p3_s8_1_1_eval_benchmarks}
\begin{tabular}{L{0.17\textwidth} C{0.06\textwidth} C{0.11\textwidth} L{0.33\textwidth} L{0.25\textwidth}}
\toprule
\textbf{Benchmark} & \textbf{Year} & \textbf{\#~Prompts} & \textbf{Evaluation Axis} & \textbf{Scoring} \\
\midrule
DrawBench~\cite{saharia2022photorealistic}       & 2022 & 200   & Broad capability + alignment ($11$ categories) & Human pairwise \\
PartiPrompts~\cite{parti2022}     & 2022 & 1{,}600 & Capability spectrum ($12$ categories) & Human pairwise \\
VISOR~\cite{gokhale2022visor}     & 2022 & 25{,}280 & Two-object spatial relations (left/right, above/below) & Object-detector verification \\
TIFA~\cite{hu2023tifa}            & 2023 & 4{,}081 & Faithfulness (VQA over auto-generated questions) & VQA / LLM-as-judge \\
HRS-Bench~\cite{bakr2023hrsbench} & 2023 & 45{,}000 & Holistic; $13$ skills $\times$ $50$ scenarios & Automatic, mixed \\
HEIM~\cite{lee2023heim}           & 2023 & 62 scenarios & Holistic; $12$ aspects (alignment, quality, knowledge, bias, toxicity, \ldots) & Automatic $+$ human \\
GenEval~\cite{ghosh2023geneval}   & 2023 & 553   & Object-centric: single/two-obj., counting, color, position, attribute & Detector verification \\
T2I-CompBench~\cite{huang2023t2icompbench} & 2023 & 6{,}000 & Compositional: attribute binding, object relationships, complex prompts & BLIP-VQA $+$ UniDet \\
DPG-Bench~\cite{hu2024ella}       & 2024 & 1{,}065 & Dense, long compositional prompts & MLLM-based scoring \\
GenAI-Bench~\cite{li2024genaibench} & 2024 & 1{,}600 & Compositional T2V/T2I prompts & VQAScore (LLM-judge) \\
WISE~\cite{niu2025wise}           & 2025 & 1{,}000 & World-knowledge-informed semantic accuracy & LLM-based grading \\
\bottomrule
\end{tabular}
\end{table}

\paragraph{Evaluation Metrics.}
Metrics fall into four tiers.
Distributional metrics compare the generated and real image populations and say nothing about any individual prompt: the Fr\'echet Inception Distance (FID)~\cite{heusel2017gans} is the Wasserstein-2 distance between Gaussian fits of Inception-V3 features of generated and real images, the Inception Score rewards images that are individually confidently classified yet collectively diverse, and generative precision/recall~\cite{sajjadi2018assessing,kynkaanniemi2019improved} separate quality from coverage through neighborhood tests in feature space.
Alignment metrics score one image against its prompt; CLIPScore~\cite{hessel2021clipscore}, the cosine similarity between the CLIP embeddings of $\hat{\mathbf{y}}$ and $\mathbf{x}$, is cheap and reference-free but inherits CLIP's insensitivity to attribute binding, counting, and spatial relations.
Learned preference models, HPSv2~\cite{wu2023hpsv2}, ImageReward~\cite{xu2023imagereward} (fit to $137$K human comparisons), and PickScore~\cite{kirstain2023pickapic} (fit to the ${\sim}500$K real-user comparisons of Pick-a-Pic), regress a scalar that correlates with average user preference, bundling realism, alignment, and aesthetics; they have become the default automatic ranking while inheriting their raters' biases and degrading out of distribution.

The newest tier uses a vision-language model as judge.
TIFA~\cite{hu2023tifa} has a language model write verification questions from the prompt and a VQA model answer them against the image, the score being the fraction answered correctly; VQAScore~\cite{lin2024evaluating} collapses this to a single yes/no probability under a fine-tuned CLIP-FlanT5 model and reports the best human correlation on compositional suites; VIEScore~\cite{ku2023viescore} and LLMScore~\cite{lu2023llmscore} prompt a multimodal LLM with a structured rubric and return per-axis scores with brief rationales, and direct prompting of frontier VLMs with custom rubrics has become routine for ad-hoc capability probes.
Judge metrics track human preference on compositional and spatial prompts far better than CLIPScore, at the cost of a frontier-model call per image, sensitivity to rubric wording, and drift across judge-model versions that makes numbers from different papers hard to compare.
No single tier suffices, and current papers report a tuple (FID or IS, CLIPScore, one or two preference scores, one judge score) on several suites plus a small human study; cross-metric disagreement is read as a warning rather than averaged away.


\subsubsection{Challenges and Trends}
The weakness that matters most for this handbook is \textbf{spatial-relation faithfulness}, a sub-problem of compositionality.
Models frequently misplace objects relative to one another: a prompt that puts a cup to the left of a laptop is often rendered with the two swapped or stacked, and VISOR~\cite{gokhale2022visor} shows that language-specified relations are frequently ignored outright.
Data-centric remedies help; SPRIGHT~\cite{chatterjee2024spright} re-captions training images with explicit spatial descriptions and measurably improves consistency, but reliable control over where objects appear remains open.
Because layout-conditioned generation (§\ref{subsec:p3_s8_1_2}) and text-to-3D (§\ref{sec:p3_s8_2}) inherit whatever relational competence this prior provides, the limitation is foundational rather than cosmetic.
Fine structure is a second gap: legible text, hands, and faces still fail in otherwise strong models and appear to need architectural or data-specific fixes rather than scale alone.

\textbf{Evaluation reliability} is the third.
Automatic scores correlate imperfectly with human judgment precisely on the compositional, counting, and spatial-relation prompts that matter here; judge-based metrics close part of the gap but add rubric sensitivity and judge-version drift, so headline numbers across papers are not directly comparable.
As GenEval and T2I-CompBench become standard, their prompts risk leaking into web-scraped training corpora, in which case high scores increasingly measure memorization rather than capability; periodic refresh of held-out prompts, blinded judges, and explicit train/test decontamination are the obvious countermeasures and are not yet routine.

\textbf{Societal concerns} follow the data.
Models inherit demographic bias from their corpora and will produce unsafe content on request; prompt classifiers, safety filters, and fine-tuning on curated subsets reduce but do not remove the problem, and photorealistic synthesis raises misinformation and intellectual-property questions that watermarking and data curation address only in part.
The extension of these priors along the time axis, where identity and motion must stay consistent across frames, is the subject of §\ref{sec:p3_s8_3}; the next two subsections stay in the image plane and ask how the prior can be made to obey spatial constraints and local edits.



\subsection{2D Spatial-aware Image Generation}
\label{subsec:p3_s8_1_2}
\emph{2D spatial-aware image generation} synthesizes images that are faithful not only to a text prompt but to explicit spatial constraints in the image plane: where objects appear, how large they are, what shape they take, and how they relate to one another (\emph{left of}, \emph{on}, \emph{behind}).
The constraint may be a set of boxes or a scene layout, an instance or semantic mask, an edge or sketch map, a depth or normal map, 2D keypoints, a scene graph, or a phrase bound to a region, and the pipelines that honor it are almost all diffusion models augmented with a conditioning pathway or a guidance term.
The tension that organizes the subfield is three-way: \emph{adherence} to hard constraints against \emph{perceptual realism} and \emph{diversity}, and both against \emph{compositionality} when many small or occluding objects carry discrete relations.
Depth and normal maps enter here as structural guides for occlusion ordering and surface orientation, not as metric measurements; recovering true distance belongs to reconstruction (§\ref{sec:p3_s8_2}).


\subsubsection{Task Definition}
\label{subsubsec:p3_s8_1_2_task}
We inherit the latent-generator framework of §\ref{subsec:p3_s8_1_1}.
Let $\mathcal{X}$ denote tokenized prompts, $\mathcal{Z}$ the latent space induced by the autoencoder pair $(E,D)$ of Eq.~\eqref{eq:t2i_autoencoder}, and $\phi:\mathcal{X}\!\rightarrow\!\mathcal{E}$ the text encoder.
Spatial-aware generation \emph{augments} this framework with a structural condition $\mathbf{c}\in\mathcal{C}$ drawn from one or more conditioning channels (boxes, masks, edges, depth, pose, scene graphs, region phrases, reference images), encoded by a condition encoder $\kappa:\mathcal{C}\!\rightarrow\!\mathcal{C}'$.
The conditional latent sampler of Eq.~\eqref{eq:t2i_generator} is extended to a joint conditional
\begin{equation}
G_\theta:\mathcal{E}\times\mathcal{C}'\rightarrow\mathcal{Z},\qquad \hat{\mathbf{z}}\sim p_\theta\!\big(\mathbf{z}\,\big|\,\phi(\mathbf{x}),\,\kappa(\mathbf{c})\big),
\label{eq:t2i_spatial_objective}
\end{equation}
and the generated image $\hat{\mathbf{y}}=D(\hat{\mathbf{z}})$ is required to be (i) photorealistic, (ii) semantically aligned with $\mathbf{x}$, and (iii) faithful to the structural condition $\mathbf{c}$ in placement, scale, shape, and pairwise relation.
The unconditional limit $\mathbf{c}=\varnothing$ recovers pure T2I generation as in §\ref{subsec:p3_s8_1_1}; the substance of the present subsection lies in how $\kappa$ is realized for each conditioning channel and how its output is injected into $G_\theta$.
Which channel populates $\mathbf{c}$, and how densely, ranges from pixel-level maps (masks, edges, depth) through coordinate-level scaffolds (boxes, keypoints) to discrete symbols (scene graphs, region phrases); the modality survey below follows that ordering.
Interactive constraints supplied by dragging handles during generation belong to the editing regime of §\ref{subsec:p3_s8_1_3}.


\subsubsection{Methodology}
\label{subsubsec:p3_s8_1_2_methodology}
Two axes organize the methods: the \emph{conditioning modality} (what kind of structural input $\mathbf{c}$ is consumed) and the \emph{injection mechanism} (how $\kappa(\mathbf{c})$ enters the latent sampler $G_\theta$).
They are nearly independent. One mechanism, a ControlNet-style side branch, serves canny edges, pose skeletons, depth, and segmentation alike, and one modality, bounding boxes, is served by side branches, grounded cross-attention, or inference-time attention shaping, so the design space is a matrix of (modality~$\times$~mechanism) cells rather than a list of monolithic methods.
Figure~\ref{fig:p3_s8_1_2_mechanisms} draws the five mechanisms on one backbone silhouette.

\paragraph{Conditioning Modalities.}
\label{para:p3_s8_1_2_modalities}
\textit{Dense pixel-level scaffolds} supply a value at every location, which makes them the easiest channels to follow and the most restrictive on creativity.
Instance and semantic masks drive region-conditioned synthesis, with MultiDiffusion~\cite{multidiffusion} offering a training-free protocol that samples jointly over overlapping mask windows; edges, sketches, and line art, popularized by the canny and soft-edge variants of ControlNet~\cite{zhang2023adding} and the edge head of T2I-Adapter~\cite{t2iadapter2023}, fix only the location of intensity discontinuities and leave texture and appearance to the prompt; depth and normal maps, served by ControlNet's depth variant, impose coarse occlusion ordering and surface orientation without requiring a metric reconstruction.
\textit{Coordinate-level scaffolds} drop to per-region or per-point supervision.
Boxes with labels are the most studied channel, a set of $(\text{box},\text{label})$ tuples fixing location and rough scale, with LayoutDiffusion~\cite{zheng2023layoutdiffusion}, LayoutLLM~\cite{qu2023layoutllm}, and the open-set GLIGEN~\cite{li2023gligen} as representatives; 2D keypoints (typically human pose) are sparser still, a handful of joint coordinates, injected through ControlNet or handled by GLIGEN as a special case of grounding.
\textit{Discrete symbolic structure} carries no pixel- or coordinate-level information at all, and the model must map it onto the image plane itself.
Scene graphs reach back to the sg2im line of~\citet{johnson2018image} and re-emerged in the diffusion era through LLM-emitted layouts, as in LayoutLLM, where the graph is first decoded into a coordinate-level layout; region phrases and referring expressions bind a text span to a box, mask, or free-form polygon and are served by GLIGEN's grounded cross-attention and by phrase-grounding heads attached to SD and SDXL~\cite{sdxl2023}.

\paragraph{Injection Mechanisms.}
\label{para:p3_s8_1_2_mechanisms}
Five mechanism families cover the literature; production pipelines routinely stack a side branch for dense geometric channels, an adapter for style or a reference image, and inference-time guidance to sharpen adherence on adversarial prompts.

\textbf{Side-branch controllers.}
The backbone weights stay frozen and an auxiliary branch mirroring the encoder half of the base U-Net is trained, connected through \emph{zero-convolution} layers whose weights are initialized to zero so that the augmented model reproduces the base model exactly at the start of training and degrades gracefully if the branch is poorly tuned.
ControlNet~\cite{zhang2023adding} is the prototype and the first widely adopted retrofit pathway for non-text structural control: the condition (a canny edge map, an OpenPose skeleton, a depth or normal map, a segmentation mask, or a scribble) is encoded by the side branch in parallel with the noisy latent and injected into the base U-Net through additive residual connections at each scale.
One branch is trained per modality on labeled image--condition pairs, which produced a library of plug-in controllers for SD~1.5 and SDXL~\cite{sdxl2023} that compose with prompt conditioning at inference time without retraining the backbone; most later designs either shrink this recipe or fold it into a unified adapter library.

\textbf{Lightweight adapters.}
Where the side branch spends capacity (half a U-Net per modality), T2I-Adapter~\cite{t2iadapter2023}, developed concurrently with the same goal of plug-in control over a frozen backbone, spends as little as possible: a few residual blocks per modality whose multi-scale features are added directly into the corresponding encoder stages.
Several adapters can be combined at inference, an image conditioned simultaneously on an OpenPose skeleton, a depth map, and a color palette with a scalar weight on each, so the adapter is preferred for multi-channel fusion and fast iteration, and the side branch where maximum single-channel fidelity matters.

\textbf{Grounded cross-attention.}
Side branches and adapters run an external pathway alongside the backbone; the next two families intervene \emph{inside} its attention stack.
GLIGEN~\cite{li2023gligen} inserts a new \emph{gated self-attention} layer between the existing self-attention and cross-attention of each U-Net transformer block; the layer attends jointly over the visual tokens and a set of \emph{grounding tokens} that encode each $(\text{phrase},\text{box})$ pair, a learned scalar gate controls how strongly the grounding signal enters the residual stream, and only the new layer is trained, so the network falls back to ungrounded behavior when no constraint is supplied.
The problem it solves is complementary to the side branch: not how to inject a dense pixel-level scaffold, but how to introduce \emph{open-set} symbolic grounding (any category named in text, localized by a box, keypoints, or a reference image) into a frozen T2I backbone, with several grounding channels interleaved within one mechanism.

\textbf{Layout-aware cross-attention.}
LayoutDiffusion~\cite{zheng2023layoutdiffusion} is the dense-supervised cousin of grounded attention: rather than gating a new layer onto a frozen backbone, it co-trains a layout-aware cross-attention block end to end for layout-to-image (L2I) generation specifically.
A structured layout, a set of $(\text{box},\text{class label})$ tuples, becomes a sequence of region tokens fused with the noisy latent at every denoising step, and an explicit \emph{object--scene} factorization models global scene context separately from per-object features so that object identity, scale, and pairwise relations survive in cluttered scenes.
It reports stronger layout IoU and relation satisfaction than ControlNet or T2I-Adapter on multi-object COCO benchmarks, at the cost of being specialized to box-and-label input rather than the broader modality library.

\textbf{Inference-time guidance.}
The fifth family adds no parameters and modifies only sampling.
BoxDiff~\cite{xie2023boxdiff} is the canonical training-free spatial controller: it shapes the cross-attention maps so that activations of box-localized prompt tokens concentrate inside their boxes while the backbone stays untouched.
MultiDiffusion~\cite{multidiffusion} fuses overlapping mask windows at sampling time, and composing scores from several conditional models or tuning the classifier-free guidance scale per region are related strategies.
Because nothing is retrained, these methods suit rapid prototyping, backbones whose weights are unavailable for fine-tuning, and per-prompt sharpening on top of the trained mechanisms above.

\begin{figure}[!ht]
\centering
\includegraphics[width=\linewidth]{figures/part3/section8/Five_injection_mechanisms_for_spatial-aware_T2I.pdf}
\caption{Five mechanisms for injecting a spatial condition $\kappa(\mathbf{c})$ into a text-to-image denoiser: (a) side-branch control (ControlNet~\cite{zhang2023adding}), (b) lightweight adapters (T2I-Adapter~\cite{t2iadapter2023}), (c) gated grounding (GLIGEN~\cite{li2023gligen}), (d) layout-aware cross-attention (LayoutDiffusion~\cite{zheng2023layoutdiffusion}), and (e) inference-time guidance (BoxDiff~\cite{xie2023boxdiff}). The corgi example is illustrative, not a comparison of outputs.}
\label{fig:p3_s8_1_2_mechanisms}
\end{figure}


\subsubsection{Benchmarks}
\label{subsubsec:p3_s8_1_2_benchmarks}
Two things change relative to §\ref{subsec:p3_s8_1_1}.
The ground truth is a \emph{structural annotation} (a layout, a mask, a keypoint set) rather than a reference image, so reconstruction-style metrics give way to adherence metrics that score how closely the generated image honors $\mathbf{c}$; and every channel admits its own adherence metric, so a meaningful report combines a channel-specific score with the realism and alignment scores inherited from the previous subsection.

\paragraph{Training Corpora.}
Spatial-aware models are typically trained on corpora that pair images with rich structural annotations: COCO Captions~\cite{lin2014microsoft} for layouts and captions; ADE20K~\cite{zhou2017ade20k} ($20$K images, $150$ classes) for semantic and instance masks; Visual Genome~\cite{krishna2017visual} for object boxes, attributes, and pairwise relations supporting scene-graph conditioning; RefCOCO / RefCOCO+ / RefCOCOg~\cite{yu2016modeling} for phrase grounding; and PhraseCut~\cite{wu2020phrasecut} for region-prompt mask supervision.
Additional structure (canny edges, depth, normals, OpenPose skeletons) is typically derived on the fly from these base corpora by running pretrained predictors over each image, producing the paired image--condition data on which ControlNet-style and adapter-style pipelines are trained.

\paragraph{Evaluation Benchmark Suites.}
Most of the suites surveyed in §\ref{subsec:p3_s8_1_1} (Table~\ref{tab:p3_s8_1_1_eval_benchmarks}) include axes that probe spatial-aware generation directly and serve as its primary out-of-the-box benchmarks.
GenEval~\cite{ghosh2023geneval} exposes a position subset that scores whether two objects appear in a specified relation (\emph{left of}, \emph{above}); T2I-CompBench~\cite{huang2023t2icompbench} provides explicit \emph{spatial} and \emph{non-spatial relationship} categories with BLIP-VQA and UniDet scoring; HRS-Bench~\cite{bakr2023hrsbench} includes a \emph{spatial} skill among its thirteen; GenAI-Bench~\cite{li2024genaibench} offers compositional prompts on which spatial relations are a frequent failure of CLIP-based scoring; and WISE~\cite{niu2025wise} measures knowledge-grounded location accuracy.
Beyond these, the L2I literature reports task-specific evaluations on COCO layouts: detection-based AP for box-localized objects alongside FID and CLIPScore.

\paragraph{Constraint-specific Metrics.}
Layout adherence is box IoU between the conditioning layout and the boxes recovered by a frozen detector on $\hat{\mathbf{y}}$, often with per-object center distance; mask adherence uses mask IoU and boundary $F$-score against the conditioning mask; edge and sketch adherence report recall (sometimes precision) of edges re-extracted from $\hat{\mathbf{y}}$; geometry adherence reports depth RMSE and normal angular error where ground-truth geometry exists; pose adherence reports keypoint reprojection error or PCK; and relational adherence reports relation-satisfaction accuracy and per-entity object AP for scene-graph and phrase-grounded settings.
These numbers do not aggregate: averaging layout IoU with mask IoU has no operational meaning, since the two live on different scales and reward different shape regimes, and most papers report a different subset of them, which is the main obstacle to cross-paper comparison in this subfield.


\subsubsection{Challenges and Trends}
\label{subsubsec:p3_s8_1_2_challenges}
Spatial-aware T2I inherits the compositionality and evaluation problems of §\ref{subsec:p3_s8_1_1} and sharpens them.

\textbf{Faithful multi-object composition.}
Crowded scenes with many small or partially occluded objects and hard relations (support, containment, depth ordering) fail in three recognizable ways: \emph{attention dilution}, in which a prompt token's cross-attention spreads across the canvas and no object is well localized; \emph{attribute leakage}, in which color, material, or count attaches to the wrong object (the canonical ``red cube on a blue sphere'' confusion); and \emph{object dropping}, in which small or rare entities in a multi-object layout are simply omitted.
On GenEval the position and counting subsets remain the weakest axes across SD~1.5, SDXL, and SD~3~\cite{sd3rft2024}.
Remedies work at three levels: Attend-and-Excite~\cite{chefer2023attend} reweights cross-attention at inference so that every subject token is strongly activated somewhere in the canvas; StructureDiffusion~\cite{feng2023training} uses prompt syntax to constrain attention structure without training; and LLM-emitted layouts, LayoutGPT~\cite{feng2023layoutgpt} and the multimodal-LLM-driven RPG of~\citet{yang2024mastering}, plan a structured layout from the free-form prompt and hand it to an L2I renderer.
Whether the gap closes through a stronger compositional prior in the backbone or through planning as a separate stage is unresolved, and the two are not exclusive.

\textbf{Constraint satisfaction against creativity and diversity.}
Turning guidance up tightens adherence until, past a model-dependent threshold, edges oversharpen, textures flatten, and variation across seeds collapses; the same curve appears for the ControlNet conditioning weight, GLIGEN's gating strength, and the classifier-free guidance scale, leaving practitioners to tune two or three knobs at once.
Annealing the control weight along the denoising trajectory, per-region or per-token guidance scales, and down-weighting guidance where the base model is already confident are the current responses; how to expose such schedules to end users without overwhelming the interface is open.
Chaining spatial-aware generation with repeated local edits adds a drift problem of its own, treated in §\ref{subsec:p3_s8_1_3}.

\textbf{Geometry-aware realism in 2D.}
Depth and normal conditioning fix coarse occlusion but introduce their own artifacts: \emph{boundary banding} where a sharp depth discontinuity meets a smooth gradient in the generated image, \emph{scale ambiguity} because a monocular depth map carries no absolute scale and two objects at the same depth value can be rendered at incompatible physical sizes, and \emph{lighting incoherence} because the map fixes geometry but not the light field, so shadows disagree with the implied surface orientations.
Better estimators are the first lever, with the diffusion-prior Marigold~\cite{ke2024repurposing} producing sharper and more globally consistent depth than CNN baselines and Depth Anything~\cite{DepthAnything} staying robust on out-of-distribution conditioning images; differentiable geometry proxies that translate depth into shading constraints are the second, explored at the boundary with full 3D synthesis (§\ref{sec:p3_s8_2}) but not yet integrated into general T2I backbones.
Region grounding also changes the safety picture, since a harmful edit can be confined to a footprint that whole-image forensics miss; that concern belongs to editing and is taken up in §\ref{subsec:p3_s8_1_3}.


\subsection{2D Spatial Inpainting \& Editing}
\label{subsec:p3_s8_1_3}

\emph{2D spatial inpainting \& editing} modifies an image inside a specified region while leaving the rest intact: filling holes and occlusions, removing or inserting objects, changing attributes (color, material, texture, text), adjusting layout, or harmonizing pasted content with its new background.
The intent arrives as a mask, a structural hint (edges, sketches), a geometry proxy (depth, normals), or a natural-language instruction, and the tensions are between \textbf{local editability} and \textbf{global consistency} of lighting, perspective, and shadow, between \textbf{constraint adherence} and \textbf{realism}, and between \textbf{semantic plausibility} and \textbf{identity preservation} for person- or product-centric edits.
Editing holds the model to its strongest geometric account: the \emph{unedited} region of $I$ supplies ground-truth pixels that any plausible edit must respect, exposing failure modes that pure synthesis can hide, among them color drift, lighting reflow, identity loss, and attribute leakage across the mask boundary.


\subsubsection{Task Definition}
\label{subsubsec:p3_s8_1_3_task}
We inherit the latent-generator framework of §\ref{subsec:p3_s8_1_1}.
Let $(E,D)$ denote the autoencoder pair of Eq.~\eqref{eq:t2i_autoencoder}, $\phi:\mathcal{X}\!\rightarrow\!\mathcal{E}$ the text encoder, and $G_\theta$ the conditional latent sampler of Eq.~\eqref{eq:t2i_generator}.
Inpainting and editing depart from de-novo generation in one crucial respect: the input is no longer the empty set but an existing image $I\in\mathcal{Y}$, encoded into $\mathbf{z}^{\text{in}}=E(I)$.
A mask $\mathcal{M}\in\{0,1\}^{H\times W}$ designates the editable region, and an editing control set $\mathbf{c}\in\mathcal{C}$ (text prompts, structural cues, exemplar patches, or drag-style handles) populates the conditioning channel, encoded by $\kappa:\mathcal{C}\!\rightarrow\!\mathcal{C}'$.
The conditional latent sampler is extended to
\begin{equation}
G_\theta:\mathcal{Z}\times\{0,1\}^{H\times W}\times\mathcal{E}\times\mathcal{C}'\!\rightarrow\!\mathcal{Z},\qquad
\hat{\mathbf{z}}\sim p_\theta\!\big(\mathbf{z}\,\big|\,\mathbf{z}^{\text{in}},\,\mathcal{M},\,\phi(\mathbf{x}),\,\kappa(\mathbf{c})\big),
\label{eq:t2i_edit_objective}
\end{equation}
and the edited image $\hat{I}=D(\hat{\mathbf{z}})$ is required to (i) realize the requested change \emph{inside} $\mathcal{M}$, (ii) preserve $I$ \emph{outside} $\mathcal{M}$ up to imperceptible deviations, and (iii) remain globally coherent in color, lighting, geometry, and semantics.
Setting $\mathbf{z}^{\text{in}}=\varnothing$ and $\mathcal{M}=\mathbf{1}$ recovers the spatial-aware synthesis of §\ref{subsec:p3_s8_1_2} (Eq.~\eqref{eq:t2i_spatial_objective}); the distinctive substance of editing therefore lies in the interplay between the \emph{preservation} constraint outside $\mathcal{M}$ and the \emph{change} constraint inside it.
The mask may be hard, soft, inferred from a referring expression, or implicit (whole-image instruction editing with $\mathcal{M}=\mathbf{1}$).

\paragraph{Variants by Task Type.}
\label{para:p3_s8_1_3_variants}
The tasks differ in what populates $\mathbf{c}$ and $\mathcal{M}$ and in whether ground-truth pixels exist for evaluation.
Reconstruction-style completion carries no semantic intent: hole filling repairs scratches and occlusions from $(I,\mathcal{M})$ alone, and object removal reconstructs the background behind a specified instance.
Constraint-driven local synthesis adds intent inside $\mathcal{M}$: text-guided local editing follows an instruction scoped to a region (``make the \emph{left cup} red''), structure-guided inpainting takes edges, sketches, or semantic maps as a scaffold for shape and boundary, and geometry-guided inpainting conditions on depth or normals to respect occlusion ordering.
Object-level edits treat the region as an object whose identity or embedding in the scene is the unit of change: insertion with harmonization of color, grain, lighting, shadows, and reflections; attribute and appearance editing that changes material or style while preserving identity (garments, products, logos); and exemplar-guided transfer that drives content or style from a reference patch.
Interactive workflows replace declarative constraints with actions taken at edit time: point- or handle-based reshaping, and multi-round, multi-region sessions in which stability across rounds is itself the requirement.
The metrics appropriate to each are collected in Benchmarks below.


\subsubsection{Methodology}
\label{subsubsec:p3_s8_1_3_methodology}
Methods are organized by \emph{editing mechanism}, that is, by how $G_\theta$ of Eq.~\eqref{eq:t2i_edit_objective} is realized so that change is confined to $\mathcal{M}$ and preservation is enforced outside it.
The mechanism is largely independent of the task type above, since text-guided local editing can be served by masked denoising, by inversion and rewrite, or by an instruction-tuned editor, and one mechanism such as inversion serves both attribute edits and structure-guided ones.
Figure~\ref{fig:p3_s8_1_3_mechanisms} sets the five mechanisms against a single example edit.

\begin{figure}[!ht]
\centering
\includegraphics[width=\linewidth]{figures/part3/section8/Five_editing_mechanisms_for_2D_spatial_inpainting_and_editing.pdf}
\caption{Local 2D editing with context preservation (top) and five editing mechanisms: (a) masked denoising, where $M=1$ marks the editable region; (b) inversion and rewrite; (c) instruction-tuned editing; (d) exemplar / reference-guided editing; and (e) drag-based optimization. The cup images are illustrative, not outputs of the surveyed methods.}
\label{fig:p3_s8_1_3_mechanisms}
\end{figure}

\paragraph{Editing Mechanisms.}
\label{para:p3_s8_1_3_mechanisms}

\textbf{Masked latent denoising.}
The most direct realization of Eq.~\eqref{eq:t2i_edit_objective} runs the diffusion sampler only inside $\mathcal{M}$ while clamping the latent outside $\mathcal{M}$ to a noise-perturbed copy of $\mathbf{z}^{\text{in}}$ at every step.
RePaint~\cite{lugmayr2022repaint} does this with a \emph{pretrained, unconditional} DDPM and no retraining, replacing the known-region pixels with a forward-noised sample of the input at each reverse step and adding a \emph{resampling} schedule that re-injects noise periodically so the in-fill can re-harmonize with its context over several rounds.
Blended Diffusion~\cite{avrahami2022blended} made the same scheme text-driven by interleaving a CLIP-guided gradient step that steers the content inside $\mathcal{M}$ toward a target prompt with the context-clamping step outside it, the first widely cited text-guided local editor.
LaMa~\cite{suvorov2022lama} predates the diffusion line and is instead a feed-forward inpainter whose fast Fourier convolutions give every layer an image-wide receptive field, a Fourier-domain inductive bias that stays robust on the very large holes that defeat purely local CNN backbones.
Region-aware noise schedules and classifier-free guidance are standard add-ons, and ControlNet-inpaint and the inpainting variants of T2I-Adapter (§\ref{para:p3_s8_1_2_mechanisms}) compose with masked denoising by supplying the inside-$\mathcal{M}$ condition.

\textbf{Inversion + rewrite.}
A complementary strategy first recovers the noise trajectory that the model would have used to generate $I$ (\emph{inversion}) and then re-runs the reverse process under an edited prompt while reusing or swapping the attention maps of the original (\emph{rewrite}).
Prompt-to-Prompt~\cite{prompt-to-prompt} reuses or reweights source cross-attention maps during generation under the edited prompt, which preserves spatial structure while changing selected semantic content; for a real image an inversion step first supplies a suitable noise trajectory.
Imagic~\cite{kawar2023imagic} jointly optimizes the text embedding and lightly fine-tunes the model to reconstruct $I$, then interpolates the embedding toward the target prompt, preserving identity-level features (faces, garments) under non-rigid edits.
Null-Text Inversion~\cite{mokady2023nulltext} fixes the practical failure that plain DDIM inversion drifts away from $I$ under classifier-free guidance, by optimizing the \emph{null-text} embeddings used in the inverse pass so that exact reconstruction holds at ordinary guidance scales.
The family needs no full backbone fine-tune for Prompt-to-Prompt or Null-Text and dominates attribute, style, and structure-preserving real-image edits.

\textbf{Instruction-tuned editing.}
A third family avoids inversion entirely by retraining a diffusion editor on paired (before, instruction, after) triplets, mapping $I$ directly to $\hat{I}$ from a natural-language instruction with no explicit mask.
InstructPix2Pix~\cite{brooks2023instructpix2pixlearningfollowimage} built its corpus by pairing GPT-3-generated edit instructions with Prompt-to-Prompt edits of Stable Diffusion samples and fine-tuned Stable Diffusion on the triplets, so that at inference the user supplies only $(I,\text{instruction})$ and the editor localizes the change itself.
MagicBrush~\cite{zhang2024magicbrush} replaced the synthetic data with about 10K manually annotated real-image edits across multi-turn sessions, yielding both a stronger supervised editor and the standard instruction-editing test set.
Emu Edit~\cite{sheynin2024emu} trained a single model on a heterogeneous mixture of \emph{recognition} (segmentation, detection) and \emph{generation} (insertion, removal, style change) tasks with learned task embeddings that route each input, and leads its own benchmark on precision and instruction following.
Instruction editors are the most user-friendly mechanism, no mask and no inversion, at the cost of large supervised retraining.

\textbf{Exemplar / reference-guided editing.}
A visual \emph{exemplar}, an example image or patch encoding the desired content or style, replaces the textual instruction and enters through an auxiliary cross-attention pathway beside the standard text conditioning.
Paint-by-Example~\cite{yang2023paintbyexample} fine-tunes Stable Diffusion to fill $\mathcal{M}$ with the content of an exemplar whose identity is encoded by a CLIP image encoder, with strong exemplar augmentation (cropping, color jitter) at training time to prevent copy-paste behavior and force semantic rather than pixel-level transfer.
IP-Adapter~\cite{ye2023ipadapter} generalizes this into a \emph{plug-in} reference channel: a lightweight decoupled cross-attention head consumes a CLIP image embedding of the exemplar and is inserted into a frozen SD or SDXL backbone, so that any number of reference images can be combined with text prompts and with structural adapters such as ControlNet at inference time.
It is the standard route for subject personalization and style transfer when the desired content is hard to describe in words.

\textbf{Drag-based latent optimization.}
The most interactive family takes user-supplied (source point, target point) pairs and iteratively optimizes the latent so that the features at each source point migrate to the corresponding target while a locality regularizer pins the latent outside the user-marked region.
DragGAN~\cite{pan2023draggan} introduced the recipe on a StyleGAN backbone: a motion-supervision loss tracks the source feature and pulls it toward the target by gradient descent in the $\mathcal{W}^+$ latent space, and a point-tracking step re-localizes the moved feature at each iteration to prevent drift.
DragDiffusion~\cite{shi2023dragdiffusion} ported the recipe to the LDM latent space, running a per-timestep latent-update loop together with an identity-preserving LoRA fine-tune of the U-Net so that both the dragged subject and the unmarked context survive.
Both are inference-time optimizers with latency proportional to the number of drag iterations, and no other family yet offers the same directness of spatial control.

Read across the five families, the mechanisms track the generative substrate: feed-forward inpainters such as LaMa gave way to inversion-based diffusion editors (Prompt-to-Prompt, Imagic), which in turn yielded to supervised instruction-tuned editors (InstructPix2Pix, Emu Edit) as paired data became available at scale.
Orthogonal to that migration, the structural conditioning interfaces of §\ref{subsec:p3_s8_1_2} compose with every mechanism here, so a production pipeline pairs a mechanism with a structural channel rather than choosing one: Blended Diffusion with ControlNet-depth for geometry-aware text-guided edits, or IP-Adapter with DragDiffusion for exemplar-driven drag edits.


\subsubsection{Benchmarks}
\label{subsubsec:p3_s8_1_3_benchmarks}
Editing is evaluated in two regimes: \emph{reference-based}, where a held-out ground-truth region exists (holes synthesized over real images), and \emph{reference-free}, where only constraint satisfaction and human preference are available (instruction editing).
No single suite covers both, so the field keeps a small library of task-specific benchmarks next to the holistic T2I suites of §\ref{subsec:p3_s8_1_1}.

\paragraph{Training Corpora.}
Inpainting draws on Places2~\cite{zhou2017places} and Paris StreetView for scene-level masks, DIV2K for high-resolution patches, and the COCO / ADE20K mask supervision of §\ref{subsec:p3_s8_1_2}~\cite{lin2014microsoft,zhou2017ade20k}.
Constraint-driven local synthesis reuses the structural channels of the previous subsection, RefCOCO / RefCOCO+ / RefCOCOg~\cite{yu2016modeling} and PhraseCut~\cite{wu2020phrasecut} for referring-expression masks, SA-1B masks~\cite{kirillov2023sam}, and COCO edge and depth predictions, with the convention that the mask is supplied or derived from the expression; geometry-guided editing adds depth and normal supervision from NYUv2~\cite{silberman2012nyudv2}, Replica~\cite{straub2019replica}, and synthetic corpora, and insertion with harmonization uses composite datasets such as iHarmony4.
Instruction editing has its own data ecology: the synthetic InstructPix2Pix corpus of Prompt-to-Prompt-generated triplets, the manually annotated MagicBrush triplets~\cite{zhang2024magicbrush}, and the multi-task Emu Edit training mixture~\cite{sheynin2024emu}.

\paragraph{Evaluation Benchmark Suites.}
The holistic suites of §\ref{subsec:p3_s8_1_1} (Table~\ref{tab:p3_s8_1_1_eval_benchmarks}) carry over in part: T2I-CompBench~\cite{huang2023t2icompbench} and GenAI-Bench~\cite{li2024genaibench} score the attribute and compositional fidelity that edited images must also satisfy, and HEIM~\cite{lee2023heim} includes alignment and quality axes that apply to edits.
The dedicated suites are EditBench~\cite{wang2023editbench}, released with Imagen Editor, a curated set of text-guided inpainting prompts over object, attribute, and scene categories that serves as the standard evaluation for text-driven mask inpainting; the MagicBrush test set~\cite{zhang2024magicbrush}, $1{,}053$ manually annotated single- and multi-turn triplets, which is the standard for instruction-tuned editors; and the Emu Edit benchmark~\cite{sheynin2024emu}, which extends instruction editing to per-task subsets (insertion, removal, style change, color change, global change).
Drag-based methods evaluate on small custom test sets curated by~\citet{pan2023draggan} and~\citet{shi2023dragdiffusion}; a community-wide drag-editing benchmark has not yet emerged.

\paragraph{Constraint-specific Metrics.}
Reconstruction quality (PSNR / SSIM / LPIPS against ground truth, with FID / KID for diverse-completion settings) is defined only where held-out content exists.
Region adherence uses mask-restricted CLIPScore or region-level VQAScore to ask whether the edit honors the prompt \emph{inside} $\mathcal{M}$; edit locality measures $\ell_1$ or LPIPS \emph{outside} $\mathcal{M}$ against $I$, lower being better; harmonization is scored by color-statistic divergence, shadow-direction consistency, and the rate at which an off-the-shelf forgery detector flags the composite; identity preservation uses feature-cosine similarity over face or landmark regions for people and CLIP-image similarity for products and characters; and instruction following is paired accuracy, whether $\hat{I}$ honors the instruction better than $I$, judged by human raters or by a VLM over a structured rubric, as is now standard on MagicBrush and Emu Edit.
Change inside $\mathcal{M}$ and preservation outside it are reported together; a method that wins on one and loses on the other has not solved the task.


\subsubsection{Challenges and Trends}
\label{subsubsec:p3_s8_1_3_challenges}
Several problems are specific to editing an existing image.

\textbf{Structure-faithful hallucination in large holes.}
Large irregular holes and rare structures remain the hardest case: inpainters tile a salient local texture across the hole instead of synthesizing appropriate content, straight lines, perspective cues, and silhouettes break at the mask seam, and in-fills that are locally plausible contradict the global scene (a patch of grass inside a room).
Quality degrades sharply with hole size.
LaMa's image-wide Fourier convolutions and RePaint's resampling schedule (§\ref{para:p3_s8_1_3_mechanisms}) are the established responses; retrieval-augmented priors, which search a paired corpus for candidate completions and use them as soft conditioning, are an emerging one at the intersection of editing and retrieval.

\textbf{Edit locality against drift across rounds.}
A single well-calibrated edit does not guarantee a stable session.
Across rounds the unedited background shifts in white balance or saturation (\emph{color drift}), shadows and highlights migrate as if the edit had changed the global illumination (\emph{lighting reflow}), and salient textures from the edited region reappear in nearby unedited regions through attention leakage (\emph{texture replication}); the errors compound even when each round individually meets its locality budget.
Hard attention masking that confines non-edit tokens to $\mathcal{M}$, mask-aware noise schedules that anneal noise inside $\mathcal{M}$ only, and edit-history-aware sampling that conditions on the trajectory of prior edits are the active mitigations.
Inversion-based editors (Null-Text Inversion~\cite{mokady2023nulltext}, Imagic~\cite{kawar2023imagic}) carry a variant of the problem: any residual reconstruction error in the inverted trajectory leaks into the edit and accumulates.

\textbf{Harmonization under complex illumination.}
Inserted content betrays itself when the host scene has structured light: its shadow disagrees with the dominant light direction recoverable from the scene, there is no interreflection between the object and adjacent surfaces, and its albedo is tone-mapped inconsistently with the host, producing the cut-and-paste look that off-the-shelf forgery detectors flag.
Differentiable rendering proxies that turn depth and normals into shading constraints and back-propagate harmonization losses into the editor, illumination maps as an additional conditioning channel, and explicit shadow-synthesis modules are being explored; integration with the 3D-aware methods of §\ref{sec:p3_s8_2} is the natural next step and remains nascent.

\textbf{Language grounding and safety.}
Region-level instructions must be grounded reliably: an edit aimed at one object spills onto neighbors with similar attributes, spatial referring expressions are misparsed (``the left cup'' lands on the right one), and mildly perturbed instructions produce large deviations in the edit trajectory.
Borrowing GLIGEN's gated grounding~\cite{li2023gligen} (§\ref{para:p3_s8_1_2_mechanisms}) to confine edit attention, and using LLM-emitted layouts to harden the grounding before the edit is applied, are the current responses.
The same precision cuts the other way for safety.
A harmful edit (an identity swap inside one bounding box, watermark removal inside one mask) can be confined to a footprint that whole-image forensic detectors miss; large generative models reproduce individual training images under targeted prompting~\cite{somepalli2023diffusion}, which becomes acute when the training set contains copyrighted artwork or identifiable faces and motivates artist-side defenses such as Glaze~\cite{shan2023glaze}, which perturbs source artworks so that downstream fine-tunes mimic a different style; and most pipelines still lack robust per-image watermarking, so a photograph with a small synthesized region cannot be told from a wholly synthetic one.
Watermarking, region-aware safety filters, and licensing frameworks for the underlying corpora are all needed, and none is settled.

\textbf{Evaluation under plausibility ambiguity.}
Ground-truth pixels are often not the only plausible completion: a removed pedestrian leaves a hole that could contain road, curb, or shadow, and any reconstruction metric that compares against one ground truth over-penalizes good alternatives, while PSNR and SSIM rate sharp but contextually wrong in-fills highly.
The field is converging on \emph{paired evaluation}, asking whether $\hat{I}$ honors the instruction better than $I$ under a held-out judge, now standard on MagicBrush~\cite{zhang2024magicbrush} and Emu Edit~\cite{sheynin2024emu}; the judge-prompt sensitivity and benchmark-contamination caveats of §\ref{subsec:p3_s8_1_1} apply here as well.

\textbf{Scalability and interactive latency.}
Inversion-based methods store cross-attention activations at every timestep, multi-step sampling at $4$K and above scales with image area for standard U-Net backbones, and drag editors run an inner optimization loop per operation, so interactive editing lags behind the user's hand by hundreds of milliseconds.
Consistency-style distillation~\cite{song2023consistency,lcm2023} (§\ref{subsec:p3_s8_1_1}) compresses the reverse trajectory to a handful of steps, memory-efficient adaptive patching runs the denoiser only in a window around $\mathcal{M}$, and progressive-refinement schedules return a coarse preview within a perceptual-latency budget and refine in the background.
Closing the gap to the responsiveness of a photo editor is increasingly a systems problem shared with production T2I inference stacks.


% \subsection{Summary}
% \label{subsec:p3_s8_1_summary}

% Table~\ref{tab:p3_s8_1_subfield_comparison} places the three subfields side by side: the conditioning input each consumes, the sampler it instantiates, the mechanisms it surveys, and the benchmarks on which it reports.

% \begin{table}[t]
%     \centering
%     \fontsize{8}{10}\selectfont
%     \setlength{\tabcolsep}{1.0mm}
%     \caption{The three subfields of 2D spatial generation; each sampler is a successive specialization of Eq.~\eqref{eq:t2i_generator}.}
%     \label{tab:p3_s8_1_subfield_comparison}
%     \begin{tabular}{L{0.14\textwidth} L{0.19\textwidth} L{0.19\textwidth} L{0.19\textwidth} L{0.21\textwidth}}
%         \toprule
%         \textbf{Subfield} & \textbf{Input / sampler} & \textbf{Mechanisms} & \textbf{Canonical methods} & \textbf{Benchmarks / metrics} \\
%         \midrule
%         §\ref{subsec:p3_s8_1_1} Text-to-image generation & Text $\mathbf{x}$; $G_\theta:\mathcal{E}\!\to\!\mathcal{Z}$ (Eq.~\eqref{eq:t2i_generator}) & GAN; autoregressive; diffusion / flow matching; backbone, acceleration, and proprietary refinements & StackGAN, AttnGAN; DALL-E, Parti; LDM, Imagen, DiT, SD3 & FID, CLIPScore, HPSv2 / PickScore, VQAScore; GenEval, T2I-CompBench, HRS-Bench, HEIM, WISE \\
%         §\ref{subsec:p3_s8_1_2} Spatial-aware image generation & Text + structure $\mathbf{c}$; $G_\theta:\mathcal{E}\!\times\!\mathcal{C}'\!\to\!\mathcal{Z}$ (Eq.~\eqref{eq:t2i_spatial_objective}) & Side branch; adapter; grounded cross-attention; layout-aware cross-attention; inference-time guidance & ControlNet, T2I-Adapter, GLIGEN, LayoutDiffusion, BoxDiff & Channel adherence (IoU, edge recall, depth RMSE, PCK, relation accuracy); GenEval position; T2I-CompBench spatial \\
%         §\ref{subsec:p3_s8_1_3} Spatial inpainting and editing & Image $I$, mask $\mathcal{M}$, prompt, control; $G_\theta:\mathcal{Z}\!\times\!\{0,1\}^{H\times W}\!\times\!\mathcal{E}\!\times\!\mathcal{C}'\!\to\!\mathcal{Z}$ (Eq.~\eqref{eq:t2i_edit_objective}) & Masked denoising; inversion + rewrite; instruction-tuned; exemplar-guided; drag-based & RePaint, Blended Diffusion; Prompt-to-Prompt, Null-Text; InstructPix2Pix, Emu Edit; IP-Adapter; DragGAN & PSNR / LPIPS, edit locality, harmonization, identity preservation; EditBench, MagicBrush, Emu Edit \\
%         \bottomrule
%     \end{tabular}
% \end{table}

% One move structures the whole section: specialize the conditioning channel of $G_\theta$, from text alone (§\ref{subsec:p3_s8_1_1}) to text plus a structural scaffold (§\ref{subsec:p3_s8_1_2}) to scaffold plus a source image and mask (§\ref{subsec:p3_s8_1_3}).
% Each step is conservative, since the unconditional limit ($\mathbf{c}=\varnothing$ for §\ref{subsec:p3_s8_1_2}; $\mathbf{z}^{\text{in}}=\varnothing,\mathcal{M}=\mathbf{1}$ for §\ref{subsec:p3_s8_1_3}) recovers the preceding subfield, so the three sampler equations form a containment chain rather than a horizontal partition.
% The mechanism axes of the second and third subsections are, moreover, largely orthogonal to the generative core: side branches, adapters, grounded cross-attention, masked denoising, inversion, and drag-based optimization compose with LDM, DiT, and flow-matching backbones interchangeably, and production pipelines routinely stack a spatial-aware injection (Eq.~\eqref{eq:t2i_spatial_objective}) on top of an editing mechanism (Eq.~\eqref{eq:t2i_edit_objective}), so the three subfields are best read as one composable toolbox rather than as three disjoint task families.
% Evaluation has followed the same specialization, from a single FID against COCO for §\ref{subsec:p3_s8_1_1} circa~2020 to a task-specific tuple that adds channel adherence in §\ref{subsec:p3_s8_1_2} and reconstruction quality, edit locality, harmonization, and identity preservation in §\ref{subsec:p3_s8_1_3}.

% The extension of these 2D priors to \emph{volumetric} synthesis (turning a text prompt, an image, or a mask into a textured mesh, a NeRF, or a 3DGS scene) is the subject of §\ref{sec:p3_s8_2}. The conditional latent sampler of Eq.~\eqref{eq:t2i_generator} continues to play a central role there, but now as a \emph{prior to be distilled} via score-distillation sampling, multi-view consistency objectives, or volume-rendering supervision, rather than as the terminal generator.

% \newpage

\section{3D Spatial Reconstruction and Generation}
\label{sec:p3_s8_2}

\FramedInitial{T}{he} 3D spatial reconstruction and generation encompasses the tasks that turn images, video, sensor readings, or high-level prompts into an explicit 3D representation (a textured mesh, a signed-distance or occupancy field, a neural radiance field, or a 3D Gaussian field) of an object, a scene, or a complete navigable environment.
Whereas §\ref{sec:p3_s8_1} taught a 2D semantic-to-pixel before honor spatial structure on the image plane, this section lifts that prior off the plane and into the volumetric substrate on which embodied perception, simulation, and interaction ultimately operate.
This is where appearance and geometry are forced into a single, view-consistent representation: a model must not only produce plausible pixels but place matter correctly in $\mathbb{R}^3$ so that the result can be rendered from arbitrary viewpoints, physically simulated, navigated, and edited. It is the bridge between the flat generative priors of §\ref{sec:p3_s8_1} and the embodied world models, simulators, and digital twins that consume 3D structure downstream.

The five subsections are organized by the direction and source of the map into a shared 3D-representation space $\mathcal{R}$.
§\ref{subsec:p3_s8_2_1} starts with a single asset: 3D object synthesis turns the 2D prior of §\ref{sec:p3_s8_1} into a volumetric generator through 3D-native training, score distillation, multi-view lifting, or feed-forward reconstruction.
§\ref{subsec:p3_s8_2_2} scales that generator to a structured scene of placed objects and a room envelope.
§\ref{subsec:p3_s8_2_3} reverses the direction: reconstruction recovers a representation faithful to sensor observations, and meets synthesis exactly where the observations run out.
§\ref{subsec:p3_s8_2_4} routes through a 360-degree panorama as a structured 2D intermediate that is then lifted to 3D, and §\ref{subsec:p3_s8_2_5} closes the loop by editing an existing representation under user control.
% A Summary (§\ref{subsec:p3_s8_2_summary}) compares the five along shared axes.


\subsection{3D Object Synthesis}
\label{subsec:p3_s8_2_1}


3D object synthesis generates a complete, usable 3D asset (geometry and, optionally, appearance) from \textit{text}, \textit{a single image}, or \textit{a sparse set of views}.
Where reconstruction recovers an object consistent with given observations, synthesis additionally asks for plausible completion of what was never observed, for generative diversity, and for usability in downstream graphics and embodied pipelines (\eg, game engines, simulators, manipulation); the output is a textured mesh, a radiance field, or a set of 3D Gaussians depending on which of fidelity, editability, rendering speed, and toolchain compatibility matters most.
It is also the point at which the \emph{semantic-to-pixel} prior of §\ref{sec:p3_s8_1} (the conditional latent sampler $G_\theta$ of Eq.~\eqref{eq:t2i_generator}) becomes a \emph{semantic-to-3D} prior $G_\theta^{\text{3D}}$: 3D-native generators train on 3D corpora directly, score-distillation pipelines invert the frozen 2D prior under rendering supervision, multi-view pipelines fine-tune it to emit consistent views and lift them, and feed-forward reconstructors and 3D-native latent diffusion exploit the recent object-scale 3D corpora to learn one-shot text- or image-to-3D maps.
Figure~\ref{fig:3D_object_generation} illustrates the text- or single-image-conditioned setting.

\begin{figure}[!ht]
    \centering
    \includegraphics[width=0.89\linewidth]{figures/part3/section8/3D-objects-gen-1.png}
    \caption{Illustration of 3D object generation conditioned on text or a single image (adapted from~\cite{xiang2025trellis}).}
    \label{fig:3D_object_generation}
\end{figure}


\subsubsection{Task Definition}
\label{subsubsec:p3_s8_2_1_task}

We cast 3D object synthesis as a paradigm-agnostic extension of the latent-generator framework of §\ref{subsec:p3_s8_1_1}.
Let $\mathcal{R}$ denote a 3D-representation space taken from $\{\,\textit{mesh},\,\textit{SDF},\,\textit{voxel},\,\textit{triplane},\,\textit{NeRF},\,\text{\textit{3DGS}},\,\textit{structured latent}\,\}$, and let $\mathcal{C}$ denote the conditioning-input space comprising a text prompt $\mathbf{x}\!\in\!\mathcal{X}$ embedded via $\phi$ (as in Eq.~\eqref{eq:t2i_generator}), a single image $I\!\in\!\mathcal{Y}$, or a multi-view set $\{(I_v,\pi_v)\}_{v=1}^V$ pairing images with camera poses $\pi_v$.
A condition encoder $\kappa:\mathcal{C}\!\to\!\mathcal{C}'$ maps each input into a shared conditioning embedding.
The core learning problem is to train a conditional 3D generator
\begin{equation}
G_\theta^{\text{3D}}:\mathcal{C}'\rightarrow \mathcal{R},\qquad \hat{\mathbf{r}}\sim p_\theta^{\text{3D}}\!\big(\mathbf{r}\,\big|\,\kappa(\mathbf{c})\big),
\label{eq:p3_s8_2_1_objective}
\end{equation}
so that the sampled asset $\hat{\mathbf{r}}\in\mathcal{R}$ admits high-fidelity geometry, plausibly completes unobserved regions, and renders consistently across viewpoints.
Two limiting cases locate Eq.~\eqref{eq:p3_s8_2_1_objective} in the literature.
When $\kappa(\mathbf{c})=\varnothing$ the sampler degenerates to \emph{unconditional} 3D generation, the regime of early GAN- and diffusion-based 3D generators trained on category-specific corpora.
At the other extreme, when $\mathbf{c}$ is text-only and $G_\theta^{\text{3D}}$ is realized by distilling the 2D sampler $G_\theta$ of Eq.~\eqref{eq:t2i_generator} (via score-distillation sampling, multi-view consistency objectives, or volume-rendering supervision), one recovers the DreamFusion-style text-to-3D pipeline that has dominated open-vocabulary 3D synthesis since 2022.


\subsubsection{Methodology}
\label{subsubsec:p3_s8_2_1_methodology}

The four paradigms below are four ways of parameterizing and optimizing the sampler $G_\theta^{\text{3D}}$ of Eq.~\eqref{eq:p3_s8_2_1_objective}, taken in the order in which each became prominent.
Figure~\ref{fig:p3_s8_2_1_paradigms} contrasts them, making explicit which components are trained versus frozen and where the 3D representation $\mathcal{R}$ is instantiated.

\begin{figure}[!ht]
    \centering
    \includegraphics[width=0.99\linewidth]{figures/part3/section8/3D_object_synthesis.pdf}
    \caption{Four paradigm families for 3D object synthesis: (P1) 3D-native generation; (P2) distillation from a frozen 2D diffusion prior (SDS); (P3) multi-view generation with a fine-tuned prior, then lifting to 3D; and (P4) learned direct-to-3D routes, by one-pass reconstruction or iterative 3D latent sampling. The families can overlap, notably P1 and the latent route of P4; chair images are illustrative.}
    \label{fig:p3_s8_2_1_paradigms}
\end{figure}

\paragraph{3D-Native Generators.}
\label{para:p3_s8_2_1_native}
The earliest 3D generators instantiate $G_\theta^{\text{3D}}$ as a sampler trained \emph{directly} on a 3D representation, by denoising score matching on voxels, SDFs, or triplanes or, in earlier work, by an adversarial fit to a category-specific 3D corpus $p_{\text{data}}^{\text{3D}}(\mathbf{r})$,
\begin{equation}
\mathcal{L}_{\text{native}}(\theta)
=
\mathbb{E}_{\mathbf{r}_0\sim p_{\text{data}}^{\text{3D}},\,t,\,\epsilon}\Big[\big\|\epsilon - \epsilon_\theta^{\text{3D}}\!\big(\mathbf{r}_t,\,t,\,\kappa(\mathbf{c})\big)\big\|^2\Big],
\quad
\mathbf{r}_t=\alpha_t \mathbf{r}_0 + \sigma_t\,\epsilon,
\label{eq:p3_s8_2_1_native_objective}
\end{equation}
where the noise predictor $\epsilon_\theta^{\text{3D}}$ operates directly on $\mathbf{r}\!\in\!\mathcal{R}$ with no 2D intermediary, so the supervision signal scales with the size of the 3D corpus rather than with 2D image--text pairs.
Because 3D data is heterogeneous, these generators are themselves diverse, having been instantiated on signed distance fields~\cite{gao2022get3d}, voxel grids~\cite{wu2016learning, yang20173d, zhu2018visual}, point clouds~\cite{lin2018learning, huang2020pf}, view sets~\cite{park2017transformation, nguyen2019hologan}, neural radiance fields~\cite{chan2022efficient, chan2021pi, deng2022gram, niemeyer2021giraffe, schwarz2020graf}, triplanes~\cite{chen2023singlestage, shue20233dtriplanetriplane, wang2023rodin}, and 3D Gaussian Splatting~\cite{wewer2024latentsplat, roessle2024l3dg, zou2024triplane}.
The scarcity and monotony of 3D corpora remain the bottleneck: even with sophisticated architectures, generated geometry and texture lack fidelity and stay confined to a handful of canonical categories, which is what pushed the field, around 2022--2023, toward paradigms that use the much larger \emph{2D} corpora as an indirect supervision signal for 3D structure.

\paragraph{2D-Prior Distillation (SDS Family).}
\label{para:p3_s8_2_1_sds}
The second paradigm sidesteps the 3D-data bottleneck by realizing $G_\theta^{\text{3D}}$ as a \emph{per-shape optimization} that distills gradients from a pretrained 2D generator (a CLIP model~\cite{radford2021learning} or a 2D diffusion model~\cite{rombach2022ldm}) trained on web-scale image data~\cite{schuhmann2022laion5b}.
Given the \emph{frozen} 2D noise predictor $\epsilon_\phi$ of Eq.~\eqref{eq:t2i_generator} and a differentiable renderer $\mathrm{Render}(\mathbf{r}_\theta,\pi)$ that images a 3D representation $\mathbf{r}_\theta\!\in\!\mathcal{R}$ under camera pose $\pi$, score distillation updates only the \emph{per-shape} 3D parameters $\theta$ via the gradient
\begin{equation}
\nabla_\theta \mathcal{L}_{\text{SDS}}
=
\mathbb{E}_{t,\,\epsilon,\,\pi}\Big[\,w(t)\,\big(\epsilon_\phi(\mathbf{y}_t^\pi,\,t,\,\phi(\mathbf{x})) - \epsilon\big)\,\tfrac{\partial \mathbf{y}^\pi}{\partial \theta}\,\Big],
\quad
\mathbf{y}^\pi = \mathrm{Render}(\mathbf{r}_\theta,\pi),
\label{eq:p3_s8_2_1_sds_objective}
\end{equation}
where $w(t)$ is a step-dependent weighting and $\mathbf{y}_t^\pi$ is the rendered view perturbed by the diffusion forward process.
The contrast with Eq.~\eqref{eq:p3_s8_2_1_native_objective} is sharp: the 2D prior of §\ref{sec:p3_s8_1} is held \emph{fixed} and no 3D corpus is consumed, while the trainable parameters live in the representation of a \emph{single} asset, hence the per-shape optimization cost paid at inference.
A line of work~\cite{jain2022zerodreamfields, poole2022dreamfusion, lin2023magic3d, wang2023score} applies score-distillation sampling (SDS) or its variants to a NeRF or mesh; DreamFusion~\cite{poole2022dreamfusion} obtains zero-shot text-to-3D by optimizing a NeRF whose multi-view renderings the 2D diffusion prior scores, and GaussianDreamer~\cite{yi2024gaussiandreamer} moves the same recipe onto a 3DGS representation.
These methods need no 3D training data, but because the optimization is rerun for every prompt they are slow, prone to over-saturated color, and subject to the ``multi-face'' (Janus) artifact in which the 2D prior fails to disambiguate the front and back of an object; variational score distillation~\cite{wang2023prolificdreamer} mitigates the over-saturation but not the per-prompt cost.

\paragraph{Multi-View Generation and Lifting.}
\label{para:p3_s8_2_1_mv}
The third paradigm realizes $G_\theta^{\text{3D}}$ as a two-stage pipeline: a fine-tuned 2D diffusion model generates a small set of geometrically consistent views, and a separate reconstructor lifts them to a 3D asset, so that Eq.~\eqref{eq:p3_s8_2_1_objective} factorizes as
\begin{equation}
\hat{\mathbf{r}}
=
\mathcal{F}_\psi\Big(\big\{\hat{\mathbf{y}}_v,\,\pi_v\big\}_{v=1}^V\Big),
\qquad
\hat{\mathbf{y}}_v \sim p_\theta\big(\mathbf{y}_v \,\big|\, \mathbf{y}_{\text{ref}},\,\pi_v,\,\phi(\mathbf{x})\big),
\label{eq:p3_s8_2_1_mv_objective}
\end{equation}
where the view-conditional 2D sampler $p_\theta(\cdot\mid\cdot)$ is a fine-tuned variant of Eq.~\eqref{eq:t2i_generator} that additionally consumes a reference image or text $\mathbf{y}_{\text{ref}}$ and a target camera pose $\pi_v$, and $\mathcal{F}_\psi$ is a (typically feed-forward) multi-view-to-3D reconstructor.
Two differences from Eq.~\eqref{eq:p3_s8_2_1_sds_objective} follow: the 2D prior is \emph{fine-tuned} rather than frozen, and the 3D representation appears only inside the deterministic lift $\mathcal{F}_\psi$, with no per-shape inner loop at inference, which is what lets this paradigm run one to two orders of magnitude faster than SDS at the price of depending on the view-consistency of the synthesized $\hat{\mathbf{y}}_v$.
The first step was to make a pretrained diffusion model pose-aware. Following fine-tuning recipes such as ControlNet~\cite{zhang2023adding} and LoRA~\cite{hu2022lora}, several works~\cite{raj2023dreambooth3d, seo2023let3DFuse} fine-tune for cross-view consistency; Zero-1-to-3~\cite{liu2023zero} opened the open-world single-image setting by training a diffusion model to synthesize a novel view conditioned on an input view and a relative camera pose; MVDream~\cite{shi2023mvdream} and ImageDream~\cite{wang2023imagedream} fix a world-camera coordinate convention from which geometry can be recovered; Magic123~\cite{qian2023magic123} combines the Zero-1-to-3 prior with the 2D prior of Stable Diffusion; One-2-3-45~\cite{liu2023one2345} lifts Zero-1-to-3 views to a mesh; and Consistent123~\cite{lin2024consistent123} and MVD-Fusion~\cite{hu2024mvdfusion} add boundary and depth priors.
Attention then turned to consistency \emph{across} the generated views: SyncDreamer~\cite{liu2023syncdreamer} correlates features across views with a 3D-aware attention, CAT3D and Free3D~\cite{gao2024cat3d, zheng2024free3d} model the multi-view distribution jointly, and Zero123++~\cite{shi2023zero123++} tiles six surrounding views into a single grid image, a trick inherited by One-2-3-45++~\cite{liu2024one2345++}, Instant3D~\cite{li2023instant3d}, InstantMesh~\cite{xu2024instantmesh}, and AR-1-to-3~\cite{zhang2025ar123}.

\paragraph{Feed-Forward Reconstruction and 3D-Native Latent Diffusion.}
\label{para:p3_s8_2_1_lrm}
The fourth group amortizes generation across object-scale 3D corpora rather than optimizing each asset against a frozen 2D prior, in two instantiations.
A \emph{feed-forward reconstructor} trains $G_\theta^{\text{3D}}$ end-to-end to map the conditioning embedding $\kappa(\mathbf{c})$ directly to a 3D representation $\hat{\mathbf{r}}\!\in\!\mathcal{R}$, supervised by differentiable rendering against $V$ ground-truth views $\{\mathbf{y}_v,\pi_v\}$ drawn from the corpus,
\begin{equation}
\mathcal{L}_{\text{LRM}}(\theta)
=
\mathbb{E}_{(\mathbf{c},\,\{\mathbf{y}_v,\pi_v\}_v)\sim p_{\text{data}}^{\text{3D}}}\,
\sum_{v=1}^V \big\|\mathrm{Render}\big(G_\theta^{\text{3D}}(\kappa(\mathbf{c})),\,\pi_v\big) - \mathbf{y}_v\big\|^2.
\label{eq:p3_s8_2_1_lrm_objective}
\end{equation}
A \emph{3D-native latent diffusion} model instead learns a 3D autoencoder pair $(E^{\text{3D}}, D^{\text{3D}})$ with $\mathbf{z}^{\text{3D}}\!=\!E^{\text{3D}}(\mathbf{r})$, trains a denoising model on the compact latent, and decodes via $\hat{\mathbf{r}}=D^{\text{3D}}(\hat{\mathbf{z}}^{\text{3D}})$,
\begin{equation}
\mathcal{L}_{\text{3DL}}(\theta)
=
\mathbb{E}_{\mathbf{z}^{\text{3D}}_0,\,t,\,\epsilon}\Big[\big\|\epsilon - \epsilon_\theta\big(\mathbf{z}^{\text{3D}}_t,\,t,\,\kappa(\mathbf{c})\big)\big\|^2\Big],
\label{eq:p3_s8_2_1_latdiff_objective}
\end{equation}
which is the 2D LDM objective Eq.~\eqref{eq:t2i_ldm_objective} of §\ref{subsec:p3_s8_1_1} with the autoencoder moved from pixels to $\mathcal{R}$.
Both variants avoid the per-shape optimization of Eq.~\eqref{eq:p3_s8_2_1_sds_objective}, but only the reconstructor uses a single prediction pass; latent diffusion still samples iteratively.
The unlocking factor was the arrival of Objaverse and Objaverse-XL~\cite{deitke2023objaverse, deitke2023objaversexl}.
LRM~\cite{hong2023lrm} and its follow-ups~\cite{tang2024lgm, wang2024crm, wang2025moving} show that large transformers trained purely with rendering losses capture object geometry and texture in detail and generalize to unseen images, and a parallel line introduces compact \emph{3D-native latent codes}~\cite{chen20253dtopia, zhang20233dshape2vecset, ye2022structured, yang2025prometheus} whose distribution a latent diffusion model can learn; TRELLIS~\cite{xiang2025trellis}, CLAY~\cite{zhang2024clay}, HiGS-RGBN~\cite{shen2025highGSRGBN}, and Hunyuan3D~\cite{yang2024hunyuan3d} build on this to bring generated assets close to production grade while keeping the open-vocabulary reach that the SDS family pioneered.

Two design choices cut across the four paradigms.
The representation $\mathcal{R}$ (\textit{mesh}, \textit{SDF}, \textit{voxel}, \textit{triplane}, \textit{NeRF}, \textit{3DGS}, or \textit{structured latent}) is largely independent of the paradigm: SDS has been run on NeRFs and on 3DGS, feed-forward reconstructors output triplanes, meshes, or Gaussians, and latent-diffusion stacks decode to almost any explicit representation through a learned VAE.
So is the conditioning modality: the same backbone takes text, a single image, or sparse views by changing only the encoder $\kappa$ and the camera-pose embedding, and production systems chain them (image-to-3D conditioned on a text-to-image preview).
The succession of paradigms itself tracks the accumulation of supervision at each level, from 3D-native corpora to 2D image--text pairs to Zero-1-to-3-style multi-view priors to Objaverse-XL-scale 3D, with the most recent industrial systems sitting in the fourth paradigm while inheriting the multi-view-consistency tricks of the third.


\subsubsection{Benchmarks}
\label{subsubsec:p3_s8_2_1_benchmarks}

Evaluation draws on object-scale training corpora, on dedicated benchmark suites that curate prompts and human ratings, and on metrics that probe alignment, geometry, and appearance separately.

\begin{figure}[!ht]
    \centering
    \includegraphics[width=0.7\linewidth]{figures/part3/section8/objaverse.png}
    \caption{Demonstration of the Objaverse~\cite{deitke2023objaversexl} dataset.}
    \label{fig:objaverse}
\end{figure}

\paragraph{Training Corpora.}
Object-scale 3D corpora drive the feed-forward and latent-diffusion paradigms; Table~\ref{tab:p3_s8_2_1_corpora} lists the ones most commonly used, and Figure~\ref{fig:objaverse} shows representative Objaverse-XL assets.
They span the classical CAD repository ShapeNet~\cite{shapenet2015}, category-bounded real-world scans with accurate geometry and PBR textures (Google Scanned Objects~\cite{downs2022googleGSO}, OmniObject3D~\cite{wu2023omniobject3d} with dense multi-view imagery), the open-vocabulary web-scraped Objaverse and Objaverse-XL~\cite{deitke2023objaverse, deitke2023objaversexl} that unlocked the fourth paradigm, and G-Objaverse~\cite{zuo2024gobjavers}, which adds text-driven synthesized assets between the two.

\begin{table}[t]
    \centering
    \fontsize{8}{10}\selectfont
    \setlength{\tabcolsep}{1.0mm}
    \caption{Representative training corpora for 3D object synthesis. ``Scale'' gives order-of-magnitude object counts as reported in each release.}
    \label{tab:p3_s8_2_1_corpora}
    \begin{tabular}{L{0.22\textwidth} C{0.07\textwidth} C{0.14\textwidth} L{0.20\textwidth} L{0.29\textwidth}}
        \toprule
        \textbf{Dataset} & \textbf{Year} & \textbf{Scale (\#objects)} & \textbf{Acquisition} & \textbf{Primary annotation / modality} \\
        \midrule
        ShapeNet~\cite{shapenet2015} & 2015 & ${\sim}51\mathrm{K}$ & Curated CAD & Semantic category, part labels \\
        Google Scanned Objects~\cite{downs2022googleGSO} & 2022 & ${\sim}1\mathrm{K}$ & Real-world scans & High-fidelity geometry + PBR textures \\
        OmniObject3D~\cite{wu2023omniobject3d} & 2023 & ${\sim}6\mathrm{K}$ & Real-world scans & Dense multi-view + reconstructions \\
        Objaverse~\cite{deitke2023objaverse} & 2023 & ${\sim}800\mathrm{K}$ & Web-scraped & Open-vocabulary text tags \\
        Objaverse-XL~\cite{deitke2023objaversexl} & 2023 & ${\sim}10\mathrm{M}$ & Web-scraped at scale & Open-vocabulary text tags \\
        G-Objaverse~\cite{zuo2024gobjavers} & 2024 & ${\sim}280\mathrm{K}$ & Synthesized + curated & Text-driven assets, rendered views \\
        \bottomrule
    \end{tabular}
\end{table}

\paragraph{Evaluation Benchmark Suites.}
Recent benchmarks have moved to \emph{multi-dimensional} assessment beyond plain text--image alignment.
GT23D-Bench~\cite{cai2024gt23d} curates a high-fidelity dataset and a suite of 3D-aware metrics spanning alignment and visual quality; MATE-3D~\cite{zhang2025mate3d} provides designed prompt categories with large-scale human annotations across several perceptual dimensions and introduces HyperScore as a learned evaluator; T23D-QA~\cite{li2025t23d} adopts a question-answer protocol that scores geometry, texture, and alignment independently; and Gen3DEval~\cite{maiti2025gen3deval} fine-tunes a vision-language model for 3D quality assessment, scoring text fidelity, appearance, and surface quality from rendered surface normals.
Together with the 3D-relevant subsets of the 2D suites collected in §\ref{subsec:p3_s8_1_1} (Table~\ref{tab:p3_s8_1_1_eval_benchmarks}), they form the de facto evaluation protocol for 3D object synthesis.

\paragraph{Constraint-specific Metrics.}
Because an asset is judged from many viewpoints and under different downstream uses, evaluation reports complementary metrics along three axes.
\textbf{Text-3D alignment} renders multi-view images from the asset and measures text--image similarity with CLIP-like encoders, aggregated across views (mean, worst-case, or view-consistency variants); pure CLIP alignment can be misled by texture ``shortcuts'' and does not guarantee geometric correctness, which motivates the 3D-aware alignment scores of GT23D-Bench~\cite{cai2024gt23d} and learned evaluators such as HyperScore~\cite{zhang2025mate3d}.
\textbf{Geometry quality} is measured against ground-truth shapes, where they exist, by Chamfer Distance (CD)~\cite{huang2025spar3d}, F-score at varying thresholds, normal consistency, and volumetric IoU for watertight meshes; for text-to-3D without ground truth it falls back on multi-view depth and normal consistency, silhouette consistency, and learned 3D-quality predictors.
\textbf{Appearance and multi-view consistency} are measured by perceptual distances (LPIPS) between rendered views and references where available, by view-consistency measures (variance across views under fixed lighting and camera trajectories), and by human studies, which recent benchmarks collect as multi-dimensional ratings to calibrate the automatic metrics~\cite{zhang2025mate3d, li2025t23d}.

The paradigm succession of §\ref{subsubsec:p3_s8_2_1_methodology} shows up in two measurable quantities.
Generation time has fallen by about three orders of magnitude: DreamFusion~\cite{poole2022dreamfusion} needs roughly 1.5 hours per asset and Magic3D~\cite{lin2023magic3d} about 40 minutes, whereas LRM~\cite{hong2023lrm} reconstructs a shape in about 5 seconds and InstantMesh~\cite{xu2024instantmesh} completes image-to-3D in around 10 seconds.
Image-to-3D geometry is now routinely scored on Google Scanned Objects~\cite{downs2022googleGSO} with Chamfer Distance (lower is better) and F-score at $\tau{=}0.2$ (higher is better), where InstantMesh~\cite{xu2024instantmesh} reaches CD~$\approx 0.18$ and FS~$\approx 0.88$, ahead of TripoSR~\cite{tochilkin2024triposr} ($0.217$ / $0.843$), CRM~\cite{wang2024crm} ($0.252$ / $0.787$), and LGM~\cite{tang2024lgm} ($0.345$ / $0.671$).
Coarse text-to-3D alignment, by contrast, is saturated: CLIP-style scores no longer separate the strongest systems, so Hunyuan3D~2.0~\cite{hunyuan3d22025tencent} reports ULIP- and Uni3D-based image- and text-to-shape similarity backed by a study in which 50 participants rated 300 unselected results, and learned evaluators (HyperScore on MATE-3D~\cite{zhang2025mate3d}, T23D-QA~\cite{li2025t23d}) score geometry and texture on separate axes.
The reconstruction metrics carry their own caveat: CD and F-score saturate on coarse geometry and can reward visually plausible but physically thin assets, so reported gains are credible mainly where geometric, perceptual, and human ratings agree.
Commercial text- and image-to-3D services now ship these capabilities with PBR export to game engines, alongside research systems such as Meta 3D Gen~\cite{bensadoun2024meta3dgen} and the open-vocabulary 3D-native-latent systems CLAY~\cite{zhang2024clay}, TRELLIS~\cite{xiang2025trellis}, and Hunyuan3D~\cite{yang2024hunyuan3d, hunyuan3d22025tencent}.


\subsubsection{Challenges and Trends}
\label{subsubsec:p3_s8_2_1_challenges}

Four problems remain open, and they depend on one another: the multi-view consistency of the asset, the representation it lands in, whether progress on either can be measured reliably, and the data that ultimately drives all three.

\paragraph{Multi-View and 3D Consistency under Open-Vocabulary Prompts.}
Enforcing consistent geometry and texture across viewpoints while preserving fine detail remains difficult in all four paradigms, and the failures have names.
The \emph{Janus artifact} renders both the front and the back of an object as a frontal face, most acutely in SDS pipelines, where the 2D prior cannot disambiguate camera direction; \emph{view drift} makes adjacent views of a multi-view stack disagree on geometry, a residual concern in Zero-1-to-3-style cascades that SyncDreamer's cross-view attention and Zero123++'s tiled six-view trick only partly remove; and \emph{texture seams} appear where neighboring views' color and lighting cannot be merged into a globally coherent UV map, a persistent weakness of feed-forward reconstructors.
These failures motivate 3D-native backbones with \emph{structured latents}~\cite{xiang2025trellis} that share parameters across views by construction, and benchmarks that score multi-view consistency rather than a single rendered view~\cite{zhang2025mate3d}.

\paragraph{Representation Trade-offs between Fidelity, Editability, and Deployability.}
Even a consistent asset must \emph{land somewhere} a downstream tool can consume, and the choice of representation is not neutral.
Meshes with UV/PBR textures are what game engines and simulators want but lose detail to discretization; radiance fields and Gaussians give high visual fidelity and fast rendering but require conversion, which introduces its own failures: \emph{topology breakage} during marching-cubes-style meshing, \emph{baked illumination} that fuses lighting into albedo when converting from a NeRF or 3DGS, and \emph{texture bleeding} across UV seams.
The response is \emph{unified latent representations} that decode to several formats~\cite{xiang2025trellis}, so that a single generator serves diverse downstream requirements and the choice of $\mathcal{R}$ is deferred to inference.

\paragraph{Scalable, Reliable Evaluation and Human-Aligned Metrics.}
These representation choices complicate evaluation, because the metric panel a paper reports tends to depend on which format its generator targets, and the automatic metrics remain imperfect for geometry plausibility and asset usability.
CLIP-style alignment is fooled by \emph{texture shortcuts} that ignore the underlying shape; Chamfer Distance saturates on coarse geometry while penalizing semantically irrelevant fine detail; and LPIPS on rendered views cannot distinguish a physically plausible asset from a visually plausible cardboard cut-out.
Multi-dimensional evaluators trained on large human-annotation sets (\eg, HyperScore on MATE-3D~\cite{zhang2025mate3d}) and diagnostic benchmarks (\eg, GT23D-Bench~\cite{cai2024gt23d}, T23D-QA~\cite{li2025t23d}) that score alignment, geometry, and texture on independent axes are the current answer, so that progress on one axis is no longer masked by averaging.

\paragraph{Data Scaling and Asset Readiness.}
Underlying consistency, representation, and evaluation alike is the supply of training data and the industrial scale at which it can be exploited.
Large-scale corpora such as Objaverse-XL~\cite{deitke2023objaversexl} and industrial-scale systems such as Hunyuan3D~\cite{hunyuan3d22025tencent, hunyuan3d2025hunyuan3domni} and Meta 3D Gen~\cite{bensadoun2024meta3dgen} point toward \emph{high-resolution, engine-ready assets} with controllable materials, relightable textures, and post-processing for watertightness and topology.
The question has shifted from \emph{whether} a paradigm can be trained at scale, which the fourth paradigm has settled, to \emph{which slice} of the 3D-asset distribution the available corpora still under-represent: articulated objects, semi-transparent or refractive materials, and topologically non-trivial assemblies are visibly thinner in current open corpora than in production pipelines, and closing this gap is the explicit target of recent industrial releases.




\subsection{3D Scene Synthesis}
\label{subsec:p3_s8_2_2}


Scene-level synthesis extends 3D generation from a single asset $\mathbf{r}\!\in\!\mathcal{R}$ to an entire environment: a clutter of objects whose positions, orientations, and pairwise relations together compose a navigable room or open-world scene.
Where §\ref{subsec:p3_s8_2_1} models the conditional distribution of one object given a prompt, scene synthesis must jointly model (i) a variable number of objects, (ii) their global $\mathrm{SE}(3)$ placements, (iii) their pairwise spatial and functional relations, and (iv) the envelope geometry (floor, walls, ceiling) that closes the scene; the output is a structured set rather than a vector in $\mathcal{R}$, and the conditioning broadens from text to layouts, scene graphs, sparse calibrated views, and full camera trajectories.
The result must be navigable, not only renderable, which makes the problem harder than object synthesis along three axes at once: geometric coherence across the clutter, semantic consistency between related objects, and view generalization along an $\mathrm{SE}(3)$ trajectory rather than a fixed look-at orbit.
Four families answer it: \emph{layout- and scene-graph-conditioned} generators that produce the structured set from a symbolic prior; \emph{view-to-3D} synthesizers that regress or sample a continuous radiance or Gaussian field from sparse calibrated views; \emph{view-to-view} synthesizers that skip the explicit 3D representation and learn a conditional map from observed to novel views; and \emph{video-model-based} synthesizers that treat a pretrained video diffusion model as an implicit 3D prior with the camera trajectory as an additional conditioning channel.


\subsubsection{Task Definition}
\label{subsubsec:p3_s8_2_2_task}

We cast 3D scene synthesis as a scene-level extension of the object-synthesis framework of §\ref{subsec:p3_s8_2_1}.
Let $\mathcal{S}$ denote the scene-output space whose elements
$\mathbf{s} = \big(\{\mathbf{r}_i\}_{i=1}^N,\, \{\mathbf{T}_i\}_{i=1}^N,\, \mathbf{G},\, \mathcal{E}_{\text{room}}\big)$
package $N$ object representations $\mathbf{r}_i\!\in\!\mathcal{R}$ (with $\mathcal{R}$ the per-object space of §\ref{subsec:p3_s8_2_1}), their global poses $\mathbf{T}_i\!\in\!\mathrm{SE}(3)$, a relational graph $\mathbf{G}$ over the object set, and an envelope geometry $\mathcal{E}_{\text{room}}$ encoding the room shell (walls, floor, ceiling).
The conditioning space broadens to
$\mathcal{C}^{\text{scene}} \supseteq \{\,\text{text }\mathbf{x},\,\text{layout }\mathbf{L},\,\text{scene graph }\mathbf{G}_{\text{in}},\,\text{sparse views }\{(I_v,\pi_v)\}_v,\,\text{camera trajectory }\{\pi_v\}_v\,\}$,
with a condition encoder $\kappa^{\text{scene}}:\mathcal{C}^{\text{scene}}\!\to\!\mathcal{C}^{\text{scene}\prime}$ mapping each input mode into a shared embedding.
The core learning problem is to train a conditional scene generator
\begin{equation}
G_\theta^{\text{scene}}:\mathcal{C}^{\text{scene}\prime}\rightarrow \mathcal{S},
\qquad
\hat{\mathbf{s}}\sim p_\theta^{\text{scene}}\!\big(\mathbf{s}\,\big|\,\kappa^{\text{scene}}(\mathbf{c})\big),
\label{eq:p3_s8_2_2_objective}
\end{equation}
so that the sampled scene $\hat{\mathbf{s}}\!\in\!\mathcal{S}$ admits coherent inter-object geometry, semantically plausible relations, and view-consistent renderings under arbitrary $\mathrm{SE}(3)$ trajectories.
Two limiting cases fix the scope of Eq.~\eqref{eq:p3_s8_2_2_objective}.
When $N\!=\!1$, $\mathbf{G}\!=\!\varnothing$, and $\mathcal{E}_{\text{room}}\!=\!\varnothing$, the scene degenerates to a single asset and the sampler recovers Eq.~\eqref{eq:p3_s8_2_1_objective}: scene synthesis is a strict generalization of object synthesis.
At the opposite extreme, when $\mathcal{S}$ is instantiated as an ordered sequence of views $\{I_v\}_{v=1}^V$ \emph{without an explicit 3D substrate}, the sampler reduces to a camera-trajectory-conditioned video generator (the regime of §\ref{para:p3_s8_2_2_video} below), so that ``scene synthesis'' and ``consistent video generation'' are two boundary regimes of the same problem.


\subsubsection{Methodology}
\label{subsubsec:p3_s8_2_2_methodology}

Each of the four paradigms is a different factorization of $G_\theta^{\text{scene}}$ in Eq.~\eqref{eq:p3_s8_2_2_objective}, from sparse structural controls (layouts and scene graphs) through continuous 3D representations to video-based synthesis; Figure~\ref{fig:p3_s8_2_2_paradigms} previews them against a common scene.

\begin{figure}[!ht]
    \centering
    \includegraphics[width=\linewidth]{figures/part3/section8/Four_paradigm_families_for_3D_scene_synthesis.pdf}
    \caption{Four paradigm families for 3D scene synthesis (Eq.~\eqref{eq:p3_s8_2_2_objective}): (a) layout- or scene-graph-conditioned placement of objects, (b) view-to-3D reconstruction into a radiance field or 3D Gaussians, (c) view-to-view synthesis with no explicit 3D, and (d) video-model-based synthesis along a camera path. The families are not mutually exclusive, and view consistency does not guarantee a valid 3D asset; room images are illustrative.}
    \label{fig:p3_s8_2_2_paradigms}
\end{figure}

\paragraph{Layout and Scene-graph Conditioned Methods.}
\label{para:p3_s8_2_2_layout}
A major line of research generates 3D environments from high-level structural priors such as layouts $\mathbf{L}$ or scene graphs $\mathbf{G}_{\text{in}}$, instantiating Eq.~\eqref{eq:p3_s8_2_2_objective} by materializing the structured tuple as an explicit object set,
\begin{equation}
\hat{\mathbf{s}} = \big\{(\hat{\mathbf{r}}_i,\,\hat{\mathbf{T}}_i)\big\}_{i=1}^N,
\qquad
\hat{\mathbf{s}}\sim p_\theta\!\Big(\{(\mathbf{r}_i,\mathbf{T}_i)\}_i \,\Big|\, \kappa(\mathbf{L} \text{ or } \mathbf{G}_{\text{in}}),\,\phi(\mathbf{x})\Big),
\label{eq:p3_s8_2_2_layout_objective}
\end{equation}
where the joint over the $N$ slots and their poses is factorized autoregressively (a transformer over object slots) or as a denoising diffusion over the slot tensor.
SceneFormer~\cite{wang2021sceneformer} explored autoregressive modeling of room layouts and furniture placement, and DiffuScene~\cite{tang2024diffuscene} introduced diffusion over object arrangements, improving realism and diversity.
Recent work adds language and visual reasoning for compositional control: LayoutGPT~\cite{feng2023layoutgpt} plans the layout with an LLM before neural scene construction, and LayoutVLM~\cite{sun2025layoutvlm} optimizes layout consistency under multimodal alignment.
Scene-graph approaches (GraphDreamer~\cite{gao2024graphdreamer}, CommonScenes~\cite{zhai2023commonscenes}, EchoScene~\cite{zhai2024echoscene}) use relational reasoning to keep placements plausible, and InstructScene~\cite{lin2024instructscene} adds natural-language conditions for controllable scene creation.
Beyond indoor settings, CityCraft~\cite{deng2024citycraft} scales layout synthesis to the city level, and GALA3D~\cite{zhou2024gala3d} couples layout-guided diffusion with Gaussian-based rendering for photorealistic multi-object environments.

\paragraph{View-to-3D Representation Synthesis.}
\label{para:p3_s8_2_2_view3d}
The second paradigm trades the explicit structured set for a \emph{continuous} 3D representation, recovering the scene as a radiance field or 3D Gaussian collection from sparse calibrated views; the scene tuple collapses into a single field,
\begin{equation}
\hat{\mathcal{R}}^{\text{scene}} = \mathcal{F}_\psi\!\Big(\big\{(I_v,\pi_v)\big\}_{v=1}^V\Big),
\qquad \hat{\mathbf{s}}\equiv \hat{\mathcal{R}}^{\text{scene}},
\label{eq:p3_s8_2_2_view3d_objective}
\end{equation}
where $\mathcal{F}_\psi$ is either a deterministic feed-forward regressor or a generative sampler over the field.
This scales the multi-view-and-lift recipe of §\ref{para:p3_s8_2_1_mv} from one asset to a full scene by replacing the per-object reconstructor with a scene-scale one.
Deterministic regressors such as PixelNeRF~\cite{yu2021pixelnerf}, PixelSplat~\cite{charatan2024pixelsplat}, DepthSplat~\cite{xu2024depthsplat}, and MVSplat~\cite{chen2024mvsplat} / MVSplat360~\cite{chen2024mvsplat360} stay strictly faithful to the input views and are reconstruction methods, treated in §\ref{para:p3_s8_2_3_feedforward}.
The generative variant samples the field instead: DreamScene~\cite{li2024dreamscene} introduces a diffusion prior over Gaussian fields for text-conditioned scene creation, and Generative 3DGS~\cite{Schwarz2025ggs} generates 3D Gaussians for scenes with video priors.

\paragraph{View-to-View Synthesis (Regression and Generative).}
\label{para:p3_s8_2_2_viewview}
The third paradigm skips the explicit 3D representation altogether and formulates novel view synthesis as a direct conditional 2D-to-2D map from observed to target views,
\begin{equation}
\hat{I}_{\pi^*}\sim p_\theta\!\Big(I_{\pi^*}\,\Big|\,\big\{(I_v,\pi_v)\big\}_{v=1}^V,\,\pi^*\Big),
\label{eq:p3_s8_2_2_viewview_objective}
\end{equation}
where $\pi^*$ is the target camera pose and no intermediate 3D representation is materialized.
Regression-based approaches learn a deterministic map: SRT~\cite{sajjadi2022scene} encodes the scene into latent tokens that support interpolation and extrapolation across viewpoints, and LVSM~\cite{jin2025lvsm} predicts novel views with a large feed-forward transformer and minimal 3D inductive bias.
Generative approaches model the conditional distribution over unseen views, which permits stochastic synthesis under ambiguity and large viewpoint changes: Photoconsistent-NVS~\cite{yu2023long} and SEVA~\cite{zhou2025stable} use diffusion to improve diversity and robustness along extended camera trajectories, and CausNVS~\cite{kong2025causnvs} generates views autoregressively and causally, modeling inter-view dependencies explicitly to stay coherent across complex viewpoint transitions.

\paragraph{Video-model Based Scene Synthesis.}
\label{para:p3_s8_2_2_video}
The fourth paradigm adds \emph{temporal} coherence to the viewpoint variety of the view-to-view family by exploiting large video generation models as spatiotemporal priors; the scene tuple is materialized as a frame sequence,
\begin{equation}
\{\hat{I}_v\}_{v=1}^V \sim p_\theta^{\text{video}}\!\Big(\{I_v\}_{v=1}^V \,\Big|\, \{\pi_v\}_{v=1}^V,\,\phi(\mathbf{x})\Big),
\label{eq:p3_s8_2_2_video_objective}
\end{equation}
treating the scene as the implicit 3D structure of a coherent video and the camera trajectory $\{\pi_v\}$ as an explicit conditioning channel.
Where the feed-forward reconstructors of §\ref{para:p3_s8_2_1_lrm} output a static asset in one pass, the video paradigm trades static 3D for temporally coherent multi-frame output that mirrors the navigability requirement of a scene more directly.
CameraCtrl~\cite{he2024cameractrl} and MotionCtrl~\cite{wang2024motionctrl} condition video generation on camera trajectories directly; ViewCrafter~\cite{yu2024viewcrafter} and Gen3C~\cite{ren2025gen3c} add geometric guidance by unprojecting observed views into an intermediate 3D representation and reprojecting it to the target viewpoints, which improves pose fidelity; and CAT4D~\cite{wu2025cat4d} and 4DiM~\cite{watson2024controlling} train jointly on posed static multi-view data and on temporal video data, so that viewpoint variation and temporal dynamics are modeled together for 4D scene synthesis.


\subsubsection{Benchmarks}
\label{subsubsec:p3_s8_2_2_benchmarks}

The four paradigms draw on different training corpora and report on different metric panels, each inherited from a precursor task.

\paragraph{Training Corpora.}
Table~\ref{tab:p3_s8_2_2_corpora} lists the corpora in use.
Layout-conditioned methods rely overwhelmingly on 3D-FRONT~\cite{fu20213dfront}, usually combined with its companion 3D-FUTURE~\cite{fu20213dfuture} furniture library for instance-level geometry.
View-based methods (P2 / P3) draw on real-world scans of residential interiors (ScanNet~\cite{dai2017scannet}, Matterport3D~\cite{chang2017matterport3d}), on photorealistic synthetic corpora with ground-truth depth and normals (Hypersim~\cite{roberts2021hypersim}), and on Replica~\cite{straub2019replica} for compact high-fidelity indoor benchmarks.
Video-model-based methods (P4) inherit the web-scale supervision of their underlying video prior, far larger than any 3D corpus but without explicit 3D supervision, supplemented by camera-pose-annotated subsets distilled from ScanNet, RealEstate10K~\cite{zhou2018stereo}, or hand-labeled trajectories.

\begin{table}[t]
    \centering
    \fontsize{8}{10}\selectfont
    \setlength{\tabcolsep}{1.0mm}
    \caption{Representative training corpora for 3D scene synthesis.}
    \label{tab:p3_s8_2_2_corpora}
    \begin{tabular}{L{0.20\textwidth} C{0.07\textwidth} C{0.14\textwidth} L{0.22\textwidth} L{0.29\textwidth}}
        \toprule
        \textbf{Dataset} & \textbf{Year} & \textbf{Scale} & \textbf{Acquisition} & \textbf{Primary use / paradigm} \\
        \midrule
        ScanNet~\cite{dai2017scannet} & 2017 & ${\sim}1.5\mathrm{K}$ scenes & Real-world RGB-D scans & Sparse-view supervision (P2 / P3) \\
        Matterport3D~\cite{chang2017matterport3d} & 2017 & ${\sim}90$ buildings & Real-world RGB-D scans & Whole-house multi-view (P2 / P3) \\
        Replica~\cite{straub2019replica} & 2019 & ${\sim}18$ scenes & Real-world high-fidelity scans & Photorealistic indoor benchmarks (P2 / P3) \\
        Hypersim~\cite{roberts2021hypersim} & 2021 & ${\sim}460$ scenes & Photorealistic synthetic & Ground-truth depth / normals (P2 / P3) \\
        3D-FRONT~\cite{fu20213dfront} & 2021 & ${\sim}19\mathrm{K}$ rooms & Synthetic indoor layouts & Layout / object-arrangement supervision (P1) \\
        3D-FUTURE~\cite{fu20213dfuture} & 2021 & ${\sim}10\mathrm{K}$ assets & Synthetic furniture library & Companion to 3D-FRONT (P1) \\
        \bottomrule
    \end{tabular}
\end{table}

\paragraph{Evaluation Benchmark Suites.}
Layout-conditioned methods report the metrics introduced with ATISS on 3D-FRONT (object-category frequency, pairwise-relation accuracy, and FID on rendered top-down or perspective views).
View-to-3D and view-to-view methods adopt the NeRF-era panel of PSNR / SSIM / LPIPS on held-out views of ScanNet, RealEstate10K~\cite{zhou2018stereo}, or DL3DV-10K~\cite{ling2024dl3dv}.
Video-model-based methods add FVD and a camera-trajectory-adherence metric, the mean translation and rotation error between requested and recovered camera poses.

\paragraph{Constraint-specific Metrics.}
Four axes mirror the scene tuple of Eq.~\eqref{eq:p3_s8_2_2_objective}: \textbf{layout adherence} (object-count agreement, pairwise-relation accuracy on $\mathbf{G}$, and bounding-box IoU between requested and synthesized placements); \textbf{geometric fidelity} where ground truth exists (Chamfer Distance, F-score, and depth and normal consistency against held-out views); \textbf{photometric quality} (PSNR / SSIM / LPIPS against ground-truth views, FID / KID against held-out scene distributions, and FVD for video outputs); and \textbf{camera-pose adherence}, most relevant for P4 (translation and rotation error of poses recovered by structure-from-motion or a learned pose regressor).
A paper reports the subset of axes its paradigm can be scored on, which is why comparison across paradigms remains difficult.
% (§\ref{subsec:p3_s8_2_summary}).


\subsubsection{Challenges and Trends}
\label{subsubsec:p3_s8_2_2_challenges}

Scene synthesis meets the problems of object synthesis on a structured set: more conditioning variables, more output components, and longer trajectories enter the cost at once.

\paragraph{Cross-Paradigm Inconsistency.}
The four paradigms solve subtly different problems, and each leaves a gap when used alone: a layout-conditioned generator can satisfy the symbolic prior $\mathbf{G}_{\text{in}}$ yet produce a non-traversable static scene, and a view-to-3D generator can produce a navigable Gaussian field that ignores the requested object semantics.
The gaps surface as \emph{layout-adherence drift}, where the generated NeRF or 3DGS no longer satisfies the input layout once the user roams outside the trained view envelope; \emph{object drop}, where small or thin objects from $\mathbf{L}$ are silently omitted during continuous-field synthesis; and \emph{geometry flattening}, where a multi-object composition collapses into a single textured slab when the 2D or video prior dominates.
Hybrids such as GALA3D~\cite{zhou2024gala3d} and ViewCrafter~\cite{yu2024viewcrafter} chain a layout or geometry prior with a Gaussian or video renderer to keep both halves honest, and future systems will more likely compose the four paradigms than choose among them.

\paragraph{Scaling Beyond Indoor Rooms.}
Almost all paradigms remain anchored to furnished indoor rooms because that is where the training corpora live: 3D-FRONT~\cite{fu20213dfront} alone supplies the layout supervision for the dominant P1 methods, and ScanNet, Matterport3D, and Hypersim dominate P2 / P3.
The result is an outdoor and urban scale gap that CityCraft~\cite{deng2024citycraft} is among the few systems to address, a category imbalance with many furniture classes but few terrain, vegetation, or industrial ones, and an open-world domain shift under which methods trained on residential interiors fail on industrial sites, excavations, or invented environments.
Web-scraped open-vocabulary scene corpora and hybrid synthetic-real pipelines that blend procedurally generated environments with captured scans are the active responses, both for breadth and to control the long tail of the scene-class distribution.

\paragraph{Camera-Trajectory Coherence under Long-Horizon Synthesis.}
The video paradigm (§\ref{para:p3_s8_2_2_video}) excels at short camera moves but degrades on long, unbounded trajectories that require persistent world memory.
Long trajectories expose \emph{temporal blur}, where textures shimmer or drift between frames; \emph{scene amnesia}, where objects reappear with different geometry when the camera revisits a location; and \emph{pose-conditioning collapse}, where the model learns to ignore the camera signal in favor of the marginal video distribution.
ViewCrafter~\cite{yu2024viewcrafter}'s unprojection and reprojection scaffold, Gen3C~\cite{ren2025gen3c}'s 3D cache, and CAT4D~\cite{wu2025cat4d}'s joint multi-view and video training all inject an explicit 3D inductive bias back into the video prior, but the trade-off between the video prior's perceptual quality and the geometric prior's rigidity remains an open design tension.

\paragraph{Evaluation.}
A P1 layout-adherence score and a P3 PSNR cannot be ranked against each other, no single dataset evaluates layout faithfulness, geometric fidelity, photometric quality, and pose adherence jointly, and no human-aligned evaluator analogous to HyperScore~\cite{zhang2025mate3d} has been validated on scene outputs.
% ; the cross-subfield picture is drawn in §\ref{subsec:p3_s8_2_summary}.

\paragraph{Trends: Toward 4D and World Models.}
Scene synthesis targets static geometry conditioned on camera poses, but the same formulation is being extended toward 4D synthesis and action-conditioned world modeling: Dreamer~\cite{hafner2019dreamer}, GameNGen~\cite{valevski2024diffusion}, Diamond~\cite{alonso2024diffusionworldmodelingvisual}, RTFM~\cite{worldlabs_rtfm}, and Genie3~\cite{genie3} move from static view or video synthesis toward open-ended environment generation in which spatial structure, temporal evolution, and interaction are modeled together.
From the standpoint of this subsection, scalable, temporally coherent video generation with explicit viewpoint control is the foundation those systems build on; they are treated as world models in §\ref{sec:p3_s8_6}.



\subsection{3D Scene Reconstruction}
\label{subsec:p3_s8_2_3}

3D reconstruction recovers a digital 3D representation of a scene or object from sensor observations such as RGB images, video, depth maps, or point clouds, and has long been a foundational problem in computer vision, graphics, and immersive systems such as augmented reality (AR), virtual reality (VR), and digital twins.
Its difference from the synthesis tasks of §\ref{subsec:p3_s8_2_1} and §\ref{subsec:p3_s8_2_2} is one of \emph{intent}: synthesis is open-ended generation in which plausibility and diversity dominate, whereas reconstruction must explain the input views faithfully before any inference about unobserved regions is attempted.
Reconstruction is where perception becomes geometry: it turns raw sensor streams into an explicit, metric 3D substrate that simulators, AR / VR systems, digital twins, and embodied agents can reason over, manipulate, and render, and it supplies the ground-truth geometry against which the generative tasks of §\ref{sec:p3_s8_2} are ultimately measured.
Three families follow: \emph{iterative reconstruction}, which solves a per-scene optimization against the input views; \emph{feed-forward reconstruction}, which amortizes that cost into a learned map from observations to a 3D representation; and \emph{complete reconstruction}, which generatively fills the occluded regions and so makes the boundary with synthesis explicit.
Figure~\ref{fig:3D_reconstruction} previews representative reconstruction tasks.

\begin{figure}[!ht]
    \centering
    \includegraphics[width=0.99\linewidth]{figures/part3/section8/3D-reconstruction.png}
    \caption{Illustrations of 3D reconstruction, adapted from~\cite{zhang2025advances}.}
    \label{fig:3D_reconstruction}
\end{figure}


\subsubsection{Task Definition}
\label{subsubsec:p3_s8_2_3_task}

We cast 3D scene reconstruction as the observation-conditioned counterpart of the generative frameworks of §\ref{subsec:p3_s8_2_1} and §\ref{subsec:p3_s8_2_2}.
Let $\mathcal{O}$ denote the space of multi-sensor observations.
Each observation $\mathbf{O}\in\mathcal{O}$ is represented as a finite collection of sensor measurements.
Depending on the sensing setup, $\mathbf{O}$ may consist of multi-view RGB images $\{I_v\}_{v=1}^{V}$, posed images $\{(I_v,\pi_v)\}_{v=1}^{V}$, RGB-D or RGB-normal observations $\{(I_v,d_v)\}_{v=1}^{V}$, point clouds, LiDAR scans, or their combinations.
Here, $I_v$ denotes the RGB image from view $v$, $\pi_v$ denotes the corresponding calibrated camera pose, and $d_v$ denotes an associated depth or normal map.
Let $\mathcal{R}^{\text{rec}}$ denote the output 3D-representation space, taken from the same $\mathcal{R}$ as in §\ref{subsec:p3_s8_2_1}, including mesh, SDF, voxel, triplane, NeRF, 3DGS, point map, structured latent.
The core learning problem is to obtain a reconstruction operator
\begin{equation}
G_\theta^{\text{rec}}:\mathcal{O}\rightarrow \mathcal{R}^{\text{rec}},
\qquad
\hat{\mathbf{r}}^{\text{rec}} = G_\theta^{\text{rec}}(\mathbf{O}) \quad \text{(deterministic)}
\quad \text{or} \quad
\hat{\mathbf{r}}^{\text{rec}}\sim p_\theta^{\text{rec}}\!\big(\mathbf{r}^{\text{rec}}\,\big|\,\mathbf{O}\big) \quad \text{(generative)},
\label{eq:p3_s8_2_3_objective}
\end{equation}
that minimizes a reconstruction loss against the input views,
\begin{equation}
\mathcal{L}_{\text{rec}}\big(\hat{\mathbf{r}}^{\text{rec}},\mathbf{O}\big)
\;=\;
\sum_{(I_v,\pi_v)\in\mathbf{O}} \big\|\mathrm{Render}(\hat{\mathbf{r}}^{\text{rec}},\pi_v) - I_v\big\|^2
\;+\; \lambda\,\Omega(\hat{\mathbf{r}}^{\text{rec}}),
\label{eq:p3_s8_2_3_loss}
\end{equation}
where $\Omega$ is an optional geometric or smoothness regularizer and $\mathrm{Render}$ is a differentiable renderer of $\mathcal{R}^{\text{rec}}$.
The structural contrast with §\ref{subsec:p3_s8_2_1} is sharp: there, the conditioning $\kappa(\mathbf{c})$ is a high-level prior (text, single image, multi-view) and the objective is a \emph{generative} likelihood; here, the conditioning $\mathbf{O}$ is sensor data and the objective is an \emph{observation-fidelity} loss.
% Two limiting cases anchor Eq.~\eqref{eq:p3_s8_2_3_objective} to the broader literature.
% First, when $\mathbf{O}$ is a dense calibrated multi-view set, the deterministic form of Eq.~\eqref{eq:p3_s8_2_3_objective} reduces to the classical SfM\,+\,MVS pipeline or to per-scene NeRF / 3DGS optimization: the regime of \S\ref{para:p3_s8_2_3_iterative} below.
% Second, when $\mathbf{O}$ is sparse or even a single image and large portions of the scene must be inferred, the generative form of Eq.~\eqref{eq:p3_s8_2_3_objective} merges into the synthesis regime of Eq.~\eqref{eq:p3_s8_2_1_objective}: this is precisely the regime of \S\ref{para:p3_s8_2_3_complete}, where reconstruction and synthesis become two formulations of the same underlying problem.


\subsubsection{Methodology}
\label{subsubsec:p3_s8_2_3_methodology}

The three paradigms are three parameterizations of the reconstruction operator $G_\theta^{\text{rec}}$ of Eq.~\eqref{eq:p3_s8_2_3_objective}, taken in the order in which each became dominant.
Figure~\ref{fig:p3_s8_2_3_paradigms} contrasts them, making explicit whether the 3D representation is optimized per scene, regressed in a single forward pass, or sampled from a generative posterior over the unobserved regions.

\begin{figure}[!ht]
    \centering
    \includegraphics[width=0.99\linewidth]{figures/part3/section8/3D-scene-reconstruction.pdf}
    \caption{Three paradigm families for 3D scene reconstruction (Eq.~\eqref{eq:p3_s8_2_3_objective}): (P1) iterative per-scene optimization, (P2) feed-forward reconstruction with weights shared across scenes, and (P3) generative completion of unseen (hatched) regions from a learned prior.}
    \label{fig:p3_s8_2_3_paradigms}
\end{figure}

\paragraph{Iterative 3D Reconstruction.}
\label{para:p3_s8_2_3_iterative}
The classical paradigm casts reconstruction as a \emph{per-scene} optimization: there is no learned $\theta$ shared across scenes; the 3D representation $\mathbf{r}$ itself is the optimization variable, and Eq.~\eqref{eq:p3_s8_2_3_loss} is minimized from scratch for each new input,
\begin{equation}
\hat{\mathbf{r}}
\;=\;
\arg\min_{\mathbf{r}\in\mathcal{R}^{\text{rec}}}
\;\sum_{(I_v,\pi_v)\in\mathbf{O}} \big\|\mathrm{Render}(\mathbf{r},\pi_v) - I_v\big\|^2
\;+\; \lambda\,\Omega(\mathbf{r}).
\label{eq:p3_s8_2_3_iterative_objective}
\end{equation}
Classical image-based reconstruction is rooted in Structure-from-Motion (SfM)~\cite{bolles1987epipolar, tomasi1992shape, fitzgibbon1998automatic, schonberger2016structure}, which estimates camera poses from unposed images, followed by Multi-View Stereo (MVS)~\cite{ullman1979interpretation, hartley2003multiple, Agarwal2010BAL, snavely2006photo, furukawa2015multi, schonberger2016pixelwise}, which recovers dense geometry.
Decades of refinement have made SfM\,+\,MVS a mature and widely used pipeline (see~\cite{hartley2003multiple} for a comprehensive treatment), yet it still struggles with textureless surfaces, reflective objects, and large illumination changes.

A more recent line represents the scene as a continuous function parameterized by a neural network: signed distance functions~\cite{park2019deepsdf, chibane2020neural, wang2021neus, yu2022monosdf}, occupancy fields~\cite{mescheder2019occupancy, peng2020convolutional}, and neural radiance fields~\cite{mildenhall2021nerf, barron2021mip, wang2021nerf, barron2022mip}.
The original formulations were expensive to train and render and limited in quality; PlenOctrees~\cite{yu2021plenoctrees}, Plenoxels~\cite{fridovich2022plenoxels}, Instant-NGP~\cite{muller2022instant}, RefNeRF~\cite{verbin2022ref}, TensoRF~\cite{chen2022tensorf}, and ZipNeRF~\cite{barron2023zip} are among the many follow-ups that improved efficiency and fidelity, but all of them still require a large number of input images per scene.
Sparse-view variants add auxiliary information such as regularization~\cite{niemeyer2022regnerf}, depth priors~\cite{roessle2022dense, deng2022depth}, point clouds~\cite{xu2022point}, and surface priors~\cite{oechsle2021unisurf}, or distill a 2D diffusion prior into the per-scene representation, the SDS recipe of §\ref{para:p3_s8_2_1_sds} applied to reconstruction from one or a few images~\cite{melas2023realfusion, metzer2023latent, tang2023make}.
3D Gaussian Splatting~\cite{kerbl20233d} then made per-scene optimization fast, training within minutes and rendering in real time, with subsequent work adding depth regularization, anti-aliasing, mesh extraction, and surface alignment~\cite{chung2024depth, yu2024mip, huang20242d, guedon2024sugar}.
All of these systems, however, work well only over nearby views or a small scene and must be re-optimized from scratch for every new input, which motivates the amortized paradigm.

\paragraph{Feed-forward 3D Reconstruction.}
\label{para:p3_s8_2_3_feedforward}
The feed-forward paradigm amortizes the per-scene optimization of Eq.~\eqref{eq:p3_s8_2_3_iterative_objective} across a learned prior: a network $G_\theta^{\text{rec}}$ is trained once on a large multi-view corpus, then deployed in a \emph{single forward pass} on unseen inputs, minimizing
\begin{equation}
\mathcal{L}_{\text{FF}}(\theta)
\;=\;
\mathbb{E}_{(\mathbf{O},\,\mathbf{r}_{\text{gt}})\sim p_{\text{data}}}\Big[\mathcal{D}\big(G_\theta^{\text{rec}}(\mathbf{O}),\,\mathbf{r}_{\text{gt}}\big)\Big],
\label{eq:p3_s8_2_3_feedforward_objective}
\end{equation}
where $\mathcal{D}$ is a representation-specific distance (Chamfer Distance for point maps, a rendering loss for radiance fields) and $(\mathbf{O},\mathbf{r}_{\text{gt}})$ are observation--asset pairs drawn from a 3D corpus.
This is the reconstruction-side analogue of the view-to-3D paradigm of §\ref{para:p3_s8_2_2_view3d}: sparse calibrated views are mapped feed-forward to a continuous scene representation, but the supervision is a reconstruction loss against ground-truth 3D rather than a rendering loss against held-out views alone.

Early feed-forward work applied 2D networks to individual 3D representations, for depth estimation~\cite{eigen2014depth, liu2015learning}, voxel reconstruction~\cite{choy20163d, wu20153d}, point-cloud reconstruction~\cite{fan2017point, qi2017pointnet}, and mesh deformation~\cite{wang2018pixel2mesh, gkioxari2019mesh}, but remained limited to single objects or small datasets.
Learned components then replaced individual stages of the SfM\,+\,MVS pipeline, keypoint detection and description~\cite{detone2018superpoint, dusmanu2019d2, tyszkiewicz2020disk}, image matching~\cite{sarlin2020superglue, lindenberger2023lightglue}, pose estimation~\cite{kendall2015posenet, yin2018geonet, kendall2017geometric}, and depth prediction~\cite{zhou2017unsupervised, yin2018geonet}, proving more robust and accurate than their hand-designed counterparts and being integrated into the classical pipeline~\cite{vzbontar2016stereo, yao2018mvsnet, yao2019recurrent}.
DUSt3R~\cite{wang2024dust3r} and MASt3R~\cite{mast3r} then reconstructed large scenes from unposed images end to end: similar in spirit to MVS but without explicit keypoint detection, matching, and pose estimation, the model predicts point maps in the reference view directly and recovers camera poses in post-processing, and follow-ups~\cite{tang2025mv, wang2025Cut3R, yang2025fast3r, zhang2025FLARE, wang2025vggt} have improved quality and efficiency at scale.

Beyond geometry alone, a parallel line recovers geometry and appearance together through differentiable rendering.
Scene Representation Networks~\cite{sitzmann2019scene} and Light Field Networks~\cite{sitzmann2021light} introduced ray-based querying for novel-view synthesis and were followed by epipolar attention~\cite{suhail2022light, suhail2022generalizable}, large transformers~\cite{jin2025lvsm, jiang2025rayzer}, and integration with pretrained generative models~\cite{zheng2024free3d, gao2024cat3d}; PixelNeRF~\cite{yu2021pixelnerf} pioneered pixel-aligned features for NeRF-based reconstruction and inspired 3D volume representations~\cite{chen2021mvsnerf, xu2024murf}, feature-matching fusion~\cite{chen2025explicit}, and large transformers~\cite{du2023learning, hong2023lrm}.
3D Gaussian Splatting~\cite{kerbl20233d} has since become the representation of choice for feed-forward reconstruction, in PixelSplat~\cite{charatan2024pixelsplat}, Splatter Image~\cite{szymanowicz2024splatter}, MVSplat~\cite{chen2024mvsplat}, Flash3D~\cite{szymanowicz2024flash3d}, LatentSplat~\cite{wewer2024latentsplat}, Splatt3R~\cite{smart2024splatt3r}, GS-LRM~\cite{zhang2024gs}, NoPoSplat~\cite{ye2024no}, and DepthSplat~\cite{xu2024depthsplat}.
These methods synthesize novel views with high appearance fidelity, but the views remain confined to small camera changes within the observed envelope, which motivates the generative paradigm below.

\paragraph{Complete 3D Reconstruction.}
\label{para:p3_s8_2_3_complete}
A harder regime is to reconstruct a \emph{complete} scene from an observation $\mathbf{O}$ that covers only part of it, which requires inferring both \emph{visible} and \emph{occluded} regions; the operator shifts from a deterministic map to a generative posterior over the unobserved regions,
\begin{equation}
\hat{\mathbf{r}}^{\text{rec}}\sim p_\theta^{\text{rec}}\!\big(\mathbf{r}^{\text{rec}}\,\big|\,\mathbf{O}\big),
\qquad
\mathbf{r}^{\text{rec}}\big|_{\text{visible}}\approx \text{observation-consistent},
\quad
\mathbf{r}^{\text{rec}}\big|_{\text{occluded}}\sim p_\theta(\,\cdot\,\big|\,\mathbf{O}),
\label{eq:p3_s8_2_3_complete_objective}
\end{equation}
which shares the formal structure of the synthesis sampler of Eq.~\eqref{eq:p3_s8_2_1_objective} with the conditioning replaced by a strong observation constraint $\mathbf{O}$.

Complete reconstruction has been studied extensively for single objects~\cite{vahdat2022lion, zhang20233dshape2vecset, zhao2023michelangelo, zhang2024clay, ren2024xcube, xiang2025trellis, tochilkin2024triposr, yang2024hunyuan3d, hunyuan3d22025tencent, li2025triposg} but remains under-explored for large scenes.
Classical scene-scale methods fuse depth maps from multiple views into a volumetric representation~\cite{curless1996volumetric, newcombe2011kinectfusion, whelan2015elasticfusion, dai2017bundlefusion, niessner2013real}, but they are sensitive to depth errors and produce noisy, incomplete reconstructions.
A newer line fine-tunes large pretrained diffusion models~\cite{rombach2022ldm, blattmann2023stable} for the task, conditioning on controlled camera trajectories~\cite{wang2024motionctrl, sargent2024zeronvs, wu2024reconfusion, gao2024cat3d, wallingford2024image, zhou2025stable} or on auxiliary signals~\cite{liu2024reconx, yu2024viewcrafter, chen2024mvsplat360, yu2025wonderworld}; these methods synthesize observation-consistent views but leave the underlying geometry implicit.
A small but growing body of work addresses geometry directly: WVD~\cite{zhang2024world} and Bolt3D~\cite{szymanowicz2025bolt3d} generate a hybrid RGB\,+\,point-map output that recovers geometry and appearance rather than novel views alone, and LaRI~\cite{li2025lari} reconstructs a complete scene from a single image, still limited to small indoor scenes.

Two design choices cut across the three paradigms.
The density of the observation decides which paradigm applies: a dense calibrated multi-view set calls for per-scene optimization, a sparse uncalibrated capture for a feed-forward map, and a single image for a generative posterior, since progressively more of the scene must be inferred rather than observed.
The output representation $\mathcal{R}^{\text{rec}}$, by contrast, is largely independent of the paradigm and driven by the downstream consumer: a point map (DUSt3R), a mesh (classical MVS post-processing), a NeRF (per-scene optimization, PixelNeRF), a 3DGS field (3DGS, PixelSplat), or a hybrid RGB\,+\,point map (Bolt3D, WVD).


\subsubsection{Benchmarks}
\label{subsubsec:p3_s8_2_3_benchmarks}

Datasets span object-centric collections, indoor scenes, and outdoor environments, and partition further by acquisition (synthetic or real-world) and dynamics (static or dynamic); Table~\ref{tab:p3_s8_2_3_corpora} summarizes those most commonly used.

\paragraph{Training Corpora.}
Synthetic corpora such as Objaverse~\cite{deitke2023objaverse} supply large and diverse 3D assets, and MegaSynth~\cite{jiang2025megasynth} and Zeroverse~\cite{xie2024lrmzero} explore procedurally generated data for scalable supervision.
Real-world static corpora include ARKitScenes~\cite{baruch2021arkitscenes} and ScanNet++~\cite{yeshwanth2023scannet++} for indoor capture, dynamic corpora such as KITTI-360~\cite{liao2022kitti360} capture outdoor scenes with temporal variation, and TartanAir~\cite{TartanAir} contains both static and dynamic scenes, which permits evaluation across both settings.

\begin{table}[t]
    \centering
    \fontsize{8}{10}\selectfont
    \setlength{\tabcolsep}{1.0mm}
    \caption{Representative training corpora for 3D scene reconstruction. ``Type'' indicates synthetic / real-world acquisition; ``Dynamics'' indicates whether the corpus contains temporal variation.}
    \label{tab:p3_s8_2_3_corpora}
    \begin{tabular}{L{0.17\textwidth} C{0.06\textwidth} C{0.14\textwidth} L{0.13\textwidth} C{0.09\textwidth} L{0.33\textwidth}}
        \toprule
        \textbf{Dataset} & \textbf{Year} & \textbf{Scale} & \textbf{Type} & \textbf{Dynamics} & \textbf{Primary use} \\
        \midrule
        Objaverse~\cite{deitke2023objaverse} & 2023 & ${\sim}800\mathrm{K}$ objects & Synthetic (CAD) & Static & Object-centric multi-view supervision \\
        Zeroverse~\cite{xie2024lrmzero} & 2024 & Procedural & Synthetic & Static & LRM-zero scalable supervision \\
        MegaSynth~\cite{jiang2025megasynth} & 2025 & Procedural & Synthetic & Static & Scalable multi-view scene supervision \\
        ARKitScenes~\cite{baruch2021arkitscenes} & 2021 & ${\sim}5\mathrm{K}$ scans & Real-world & Static & Indoor RGB-D / mesh supervision \\
        ScanNet++~\cite{yeshwanth2023scannet++} & 2023 & ${\sim}460$ scenes & Real-world & Static & High-fidelity indoor multi-modal scans \\
        KITTI-360~\cite{liao2022kitti360} & 2022 & ${\sim}300\mathrm{K}$ frames & Real-world & Dynamic & Outdoor driving scenes \\
        TartanAir~\cite{TartanAir} & 2020 & ${\sim}1\mathrm{M}$ frames & Synthetic & Mixed & Static + dynamic mixed-domain evaluation \\
        \bottomrule
    \end{tabular}
\end{table}

\paragraph{Evaluation Benchmark Suites.}
Iterative methods (P1) are evaluated on DTU~\cite{Aanaes2016DTU} and Tanks-and-Temples~\cite{knapitsch2017tanks} for geometric fidelity and on the NeRF-era benchmarks LLFF~\cite{mildenhall2019llff}, MipNeRF360~\cite{barron2022mip}, and RealEstate10K~\cite{zhou2018stereo} for novel-view quality.
Feed-forward methods (P2) adopt the same panel together with ScanNet~\cite{dai2017scannet} subsets that supply ground-truth point maps for direct geometry supervision.
Complete-reconstruction methods (P3) have no standardized suite and typically report on the precursor benchmark of their underlying 2D or video prior.

\paragraph{Constraint-specific Metrics.}
Four axes mirror the rendering-versus-geometry decomposition of Eq.~\eqref{eq:p3_s8_2_3_loss}.
\textbf{Photometric quality} is measured by PSNR, SSIM~\cite{wang2004image}, and LPIPS~\cite{zhang2018unreasonable}, covering pixel-level fidelity, structural similarity, and perceptual consistency of novel views.
\textbf{Camera-pose adherence} uses Relative Translation Accuracy (RTA), Relative Rotation Accuracy (RRA), and the Area Under the Curve (AUC) that summarizes pose accuracy across angular thresholds.
\textbf{Geometric fidelity} of point maps and point clouds is measured by Accuracy (precision), Completeness (recall), and Chamfer Distance, where Accuracy reflects geometric precision, Completeness measures surface coverage, and Chamfer Distance aggregates both.
\textbf{Depth quality} for monocular estimation uses the absolute relative error $|y - \hat{y}| / y$ and the inlier ratio $\delta\!<\!T$ (the fraction of predictions within a multiplicative threshold of the ground truth), with median scaling or least-squares fitting applied beforehand because monocular depth is scale-ambiguous; multi-view depth reports Accuracy, Completeness, and Overall, the mean of the first two and equivalent to Chamfer Distance.

The paradigm succession of §\ref{subsubsec:p3_s8_2_3_methodology} is visible in two quantities.
Per-scene optimization has moved from NeRF~\cite{mildenhall2021nerf}, which trains for hours and renders at a few frames per second, to 3D Gaussian Splatting~\cite{kerbl20233d}, which trains a scene in minutes and renders in real time (${\geq}30$~fps at 1080p); the feed-forward paradigm removes per-scene optimization altogether, with DUSt3R / MASt3R~\cite{wang2024dust3r, mast3r} and VGGT~\cite{wang2025vggt} recovering point maps, depth, and camera poses in a single sub-second forward pass and Fast3R~\cite{yang2025fast3r} processing over a thousand images at once.
On pose, learned reconstructors now rival classical pipelines: VGGT~\cite{wang2025vggt} reports a camera-pose AUC@30 of 85.3 on RealEstate10K in pure feed-forward mode and 93.5 with bundle-adjustment refinement.
Two caveats apply. There is no cross-paradigm leaderboard, since P1 is scored on novel-view PSNR / SSIM / LPIPS, P2 adds geometric Accuracy / Completeness / Chamfer Distance and pose AUC, and P3 inherits whichever protocol its 2D or video prior used; and on raw depth-map accuracy feed-forward models still trail tuned MVS pipelines, with quality degrading under free-viewpoint extrapolation (§\ref{subsubsec:p3_s8_2_3_challenges}), so the strongest numbers hold inside the observed view envelope.
Phone-based NeRF and Gaussian-Splatting capture apps and the open-source Nerfstudio stack have brought per-scene reconstruction to commodity devices, while classical photogrammetry suites remain the production baseline for metric accuracy.


\subsubsection{Challenges and Trends}
\label{subsubsec:p3_s8_2_3_challenges}

Four problems span the pipeline end to end: the data the model trains on, the accuracy it can deliver, the quality of the views it renders, and the compute envelope that constrains all three at long input lengths; reconstruction shares the first and third with its synthesis counterparts but sharpens the other two because it commits to recovering ground truth from raw sensor data.

\paragraph{Limited Modality in Datasets.}
Most reconstruction and view-synthesis datasets cover a narrow band of modalities.
Widely used benchmarks such as RealEstate10K~\cite{zhou2018stereo} and MipNeRF360~\cite{barron2022mip} are RGB-only, without depth, LiDAR, or semantic annotation, and even Objaverse-XL~\cite{deitke2023objaversexl}, with 10.2M objects, is primarily synthetic mesh data without real-world sensor modalities.
Merging datasets of different modalities is the usual workaround, at the cost of domain shift and annotation inconsistency, and the gap is widest for dynamic scenes: datasets with synchronized RGB, depth, optical flow, and 3D motion remain far fewer than static ones, and most dynamic datasets record either camera motion or object motion rather than both with a full sensor suite.
The open problem is a scalable, modality-rich corpus that combines the diversity of Objaverse~\cite{deitke2023objaverse} with the multi-sensor completeness of ScanNet++~\cite{yeshwanth2023scannet++}, collected under one protocol so that aligned RGB, depth, and semantic signals supervise the same model.

\paragraph{Reconstruction Accuracy.}
Feed-forward models still trail traditional MVS pipelines~\cite{zhang2023geomvsnet, zhang2025advances} on raw depth-map accuracy, because MVS exploits known camera parameters and a hypothesized depth set to build cost volumes from which depth is predicted.
Whether feed-forward reconstructors~\cite{yang2025fast3r, wang2025vggt} learn an approximation of such cost volumes implicitly in their attention layers is an open question; injecting explicit camera parameters or other priors into the backbone, for instance through Diffusion Transformers~\cite{peebles2023dit}, is one route to closing the gap.

\paragraph{Free-viewpoint Rendering.}
The accuracy gap is most visible at \emph{free-viewpoint} rendering, where the model must produce novel views far from the input camera distribution.
Disocclusions, geometric uncertainty, and limited generalization conspire: regions never observed come out blurred, ghosted, or geometrically wrong, because current methods rely on local consistency and cannot infer plausible content for occluded areas, and view-dependent effects and complex light transport demand reasoning beyond interpolation.
Hallucinating missing detail while keeping novel views consistent is exactly the generative regime of §\ref{para:p3_s8_2_3_complete}.

\paragraph{Long Context Input.}
Reconstruction from many views meets a hard computational ceiling: methods for geometry reasoning and novel-view synthesis rely on full attention, whose compute grows quadratically with the number of tokens; exact-attention kernels such as FlashAttention-2~\cite{dao2023flashattention2} remove the quadratic memory of the attention matrix but not the quadratic compute.
VGGT~\cite{wang2025vggt}, for example, needs about 21\,GB of GPU memory for 100 input frames and about 41\,GB for 200, so inference over long sequences and training on many views quickly exceed a single GPU.
Recurrent alternatives such as CUT3R~\cite{wang2025Cut3R} integrate new views incrementally while maintaining a memory state, which keeps inference memory near constant, but they forget earlier views and degrade as the number of inputs grows.
Reasoning efficiently over hundreds or thousands of views while keeping memory and compute manageable remains open, with recurrent and sparse-attention architectures the main candidates.



\subsection{3D Panorama Synthesis}
\label{subsec:p3_s8_2_4}

3D panorama synthesis occupies a distinctive position in spatial generation: rather than producing a 3D representation directly, it generates a complete 360-degree omnidirectional image as a structured \emph{intermediate} from which a full 3D environment can then be lifted.
The panorama gives embodied agents, AR / VR pipelines, and world models a holistic spatial context with minimal occlusion anchored at one observation point, and it bridges the mature 2D generative prior of §\ref{sec:p3_s8_1} with the 3D substrate on which the rest of §\ref{sec:p3_s8_2} operates.
Four method families trace the path from a single omnidirectional snapshot to an explorable world: panoramic image synthesis, panoramic video generation, panoramic 3D reconstruction, and panorama-based 3D scene generation.


\subsubsection{Task Definition}
\label{subsubsec:p3_s8_2_4_task}

Panorama synthesis generates a complete 360-degree environment conditioned on text, a reference panorama or NFoV image, a spatial layout, or a specified viewpoint; unlike object-centric or narrow-field-of-view (NFoV) generation, it covers the whole environment around the observer at once rather than a fragment of it.
Formally, we cast it as the omnidirectional-intermediate counterpart of the generative frameworks of §\ref{subsec:p3_s8_2_1} and §\ref{subsec:p3_s8_2_2}.
Let $\mathcal{P}$ denote the space of panoramic images (functions defined on the sphere $\mathbb{S}^2$ and encoded through a projection operator $\Pi\in\{\text{ERP},\text{cubemap},\text{tangent},\text{fisheye}\}$ onto a processable 2D array) and let $\mathbf{c}\in\mathcal{C}\supseteq\{\text{text},\text{reference panorama / NFoV image},\text{layout},\text{viewpoint / camera trajectory}\}$ be the conditioning.
Panorama synthesis first learns a panoramic generator and then optionally lifts its output into 3D,
\begin{equation}
G_\theta^{\text{pano}}:\mathcal{C}'\rightarrow\mathcal{P},
\qquad
\hat{P}\sim p_\theta^{\text{pano}}\!\big(P\,\big|\,\kappa(\mathbf{c})\big),
\qquad
\hat{\mathbf{r}}^{\text{rec}} = \Lambda\big(\hat{P}\big),
\label{eq:p3_s8_2_4_objective}
\end{equation}
where $\kappa$ is the condition encoder of §\ref{subsec:p3_s8_2_1}, $\hat{P}\in\mathcal{P}$ is the synthesized omnidirectional image, and $\Lambda:\mathcal{P}\rightarrow\mathcal{R}^{\text{rec}}$ is an optional \emph{lifting operator} mapping the panorama into the 3D-representation space $\mathcal{R}^{\text{rec}}$ of §\ref{subsec:p3_s8_2_3} (drawn from $\mathcal{R}$ of §\ref{subsec:p3_s8_2_1}, \eg NeRF, 3DGS, mesh, point map).
Equation~\eqref{eq:p3_s8_2_4_objective} makes the role of the panorama as a structured \emph{intermediate} explicit: the hard generative work is performed on the sphere by $G_\theta^{\text{pano}}$ (inheriting the 2D generator $G_\theta$ of Eq.~\eqref{eq:t2i_generator} with its outputs defined on $\mathbb{S}^2$ rather than the plane), while $\Lambda$ reuses the reconstruction machinery of §\ref{subsec:p3_s8_2_3}.
This is also what separates panorama synthesis from the scene synthesis of §\ref{subsec:p3_s8_2_2}, which produces 3D representations directly from arbitrary viewpoints distributed over $\mathrm{SE}(3)$: here every direction around one observation point is represented coherently first, and 3D structure is recovered from that complete view.
Two limiting cases locate Eq.~\eqref{eq:p3_s8_2_4_objective} within §\ref{sec:p3_s8_2}.
When $\Lambda=\mathrm{id}$ (no lifting), the task collapses to omnidirectional 2D image generation, a spatially aware variant of the text-to-image sampler of §\ref{sec:p3_s8_1} defined on the sphere rather than the plane.
When $\mathbf{c}$ is a set of sparse posed panoramas and $\Lambda$ is a per-scene optimization, Eq.~\eqref{eq:p3_s8_2_4_objective} reduces to panoramic reconstruction, recovering the observation-conditioned operator of Eq.~\eqref{eq:p3_s8_2_3_objective}.

\paragraph{Panoramic Image Formats.}
The projection $\Pi$ fixes how the sphere is laid out for a 2D network and which distortions the generator must absorb.
Equirectangular projection (ERP) maps the sphere onto a rectangle in latitude--longitude coordinates; it stretches content near the poles but is continuous and directly compatible with standard 2D pipelines and convolutional networks, which has made it the dominant format for learning-based synthesis.
Cubemap projection represents the sphere as six perspective faces (front, back, left, right, top, bottom), removing polar distortion at the cost of discontinuities along face boundaries, and is the native format of real-time rendering and VR / AR systems, where each face is processed as an ordinary perspective image.
Tangent images and perspective patches decompose the panorama into overlapping tangent-plane projections so that models trained on ordinary perspective images can be reused.
Fisheye and dual-fisheye projections capture 180--220 degrees per lens with radial distortion; most consumer 360-degree cameras capture two hemispheres this way and stitch them, so dual-fisheye is the native acquisition format even when methods operate on ERP.


\subsubsection{Methodology}
\label{subsubsec:p3_s8_2_4_methodology}

The four families are four instantiations of the panoramic generator $G_\theta^{\text{pano}}$ and lifting operator $\Lambda$ of Eq.~\eqref{eq:p3_s8_2_4_objective}: image synthesis stops at a static panorama, video generation extends it over time, panoramic reconstruction applies $\Lambda$ to observed panoramas, and panorama-based scene generation composes the two.
Figure~\ref{fig:p3_s8_2_4_paradigms} schematizes them and the panorama-as-intermediate pipeline that connects them.

\begin{figure}[!ht]
    \centering
    \includegraphics[width=0.99\linewidth]{figures/part3/section8/Four_method_families_for_3D_panorama_synthesis.pdf}
    \caption{The four method families of 3D panorama synthesis (Eq.~\eqref{eq:p3_s8_2_4_objective}), linked by a panoramic intermediate (top): (a) panoramic image synthesis, (b) panoramic video generation, (c) panoramic 3D reconstruction, and (d) panorama-based 3D scene generation. Room imagery is illustrative.}
    \label{fig:p3_s8_2_4_paradigms}
\end{figure}

\paragraph{Panoramic Image Synthesis.}
Two task formulations differ in their conditioning: \emph{text-to-panorama generation}, which leans on the prior of a pretrained generator and is judged on semantic alignment and diversity, and \emph{image-conditioned panorama generation}, which expands an NFoV image or an incomplete panorama by inpainting or outpainting and must preserve the content, style, and panoramic distortion of its input.
Formally, this family is the base case of Eq.~\eqref{eq:p3_s8_2_4_objective} with the lifting operator disabled: it samples a single static omnidirectional image and stops,
\begin{equation}
\hat{P} \sim p_\theta^{\text{pano}}\!\big(P \mid \kappa(\mathbf{c})\big),
\qquad
\Lambda = \mathrm{id},
\label{eq:p3_s8_2_4_image}
\end{equation}
where the generative work is carried out under a chosen projection $\Pi\in\{\text{ERP},\text{cubemap},\text{tangent}\}$ that dictates how spherical distortion is handled; Figure~\ref{fig:p3_s8_2_4_paradigms}(a) depicts this static-panorama pathway.

Early approaches~\cite{chen2022text2light, may2023cubegan} relied on GANs~\cite{goodfellow2014gan} and were constrained by limited training data and mode collapse; diffusion models~\cite{ho2020denoising, rombach2022ldm} brought diversity and fidelity, and current methods divide by projection format and architecture.
The three strategies below are three realizations of the sampler of Eq.~\eqref{eq:p3_s8_2_4_image}, differing in how the spherical output is factorized before being mapped onto the sphere,
\begin{equation}
\hat{P} =
\begin{cases}
\Pi_{\text{ERP}}^{-1}(\hat{Z}),\ \ \hat{Z}\sim p_\theta\big(Z\mid\kappa(\mathbf{c})\big) & \text{(direct ERP)},\\[3pt]
\Pi_{\text{cube}}^{-1}\big(\{\hat{F}_i\}_{i=1}^{6}\big),\ \ \{\hat{F}_i\}\sim p_\theta\big(\{F_i\}\mid\kappa(\mathbf{c})\big) & \text{(cubemap)},\\[3pt]
\mathcal{S}\big(\{\hat{I}_k\}_{k=1}^{K}\big),\ \ \hat{I}_k\sim p_\theta^{\text{2D}}\big(I\mid\kappa(\mathbf{c})\big) & \text{(perspective + stitch)},
\end{cases}
\label{eq:p3_s8_2_4_image_proj}
\end{equation}
where $\Pi^{-1}$ lifts the generated representation onto the sphere and $\mathcal{S}$ is a spherical stitch-and-blend operator. The first two cases synthesize the panorama \emph{directly} (and so require panoramic training data), whereas the third reuses a frozen perspective generator $p_\theta^{\text{2D}}$ and pushes the cross-view-consistency burden into $\mathcal{S}$.

$\blacktriangleright$ \textit{Equirectangular projection (ERP) methods.} The direct-ERP case generates on the latitude--longitude grid and must handle polar distortion and the wrap-around seam architecturally.
For text-to-panorama generation, DiffPano~\cite{ye2024diffpano} fine-tunes a diffusion model with LoRA on ERP, SMGD~\cite{sun2025spherical} guides the diffusion process with spherical manifold convolutions that respect geometry across the sphere, and PAR~\cite{wang2025conditional} replaces diffusion with autoregressive modeling to avoid the mismatch between ERP and the noise assumptions of standard diffusion.
The seam is treated at the denoising or decoding stage: PanoDiff~\cite{panodiff} uses geometric augmentation when completing a panorama from unregistered NFoV images, PanoDiffusion~\cite{wupanodiffusion} generates depth jointly and rotates the latent circularly to enforce left--right continuity, StitchDiffusion~\cite{stitchdiffusion} re-denoises the wrap-around region at every step together with global cropping, and PanFusion~\cite{zhang2024taming} pairs an ERP branch with a multi-view perspective branch to draw on both representations.
For image-conditioned generation, OmniDreamer~\cite{omnidreamer2022} introduced mask-based outpainting from NFoV images with GANs, AOG~\cite{aog} and ODIS~\cite{odis} explore alternative conditioning for inpainting and outpainting, and PanoDecouple~\cite{zheng2025panorama} separates distortion guidance from content completion to counter the ``visual cheating'' in which a model buys image quality at the cost of geometric accuracy.

$\blacktriangleright$ \textit{Cubemap-based methods.} The second case of Eq.~\eqref{eq:p3_s8_2_4_image_proj} generates six perspective faces, which removes polar distortion but moves the problem to the face boundaries.
CubeDiff~\cite{kalischek2025cubediff} introduced the cubemap representation for text-to-panorama generation, and DreamCube~\cite{huang2025dreamcube} synchronizes the six planes to suppress discontinuous seams and inconsistent color tone; overlapping the fields of view of adjacent faces reduces seams further but duplicates generation in the overlap and lowers the effective resolution.
JoPano~\cite{feng2025jopano} models the faces jointly with a DiT and a Joint-Face Adapter, applies Poisson blending along the boundaries, and introduces the seam metrics (Seam-SSIM, Seam-Sobel) used in §\ref{subsubsec:p3_s8_2_4_benchmarks}.

$\blacktriangleright$ \textit{Perspective image-based methods.} The third case leaves a perspective text-to-image model untouched and pushes the consistency burden into the stitch operator $\mathcal{S}$, with or without panorama-specific training.
\emph{Training-free methods} apply a pretrained perspective model directly: MultiDiffusion~\cite{multidiffusion} denoises overlapping patches independently and blends them, PanoGen~\cite{li2023panogen} outpaints recursively from an initial view until the sphere is covered, DiffCollage~\cite{zhang2023diffcollage} casts the panorama as a factor graph whose factor nodes are portions and whose variable nodes are their overlaps so that portions are generated in parallel rather than autoregressively, PanoFree~\cite{panofree} iterates warping and inpainting over multi-view images, and SphereDiff~\cite{park2025spherediff} moves MultiDiffusion onto a spherical latent with dynamic spherical sampling and distortion-aware weighted averaging so that quality holds up to the poles.
These methods ride on advances in perspective generation but struggle with global consistency across the full sphere.
\emph{Training-based methods} fine-tune on panoramic data: MVDiffusion~\cite{tang2023MVDiffusion} generates eight perspective views with consistent geometry and reprojects them into a closed-loop panorama, and TanDiT~\cite{tandit2025} decomposes the panorama into a structured grid of tangent-plane images that a Diffusion Transformer generates jointly in one denoising pass under spatial attention.


\paragraph{Panoramic Video Generation.}
Extending synthesis over time adds temporal coherence to omnidirectional consistency.
Formally, this family extends the static generator of Eq.~\eqref{eq:p3_s8_2_4_image} over time, sampling a sequence under a joint spatial-and-temporal coherence constraint,
\begin{equation}
\{\hat{P}_t\}_{t=1}^{T} \sim p_\theta^{\text{pano}}\!\big(\{P_t\}_{t=1}^{T} \mid \kappa(\mathbf{c})\big),
\qquad
\text{s.t.}\;\; \hat{P}_t \approx \mathcal{W}^{\text{t}}(\hat{P}_{t-1}),
\label{eq:p3_s8_2_4_video}
\end{equation}
where $\mathcal{W}^{\text{t}}$ enforces frame-to-frame motion consistency on top of the per-frame 360-degree wrap-around continuity; Figure~\ref{fig:p3_s8_2_4_paradigms}(b) depicts this temporally-extended pathway.
The methods share this generator but differ in the conditioning argument $\mathbf{c}$ supplied to Eq.~\eqref{eq:p3_s8_2_4_video}, which separates them into three input modalities,
\begin{equation}
\mathbf{c}=
\begin{cases}
\mathbf{x} & \text{(text-to-video)},\\
P_0 & \text{(image-to-video)},\\
\{(I_t,\mathcal{M}_t)\}_{t=1}^{T} & \text{(perspective-to-video)},
\end{cases}
\label{eq:p3_s8_2_4_video_cond}
\end{equation}
where $\mathbf{x}$ is a text prompt, $P_0$ a static panorama anchor reused as the first frame, and $(I_t,\mathcal{M}_t)$ an NFoV frame paired with its visibility mask. The difficulty grows down the list: as the conditioning becomes spatially sparser, the model must hallucinate progressively more of the sphere while preserving the temporal coherence of Eq.~\eqref{eq:p3_s8_2_4_video}.

$\blacktriangleright$ \textit{Text-to-panoramic video generation.} 360DVD~\cite{wang2024360dvd} adapts a pretrained text-to-video model to the ERP domain through adapter modules, keeping wrap-around consistency across frames but struggling to balance text alignment against spherical geometry and to vary style under several control conditions.
TiP4GEN~\cite{xing2025tip4gen} adopts a dual-branch panoramic-and-perspective architecture in the spirit of PanFusion and lifts the generated video into a geometry-aligned 4D Gaussian scene.
PanoWan~\cite{xia2025panowan} and PanoDiT~\cite{zhang2025panodit} fine-tune large video diffusion transformers on panoramic video data and gain motion quality from scale, and PanoLora~\cite{dong2025panolora} treats the task as a projection adaptation from perspective views that LoRA can learn with little data.
Training-free alternatives carry the static-panorama machinery into time: SphereDiff~\cite{park2025spherediff} extends its spherical latent to video, and DynamicScaler~\cite{liu2025dynamicscaler} denoises through a seamless rotating window with an offset-shifting scheme and Global Motion Guidance, which yields arbitrary output resolution at constant memory.
Text alone, however, offers little control over layout and motion, which motivates the richer conditioning below.

$\blacktriangleright$ \textit{Image-to-panoramic video generation.} Animating a static panorama anchors the video spatially and asks for plausible omnidirectional motion that respects the input geometry.
4K4DGen~\cite{li20244k4dgen} animates user-selected regions with an image-to-video model under overlapping windows, at the cost of fixed-window artifacts and limited motion.
ViewPoint~\cite{fang2025panoramic} introduces a panorama representation that keeps global spatial continuity while Pano-Perspective attention transfers a perspective video prior, and HoloTime~\cite{zhou2025holotime} reconstructs the result as 4D Gaussian Splatting so that the animated scene can be explored under viewpoint changes.
GenEx~\cite{lu2024genex} generates a 3D-consistent imaginative environment from a single image as a panoramic video stream with loop consistency over long trajectories, which makes it usable for embodied exploration.
For explicit control, PanFlow~\cite{panflow} conditions on synthetic optical flow as a scene-agnostic motion signal and warps noise spherically so that camera and object motion can be controlled separately while the seam stays consistent; OmniDrag~\cite{li2026omnidrag} adds drag-style interaction at scene and object level through a spherical motion estimator; and PanoWorld-X~\cite{yin2025panoworld} reprojects ERP features onto the sphere inside a Sphere-Aware Diffusion Transformer to follow diverse camera trajectories.

$\blacktriangleright$ \textit{Perspective-to-panoramic video generation.} Expanding an NFoV video to a full 360-degree video requires coherent hallucination in space and time at once, under visibility masks that change with camera motion.
VidPanos~\cite{ma2024vidpanos} frames the task as space-time outpainting and uses a generative video model to fill the regions a panning camera never observed; Imagine360~\cite{tan2024imagine360} combines a global panoramic branch with local perspective priors and an elevation-aware design for the moving masks; and Beyond the Frame~\cite{beyondtheframe} adds geometry-aware modules and motion-aware conditioning to improve realism and motion consistency in the outpainted regions.


\paragraph{Panoramic 3D Reconstruction and Generation.}
Reconstruction bypasses the generator of Eq.~\eqref{eq:p3_s8_2_4_objective}: rather than generating $\hat{P}$, it takes \emph{observed} posed panoramas and applies the lifting operator directly,
\begin{equation}
\hat{\mathbf{r}}^{\text{rec}} = \Lambda\big(\{P_v, \pi_v\}_{v=1}^{V}\big),
\label{eq:p3_s8_2_4_recon}
\end{equation}
where $\Lambda$ is realized either as a per-scene optimization (iterative) or as an amortized network (feed-forward), mirroring the split of §\ref{subsec:p3_s8_2_3}; Figure~\ref{fig:p3_s8_2_4_paradigms}(c) depicts this reconstruction pathway.
Generation then composes the two directions, feeding a synthesized panorama into the same $\Lambda$.

$\blacktriangleright$ \textit{Iterative optimization-based panoramic reconstruction.} Realizing $\Lambda$ of Eq.~\eqref{eq:p3_s8_2_4_recon} as a \emph{per-scene} optimization,
\begin{equation}
\hat{\mathbf{r}}^{\text{rec}} = \arg\min_{\mathbf{r}} \sum_{v} \big\| \mathrm{Render}(\mathbf{r},\pi_v) - P_v \big\|^2 + \lambda\,\Omega(\mathbf{r}),
\label{eq:p3_s8_2_4_recon_iter}
\end{equation}
re-solved from scratch for each scene, these methods fit a radiance field or an explicit representation to multi-view panoramas, the panoramic analogue of §\ref{para:p3_s8_2_3_iterative}.
NeRF-based work includes 360Roam~\cite{huang2022360roam} for geometry-aware indoor reconstruction, 360FusionNeRF~\cite{kulkarni2023fusion} for novel views from a single panorama under geometric and semantic supervision, EgoNeRF~\cite{choi2023egonerf} with spherical grids for egocentric capture, and OmniNeRF~\cite{gu2022omninerf}, which processes raw dual-fisheye input and its lens distortion.
Gaussian-splatting alternatives are faster: 360GS~\cite{bai2024360gs} projects Gaussians onto spherical tangent planes, ODGS~\cite{lee2024odgs} rasterizes on tangent planes with elevation-aware densification, OP43DGS~\cite{wang2024op43dgs} corrects the projection bias of the local affine approximation across camera models, OmniGS~\cite{li2024omnigs} exploits omnidirectional optical properties, and Seam360GS~\cite{seam360gs} calibrates dual-fisheye input by jointly optimizing translation vectors and learnable angular-distortion embeddings.
Scene4U~\cite{scene4u} reconstructs from a single panorama in depth layers: an open-vocabulary segmenter and an LLM decompose it semantically, diffusion inpainting restores occluded regions, and 3D Gaussians are fitted per layer.

$\blacktriangleright$ \textit{Feed-forward generalization-based panoramic reconstruction.} Amortizing $\Lambda$ of Eq.~\eqref{eq:p3_s8_2_4_recon} into a trained network applied in a single forward pass,
\begin{equation}
\hat{\mathbf{r}}^{\text{rec}} = \Lambda_\psi\big(\{P_v, \pi_v\}_{v=1}^{V}\big),
\label{eq:p3_s8_2_4_recon_ff}
\end{equation}
where $\Lambda_\psi$ is trained once over a panoramic corpus (the analogue of §\ref{para:p3_s8_2_3_feedforward}), trades the per-scene cost of Eq.~\eqref{eq:p3_s8_2_4_recon_iter} for a learned prior that reconstructs from sparse panoramas without scene-specific training.
PanoGRF~\cite{chen2023panogrf} pioneered feed-forward reconstruction of wide-baseline panoramas with depth guidance; PanoSplatt3R~\cite{panosplatt3r} extends DUSt3R~\cite{wang2024dust3r} to unposed wide-baseline panoramas through RoPE rolling, a rotary positional embedding rolled across attention heads so that ERP periodicity and correspondences across the seam are modeled; and PanSplat~\cite{pansplat} builds spherical 3D Gaussian pyramids on a Fibonacci lattice to reach 4K novel views with feed-forward inference and cubemap rendering for memory efficiency.

$\blacktriangleright$ \textit{Panorama-based 3D generation and world modeling.} This family realizes the full composition of Eq.~\eqref{eq:p3_s8_2_4_objective}: it first \emph{generates} the panorama and then lifts it,
\begin{equation}
\hat{\mathbf{r}}^{\text{rec}} = \Lambda\big(\hat{P}\big),
\qquad
\hat{P} \sim p_\theta^{\text{pano}}\!\big(P \mid \kappa(\mathbf{c})\big),
\label{eq:p3_s8_2_4_gen}
\end{equation}
so that the synthesized omnidirectional image of Eq.~\eqref{eq:p3_s8_2_4_image} becomes the structured intermediate fed to the reconstruction operator of Eq.~\eqref{eq:p3_s8_2_4_recon}; Figure~\ref{fig:p3_s8_2_4_paradigms}(d) depicts this generate-then-lift pathway.
Text $\rightarrow$ panorama $\rightarrow$ 3D pipelines generate the omnidirectional image with a diffusion model~\cite{ye2024diffpano,kalischek2025cubediff} before reconstructing a NeRF, Gaussian, or mesh scene.
LayerPano3D~\cite{yang2024layerpano3d} parses the panorama into hierarchical depth layers and reconstructs each with 3D Gaussian Splatting, DreamScene360~\cite{zhou2025dreamscene360} fits panoramic Gaussians for free-roaming text-conditioned scenes, Matrix-3D~\cite{yang2025matrix} couples trajectory-guided panoramic video generation with a dual-pathway reconstruction that conditions on scene-mesh renders, and HunyuanWorld~\cite{team2025hunyuanworld} treats the panorama as a 360-degree world proxy for text- and image-conditioned explorable scenes with mesh export and disentangled object representations.
The panoramic video generators GenEx~\cite{lu2024genex} and PanoWorld-X~\cite{yin2025panoworld} (above) turn the same intermediate into an action- or trajectory-conditioned observation space, which is where panorama synthesis meets the world models of §\ref{sec:p3_s8_6}.


\subsubsection{Benchmarks}
\label{subsubsec:p3_s8_2_4_benchmarks}

Panoramic datasets must supply geometric consistency over the full sphere, a declared projection format, and, for reconstruction, camera poses and depth; Table~\ref{tab:p3_s8_2_4_corpora} lists the corpora in use, grouped by modality.

\paragraph{Training Corpora.}
For static synthesis, Structured3D~\cite{zheng2020structured3d} (synthetic indoor scenes with controlled geometry) and SUN360~\cite{xiao2012recognizing} (real-world indoor and outdoor panoramas) are the standard benchmarks for distortion handling and seam artifacts, in the splits adopted by JoPano, PanoDecouple, and DreamCube~\cite{feng2025jopano,zheng2025panorama,huang2025dreamcube}; Matterport3D~\cite{chang2017matterport3d}, with RGB-D panoramas, camera poses, and mesh geometry across 90 buildings, is the most widely used corpus for both generation and reconstruction~\cite{pansplat,yang2024layerpano3d,team2025hunyuanworld,tang2023MVDiffusion,zhang2024taming}, and Upright360~\cite{yang2024layerpano3d} filters Matterport3D, web, and Blockade Labs panoramas with GeoCalib~\cite{veicht2024geocalib} to near-zero pitch and roll for realistic indoor synthesis.
Panoramic video with geometric annotation has been the bottleneck.
360+x~\cite{chen2023360plusx} and WEB360~\cite{wang2024360dvd} are small real-world collections (the latter captioned for text conditioning), 360-1M~\cite{wallingford2024image} scales to over a million videos for world-model pretraining, PanFlow~\cite{panflow} adds optical flow and camera poses on motion-rich sequences, PANOVID~\cite{xia2025panowan} aggregates these with Imagine360~\cite{tan2024imagine360} and public VR videos but lacks poses and depth, and Matrix-Pano~\cite{yang2025matrix} fills that gap with Unreal Engine 5 renders carrying depth maps and camera trajectories; embodied work renders action-conditioned observations in Unreal Engine directly~\cite{lu2024genex,yin2025panoworld}.
For reconstruction, Replica~\cite{straub2019replica} and Residential~\cite{habtegebrial2022somsi} serve as synthetic indoor benchmarks at baselines of about 1.0\,m and 0.3\,m respectively~\cite{pansplat}, 360Loc~\cite{huang2024360loc} supplies posed panorama sequences used for 4K fine-tuning~\cite{pansplat}, and the scarcity of posed real panoramic video led PanSplat~\cite{pansplat} to capture its own dual-fisheye footage with poses from visual SLAM.

\begin{table}[t]
    \centering
    \fontsize{8}{10}\selectfont
    \setlength{\tabcolsep}{1.0mm}
    \caption{Representative training corpora for 3D panorama synthesis, grouped by primary modality. ``Modality'' indicates the data type and synthetic / real acquisition. Structured3D and SUN360 scales are those of the splits used by recent panorama generators (SUN360 in full has $67{,}583$ panoramas).}
    \label{tab:p3_s8_2_4_corpora}
    \begin{tabular}{L{0.17\textwidth} C{0.06\textwidth} C{0.15\textwidth} L{0.22\textwidth} L{0.32\textwidth}}
        \toprule
        \textbf{Dataset} & \textbf{Year} & \textbf{Scale} & \textbf{Modality} & \textbf{Primary use} \\
        \midrule
        Structured3D~\cite{zheng2020structured3d} & 2020 & ${\sim}21\mathrm{K}$ panoramas & Synthetic image & Indoor synthesis; distortion / seam benchmark \\
        SUN360~\cite{xiao2012recognizing} & 2012 & ${\sim}29\mathrm{K}$ panoramas & Real image & Real-scene generation benchmark \\
        Matterport3D~\cite{chang2017matterport3d} & 2017 & $10{,}800$ RGB-D panoramas & Real image + depth & Generation and reconstruction supervision \\
        Upright360~\cite{yang2024layerpano3d} & 2024 & ${\sim}9.4\mathrm{K}$ panoramas & Mixed image & Geometrically upright indoor synthesis \\
        \midrule
        360+x~\cite{chen2023360plusx} & 2024 & $2{,}152$ videos & Real multimodal video & Panoptic scene understanding \\
        WEB360~\cite{wang2024360dvd} & 2024 & ${\sim}2\mathrm{K}$ videos & Real captioned video & Text-to-panoramic video \\
        360-1M~\cite{wallingford2024image} & 2024 & $>1\mathrm{M}$ videos & Real video & Large-scale world-model pretraining \\
        PANOVID~\cite{xia2025panowan} & 2025 & Aggregated & Real video & Text-to-panoramic video training \\
        Matrix-Pano~\cite{yang2025matrix} & 2025 & $116{,}759$ sequences & Synthetic video + depth + trajectory & Trajectory-conditioned generation \\
        \midrule
        Replica~\cite{straub2019replica} & 2019 & $18$ scenes & Synthetic (rendered scans) & Indoor reconstruction benchmark \\
        Residential~\cite{habtegebrial2022somsi} & 2022 & Indoor scenes & Synthetic image & Indoor reconstruction benchmark \\
        360Loc~\cite{huang2024360loc} & 2024 & $18$ posed sequences & Real reconstruction & 4K reconstruction fine-tuning \\
        \bottomrule
    \end{tabular}
\end{table}

\paragraph{Evaluation Metrics.}
Panorama evaluation combines general image and video quality metrics with panorama-specific measures of distortion, seam consistency, and spherical topology.

$\blacktriangleright$ \textit{Image fidelity and semantic alignment.} FID measures distributional similarity between generated and real panoramas~\cite{feng2025jopano,yang2024layerpano3d,zheng2025panorama,tandit2025}; PanFusion~\cite{zhang2024taming} shows that FAED (Fréchet Auto-Encoder Distance), computed in the latent space of an auto-encoder trained on the target data, tracks panorama quality more closely than Inception features do.
Text-to-panorama generation reports CLIP-Score between prompt and output~\cite{team2025hunyuanworld,feng2025jopano,yang2024layerpano3d,tandit2025}, image-conditioned tasks report CLIP-I between input and output~\cite{team2025hunyuanworld}, and no-reference quality is scored with BRISQUE and NIQE~\cite{team2025hunyuanworld,yang2024layerpano3d,yang2025matrix}, complemented by aesthetic scorers and human preference studies, for instance the Average User Ranking protocol of LayerPano3D~\cite{yang2024layerpano3d}.
Layout-conditioned synthesis measures adherence by 3D and 2D IoU against layouts recovered with HorizonNet~\cite{sun2019horizonnet}, following PanFusion~\cite{zhang2024taming}.

$\blacktriangleright$ \textit{Panorama-specific consistency.} Projection-induced distortion and boundary discontinuities call for metrics of their own.
JoPano~\cite{feng2025jopano} scores cubemap seams with Seam-SSIM (SSIM between adjacent edge pixels of neighboring faces) and Seam-Sobel (gradient discontinuity along the seam); PanoDecouple~\cite{zheng2025panorama} introduces Distort-CLIP, a CLIP variant trained to detect panoramic distortion, so that the ``visual cheating'' of models that improve image quality while degrading geometry becomes measurable; and loop consistency across the 360-degree wrap-around is assessed by PSNR in overlapping regions~\cite{tang2023MVDiffusion} or across the left--right boundary~\cite{panflow,xia2025panowan}.

$\blacktriangleright$ \textit{Reconstruction and novel-view synthesis.} Because pixel density on the sphere is non-uniform, WS-PSNR (Weighted Spherical PSNR) weights pixels by latitude and replaces plain PSNR for equirectangular outputs~\cite{pansplat,chen2023panogrf}, complemented by SSIM and LPIPS~\cite{yang2024layerpano3d,pansplat,yang2025matrix}.

$\blacktriangleright$ \textit{Video quality and spatiotemporal dynamics.} Panoramic videos are often projected to several perspective views (\eg, four views at azimuths $\{0\degree, 90\degree, 180\degree, 270\degree\}$ with a $90\degree$ field of view) so that human raters and perspective-domain tools can be used~\cite{team2025hunyuanworld,yang2024layerpano3d,tan2024imagine360}.
Q-Align~\cite{wu2023qalign}, a large multimodal model trained to score visual quality on discrete text-defined levels, is increasingly used for images and videos alike~\cite{team2025hunyuanworld,tan2024imagine360}; FVD~\cite{unterthiner2018fvd}, computed in the feature space of an inflated 3D ConvNet, is the standard distributional metric for video~\cite{panflow,xia2025panowan}; VBench~\cite{huang2023vbench} adds imaging quality, aesthetic quality, motion smoothness, and subject consistency over projected views; and trajectory-conditioned generation is scored by frame-to-frame consistency and by the range of camera motion it sustains~\cite{yin2025panoworld}.
Two limits remain. Ground truth is scarce for dynamic panoramic video with depth and camera annotation, so evaluation falls back on projected perspective views or user studies; and world-level properties such as loop consistency over long trajectories and action-conditioned spatial coherence have no dedicated benchmark yet.


\subsubsection{Challenges and Trends}
\label{subsubsec:p3_s8_2_4_challenges}

\paragraph{Unique Challenges in Panoramic Synthesis.}
Modeling a complete omnidirectional environment imposes constraints that NFoV synthesis never meets.
The first is \textbf{distortion-aware learning}: ERP stretches content near the poles, so visual information is non-uniformly distributed over the image and convolutional kernels designed for translation invariance on the plane struggle with the varying distortion, which motivates spherical manifold convolutions and latitude-aware architectures~\cite{sun2025spherical,zheng2025panorama}.
The second is \textbf{seamless boundary consistency}: wrap-around coherence across the longitudinal boundary must be preserved to avoid visible seams, typically through circular padding or latent rotation during denoising~\cite{panflow,xia2025panowan}.
The third is \textbf{scaling structured generation} to 4K and beyond while keeping global geometric and semantic consistency over the whole sphere, which pushes toward memory-efficient attention and hierarchical spherical representations~\cite{pansplat,liu2025dynamicscaler}.

\paragraph{From Panorama to World Models.}
The clearest trend is the move from static spherical snapshots toward \textbf{interactive 3D world modeling}, in which the panorama is the interface through which an embodied agent reasons about a complete environment rather than a fragment of one.
The systems that lift a generated panorama into an explorable world (HunyuanWorld~\cite{team2025hunyuanworld}, Matrix-3D~\cite{yang2025matrix}) and World Labs' Marble~\cite{Feifei2025Marble} show what panoramic priors make possible, and \textbf{temporal panoramic modeling} supplies the 4D half: generating temporally coherent panoramic video requires motion dynamics that respect spherical topology and occlusion boundaries~\cite{panflow,yin2025panoworld}.
Integrated with interactive systems, such panoramic world simulations become a basis for embodied interaction, long-range spatial reasoning, and environment simulation in robotics and immersive media.

\paragraph{Representation Innovations.}
Efficiency and fidelity both hinge on the omnidirectional representation.
\textbf{Efficient omnidirectional representations} such as the Fibonacci-lattice Gaussian pyramids of PanSplat~\cite{pansplat} render at high fidelity without the redundancy of uniform grids, and \textbf{multi-modal panoramic fusion}, which binds 360-degree RGB to depth, semantic labels, and geometric proxies such as surface normals~\cite{yang2024layerpano3d,scene4u}, makes the synthesized environment both more realistic and more editable, closing part of the gap between 2D generative priors and full 3D scene intelligence.




\subsection{3D Scene Editing}
\label{subsec:p3_s8_2_5}

The preceding subsections build 3D content; 3D scene editing takes an existing asset and modifies it under user control while preserving everything the edit should leave untouched.
Such incremental control is exactly what spatial intelligence asks of a 3D world model: an agent that must correct a mistake, personalize an environment, or keep a scene consistent as its task unfolds cannot afford to regenerate the world from scratch every time it changes.
Technically, 3D editing lifts the 2D inpainting-and-editing machinery of §\ref{subsec:p3_s8_1_3} from the image plane to a multi-view-consistent 3D representation, adding a cross-view coherence requirement that has no analogue in the 2D case; the five method families below are distinguished by how the edit signal is injected into that representation.


\begin{figure}[!ht]
    \centering
    \includegraphics[width=0.99\linewidth]{figures/part3/section8/3d_editing_tasks.png}
    \caption{Demonstrations of 3D editing tasks. (a) appearance editing, (b) object transformation, (c) shape deformation, (d) scene inpainting, and (e) creative editing. Figure adapted from~\cite{zhu2026survey3D}.}
    \label{fig:3D_editing_tasks}
\end{figure}

\subsubsection{Task Definition}
\label{subsubsec:p3_s8_2_5_task}

3D scene (and object) editing modifies an existing 3D asset or reconstructed scene according to a user intent (typically language or reference images, as in Figure~\ref{fig:3D_editing_tasks}) while preserving 3D consistency for novel-view rendering.
Unlike text-to-3D generation, editing operates on a given representation (\eg, NeRF, 3D Gaussian Splatting, or meshes) and requires (i) locality and identity preservation for unedited regions, (ii) multi-view (and potentially temporal) coherence, and (iii) controllability over the edited attributes, whether geometry, appearance, material, or semantics.

Formally, 3D editing is the 3D-lifted counterpart of the 2D inpainting-and-editing operator of §\ref{subsec:p3_s8_1_3} (Eq.~\eqref{eq:t2i_edit_objective}).
Given an input 3D representation $\mathbf{r}^{\text{in}}\in\mathcal{R}$ drawn from the same representation space $\mathcal{R}$ as §\ref{subsec:p3_s8_2_1}, an edit intent $\mathbf{c}\in\mathcal{C}\supseteq\{\text{text instruction},\text{reference image},\text{mask / region},\text{scribble}\}$, and an optional 3D region of interest $\mathcal{M}\subseteq\mathbb{R}^3$, the goal is an editing operator
\begin{equation}
E_\theta:\mathcal{R}\times\mathcal{C}'\rightarrow\mathcal{R},
\qquad
\hat{\mathbf{r}} = E_\theta\big(\mathbf{r}^{\text{in}},\,\kappa(\mathbf{c})\big),
\label{eq:p3_s8_2_5_objective}
\end{equation}
trained or optimized to minimize an objective that couples three competing terms,
\begin{equation}
\mathcal{L}_{\text{edit}}
=
\underbrace{\mathcal{L}_{\text{intent}}\big(\hat{\mathbf{r}},\mathbf{c};\mathcal{M}\big)}_{\text{edit-region adherence}}
\;+\;
\lambda_{\text{p}}\underbrace{\mathcal{L}_{\text{preserve}}\big(\hat{\mathbf{r}},\mathbf{r}^{\text{in}};\mathbb{R}^3\!\setminus\!\mathcal{M}\big)}_{\text{identity preservation}}
\;+\;
\lambda_{\text{c}}\underbrace{\mathcal{L}_{\text{consist}}\big(\{\mathrm{Render}(\hat{\mathbf{r}},\pi_v)\}_v\big)}_{\text{multi-view coherence}},
\label{eq:p3_s8_2_5_loss}
\end{equation}
where $\kappa$ is the condition encoder of §\ref{subsec:p3_s8_2_1}, $\mathrm{Render}$ is the differentiable renderer of §\ref{subsec:p3_s8_2_3}, and $\lambda_{\text{p}},\lambda_{\text{c}}$ trade off preservation and consistency against edit strength.
The limiting case that anchors Eq.~\eqref{eq:p3_s8_2_5_objective} to its 2D origin is instructive: when $\mathcal{R}$ collapses to a single image ($V=1$, no rendering), $E_\theta$ reduces exactly to the 2D editing operator of Eq.~\eqref{eq:t2i_edit_objective} and the multi-view-consistency term vanishes, so 3D editing is 2D editing plus a cross-view coherence constraint.


\subsubsection{Methodology}
\label{subsubsec:p3_s8_2_5_methodology}

Recent 3D object and scene editing methods can be categorized by \emph{how} the edit signal of Eq.~\eqref{eq:p3_s8_2_5_objective} is injected into the 3D representation, from lifting 2D edits view by view, through enforcing cross-view consistency, to operating directly on 3D primitives, reasoning with language models, and amortizing the edit into a single feed-forward pass.
Figure~\ref{fig:p3_s8_2_5_mechanisms} schematizes the five families and how each realizes the operator $E_\theta$.

\begin{figure}[!ht]
    \centering
    \includegraphics[width=0.99\linewidth]{figures/part3/section8/Five_method_families_for_3D_scene_editing.pdf}
    \caption{Five method families for 3D scene editing: (a) 2D-lifted iterative editing, (b) multi-view-consistent guidance, (c) 3D-native primitive editing, (d) LLM-integrated editing, and (e) efficient and training-free editing. The chair edit (top) is illustrative, not a measured result.}
    \label{fig:p3_s8_2_5_mechanisms}
\end{figure}

\paragraph{2D-lifted Iterative Editing.}
A dominant paradigm inherits from instruction-based NeRF editing~\cite{bao2023sine,khalid2024latenteditor,mikaeili2023sked,zhou2023repaintnerf}: the method iterates between rendering views from the current scene, applying a 2D diffusion-based editor to them, and updating the 3D representation to match the edited guidance under geometric regularization.
Formally, this family instantiates $E_\theta$ of Eq.~\eqref{eq:p3_s8_2_5_objective} as a \emph{per-scene} optimization that drives the rendered views toward the output of a \emph{frozen} 2D editor $\mathcal{E}_{\text{2D}}$ (the InstructPix2Pix-style operator of §\ref{subsec:p3_s8_1_3}),
\begin{equation}
\hat{\mathbf{r}} = \arg\min_{\mathbf{r}\in\mathcal{R}} \sum_{v} \big\| \mathrm{Render}(\mathbf{r},\pi_v) - \mathcal{E}_{\text{2D}}\!\big(\mathrm{Render}(\mathbf{r}^{\text{in}},\pi_v),\,\mathbf{c}\big) \big\|^2 + \lambda\,\Omega(\mathbf{r}),
\label{eq:p3_s8_2_5_lift}
\end{equation}
so that the 2D edits are \emph{lifted} into the 3D representation and the $\arg\min$ is re-solved from scratch for every scene; Figure~\ref{fig:p3_s8_2_5_mechanisms}(a) depicts this render-edit-update loop.
Instruct-NeRF2NeRF~\cite{haque2023instructnerf2nerf}, EditNeRF~\cite{liu2021editingnerf}, Free-Editor~\cite{karim2024freeeditor}, and CTRL-D~\cite{he2025ctrld} lift instruction-driven 2D edits into the radiance field this way, and with the rise of 3D Gaussian Splatting (3DGS)~\cite{kerbl20233d} the same principle moved to Gaussian primitives: Instruct-GS2GS~\cite{vachha2024instructgs2gs} iteratively edits the training images with InstructPix2Pix~\cite{brooks2023instructpix2pixlearningfollowimage} and re-optimizes the 3DGS to reflect the instruction while preserving 3D structure~\cite{wang2024viewedit,wynn2025morpheus}.
The recipe is simple and borrows the full strength of 2D editors, but it is slow, prone to view-dependent artifacts, and only as good as the stability of the 2D editor and the coverage of the rendered guidance views.

\paragraph{Multi-view-consistent Guidance.}
Because independently edited views can disagree (inconsistent textures, hallucinated geometry), recent work~\cite{chen2024dge} enforces multi-view consistency at the \emph{guidance} level rather than relying on the 3D update to reconcile conflicting edits.
Formally, rather than editing each view independently as in Eq.~\eqref{eq:p3_s8_2_5_lift}, this family makes the guidance images mutually consistent before the lift,
\begin{equation}
\{\hat{I}_v\} = \arg\min_{\{I_v\}} \sum_{v} \mathcal{L}_{\text{intent}}(I_v,\mathbf{c}) \;+\; \beta\!\!\sum_{(u,v)} \big\| I_u - \mathcal{W}_{u\to v}(I_v) \big\|^2,
\label{eq:p3_s8_2_5_mvc}
\end{equation}
where $\mathcal{W}_{u\to v}$ warps view $v$ into view $u$ using geometry or depth derived from the current 3D estimate, penalizing cross-view disagreement; the consolidated $\{\hat{I}_v\}$ then drive the 3D update of Eq.~\eqref{eq:p3_s8_2_5_lift}, as in Figure~\ref{fig:p3_s8_2_5_mechanisms}(b).
GaussCtrl~\cite{wu2024gaussctrl} edits several rendered views in a coupled fashion rather than sequentially, improving multi-view consistency for text-driven 3DGS editing~\cite{wang2024viewedit,xia2025towardsConsistent}, and VcEdit~\cite{wang2024viewedit} integrates the 3DGS into the image-editing process itself with consistency modules that counter cross-view disagreement and mode collapse.
A complementary line transfers a \emph{single-view} edit to other views via geometry: \citet{gomel2024diffusionBasedAttention} extracts diffusion attention features from a reference edit, warps them across views using depth from the Gaussian splats, and injects the warped features to propagate a coherent edit.
These methods preserve 3D coherence more robustly but still struggle with large non-rigid shape changes or severe occlusion, where reliable correspondence and geometry are hard to come by.

\paragraph{3D-native Primitive Editing.}
To reduce reliance on unstable 2D guidance and enable interactive use, 3D-native approaches operate directly on explicit scene primitives.
Formally, this family bypasses the rendering loop of Eq.~\eqref{eq:p3_s8_2_5_lift} entirely: a grounding or segmentation operator $\mathcal{G}$ selects the target subset of primitives and a direct geometric transform $\mathcal{T}$ applies the edit,
\begin{equation}
\hat{\mathbf{r}} = \mathcal{T}\!\big(\mathbf{r}^{\text{in}};\,\mathcal{G}(\mathbf{r}^{\text{in}},\mathbf{c})\big), \qquad \mathcal{T}\in\{\text{add},\,\text{delete},\,\text{move},\,\text{recolor},\,\text{replace}\},
\label{eq:p3_s8_2_5_native}
\end{equation}
acting on explicit Gaussians or mesh facets so that the unedited primitives are preserved \emph{by construction}: the identity-preservation term of Eq.~\eqref{eq:p3_s8_2_5_loss} is satisfied structurally rather than by a penalty; Figure~\ref{fig:p3_s8_2_5_mechanisms}(c) shows this pathway.
3DSceneEditor~\cite{yan20243dsceneeditor} runs a fully 3D pipeline, combining instance segmentation and zero-shot grounding to localize the target and then performing Gaussian-level addition, deletion, relocation, recoloring, and replacement with strong controllability in real time, and InterGSEdit~\cite{wen2025intergsedit} adds user-in-the-loop key-view selection and a geometry-consistent attention prior for iterative, preference-guided 3DGS editing with control over the edit degree.
More broadly, GaussianEditor and 3DSceneEditor~\cite{chen2024gaussianeditor,yan20243dsceneeditor} impose higher-level structure on the Gaussians, grouping them by object identity or semantic coherence (typically through zero-shot grounding or segmentation~\cite{hu2024semanticanythinggaussians,cen2025segmentany3dgaussians} followed by anchoring) so that an edit can be formulated at the level of an instance rather than of individual primitives, which localizes the target more precisely and protects the background.

\paragraph{LLM-integrated Editing.}
Large language models have been integrated into 3D editing pipelines not merely as text encoders but as high-level reasoning and control modules.
Formally, the LLM decomposes the instruction into an ordered program of geometry-aware operators applied to the asset,
\begin{equation}
\hat{\mathbf{r}} = (o_K \circ \cdots \circ o_1)(\mathbf{r}^{\text{in}}), \qquad (o_1,\dots,o_K) = \mathrm{LLM}\big(\mathbf{c},\,\mathbf{r}^{\text{in}}\big),
\label{eq:p3_s8_2_5_llm}
\end{equation}
so that high-level reasoning (attribute inference, edit planning, dependency resolution) is separated from the low-level operators $o_k$, often the primitive transforms of Eq.~\eqref{eq:p3_s8_2_5_native}, that it dispatches; Figure~\ref{fig:p3_s8_2_5_mechanisms}(d) depicts this plan-then-execute pathway.
ShapeLLM-Omni~\cite{ye2025shapellmOmni} casts 3D editing as language-grounded reasoning in which the LLM performs shape understanding, attribute inference, and edit planning before invoking geometry-aware operators, and LL3M~\cite{lu2025ll3m} extends this to multimodal 3D manipulation requiring compositional understanding of geometry, appearance, and spatial relations.
A complementary line~\cite{ma2025feedforward3DEditing} uses the LLM to guide or parameterize image-to-3D or text-to-3D backbones in a feed-forward setting rather than a per-scene optimization, and LLMs also serve as coordination agents for global consistency~\cite{huang2024blenderalchemy,zhang2025scenelanguage}: \citet{xia2025towardsConsistent} shows that reasoning over object identities, relations, and edit dependencies reduces the conflicts that arise in purely diffusion- or optimization-based pipelines when several objects and attributes must be edited at once.

\paragraph{Efficient and Training-free Editing.}
A final trend reduces test-time cost through feed-forward reconstruction or reuse of pretrained models~\cite{erkocc2025preditor3d,li2025voxhammer}. These strategies are distinct: a training-free editing pipeline can still involve iterative diffusion inference. The amortized feed-forward case evaluates the editing operator of Eq.~\eqref{eq:p3_s8_2_5_objective} directly,
\begin{equation}
\hat{\mathbf{r}} = E_\theta\big(\mathbf{r}^{\text{in}},\,\kappa(\mathbf{c})\big) \quad \text{(amortized, no per-scene optimization)},
\label{eq:p3_s8_2_5_ff}
\end{equation}
eliminating the per-scene $\arg\min$ of Eq.~\eqref{eq:p3_s8_2_5_lift} by amortizing inference into a trained network. Training-free alternatives instead reuse pretrained components, potentially through single-view guidance, mesh proxies, or iterative native-3D latent editing; Figure~\ref{fig:p3_s8_2_5_mechanisms}(e) distinguishes these routes from the per-scene optimization loop in panel~(a).
Free-Editor~\cite{karim2024freeeditor} edits a NeRF zero-shot from text without scene-specific retraining, countering multi-view style inconsistency with a ``single-view editing'' scheme and an Edit Transformer; on the scene-operation side (insert, replace, delete), FreeEdit~\cite{madhavaram2025towardsTrainingFree} uses a mesh proxy to perform training-free edits guided by foundation models at near real-time rates in complex indoor scenes.
These methods improve usability but typically trade away edit fidelity and coverage, especially for fine-grained geometry, and depend strongly on robust grounding and localization in 3D.


\subsubsection{Benchmarks}
\label{subsubsec:p3_s8_2_5_benchmarks}

Several benchmarks have recently been introduced for systematic evaluation.
3D-Alpaca-Editing~\cite{ye2025shapellmOmni} focuses on instruction-following 3D editing, assessing whether a model interprets natural-language commands correctly while preserving unedited content.
CMD~\cite{li2025cmd} emphasizes compositional and multi-step instructions, where reasoning and planning matter.
Edit3D-Bench~\cite{li2025voxhammer} is a general-purpose suite covering geometry, appearance, and object-level operations, with 3D edit regions annotated on 100 source models (300 edits), which allows fair comparison across editing paradigms.
3DEditVerse~\cite{xia2025towardsConsistent} is a large-scale paired benchmark built for robustness and generalization across varied objects, scenes, and intents.
Together they standardize evaluation in terms of instruction compliance, edit locality, and scalability; Table~\ref{tab:p3_s8_2_5_datasets} summarizes their scale and the evaluation axes each one annotates.

\begin{table}[!ht]
\centering
\fontsize{8}{10}\selectfont
\setlength{\tabcolsep}{1.0mm}
\caption{Representative 3D editing benchmarks. ``Edit region'' indicates whether a localized edit mask / region is annotated; ``Consistency'' and ``Harmony'' indicate whether multi-view consistency and edit harmonization are evaluated. The criteria follow the four requirements set out with 3DEditVerse~\cite{xia2025towardsConsistent}.}
\label{tab:p3_s8_2_5_datasets}
\begin{tabular}{L{0.24\textwidth} C{0.11\textwidth} C{0.10\textwidth} C{0.11\textwidth} C{0.12\textwidth} C{0.12\textwidth} C{0.10\textwidth}}
\toprule
\textbf{Dataset} & \textbf{Training size} & \textbf{Testing size} & \textbf{Edit region} & \textbf{Train / test split} & \textbf{Consistency} & \textbf{Harmony} \\
\midrule
3D-Alpaca-Editing~\cite{ye2025shapellmOmni} & ${\sim}70{,}000$ & -- & \cmark & \crossmark & \crossmark & \cmark \\
CMD~\cite{li2025cmd} & ${\sim}40{,}000$ & 50 & \crossmark & \cmark & \cmark & \crossmark \\
Edit3D-Bench~\cite{li2025voxhammer} & -- & 300 & \cmark & \crossmark & \cmark & \cmark \\
3DEditVerse~\cite{xia2025towardsConsistent} & 116,309 & 1,500 & \cmark & \cmark & \cmark & \cmark \\
\bottomrule
\end{tabular}
\end{table}

\paragraph{Evaluation Metrics.}
Evaluation considers geometric fidelity, appearance quality, and semantic alignment with the instruction.
Geometry is measured by Chamfer Distance (CD)~\cite{ma2025feedforward3DEditing}, lower being better, and by F-score~\cite{xia2025towardsConsistent}, which captures surface overlap at a distance threshold.
Texture and appearance are assessed by comparing renderings of the edited representation against target or inpainted images with SSIM~\cite{wang2004image}, LPIPS~\cite{zhang2018unreasonable}, and DreamSim~\cite{fu2023dreamsim}~\cite{gao20253dmeshedit}.
Visual--text alignment is scored with CLIP- or DINO-based embedding similarity~\cite{shu2025gaussedit}, and GPTEval3D~\cite{wu2024gpt4vt23dgen} uses GPT-4V to compare multi-view renderings of edited shapes on text alignment, 3D plausibility, and texture detail.
To measure the edit relative to the input rather than in isolation, directional CLIP (CLIP$_{\text{dir}}$)~\cite{wynn2025morpheus} checks that the change in image embedding, computed over multi-view renderings of input and output, follows the change in text embedding.

Edit cost is the clearest axis separating the five families.
Instruct-NeRF2NeRF~\cite{haque2023instructnerf2nerf} takes on the order of 30--90 minutes per edit; GaussianEditor~\cite{chen2024gaussianeditor} brings this under 20 minutes by operating on explicit Gaussians; PrEditor3D~\cite{erkocc2025preditor3d} edits a shape in a few minutes, and VoxHammer~\cite{li2025voxhammer} performs precise local edits training-free in a native 3D latent space.
On quality the field reports DINO-I and CLIP-T on Edit3D-Bench~\cite{li2025voxhammer}, directional CLIP~\cite{wynn2025morpheus}, and GPT-4V judgment~\cite{wu2024gpt4vt23dgen}; the 3D-native and feed-forward families lead on locality and speed, while 2D-lifted methods keep an edge on open-vocabulary appearance edits where a strong 2D editor does the work.
The speed of feed-forward editing comes at a cost in fine-grained geometry control and a dependence on robust 3D grounding, taken up below.


\subsubsection{Challenges and Trends}
\label{subsubsec:p3_s8_2_5_challenges}

Several problems remain genuinely open and define the current frontier of 3D editing.

\paragraph{Geometry-level, not merely appearance, editing.}
Most current methods edit color and texture well but treat geometry as fixed: precise, instruction-driven changes to shape, including thin, transparent, or high-frequency structures, remain unreliable, and edits that should alter the underlying surface often only repaint it. Closing this gap requires editors that operate on geometry and appearance jointly rather than deferring geometry to a frozen reconstruction.

\paragraph{Large non-rigid and structural edits.}
Adding, removing, or substantially deforming structure while preserving the identity of the rest of the asset is still difficult: drag-driven, as-rigid-as-possible Gaussian editors~\cite{han2025arapgs} handle moderate deformation, but large non-rigid motion and topology changes frequently lose original detail. Robust structural and compositional editing, including relation-aware multi-object edits, remains an active frontier.

\paragraph{Material--lighting disentanglement and relightable editing.}
Because most representations bake illumination into appearance, an edit that changes an object's material or color is not guaranteed to stay consistent when the scene is relit. Relightable representations that decompose geometry, BRDF, and incident light~\cite{gao2024relightable} are a prerequisite for edits that survive relighting and for inserting edited content into new lighting, yet stable material--light decomposition under editing is not solved.

\paragraph{The consistency bottleneck has moved to 3D grounding.}
Now that native-3D editing handles appearance coherence, the limiting factor is localization: methods such as VoxHammer~\cite{li2025voxhammer} are highly sensitive to 3D mask accuracy and degrade rapidly when the target region is imprecisely grounded. Robust open-vocabulary 3D grounding and segmentation~\cite{hu2024semanticanythinggaussians,cen2025segmentany3dgaussians}, ideally without user-provided masks~\cite{xia2025towardsConsistent}, is therefore the new gate on edit quality, and it is most acute in dynamic scenes, where the edited region must be tracked with temporal identity preserved across frames~\cite{he2025ctrld}.

Two pressures cut across all four. Evaluation remains immature: the benchmarks of §\ref{subsubsec:p3_s8_2_5_benchmarks} are small (Edit3D-Bench annotates 100 source models and 300 edits) and no metric panel is yet standard, so progress on the harder edits above is difficult to measure reliably. And the strongest forward trend is unification, bringing generation, segmentation, grounding, and editing into a single structured pipeline~\cite{park2026generative} with interfaces that go beyond text to reference images~\cite{shu2025gaussedit}, masks~\cite{barda2025instant3dit}, sketches~\cite{mikaeili2023sked}, and object-level operations, so that an agent can specify and verify a geometry-, material-, and relighting-consistent edit rather than only a repaint.



% \subsection{Summary}
% \label{subsec:p3_s8_2_summary}

% Table~\ref{tab:p3_s8_2_subfield_comparison} places the five subfields side by side: the input each consumes, the operator it instantiates over the shared 3D-representation space $\mathcal{R}$, the paradigms it surveys, the methods anchoring each, and the benchmarks on which it reports.

% \begin{table}[t]
%     \centering
%     \fontsize{8}{10}\selectfont
%     \setlength{\tabcolsep}{1.0mm}
%     \caption{The five subfields of 3D spatial reconstruction and generation; each operator acts on the shared 3D-representation space $\mathcal{R}$.}
%     \label{tab:p3_s8_2_subfield_comparison}
%     \begin{tabular}{L{0.13\textwidth} L{0.20\textwidth} L{0.18\textwidth} L{0.19\textwidth} L{0.22\textwidth}}
%         \toprule
%         \textbf{Subfield} & \textbf{Input / operator} & \textbf{Dominant paradigms} & \textbf{Canonical methods} & \textbf{Benchmarks / metrics} \\
%         \midrule
%         §\ref{subsec:p3_s8_2_1} 3D object synthesis & Text / image / multi-view $\to$ asset; $G_\theta^{\text{3D}}:\mathcal{C}'\!\to\!\mathcal{R}$ (Eq.~\eqref{eq:p3_s8_2_1_objective}) & 3D-native generation; SDS distillation; multi-view generation and lifting; feed-forward reconstruction and 3D latent diffusion & GET3D; DreamFusion; Zero-1-to-3, MVDream; LRM, TRELLIS, Hunyuan3D & GT23D-Bench, MATE-3D, Gen3DEval; CLIP alignment, Chamfer / F-score, LPIPS \\
%         §\ref{subsec:p3_s8_2_2} 3D scene synthesis & Text / layout / graph / sparse views / trajectory $\to$ scene tuple; $G_\theta^{\text{scene}}$ (Eq.~\eqref{eq:p3_s8_2_2_objective}) & Layout / scene-graph conditioning; view-to-3D; view-to-view; video-model-based & DiffuScene, GALA3D; DreamScene; LVSM, SEVA; CameraCtrl, ViewCrafter & 3D-FRONT, ScanNet, RealEstate10K; layout adherence, PSNR / LPIPS, FVD, pose error \\
%         §\ref{subsec:p3_s8_2_3} 3D scene reconstruction & Posed or unposed images, depth, point clouds $\to$ 3D; $G_\theta^{\text{rec}}:\mathcal{O}\!\to\!\mathcal{R}^{\text{rec}}$ (Eq.~\eqref{eq:p3_s8_2_3_objective}) & Iterative (SfM--MVS, NeRF, 3DGS); feed-forward; generative complete reconstruction & NeRF, 3DGS; DUSt3R, MASt3R, VGGT; ReconFusion, Bolt3D & DTU, Tanks-and-Temples, Mip-NeRF360; PSNR / SSIM / LPIPS, Chamfer, pose AUC \\
%         §\ref{subsec:p3_s8_2_4} 3D panorama synthesis & Text / NFoV image / layout / viewpoint $\to$ panorama $\to$ 3D; $G_\theta^{\text{pano}}$ $+$ lift $\Lambda$ (Eq.~\eqref{eq:p3_s8_2_4_objective}) & Panoramic image synthesis; panoramic video; panoramic reconstruction; panorama-based scene generation & PanFusion, CubeDiff; 360DVD, Imagine360; PanSplat; LayerPano3D, HunyuanWorld & Structured3D, SUN360, Matterport3D; FID / FAED, WS-PSNR, seam and loop consistency \\
%         §\ref{subsec:p3_s8_2_5} 3D scene editing & Existing asset $+$ edit intent $\to$ edited asset; $E_\theta:\mathcal{R}\!\times\!\mathcal{C}'\!\to\!\mathcal{R}$ (Eq.~\eqref{eq:p3_s8_2_5_objective}) & 2D-lifted iterative; multi-view-consistent guidance; 3D-native primitive; LLM-integrated; feed-forward / training-free & Instruct-NeRF2NeRF; GaussCtrl, VcEdit; 3DSceneEditor; ShapeLLM-Omni; VoxHammer & Edit3D-Bench, 3DEditVerse; Chamfer / F-score, CLIP$_{\text{dir}}$, GPTEval3D \\
%         \bottomrule
%     \end{tabular}
% \end{table}

% The subfields partition not by output domain, since all five produce an element of the same $\mathcal{R}$, but by the direction and source of the map: synthesis (§\ref{subsec:p3_s8_2_1}, §\ref{subsec:p3_s8_2_2}) maps high-level priors into $\mathcal{R}$, reconstruction (§\ref{subsec:p3_s8_2_3}) maps sensor observations into it, panorama synthesis (§\ref{subsec:p3_s8_2_4}) inserts a 2D intermediate before the map, and editing (§\ref{subsec:p3_s8_2_5}) maps $\mathcal{R}$ back into itself.
% Synthesis and reconstruction are not disjoint. As observations grow sparse, the reconstruction operator of Eq.~\eqref{eq:p3_s8_2_3_objective} must hallucinate the unobserved regions and collapses into the generative posterior of Eq.~\eqref{eq:p3_s8_2_1_objective}, so the complete-reconstruction regime of §\ref{subsec:p3_s8_2_3} and the generative regime of §\ref{subsec:p3_s8_2_1} approach one problem from opposite sides, and the formal distinction between the two operators dissolves into a single estimator that consumes whatever observations are available and falls back on a learned prior wherever they are not.
% Across all five subfields the 2D prior of §\ref{sec:p3_s8_1} is reused as a prior to be distilled rather than as a terminal generator (score distillation, multi-view and panorama diffusion, frozen 2D editors lifted into 3D), and all inherit the same failure mode, the 2D$\leftrightarrow$3D consistency gap (Janus artifacts, view drift, cubemap seams, edit flicker) whose suppression motivates the multi-view-consistency machinery surveyed throughout the section.
% Evaluation, finally, has not followed the convergence seen in §\ref{sec:p3_s8_1}: no cross-subfield suite exists, each subfield reports a paradigm-specific tuple over geometry, appearance, and constraint adherence, and a layout-adherence score from §\ref{subsec:p3_s8_2_2} cannot be ranked against a novel-view PSNR from §\ref{subsec:p3_s8_2_3}; a human-aligned protocol that spans the subfields is the clearest open need.

% The temporal extension of these volumetric priors (generating spatially \emph{and} temporally coherent 3D content, where a camera or the scene itself moves through time) is the subject of video-based spatial generation in §\ref{sec:p3_s8_3}.

% \newpage
\section{Video Spatial Generation}
\label{sec:p3_s8_3}

\FramedInitial{V}{ideo} spatial generation is the task of synthesizing a temporally coherent image sequence (a video clip) from a high-level prompt together with optional spatial, temporal, and other conditioning signals.
It has advanced substantially in recent years, evolving from early generative adversarial network (GAN)-based approaches~\cite{goodfellow2014gan} through autoregressive token models to the diffusion- and diffusion-transformer-based systems that now dominate the field~\cite{polyak2024movie,hong2022cogvideo,wan2025wan,sora2024}.
This rapid progress has stimulated a wide range of downstream tasks, extending beyond basic text-to-video generation to encompass fine-grained control mechanisms (spatial, \eg, image, pose, depth, landmark, sketch, and bounding box; temporal, \eg, optical flow, trajectory, camera, and motion; and other conditional controls, \eg, style, audio, bird's-eye view) as well as sophisticated video editing capabilities.

Within the arc of §\ref{sec:generation}, video spatial generation is the \emph{temporal} extension of the priors built in the preceding two sections.
Whereas §\ref{sec:p3_s8_1} learned a 2D semantic-to-pixel prior on a single image plane and §\ref{sec:p3_s8_2} lifted that prior into a static volumetric representation, this section adds the time axis: the model must place matter and appearance correctly \emph{and} make them evolve plausibly as the camera or the scene itself moves through time.
Appearance, geometry, and dynamics must now be \emph{jointly} consistent: a generated clip is judged not only frame by frame but on object permanence, motion plausibility, and multi-view coherence under camera motion, which makes video the natural substrate for the dynamic world models and simulators that consume spatial structure downstream.

Formally, we cast every system in this section as a conditional spatio-temporal sampler that extends the latent-generator framework of §\ref{subsec:p3_s8_1_1}.
Let $\mathcal{V}=[0,1]^{T\times H\times W\times 3}$ be the space of $T$-frame clips and $\mathcal{Z}^{\text{vid}}$ a compressed spatio-temporal latent space supplied by a video autoencoder that extends the image autoencoder of Eq.~\eqref{eq:t2i_autoencoder} along the time axis.
Writing $\phi$ for the text encoder of Eq.~\eqref{eq:t2i_generator} and $\kappa:\mathcal{C}\!\to\!\mathcal{C}'$ for a generic encoder of an auxiliary control signal $\mathbf{c}\in\mathcal{C}$, the core learning problem is to train
\begin{equation}
G_\theta^{\text{vid}}:\mathcal{E}\times\mathcal{C}'\rightarrow\mathcal{Z}^{\text{vid}},\qquad
\hat{\mathbf{z}}^{1:T}\sim p_\theta\!\big(\mathbf{z}^{1:T}\,\big|\,\phi(\mathbf{x}),\,\kappa(\mathbf{c})\big),
\label{eq:p3_s8_3_objective}
\end{equation}
so that the decoded clip $\hat{\mathbf{v}}=D(\hat{\mathbf{z}}^{1:T})$ is photorealistic, temporally coherent, and faithful to both the prompt $\mathbf{x}$ and the control $\mathbf{c}$.
Two limiting cases anchor the framework: with no auxiliary control ($\kappa=\varnothing$) Eq.~\eqref{eq:p3_s8_3_objective} reduces to text-to-video generation, and with a single frame ($T=1$) it recovers the image generator of Eq.~\eqref{eq:t2i_generator} exactly, so video spatial generation is precisely the image prior of §\ref{sec:p3_s8_1} unrolled in time and equipped with a richer conditioning interface.

The subsections are organized by the axis of control each introduces into Eq.~\eqref{eq:p3_s8_3_objective}.
Text-to-video generation (§\ref{subsec:p3_s8_3_1}) is the base case without auxiliary control; spatial control (§\ref{subsec:p3_s8_3_2}) adds a per-frame scaffold such as a pose or depth sequence; temporal control (§\ref{subsec:p3_s8_3_3}) adds a motion or camera signal that governs dynamics rather than appearance; §\ref{subsec:p3_s8_3_4} covers the remaining modalities and the frameworks that fuse several at once; and video editing (§\ref{subsec:p3_s8_3_5}) modifies an existing clip, lifting the 2D editing of §\ref{subsec:p3_s8_1_3} into time.
§\ref{subsec:p3_s8_3_6} compares the five and discusses the challenges they share.



\subsection{Text-to-Video Generation}
\label{subsec:p3_s8_3_1}

Text-to-video (T2V) generation~\cite{sora2024,polyak2024movie,hong2022cogvideo,wan2025wan} synthesizes photorealistic, temporally coherent video from a natural-language description.
It extends text-to-image (T2I) synthesis along time, so the model must produce plausible motion and state transitions as well as high-fidelity frames, and it does so with no explicit spatial scaffold: object permanence, motion plausibility, and view coherence under camera motion cannot be supplied from outside but must \emph{emerge} from the generator itself.
That is why the unconditional T2V model is the substrate on which the explicit control axes of §\ref{subsec:p3_s8_3_2}--§\ref{subsec:p3_s8_3_5} are built.
Its architecture has followed the image case, from generative adversarial networks (GANs) through autoregressive models to the diffusion models that dominate today, with progress driven as much by data as by architecture~\cite{chen2024panda70mcaptioning70mvideos,NEURIPS2024_57f6683e}.

\subsubsection{Task Definition}
\label{subsubsec:p3_s8_3_1_task}
Text-to-video generation is the base case of Eq.~\eqref{eq:p3_s8_3_objective} with no auxiliary control ($\kappa=\varnothing$): the generator samples a latent from the prompt embedding alone,
\begin{equation}
\hat{\mathbf{z}}^{1:T}\sim p_\theta\!\big(\mathbf{z}^{1:T}\,\big|\,\phi(\mathbf{x})\big),
\label{eq:p3_s8_3_1_objective}
\end{equation}
which is exactly the latent sampler $G_\theta$ of Eq.~\eqref{eq:t2i_generator} unrolled over $T$ frames in the spatio-temporal latent space $\mathcal{Z}^{\text{vid}}$.

\subsubsection{Methodology}
\label{subsubsec:p3_s8_3_1_methodology}
The three paradigms factorize the clip distribution $p_\theta(\mathbf{z}^{1:T}\mid\phi(\mathbf{x}))$ adversarially, causally over discrete tokens, or by iterative spatio-temporal denoising, and each changes its image counterpart of §\ref{subsec:p3_s8_1_1} in one structural respect; Figure~\ref{fig:p3_s8_3_1_paradigms} contrasts how each handles the temporal dimension.

\begin{figure}[!ht]
\centering
\includegraphics[width=\linewidth]{figures/part3/section8/text-to-video_generation.pdf}
\caption{Three mechanisms for text-to-video generation: (a) a GAN with spatial and temporal critics, a representative rather than universal design; (b) autoregressive prediction of discrete space--time tokens, with pale cells as generated context and the blue cell as the next token; and (c) diffusion with factorized spatial / temporal attention, where $s$ is the sampling step and $t$ is video time. The ball sequence is illustrative.}
\label{fig:p3_s8_3_1_paradigms}
\end{figure}

\noindent \textbf{Generative Adversarial Networks (GANs)}. The first T2V systems extended the image GANs of §\ref{para:p3_s8_1_1_gan} to clips. The video-specific change is in the discriminator: rather than scoring a single image, it must judge both per-frame realism \emph{and} the plausibility of motion across the sequence, which is naturally realized as a two-critic objective
\begin{equation}
\min_\theta\ \max_{\psi_s,\psi_t}\ \mathbb{E}_{\mathbf{v}\sim p_\text{data}}\!\big[\log D_{\psi_s}(\mathbf{v})+\log D_{\psi_t}(\mathbf{v})\big] + \mathbb{E}_{\mathbf{n}}\!\big[\log\!\big(1-D_{\psi_s}(\hat{\mathbf{v}})\big)+\log\!\big(1-D_{\psi_t}(\hat{\mathbf{v}})\big)\big],
\label{eq:p3_s8_3_1_gan}
\end{equation}
where $\hat{\mathbf{v}}=D\big(G_\theta(\mathbf{n},\phi(\mathbf{x}))\big)$ is the decoded clip, the spatial critic $D_{\psi_s}$ scores individual frames, and the temporal critic $D_{\psi_t}$ scores motion coherence over the clip, the term absent from the single-discriminator image objective of Eq.~\eqref{eq:t2i_gan_objective}. TGAN~\cite{saito2017temporal} and TGAN-v2~\cite{saito2020train} used a temporal generator to produce latent motion vectors from which an image generator rendered frames; MoCoGAN~\cite{tulyakov2018mocogan} and MoCoGAN-HD~\cite{tian2021good} disentangled content from motion, generating the motion code recurrently; DVD-GAN~\cite{clark2019adversarial} instantiated Eq.~\eqref{eq:p3_s8_3_1_gan} directly with separate spatial and temporal discriminators; and DIGAN~\cite{yu2022generating} used implicit neural representations for longer sequences. Even with architectures such as StyleGAN-V~\cite{skorokhodov2022stylegan}, GANs suffered from flicker, training instability, and mode collapse, which limited the diversity of their motion and pushed the field toward other paradigms.


\noindent \textbf{Autoregressive Models}. Autoregressive (AR) models recast T2V as the sequential prediction of §\ref{para:p3_s8_1_1_ar}, tokenizing a clip with a 3D Vector Quantized Variational Autoencoder (3D~VQ-VAE) and predicting each token from its predecessors with a large transformer. The video-specific change is the \emph{order} of the factorization: where the image model of Eq.~\eqref{eq:t2i_ar_factorisation} runs over the spatial tokens of a single frame, the video model factorizes over space \emph{and} time,
\begin{equation}
p_\theta\!\big(\mathbf{z}^{1:T}\,\big|\,\phi(\mathbf{x})\big)\;=\;\prod_{t=1}^{T}\prod_{n=1}^{N} p_\theta\!\big(z_{t,n}\,\big|\,\mathbf{z}^{<t},\,z_{t,<n},\,\phi(\mathbf{x})\big),
\label{eq:p3_s8_3_1_ar}
\end{equation}
so that each frame is conditioned on all previously generated frames $\mathbf{z}^{<t}$. This inter-frame causal dependency is the source of both the paradigm's strength (global temporal coherence) and its characteristic failure, the accumulation of per-step error over long roll-outs. VideoGPT~\cite{yan2021videogpt} demonstrated the approach for unconditional generation, N\"{U}WA~\cite{wu2022nuwa} and CogVideo~\cite{hong2022cogvideo} adapted it to text-to-video synthesis, and TATS~\cite{ge2022long} addressed quality degradation in long sequences. The paradigm has since resurged in large multimodal backbones: Emu3~\cite{wang2024emu3} generates video inside a language-model backbone by treating frames as discrete visual tokens, and CausVid~\cite{yin2025slow} and MAGI-1~\cite{ai2025magi1autoregressivevideogeneration} scale to long-form generation with autoregressive diffusion transformers. Sequential decoding nevertheless remains expensive and prone to the error accumulation that Eq.~\eqref{eq:p3_s8_3_1_ar} makes explicit.


\noindent \textbf{Diffusion Models}. Diffusion models, which now dominate T2V, lift the latent-denoising objective of §\ref{para:p3_s8_1_1_diffusion} from a single image to the clip latent $\mathbf{z}^{1:T}$. The forward and reverse processes are those of Eq.~\eqref{eq:t2i_ldm_objective}; the video-specific change is that the denoiser $\epsilon_\theta$ operates on a spatio-temporal tensor and is trained to predict the noise added to the whole clip at sampling step $s$,
\begin{equation}
\min_\theta\ \mathbb{E}_{s,\,\mathbf{z}_0^{1:T},\,\epsilon}\!\left[\big\|\epsilon-\epsilon_\theta\!\big(\mathbf{z}_s^{1:T},\,s,\,\phi(\mathbf{x})\big)\big\|_2^2\right],
\label{eq:p3_s8_3_1_diff}
\end{equation}
where $\epsilon_\theta$ models dependencies across both space and time, and $s$ is distinct from the frame index $t$. One widely used design \emph{factorizes} attention into spatial attention within each frame and temporal attention across frames, as illustrated in Figure~\ref{fig:p3_s8_3_1_paradigms}(c); other architectures use joint spacetime attention.
VDM~\cite{ho2022video} applied the idea in pixel space with a 3D U-Net; the breakthrough was moving diffusion into the compressed latent space of a video-specific autoencoder, as in VideoLDM~\cite{blattmann2023align}, LVDM~\cite{he2022latent}, and the open-source Stable Video Diffusion (SVD)~\cite{blattmann2023stable}, which cut training and inference cost enough to make high-resolution generation practical.

A second recipe adapts pretrained T2I models. Imagen Video~\cite{ho2022imagen} and Make-A-Video~\cite{singer2022make} freeze a T2I model and insert trainable temporal modules (\eg, temporal attention layers or 1D convolutions) between its spatial layers, which are then trained on video to learn motion while the image prior is inherited intact. VideoCrafter~\cite{chen2023videocrafter1,chen2024videocrafter2} and LaVie~\cite{wang2025lavie} built high-quality generators on this recipe, Tune-A-Video~\cite{wu2023tune} adapted it to one-shot tuning on a single clip, and AnimateDiff~\cite{guo2023animatediff} turned the temporal modules into plug-and-play motion modules that animate personalized T2I models without further fine-tuning.

Diffusion Transformers (DiTs) have since replaced the U-Net backbone, operating on ``spacetime patches'' (unified tokens representing blocks of video) and scaling well. Sora~\cite{sora2024} demonstrated long-duration (up to one minute), high-resolution clips and presented the approach as a step toward world simulation, and a wave of DiT models followed, including Open-Sora~\cite{zheng2024open}, LATTE~\cite{ma2024latte}, CogVideoX~\cite{yang2024cogvideox}, StepVideo~\cite{ma2025stepvideot2vtechnicalreportpractice}, Movie Gen~\cite{polyak2024movie}, Wan~\cite{wan2025wan}, and Veo~\cite{DeepMindVeo}.
Large open DiT foundation models have substantially narrowed the gap to proprietary systems. HunyuanVideo~\cite{kong2024hunyuanvideo}, at $13$B parameters the largest open T2V model at release, couples a ``dual-stream-to-single-stream'' full-attention DiT (text and video tokens processed in separate streams before fusion) with a multimodal-LLM text encoder and a 3D causal VAE, and reports human-evaluated quality competitive with leading closed systems. Mochi~1~\cite{genmo2024mochi} introduces an asymmetric DiT (AsymmDiT) that allocates most capacity to the visual stream, and LTX-Video~\cite{hacohen2024ltxvideo} uses a high-compression video VAE to run latent video diffusion in near real time. NVIDIA's Cosmos~\cite{nvidia2025cosmos} reframes the video generator as a world-foundation model for ``Physical AI'', trained so that embodied and robotics systems can use it as a learned simulator, which ties video generation to the world models of §\ref{sec:p3_s8_6}.
Across these shifts the T2V trajectory has trailed the T2I trajectory of §\ref{subsec:p3_s8_1_1} by roughly a year, each image paradigm being inflated along the time axis once it matured.

\subsubsection{Benchmarks}
\label{subsubsec:p3_s8_3_1_benchmarks}
Video adds a requirement at every level of evaluation: corpora must pair clips with captions at scale, suites must stratify by temporal phenomena (motion, object permanence, dynamic degree), and metrics must measure cross-frame coherence as well as per-frame quality.

\paragraph{Training Corpora.}
Caption--clip scale has been the principal driver of T2V quality, and high-quality paired video--text data is far scarcer than image--text data. Early work repurposed action-recognition corpora such as UCF-101~\cite{soomro2012ucf101} and Kinetics~\cite{carreira2017quo}, which supply short, class-labeled clips but no captions. WebVid-10M~\cite{bain2021frozen}, the first web-scale captioned corpus, underpinned most 2022--2023 systems despite noisy alt-text captions and watermarks; later corpora raised caption density and clip length, with Panda-70M~\cite{chen2024panda70mcaptioning70mvideos} captioned automatically by cross-modality teachers, MiraData~\cite{NEURIPS2024_57f6683e} targeting long clips with dense structured captions, and Koala-36M~\cite{wang2025koala36m} explicitly optimizing the consistency between caption and content, the property that most directly limits a generator's controllability. Table~\ref{tab:p3_s8_3_1_corpora} gives their scale.

\begin{table}[!ht]
\centering
\caption{Representative training corpora for text-to-video generation. Action-recognition sets (top) supply class labels rather than captions; web-scale captioned corpora (bottom) supply the clip--text pairs on which modern generators are fit. MiraData scale is that of its largest release.}
\label{tab:p3_s8_3_1_corpora}
\fontsize{8}{10}\selectfont
\setlength{\tabcolsep}{1.0mm}
\begin{tabular}{L{0.22\textwidth} C{0.08\textwidth} L{0.24\textwidth} L{0.38\textwidth}}
\toprule
\textbf{Corpus} & \textbf{Year} & \textbf{Scale} & \textbf{Annotation / Notes} \\
\midrule
UCF-101~\cite{soomro2012ucf101} & 2012 & 13{,}320 clips & 101 action classes; labels, no captions \\
Kinetics~\cite{carreira2017quo} & 2017 & $\sim$650K clips & Action classes; labels, no captions \\
WebVid-10M~\cite{bain2021frozen} & 2021 & 10M clip--text pairs & Web alt-text captions; noisy, watermarked \\
Panda-70M~\cite{chen2024panda70mcaptioning70mvideos} & 2024 & 70M clip--text pairs & Auto-captioned by cross-modality teachers \\
MiraData~\cite{NEURIPS2024_57f6683e} & 2024 & 330K clips & Long-duration; dense structured captions \\
Koala-36M~\cite{wang2025koala36m} & 2025 & 36M clip--text pairs & Fine-grained structured captions; condition--content consistency \\
\bottomrule
\end{tabular}
\end{table}

\paragraph{Evaluation Benchmark Suites.}
Classical T2V evaluation reported zero-shot distributional scores on the held-out splits of UCF-101 and MSR-VTT~\cite{xu2016msrvtt}, which probe frame quality and coarse text relevance but say nothing about temporal or compositional faithfulness. Recent suites are purpose-built: VBench~\cite{huang2023vbench} decomposes quality into fine-grained dimensions (subject and background consistency, motion smoothness, dynamic degree, temporal flicker, per-frame aesthetic and imaging quality), each with an automatic protocol validated against human preference; EvalCrafter~\cite{liu2024evalcrafter} aggregates objective metrics over prompts derived from real-world usage; and T2V-CompBench~\cite{sun2024t2vcompbench} targets compositional faithfulness over time (attribute binding, spatial relations, motion binding, interactions), the temporal analogue of the T2I-CompBench probe of §\ref{subsec:p3_s8_1_1}. Table~\ref{tab:p3_s8_3_1_eval} lists them.

\begin{table}[!ht]
\centering
\caption{Representative evaluation resources for text-to-video generation. Distributional resources (top) reuse captioned-clip splits for zero-shot FVD/IS/CLIPSIM; multi-dimensional suites (bottom) score temporal coherence, dynamics, and compositional faithfulness directly.}
\label{tab:p3_s8_3_1_eval}
\fontsize{8}{10}\selectfont
\setlength{\tabcolsep}{1.0mm}
\begin{tabular}{L{0.22\textwidth} C{0.10\textwidth} L{0.36\textwidth} L{0.24\textwidth}}
\toprule
\textbf{Suite} & \textbf{Year} & \textbf{Evaluation Axis} & \textbf{Scoring} \\
\midrule
UCF-101 / MSR-VTT~\cite{soomro2012ucf101,xu2016msrvtt} & 2012--16 & Frame quality $+$ coarse text relevance (MSR-VTT: 10K captioned clips) & Zero-shot FVD / IS / CLIPSIM \\
VBench~\cite{huang2023vbench} & 2024 & 16 dimensions (consistency, motion, dynamics, flicker, quality) & Per-dimension automatic, human-validated \\
EvalCrafter~\cite{liu2024evalcrafter} & 2024 & 17 objective metrics over 700 structured prompts & Aggregated automatic, fitted to human opinion \\
T2V-CompBench~\cite{sun2024t2vcompbench} & 2024 & Compositional: attribute/motion binding, relations, interactions & MLLM- / detector-based \\
\bottomrule
\end{tabular}
\end{table}

\paragraph{Evaluation Metrics.}
The metric that defines video evaluation is the Fr\'echet Video Distance (FVD)~\cite{unterthiner2018fvd}, the spatio-temporal analogue of FID: it compares distributions of features extracted from generated and real clips by a pretrained \emph{video} backbone (I3D), and is therefore sensitive to motion realism in a way that per-frame FID~\cite{heusel2017gans} and Inception Score are not. Text--clip alignment is measured by CLIPSIM, the per-frame CLIPScore~\cite{hessel2021clipscore} averaged over the clip, which captures topical relevance but is blind to whether the \emph{described} motion actually occurs. Because none of these scalars captures coherence or dynamics, papers report several VBench dimensions alongside them, which exposes a trade-off a single number hides: \emph{dynamic degree} (how much genuine motion a clip contains) against \emph{temporal consistency} (how stable it stays), since a near-static clip trivially maximizes consistency while failing the prompt. Headline claims are corroborated with human preference.

\subsubsection{Challenges and Trends}
\label{subsubsec:p3_s8_3_1_challenges}
Three problems belong to the unconditional base case. \textbf{Long-horizon coherence from text alone} is the first: with no scaffold, object permanence and motion plausibility must emerge implicitly, and all three paradigms degrade as clips lengthen (GANs flicker, autoregressive roll-outs accumulate the error of Eq.~\eqref{eq:p3_s8_3_1_ar}, diffusion models lose consistency beyond their trained window), so coherence beyond a few seconds remains the binding constraint. \textbf{Motion realism without a control signal} is the second: prompts under-specify dynamics, so the generator must sample from a highly multi-modal motion posterior, and the easy optimum is near-static or physically implausible motion that still scores well on per-frame metrics. The \textbf{video--text data bottleneck} is the third: paired data is scarce relative to images, web captions are short and weakly aligned, and this caps both controllability and the reliability of the automatic captioning used to scale corpora, which is what the dense-caption efforts (Panda-70M, MiraData, Koala-36M) address. The explicit control interfaces of the following subsections relieve the first two by moving structure and dynamics from the prompt into the input.


\subsection{Spatial Controlled Video Generation}
\label{subsec:p3_s8_3_2}

While textual descriptions provide high-level semantic guidance, they often lack the granularity required for precise control over object placement, scene geometry, or character appearance. Spatial-controlled video generation addresses this limitation by augmenting the generation process with explicit spatial conditioning signals, affording users more direct and deterministic influence over the synthesized visual output.
The geometry, articulation, and layout that the text-to-video base case of §\ref{subsec:p3_s8_3_1} had to \emph{hallucinate} from the prompt alone are here \emph{externalized} into a per-frame scaffold: the spatial structure becomes an input the model must honor rather than an emergent property it must invent, which is precisely what makes the output controllable and verifiable.

\subsubsection{Task Definition}
\label{subsubsec:p3_s8_3_2_task}
Spatial control instantiates Eq.~\eqref{eq:p3_s8_3_objective} with a control encoder $\kappa_{\text{s}}$ acting on a spatial condition $\mathbf{c}_{\text{s}}^{1:T}$ (a static or sparse reference image, a pose sequence, a depth-map sequence, facial landmarks, a sketch sequence, or a bounding-box track),
\begin{equation}
\hat{\mathbf{z}}^{1:T}\sim p_\theta\!\big(\mathbf{z}^{1:T}\,\big|\,\phi(\mathbf{x}),\,\kappa_{\text{s}}(\mathbf{c}_{\text{s}}^{1:T})\big),
\label{eq:p3_s8_3_2_objective}
\end{equation}
so that the scaffold pins down the per-frame spatial layout while the prompt $\mathbf{x}$ and the temporal modules supply appearance and motion. The conditioning channels surveyed below differ only in what $\mathbf{c}_{\text{s}}$ encodes (appearance for an image, articulation for pose or landmarks, 3D structure for depth, free-form layout for a sketch, or object location for a bounding box), and recover the text-to-video base case of Eq.~\eqref{eq:p3_s8_3_1_objective} exactly when $\mathbf{c}_{\text{s}}=\varnothing$.

\subsubsection{Methodology}
\label{subsubsec:p3_s8_3_2_methodology}
The six channels carry different spatial information, while the injection mechanism depends on the model. A common design is a ControlNet-style adapter~\cite{zhang2023adding}: a trainable control branch processes the encoded condition and injects a residual into intermediate features of a video backbone, schematically at layer $\ell$
\begin{equation}
h_\ell' = h_\ell + \Delta h_\ell,
\qquad
\Delta h_\ell = \mathcal{A}_{\psi,\ell}\!\big(\kappa_{\text{s}}(\mathbf{c}_{\text{s}}^{1:T})\big),
\label{eq:p3_s8_3_2_inject}
\end{equation}
where $h_\ell$ denotes backbone features and $\mathcal{A}_{\psi,\ell}$ a trainable control branch (dependence on the noisy latent, diffusion step, and other conditioning is omitted, and the backbone may be frozen while the adapter trains); latent concatenation and attention-based conditioning are the other common implementations.
Figure~\ref{fig:p3_s8_3_2_mechanisms} shows the six channels alongside this adapter pathway: a reference image acts as a static appearance anchor, the other five are usually frame-indexed sequences, and compatible channels can be combined.

\begin{figure}[!ht]
\centering
\includegraphics[width=\linewidth]{figures/part3/section8/spatial_control_channels_for_video_generation.pdf}
\caption{Six spatial-control channels for video generation, (a)--(f), which can be selected or combined, and an example injection route (right) in which a trainable adapter adds residual features to a frozen video backbone (Eq.~\eqref{eq:p3_s8_3_2_inject}). Frames and control maps are illustrative.}
\label{fig:p3_s8_3_2_mechanisms}
\end{figure}

\noindent\textbf{Image Control.} A reference image (Figure~\ref{fig:p3_s8_3_2_mechanisms}a) specifies the style, content, or initial state of a video, which is the basis of animating a still and of keeping a subject's appearance consistent. I2VGen-XL~\cite{zhang2023i2vgen}, Stable Video Diffusion (SVD)~\cite{blattmann2023stable}, VideoCrafter1~\cite{chen2023videocrafter1}, and Lumiere~\cite{bar2024lumiere} use the reference image as an initial condition, often by concatenating its latent with the initial noise, and ConsistI2V~\cite{ren2024consisti2v} refines this for consistency. The model then generates a plausible motion sequence that evolves from the image, typically under a text prompt, as in DynamiCrafter~\cite{xing2024dynamicrafter}.

\noindent\textbf{Pose Control.} Pose-guided generation (Figure~\ref{fig:p3_s8_3_2_mechanisms}b) animates articulated subjects, predominantly humans, along a predefined sequence of skeletal poses, the video counterpart of pose-guided image generation~\cite{ju2023humansd,zhang2023adding,Ju_2023_CVPR} and the core of avatar animation and character generation. The pose keypoints (\eg, from MMPose~\cite{mmpose2020}) are encoded and injected at each denoising step, frequently through a ControlNet-style adapter~\cite{zhang2023adding}. MagicAnimate~\cite{xu2024magicanimate}, Animate Anyone~\cite{hu2024animate}, Follow-Your-Pose~\cite{ma2024followyourpose}, and DreamPose~\cite{karras2023dreampose} pair a reference image for appearance with a pose sequence for motion, DisPose~\cite{li2024dispose} disentangles the pose guidance, and Champ~\cite{zhu2024champ} and MagicPose~\cite{chang2023magicpose} target high-fidelity motion, so that the output keeps the subject's identity while following the target poses.

\noindent\textbf{Depth Control.} A sequence of depth maps (Figure~\ref{fig:p3_s8_3_2_mechanisms}c), estimated from an existing video or rendered from a 3D model, specifies dense 3D structure and so enforces consistent object layout and realistic occlusion. ControlVideo~\cite{zhao2025controlvideo}, Make-Your-Video~\cite{xing2024make}, Control-A-Video~\cite{chen2023control}, and ControlNeXt~\cite{peng2024controlnext} adapt the ControlNet framework to video, processing each frame's depth map with a trainable adapter that guides the denoiser; DreamDance~\cite{pang2024dreamdance} and SparseCtrl~\cite{guo2024sparsectrl} also use depth for dynamic scene generation.

\noindent\textbf{Landmark Control.} Facial landmarks (Figure~\ref{fig:p3_s8_3_2_mechanisms}d) give fine-grained control over expression and are the standard interface for talking-head animation. EchoMimic~\cite{chen2024echomimic}, HunyuanPortrait~\cite{xu2025hunyuanportrait}, and LivePortrait~\cite{guo2024liveportrait} animate a static portrait from landmarks, often themselves driven by audio, so that lip movements, blinks, and subtle expressions stay synchronized; TASTE-Rob~\cite{zhao2025taste}, Takin-ADA~\cite{lin2024takin}, and Follow-Your-Emoji~\cite{ma2024follow} extend landmarks to more expressive control.

\noindent\textbf{Sketch Control.} Sketches (Figure~\ref{fig:p3_s8_3_2_mechanisms}e) let users define compositional layout and motion with simple line drawings, an intuitive interface for animation. ToonCrafter~\cite{xing2024tooncrafter}, VidSketch~\cite{jiang2025vidsketch}, AniDoc~\cite{meng2025anidoc}, and SketchBetween~\cite{loftsdottir2022sketchbetween} take a sequence of sketches to guide interpolation or generation, encoding the line art and injecting it as a condition; LayerAnimate~\cite{yang2025layeranimate} animates hand-drawn characters and scenes layer by layer, and MagicStick~\cite{ma2025magicstick} uses sketch-like structural cues for editing.

\noindent\textbf{Bounding Box (BBox) Control.} Bounding-box tracks (Figure~\ref{fig:p3_s8_3_2_mechanisms}f) specify where objects are and how they move over time, and dominate autonomous-driving simulation: MagicDrive~\cite{gao2023magicdrive}, DriveDreamer~\cite{wang2024drivedreamer}, DriveScape~\cite{wu2024drivescape}, and DreamForge~\cite{mei2024dreamforge} use 3D boxes to control the position and motion of vehicles and pedestrians. The box information enters either through cross-attention, as in DiVE~\cite{jiang2024dive}, or as a spatial feature map, and Boximator~\cite{wang2024boximator} and Ctrl-V~\cite{luo2025ctrl} generate and control object motion from box sequences.

Nearly every structural channel (pose, depth, sketch, box) is paired with an appearance reference image, since structure alone cannot fix identity or texture, and the channels cluster by application: image and landmark control in portrait and talking-head synthesis, pose and sketch control in character animation, depth and box control in geometry-faithful and driving-simulation settings.

\subsubsection{Benchmarks}
\label{subsubsec:p3_s8_3_2_benchmarks}
Unlike the text-to-video base case of §\ref{subsec:p3_s8_3_1}, spatial control has no single canonical corpus or suite: each channel is trained and evaluated on data specialized to its application. Models are typically built on the general web-scale corpora of §\ref{subsec:p3_s8_3_1} (WebVid-10M~\cite{bain2021frozen} and successors) for appearance priors, then specialized on task data: human-animation channels on dance- and fashion-video collections paired with extracted pose keypoints (\eg, via OpenPose~\cite{cao2017realtime} or MMPose~\cite{mmpose2020}), talking-head landmark control on portrait-video sets, and bounding-box driving simulation on annotated driving datasets such as nuScenes~\cite{caesar2020nuscenes}, whose 3D box and trajectory labels supply the control signal directly.

\paragraph{Control-fidelity Metrics.}
Spatial control is judged on two axes simultaneously: the generic video-quality metrics of §\ref{subsec:p3_s8_3_1} (FVD~\cite{unterthiner2018fvd}, CLIPSIM, and VBench~\cite{huang2023vbench} dimensions) and channel-specific \emph{adherence} metrics that measure how faithfully the output respects $\mathbf{c}_{\text{s}}$. The latter are necessarily per-channel: pose and landmark control report keypoint accuracy (re-detected-keypoint distance to the target skeleton or face mesh); depth control reports depth consistency between the re-estimated and the conditioning maps; image control reports identity and appearance preservation (face-/subject-similarity to the reference); and bounding-box control reports detection mAP and trajectory IoU of the generated objects against the conditioning boxes. Neither axis suffices alone, so methods report a quality-plus-adherence tuple: a model can maximize adherence by ignoring the prompt (rigid, texture-collapsed output) or maximize quality by ignoring the control. The same holds for every controlled axis that follows.

\subsubsection{Challenges and Trends}
\label{subsubsec:p3_s8_3_2_challenges}
Single-channel spatial control has three problems of its own; those that arise once several conditions are fused are discussed in §\ref{subsec:p3_s8_3_6}. Because the structural scaffold and the appearance reference pass through the same denoiser, strong control leaks structure into texture: pose and depth conditioning induce identity drift or bake the scaffold's silhouette into the surface, and disentangling what moves from what it looks like (as DisPose~\cite{li2024dispose} attempts for pose) remains open. Dense scaffolds (image, depth) over-constrain and suppress plausible motion, whereas sparse ones (boxes, coarse sketches) leave appearance and fine geometry ambiguous, so no channel is uniformly preferable and the right granularity depends on the task. And adapters trained on, say, frontal human dance transfer poorly to unusual viewpoints or non-human subjects, which is the data cost of per-channel specialization and one motivation for the multi-condition frameworks of §\ref{subsec:p3_s8_3_4}.


\subsection{Temporal Controlled Video Generation}
\label{subsec:p3_s8_3_3}

Temporal control regulates how motion and timing evolve across frames, the lever behind cinematic camera work, narrative pacing, and believable physical interaction. It is the complement of the spatial scaffold of §\ref{subsec:p3_s8_3_2}: where spatial control fixes \emph{where things are} in each frame, temporal control fixes \emph{how they move between} frames, externalizing the dynamics (object motion, camera path, timing) that the text-to-video base case of §\ref{subsec:p3_s8_3_1} otherwise had to sample blindly from a highly multi-modal motion posterior.

\subsubsection{Task Definition}
\label{subsubsec:p3_s8_3_3_task}
Temporal control instantiates Eq.~\eqref{eq:p3_s8_3_objective} with an encoder $\kappa_{\text{t}}$ of a motion or timing signal $\mathbf{m}^{1:T}$ (an optical-flow field, an object trajectory, a camera-pose path, or a learned motion embedding),
\begin{equation}
\hat{\mathbf{z}}^{1:T}\sim p_\theta\!\big(\mathbf{z}^{1:T}\,\big|\,\phi(\mathbf{x}),\,\kappa_{\text{t}}(\mathbf{m}^{1:T})\big).
\label{eq:p3_s8_3_3_objective}
\end{equation}
The contrast with spatial control (Eq.~\eqref{eq:p3_s8_3_2_objective}) is that $\mathbf{m}^{1:T}$ constrains the \emph{inter-frame dynamics} (how content moves between frames) rather than the per-frame appearance, so $\kappa_{\text{t}}$ typically steers the temporal-attention or latent-warping pathway of the denoiser. The channels below differ in how the motion is specified: densely (optical flow), as object paths (trajectory), as virtual-camera parameters (camera), or as a transferable behavioral style (motion).

\subsubsection{Methodology}
\label{subsubsec:p3_s8_3_3_methodology}
The four channels specify inter-frame dynamics through different signals, and no single injection architecture serves them all. One implementation uses a trainable temporal-control branch $\mathcal{T}_\psi$ to modulate features of a pretrained video backbone,
\begin{equation}
h_\ell' = h_\ell + \Delta h_\ell,
\qquad
\Delta h_\ell=\mathcal{T}_{\psi,\ell}\!\big(\kappa_{\text{t}}(\mathbf{m}^{1:T})\big),
\label{eq:p3_s8_3_3_inject}
\end{equation}
with dependence on noisy latents, diffusion step, and text omitted.
More broadly, two routes recur: explicit transformations that use the motion to warp noise, latents, or attention features, and encoded conditioning that steers generation through adapters, attention, or learned motion features. They are not tied to particular channels (flow and trajectories can drive learned conditioning, and camera control can use geometry-aware warping), and Figure~\ref{fig:p3_s8_3_3_mechanisms} keeps the four input channels separate from the two routes.

\begin{figure}[!ht]
\centering
\includegraphics[width=\linewidth]{figures/part3/section8/temporal_control_channels_for_video_generation.pdf}
\caption{Four temporal-control channels for video generation, (a)--(d), and two alternative ways of applying them (right): warping noise or latents, and encoding the signal as conditioning features. Frames are illustrative.}
\label{fig:p3_s8_3_3_mechanisms}
\end{figure}

\noindent\textbf{Flow Control.} The most explicit motion signal (Figure~\ref{fig:p3_s8_3_3_mechanisms}a) is an optical-flow field, which records the per-pixel displacement between consecutive frames and so specifies dynamics at the finest granularity; conditioning on a sequence of such fields yields specific, coherent motion at the cost of authoring or estimating a dense signal. Motion-I2V~\cite{shi2024motion} integrates a motion-field predictor into its diffusion framework so that generation follows the supplied flow, I2VControl~\cite{feng2024i2vcontrol} disentangles flow for unified control, and MOFA-Video~\cite{niu2024mofa} animates a still image from user-painted motion fields. Go-with-the-Flow~\cite{burgert2025gowiththeflow} dispenses with the adapter altogether, warping the diffusion \emph{noise prior} along the target flow so that the motion constraint lives in the sampling process rather than in an added branch.

\noindent\textbf{Trajectory Control.} Sparsifying the dense field into a few object paths (Figure~\ref{fig:p3_s8_3_3_mechanisms}b) trades pixel-level precision for an interface a user can actually author: a sequence of 2D or 3D coordinates, or a path drawn interactively, directs character movement or object kinematics, entering the denoiser either by modifying the latents or by steering cross-attention maps. MotionCanvas~\cite{xing2025motioncanvas}, MagicMotion~\cite{li2025magicmotion}, and Tora~\cite{zhang2025tora} consume a predefined curve, FreeTraj~\cite{qiu2024freetraj} relaxes this to free-form trajectories, Peekaboo~\cite{jain2024peekaboo} also schedules when objects appear and disappear along the path, and DragNUWA~\cite{yin2023dragnuwa} and DragAnything~\cite{wu2024draganything} let the user drag regions through the frame interactively. MotionCtrl~\cite{wang2024motionctrl} handles object trajectory and camera motion in one controller, treating an object path and a camera path as two facets of the same inter-frame displacement.

\noindent\textbf{Camera Control.} Moving that displacement from objects to the observer (Figure~\ref{fig:p3_s8_3_3_mechanisms}c), camera-guided generation manipulates the virtual camera's extrinsics (translation, rotation) and intrinsics (\eg, focal length) to produce pans, tilts, zooms, and dolly moves. CameraCtrl~\cite{he2024cameractrl} and CameraCtrl~II~\cite{he2025cameractrl} encode the camera trajectory as a conditioning feature; Direct-a-Video~\cite{yang2024direct}, ViewCrafter~\cite{yu2024viewcrafter}, and Cavia~\cite{xu2024cavia} couple it with novel-view synthesis; and 3DTrajMaster~\cite{fu20243dtrajmaster}, MotionFlow~\cite{meral2024motionflow}, and CineMaster~\cite{wang2025cinemaster} control camera and objects jointly in 3D. Two directions now lead. ReCamMaster~\cite{bai2025recammaster} and TrajectoryCrafter~\cite{yu2025trajectorycrafter} re-render an \emph{existing} monocular clip under a new camera path rather than generating from scratch, and Diffusion-as-Shader~\cite{gu2025diffusionshader3dawarevideo} and GS-DiT~\cite{bian2025gsditadvancingvideogeneration} route the camera signal through a 3D proxy (a rendered shader pass or a pseudo-4D Gaussian field), which sharpens control by making it geometry-aware.

\noindent\textbf{Motion Control.} At the abstract end (Figure~\ref{fig:p3_s8_3_3_mechanisms}d), the signal is no longer a geometric displacement but a \emph{behavioral concept} (walking, dancing, jumping) transferred from a reference video as a pose sequence or a learned motion embedding. Unlike the pose control of §\ref{subsec:p3_s8_3_2}, it captures the dynamic style and cadence of a behavior rather than a frame-exact skeleton, so the central technical problem is disentanglement: extracting the motion while discarding the reference's appearance. MotionDirector~\cite{zhao2024motiondirector}, MotionClone~\cite{ling2024motionclone}, and VMC~\cite{jeong2024vmc} learn this motion--appearance separation; Motion Prompting~\cite{geng2025motion}, MotionAgent~\cite{liao2025motionagent}, and MotionMatcher~\cite{wu2025motionmatcher} refine its granularity through trajectory- or feature-level motion descriptors; and DreamMotion~\cite{jeong2024dreammotion} and Customize-A-Video~\cite{ren2024customize} apply a reference clip's dynamics to a new subject or scene fixed by the text prompt.

The channels trade precision against ease of specification: dense flow gives pixel-accurate control but is laborious to author, a motion embedding is effortless to supply but steers the result only loosely, and the most usable systems therefore pair an abstract signal (a drawn path, a reference clip) with a learned predictor that densifies it into the flow- or latent-level displacement the backbone actually consumes.

\subsubsection{Benchmarks}
\label{subsubsec:p3_s8_3_3_benchmarks}
As with spatial control (§\ref{subsec:p3_s8_3_2}), temporal control has no single canonical suite: each channel is evaluated against data specialized to its signal.
Models build on the general corpora of §\ref{subsec:p3_s8_3_1} for appearance and motion priors, then evaluate on task-specific resources: camera control on the posed-video collection RealEstate10K~\cite{zhou2018stereo}, whose per-frame camera poses supply both the conditioning signal and the evaluation ground truth, increasingly supplemented by the larger and more scene-diverse DL3DV-10K~\cite{ling2024dl3dv} to test generalization beyond real-estate walkthroughs; flow control against the MPI-Sintel optical-flow benchmark~\cite{sintel} or flow re-estimated from generated clips with RAFT~\cite{teed2020raft}; and motion- and trajectory-control quality through the motion-smoothness and dynamic-degree dimensions of VBench~\cite{huang2023vbench}.

\paragraph{Control-fidelity Metrics.}
Alongside the generic video-quality metrics of §\ref{subsec:p3_s8_3_1}, each channel measures adherence by re-extracting the control signal from the generated clip and comparing it with the conditioning input; the estimator and the error differ per channel.

\noindent\textit{Flow adherence.} The optical flow of the generated clip is re-estimated with RAFT~\cite{teed2020raft} and compared to the conditioning flow by the \emph{end-point error} (EPE), the per-pixel $L_2$ displacement difference averaged over pixels and frames. A complementary \emph{flow-warping error} warps each generated frame onto the next with the \emph{target} flow and reports the photometric residual, which tests whether the rendered motion realizes the commanded displacement rather than merely correlating with it.

\noindent\textit{Trajectory adherence.} The controlled object is re-localized in the output with a dense point tracker such as CoTracker~\cite{karaev2024cotracker} or a detector, and its tracked center is compared to the target path by the mean $L_2$ distance, normalized by frame size and averaged over time and instances; this is the object motion-control error reported by Tora~\cite{zhang2025tora} and DragNUWA~\cite{yin2023dragnuwa}, and it penalizes both spatial deviation and mistimed arrival.

\noindent\textit{Camera adherence.} The camera trajectory of the generated clip is recovered by structure-from-motion (COLMAP~\cite{schonberger2016structure}) and aligned to the target up to a similarity transform, since pose is recoverable only up to scale. Following CameraCtrl~\cite{he2024cameractrl}, a \emph{rotation error} (the geodesic angle between estimated and target rotations, summed over frames) and a \emph{translation error} (the $L_2$ distance between scale-normalized camera centers) are reported; failure to reconstruct a stable trajectory at all is itself diagnostic of an incoherent output.

\noindent\textit{Motion-transfer fidelity.} With no ground-truth video, fidelity is measured from two sides: a \emph{motion-similarity} term compares the dynamics of output and reference (the correlation of tracked point trajectories or temporal feature statistics), and an \emph{appearance-divergence} term (CLIP similarity to the new prompt together with low feature similarity to the reference's appearance) confirms that the subject was re-synthesized rather than copied.

As for spatial control, these adherence errors are reported together with generic quality, since a near-static clip trivially minimizes EPE and pose error while ignoring the prompt.

\subsubsection{Challenges and Trends}
\label{subsubsec:p3_s8_3_3_challenges}
Three problems are specific to temporal control; fusing it with other conditions is discussed in §\ref{subsec:p3_s8_3_6}. Because the motion signal passes through the same backbone that renders appearance, a transferred motion often drags the source clip's identity or texture with it, the leakage that MotionDirector~\cite{zhao2024motiondirector} and VMC~\cite{jeong2024vmc} try to disentangle and that remains unsolved for large appearance gaps between reference and target. Dense flow over-constrains and suppresses generative freedom, whereas a single trajectory or a short reference clip under-specifies the global dynamics, which is why methods increasingly interpose a learned predictor that densifies an abstract signal into a flow- or latent-level field. And a commanded camera path or object trajectory can demand views or dynamics inconsistent with the scene geometry, producing parallax errors and object distortion; 3D-aware camera control (ViewCrafter~\cite{yu2024viewcrafter}, CineMaster~\cite{wang2025cinemaster}) routes the control through an explicit geometric proxy for exactly this reason, which links temporal control back to the volumetric representations of §\ref{sec:p3_s8_2}.





\subsection{Other Controlled Video Generation}
\label{subsec:p3_s8_3_4}

Beyond the spatial scaffolds of §\ref{subsec:p3_s8_3_2} and the motion signals of §\ref{subsec:p3_s8_3_3}, a video can be steered by further, application-specific conditions: an artistic style reference, a handful of draggable points, a bird's-eye-view road layout, a target text string to render, or a synchronizing audio track.
Each externalizes an attribute that text specifies only loosely, extending controllable generation to creative, interactive, and domain-specific settings; together they raise the question the single-axis interfaces leave open, how to honor \emph{several} conditions at once, which the universal frameworks at the end of this subsection try to answer.

\subsubsection{Task Definition}
\label{subsubsec:p3_s8_3_4_task}
Taken one at a time, each ``other'' modality instantiates the generic conditional sampler of Eq.~\eqref{eq:p3_s8_3_objective} with an encoder $\kappa_{\text{o}}$ of a heterogeneous control signal $\mathbf{c}$ (a style reference, a set of draggable points, a BEV layout, a target text string, or an audio track). The distinctive task of this subsection, however, is to honor \emph{several} such conditions simultaneously: universal and multi-task frameworks fuse $K$ encoded signals into a single sequence,
\begin{equation}
\hat{\mathbf{z}}^{1:T}\sim p_\theta\!\big(\mathbf{z}^{1:T}\,\big|\,\phi(\mathbf{x}),\,\kappa_1(\mathbf{c}_1),\dots,\kappa_K(\mathbf{c}_K)\big),
\label{eq:p3_s8_3_4_objective}
\end{equation}
so that the $K$ encoded conditions are processed jointly rather than one axis at a time, with $K=1$ recovering the single-modality case and $\kappa_k=\varnothing$ for all $k$ recovering the text-to-video base case of §\ref{subsec:p3_s8_3_1}.

\subsubsection{Methodology}
\label{subsubsec:p3_s8_3_4_methodology}
The methods fall into single-modality channels (style, points, BEV, text rendering, audio) and multi-task frameworks that combine compatible conditions. Writing $u_k=\kappa_k(\mathbf c_k)$, the common interface is
\begin{equation}
\hat{\mathbf z}^{1:T}=G_\theta\!\big(\mathbf n,\phi(\mathbf x);\mathcal F_\psi(u_1,\dots,u_K)\big),
\label{eq:p3_s8_3_4_inject}
\end{equation}
where $\mathcal F_\psi$ is a model-specific conditioning interface that may act at several layers rather than produce one fixed-size embedding. Three implementations are representative and can be combined: summed adapter residuals $\sum_k\Delta h_k$, concatenation of the encoded channels with the model input or features, and unified attention over condition tokens. Figure~\ref{fig:p3_s8_3_4_mechanisms} shows the five input modalities and these three fusion designs; $K=1$ denotes one selected modality and $K>1$ a supported combination.

\begin{figure}[!ht]
\centering
\includegraphics[width=\linewidth]{figures/part3/section8/other_control_modalities_for_video_generation.pdf}
\caption{Further control modalities for video generation, (a)--(e), each with its own encoder $E_k$ implementing $\kappa_k$, and (f) alternative ways of fusing $K\geq 1$ encoded conditions. The watercolor driving sequence is illustrative.}
\label{fig:p3_s8_3_4_mechanisms}
\end{figure}

\noindent\textbf{Style Control.} Style-guided generation transfers the artistic style of a reference image or video to a target video. Per-frame style transfer flickers; diffusion-based approaches such as Rerender A Video~\cite{yang2023rerender}, StyleCrafter~\cite{liu2023stylecrafter}, FRESCO~\cite{yang2024fresco}, and StyleMaster~\cite{ye2025stylemaster} instead integrate the style into generation, often by guiding it with the CLIP embedding of the style reference so that the style holds across frames, and UniVST~\cite{song2024univst} generalizes this into a universal framework. FrameGuidance~\cite{jang2025frameguidancetrainingfreeguidance} and OmniVDiff~\cite{xi2025omnivdiffomnicontrollablevideo} add training-free frame-level guidance and omni-controllable diffusion, respectively.

\noindent\textbf{Point Control.} Point control fixes \emph{motion} through a sparse, interactive interface: the user drags a few points on an object from an initial frame to a new position in a later frame, and the model infers a plausible motion for the whole object. Drag-A-Video~\cite{teng2023drag} introduced the approach by optimizing latent codes to satisfy the point constraints, Track4Gen~\cite{jeong2025track4genteachingvideodiffusion} focuses on tracking the user-specified points robustly, and the 3D-aware camera controllers of §\ref{subsec:p3_s8_3_3}~\cite{gu2025diffusionshader3dawarevideo,bian2025gsditadvancingvideogeneration} also accept sparse point input for non-rigid manipulation.

\noindent\textbf{Bird's-Eye View (BEV) Control.} BEV-guided generation serves autonomous-driving simulation: a sequence of top-down semantic maps dictates the road layout, the positions of other vehicles, and the movements of pedestrians. DriveDreamer-2~\cite{zhao2024drivedreamer2llmenhancedworldmodels}, MagicDrive-V2~\cite{gao2025magicdrivev2highresolutionlongvideo}, DiVE~\cite{jiang2024dive}, DreamForge~\cite{mei2024dreamforge}, and Panacea~\cite{wen2023panaceapanoramiccontrollablevideo} take these maps as input and generate first-person driving video consistent with the specified layout and dynamics.

\noindent\textbf{Text Rendering.} Many generators cannot render coherent, legible text inside a scene, so a dedicated channel treats the target string itself as the control signal. Text-Animator~\cite{liu2024textanimatorcontrollablevisualtext} adds text-embedding injection and refinement modules that improve the structure and temporal stability of rendered text, and Wan~\cite{wan2025wan} reports being the first video foundation model able to render legible Chinese and English text.

\noindent\textbf{Audio Control.} Audio-controlled generation synthesizes video synchronized with an input audio signal, in two sub-domains.
\emph{Voice-guided generation} produces talking heads: EMO~\cite{tian2024emoemoteportraitalive}, VASA-1~\cite{xu2024vasa1lifelikeaudiodriventalking}, DreamTalk~\cite{ma2023dreamtalk}, and AniPortrait~\cite{wei2024aniportrait} take a reference image and a speech clip, extract audio features with a pretrained encoder, and drive facial motion with accurate lip-sync and expression; Say Anything~\cite{ma2025sayanythingaudiodrivenlipsynchronization}, MuseTalk~\cite{zhang2025musetalkrealtimehighfidelityvideo}, and CAFE-TALK~\cite{chen2025cafetalkgenerating3dtalking} focus on high-fidelity lip synchronization, and V-Express~\cite{wang2024vexpressconditionaldropoutprogressive} and MegActor~\cite{yang2024megactorharnesspowerraw} balance audio against other control signals.
\emph{Sound-guided generation} follows more general audio such as music or ambient sound, using its rhythm, pitch, and emotional tone to drive visual elements, motion, and scene transitions: TA2V~\cite{10420468}, MotionCraft~\cite{aira2024motioncraftphysicsbasedzeroshotvideo}, AV-Link~\cite{hajiali2025avlinktemporallyaligneddiffusionfeatures}, and ASVA~\cite{linz2024asva} generate music visualizations or scenes that follow the energy and structure of a track, and MagicInfinite~\cite{yi2025magicinfinitegeneratinginfinitetalking} and Sound2Sight~\cite{cherian2020sound2sightgeneratingvisualdynamics} explore more open-ended audio-driven generation.

\noindent\textbf{Universal and Multi-Task Control.} The single-modality channels each add one encoder; universal frameworks realize the fusion of Eq.~\eqref{eq:p3_s8_3_4_objective} in a single model, along two lines. Autoregressive models built on multimodal large language models generate video from discrete visual tokens, as in VideoPoet~\cite{kondratyuk2024videopoetlargelanguagemodel} and Emu3~\cite{wang2024emu3}. Diffusion models are being adapted to many conditions at once: VideoComposer~\cite{wang2023videocomposercompositionalvideosynthesis} integrates depth, pose, motion, and other conditions by channel concatenation, OmniHuman~\cite{lin2025omnihuman1rethinkingscalinguponestage} generates human video under multiple controls by mixing cross-attention and concatenation, and FullDiT~\cite{ju2025fullditmultitaskvideogenerative} processes all conditions with unified attention, the third realization of $\mathcal{F}_\psi$.

\subsubsection{Benchmarks}
\label{subsubsec:p3_s8_3_4_benchmarks}
Being heterogeneous, the ``other'' modalities share no single suite: each is evaluated against data and metrics specialized to its signal, with the general corpora of §\ref{subsec:p3_s8_3_1} supplying the appearance and motion priors. BEV control is evaluated on driving collections such as nuScenes~\cite{caesar2020nuscenes}, whose paired BEV layouts and front-camera footage furnish both the conditioning signal and the ground truth. Voice-guided talking-head generation is evaluated on portrait-video collections (HDTF~\cite{zhang2021hdtf} and VFHQ~\cite{xie2022vfhq} are standard) and sound-guided generation on audio--video corpora with aligned music or ambient tracks. Style and point control, having no canonical benchmark, are assessed on the general corpora plus qualitative and user studies.

\paragraph{Control-fidelity Metrics.}
Beyond the generic video-quality metrics of §\ref{subsec:p3_s8_3_1} (FVD~\cite{unterthiner2018fvd}, CLIPSIM, and VBench~\cite{huang2023vbench} dimensions), each modality adds an \emph{adherence} metric tailored to its signal, reported together with generic quality as in §\ref{subsec:p3_s8_3_2}.
\emph{Style adherence} combines a CLIP style-similarity to the reference with a temporal-consistency term that penalizes flicker.
\emph{Point adherence} re-localizes the dragged keypoints in the output and reports their $L_2$ deviation from the target path, as in the trajectory metric of §\ref{subsec:p3_s8_3_3}.
\emph{BEV adherence} renders or detects objects in the generated clip and reports detection-mAP and layout-IoU against the input map.
\emph{Text adherence} runs OCR on the rendered frames and reports recognition accuracy and legibility.
\emph{Audio adherence} measures lip-sync quality with SyncNet-style confidence and landmark error for voice, and audio--visual synchrony for sound.

\subsubsection{Challenges and Trends}
\label{subsubsec:p3_s8_3_4_challenges}
The central difficulty here, fusing conditions of incommensurate dimensionality and strength without one washing out or dominating the others, is shared with the spatial and temporal axes and discussed in §\ref{subsec:p3_s8_3_6}. Two problems are specific to these modalities. Text rendering must keep glyphs legible and stable at once, and per-frame legibility and temporal coherence pull against each other as the camera or the text moves. BEV control covers a narrow domain: models trained on a handful of driving datasets inherit their sensor rigs, cities, and traffic statistics, and generalize poorly beyond them.


\subsection{Video Editing}
\label{subsec:p3_s8_3_5}

Video editing takes an \emph{existing} clip and modifies it under user instruction, rather than generating one from scratch (§\ref{subsec:p3_s8_3_1}) or steering that generation with controls (§\ref{subsec:p3_s8_3_2}--§\ref{subsec:p3_s8_3_4}). It is the temporal lift of the image editing of §\ref{subsec:p3_s8_1_3}, and its defining demand is selective preservation: the operator must change only the targeted content while holding the source's geometry, layout, and motion fixed across every frame, so the central technical problem is to balance editability against spatio-temporal consistency without artifacts or unintended drift.

\subsubsection{Task Definition}
\label{subsubsec:p3_s8_3_5_task}
Given a source clip encoded as $\mathbf{z}^{\text{in},1:T}=E(\mathbf{v}^{\text{in}})$ and an edit instruction $\mathbf{c}$, an editing operator
\begin{equation}
\hat{\mathbf{z}}^{1:T}=E_\theta^{\text{vid}}\!\big(\mathbf{z}^{\text{in},1:T},\,\phi(\mathbf{x}),\,\kappa(\mathbf{c})\big)
\label{eq:p3_s8_3_5_objective}
\end{equation}
returns an edited clip that applies the requested change while preserving unedited content and the source motion, recovering the image editing operator of Eq.~\eqref{eq:t2i_edit_objective} exactly at $T=1$. Eq.~\eqref{eq:p3_s8_3_5_objective} encodes two competing objectives: \emph{editability} (the output must realize the instruction $\mathbf{c}$) and \emph{preservation} (the unedited regions and the source dynamics must survive untouched).

\subsubsection{Methodology}
\label{subsubsec:p3_s8_3_5_methodology}
The five paradigms are overlapping choices of training strategy and preservation mechanism: training-free editing keeps the weights fixed, one-shot tuning updates selected parameters on the source clip, trained instruction editors learn across editing data, attention-based editing reuses or modifies source attention, and structural guidance conditions the edit on source-derived depth, edges, or poses. Many of them share a diffusion recipe that inverts the source clip and re-denoises it under the instruction; writing $\mathrm{Inv}$ for DDIM inversion of the pretrained or adapted backbone $\epsilon_\theta^{0}$ of §\ref{subsec:p3_s8_3_1},
\begin{equation}
\mathbf{z}_S^{1:T}=\mathrm{Inv}\big(\mathbf{z}^{\text{in},1:T};\,\phi(\mathbf{x}^{\text{src}})\big),\qquad
\hat{\mathbf{z}}^{1:T}=\mathrm{Denoise}\big(\mathbf{z}_S^{1:T};\,\phi(\mathbf{x}),\,\kappa(\mathbf{c})\big),
\label{eq:p3_s8_3_5_invert}
\end{equation}
where $S$ denotes the terminal diffusion step, distinct from the clip length $T$. Inversion supplies a source-correlated initialization; how much of the source survives then depends on the sampling, feature reuse, and guidance applied during re-denoising. Figure~\ref{fig:p3_s8_3_5_mechanisms} schematizes the five paradigms, with panels (a)--(e) keyed to the paragraphs below.

\begin{figure}[!th]
\centering
\includegraphics[width=\linewidth]{figures/part3/section8/five_video_editing_paradigms.pdf}
\caption{Five overlapping video-editing paradigms, (a)--(e), applied to one task (top): restyling a driving clip while preserving its layout and motion. The dashed link in (c) marks a paired training example; scene sequences are illustrative.}
\label{fig:p3_s8_3_5_mechanisms}
\end{figure}

\noindent\textbf{Training-Free (Zero-Shot) Methods.} The lightest realization of Eq.~\eqref{eq:p3_s8_3_5_invert} (Figure~\ref{fig:p3_s8_3_5_mechanisms}a) edits a video without any fine-tuning, analogous to zero-shot image editing~\cite{ju2023directinversionboostingdiffusionbased}: DDIM inversion encodes the source into latent codes and noise maps, and the denoising process is rerun on them under the edited prompt. Text2Video-Zero~\cite{Khachatryan_2023_ICCV} and Pix2Video~\cite{ceylan2023pix2videovideoeditingusing} established the category, and vid2vid-zero~\cite{wang2024zeroshotvideoeditingusing}, I2VEdit~\cite{ouyang2024i2veditfirstframeguidedvideoediting}, RAVE~\cite{kara2024rave}, and FastVideoEdit~\cite{zhang2024fastvideoeditleveragingconsistencymodels} add latent-code warping and feature injection to improve temporal consistency without fine-tuning.

\noindent\textbf{One-Shot Tuning Methods.} Trading a little training for sharper fidelity (Figure~\ref{fig:p3_s8_3_5_mechanisms}b), Tune-A-Video~\cite{wu2023tune} fine-tunes a limited subset of a pretrained text-to-image diffusion model's parameters on the frames of a single source video, so that the model learns its content and motion and can regenerate it under a new prompt. Dreamix~\cite{molad2023dreamixvideodiffusionmodels} builds on one-shot fine-tuning to improve both editability and fidelity to the original motion.

\noindent\textbf{Training-Based Methods.} The limited control and quality of inference-time editing have motivated editors that learn a reusable operator from editing data (Figure~\ref{fig:p3_s8_3_5_mechanisms}c), particularly for instruction-based editing. InsV2V~\cite{cheng2023consistentvideotovideotransferusing} extends the InstructPix2Pix~\cite{brooks2023instructpix2pixlearningfollowimage} framework directly to video; GenProp~\cite{Liu_2025_CVPR} propagates an edit from a user-modified first frame to the rest of the sequence; and VideoPainter~\cite{videopainter} uses a video inpainting model as an intermediary. UNIC~\cite{ye2025unicunifiedincontextvideo} handles many editing tasks in one model by concatenating the conditions into a sequence with task-aware positional embeddings, and EditVerse~\cite{ju2025editverseunifyingimagevideo} unifies video generation and editing by representing text, image, and video as one token sequence processed with self-attention, which permits in-context learning and flexible inputs.

\noindent\textbf{Attention-Based Editing.} To protect preservation within the inference-time recipe (Figure~\ref{fig:p3_s8_3_5_mechanisms}d), many techniques manipulate the cross-attention maps directly: following Prompt-to-Prompt in images, they blend the attention maps of the original and edited prompts so that the inverted source structure of Eq.~\eqref{eq:p3_s8_3_5_invert} survives the re-denoising. FateZero~\cite{qi2023fatezerofusingattentionszeroshot} and Video-P2P~\cite{liu2023videop2p} were early examples, and Edit-A-Video~\cite{shin2023editavideosinglevideoediting}, TokenFlow~\cite{tokenflow2023}, UniEdit~\cite{bai2024uniedit}, and AnyV2V~\cite{ku2024anyv2vtuningfreeframeworkvideotovideo} improve temporal consistency at the token or feature level, enabling localized edits such as changing an object's material or color while the scene structure is kept.

\noindent\textbf{Editing with Structural Guidance.} At the explicitly guided end (Figure~\ref{fig:p3_s8_3_5_mechanisms}e), editing is conditioned on spatial signals extracted from the source video, the conditioning machinery of §\ref{subsec:p3_s8_3_2} repurposed for preservation rather than generation. Rerender A Video~\cite{yang2023rerender} feeds the source's depth or edge maps to a ControlNet-like module during editing so that geometry and layout stay intact, and MagicEdit~\cite{liew2023magicedit} accepts a variety of spatial inputs (\eg, pose, sketches, segmentation maps) for high structural fidelity, which suits style transfer and object replacement, where preserving structure is what prevents distortion and keeps motion coherent.

The paradigms trade cost against control: training-free editing avoids per-clip weight updates but may still need iterative inference, whereas a trained instruction editor is expensive to build but generalizes across edits.

\subsubsection{Benchmarks}
\label{subsubsec:p3_s8_3_5_benchmarks}
Video editing has no single canonical suite, and evaluation has moved from shared source-video collections toward purpose-built, fine-grained editing benchmarks. Early evaluations report on DAVIS~\cite{pont20172017}, the de facto source set for object-centric edits, and on the LOVEU text-guided video-editing (TGVE) benchmark~\cite{wu2023loveutgve}, which supplies prompt--clip pairs with human ratings; VPBench, released with VideoPainter~\cite{videopainter}, adds masked clips with dense captions for inpainting-style edits. The most recent suites target instruction-guided editing and pair richer data with model-based scoring. FiVE~\cite{li2025five} contributes $74$ real clips (drawn from DAVIS) and $26$ generated clips across six fine-grained editing types, with $420$ object-level prompt pairs and masks, evaluated under a $15$-metric protocol spanning background preservation, text--video similarity, temporal consistency, video quality, and runtime. IVEBench~\cite{chen2026ivebench} scales this to $600$ source videos spanning seven semantic dimensions and clip lengths from $32$ to $1{,}024$ frames, across eight task categories ($35$ subcategories) whose prompts are generated by LLMs and refined by experts, and adopts a three-dimensional protocol (video quality, instruction compliance, and video fidelity) that fuses traditional metrics with multimodal-LLM-based assessment.

\paragraph{Editing-fidelity Metrics.}
Editing is judged on three axes that mirror the editability--preservation tension of Eq.~\eqref{eq:p3_s8_3_5_objective}. The classic metrics are, for \emph{edit adherence}, the CLIP similarity between each edited frame and the target prompt; for \emph{temporal consistency}, frame-to-frame CLIP similarity together with a warping error that flow-aligns adjacent frames and reports the photometric residual; and for \emph{structure and background preservation}, PSNR / SSIM / LPIPS on the unedited regions against the source, atop the generic FVD~\cite{unterthiner2018fvd} quality score. These scalars correlate only loosely with human judgment on fine-grained edits (an object recolored but mislocalized can still score well), so the latest benchmarks add vision-language judges: FiVE-Acc~\cite{li2025five} uses a VLM to test whether the intended fine-grained edit succeeded, and IVEBench~\cite{chen2026ivebench} scores instruction compliance and fidelity with multimodal LLMs alongside the traditional metrics. Results are reported as an edit-versus-preservation pair, since ignoring the instruction trivially maximizes preservation.

\subsubsection{Challenges and Trends}
\label{subsubsec:p3_s8_3_5_challenges}
Three problems are specific to editing. Pushing the edit harder erodes the unedited regions and the source motion, while protecting them too strongly suppresses the requested change; this editability--preservation trade-off of Eq.~\eqref{eq:p3_s8_3_5_objective} is what every paradigm above negotiates differently. The invert-and-re-denoise recipe of Eq.~\eqref{eq:p3_s8_3_5_invert} is imperfect, and its reconstruction error compounds over long clips, so edits that are clean over a few seconds degrade over minute-scale video. And training-based editors follow only the edit types seen in their training pairs, so a flexible instruction editor that handles arbitrary edits remains open; the unified generate-and-edit models UNIC~\cite{ye2025unicunifiedincontextvideo} and EditVerse~\cite{ju2025editverseunifyingimagevideo}, which treat the source clip as one stream among many in shared attention, are the current attempt.







\subsection{Summary and Critical Challenges and Future Directions}
\label{subsec:p3_s8_3_6}
\label{subsec:p3_s8_3_summary}

Table~\ref{tab:p3_s8_3_subfield_comparison} places the five control axes side by side: the control signal each introduces, the operator it instantiates over the shared sampler of Eq.~\eqref{eq:p3_s8_3_objective}, the paradigms it surveys, the methods anchoring each, and the benchmarks on which it reports.

\begin{table}[t]
    \centering
    \fontsize{8}{10}\selectfont
    \setlength{\tabcolsep}{1.0mm}
    \caption{The five control axes of video spatial generation; each operator specializes the spatio-temporal sampler of Eq.~\eqref{eq:p3_s8_3_objective}.}
    \label{tab:p3_s8_3_subfield_comparison}
    \begin{tabular}{L{0.12\textwidth} L{0.22\textwidth} L{0.18\textwidth} L{0.19\textwidth} L{0.21\textwidth}}
        \toprule
        \textbf{Control axis} & \textbf{Control signal / operator} & \textbf{Dominant paradigms} & \textbf{Canonical methods} & \textbf{Benchmarks / metrics} \\
        \midrule
        §\ref{subsec:p3_s8_3_1} Text-to-video & Text only, $\kappa=\varnothing$; $G_\theta^{\text{vid}}:\mathcal{E}\!\to\!\mathcal{Z}^{\text{vid}}$ (Eq.~\eqref{eq:p3_s8_3_1_objective}) & GAN; autoregressive token; diffusion / diffusion transformer & CogVideo, Sora, Wan, HunyuanVideo, Cosmos & UCF-101, MSR-VTT, VBench, EvalCrafter; FVD, CLIPSIM \\
        §\ref{subsec:p3_s8_3_2} Spatial control & Per-frame scaffold $\kappa_{\text{s}}(\mathbf{c}_{\text{s}}^{1:T})$ (Eq.~\eqref{eq:p3_s8_3_2_objective}); residual injection (Eq.~\eqref{eq:p3_s8_3_2_inject}) & Image, pose, depth, landmark, sketch, bounding box & Animate Anyone, MagicAnimate, ControlVideo, EchoMimic, MagicDrive & nuScenes, pose-annotated video; keypoint / depth / identity / detection adherence \\
        §\ref{subsec:p3_s8_3_3} Temporal control & Motion signal $\kappa_{\text{t}}(\mathbf{m}^{1:T})$ (Eq.~\eqref{eq:p3_s8_3_3_objective}); temporal-branch injection (Eq.~\eqref{eq:p3_s8_3_3_inject}) & Flow, trajectory, camera, motion; latent warping or encoded conditioning & Motion-I2V, Tora, DragNUWA, CameraCtrl, MotionDirector & RealEstate10K, DL3DV-10K, MPI-Sintel; EPE, path / pose error, transfer fidelity \\
        §\ref{subsec:p3_s8_3_4} Other control & Heterogeneous $\kappa_{\text{o}}$ fused by $\mathcal{F}_\psi$ (Eqs.~\eqref{eq:p3_s8_3_4_objective}, \eqref{eq:p3_s8_3_4_inject}) & Style, point, BEV, text, audio; summed adapters, concatenation, unified attention & StyleCrafter, DriveDreamer-2, EMO, VideoComposer, FullDiT & nuScenes, HDTF, VFHQ; style / point / BEV / OCR / lip-sync adherence \\
        §\ref{subsec:p3_s8_3_5} Video editing & Existing clip $+$ edit intent $\to$ edited clip; $E_\theta^{\text{vid}}$ (Eq.~\eqref{eq:p3_s8_3_5_objective}) & Training-free, one-shot tuning, trained instruction editors, attention reuse, structural guidance & Tune-A-Video, TokenFlow, InsV2V, UNIC, EditVerse & DAVIS, LOVEU-TGVE, FiVE, IVEBench; edit-vs-preservation tuple, VLM / MLLM judges \\
        \bottomrule
    \end{tabular}
\end{table}

The axes partition not by output domain, since all five decode a clip from the same $\mathcal{Z}^{\text{vid}}$, but by the control fed to one shared sampler: text-to-video is the $\kappa=\varnothing$ limit, spatial, temporal, and other control progressively enrich $\kappa$, and editing alone reverses the direction of the map, taking an existing clip in place of pure noise.
The injection mechanism converges across axes: the spatial residual branch (Eq.~\eqref{eq:p3_s8_3_2_inject}), the temporal branch (Eq.~\eqref{eq:p3_s8_3_3_inject}), and the multi-condition interface (Eq.~\eqref{eq:p3_s8_3_4_inject}) are one idea, a trainable control pathway added to a pretrained backbone, generalized from one condition to many, and the most recent systems (VideoComposer, FullDiT, UNIC, EditVerse) go further by folding every axis, editing included, into a single token sequence processed by shared attention.
Evaluation has not converged: no suite spans the five axes, each reports generic quality together with a channel-specific adherence score, and the newest editing benchmarks (FiVE, IVEBench) add vision-language and multimodal-LLM judges to improve human alignment.

Video generation inherits every difficulty of the 2D and 3D priors of §\ref{sec:p3_s8_1} and §\ref{sec:p3_s8_2} and adds the time axis, which turns several already-hard problems into open ones; the four below cut across all five axes.

\noindent\textbf{Temporal coherence and long-horizon drift.} The defining difficulty of video over image generation is keeping content stable across frames. The signature artifact is \emph{flicker}, the high-frequency appearance instability between adjacent frames that plagues per-frame and early GAN-based methods (§\ref{subsec:p3_s8_3_1}); over longer spans it gives way to \emph{identity drift}, the gradual mutation of a subject's appearance, color, or fine texture across a clip; and in autoregressive roll-outs to \emph{error accumulation}, the compounding of small per-step mistakes that caps usable clip length. Latent video autoencoders with temporal compression~\cite{blattmann2023stable,blattmann2023align} and dedicated spatio-temporal attention suppress flicker, while autoregressive--diffusion hybrids such as CausVid~\cite{yin2025slow} and MAGI-1~\cite{ai2025magi1autoregressivevideogeneration} combine sequential consistency with diffusion fidelity to push toward minute-scale generation; coherence over arbitrarily long horizons remains unsolved and is the binding constraint on the world-model applications that motivate the section.

\noindent\textbf{Controllability versus fidelity under multi-condition fusion.} As §\ref{subsec:p3_s8_3_2}--§\ref{subsec:p3_s8_3_4} stack heterogeneous conditions (spatial scaffolds, camera paths, audio, style) onto Eq.~\eqref{eq:p3_s8_3_objective}, the conditions begin to conflict. Over-strong conditioning yields rigidity and texture collapse; weak conditioning ignores the control; and naive fusion produces \emph{control-signal dilution} (summed adapters wash one another out), \emph{condition leakage} across frames, and \emph{collapse to the dominant condition} when one modality (\eg, audio in talking-head generation) overwhelms the rest. Compositional and universal frameworks (VideoComposer~\cite{wang2023videocomposercompositionalvideosynthesis}, OmniHuman~\cite{lin2025omnihuman1rethinkingscalinguponestage}, and the unified-attention design of FullDiT~\cite{ju2025fullditmultitaskvideogenerative}) replace per-axis injection with joint token fusion and task-aware positional embeddings, but a principled account of how to weight and arbitrate competing conditions is still missing.

\noindent\textbf{Physical and 3D plausibility.} Even a temporally coherent, well-controlled clip may depict an impossible world, which is the bridge back to §\ref{sec:p3_s8_2}. Generated motion frequently violates object permanence (objects teleport, morph, or vanish), contact and rigid-body dynamics (rigid objects warp non-rigidly; feet slide; occlusions resolve implausibly), and \emph{multi-view consistency} under camera motion, where the same scene is rendered inconsistently from successive viewpoints. Systems presented as steps toward ``world simulators'', Sora~\cite{sora2024} among them, still exhibit these failures on closer inspection. Camera-control methods~\cite{he2024cameractrl,yu2024viewcrafter} and explicitly 3D-aware video diffusion (Diffusion-as-Shader~\cite{gu2025diffusionshader3dawarevideo} and GS-DiT~\cite{bian2025gsditadvancingvideogeneration}) narrow the gap by conditioning generation on geometry, which suggests that the more promising route to physical plausibility runs through the volumetric representations of §\ref{sec:p3_s8_2} rather than through scale alone.

\noindent\textbf{Compute, latency, and evaluation.} Underlying all three is cost: a spatio-temporal latent is an order of magnitude larger than an image latent, so long, high-resolution clips strain memory and make iterative sampling slow, which in turn limits the context length available for coherence and the conditioning budget available for control. Consistency distillation and few-step samplers, together with sparse and causal attention, are the active responses. Evaluation compounds the problem: distributional scores such as FVD and per-frame FID are blind to temporal and control adherence, and automatic metrics correlate weakly with human judgment, so reported progress is hard to compare across methods; benchmarks and human-aligned protocols that jointly score fidelity, temporal coherence, and per-axis control adherence are a central near-term need.

Coherence and physical plausibility both demand longer temporal context, which inflates compute, and every added condition enlarges the evaluation surface, so progress on any one of the four is bounded by the others. Together they mark the last step of the progression traced across §\ref{sec:generation}, from the static 2D prior of §\ref{sec:p3_s8_1} through the static volumetric prior of §\ref{sec:p3_s8_2} to dynamic, navigable models that consume spatial structure over time.

% \newpage


\section{Motion Generation}
\label{sec:p3_s8_4}

\FramedInitial{M}{otion} generation is the temporal, embodied output channel of spatial generation: where §\ref{sec:p3_s8_1}–§\ref{sec:p3_s8_3} synthesize static images, volumetric assets, and controllable video, motion generation produces the \emph{articulated dynamics} that an avatar, character, or robot can actually enact.
A generated motion is a time series of body poses that must be believable in its dynamics, physically plausible (balance, ground contact, non-penetration), and faithful to its driving signal, whether a text description, audio, a scene, or a key frame.
The conditional-sampler template that turns a prompt into pixels (Eq.~\eqref{eq:t2i_generator}) is here instantiated over a sequence of poses in a structured kinematic space, so geometric priors meet physical constraints: the output must respect the biomechanics of an articulated body as well as the semantics of the condition.

The subsections are indexed by the moving agent and then by one cross-cutting operation: human motion (§\ref{subsec:p3_s8_4_1}), animal motion with its cross-species morphology (§\ref{subsec:p3_s8_4_2}), vehicle motion planning, where the body is rigid and the difficulty moves into the multi-agent traffic scene (§\ref{subsec:p3_s8_4_3}), and motion editing (§\ref{subsec:p3_s8_4_4}).
% , with a comparison in §\ref{subsec:p3_s8_4_summary}.




\subsection{Human Motion Generation}
\label{subsec:p3_s8_4_1}
Human motion generation produces natural, varied sequences of human body poses from intentions, environmental stimuli, or explicit conditions, and supports animation, virtual characters in film and games, AR/VR, and human--robot interaction.

\subsubsection{Task Definition}
\label{subsubsec:p3_s8_4_1_task}
Let a motion be a temporal sequence of body poses $\mathbf{m}^{1:T}=(\mathbf{m}^1,\dots,\mathbf{m}^T)$, where each frame $\mathbf{m}^t\in\mathcal{P}$ is a pose in a chosen representation space $\mathcal{P}$ and $\mathcal{M}=\mathcal{P}^T$ is the resulting motion space. Writing $\phi$ for the text encoder of Eq.~\eqref{eq:t2i_generator} and $\kappa:\mathcal{C}\!\to\!\mathcal{C}'$ for an encoder of a driving signal $\mathbf{c}\in\mathcal{C}$ (an action label, a text description, a music or speech track, or a scene), human motion generation learns a conditional sampler
\begin{equation}
\label{eq:p3_s8_4_1_objective}
G_\theta^{\text{mot}}:\mathcal{C}'\!\to\!\mathcal{M},\qquad
\hat{\mathbf{m}}^{1:T}\sim p_\theta\!\left(\mathbf{m}^{1:T}\mid\kappa(\mathbf{c})\right).
\end{equation}
This is the conditional-sampler template of Eq.~\eqref{eq:t2i_generator} instantiated over a temporal pose sequence rather than an image latent: with no driving signal ($\kappa=\varnothing$) it reduces to unconditional motion synthesis, and the choice of $\mathbf{c}$ indexes the five conditioning settings surveyed in §\ref{subsubsec:p3_s8_4_1_method}.

Figure~\ref{fig:p3_s8_4_1_repr} illustrates the pose representations that instantiate $\mathcal{P}$. Broadly, representations fall into keypoint-based and rotation-based categories, and the two are convertible: joint rotations map to spatial keypoints via forward kinematics (FK), and keypoints recover joint rotations via inverse kinematics (IK).

\begin{figure}
    \centering
    \includegraphics[width=\linewidth]{figures/part3/section8/human_motion_repr.png}
    \includegraphics[width=\linewidth]{figures/part3/section8/human_motion_repr_smpl.png}
    \caption{Human motion representations. Top: keypoint-based, as (a) joint positions, (b) joints grouped by kinematic chain, and (c) dense surface markers; and rotation-based, as (d) per-joint rotations in $\mathrm{SO}(3)$ on a kinematic tree. Bottom: the SMPL body model, from (a) the template mesh and skinning weights through (b) shape and (c) pose blend shapes to (d) the skinned, posed mesh.}
    \label{fig:p3_s8_4_1_repr}
\end{figure}

\paragraph{Keypoint-based Representation.} The body is modeled as a collection of keypoints, each a distinct anatomical landmark (joints or salient surface locations) given by its 2D or 3D coordinates in a chosen frame (pixel space or world space), and motion is encoded as a temporal series of such configurations.
Keypoint representations are directly available from many motion-capture systems and are highly interpretable, but many downstream uses (character animation, robotic control) require converting keypoints to joint rotations through the IK problem.
Some works instead augment or replace sparse skeletal keypoints with dense body markers distributed across the body surface~\cite{WuWZZHYT22, TaheriCBT22}, which carry richer information about body shape, limb twisting, and fine geometric detail.

\paragraph{Rotation-based Representation.} A pose is represented by joint angles, the rotations of body segments relative to their parent in a hierarchical kinematic tree. Most studies model these in three dimensions, treating each non-root joint as an element of $\mathrm{SO}(3)$, parameterized as Euler angles, axis--angle vectors, or unit quaternions.

\paragraph{Skinned Multi-Person Linear Model.} Beyond pure joint rotation, statistical mesh models jointly capture pose and body shape, including non-rigid deformations associated with articulation.
The Skinned Multi-Person Linear (SMPL) model~\cite{loper2015smpl} defines a template mesh with fixed topology (6{,}890 vertices, 23 joint rotations plus a root) and uses linear blend skinning augmented by learned blend shapes: identity (shape) blend shapes for inter-subject variation and pose blend shapes for rotation-dependent deformations such as muscle bulging and twisting.



\subsubsection{Methodology}
\label{subsubsec:p3_s8_4_1_method}
Methods for human motion generation are organized along two axes. The first is the \emph{learning paradigm}: regression-based techniques learn a deterministic mapping from conditions to a single target motion, whereas generative approaches model the full conditional distribution of motion using GANs, VAEs, normalizing flows, or diffusion models (the generative paradigms themselves are detailed in §\ref{subsec:p3_s8_1_1}).
Formally, a regression model realizes Eq.~\eqref{eq:p3_s8_4_1_objective} as a point estimate,
\begin{equation}
\label{eq:p3_s8_4_1_regression}
\hat{\mathbf{m}}^{1:T}=f_\theta\!\left(\kappa(\mathbf{c})\right),\qquad
\min_\theta\ \big\|\hat{\mathbf{m}}^{1:T}-\mathbf{m}^{1:T}\big\|,
\end{equation}
which is simple and stable but collapses the inherently one-to-many condition--motion relationship onto an averaged trajectory. Generative models instead sample $\hat{\mathbf{m}}^{1:T}\!\sim\!p_\theta(\mathbf{m}^{1:T}\mid\kappa(\mathbf{c}))$; the now-dominant diffusion instantiation trains a denoiser $G_\theta$ to recover the clean motion at every noise level $n$,
\begin{equation}
\label{eq:p3_s8_4_1_diffusion}
\mathcal{L}=\mathbb{E}_{\mathbf{m}^{1:T},\,n}\,\big\|\,\mathbf{m}^{1:T}-G_\theta(\mathbf{m}^{1:T}_n,\,n,\,\kappa(\mathbf{c}))\,\big\|^2,
\end{equation}
which preserves diversity and supports flexible conditioning, at higher sampling cost.
The second axis is the \emph{conditioning modality} $\mathbf{c}$, which indexes the five settings below; each setting instantiates the same Eq.~\eqref{eq:p3_s8_4_1_objective} and has migrated, broadly, from regression and adversarial models toward diffusion. Figure~\ref{fig:p3_s8_4_1_methods} previews this two-axis organization.



\begin{figure}
    \centering
    \includegraphics[width=\linewidth]{figures/part3/section8/two_axis_human_motion_generation.pdf}
    \caption{Two-axis organization of human motion generation: the conditioning modality (left), encoded by $\kappa$, and the learning paradigm (right), either a regressor that returns one estimate (Eq.~\eqref{eq:p3_s8_4_1_regression}) or a generative model that returns alternative samples, shown with diffusion (Eq.~\eqref{eq:p3_s8_4_1_diffusion}). Branches are alternatives, not consecutive stages.}
    \label{fig:p3_s8_4_1_methods}
\end{figure}

\paragraph{Action-to-Motion.} Given a discrete action label (\eg, ``walk'', ``throw''), these methods generate a corresponding sequence. Early work used GANs with graph convolution and self-attention (SA-GAN~\cite{YuZLYC20}, Kinetic-GAN~\cite{DegardinNLBY022}) for realism and diversity, while VAE/CVAE frameworks (Action2Motion~\cite{GuoZWZSDG020}, ACTOR~\cite{PetrovichBV21}) produced probabilistic, controllable motions. More recent models adopt contrastive latent learning (ODMO~\cite{LuZLR22}), discrete autoregressive sequences (PoseGPT~\cite{LucasBWR22}), and diffusion (MDM~\cite{TevetRGSCB23}, the latent-diffusion hybrid MLD~\cite{ChenJLHFCY23}). As most methods excel only at single actions, MultiAct~\cite{LeeML23} recurrently composes multiple action motions into long sequences via a canonicalization scheme.

\paragraph{Text-to-Motion.} Conditioning on free-form language exploits richer semantic cues than fixed labels allow. Early methods used GANs (Text2Action~\cite{AhnHCYO18}) or joint text--motion embeddings (JL2P~\cite{AhujaM19}, two-stream encoder--decoders~\cite{GhoshCOTS21}); later ones added variational modeling and attention (TEMOS~\cite{PetrovichBV22}, TEACH~\cite{AthanasiouPBV22}).
For zero-shot generalization, MotionCLIP~\cite{TevetGHBC22} and AvatarCLIP~\cite{HongZPCYL22} align the motion latent space with CLIP.
The current frontier combines VQ-VAE tokenization~\cite{GuoZWC22}, autoregressive transformers (T2M-GPT~\cite{zhang2023generating}), and diffusion (FLAME~\cite{KimKC23}), with multimodal conditioning extensions such as MoFusion~\cite{DabralMGT23} fusing text and audio. More recent masked-generative models (MoMask~\cite{GuoMJW024}) push quality and inference speed further by predicting quantized motion tokens with a bidirectional, iteratively-refined masking schedule rather than left-to-right.

\paragraph{Music-to-Dance.} The goal is coherent, expressive dance aligned to a music track. Early approaches framed this as supervised sequence-to-sequence translation via LSTM autoencoders~\cite{TangJM18}, transformers (AI Choreographer~\cite{LiYRK21}), and multi-person frameworks (GroupDancer~\cite{Wang0WXCMCW22}, GDancer~\cite{LePDTT023}).
Because music--dance is inherently one-to-many, generative models model the diverse distribution: MNET~\cite{0001OKT022} with style conditioning, Transflower~\cite{PerezHBHOA21}, Bailando~\cite{SiyaoYGLW0L022} (VQ-VAE + GPT), and EDGE~\cite{TsengCL23} (diffusion). Graph-based methods (ChoreoMaster~\cite{ChenTLZGZ021}, PC-Dance~\cite{GaoPZSZ22}, ChoreoGraph~\cite{AuCJG22}) use motion graphs and shared music--motion embeddings for rhythmic coherence, while curriculum learning, dynamics-prior denoising, and hierarchical pose/motif/choreography modeling target long-term, genre-consistent full-song dances~\cite{HuangHWSZJ21, SunWHLJH22, AristidouYACSC23}.

\paragraph{Speech-to-Gesture.} Co-speech gesture synthesis generates upper-body or full gesture sequences from speech (audio, sometimes text).
Early adversarial methods modeled person-specific gesture patterns~\cite{GinosarBKCOM19}; subsequent work disentangled shared versus motion-specific variation~\cite{KucherenkoHHKK19} and decomposed body, hands, and face~\cite{HabibieXMLSPET21}.
Recent methods employ diffusion~\cite{ZhuLLQLY23, AlexandersonNBH23, AoZL23}, VQ-VAE quantization~\cite{YiLLCWBTB23}, and style or speaker embeddings~\cite{AhujaLNM20} to handle diversity, identity, rhythm, and expression more robustly.

\paragraph{Scene-to-Motion.} Here the condition is a 3D scene, and most methods adopt a multi-stage pipeline that predicts a goal or interaction anchor, plans a trajectory, then generates motion along it. Early works predict a goal from 2D/RGB scene inputs and synthesize motion through trajectory-guided VAEs~\cite{CaoGMCVM20, WangYDL21}, sometimes with scene-aware discriminators. SAMP~\cite{HassanCVSYZB21} models 3D paths and human--object interactions; later methods refine interaction with per-joint contact maps~\cite{D00010HS22} or path-search ($A^\ast$) for trajectory diversity~\cite{0003RLYLD22}, and GOAL~\cite{TaheriCBT22} and SAGA~\cite{WuWZZHYT22} target grasping via two-step goal-to-motion VAEs.
Some assume given goals and stitch sub-goal motions hierarchically for long-term synthesis~\cite{WangXXL021}, while refinement-based frameworks such as CIRCLE~\cite{AraujoLVAWGCL23} initialize by interpolation and refine with scene-geometry-aware modules.


\subsubsection{Benchmarks}
\label{subsubsec:p3_s8_4_1_benchmarks}
Human motion datasets group naturally by conditioning modality; Table~\ref{tab:p3_s8_4_1_datasets} gives their representation, scale, and condition.
Text- and action-conditioned corpora dominate, led by the language-annotated KIT Motion-Language~\cite{PlappertMA16} and HumanML3D~\cite{GuoZZ0JL022}, BABEL~\cite{PunnakkalCAQB21}'s sequence- and frame-level labels over the AMASS~\cite{mahmood2019amass} aggregate, and the action taxonomies of UESTC~\cite{JiXYSSZ18}, NTU-RGB+D~120~\cite{LiuSPWDK20}, and HumanAct12~\cite{GuoZWZSDG020}.
Audio-conditioned resources split into music-to-dance sets (AIST++~\cite{LiYRK21}, AIOZ-GDANCE~\cite{LePDTT023}, and several genre-specific collections) and co-speech gesture sets (Trinity~\cite{FerstlM18}, BEAT~\cite{LiuZIPLZBZ22}, PATS~\cite{AhujaLNM20}), many built by fitting parametric body models to in-the-wild video.
Scene-conditioned datasets (PROX~\cite{HassanCTB19}, GRAB~\cite{TaheriGBT20}, HUMANISE~\cite{WangCLZLH22}, CIRCLE~\cite{AraujoLVAWGCL23}) capture human--scene and human--object interaction with synchronized geometry, from Vicon and IMU rigs to synthetic engines such as GTA-IM~\cite{CaoGMCVM20}. Across all three groups, the cost of high-quality capture forces a trade-off between scale and annotation fidelity.

\begin{table}[htbp]
\centering
\fontsize{8}{10}\selectfont
\setlength{\tabcolsep}{1.2mm}
\caption{Benchmarks for human motion generation, grouped by condition (none, text / action, music, speech, scene/object).}
\label{tab:p3_s8_4_1_datasets}
\begin{tabular}{L{0.18\textwidth} C{0.115\textwidth} C{0.07\textwidth} C{0.08\textwidth} C{0.06\textwidth} C{0.06\textwidth} C{0.13\textwidth} L{0.17\textwidth}}
\toprule
\textbf{Dataset} & \textbf{Representation} & \textbf{Subjects} & \textbf{Sequences} & \textbf{Frames} & \textbf{Length} & \textbf{Condition} & \textbf{Remarks} \\
\midrule
Human3.6M \cite{IonescuPOS14} & Keypoint & 11 & - & 3.6M & 5h & - & 15 actions \\
AMASS \cite{mahmood2019amass} & Rotation & 344 & 11,265 & - & 40h & - & Unifies 15 marker-based MoCap datasets \\
HuMMan \cite{CaiRZLYWFGYPHZL22} & Rotation & 1000 & 400K & 60M & - & - & 500 actions \\
KIT Motion Language \cite{PlappertMA16} & Keypoint & 111 & 3,911 & - & 10.3h & Text & 6.3K text descriptions \\
UESTC \cite{JiXYSSZ18} & Keypoint & 118 & 25.6K & - & 83h & Action & 40 action classes \\
NTU-RGB+D 120 \cite{LiuSPWDK20} & Keypoint & 106 & 114.5K & - & 74h & Action & 120 action classes \\
HumanAct12 \cite{GuoZWZSDG020} & Keypoint & 12 & 1,191 & 90K & 6h & Action & 12 action classes \\
BABEL \cite{PunnakkalCAQB21} & Rotation & 344 & - & - & 43.3h & Text & 260 action classes \\
HumanML3D \cite{GuoZZ0JL022} & Keypoint + Rotation & 344 & 14.6K & - & 28.5h & Text & 44.9K text descriptions \\
DanceWithMelody \cite{TangJM18} & Keypoint & - & 61 & 907K & 1.6h & Music & 4 genres \\
Lee \cite{LeeY0WLYK19} & Keypoint & - & 361K & - & 71h & Music & 3 genres \\
Huang \cite{HuangHWSZJ21} & Keypoint & - & 790 & - & 12h & Music & 3 genres\\
AIST++ \cite{LiYRK21} & Rotation & 30 & 1,408 & 10.1M & 5.2h & Music & 10 genres \\
PMSD \cite{PerezHBHOA21} & Keypoint & 8 & - & - & 3.8h & Music & 4 genres \\
ShaderMotion \cite{PerezHBHOA21} & Keypoint & 8 & - & - & 10.2h & Music & 2 genres \\
Chen \cite{ChenTLZGZ021} & Rotation & - & - & 160K & 9.9h & Music & 9 genres \\
PhantomDance \cite{LiZZS22} & Rotation & - & 260 & 795K & 3.7h & Music & 13 genres \\
MMD-ARC \cite{GaoPZSZ22} & Rotation & - & 213 & - & 11.3h & Music & - \\
MDC \cite{Wang0WXCMCW22} & Rotation & - & 798 & - & 3.5h & Music & 2 genres \\
Aristidou \cite{AristidouYACSC23} & Rotation & 32 & - & - & 2.4h & Music & 3 genres \\
AIOZ-GDANCE \cite{LePDTT023} & Rotation & >4000 & - & - & 16.7h & Music & 7 styles, 16 genres \\
Trinity \cite{FerstlM18} & Keypoint & 1 & 23 & - & 4.1h & Speech & Casual talks\\
TED-Gesture \cite{YoonKJLKL19} & Keypoint & - & 1,295 & - & 52.7h & Text & TED talks \\
Speech2Gesture \cite{GinosarBKCOM19} & Keypoint & 10 & - & - & 144h & Speech & TV shows, lectures \\
TED-Gesture++ \cite{YoonCLJLKL20} & Keypoint & - & 1,766 & - & 97.0h & Speech, Text & Extension of \cite{YoonKJLKL19} \\
PATS \cite{AhujaLNM20} & Keypoint & 25 & - & - & 251h & Speech, Text & Extension of \cite{GinosarBKCOM19} \\
Speech2Gesture-3D \cite{HabibieXMLSPET21} & Keypoint & 6 & - & - & 33h & Speech & Videos from \cite{GinosarBKCOM19}\\
BEAT \cite{LiuZIPLZBZ22} & Rotation & 30 & 2,580 & 30M & 76h & Speech, Text, Emotion & 8 emotions, 4 languages \\
Chinese Gesture \cite{AoGLCL22} & Rotation & 5 & - & - & 4h & Speech, Text & Chinese \\
ZEGGS \cite{GhorbaniFHTC23} & Rotation & 1 & 67 & - & 2.3h & Speech, Style & 19 styles \\
SHOW \cite{YiLLCWBTB23} & Rotation & - & - & - & 27h & Speech & Videos from \cite{GinosarBKCOM19} \\
WBHM \cite{ManderyTDVA15} & Rotation & 43 & 3,704 & 691K & 7.68h & Object & 41 objects \\
PiGraph \cite{Hanrahan16a} & Keypoint & 5 & 63 & 0.1M & 2h & Scene, Object & 30 scenes, 19 objects \\
PROX \cite{HassanCTB19} & Rotation & 20 & 60 & 0.1M & 1h & Scene, Object & 12 indoor scenes\\
i3DB \cite{MonszpartGCYM19} & Keypoint & 1 & - & - & - & Scene, Object & 15 scenes\\
GTA-IM \cite{CaoGMCVM20} & Keypoint & 50 & 119 & 1M & - & Scene & Synthetic, 10 indoor scenes \\
GRAB \cite{TaheriGBT20} & Rotation & 10 & 1,334 & 1.6M & - & Object & 51 objects \\
HPS \cite{GuzovMSP21} & Rotation & 7 & - & 300K & - & Scene & 8 large scenes\\
SAMP \cite{HassanCVSYZB21} & Rotation & 1 & - & 185K & 0.83h & Scene, Object & 7 objects\\
COUCH \cite{ZhangBSGP22} & Rotation & 6 & >500 & - & 3h & Scene, Chairs & 3 chairs\\
HUMANISE \cite{WangCLZLH22} & Rotation & - & 19.6K & 1.2M & - & Scene, Object, Text & 643 scenes \\
CIRCLE \cite{AraujoLVAWGCL23} & Rotation & 5 & >7K & 4.3M & 10h & Scene & 9 scenes \\
\bottomrule
\end{tabular}
\end{table}

\paragraph{Evaluation Metrics.} Evaluating synthesized motion is hard because the condition--motion map is one-to-many, judgments are partly subjective, and quality depends on high-level conditional cues, so the literature reports an ensemble of metrics along three axes.
\textbf{Fidelity} captures naturalness, smoothness, and plausibility: distance-based measures such as NPSS~\cite{GopalakrishnanM19} and NDMS~\cite{TankeZG21}, keypoint accuracy (PCK~\cite{YangR13}), motion-space statistics (AVE, Hellinger distance over speed histograms~\cite{yang2023qpgesture}, Fréchet Distance over pose rotations~\cite{HassanCVSYZB21, 0003RLYLD22}, PSKL~\cite{RuizGM19}), feature-space distributional distances led by FID (with IS~\cite{AhujaLNM20}, MMD~\cite{YuZLYC20}, and related scores), and physical-plausibility proxies for foot sliding and foot--ground contact (FCS~\cite{TsengCL23}, non-collision ratios~\cite{WuWZZHYT22}, contact accuracy~\cite{ZhangT22}), the last still lacking a standardized definition.
\textbf{Diversity} guards against the long-sequence ``freezing'' failure via intra-motion variation and freezing rate~\cite{SunWHLJH22}, inter-motion manifold coverage, and multimodality metrics (n-APD, ADE/FDE and their multimodal variants~\cite{YuanK20}) computed by sampling many motions per condition.
\textbf{Condition consistency} measures alignment to the driving signal: recognition accuracy and R-Precision for text/action, feature-level text--motion distance (Multimodal Distance~\cite{GuoZZ0JL022}, mCLIP and MID~\cite{KimKC23}), audio--motion beat metrics (beat coverage, hit rate~\cite{LeeY0WLYK19}, beat alignment~\cite{LiYRK21}, SRGR~\cite{LiuZIPLZBZ22}, Semantic Score~\cite{AoZL23}), and, for scene conditioning, non-collision and human--scene contact scores~\cite{WuWZZHYT22, ZhangBSGP22}. No single metric is decisive, so methods are compared on a task-appropriate composition of the three.


\subsubsection{Challenges and Trends}
\label{subsubsec:p3_s8_4_1_challenges}
\textbf{Data scarcity} is the most basic constraint: high-quality motion capture is expensive and slow, so dataset scale trades against annotation fidelity, and heterogeneous pose representations and conditioning signals fragment the data that exist; fusing heterogeneous sources through weak supervision and multimodal foundation models is the main response.
\textbf{Semantic alignment} is the second. Motion encodes culture, communication, and context, not only mechanics, and a model that fits joint trajectories while missing their relation to language, music, or environment produces motion that is kinematically clean but perceptually empty; semantic priors from language and foundation models help, but how to measure the result is unresolved.
That is the \textbf{evaluation} problem: every objective metric above has blind spots, none replaces a user study, and the field still needs metrics that are both interpretable and better correlated with human perception.
\textbf{Controllability} matters for practical use as much as in image generation, and joint masking, attribute prompts, and style conditioning are moving toward editing tools with textual, spatial, and temporal handles (§\ref{subsec:p3_s8_4_4}).
\textbf{Interactivity} remains underexplored: most methods animate one person in a static scene, and socially coupled groups (conversation, group dance) and motion grounded in dynamic, actionable scenes need both new corpora and new models.





\subsection{Animal Motion Generation}
\label{subsec:p3_s8_4_2}
Most motion-generation work targets the human body and ignores the diversity of animal behavior, although synthesizing animal motion matters for wildlife conservation, ecology, and biomechanics.
Relative to §\ref{subsec:p3_s8_4_1}, animal motion is harder on two fronts that the human setting never confronts: species vary widely in skeletal topology and morphology, so a single fixed-skeleton representation no longer suffices, and the same textual description (``walk stealthily'') maps to radically different gaits across species.
These two pressures, \ie, data scarcity and morphological heterogeneity, shape every method below, which is why animal motion generation borrows the architectures, datasets, and training recipes of human motion generation while adding species-conditioning machinery on top.



\subsubsection{Task Definition}
\label{subsubsec:p3_s8_4_2_task}
Animal motion generation extends the human sampler of Eq.~\eqref{eq:p3_s8_4_1_objective} with an explicit species (morphology) descriptor. Let $s\in\mathcal{S}$ index a species and let $\mathcal{P}_s$ be its species-dependent pose space (its skeleton topology and joint count depend on $s$), so that $\mathcal{M}_s=\mathcal{P}_s^{T}$ is the corresponding motion space. With $\kappa$ encoding a driving signal $\mathbf{c}\in\mathcal{C}$ (an action class, a text description, an anchor pose, or past motion), the task is to learn
\begin{equation}
\label{eq:p3_s8_4_2_objective}
G_\theta^{\text{amot}}:\mathcal{C}'\!\times\!\mathcal{S}\!\to\!\mathcal{M}_s,\qquad
\hat{\mathbf{m}}^{1:T}\sim p_\theta\!\left(\mathbf{m}^{1:T}\mid\kappa(\mathbf{c}),\,s\right).
\end{equation}
Restricting $\mathcal{S}$ to the single human morphology recovers Eq.~\eqref{eq:p3_s8_4_1_objective} exactly; the open difficulty is that $s$ ranges over a heterogeneous, long-tailed set of skeletons for which paired text--motion data are scarce, so the species condition must do double duty as both a generation control and a regularizer that shares statistical strength across morphologies.


\subsubsection{Methodology}
\label{subsubsec:p3_s8_4_2_method}
Almost all animal motion generation is text-driven, and the methods share a common response to the data-scarcity and morphology problems of Eq.~\eqref{eq:p3_s8_4_2_objective}: transfer representations learned from abundant human motion, then specialize them to a species through \emph{feature modulation}. Concretely, given a backbone feature $\mathbf{h}$ at some layer, species-aware modulation applies a FiLM-style affine transform conditioned on $s$,
\begin{equation}
\label{eq:p3_s8_4_2_modulation}
\tilde{\mathbf{h}}=\gamma(s)\odot\mathbf{h}+\beta(s),
\end{equation}
where $\gamma(s),\beta(s)$ are predicted from the species descriptor, allowing one shared network to adapt its internal representation to diverse anatomical forms. The methods differ in how they obtain the shared representation and how strictly they enforce cross-species plausibility, as Figure~\ref{fig:p3_s8_4_2_methods} summarizes.

\textbf{OmniMotionGPT}~\cite{YangZSWXHBQ024} tackles limited animal data by jointly training motion autoencoders on human and animal corpora and aligning their latent encodings to text via CLIP embeddings; it relies most directly on transfer, and at inference it still depends on human-motion priors, which prevents truly text-only animal generation.
\textbf{Motion Avatar}~\cite{0006W0CZH0F0024} uses SinMDM~\cite{RaabLTABC24} for data augmentation in its ZooGen pipeline (the route by which the Zoo-300K dataset is built) and extends the masked-motion modeling of MoMask~\cite{GuoMJW024} to animals, adding an LLM planner that orchestrates joint motion and 3D-avatar mesh generation; it makes little use of species-specific structure, which is what Eq.~\eqref{eq:p3_s8_4_2_modulation} addresses.
\textbf{AniMo}~\cite{WangRZW25} makes that remedy central, injecting species-aware modulation into a joint-aware spatiotemporal tokenizer and learning, through masked modeling, a mapping from text plus species attributes to motion tokens; its AniMo4D dataset (78{,}149 sequences over 114 species) yields markedly better morphological and behavioral fidelity than transfer-only baselines.
\textbf{X-MoGen}~\cite{abs-2508-05162} pushes toward a single model spanning humans and animals: it embeds motion into a shared space with graph-based VAEs and latent autoencoders, generates via masked modeling, and adds a morphological-consistency module that enforces skeletal plausibility across species, supported by the UniMo4D corpus (${\sim}115$ species, ${\sim}119$K sequences). OmniMotionGPT and Motion Avatar still lean on human priors, whereas AniMo and X-MoGen promote species morphology to a first-class, explicitly regularized conditioning variable, so the four together trace a progression from transfer toward native cross-species modeling.

\begin{figure}[!ht]
    \centering
    \includegraphics[width=\linewidth]{figures/part3/section8/cross_species_animal_motion_generation.pdf}
    \caption{Abstract cross-species motion generation: human and animal corpora support a shared representation, species-dependent modulation (Eq.~\eqref{eq:p3_s8_4_2_modulation}) adapts features, and an action-conditioned sampler produces motion for the selected skeleton. Dog, bird, and horse sequences illustrate alternative target morphologies, not simultaneous outputs. The dashed constraint denotes optional morphology regularization; implementations vary across methods.}
    \label{fig:p3_s8_4_2_methods}
\end{figure}


\subsubsection{Benchmarks}
\label{subsubsec:p3_s8_4_2_benchmarks}
Animal motion resources fall into two groups (Table~\ref{tab:p3_s8_4_2_datasets}).
\textbf{Capture and pose-estimation corpora} supply ground-truth articulated motion but no language: RGBD-Dog~\cite{Kearney0PKC20} and 3DDogs-Wild~\cite{abs-2406-14412} pair RGB-D or RGB video with marker-based mocap for canine pose, and DeformingThings4D~\cite{LiTTZN21} provides synthetic 4D sequences with dense mesh-deformation annotations across humanoid and animal categories.
\textbf{Text-to-motion corpora} drive recent progress: AnimalML3D~\cite{YangZSWXHBQ024} is the first text--animal-motion dataset, and the field then scales sharply through Zoo-300K~\cite{0006W0CZH0F0024} (built by the ZooGen pipeline with SinMDM and LLM annotators), AniMo4D~\cite{WangRZW25} (normalized to a 30-joint cross-species skeleton), and UniMo4D~\cite{abs-2508-05162}, which merges AniMo4D with HumanML3D and KIT-ML. Even the largest text-conditioned sets remain smaller and more morphologically fragmented than their human counterparts of §\ref{subsubsec:p3_s8_4_1_benchmarks}, which is what forces the transfer-and-modulate methodology above.



\begin{table}[!ht]
\centering
\fontsize{8}{10}\selectfont
\setlength{\tabcolsep}{1.0mm}
\caption{Benchmarks for animal motion generation. ``Subjects / species'' counts individual animals for the canine capture sets and species or categories elsewhere; ``Texts / frames'' gives text descriptions for text-conditioned corpora and frames for capture corpora.}
\label{tab:p3_s8_4_2_datasets}
\begin{tabular}{L{0.21\textwidth} L{0.13\textwidth} L{0.13\textwidth} C{0.15\textwidth} C{0.14\textwidth} C{0.14\textwidth}}
\toprule
\textbf{Dataset} & \textbf{Type} & \textbf{Modality} & \textbf{Subjects / species} & \textbf{Sequences} & \textbf{Texts / frames} \\
\midrule
RGBD-Dog \cite{Kearney0PKC20} & Canine & RGB-D, mocap & 5 dogs & 5 motion types & 8,346 frames \\
3DDogs-Wild \cite{abs-2406-14412} & Canine & RGB, mocap & 37 dogs & 143 recordings & - \\
DeformingThings4D \cite{LiTTZN21} & Animal / humanoid & 4D mesh & 31 categories & 1,972 & 122,365 frames \\
AnimalML3D \cite{YangZSWXHBQ024} & Animal / humanoid & 4D & 36 identities & 1,240 & - \\
Zoo-300K \cite{0006W0CZH0F0024} & Animal & 3D mesh & 65 categories & ${\sim}$300K text--motion pairs & - \\
AniMo4D \cite{WangRZW25} & Animal & 4D & 114 species & 78,149 & 185,435 texts \\
UniMo4D \cite{abs-2508-05162} & Animal + human & 4D & ${\sim}$115 species & ${\sim}$119K & - \\
\bottomrule
\end{tabular}
\end{table}

\paragraph{Evaluation Metrics.} Animal motion evaluation reuses the three-axis human protocol of §\ref{subsubsec:p3_s8_4_1_benchmarks}, adapted for species-specific biomechanics.
\emph{Fidelity} is measured by FID against real movement distributions and by MPJPE against ground-truth poses, augmented with physical-plausibility checks for gait-cycle regularity, ground contact, and the absence of foot sliding or limb penetration.
\emph{Condition alignment} carries over R-Precision and Multimodal Distance to score text-/action-to-motion correspondence, and \emph{diversity / multimodality} metrics quantify behavioral variety within and across species. The animal-specific addition is perceptual and biomechanical judgment of naturalness, balance, and anatomical consistency, which matters more for quadrupeds and multi-limbed morphologies than for the single human skeleton.



\subsubsection{Challenges and Trends}
\label{subsubsec:p3_s8_4_2_challenges}
\textbf{Physical realism} is the most basic gap: most methods do not simulate ground-contact forces, quadruped balance, or the fluid dynamics of swimming and flight, so artifacts the human setting has largely learned to suppress reappear, and differently for every body plan.
\textbf{Behavioral diversity} is harder than for humans because one instruction such as ``walk stealthily'' realizes into different gaits, postures, and speeds across species, so the conditional distribution to be captured is far wider while the data to capture it are scarcer.
\textbf{Unseen species and radical morphologies} remain out of reach: species-aware modulation (Eq.~\eqref{eq:p3_s8_4_2_modulation}) and morphological-consistency constraints generalize across the skeletons seen in training, but insects and winged or aquatic creatures violate the shared-skeleton, fixed-topology assumptions on which current methods rest.
AniMo and X-MoGen already move toward treating morphology as a variable rather than a template, and topology-flexible representations are the direction this points.




\subsection{Vehicle Motion Planning}
\label{subsec:p3_s8_4_3}
Where §\ref{subsec:p3_s8_4_1} and §\ref{subsec:p3_s8_4_2} generate the motion of an articulated body, vehicle motion planning generates the trajectory of a (near-)rigid agent moving through a traffic scene.
The body is far simpler, since a vehicle has essentially no internal articulation, but the difficulty moves outward into the environment: a planned trajectory must obey vehicle dynamics, stay on the drivable surface dictated by an HD map, and remain socially compatible with the future motion of many surrounding agents whose intentions are unobserved.
This is the setting where motion generation becomes inseparable from \emph{multi-agent} spatial reasoning and physical feasibility: trajectory forecasting (predicting how others will move) and planning (deciding how the ego vehicle should move) are two readings of the same generative object.

\subsubsection{Task Definition}
\label{subsubsec:p3_s8_4_3_task}
Let a trajectory be a sequence of agent states $\boldsymbol{\tau}^{1:H}=(\mathbf{s}^1,\dots,\mathbf{s}^H)$ over a horizon $H$, where each state $\mathbf{s}^t$ collects position, heading, and velocity.
The driving signal is now a structured scene context $\mathbf{c}=(\mathbf{s}^{-P:0},\,\mathcal{G},\,\mathcal{A})$ comprising the agent's own past states, an HD-map / lane graph $\mathcal{G}$, and the observed histories of neighboring agents $\mathcal{A}$.
With $\kappa$ encoding this context, vehicle motion planning learns
\begin{equation}
\label{eq:p3_s8_4_3_objective}
G_\theta^{\text{veh}}:\mathcal{C}'\!\to\!\mathcal{T},\qquad
\hat{\boldsymbol{\tau}}^{1:H}\sim p_\theta\!\left(\boldsymbol{\tau}^{1:H}\mid\kappa(\mathbf{c})\right)
\quad\text{s.t.}\quad
\mathbf{s}^{t+1}=f_{\text{dyn}}(\mathbf{s}^{t},\mathbf{u}^{t}),
\end{equation}
the conditional-sampler template of Eq.~\eqref{eq:p3_s8_4_1_objective} with the pose sequence replaced by a state trajectory and an added kinematic-feasibility constraint $f_{\text{dyn}}$ (\eg, a bicycle model) that distinguishes this setting from articulated-body generation.
Dropping $\mathcal{A}$ recovers single-agent forecasting; supplying a goal or reward in place of $\mathbf{c}$ turns prediction into planning.


\subsubsection{Methodology}
\label{subsubsec:p3_s8_4_3_method}
Prediction methods vary along two axes: how the scene context is \emph{encoded} and how the multimodal future is \emph{decoded}.
Early encoders rasterize the scene into a bird's-eye-view image, but the field has converged on \emph{vectorized} encoders that represent lanes and agent tracks as polylines: VectorNet~\cite{gao2020vectornet} and LaneGCN~\cite{liang2020lanegcn} build graph representations of the HD map, and transformer encoders such as Wayformer~\cite{nayakanti2023wayformer} and the Motion Transformer (MTR)~\cite{shi2022mtr} attend jointly over map and agents; graph-attention predictors~\cite{ZhangFWH22} model spatial and temporal interaction in the same spirit.
On the decoding side, the defining property is multimodality, since a single context admits many plausible futures, so models predict a set of $K$ modes with probabilities,
\begin{equation}
\label{eq:p3_s8_4_3_multimodal}
p_\theta(\boldsymbol{\tau}\mid\kappa(\mathbf{c}))=\sum_{k=1}^{K}\pi_k\,\mathcal{N}\!\left(\boldsymbol{\tau};\,\mu_k(\kappa(\mathbf{c})),\,\Sigma_k\right),
\end{equation}
rather than the single point estimate of Eq.~\eqref{eq:p3_s8_4_1_regression}, which here would average distinct maneuvers (turning \vs going straight) into an infeasible mean.
The earliest learned predictors modeled social interaction directly in agent-coordinate space, Social-LSTM~\cite{alahi2016social} through pooled recurrent states and Social-GAN~\cite{gupta2018social} through adversarial sampling, and Trajectron++~\cite{salzmann2020trajectron} added dynamic feasibility to a graph-structured latent-variable model.
Goal-based decoders make the modes explicit: DenseTNT~\cite{gu2021densetnt} samples dense candidate endpoints on the map and scores trajectories conditioned on them, exposing the maneuver structure that Eq.~\eqref{eq:p3_s8_4_3_multimodal} captures only implicitly.
Generative decoders replace the fixed mixture with a learned sampler: MotionDiffuser~\cite{jiang2023motiondiffuser} denoises whole multi-agent trajectory sets jointly.

Planning turns the same machinery on the ego vehicle. ChauffeurNet~\cite{bansal2019chauffeurnet} learns the ego policy by imitating expert driving and augments the demonstrations with synthesized perturbations and worst cases, so that the policy stays robust when it drifts off the data manifold.
Figure~\ref{fig:p3_s8_4_3_methods} situates these families along the encode/decode axes.

\begin{figure}[!ht]
    \centering
    \includegraphics[width=\linewidth]{figures/part3/section8/vehicle_motion_planning.pdf}
    \caption{Vehicle motion planning as conditional trajectory generation: a vectorized encoder represents map polylines and agent interactions, and a multimodal decoder generates candidate futures (Eq.~\eqref{eq:p3_s8_4_3_multimodal}). Gray intensity indicates mode weight; blue marks a plan selected using a goal or cost and feasibility constraints, not necessarily the highest-probability mode. Road geometry and trajectories are illustrative.}
    \label{fig:p3_s8_4_3_methods}
\end{figure}

\subsubsection{Benchmarks}
\label{subsubsec:p3_s8_4_3_benchmarks}
Progress is measured on large real-world driving logs that pair agent tracks with HD maps: nuScenes~\cite{caesar2020nuscenes}, Argoverse and Argoverse~2~\cite{Argoverse-ChangLSSBHW0LRH19, wilson2023argoverse}, and the Waymo Open Motion Dataset~\cite{ettinger2021womd}, the last two offering the scale and interaction density that current transformer models require.
Evaluation centers on the multimodal-aware displacement metrics that match Eq.~\eqref{eq:p3_s8_4_3_multimodal}: minimum average and final displacement error over the $K$ predicted modes (minADE$_K$, minFDE$_K$), miss rate (the fraction of scenes where no mode lands within a tolerance of the ground truth), and the soft mean-average-precision of the Waymo protocol, which jointly scores accuracy and the calibration of mode probabilities $\pi_k$.
Off-road and collision rates further test the feasibility constraint $f_{\text{dyn}}$ that purely displacement-based scores ignore.

\subsubsection{Challenges and Trends}
\label{subsubsec:p3_s8_4_3_challenges}
A planner must assign trustworthy probabilities to the futures it enumerates, not merely list them: an over-confident missed mode of Eq.~\eqref{eq:p3_s8_4_3_multimodal} is a safety failure rather than a metric error, so \textbf{calibrated multimodality} is the first open problem.
\textbf{Feasibility and safety} are the second: displacement metrics reward trajectories a vehicle cannot physically execute, so feasibility has to be enforced inside the model ($f_{\text{dyn}}$) rather than checked afterward, and the long tail of near-collision scenarios is exactly where data are scarcest.
The third is \textbf{closed-loop distribution shift}: models trained to imitate logged, open-loop futures drift once their own predictions feed back into the scene, the difficulty ChauffeurNet's worst-case synthesis was built to expose.
Diffusion- and transformer-based scene-level samplers that forecast all agents jointly and run closed-loop are the current direction, and they connect directly to the controllable-video and world-model lines of §\ref{sec:p3_s8_3} and §\ref{sec:p3_s8_6}, in which a driving world model both predicts and is conditioned on the ego plan.







\subsection{Motion Editing}
\label{subsec:p3_s8_4_4}
The subsections above \emph{synthesize} motion from scratch; motion editing instead modifies an \emph{existing} motion so that it satisfies a new constraint while leaving the rest intact. It is the temporal, articulated analogue of the image editing of §\ref{subsec:p3_s8_1_3} and the video editing of §\ref{subsec:p3_s8_3_5}, and its demand is again selective preservation: change a targeted aspect (a trajectory, a joint, a style, a semantic action) while holding the body's identity, timing, and untouched portions fixed, which makes editing the controllability layer that closes the generate-and-refine loop for motion.



\subsubsection{Task Definition}
\label{subsubsec:p3_s8_4_4_task}
Given a source motion $\mathbf{m}^{\text{in}}\!\in\!\mathcal{M}$, an edit region $\mathbf{M}\!\in\!\{0,1\}^{T\times J}$ over time $T$ and joints $J$, and an edit control $\mathbf{c}$ (a text instruction, target keyframes, a desired root trajectory, or a target joint), motion editing learns an operator
\begin{equation}
\label{eq:p3_s8_4_4_objective}
E_\theta:\mathcal{M}\!\times\!\{0,1\}^{T\times J}\!\times\!\mathcal{C}'\!\to\!\mathcal{M},\qquad
\hat{\mathbf{m}}=E_\theta\!\left(\mathbf{m}^{\text{in}},\mathbf{M},\kappa(\mathbf{c})\right),
\end{equation}
which lifts the image-editing operator of Eq.~\eqref{eq:t2i_edit_objective} into the motion domain. As in image and video editing, Eq.~\eqref{eq:p3_s8_4_4_objective} encodes two competing objectives, \emph{editability} (the output must realize the control $\mathbf{c}$ inside $\mathbf{M}$) and \emph{preservation} (it must match $\mathbf{m}^{\text{in}}$ outside $\mathbf{M}$ and remain a plausible, smooth motion), and methods differ chiefly in how they trade these off.



\subsubsection{Methodology}
\label{subsubsec:p3_s8_4_4_method}
Because a pretrained motion diffusion model already samples plausible motion, most editors realize Eq.~\eqref{eq:p3_s8_4_4_objective} without retraining, by \emph{conditioning the sampler} of Eq.~\eqref{eq:p3_s8_4_1_diffusion} rather than learning a dedicated $E_\theta$. The simplest mechanism is \textbf{masked infilling}: at every denoising step the observed entries are overwritten with the (noised) source while only the masked entries are generated,
\begin{equation}
\label{eq:p3_s8_4_4_inpaint}
\hat{\mathbf{m}}=(1-\mathbf{M})\odot\mathbf{m}^{\text{in}}+\mathbf{M}\odot G_\theta\!\left(\mathbf{m}_n,\,n,\,\kappa(\mathbf{c})\right),
\end{equation}
which directly enforces preservation by construction. This is the basis of motion in-betweening, where $1-\mathbf{M}$ marks given keyframes: robust in-betweening~\cite{harvey2020robust} predates the diffusion era, and CondMDI~\cite{cohan2024condmdi} casts flexible keyframe-conditioned in-betweening as masked diffusion, while PriorMDM~\cite{shafir2024priormdm} reuses a frozen motion-diffusion prior for editing and composition tasks such as joining and in-filling. A second family handles \emph{spatial} controls that masking alone cannot express (a target root path, or a specific joint at specific times) through classifier-style guidance: GMD~\cite{karunratanakul2023gmd} and OmniControl~\cite{xie2024omnicontrol} steer the denoiser toward arbitrary spatial constraints at any joint and time. A third, optimization-based family treats the initial diffusion noise itself as the free variable, solving $\min_{\mathbf{z}}\,\mathcal{L}_{\text{edit}}(\mathcal{D}(\mathbf{z}))$ so that the decoded motion satisfies an arbitrary differentiable criterion: DNO~\cite{karunratanakul2024dno} shows this single mechanism serves editing, in-betweening, and refinement as one universal motion prior. Finally, \emph{semantic} editors take a natural-language instruction in place of a geometric mask: MotionFix~\cite{athanasiou2024motionfix} learns text-driven edits from paired source--target motions, and LLM-driven approaches~\cite{goel2024iterative} decompose an instruction into a program of primitive edits applied iteratively. Figure~\ref{fig:p3_s8_4_4_methods} groups these four mechanisms by how they obtain $E_\theta$, paralleling the editing taxonomies of §\ref{subsec:p3_s8_1_3} and §\ref{subsec:p3_s8_3_5}.

\begin{figure}[!ht]
    \centering
    \includegraphics[width=\linewidth]{figures/part3/section8/four_motion_editing_mechanisms.pdf}
    \caption{Four motion-editing mechanisms: (a) reinsert known states while infilling missing frames; (b) guide denoising with root-path or joint targets; (c) optimize initial noise through a fixed decoder using an edit loss; and (d) edit source motion from text, using a learned editor or iterative primitive edits. Rendered people illustrate edits rather than measured model outputs.}
    \label{fig:p3_s8_4_4_methods}
\end{figure}

\subsubsection{Benchmarks}
\label{subsubsec:p3_s8_4_4_benchmarks}
Motion editing inherits the corpora of §\ref{subsubsec:p3_s8_4_1_benchmarks} (HumanML3D~\cite{GuoZZ0JL022} and BABEL~\cite{PunnakkalCAQB21} over AMASS), used in editing protocols where keyframes, joints, or trajectories are held out as the constraint, together with the dedicated MotionFix~\cite{athanasiou2024motionfix} benchmark of paired source--edit--instruction triples that supplies the ground truth most editing tasks lack.
Evaluation must report the editability--preservation pair jointly: editability through control adherence (instruction-conditioned R-Precision, target-joint or trajectory error) and overall realism (FID, foot-skating), and preservation through the deviation from the source outside the edit region.
As in video editing, single scalars correlate weakly with perceived edit quality, so reporting a tuple, and increasingly a learned or VLM-style judge of instruction compliance, is becoming standard.

\subsubsection{Challenges and Trends}
\label{subsubsec:p3_s8_4_4_challenges}
The central tension is the \textbf{editability--preservation trade-off} of Eq.~\eqref{eq:p3_s8_4_4_objective}: pushing harder on a control propagates unwanted change through the kinematic chain (editing one arm disturbs balance and foot contact), so a good editor must respect the body's global coordination, not just the masked entries.
This is compounded by \textbf{sparse-control ambiguity}: a single keyframe or a root path under-determines the full-body motion, and the model must fill the gap plausibly, which the guidance and noise-optimization families address with differing robustness.
Underlying both is an \textbf{evaluation} gap sharper than in synthesis: because the ``correct'' edit is rarely unique, ground-truth-based metrics under-reward valid alternatives, motivating the paired-triplet and judge-based protocols above.
The trend mirrors the rest of §\ref{sec:p3_s8_4} and §\ref{sec:p3_s8_3}: from per-mechanism, training-free editors toward unified models that generate, control, and edit motion within a single conditional sampler, increasingly driven by natural-language instructions rather than geometric masks.








% \subsection{Summary}
% \label{subsec:p3_s8_4_summary}
% Table~\ref{tab:p3_s8_4_subfield_comparison} compares the four subfields as specializations of the human-motion sampler of Eq.~\eqref{eq:p3_s8_4_1_objective}: by the moving object and its condition, the sampler or operator each instantiates, the dominant method families, and the benchmarks on which they report.

% \begin{table}[!ht]
% \centering
% \fontsize{8}{10}\selectfont
% \setlength{\tabcolsep}{1.0mm}
% \caption{The four subfields of motion generation; each sampler or operator specializes Eq.~\eqref{eq:p3_s8_4_1_objective}.}
% \label{tab:p3_s8_4_subfield_comparison}
% \begin{tabular}{L{0.13\textwidth} L{0.19\textwidth} L{0.16\textwidth} L{0.22\textwidth} L{0.22\textwidth}}
% \toprule
% \textbf{Subfield} & \textbf{Moving object \& condition} & \textbf{Sampler / operator} & \textbf{Dominant method families} & \textbf{Benchmarks / metrics} \\
% \midrule
% §\ref{subsec:p3_s8_4_1} Human motion & Articulated body; action / text / music / speech / scene & $G_\theta^{\text{mot}}$ (Eq.~\eqref{eq:p3_s8_4_1_objective}) & Regression or generative (GAN / VAE / diffusion); per-modality conditioning & HumanML3D, KIT-ML, AIST++, BEAT; FID, R-Precision, diversity, beat / contact \\
% §\ref{subsec:p3_s8_4_2} Animal motion & Multi-species body; text $+$ species $s$ & $G_\theta^{\text{amot}}$ (Eq.~\eqref{eq:p3_s8_4_2_objective}) & Human-prior transfer $+$ species-aware modulation; morphological consistency & AnimalML3D, Zoo-300K, AniMo4D; FID, MPJPE, biomechanical plausibility \\
% §\ref{subsec:p3_s8_4_3} Vehicle planning & Rigid agent state; HD map $+$ agent histories & $G_\theta^{\text{veh}}$, $K$-mode (Eqs.~\eqref{eq:p3_s8_4_3_objective},~\eqref{eq:p3_s8_4_3_multimodal}) & Vectorized encoders; goal-based and generative decoders; imitation planners & nuScenes, Argoverse 1/2, WOMD; minADE$_K$, minFDE$_K$, miss rate, soft-mAP \\
% §\ref{subsec:p3_s8_4_4} Motion editing & Existing motion $+$ edit region / control & $E_\theta$ (Eq.~\eqref{eq:p3_s8_4_4_objective}) & Masked infilling, spatial guidance, noise optimization, instruction editing & HumanML3D / BABEL protocols, MotionFix; editability $+$ preservation tuple \\
% \bottomrule
% \end{tabular}
% \end{table}

% The four subfields partition by \emph{what is held fixed}: human, animal, and vehicle generation fix nothing but the condition and sample a fresh trajectory, whereas editing fixes the source motion and changes only a region, standing to the other three exactly as image editing (§\ref{subsec:p3_s8_1_3}) stands to image generation. Underneath that partition the dominant failure mode is shared and physical: whether it surfaces as foot-skating and penetration in human and animal motion, kinematic infeasibility in vehicle trajectories, or balance disruption when one limb is edited, every subfield is limited by how well the generator respects the biomechanics and dynamics of the moving body, which is why physics- and contact-aware modeling recurs in all four Challenges discussions. Evaluation stays fragmented, each subfield reporting its own tuple (fidelity, diversity, condition consistency; calibrated multimodality; editability and preservation), and all are migrating toward learned or perceptual judges as scalar metrics prove weakly correlated with human assessment. Motion generation is thus the point in §\ref{sec:generation} where generative spatial modeling must answer to physical law as well as semantics, which leads directly to the dynamic 4D generation of §\ref{sec:p3_s8_5}.


% \newpage
\section{4D Spatial Generation}
\label{sec:p3_s8_5}




\FramedInitial{G}{enerative} modeling has moved outward from static images~\cite{rombach2022ldm,saharia2022photorealistic} to video~\cite{chen2024videocrafter2,chen2023videocrafter1,ho2022video,harvey2022flexible}, multi-view imagery~\cite{shi2023mvdream,liu2023zero,shi2023zero123++,zhao2024sg,zhao2024chase}, and static 3D assets~\cite{wang2023score,raj2023dreambooth3d}, yet each stops short of a scene that is at once three-dimensional and alive.
A 3D asset holds geometry but not motion; a video holds motion, but no view-consistent geometry one can navigate or simulate.
4D generation is where the two meet: the synthesis of dynamic 3D content whose geometry stays view-consistent as it moves and deforms over time~\cite{singer2023text,zheng2024unified,jiang2023consistent4d,ling2024align,wang2025moving}.
This is the representation an embodied system needs to rehearse how a scene will unfold, and it underwrites photorealistic digital humans~\cite{guo2023vid2avatar,zhao2024sg,zhao2024chase}, interactive AR/VR~\cite{zhao2024hfgs,miao2024pla4d}, and high-fidelity simulation for autonomous driving and embodied AI~\cite{gao2024magicdrive3d,zhao2024efficient}.

Two advances make this tractable.
Temporal 3D representations, deformation-aware fields over NeRFs~\cite{mildenhall2021nerf}, point clouds, and Gaussian splatting~\cite{kerbl20233d,zhao2024efficient}, supply a geometry that can bend and move without losing view consistency, while the extension of generative diffusion into this domain supplies the realism and controllability to drive it.
Together they open dynamic scene synthesis, 4D editing, and interactive agent simulation, at the cost of the difficulty this section returns to throughout: holding long-range temporal coherence while still granting flexible, user-driven control.

Formally, we cast 4D generation as the latent-generator framework of §\ref{subsec:p3_s8_1_1} unrolled along time over a shared 3D representation. Let $\mathcal{R}$ be the 3D-representation space of Eq.~\eqref{eq:p3_s8_2_1_objective} (mesh, NeRF, 3DGS, point cloud, and so on), and let $\kappa:\mathcal{C}\!\to\!\mathcal{C}'$ encode a conditioning signal $\mathbf{c}\in\mathcal{C}$ (text, an image, a video, a static 3D asset, or a combination thereof). 4D generation learns a conditional sampler
\begin{equation}
\label{eq:p3_s8_5_objective}
G_\theta^{\text{4D}}:\mathcal{C}'\!\to\!\mathcal{R}^{[0,T]},\qquad
\hat{\mathbf{r}}^{1:T}\sim p_\theta\!\left(\mathbf{r}^{1:T}\mid\kappa(\mathbf{c})\right),
\end{equation}
that emits a time-indexed sequence of 3D representations $\hat{\mathbf{r}}^{1:T}$ with view-consistent geometry and coherent motion. Two limiting cases anchor the framework: with a single time step ($T=1$) Eq.~\eqref{eq:p3_s8_5_objective} recovers the static 3D generator of Eq.~\eqref{eq:p3_s8_2_1_objective}, and rendering each frame $\mathbf{r}^t$ to a 2D view recovers the video generator of Eq.~\eqref{eq:p3_s8_3_objective}, so 4D generation is precisely the 3D generator of §\ref{sec:p3_s8_2} unrolled in time, or equivalently the video generator of §\ref{sec:p3_s8_3} lifted from per-frame pixels to a shared, renderable geometry.
The subsections are organized by the conditioning input (text, image, video, a static 3D asset, or several at once). Figure~\ref{fig:p3_s8_5_taxonomy} previews this organization together with two methodological strategies that recur across all five, per-scene score distillation and explicit ``generate-then-reconstruct'' supervision; feed-forward generators and rig-based animation of meshes appear alongside them in the subsections where they apply.

\begin{figure}[!ht]
    \centering
    \includegraphics[width=\linewidth]{figures/part3/section8/taxonomy_4d_spatial_generation.pdf}
    \caption{4D spatial generation by conditioning input, (a)--(e), with two representative strategies: (i) score distillation from a frozen diffusion prior (dashed: SDS feedback), and (ii) generating video or multiview targets, then reconstructing a dynamic 3D asset. The strategies are alternatives that may be combined; the chair sequence is illustrative.}
    \label{fig:p3_s8_5_taxonomy}
\end{figure}






\subsection{Text-to-4D Generation}
\label{subsec:p3_s8_5_1}
Text-to-4D generation synthesizes dynamic content directly from natural language, supporting text-driven asset creation and character animation~\cite{HongZPCYL22,chen2024ct4d}. Its progress has been propelled by large language models for semantic understanding and by video foundation models that supply motion priors.


\subsubsection{Task Definition}
\label{subsubsec:p3_s8_5_1_task}
Text-to-4D generation specializes the 4D sampler of Eq.~\eqref{eq:p3_s8_5_objective} to a textual condition $\mathbf{c}=\mathbf{x}$ encoded by the text encoder $\phi$ of Eq.~\eqref{eq:t2i_generator},
\begin{equation}
\label{eq:p3_s8_5_1_objective}
\hat{\mathbf{r}}^{1:T}\sim p_\theta\!\left(\mathbf{r}^{1:T}\mid\phi(\mathbf{x})\right).
\end{equation}
Because language fixes only semantics and leaves geometry and motion grossly under-determined, the sampler must hallucinate both a plausible 3D shape and its temporal evolution from the prompt alone, the source of the semantic-to-motion ambiguity examined in §\ref{subsubsec:p3_s8_5_1_challenges}.



\subsubsection{Methodology}
\label{subsubsec:p3_s8_5_1_method}
Two strategies recur throughout this section, implicit distillation and explicit supervision, implemented across various 4D representations (Figure~\ref{fig:p3_s8_5_taxonomy}).
The more common implicit strategy leverages Score Distillation Sampling (SDS) to indirectly optimize the geometry and motion of a 4D representation, using pre-trained diffusion models as frozen priors~\cite{singer2023text,zheng2024unified}. Formally, it applies a score-distillation update to per-scene parameters $\theta$~\cite{poole2022dreamfusion},
\begin{equation}
\label{eq:p3_s8_5_sds}
\nabla_\theta\mathcal{L}_{\mathrm{SDS}}
=\mathbb{E}_{s,\epsilon,\pi}\!\left[
 w(s)\left(\epsilon_\phi(\mathbf{z}_s;s,\kappa(\mathbf{c}))-\epsilon\right)
 \frac{\partial\mathbf{x}_\theta}{\partial\theta}\right],
\quad \mathbf{z}_s=\alpha_s\mathbf{x}_\theta+\sigma_s\epsilon,
\end{equation}
where $\mathbf{x}_\theta=\mathrm{Render}(\hat{\mathbf{r}}^{1:T}_\theta,\pi)$ is a differentiable rendering at camera $\pi$, $s$ indexes diffusion noise levels, and $\epsilon_\phi$ is a frozen 2D/video diffusion prior. The update differentiates through the renderer while treating the prior noise residual as the supplied gradient signal; this lifts the static score-distillation objective of §\ref{para:p3_s8_2_1_sds} to a temporal sequence.
In contrast, explicitly supervised methods initialize the 4D generation from text-to-3D or text-to-video outputs and use these generated assets as direct supervisory signals~\cite{zhao2024efficient,wu20244d,yang2024deformable,ling2024align,blattmann2023align,xu2024comp4d}. This ``generate-then-reconstruct'' strategy first synthesizes target frames $\hat{\mathbf{y}}$ with a fast generator $\mathcal{G}_{\text{2D}}$ and then fits the 4D asset to them,
\begin{equation}
\label{eq:p3_s8_5_explicit}
\hat{\mathbf{y}}=\mathcal{G}_{\text{2D}}(\mathbf{c}),\qquad
\hat{\mathbf{r}}^{1:T}=\arg\min_{\mathbf{r}^{1:T}}\sum_{\pi,t}\big\|\mathrm{Render}(\mathbf{r}^{t},\pi)-\hat{\mathbf{y}}^{t}_{\pi}\big\|^2,
\end{equation}
replacing score-distillation guidance with explicit pixel targets; reconstruction may itself require per-scene optimization, so a speed advantage depends on the implementation. The two strategies are realized across mesh, NeRF, and 3DGS representations.

\textbf{Mesh-based Methods.} Meshes offer explicit geometric control, which suits precise text-driven shape manipulation and animation.
CT4D~\cite{chen2024ct4d} implements a three-stage ``Generate-Refine-Animate'' pipeline in which a mesh is generated from text, its texture refined, and its vertices animated; TextMesh4D~\cite{dai2025textmesh4d} deforms a generated mesh through a per-face Jacobian deformation field under video-diffusion priors; and AvatarCLIP~\cite{HongZPCYL22} generates a static avatar zero-shot and animates it by registering motion to an SMPL skeleton.

\textbf{NeRF-based Methods.} NeRFs represent continuous geometry and complex temporal variation well.
MAV3D~\cite{singer2023text} introduced multi-stage SDS optimization that refines texture, geometry, and motion with priors from both T2I and T2V diffusion models; 4D-fy~\cite{bahmani20244d} blends 3D-aware, image, and video score signals in a hybrid SDS; and Dream-in-4D~\cite{zheng2024unified} first generates a static asset with a multiview diffusion model and then adds dynamics with a video diffusion model.

\textbf{Gaussian Splatting-based Methods.} 4D Gaussian Splatting (4DGS)~\cite{zhao2024hfgs,wu20244d} renders efficiently and differentiably, which speeds up text-driven dynamic scene synthesis. Early methods followed the two-stage implicit pipeline of Eq.~\eqref{eq:p3_s8_5_sds}: AYG~\cite{ling2024align} generates a static 3D Gaussian model from text via SDS and adds dynamics under a T2V prior, and Comp4D~\cite{xu2024comp4d} uses an LLM to parse a complex prompt into several dynamic objects and their trajectories for a compositional scene.
Later works moved to explicit supervision (Eq.~\eqref{eq:p3_s8_5_explicit}): PLA4D~\cite{miao2024pla4d} aligns the 4D Gaussians pixel by pixel with generated video targets, and 4Real~\cite{yu20244real} uses generated multiview video frames as the reference for photorealistic scene-level generation.
For digital humans, AvatarGO~\cite{cao2024avatargo} generates 4D human--object interactions zero-shot by animating a text-to-3D human with the SMPL-X skinning function.


\subsubsection{Benchmarks}
\label{subsubsec:p3_s8_5_1_benchmarks}
Text-to-4D, like the other subfields of this section, has no standardized benchmark, because large-scale dynamic 3D data with text descriptions is hard to capture and annotate. Training and evaluation instead adapt static 3D repositories (Table~\ref{tab:4d_datasets}): Diffusion4D~\cite{liang2024diffusion4d} filtered Objaverse and Objaverse-XL~\cite{deitke2023objaverse,deitke2023objaversexl} down to 54K high-quality animated assets, and SV4D~\cite{xie2024sv4d} built ObjaverseDy by adaptive sampling and viewpoint adjustment. These are synthetic renderings and lack the multimodal annotations, text in particular, needed to train and evaluate text-conditioned models robustly.

\begin{table}[!ht]
\centering
\fontsize{8}{10}\selectfont
\setlength{\tabcolsep}{1.0mm}
\caption{Datasets used for 4D generation across the subfields of §\ref{sec:p3_s8_5}.}
\label{tab:4d_datasets}
\begin{tabular}{L{0.20\textwidth} L{0.24\textwidth} C{0.20\textwidth} L{0.26\textwidth}}
\toprule
\textbf{Dataset} & \textbf{Content} & \textbf{Scale} & \textbf{Used for} \\
\midrule
Objaverse-XL~\cite{deitke2023objaversexl} & Static and animated 3D objects & 10M+ objects & Source pool for dynamic subsets \\
Diffusion4D~\cite{liang2024diffusion4d} & Animated assets curated from Objaverse / Objaverse-XL & 54K assets & Training image- / video-to-4D \\
ObjaverseDy~\cite{xie2024sv4d} & Animated Objaverse objects, re-rendered & - & Training and evaluating video-to-4D \\
Consistent4D~\cite{jiang2023consistent4d} & Synthetic and real monocular videos & - & Evaluating video-to-4D \\
\bottomrule
\end{tabular}
\end{table}

\paragraph{Evaluation Metrics.} Evaluation across 4D generation combines user studies with metrics borrowed from image and video generation, and the subsections that follow use this same set. Semantic alignment with a prompt is measured by CLIP Score~\cite{radford2021learning} between rendered views and the text; fidelity by PSNR~\cite{hore2010image}, SSIM~\cite{wang2004image}, and LPIPS~\cite{zhang2018unreasonable} where a reference view exists; temporal quality and motion by the Fréchet Video Distance (FVD)~\cite{unterthiner2018fvd}; and overall plausibility and aesthetics by user studies. None of these scores the two properties that make content 4D, cross-view consistency and long-range spatio-temporal coherence.


\subsubsection{Challenges and Trends}
\label{subsubsec:p3_s8_5_1_challenges}
The challenge specific to text-to-4D is semantic-to-motion ambiguity: a prompt such as ``a dancing dog'' admits endlessly many movements, styles, and trajectories, and methods that lean on generic T2V priors struggle to produce diverse, fine-grained motion that matches a nuanced description.
One response uses LLMs for motion planning, as Comp4D~\cite{xu2024comp4d} does when it decomposes a prompt into a structured spatio-temporal plan for several interacting objects; another targets human--object interaction~\cite{cao2024avatargo}, moving toward story-driven dynamic scenes.

% ======================================================================
% ======================================================================


\subsection{Image-to-4D Generation}
\label{subsec:p3_s8_5_2}
Image-to-4D generation synthesizes dynamic content from a single static image, combining the spatial information in the image with temporal priors learned by video models~\cite{yang2024diffusion,sun2024eg4d}. No large image--4D paired dataset exists, which has ruled out direct end-to-end training and pushed the field toward optimization and toward generated supervision.

\subsubsection{Task Definition}
\label{subsubsec:p3_s8_5_2_task}
Image-to-4D generation specializes the 4D sampler of Eq.~\eqref{eq:p3_s8_5_objective} to a single-image condition $\mathbf{c}=I$,
\begin{equation}
\label{eq:p3_s8_5_2_objective}
\hat{\mathbf{r}}^{1:T}\sim p_\theta\!\left(\mathbf{r}^{1:T}\mid\kappa(I)\right),
\end{equation}
where the source view at $t=0$ is fixed by $I$ while every unseen viewpoint and time step must be inferred.
Both the unseen 3D geometry and the entirely absent temporal dynamics have to be inferred from one 2D view, which requires combining the spatial cues of the input with temporal priors from large video models~\cite{ren2023dreamgaussian4d,zhao2023animate124,sun2024eg4d}.



\subsubsection{Methodology}
\label{subsubsec:p3_s8_5_2_method}
Both strategies of §\ref{subsubsec:p3_s8_5_1_method} apply, and most methods use 4D Gaussian Splatting as the representation~\cite{ren2023dreamgaussian4d,yang2024diffusion}.

\textbf{Implicit Optimization with Diffusion Priors.} These methods instantiate Eq.~\eqref{eq:p3_s8_5_sds}, typically combining a multiview diffusion model for geometric consistency with a video diffusion model for temporal coherence. Animate124~\cite{zhao2023animate124} first generates a static NeRF from the image and then refines its dynamics with a video-diffusion SDS loss, and Diffusion\textsuperscript{2}~\cite{yang2024diffusion} reconciles the scores of heterogeneous priors (multiview and video) through Variance-Reducing Sampling, improving both multiview and temporal consistency of the resulting 4D Gaussians.

\textbf{Explicit Supervision with Generated Data.} The slow, iterative nature of SDS motivates the two-stage pipeline of Eq.~\eqref{eq:p3_s8_5_explicit}, in which an image-to-video model first generates an intermediate video that then supervises the 4D reconstruction directly. DreamGaussian4D (DG4D)~\cite{ren2023dreamgaussian4d} generates a static 3D Gaussian model and learns a deformation field to animate it, using the generated video for texture refinement and dynamic supervision, and EG4D~\cite{sun2024eg4d} avoids SDS entirely by chaining image-to-video and image-to-multiview models into coherent image sequences that guide a fast reconstruction.

\textbf{Unified and Tuning-free Generation.} GenXD~\cite{zhao2024genxd} trains one model on both 3D and 4D data, using multiview-temporal modules that separate camera motion from object motion, and Free4D~\cite{liu2025free4d} builds a 4D scene from a single image without fine-tuning any model, distilling pretrained image-to-video and multiview priors into a consistent 4D representation.

\textbf{Specialized Methods for Human Generation.} Dynamic human avatars form a sub-domain in which strong body priors raise fidelity. Human4DiT~\cite{shao2024360} uses a hierarchical 4D transformer that integrates human-specific features, camera parameters, and temporal signals to generate 360-degree human animations from a single image, and Disco4D~\cite{pang2025disco4d} models the body with SMPL-X while representing clothing separately with Gaussian splats, allowing detailed animation of clothed humans.



\subsubsection{Benchmarks}
\label{subsubsec:p3_s8_5_2_benchmarks}
With no image--4D pairs, methods are evaluated on images from Objaverse~\cite{deitke2023objaverse} or on in-the-wild photographs chosen by the authors, which makes comparison across papers unreliable. The metric set is that of §\ref{subsubsec:p3_s8_5_1_benchmarks}, applied to the source view and to pseudo-ground-truth novel views where available.

\subsubsection{Challenges and Trends}
\label{subsubsec:p3_s8_5_2_challenges}
Inferring a complete shape and a plausible motion from one viewpoint is extremely ill-posed, which shows up as geometric inconsistency in unseen views (\eg, the ``Janus problem'') and as unnatural or repetitive motion, since every dynamic detail must be hallucinated.
Domain-specific priors are the main remedy: Human4DiT~\cite{shao2024360} and Disco4D~\cite{pang2025disco4d} regularize geometry and motion with parametric body models and reach much higher fidelity, at the cost of generality.


% ======================================================================
% ======================================================================


\subsection{Video-to-4D Generation}
\label{subsec:p3_s8_5_3}
Video-to-4D generation lifts a single-view video into a dynamic 4D representation that can be viewed from new angles with consistent geometry and motion, turning ordinary video into interactive 4D content. Annotated video-to-4D datasets are hard to curate, so current methods build on pretrained video diffusion priors, either as multi-view supervision or through score distillation.



\subsubsection{Task Definition}
\label{subsubsec:p3_s8_5_3_task}
Video-to-4D generation specializes the 4D sampler of Eq.~\eqref{eq:p3_s8_5_objective} to a monocular video condition $\mathbf{c}=\mathbf{v}^{1:T}$,
\begin{equation}
\label{eq:p3_s8_5_3_objective}
\hat{\mathbf{r}}^{1:T}\sim p_\theta\!\left(\mathbf{r}^{1:T}\mid\kappa(\mathbf{v}^{1:T})\right),
\end{equation}
so that re-rendering $\hat{\mathbf{r}}^{1:T}$ from the input viewpoint reproduces $\mathbf{v}^{1:T}$ while novel views remain geometrically consistent.
Unlike image-to-4D (Eq.~\eqref{eq:p3_s8_5_2_objective}), the temporal dynamics are now \emph{observed} and the residual ambiguity is purely geometric: the unseen off-axis structure.

\subsubsection{Methodology}
\label{subsubsec:p3_s8_5_3_method}
Most methods densify the monocular input into multi-view supervision and reconstruct against it (Eq.~\eqref{eq:p3_s8_5_explicit}), some add score distillation (Eq.~\eqref{eq:p3_s8_5_sds}), and the newest predict the 4D asset feed-forward.

\textbf{Mesh-based Methods.} Meshes are interpretable and efficient to animate. DreamMesh4D~\cite{li2024dreammesh4d} obtains a coarse mesh and drives it with a sparse deformation graph and linear blend skinning (LBS), reaching strong spatio-temporal consistency more cheaply than implicit volumetric methods, and V2M4~\cite{chen2025v2m4} reconstructs an animated, textured mesh asset directly from a monocular video.

\textbf{NeRF-based Methods.} NeRF-based methods model the continuous volumetric content of a video well, and they overcome the single view by generating data. Consistent4D~\cite{jiang2023consistent4d} densifies the monocular input through temporal and multiview interpolation before optimizing a NeRF with SDS, 4Diffusion~\cite{zhang20244diffusion} fine-tunes a video diffusion model with a learnable motion module on synthetic multiview data so that it produces consistent multiview video to guide generation, and SV4D~\cite{xie2024sv4d} completes missing views and frames with a diffusion model and uses the resulting multiview video to supervise the reconstruction.

\textbf{Gaussian Splatting-based Methods.} 4DGen~\cite{yin20234dgen} grounds 4D Gaussian generation in the input video with spatial-temporal pseudo labels; Efficient4D~\cite{pan2024efficient4d} converts the video into a multiview image matrix with a multiview-consistent generator (SyncDreamer) for rapid reconstruction; SC4D~\cite{wu2024sc4d} first fits sparse control Gaussians for coarse shape and motion and then drives dense Gaussians through LBS; and L4GM~\cite{ren2024l4gm} predicts the 4D Gaussians of a whole clip in a single feed-forward pass, trained on rendered animated objects.


\subsubsection{Benchmarks}
\label{subsubsec:p3_s8_5_3_benchmarks}
Video-to-4D has the one dedicated evaluation set in this section, released with Consistent4D~\cite{jiang2023consistent4d}: synthetic and real monocular videos with training and test splits. The metrics are those of §\ref{subsubsec:p3_s8_5_1_benchmarks}, with reconstruction fidelity measured at the input view and novel-view quality measured where synthetic ground-truth multiview video exists.

\subsubsection{Challenges and Trends}
\label{subsubsec:p3_s8_5_3_challenges}
The video constrains motion but not the unseen 3D structure, which leads to geometric distortion, depth ambiguity, and a ``painted-on'' look in novel views.
Explicit geometric structure is the main remedy: DreamMesh4D~\cite{li2024dreammesh4d} and SC4D~\cite{wu2024sc4d} constrain the reconstruction with meshes or sparse control points under LBS, which improves both coherence and efficiency, and combining such structure with the appearance modeling of neural representations is the obvious next step.

% ======================================================================
% ======================================================================

\subsection{3D-to-4D Generation}
\label{subsec:p3_s8_5_4}
3D-to-4D generation animates a static 3D model, so that the geometry is given and only motion and deformation must be synthesized, which makes the large stock of existing 3D assets usable for simulation, virtual reality, and character animation.
Unconditional generators of dynamic neural fields, such as HyperDiffusion~\cite{erkocc2023hyperdiffusion}, which runs diffusion directly on the weights of neural fields, sample 4D content without a given asset; 3D-to-4D instead starts from one, either through explicit mesh deformation and rigging or through neural and physics-based deformation models.

\subsubsection{Task Definition}
\label{subsubsec:p3_s8_5_4_task}
Here the condition is itself a 3D representation $\mathbf{c}=\mathbf{r}_0\in\mathcal{R}$, so only its temporal evolution must be generated. We therefore specialize Eq.~\eqref{eq:p3_s8_5_objective} to a \emph{deformation operator} that takes the static asset and an optional auxiliary driving signal $\mathbf{c}_{\text{aux}}$ (text, audio, or a motion sketch) to a sequence,
\begin{equation}
\label{eq:p3_s8_5_4_objective}
\hat{\mathbf{r}}^{1:T}=\mathcal{D}_\theta\!\left(\mathbf{r}_0,\,\kappa(\mathbf{c}_{\text{aux}})\right),\qquad \hat{\mathbf{r}}^{0}=\mathbf{r}_0,
\end{equation}
with the first frame pinned to the input asset. This places 3D-to-4D in the same family as the editing operators of §\ref{subsec:p3_s8_2_5} and §\ref{subsec:p3_s8_4_4}: it transforms an existing representation rather than synthesizing one from scratch.

\subsubsection{Methodology}
\label{subsubsec:p3_s8_5_4_method}
Methods realize the deformation operator $\mathcal{D}_\theta$ of Eq.~\eqref{eq:p3_s8_5_4_objective} either with explicit mesh structure, increasingly through learned rigging, or with neural and physics-informed deformation fields.

\textbf{Mesh-based Methods.} Meshes give explicit structure and precise control over deformation. Motion3DGAN~\cite{otberdout2022sparse} generates 3D facial animation from a static face mesh by first predicting the displacement of facial landmarks for a target expression and then deforming the full mesh, with a GAN keeping the sequence temporally coherent. A recent line makes arbitrary meshes animatable by predicting the rig itself: MagicArticulate~\cite{song2025magicarticulate} and RigAnything~\cite{liu2025riganything} predict skeletons and skinning weights feed-forward, Puppeteer~\cite{song2025puppeteer} rigs a mesh and then animates it by fitting the rig to videos generated from renderings of the mesh, and AnimateAnyMesh~\cite{wu2025animateanymesh} generates text-driven mesh animation in a single forward pass with a rectified-flow model over a dynamic-mesh latent.

\textbf{NeRF-based Methods.} ElastoGen~\cite{feng2024elastogen} generates physically grounded elastodynamic behavior: it rasterizes a static 3D model and then predicts and refines deformation fields iteratively with recurrent and 3D convolutional modules, so that the motion obeys physical principles.

\subsubsection{Benchmarks}
\label{subsubsec:p3_s8_5_4_benchmarks}
No benchmark with ground-truth dynamic sequences exists; methods are evaluated on a few individual assets, often faces or generic objects, mostly by user study and with the metrics of §\ref{subsubsec:p3_s8_5_1_benchmarks} where a reference animation is available. Physical plausibility, specific to this subfield, can be measured for physics-informed methods such as ElastoGen~\cite{feng2024elastogen} as deviation from a high-fidelity simulator or as violation of conserved quantities such as energy and momentum.

\subsubsection{Challenges and Trends}
\label{subsubsec:p3_s8_5_4_challenges}
Even with the geometry given, the space of plausible animations of a static model is unbounded, so without guidance the task is unconstrained, and current methods specialize to narrow categories such as faces.
Physical priors constrain it (ElastoGen~\cite{feng2024elastogen} grounds motion in elastodynamics), learned rigging makes arbitrary assets animatable, and pairing the asset with text, audio, or motion sketches turns open-ended generation into controllable animation.

% ======================================================================
% ======================================================================


\subsection{Multi-conditional 4D Generation}
\label{subsec:p3_s8_5_5}
Multi-conditional 4D generation steers dynamic scene synthesis with several signals at once: semantic conditions (text), appearance conditions (images, video), structural conditions (skeletal poses, explicit trajectories), and rule-based priors (differentiable physics). The central design problem is to fuse them so that each is respected without one overpowering the others.

\subsubsection{Task Definition}
\label{subsubsec:p3_s8_5_5_task}
We generalize the single-condition cases of Eqs.~\eqref{eq:p3_s8_5_1_objective}--\eqref{eq:p3_s8_5_4_objective} to a set of $K$ conditions, each mapped by its own encoder $\kappa_k$ and fused into the shared embedding:
\begin{equation}
\label{eq:p3_s8_5_5_objective}
\hat{\mathbf{r}}^{1:T}\sim p_\theta\!\left(\mathbf{r}^{1:T}\mid \kappa_1(\mathbf{c}_1),\dots,\kappa_K(\mathbf{c}_K)\right),
\end{equation}
which recovers a single-modality sampler when $K=1$.
Users can then specify not only what is generated (semantics) but how it behaves (motion path, physical behavior, artistic style), which complex applications such as character animation and driving simulation require.

\subsubsection{Methodology}
\label{subsubsec:p3_s8_5_5_method}
Methods fuse the encoders $\kappa_k$ of Eq.~\eqref{eq:p3_s8_5_5_objective}, usually combining score distillation (Eq.~\eqref{eq:p3_s8_5_sds}) with explicit supervision (Eq.~\eqref{eq:p3_s8_5_explicit}) tailored to each condition.

\textbf{Motion and Trajectory Control.} TC4D~\cite{bahmani2024tc4d} lets users define a spline trajectory that guides both the global rigid transformation and the local deformation of an object, with text supplying the semantics. For human animation, skeletons are the natural condition: Beyond Skeletons~\cite{yang2024beyond} integrates text and image conditions with the skeleton to refine motion, and STAR~\cite{chai2024star} uses skeleton-aware SDS and motion retargeting so that avatar movement is both visually coherent and anatomically accurate.

\textbf{Physics-Informed Control.} Sync4D~\cite{fu2024sync4d} and Phy124~\cite{lin2024phy124} integrate differentiable Material Point Method (MPM) simulation, so that the generated content follows realistic dynamics such as elasticity and collisions under video or text guidance.

\textbf{User-Centric and Artistic Control.} Generative Rendering~\cite{cai2024generative} uses a pre-existing dynamic mesh to guide a pretrained text-to-image model, giving precise control over the style and detail of the rendered output, and sketch-based interfaces~\cite{yang2024sketch} let users draw the dynamics they want.

\textbf{Complex, Application-Specific Control.} MagicDrive3D~\cite{gao2024magicdrive3d} generates dynamic driving scenes from 2D bounding boxes, Bird's-Eye-View (BEV) maps, agent trajectories, and text together, producing tightly controlled traffic scenarios.


\subsubsection{Benchmarks}
\label{subsubsec:p3_s8_5_5_benchmarks}
No benchmark exists; building one would require 4D assets annotated with several aligned conditions (\eg, a clip with a description, a skeletal motion sequence, and a trajectory map). Text adherence is scored with the CLIP Score of §\ref{subsubsec:p3_s8_5_1_benchmarks}, while adherence to sketches or trajectories relies on task-specific errors (\eg, distance from the specified trajectory) or on user studies, which remain the only holistic judge of how well a combination of inputs was followed.

\subsubsection{Challenges and Trends}
\label{subsubsec:p3_s8_5_5_challenges}
Semantic guidance from text, geometric constraints from skeletons, and rule-based priors from physics can contradict one another, and keeping each respected without one overpowering the rest while preserving fidelity is the open problem.
Domain-specific stacks such as MagicDrive3D~\cite{gao2024magicdrive3d} sidestep it by fixing the set of inputs for one application, and interactive interfaces such as sketches~\cite{yang2024sketch} point toward iterative creation in which several control channels are adjusted in turn.

% ======================================================================
% ======================================================================


\subsection{Summary, Overarching Challenges and Future Directions}
\label{subsec:p3_s8_5_summary}
Table~\ref{tab:p3_s8_5_subfield_comparison} compares the five subfields as specializations of the 4D sampler of Eq.~\eqref{eq:p3_s8_5_objective}: image- and video-to-4D progressively replace hallucinated structure with observed pixels, 3D-to-4D is given the geometry and generates only motion, and multi-conditional generation fuses several signals at once.

\begin{table}[!ht]
\centering
\fontsize{8}{10}\selectfont
\setlength{\tabcolsep}{1.0mm}
\caption{The five subfields of 4D spatial generation; each sampler or operator specializes Eq.~\eqref{eq:p3_s8_5_objective}.}
\label{tab:p3_s8_5_subfield_comparison}
\begin{tabular}{L{0.13\textwidth} L{0.15\textwidth} L{0.17\textwidth} L{0.25\textwidth} L{0.22\textwidth}}
\toprule
\textbf{Subfield} & \textbf{Conditioning input} & \textbf{Sampler / operator} & \textbf{Dominant method families} & \textbf{Datasets / metrics} \\
\midrule
§\ref{subsec:p3_s8_5_1} Text-to-4D & Text prompt $\mathbf{x}$ & $G_\theta^{\text{4D}}$, $\kappa=\phi$ (Eq.~\eqref{eq:p3_s8_5_1_objective}) & SDS distillation or generate-then-reconstruct; mesh / NeRF / 3DGS; LLM motion planning & Objaverse-derived; CLIP, FVD, PSNR/SSIM/LPIPS, user study \\
§\ref{subsec:p3_s8_5_2} Image-to-4D & Single image $I$ & $G_\theta^{\text{4D}}$, $\kappa(I)$ (Eq.~\eqref{eq:p3_s8_5_2_objective}) & Diffusion-prior SDS; explicit video supervision; unified / tuning-free; human-specific priors & Objaverse / in-the-wild; PSNR/SSIM/LPIPS, FVD, user study \\
§\ref{subsec:p3_s8_5_3} Video-to-4D & Monocular video $\mathbf{v}^{1:T}$ & $G_\theta^{\text{4D}}$, $\kappa(\mathbf{v}^{1:T})$ (Eq.~\eqref{eq:p3_s8_5_3_objective}) & Video-diffusion densification; explicit geometric constraints (mesh / LBS); feed-forward & Consistent4D; PSNR/SSIM/LPIPS, FVD, user study \\
§\ref{subsec:p3_s8_5_4} 3D-to-4D & Static asset $\mathbf{r}_0$ & $\mathcal{D}_\theta$ deformation (Eq.~\eqref{eq:p3_s8_5_4_objective}) & Mesh deformation and learned rigging; physics-informed deformation & Per-asset (faces / objects); FVD, physical plausibility, user study \\
§\ref{subsec:p3_s8_5_5} Multi-conditional & $K$ fused signals & $G_\theta^{\text{4D}}$, $\{\kappa_k\}$ (Eq.~\eqref{eq:p3_s8_5_5_objective}) & Trajectory / skeleton / physics / sketch control; application-specific stacks & No standard suite; CLIP, PSNR/SSIM/LPIPS, FVD, user study \\
\bottomrule
\end{tabular}
\end{table}

\paragraph{Data Scarcity and Generative Priors.}
Large, annotated 4D datasets are the field's main missing ingredient: real dynamic 3D data with multi-view consistency and multimodal annotation (text, motion capture) is expensive and technically hard to collect, which has largely ruled out end-to-end training. The field has therefore converged on distilling pretrained text-to-image, text-to-video, and text-to-3D models, the de facto way to guide 4D synthesis.

\paragraph{From Score Distillation to Explicit Supervision.}
Score distillation~\cite{poole2022dreamfusion} converges slowly, is costly and unstable, and produces multi-faced geometry and flickering textures. The generate-then-reconstruct pipeline of DreamGaussian4D~\cite{ren2023dreamgaussian4d} and EG4D~\cite{sun2024eg4d}, which first generates target videos and then reconstructs against them, is replacing per-prompt SDS, and feed-forward predictors such as L4GM~\cite{ren2024l4gm} and AnimateAnyMesh~\cite{wu2025animateanymesh} take the next step by removing per-scene optimization altogether.

\paragraph{4D Gaussian Splatting as the Default Representation.}
NeRF-based works rendered well but trained and rendered slowly. 3D Gaussian Splatting~\cite{kerbl20233d} and its temporal extension offer real-time, high-fidelity rendering with efficient gradient-based optimization, and their explicit, point-based form is easier to edit and to couple with physical simulation, so 4DGS has become the representation of choice for current methods~\cite{wu20244d, ren2023dreamgaussian4d, ling2024align}.

\paragraph{Consistency, Control, and Physical Plausibility.}
Keeping geometry, appearance, and motion coherent across all viewpoints and over long horizons remains the central difficulty, and control is moving beyond single prompts toward fine-grained manipulation of motion, appearance, and interaction. Physical plausibility is the newest goal: integrating differentiable simulators, as PhysSplat~\cite{zhao2024efficient} and Phy124~\cite{lin2024phy124} do, makes generated dynamics more realistic and predictable. Progress on all three is hard to measure while evaluation relies on small, author-curated test sets and on image and video metrics that ignore cross-view and long-range consistency; 4D-aware benchmarks and metrics are a prerequisite for the rest.

% \newpage


\section{World Modeling}
\label{sec:worldmodeling}
\label{sec:p3_s8_6}

\FramedInitial{D}{espite} its growing popularity across AI research, the notion of a ``world model'' does not yet enjoy a unified or universally agreed-upon definition.
Its conceptual roots trace back to the psychologist Kenneth Craik, who argued in his 1943 monograph \emph{The Nature of Explanation} that the mind ``carries a `small-scale model' of external reality and of its own possible actions within its head''.
In its original form, a world model referred to an individual's internal representation of external reality; contemporary AI has generalized the concept far beyond mental representations to structured models of physical and spatial environments.

Fei-Fei Li has framed the gap between current AI systems and physical understanding in these terms~\cite{Feifei2025spatial}: modern LLMs resemble ``wordsmiths in the dark'', fluent in symbols but divorced from the physical world that gives them meaning, and capable AI must instead understand, act within, and reason about three-dimensional space. On this view world models are new-generation \emph{generative} models that simulate, predict, and imagine the world's dynamics rather than only describe it, which ties world modeling directly to spatial intelligence.



Complementing Li's generative perspective, another influential formulation comes from Yann~LeCun \cite{lecun2022path}. He describes a world model as a system capable of predicting how the world will evolve given its current state and an action.
Formally, it consumes the current state $S(t)$ and an action $A(t)$, and produces a prediction of the next state $S'(t+1)$; in many cases, it also integrates external sensory observations $O(t)$ to update its internal state.
This predictive relation can be written compactly as
\begin{equation}
\label{eq:p3_s8_6_transition}
\hat{S}(t+1) = W_\theta\big(S(t),\, A(t),\, O(t)\big),
\end{equation}
where $W_\theta$ denotes the learned world model parameterized by $\theta$.
Read against the conditional-generation backbone of this chapter (Eq.~\eqref{eq:t2i_generator}), $W_\theta$ occupies the position of the generator, but its decoder is a state-transition predictor rather than a pixel or geometry renderer; this is the sense in which world modeling extends Chapter~\ref{sec:generation} from reconstruction toward prediction.
What exactly constitutes a ``state'' remains open: some work defines it as a compact latent representation in which the world model operates at an abstract semantic level, much as human cognition forecasts events without simulating every physical detail.




Although no single definition of world modeling has been universally accepted, a survey of the literature~\cite{fung2025embodied,li2025comprehensive,ding2025understanding,kong20253d} reveals three recurring paradigms, three realizations of Eq.~\eqref{eq:p3_s8_6_transition} that differ in what they predict.
\textbf{Representation world models} learn abstract latent spaces that capture the state and dynamics of the environment.
\textbf{Generative world models} explicitly simulate or generate future states of the world through data-driven or hybrid mechanisms.
\textbf{Mental world models} aim to infer and represent the mental states of humans or agents, enabling social reasoning and Theory-of-Mind (ToM) capabilities.
Figure~\ref{fig:world-modeling} contrasts them side by side.


 % \cite{}






\begin{figure*}[!t]
    \centering
    \includegraphics[width=\linewidth]{figures/web/world_modeling_tikz.pdf}
    \caption{Three paradigms of world modeling, all instances of the transition $\hat{S}(t{+}1)=W_\theta(S(t),A(t),O(t))$ of Eq.~\eqref{eq:p3_s8_6_transition}, differing in what they predict: (a) a latent state, (b) future observations, and (c) another agent's mental state and next action.}
    \label{fig:world-modeling}
\end{figure*}




\subsection{Representation World Model}
\label{subsec:p3_s8_6_rep}




\subsubsection{Definition}

A representation world model learns an internal latent description of the environment by compressing multimodal observations into structured, predictive embeddings, prioritizing the extraction of physical and semantic regularities over the generation of photorealistic pixels.
Much as LLMs learn abstract linguistic representations through self-supervised pretraining, representation world models are trained on ``world data'' to infer the principles that govern physical and semantic dynamics; the objective is to predict how the world evolves at an abstract level, not to reconstruct pixels.
The latent state is meant to encode objects and their properties (shape, size, appearance, material), their spatial relations (topology, distance, configuration), the dynamics that govern how the scene evolves, and the causal dependence of outcomes on actions. With such a state, an agent can predict over long horizons, reason, and plan without operating in high-dimensional sensory space.





\subsubsection{Multimodal Perception}

A world model draws on several sensory streams, whose encoding is the subject of Chapter~\ref{chapter:3:sec:6}. Vision dominates: images and especially video carry objects, geometry, interaction, and temporal dynamics, and modern systems rely on high-capacity encoders trained with contrastive vision--language objectives for semantic and geometric priors. Audio adds direct evidence about events and the physical state of the surroundings, which language conveys only symbolically. Touch is indispensable where vision is unreliable, occluded, or unavailable (grasping inside a bag, fine manipulation, tool use), yet it remains underexplored because tactile sensors are limited and tactile data are hard to collect at scale.




\subsubsection{World Representation Modeling}

\paragraph{Vision-oriented Modeling.}
A major line of representation world modeling builds latent spaces from visual observations: video frames are encoded into a compact latent space by self-supervised learning, and a transition model predicts future latent states from past latents and actions, which captures the dynamics of the visual world more efficiently than pixel-space prediction.
Concretely, an encoder $E_\phi$ maps observations to latent states and a predictor $P_\theta$ is trained to minimize a prediction loss defined entirely in latent space,
\begin{equation}
\label{eq:p3_s8_6_rep}
\mathcal{L}_{\text{rep}}(\theta,\phi) = \mathbb{E}\!\left[\,\big\lVert P_\theta\big(E_\phi(O_{\le t}),\, A(t)\big) - \operatorname{sg}\!\big(E_{\bar{\phi}}(O_{t+1})\big)\big\rVert^{2}\,\right],
\end{equation}
where $\operatorname{sg}(\cdot)$ denotes the stop-gradient operator and $E_{\bar{\phi}}$ a slowly updated target encoder.
This objective instantiates the transition of Eq.~\eqref{eq:p3_s8_6_transition} with the predicted state $\hat{S}(t+1)$ taken to be an abstract embedding rather than raw pixels (Figure~\ref{fig:world-modeling}(a)), which is precisely what allows representation world models to capture dynamics efficiently while discarding perceptually salient yet predictively irrelevant detail.



The \emph{World Models} architecture~\cite{ha2018world} decomposes the problem into a VAE-based vision model, an RNN-based memory model, and a lightweight controller; the vision and memory components learn latent dynamics from raw pixels without supervision, and the controller then learns its behavior inside the ``dreamed'' environment, which greatly simplifies policy learning.
The Dreamer agents carry this idea into reinforcement learning at scale: Dreamer~\cite{hafner2019dreamer} learns a latent dynamics model from pixels and trains an actor and critic entirely on imagined latent trajectories, and DreamerV3~\cite{hafner2023mastering} applies a single configuration across a wide range of domains. MuZero~\cite{schrittwieser2020mastering} learns a model that predicts only the quantities planning needs (reward, value, and policy) and plans with it by tree search, mastering Atari, Go, chess, and shogi without being given the rules.


Yann~LeCun's autonomous intelligence framework proposes the Joint Embedding Predictive Architecture (JEPA)~\cite{lecun2022path} as the core mechanism for constructing world models.
JEPA learns latent representations and predictions through self-supervised objectives rather than pixel-level reconstruction, and Eq.~\eqref{eq:p3_s8_6_rep} is its generic form.
I-JEPA~\cite{assran2023jepa} first instantiated it on images, predicting the embeddings of masked target regions from a visible context block, and V-JEPA~\cite{bardes2024vjepa} carried the same feature-prediction objective to video, predicting masked spatio-temporal embeddings \emph{in latent space} with no pixel reconstruction, text, or negative examples, and so avoiding the cost and ambiguity of pixel prediction.
V-JEPA~2~\cite{assran2025vjepa} turns this encoder into a world model. Pretrained on more than a million hours of internet video, it is extended with an action-conditioned predictor, V-JEPA~2-AC, post-trained on less than 62 hours of unlabeled robot video; deployed zero-shot on Franka arms, it picks and places objects by planning actions whose predicted embeddings reach a goal image.

The stop-gradient and slowly updated target encoder in Eq.~\eqref{eq:p3_s8_6_rep} exist only to keep the representation from collapsing, and the most recent JEPA work removes them. LeJEPA~\cite{balestriero2025lejepa} shows that an isotropic Gaussian is the optimal embedding distribution for downstream prediction and enforces it with a single sketched regularizer (SIGReg), leaving one trade-off hyperparameter and no teacher--student machinery. LeWorldModel~\cite{maes2026leworldmodel} uses the same regularizer to train an action-conditioned JEPA world model end to end from raw pixels with only two loss terms; at about 15M parameters it trains on a single GPU, plans up to 48 times faster than world models built on foundation-model encoders, and its latents encode physical quantities well enough to flag physically implausible events.
The principle has also reached language: VL-JEPA~\cite{chen2025vl} predicts continuous embeddings of the target text instead of generating tokens, and LLM-JEPA~\cite{huang2026llmjepa} adds an embedding-prediction objective to the next-token training of language models.


DINO-World~\cite{baldassarre2025back} instead freezes DINOv2 as a universal visual encoder, whose features are already semantic, geometric, and object-centric, learns a latent-space predictor from large-scale unlabeled video, and adds action-conditioned prediction to model how interventions change future visual states.



\paragraph{Hybrid World Modeling.}
Another important direction expands representation learning beyond vision by incorporating additional modalities such as LiDAR, camera trajectories, or multi-view geometric signals.
These hybrid models aim to overcome the limitations of purely visual latent spaces by injecting explicit spatial structure, depth cues, and geometric consistency constraints.
For instance, {MUVO}~\cite{bogdoll2025muvo} is a multimodal generative world model that learns jointly from synchronized camera and LiDAR streams.
By integrating LiDAR-derived depth priors into the latent world representation, MUVO achieves more reliable 3D scene prediction, improved object motion modeling, and long-horizon forecasting in autonomous driving environments.
{UniWorld}~\cite{min2023uniworld} introduces a unified pretrained architecture that merges 3D object detection, motion prediction, and scene representation learning.
Through joint training of BEV features, temporal dynamics, and multisensor context, UniWorld produces a structured latent state that simultaneously supports perception and trajectory forecasting.
{Aether}~\cite{zhu2025aether} proposes a geometry-aware unified world model where camera trajectories are treated as explicit action signals.
This allows the model to learn action-conditioned scene transitions and perform predictive visual planning.
By binding egomotion with geometric consistency constraints, Aether generates spatially coherent future predictions and supports counterfactual simulation within a 3D-aware latent space.




\paragraph{Language-oriented World Modeling.}
A final category examines language as a central medium for latent world representation.
This raises an important question: \emph{can LLMs themselves be viewed as world models?}
From a formal perspective, LLMs indeed share structural similarities with classical world models.
They maintain internal states, treat textual context as observations, accept prompts as actions that alter these states, and generate predictions that correspond to inferred future states.
Models such as Dynalang~\cite{lin2023learning} demonstrate that natural language can serve as a universal intermediate representation for encoding environment states, predicting transitions, and expressing causal relations.
LLMs also exhibit emergent knowledge of commonsense physics, intuitive causality, and human behavior derived from massive text corpora.
However, their limitations remain clear: they are trained on symbolic descriptions rather than direct interaction with the physical world; their internal dynamics reflect linguistic statistics more than physical laws; and their world representations tend to be sparse, fragmented, and disconnected from geometric or temporal continuity.
Thus, while LLMs can be interpreted as language-driven world models, their predictive capacity is primarily symbolic and semantic, rather than grounded in spatial or physical understanding.







\subsection{Generative World Model}
\label{subsec:p3_s8_6_gen}

In contrast to representation-based foundation models, generative world models focus on explicitly predicting future states or synthesizing the evolution of the environment over time.
Rather than only learning compact latent encodings, these models aim to simulate how the world might change under different conditions, capturing temporal dynamics and causal structure.
Whereas the representation paradigm of Eq.~\eqref{eq:p3_s8_6_rep} predicts in an abstract latent space, generative world models learn an explicit conditional distribution over future observations and can be unrolled autoregressively over a horizon $H$,
\begin{equation}
\label{eq:p3_s8_6_gen}
p_\theta\big(O_{t+1:t+H} \mid O_{\le t},\, A_{t:t+H-1}\big) = \prod_{k=0}^{H-1} p_\theta\big(O_{t+k+1} \mid O_{\le t+k},\, A(t+k)\big),
\end{equation}
so that sampling from Eq.~\eqref{eq:p3_s8_6_gen} realizes the transition of Eq.~\eqref{eq:p3_s8_6_transition} directly in observation space (Figure~\ref{fig:world-modeling}(b)).
This explicit synthesis is what allows generative world models to be rendered, inspected, and acted within, at the cost of having to model a high-dimensional output distribution.
Broadly, generative world models can be grouped into two main families:
\emph{physics-driven simulation}, which relies on hand-crafted rules and equations,
and \emph{learning-driven world generation}, which learns world dynamics directly from data.
The former can be viewed as the deterministic limit of the latter: it fixes the transition of Eq.~\eqref{eq:p3_s8_6_gen} to a hand-specified rule set instead of a learned distribution, trading sample diversity and coverage for exact reproducibility and verifiability.
A third route, growing quickly since 2024, sits between the two: a learned model writes the rules as code and a simulator executes them (Section~\ref{subsubsec:p3_s8_6_code}).



\subsubsection{Physics-driven Simulation}

Physics-driven simulation builds interactive virtual worlds from human-designed physical, geometric, and logical rules rather than from learned latent spaces. The environment is governed by predefined physical laws, rendering rules, collision logic, and behavioral scripts, specified top-down as explicit mathematical formulations, so the same initial conditions always produce the same outcome; this determinism, together with interpretability and verifiability, is what makes such systems suitable for safety-critical engineering.

Game engines are the canonical example. Physics engines such as Bullet, Havok, PhysX, and Box2D implement Newtonian, rigid-body, soft-body, cloth, particle, and volumetric dynamics, and the engines around them add acceleration structures (\eg, BVH trees), real-time ray tracing, and navigation meshes for path planning, which gives them a natural foothold in spatial-intelligence tasks. Film production relies on the same principle at higher fidelity: pipelines built on Pixar USD, Maya, or Houdini simulate cloth, hair, smoke, fire, and fluids with differential-equation and fluid solvers (\eg, FLIP/MantaFlow) and physically based path tracing. Engineering, architecture, and scientific visualization extend rule-based simulation to structural mechanics, electromagnetics, illumination, and material response through finite element methods, computational fluid dynamics, and photonic simulation.

For embodied AI, physics-driven simulators such as CARLA~\cite{dosovitskiy2017carla}, AirSim, and Isaac Sim~\cite{NVIDIA2021Isaac} remain the standard platforms for testing and training in autonomous driving, navigation, and manipulation, because they model lighting, collisions, dynamics, and sensor noise closely enough to yield data distributions near those of the real world. The same engines underpin digital twins built on platforms such as NVIDIA Omniverse, Unity, and Unreal Engine, which are discussed in Part~IV, Section~\ref{sec:Digital Twins}.





\subsubsection{Learning-driven World Generation}

In contrast to simulation frameworks based on explicit physical rules, learning-driven world generation seeks to acquire world structure, dynamics, and causal regularities directly from large-scale data, in a bottom-up fashion.
These models do not encode Newtonian equations, lighting models, or geometric constraints by hand; instead, world behavior emerges as statistical patterns during training.
For example, a model may never be given the analytic form of the gravity equation, yet by observing countless videos of falling objects, it can still infer the regularity that objects tend to accelerate downward.
The resulting world models are probabilistic and can exhibit complex, emergent behaviors without explicit programming.


Current learning-driven world generation follows two technical directions:
(1) video-centric generative models that learn dynamic world evolution in pixel or latent space,
and (2) 3D/4D asset generation methods that directly model spatially structured worlds in three or four dimensions.




\paragraph{Video-centric Generative World Models.}
Recent advances in video generation have renewed interest in the idea of a ``world simulator''.
Video models are intrinsically suited to encoding spatiotemporal structure, and when they produce sequences that are not only visually coherent but also consistent with basic physical intuition, they begin to resemble implicit world models.
This became especially prominent with the release of OpenAI Sora~\cite{sora2024}, whose generated videos demonstrate strong temporal consistency, long-horizon stability, and apparent physical plausibility, prompting the community to reconsider whether deep video generation might be converging towards a form of physical world modeling.
Subsequently, Google DeepMind's Genie~3~\cite{genie3} extended video generation from passive ``watchable'' content to interactive experiences: the model updates the generated world in response to user actions, enabling agent-in-the-loop, closed-loop world simulation, complementing the action-conditioned scene generators surveyed in §\ref{subsec:p3_s8_2_2}.
World Labs' RTFM~\cite{worldlabs_rtfm} further highlights real-time performance, achieving high frame-rate and low-latency prediction. Systems such as Wan~\cite{wan2025wan}, Emu3.5~\cite{cui2025emu35}, OpenAI Sora~2~\cite{sora2}, and Veo~3~\cite{veo3} continue this trajectory, gradually pushing video generation from mere visual fidelity towards world coherence, physical consistency, and interactivity.
Action-conditioned diffusion has also turned games into testbeds: GameNGen~\cite{valevski2024diffusion} simulates the game DOOM interactively in real time, and DIAMOND~\cite{alonso2024diffusionworldmodelingvisual} trains reinforcement-learning agents entirely inside a diffusion world model of Atari games, while NVIDIA's Cosmos~\cite{nvidia2025cosmos} offers pretrained video world foundation models intended to be post-trained as simulators for robotics and driving.

Video world models are particularly promising in autonomous driving. GAIA-1~\cite{hu2023gaia} proposes a generative world model tailored for driving scenarios, built around a spatiotemporal Transformer that generates future multi-view scenes based on past observations.
By implicitly capturing traffic patterns, pedestrian behavior, scene geometry, and multi-agent motion, GAIA-1 can support closed-loop evaluation, behavior prediction, and synthetic data generation, thereby reducing the cost and risk of traditional real-world testing.



\paragraph{3D/4D Asset Generation as World Modeling.}
A complementary avenue is to build world models that operate directly on structured three- or four-dimensional representations, such as 3D geometry, occupancy fields, or 4D spatiotemporal volumes.
Unlike video-based approaches that primarily model dynamics in pixel space, these methods explicitly parameterize geometry and temporal evolution, making it easier to represent interactive, spatially consistent worlds.


Genesis~\cite{Genesis} exemplifies this direction. It is not only a generative 3D/4D physics simulation platform, but also integrates a vision-language model-driven ``world-building agent'' that uses a tool API to create dynamic 4D environments (geometry plus dynamics).
These generated worlds can be rendered into multiple modalities (RGB, depth, semantics, contact events, etc.), providing rich supervision for training embodied AI systems.
Genesis thus combines a unified physics engine with generative agents for large-scale world construction.


Autonomous driving again serves as a natural application domain for 3D world models.
DriveArena~\cite{yang2025drivearena} builds a closed-loop generative simulation platform, allowing planning modules to iteratively probe and learn within synthesized 3D driving environments.
Its diffusion-model-based multi-agent scene generation can rapidly produce long, interactive driving episodes.
OccWorld~\cite{zheng2024occworld} and OccSora~\cite{wang2024occsora} adopt 4D occupancy grids as the core representation, predicting future 3D occupancy over time from continuous video and sensor streams, thereby providing a unified motion-and-structure world model.
Copilot4D~\cite{zhang2023copilot4d} tokenizes LiDAR scans and forecasts future LiDAR point clouds with discrete diffusion, learning a world model from unlabeled sensor data alone.


Generating \emph{interactive} 3D worlds brings generative world models closest to their conceptual ideal.
Marble~\cite{Feifei2025Marble} from World Labs combines action-conditioned generation, geometric consistency constraints, and multimodal inputs, so that users can act within the generated world and observe the consequences.
Its next-generation model, Atlas~\cite{worldlabs2026atlas}, is an omni world model pretrained from scratch on text, images, video, and 3D: an autoregressive diffusion transformer that conditions every generated image and depth map on an explicit camera pose and grounds them all in a shared 3D context, so that the same model generates camera-controlled video of up to a minute, reconstructs scenes from a handful of images, and exports its output as point clouds or 3D Gaussian splats. World Labs reports that it outperforms specialized open-source reconstruction models, and that it will power future versions of Marble.


\subsubsection{Code-and-Simulator World Models}
\label{subsubsec:p3_s8_6_code}
Between hand-written simulation and learned generation sits a route that has grown quickly since 2024: a language model \emph{writes} the world model as executable code, and a simulator or engine runs it. The learned component is the author of the rules, not the rules themselves,
\begin{equation}
\label{eq:p3_s8_6_code}
\mathcal{P} \sim p_\theta\big(\mathcal{P} \mid c\big), \qquad \hat{S}(t+1) = \mathrm{Sim}\big(S(t),\, A(t);\, \mathcal{P}\big),
\end{equation}
where $c$ is the conditioning context (a description, observations, or a history of interactions), $\mathcal{P}$ the generated program, and $\mathrm{Sim}$ the interpreter or engine that executes it: plain Python, a physics engine such as MuJoCo or PyChrono, or Blender. Compared with Eq.~\eqref{eq:p3_s8_6_transition}, the learned $W_\theta$ is replaced by a program that can be read, tested, and repaired, so the resulting world keeps the determinism and verifiability of physics-driven simulation while gaining the open-endedness of a generative model. The scene programs of Section~\ref{sec:p3_s8_7} rest on the same idea; the difference here is that $\mathcal{P}$ encodes how the world changes under actions, not only how it is laid out.

In one line of work, code is an agent's transition model. Code World Models~\cite{dainese2024codeworld} synthesize Python implementations of reinforcement-learning environments with a generate, improve, and fix search over programs (GIF-MCTS), evaluated on a benchmark of 18 environments, and plan with them faster and more sample-efficiently than with learned dynamics; WorldCoder~\cite{tang2024worldcoder} builds its world model as a program from its own interactions and transfers knowledge to a new environment by editing that program.

In the other, code builds a simulated 3D world. SceneCraft~\cite{hu2024scenecraft} turns a description into Blender Python, planning with a scene graph, converting spatial relations into numerical layout constraints, and checking renders with a vision-language model. GenSim~\cite{wang2023gensim} and RoboGen~\cite{wang2023robogen} have language models write robotic simulation tasks, scenes, and rewards, making the simulator an open-ended source of training data, and the world-building agent of Genesis~\cite{Genesis} follows the same pattern. Recent systems close the loop with the simulator itself. 
Code2Worlds~\cite{zhang2026code2worlds} generates 4D worlds as Blender code refined by a vision-language motion critic; Code2Worlds-2~\cite{zhang2026code2worlds2} compiles such worlds into physically valid MuJoCo scenes and revises them from penetration, contact, and stability signals; and ChronoAgentic~\cite{wang2026chronoagentic} has cooperating agents write and repair PyChrono simulations, meeting the joint semantic and physical correctness criterion on 82.5\% of PhyWorldBench demos against 52.5\% for the best of ten text-to-video models.

The price is coverage. 
A program can express only what its simulator supports, so photorealistic appearance, deformable and fluid materials, and open-ended visual variety still favor learned generators, while code-and-simulator models win where physical correctness has to be checked.


\subsection{Mental World Modeling}
\label{subsec:p3_s8_6_mental}


Physical world models enable agents to understand and predict how the external environment evolves.
However, intelligent behavior in real-world settings also requires an internal model of \emph{other minds}.
Humans naturally construct such models, \ie, often referred to as \emph{mental models}, to reason about others' beliefs, desires, intentions, emotions, and social context.
This ability, commonly associated with \emph{Theory of Mind (ToM)} \cite{astington1995theory}, underlies human social interaction, collaboration, communication, and decision-making.



Craik's ``small-scale models'' of reality are the origin of this notion too~\cite{nersessian1993theoretician,staggers1993mental}, and Johnson-Laird later argued that human reasoning, problem solving, and causal inference rely on abstract internal representations of situations rather than on replaying detailed physical dynamics.
These representations selectively capture the most relevant structure, \ie, objects, relations, causal dependencies, and possible outcomes, while abstracting away low-level perceptual details.
For example, when cycling, a human does not mentally simulate every joint torque or physical trajectory; instead, one imagines ``whether I would fall or not'' when encountering an obstacle.
This high-level, conceptual simulation is a defining property of mental world modeling.



\subsubsection{From Physical to Cognitive Mental World Models}
Human mental simulation operates by combining current observations with stored knowledge to predict plausible futures.
Perception provides the immediate sensory input, which is then transformed into an internal symbolic and probabilistic representation of the world.
Within this representation, the brain runs rapid, approximate simulations that evaluate different hypothetical actions.
These simulations are not physically precise but are cognitively meaningful: they help determine risks, anticipate others' responses, or evaluate alternative courses of action.
In this sense, mental models serve as a compact, abstract, and counterfactual ``world engine'' for decision making.


In artificial systems, the distinction between physical and mental world models mirrors that in human
cognition.
Physical world models predict state transitions in the environment (\eg, dynamics, geometry, object motion), whereas \emph{mental world models} encode the latent mental states of human users or other agents.
These models allow an AI system to reason about what humans know, believe, want, or intend to do, capabilities that embodied agents, collaborative robots, intelligent tutors, and multi-agent planners all need.



Inspired by the Belief-Desire-Intention (BDI)~\cite{rao1995bdi,georgeff1998belief} framework, a mental world model for AI typically includes a small set of structured variables.
\textbf{Beliefs.} These are representations of what the human knows, perceives, or believes about the world.
\textbf{Desires.} These capture user goals, preferences, motivations, and values.
\textbf{Intentions.} These are planned or ongoing actions that the agent expects the human (or another agent) to follow.
Such decomposition enables an AI system to explicitly reason about human rationality, goal pursuit, and action planning.
Formally, a mental world model can be cast as the inference of a latent mental state $M(t)$, comprising the beliefs, desires, and intentions above, from observed behavior, followed by a prediction of the other agent's next action,
\begin{equation}
\label{eq:p3_s8_6_mental}
\hat{M}(t) = \argmax_{M}\, p\big(M \mid O_{\le t}\big), \qquad \hat{A}^{\text{other}}(t+1) = \pi\big(\hat{M}(t)\big).
\end{equation}
This formulation mirrors the transition of Eq.~\eqref{eq:p3_s8_6_transition} but replaces the physical state with the inferred mental state of another agent (Figure~\ref{fig:world-modeling}(c)), so that the predicted quantity is no longer how the environment will evolve but how a human or peer agent is likely to act.

Machine Theory of Mind is now measured directly. ToMi~\cite{le2019tomi} tests false-belief reasoning in short stories with controls against shortcut solutions, BigToM~\cite{gandhi2023bigtom} generates controlled social scenarios with language models to probe whether models infer beliefs from percepts and actions, and MMToM-QA~\cite{jin2024mmtomqa} poses belief and goal questions about a person acting in a household from video and text together. Large language and multimodal models remain unreliable on these tests, whereas MMToM-QA's BIP-ALM, which performs Bayesian inverse planning over a symbolic state with a language model scoring the likelihoods, does better; it is an explicit instance of the inference in Eq.~\eqref{eq:p3_s8_6_mental}.

\paragraph{Dimensions of Mental State Representation.}
Beyond the belief-desire-intention core, a mental world model for embodied agents must span several further layers of internal state and social context. 
It has to read \emph{affect}, recognizing emotional cues and how they shift a person's risk tolerance and willingness to cooperate. 
It has to represent \emph{social dynamics}: the relationships, roles, group structures, and cultural norms that shape behavior. 
And it has to interpret \emph{communication}, both the verbal channel and the non-verbal one carried by tone, gesture, posture, and facial expression. 
These layers are what turn a bare goal estimate into a socially situated account of why a person acts as they do.



Human spatial reasoning is social as well as geometric, which ties mental modeling to spatial intelligence.
To assist or collaborate with humans in homes, factories, classrooms, or vehicles, an agent must infer the intention behind a spatial action (\eg, where a person intends to move next), anticipate trajectories from goals or habits, adapt navigation or manipulation to user preferences, and interpret spatial configurations that reflect what a person believes or misunderstands.


\subsubsection{A Joint Physical-Mental Formulation}
\label{subsubsec:p3_s8_6_mwm}
The benchmarks above still treat the mental state as something inferred next to a world model, not as part of it: Eq.~\eqref{eq:p3_s8_6_mental} estimates another agent's mind, but the environment it acts in is predicted separately. 
Mental World Modeling (MWM)~\cite{fei2026mental} makes the mental state part of the state that the world model of Eq.~\eqref{eq:p3_s8_6_transition} rolls forward. 
We follow its notation, in which subscripts index time, so that $s_t$ and $a_t^{\epsilon}$ play the roles of $S(t)$ and $A(t)$ above.
The case for a joint state is that neither half suffices alone. For a target agent $\epsilon$ whose true action distribution is $p^{\star}$, a physical-only representation fails whenever
\begin{equation}
\label{eq:p3_s8_6_mwm_property}
p^{\star}\big(a_t^{\epsilon} \mid s_t^{\mathrm{phy}}, s_t^{\mathrm{ment}}\big) \neq p^{\star}\big(a_t^{\epsilon} \mid s_t^{\mathrm{phy}}, \tilde{s}_t^{\mathrm{ment}}\big),
\end{equation}
that is, whenever two scenes that are physically identical but differ in what people believe or want ($s_t^{\mathrm{ment}}$ versus $\tilde{s}_t^{\mathrm{ment}}$) lead to different actions; the symmetric condition, with the mental state held fixed and the physical state changed, rules out a mental-only representation.

\paragraph{Joint State.}
The world state is the pair
\begin{equation}
\label{eq:p3_s8_6_mwm_state}
s_t = \big(s_t^{\mathrm{phy}}, s_t^{\mathrm{ment}}\big) \in \mathcal{S}^{\mathrm{phy}} \times \mathcal{S}^{\mathrm{ment}},
\qquad
s_t^{\mathrm{ment}} = \big(\{m_t^{i}\}_{i\in\mathcal{I}},\, \{m_t^{G}\}_{G\in\mathcal{G}},\, R_t^{\mathrm{ment}},\, \alpha_t\big).
\end{equation}
The physical part collects the objects, the embodied characters, their spatial and contact relations, and the environmental conditions of the scene. The mental part holds an individual state $m_t^{i}$ for every agent $i$ in the set $\mathcal{I}$, optional states $m_t^{G}$ for groups $G\in\mathcal{G}$, the social attitudes and role relations $R_t^{\mathrm{ment}}$ between agents, and the atmosphere $\alpha_t$ of the setting. Each individual state extends the belief-desire-intention core introduced above,
\begin{equation}
\label{eq:p3_s8_6_mwm_individual}
m_t^{i} = \big(\mathrm{id}_t^{i},\, b_t^{i},\, q_t^{i},\, g_t^{i},\, \iota_t^{i},\, e_t^{i},\, d_t^{i},\, n_t^{i},\, c_t^{i}\big),
\end{equation}
with identity attributes, beliefs, attention focus, goals, intentions, emotions, dispositions, norms, and behavioral constraints. 
Beliefs, goals, and intentions are the BDI variables; emotions, norms, and the relations $R_t^{\mathrm{ment}}$ turn the affect and social-dynamics layers discussed above into explicit state variables.

\paragraph{Target-Specific Observation.}
The model reasons from the perspective of one target agent $\epsilon\in\mathcal{I}$, which sees only a rendering of the joint state:
\begin{equation}
\label{eq:p3_s8_6_mwm_obs}
o_t^{\epsilon} = \Omega_{\epsilon}(s_t) = \Big(\Omega_{\epsilon}^{\mathrm{phy}}\big(s_t^{\mathrm{phy}};\, \kappa_t^{\epsilon}\big),\; \Omega_{\epsilon}^{\mathrm{ment}}\big(s_t^{\mathrm{phy}}, s_t^{\mathrm{ment}};\, \rho_t^{\epsilon}\big)\Big),
\end{equation}
where $\kappa_t^{\epsilon}$ is the target's perceptual access (position, line of sight, hearing range, available modalities) and $\rho_t^{\epsilon}$ its social and cognitive perspective (role, prior beliefs, relationships to others, awareness of norms, Theory-of-Mind capacity). The mental part of the observation contains the target's view of its own state and its estimates $\widehat{m}^{\,j,(\ell)}_{t,\epsilon}$ of every other agent $j$, nested up to a depth $L$; at the first level they are sampled as
\begin{equation}
\label{eq:p3_s8_6_mwm_tom}
\widehat{m}^{\,j,(1)}_{t,\epsilon} \sim Q^{(1)}_{\epsilon}\big(m_t^{j} \mid o_t^{\epsilon,\mathrm{phy}},\, m_t^{\epsilon},\, R_t^{\mathrm{ment}},\, \alpha_t\big),
\end{equation}
which is the inference of Eq.~\eqref{eq:p3_s8_6_mental}, made stochastic and conditioned on the target's own mind and social context. Spatial access enters the mental model directly through $\kappa_t^{\epsilon}$: what the target can see from where it stands bounds what it can believe, so a poor decision can follow from an object outside its view as much as from a mistaken belief.

\paragraph{Action and Coupled Transition.}
The target proposes $K$ candidate actions $\mathcal{A}_t^{\epsilon} = \Pi_{\epsilon}^{K}(o_t^{\epsilon}) = \{a^{\epsilon}_{t,1},\dots,a^{\epsilon}_{t,K}\}$. Each action has a physical carrier and a mental content, $a_t^{\epsilon} = (a_t^{\epsilon,\mathrm{phy}}, a_t^{\epsilon,\mathrm{ment}})$ (approaching someone, and meaning to comfort them), and the policy factorizes as $\Pi_{\epsilon}(a_t^{\epsilon,\mathrm{phy}}, a_t^{\epsilon,\mathrm{ment}} \mid o_t^{\epsilon}) = \Pi_{\epsilon}^{\mathrm{phy}}(a_t^{\epsilon,\mathrm{phy}} \mid o_t^{\epsilon})\,\Pi_{\epsilon}^{\mathrm{ment}}(a_t^{\epsilon,\mathrm{ment}} \mid a_t^{\epsilon,\mathrm{phy}}, o_t^{\epsilon})$. The world model then predicts the next joint state with a staged transition,
\begin{equation}
\label{eq:p3_s8_6_mwm_transition}
T_{\theta}\big(s_{t+1} \mid s_t, a_t^{\epsilon}\big) = T_{\theta}^{\mathrm{phy}}\big(s_{t+1}^{\mathrm{phy}} \mid s_t^{\mathrm{phy}}, s_t^{\mathrm{ment}}, a_t^{\epsilon,\mathrm{phy}}\big) \cdot T_{\theta}^{\mathrm{ment}}\big(s_{t+1}^{\mathrm{ment}} \mid s_t^{\mathrm{phy}}, s_t^{\mathrm{ment}}, a_t^{\epsilon,\mathrm{phy}}, a_t^{\epsilon,\mathrm{ment}}\big).
\end{equation}
The asymmetry is deliberate. An intention moves nothing until it is carried out physically, so the physical factor ignores $a_t^{\epsilon,\mathrm{ment}}$; the mental factor depends on both components, because the same reassuring words land differently when the speaker is approaching than when walking away. The physical factor still conditions on the mental state, since what objects, space, and visibility allow is joined by what other agents intend to do. Those other agents are not controlled: their actions $a_t^{-\epsilon}$ are marginalized under the simulator's predictive distribution,
\begin{equation}
\label{eq:p3_s8_6_mwm_marginal}
T_{\theta}\big(s_{t+1} \mid s_t, a_t^{\epsilon}\big) = \int T_{\theta}\big(s_{t+1} \mid s_t, a_t^{\epsilon}, a_t^{-\epsilon}\big)\, p\big(a_t^{-\epsilon} \mid s_t\big)\, \mathrm{d}a_t^{-\epsilon},
\end{equation}
with $p(a_t^{-\epsilon} \mid s_t) = \prod_{j\neq\epsilon} \Pi_j\big(a_t^{j} \mid \Omega_j(s_t)\big)$ for simultaneous moves, so that every agent still acts on its own partial observation. The whole pipeline, $s_t \xrightarrow{\Omega_{\epsilon}} o_t^{\epsilon} \xrightarrow{\Pi_{\epsilon}} a_t^{\epsilon} \xrightarrow{T_{\theta}} s_{t+1}$, is Eq.~\eqref{eq:p3_s8_6_transition} with a mental component added to the state and a perspective added to the observation.

\paragraph{Evaluating Candidate Futures.}
Simulating each candidate gives a branch $\hat{s}^{k}_{t+1}$, and the predicted decision is the branch of highest value,
\begin{equation}
\label{eq:p3_s8_6_mwm_select}
\hat{k} = \argmax_{k\in\{1,\dots,K\}}\, r_k .
\end{equation}
In MENTIS, the training-free and fully inspectable baseline of MWM, the branch value $r_k$ is a weighted combination of normalized scores for physical plausibility, mental consistency, and social appropriateness, and a binary safety and legality veto sets it to zero. Each formal component maps to a module: a state parser recovers $\hat{s}_t$, an observation generator implements $\Omega_{\epsilon}$, an action parser structures the candidates, a world-state transitor implements $T_{\theta}$, and an evaluator assigns $r_k$. On situated decision scenarios presented as text, image, and sounding-video stories, experiments across eight LLM-based world models indicate that the explicit mental component is what makes a person's actual decision predictable~\cite{fei2026mental}.


\subsection{Discussions and Challenges}
\label{subsec:p3_s8_6_discussion}

\paragraph{Representation-based World Models as an Alternative Foundation Model Paradigm.}
Representation world models may emerge as a parallel paradigm to current LLMs, functioning as large-scale foundational models trained through massive self-supervised representation learning to capture physical regularities, causal dynamics, and environmental constraints.
Unlike LLMs, which learn the world indirectly through linguistic abstraction, representation world models aim to model the physical world directly from multimodal observations.
Such models could serve video generation, action prediction in VLA systems, 3D and 4D understanding and generation, and decision-making in autonomous driving.



\paragraph{The Necessity of Multi-source Intelligence Integration.}
World models will likely require the integration of heterogeneous informational sources.
On the one hand, LLMs embody language-driven intelligence, while world models emphasize physically grounded intelligence; rather than being mutually exclusive, a truly capable world model must incorporate both forms of reasoning.
On the other hand, accurate world understanding requires far more than visual perception alone: audio, speech, tactile signals, force feedback, and even proprioception provide essential supplementary cues about the dynamics and structure of the physical environment. Yet the current research community focuses overwhelmingly on visual-only modeling.
Building world models that holistically integrate multimodal sensory information remains a major challenge.


\paragraph{Unifying Physical Modeling and Mental Modeling.}
Another promising direction is the integration of physical and mental modeling within a unified world-modeling framework.
Embodied agents operating in real environments must simultaneously reason about the physical evolution of the environment and the mental or behavioral responses of humans (or other agents).
Recent work in multi-agent systems, social navigation, pedagogical reasoning, human-robot collaboration, and interactive LLM agents emphasizes the growing need for AI systems to build internal hypotheses about both external dynamics and human mental states. 
Mental World Modeling~\cite{fei2026mental} is an early formulation of such a joint state (Section~\ref{subsubsec:p3_s8_6_mwm}).
Ultimately, scalable general-purpose world models must simulate not only \emph{what will happen} but also \emph{what others believe, intend, desire, or misunderstand}.


\paragraph{Balance between Generative Fidelity and Physical Correctness.}
A central challenge for learning-driven world models is reconciling generative expressiveness with physical accuracy.
High-capacity video generators often produce visually stunning but physically implausible outcomes, while models optimized for strict physical consistency may underperform in realism or diversity.
Achieving stable long-horizon predictions without temporal drift or compounding errors remains unsolved.
Designing architectures, training objectives, and constraints that jointly enforce physical laws, spatial coherence, and generative flexibility is a key open problem for the next generation of world simulators.


\paragraph{Evaluation, Benchmarks, and Verifiability of World Models.}
Unlike language models, which can be evaluated through perplexity, benchmarks, or task accuracy, world models lack standardized, reliable, and comprehensive evaluation protocols.
Physical correctness, causal consistency, interactive fidelity, multimodal coherence, and long-horizon stability are difficult to quantify in a unified manner.
Furthermore, world models that aim to simulate open-ended environments raise issues of verifiability, reproducibility, and safety: how can we assess whether a world model reliably generalizes to unseen real-world scenarios, or whether its internal dynamics exhibit hidden biases or unsafe behaviors?
Establishing principled evaluation methodologies and shared benchmarks for both representation-based and generative world models is an urgent research direction.

% \newpage
\section{Programmatic and Symbolic Spatial Generation}
\label{sec:p3_s8_7}

\FramedInitial{P}{rogrammatic} and symbolic spatial generation synthesizes 2D and 3D scene structure through explicit rules, executable programs, and symbolic geometric abstractions.
It generates the structural backbone of a scene (spatial layout, hierarchical organization, geometric constraints, executable logic) rather than its raw appearance, which gives it controllability, compositional generalization, global spatial consistency, and interpretability.

Where the largely neural samplers of §\ref{sec:p3_s8_1}–§\ref{sec:p3_s8_6} map a condition directly to pixels, voxels, or a neural latent, programmatic methods emit an intermediate \emph{program} that a deterministic renderer executes into the final scene.
This intermediate layer is what makes the structure auditable, editable, and compositionally reusable, properties that flat neural outputs lack.
Three paradigm families mark the historical and technical progression of the field: \emph{rule- and grammar-based procedural generation}, \emph{symbolic graphics programs and domain-specific languages}, and \emph{neuro-symbolic language-to-program synthesis} driven by large language models.

\subsection{Task Definition}
\label{subsubsec:p3_s8_7_task}

Let $\mathbf{c}\in\mathcal{C}$ denote a conditioning intent, typically a natural-language description, a layout specification, or a partial constraint set, and let $\kappa:\mathcal{C}\!\to\!\mathcal{C}'$ be its encoder.
A programmatic generator $G_\theta$ emits a symbolic program $\rho$ drawn from a graphics language or domain-specific language $\mathcal{L}$, which a deterministic executor $R$ evaluates into a concrete scene $\hat{\mathbf{s}}$ (a raster image, vector graphic, UI specification, or 3D shape):
\begin{equation}
\label{eq:p3_s8_7_objective}
\rho = G_\theta\big(\kappa(\mathbf{c})\big), \qquad \hat{\mathbf{s}} = R(\rho), \qquad \rho \in \mathcal{L}.
\end{equation}
This formulation inherits the conditional-sampler abstraction of Eq.~\eqref{eq:t2i_generator} but replaces its learned decoder $D$ (Eq.~\eqref{eq:t2i_autoencoder}) over a continuous latent space $\mathcal{Z}$ with an interpretable, discrete program $\rho$ and an \emph{exact, non-learned} renderer $R$.
The neural generators of §\ref{sec:generation} are recovered as the limiting case in which $R$ is itself a neural renderer and $\rho$ a continuous latent, whereas the symbolic case constrains $\rho$ to be a human-readable executable artifact.
Two properties follow directly: errors are localizable to specific program statements, and structural constraints can be enforced or verified deterministically on $\rho$ before rendering.



\subsection{Methodology}
\label{subsubsec:p3_s8_7_methodology}

Figure~\ref{fig:p3_s8_7_paradigms} schematizes the three paradigm families that instantiate Eq.~\eqref{eq:p3_s8_7_objective}, contrasting how each obtains the program $\rho$ before the shared deterministic renderer $R$ executes it.

\begin{figure}[t]
    \centering
    \includegraphics[width=\linewidth]{figures/part3/section8/programmatic_spatial_generation.pdf}
    \caption{Three paradigm families for programmatic spatial generation, all producing a program $\rho$ that a deterministic renderer $R$ executes (Eq.~\eqref{eq:p3_s8_7_objective}): (P1) rule- and grammar-based procedural generation, (P2) symbolic graphics programs and domain-specific languages, and (P3) language-to-program synthesis by an LLM planner.}
    \label{fig:p3_s8_7_paradigms}
\end{figure}

\paragraph{Rule- and Grammar-based Procedural Generation.}
The earliest line of work, procedural content generation (PCG)~\cite{parish2001procedural} in computer graphics and game development, synthesizes scenes such as dungeon maps, terrains, architectural plans, or UI layouts from hand-crafted rules or parameterized procedural programs.
Classical formalisms instantiate $G_\theta$ as a fixed production system rather than a learned model.
The Lindenmayer system (L-system)~\cite{lindenmayer1968mathematical}, originally developed to model biological growth, employs parallel rewriting rules to simulate natural growth processes and fractal patterns, making it ideal for self-similar structures such as vegetation and complex terrain.
Formally, the program evolves by repeated parallel application of a production set $\Pi$ to an axiom $\omega$,
\begin{equation}
\label{eq:p3_s8_7_grammar}
\rho_0 = \omega, \qquad \rho_{t+1} = \Pi(\rho_t), \qquad \hat{\mathbf{s}} = R(\rho_T),
\end{equation}
where the generated symbolic string is not the geometry itself but a sequence of Turtle Graphics commands that the interpreter $R$ parses to draw the final shape in 2D or 3D space.
Shape grammars, introduced by \citet{stiny1980introduction}, provide another influential symbolic mechanism by defining geometric primitives and rewrite rules for composing architectural or engineered shapes.
These early paradigms demonstrated a fundamental insight: spatial structures can be synthesized by executing symbolic generative rules, rather than by directly predicting raw geometry.

\paragraph{Symbolic Graphics Programs and Domain-Specific Languages.}
Symbolic graphics programs (SGPs) encode spatial layouts and geometric configurations as formal programs that execute deterministically into images, vector graphics, UI designs, or 3D shapes, bridging abstract scene intent and concrete visual realization. Common languages such as SVG~\cite{bing2025learning}, TikZ/LaTeX~\cite{ellis2018learning}, and HTML/CSS~\cite{chen2025designcoder} already carry rich symbolic layout specifications in real-world use.
In parallel, researchers have developed domain-specific languages (DSLs) tailored for fine-grained spatial and visual synthesis.
Examples include hierarchical GUI-code generators~\cite{zhu2018automatic} that translate visual structures into programmatic UI specifications, and ShapeAssembly~\cite{jones2020shapeassembly}, a 3D assembly language in which objects are represented as executable programs describing part composition, symmetry, and connectivity.
Such DSLs make the advantages of explicit symbolic control concrete: precise geometry, modular design, verifiable constraints, and systematic compositionality.

\paragraph{Neuro-symbolic and Language-to-Program Synthesis.}
A more recent direction couples symbolic programs with learned models, so that $G_\theta$ becomes a trained synthesizer rather than a fixed rule set. The Scene Language~\cite{zhang2025scenelanguage} unifies neural and symbolic components: a program specifying relational and hierarchical layout structure, natural-language labels summarizing entity semantics, and learned embeddings capturing visual identity.
Executable structure, rather than raw geometry, thus becomes a primary representational layer for scenes.
Large language models~\cite{brown2020language,achiam2023gpt} have accelerated the move toward language-to-program spatial generation: instead of predicting images or layouts directly, LLM-based systems act as symbolic planners that synthesize the program by approximately maximizing its conditional likelihood,
\begin{equation}
\label{eq:p3_s8_7_synthesis}
\rho^{\star} = \argmax_{\rho \in \mathcal{L}} \, p_\theta\big(\rho \mid \kappa(\mathbf{c})\big),
\end{equation}
mapping natural-language descriptions to executable layout programs~\cite{aguina2024open, hsu2025programs, yang2024holodeck}.
These programs then guide downstream renderers to produce UI designs, graphic layouts, architectural plans, or 3D scenes with precise structural fidelity, and users can edit or refine the underlying program with fine-grained control.

The three families differ along two cross-cutting axes.
The first is the \emph{representation} of $\rho$: a parallel string-rewriting system (grammar-based generation), an executable program in a general-purpose or domain-specific graphics language (SGPs and DSLs), or a hybrid structure that interleaves code, language, and learned embeddings (neuro-symbolic synthesis).
The second is the \emph{inference procedure} that produces $\rho$: hand-authored production rules, human-written or template-instantiated programs, or learned program synthesis and LLM-based planning. Across all three, the target modality (2D raster, vector graphics, UI, or 3D) is determined solely by the renderer $R$, which keeps the symbolic and rendering concerns cleanly separated.



\subsection{Benchmarks}
\label{subsubsec:p3_s8_7_benchmarks}

Evaluation in this area differs from the distributional metrics of §\ref{sec:generation} because the immediate output is a program, not an image, so quality is assessed along several complementary axes.
\textbf{Program validity} measures whether the emitted $\rho$ parses and executes without error in its target language, a prerequisite that purely neural generators do not face.
\textbf{Rendered-output fidelity} then compares the executed scene $R(\rho)$ against a reference, using pixel- or perceptual-level similarity for raster and vector targets and geometric agreement for 3D shapes.
\textbf{Constraint satisfaction} checks whether explicitly specified spatial relations, counts, alignments, or symmetries hold in the rendered result, a property that symbolic methods are designed to guarantee.
\textbf{Controllability and editability} assess how reliably localized edits to $\rho$ produce the intended localized changes in the scene without unintended global drift. Concretely, these axes are exercised across vector-graphics and SVG synthesis~\cite{bing2025learning}, UI and GUI code generation~\cite{zhu2018automatic, chen2025designcoder}, 3D program induction over part-structured shapes~\cite{jones2020shapeassembly}, evaluated against part-hierarchy benchmarks such as PartNet~\cite{mo2019partnet}, and full scene-program synthesis~\cite{zhang2025scenelanguage, yang2024holodeck}, where the executable program serves simultaneously as the prediction and as the verifiable ground for measurement.

\subsection{Challenges and Trends}
\label{subsubsec:p3_s8_7_challenges}

The field has moved from entirely hand-authored rules toward neuro-symbolic integration with large models, driven by the demand for generators that are more controllable and interpretable than end-to-end neural methods.

\textbf{Generalization.} Automatically synthesized programs still struggle with novel constraints and complex semantic specifications~\cite{mao2025neurosymbolicconcepts}. Program synthesis has to infer the structural logic of a scene rather than memorize surface patterns if the generator is to extrapolate to compositions and constraint combinations absent from training.

\textbf{Efficiency and differentiability.} Neuro-symbolic programs that are both efficient and fully differentiable~\cite{guo2023efficientsymbolicpolicy} would let a rendering-level loss propagate back through the program, so that the whole pipeline, from symbolic planning to neural rendering, could be trained and fine-tuned end to end.

\textbf{Embodied interaction.} Procedurally synthesized interactive environments already supply large-scale embodied worlds~\cite{kolve2017ai2thor, procthor, wang2024pipeline}, and the next step is for agents to generate or validate symbolic scene structure through active exploration rather than from static datasets alone, yielding symbolic representations of physical spaces that can be checked against what the agent observes.



\section{Summary and Outlook}
\label{sec:p3_s8_8}

This chapter has surveyed spatial reconstruction and generation across seven subfields that, taken together, span an axis of increasing output structure and dimensionality: 2D images (§\ref{sec:p3_s8_1}), 3D objects and scenes (§\ref{sec:p3_s8_2}), video (§\ref{sec:p3_s8_3}), embodied motion (§\ref{sec:p3_s8_4}), 4D space-time content (§\ref{sec:p3_s8_5}), world models (§\ref{sec:p3_s8_6}), and interpretable symbolic programs (§\ref{sec:p3_s8_7}). Despite their different output domains, the subfields are unified by a single conditional-generation backbone, established in §\ref{subsec:p3_s8_1_1}: an encoder $\kappa$ maps a condition $\mathbf{c}$ into a latent code, a generator $G_\theta$ samples structure from that code, and a decoder or renderer turns the sampled structure into a concrete artifact (Eq.~\eqref{eq:t2i_autoencoder}, Eq.~\eqref{eq:t2i_generator}). What distinguishes the subfields is not this backbone but the \emph{output space} and the \emph{decoder} that closes it: a pixel decoder for images, a volumetric or splatting renderer for 3D, a spatio-temporal decoder for video, a kinematic sampler for motion, a time-indexed 3D representation for 4D, a state-transition predictor for world models, and an exact, non-learned program executor for symbolic generation. Table~\ref{tab:p3_s8_chapter_comparison} consolidates this view.

\begin{table}[t]
\centering
\caption{The seven spatial-generation subfields of this chapter, all instances of the conditional latent-generator backbone of §\ref{subsec:p3_s8_1_1}.}
\label{tab:p3_s8_chapter_comparison}
\fontsize{8}{10}\selectfont
\setlength{\tabcolsep}{1.0mm}
\begin{tabular}{L{0.17\textwidth} L{0.20\textwidth} L{0.33\textwidth} L{0.22\textwidth}}
\toprule
\textbf{Subfield} & \textbf{Output space} & \textbf{Dominant paradigms} & \textbf{Decoder / renderer} \\
\midrule
§\ref{sec:p3_s8_1} 2D image & Pixel grid (raster) & GAN; autoregressive; latent diffusion; flow matching & Learned pixel decoder \\
§\ref{sec:p3_s8_2} 3D object/scene & Mesh / NeRF / 3DGS / latent & Native 3D diffusion; score distillation; multi-view lift; feed-forward reconstruction & Differentiable / volumetric renderer \\
§\ref{sec:p3_s8_3} Video & Frame sequence (clip latent) & Inflated T2I backbones; spatial / temporal control; editing & Spatio-temporal decoder \\
§\ref{sec:p3_s8_4} Motion & Pose / trajectory sequence & Regression; diffusion samplers; editing operators & Kinematic / skeletal decoder \\
§\ref{sec:p3_s8_5} 4D & Time-indexed 3D representation & Score distillation; generate-then-reconstruct & Deformable / per-frame renderer \\
§\ref{sec:p3_s8_6} World model & Latent state / future observations & Representation; generative; mental world models & Latent transition / observation predictor \\
§\ref{sec:p3_s8_7} Symbolic program & Executable program in a language $\mathcal{L}$ & Procedural grammars; SGP/DSL; LLM program synthesis & Deterministic, non-learned executor \\
\bottomrule
\end{tabular}
\end{table}

\textbf{A shared backbone, specialized by output space.} The progression from §\ref{sec:p3_s8_1} to §\ref{sec:p3_s8_7} is best read not as seven disjoint problems but as one generative abstraction instantiated against progressively richer codomains, and each section casts its own subtasks as conditional specializations of that common sampler. The symbolic case (§\ref{sec:p3_s8_7}) is the limiting instance in which the decoder degenerates to an exact program executor, trading the expressiveness of a learned renderer for interpretability and verifiable control.

\textbf{The 2D prior as a reusable engine.} The image generators of §\ref{sec:p3_s8_1} recur throughout the chapter not as terminal systems but as priors to be distilled or lifted: 3D synthesis distills frozen 2D diffusion through score distillation and multi-view consistency (§\ref{sec:p3_s8_2}), video models inflate text-to-image backbones along the temporal axis (§\ref{sec:p3_s8_3}), and 4D pipelines compose both (§\ref{sec:p3_s8_5}). A single advance in the 2D substrate therefore propagates outward across every higher-dimensional subfield, which concentrates research leverage but also propagates the substrate's failure modes.

\textbf{Physical plausibility and fragmented evaluation.} Across subfields, the recurring failure mode is not low visual fidelity but violated physics and geometry, manifest as multi-view inconsistency in 3D, temporal drift and flicker in video, foot-skating and interpenetration in motion, and dynamic implausibility in 4D. Evaluation remains correspondingly fragmented: each subfield has evolved its own metric tuples, and all are migrating toward learned and vision-language-model judges as scalar similarity scores prove insufficient for spatially structured output. A unified, physically grounded evaluation protocol that spans these subfields remains an open goal.

\textbf{Outlook.} These observations point toward convergence rather than further fragmentation. World modeling (§\ref{sec:p3_s8_6}) already acts as a conceptual capstone that absorbs reconstruction, generation, and dynamics into a single predictive substrate, while programmatic and symbolic generation (§\ref{sec:p3_s8_7}) supplies the controllability and interpretability that purely neural pipelines lack; their integration, neural generators steered and verified by symbolic structure, is a natural trajectory for the field. The capabilities surveyed here also form the generative substrate for the applications of Part~\ref{part:applications}, where embodied agents, digital twins, and interactive environments consume generated spatial content rather than producing it for its own sake.

% \input{part_3/section_4/main}


